\documentclass[11pt,a4paper]{article}
\usepackage[margin=25mm,headheight=15pt]{geometry}
\usepackage[T1]{fontenc}
\usepackage[utf8]{inputenc}
\usepackage{lmodern}
\usepackage{amsmath,amssymb,amsthm,amscd,mathtools,mathrsfs}
\usepackage{microtype,booktabs,longtable,array,xcolor,graphicx,enumitem,needspace,fancyhdr}
\usepackage[round,authoryear]{natbib}
\usepackage[unicode,colorlinks=true,linkcolor=blue!40!black,citecolor=blue!40!black,urlcolor=blue!40!black]{hyperref}
\usepackage{xurl}
\usepackage{listings}
\usepackage{tabularx,etoolbox}
\makeatletter
\patchcmd{\l@section}{\addvspace{1.0em \@plus\p@}}{\addvspace{.35em \@plus\p@}}{}{}
\patchcmd{\l@part}{\addvspace{2.25em \@plus\p@}}{\addvspace{1em \@plus\p@}}{}{}
\makeatother
\lstset{basicstyle=\ttfamily\small,breaklines=true,columns=fullflexible,keepspaces=true}
\usepackage{tikz}
\usepackage{tikz-cd}
\usetikzlibrary{arrows.meta,calc,patterns,positioning}
\usepackage{aliascnt}
\usepackage[nameinlink,noabbrev]{cleveref}
\allowdisplaybreaks[2]
\numberwithin{equation}{section}
\setcounter{MaxMatrixCols}{32}
\setlength{\emergencystretch}{3em}
\setlength{\parskip}{2pt}
\newtheorem{theorem}{Theorem}[section]
\newaliascnt{proposition}{theorem}
\newtheorem{proposition}[proposition]{Proposition}
\aliascntresetthe{proposition}
\crefname{proposition}{proposition}{propositions}
\newaliascnt{lemma}{theorem}
\newtheorem{lemma}[lemma]{Lemma}
\aliascntresetthe{lemma}
\crefname{lemma}{lemma}{lemmas}
\newaliascnt{corollary}{theorem}
\newtheorem{corollary}[corollary]{Corollary}
\aliascntresetthe{corollary}
\crefname{corollary}{corollary}{corollarys}
\newaliascnt{claim}{theorem}
\newtheorem{claim}[claim]{Claim}
\aliascntresetthe{claim}
\crefname{claim}{claim}{claims}
\newaliascnt{statement}{theorem}
\newtheorem{statement}[statement]{Statement}
\aliascntresetthe{statement}
\crefname{statement}{statement}{statements}
\theoremstyle{definition}
\newaliascnt{definition}{theorem}
\newtheorem{definition}[definition]{Definition}
\aliascntresetthe{definition}
\crefname{definition}{definition}{definitions}
\newaliascnt{remark}{theorem}
\newtheorem{remark}[remark]{Remark}
\aliascntresetthe{remark}
\crefname{remark}{remark}{remarks}
\newaliascnt{example}{theorem}
\newtheorem{example}[example]{Example}
\aliascntresetthe{example}
\crefname{example}{example}{examples}
\newaliascnt{problem}{theorem}
\newtheorem{problem}[problem]{Open problem}
\aliascntresetthe{problem}
\crefname{problem}{open problem}{open problems}
\newaliascnt{question}{theorem}
\newtheorem{question}[question]{Question}
\aliascntresetthe{question}
\crefname{question}{question}{questions}
\newaliascnt{auditrecord}{theorem}
\newtheorem{auditrecord}[auditrecord]{Verification record}
\aliascntresetthe{auditrecord}
\crefname{auditrecord}{verification record}{verification records}
\newaliascnt{test}{theorem}
\newtheorem{test}[test]{Verification}
\aliascntresetthe{test}
\crefname{test}{verification}{verifications}
\newaliascnt{conjecture}{theorem}
\newtheorem{conjecture}[conjecture]{Conjecture}
\aliascntresetthe{conjecture}
\crefname{conjecture}{conjecture}{conjectures}
\newcommand{\C}{\mathbb C}\newcommand{\R}{\mathbb R}\newcommand{\Z}{\mathbb Z}\newcommand{\Q}{\mathbb Q}\newcommand{\N}{\mathbb N}\newcommand{\F}{\mathbb F}
\newcommand{\CC}{\mathbb C}\newcommand{\RR}{\mathbb R}\newcommand{\ZZ}{\mathbb Z}\newcommand{\QQ}{\mathbb Q}\newcommand{\NN}{\mathbb N}
\newcommand{\PP}{\mathbb P}\newcommand{\Pp}{\mathbb P}\newcommand{\Pone}{\mathbb P^1}\newcommand{\HH}{\mathbb H}
\newcommand{\RP}{\mathbb{RP}}\newcommand{\CP}{\mathbb{CP}}
\newcommand{\OO}{\mathcal O}\newcommand{\cO}{\mathcal O}\newcommand{\cM}{\mathcal M}\newcommand{\cF}{\mathcal F}
\newcommand{\h}{\mathfrak h}\newcommand{\hz}{\mathfrak h_z}\newcommand{\Acon}{\mathfrak A}\newcommand{\wt}[1]{\widehat{#1}}
\newcommand{\red}{\mathrm{red}}\newcommand{\tor}{\mathrm{tor}}\newcommand{\dP}{\mathrm{dP}}
\newcommand{\gh}{\widehat\gamma}\newcommand{\uh}{\widehat u}\newcommand{\wh}{\widehat w}\newcommand{\dhat}{\widehat\delta}\renewcommand{\dh}{\widehat\delta}
\newcommand{\dd}{\mathrm d}\newcommand{\ii}{\mathrm i}\newcommand{\eps}{\varepsilon}\newcommand{\partt}{\mathrm{part}}\newcommand{\Bcirc}{B^\circ}\newcommand{\ddbar}{\partial\bar\partial}
\DeclareMathOperator{\Pic}{Pic}\DeclareMathOperator{\Alb}{Alb}\DeclareMathOperator{\Aut}{Aut}\DeclareMathOperator{\NS}{NS}
\DeclareMathOperator{\Hom}{Hom}\DeclareMathOperator{\End}{End}\DeclareMathOperator{\Ext}{Ext}\DeclareMathOperator{\Tor}{Tor}
\DeclareMathOperator{\Spec}{Spec}\DeclareMathOperator{\Proj}{Proj}\DeclareMathOperator{\Sym}{Sym}\DeclareMathOperator{\Diff}{Diff}
\DeclareMathOperator{\GL}{GL}\DeclareMathOperator{\SL}{SL}\DeclareMathOperator{\PSL}{PSL}\DeclareMathOperator{\Sp}{Sp}
\DeclareMathOperator{\im}{im}\DeclareMathOperator{\coker}{coker}\DeclareMathOperator{\rank}{rank}\DeclareMathOperator{\diag}{diag}
\DeclareMathOperator{\ord}{ord}\DeclareMathOperator{\vol}{vol}\DeclareMathOperator{\pr}{pr}\DeclareMathOperator{\id}{id}
\DeclareMathOperator{\divisor}{div}\DeclareMathOperator{\Tr}{Tr}\DeclareMathOperator{\Nm}{Nm}\DeclareMathOperator{\Sha}{Sha}
\DeclareMathOperator{\Imn}{Im}\DeclareMathOperator{\Ren}{Re}\DeclareMathOperator{\trdeg}{tr.deg}
\newcommand{\Uses}[1]{\par\smallskip\noindent\textit{Uses:} #1.\par\smallskip}
\newcommand{\Verified}{}\newcommand{\AuthorSelfCheck}{}\newcommand{\NeedsCheck}{}\newcommand{\Open}{\textbf{Open.}\ }

\newcommand{\RungZero}[1]{\par\noindent\textbf{Result and mechanism.} #1\par\smallskip}
\newcommand{\RungOne}{\par\Needspace{6\baselineskip}\noindent\textbf{Proof at the first rung.}\par\nobreak\smallskip}
\newaliascnt{cited}{theorem}
\newtheorem{cited}[cited]{Cited result}
\aliascntresetthe{cited}
\crefname{cited}{cited result}{cited results}
\newaliascnt{refuted}{theorem}
\newtheorem{refuted}[refuted]{Correction}
\aliascntresetthe{refuted}
\crefname{refuted}{correction}{corrections}
\newaliascnt{computed}{theorem}
\newtheorem{computed}[computed]{Computation}
\aliascntresetthe{computed}
\crefname{computed}{computation}{computations}
\newaliascnt{openq}{theorem}
\newtheorem{openq}[openq]{Open question}
\aliascntresetthe{openq}
\crefname{openq}{open question}{open questions}

\providecommand{\rk}{\operatorname{rk}}
\providecommand{\Ann}{\operatorname{Ann}}
\providecommand{\Log}{\operatorname{Log}}
\providecommand{\vir}{\mathrm{vir}}

\usepackage{etoolbox}
\hypersetup{bookmarksnumbered=true,bookmarksdepth=2}
\setcounter{tocdepth}{2}
\makeatletter\renewcommand{\@pnumwidth}{2.5em}\renewcommand{\@tocrmarg}{3.5em}\renewcommand*{\l@subsection}{\@dottedtocline{2}{1.5em}{3.5em}}\makeatother
\newtheorem{maintheorem}{Theorem}
\renewcommand{\themaintheorem}{\Alph{maintheorem}}
\crefname{maintheorem}{theorem}{theorems}
\newcommand{\StatementRecord}[4]{{\upshape\footnotesize[tag~\texttt{#3}]}\enspace\ignorespaces}
\newcommand{\ExternalPaperRef}[3]{#1#2 (Paper~#3)}
\newcommand{\CompendiumRef}[2]{\ifmmode\text{\upshape \href{https://github.com/KokunoYumeto/s6-complex-structure-workbench/blob/\SsixRepositoryCommit/pro_returns/S6_09_PAPERS_final_20261117T0601Z/S6_09_PAPERS_#1.pdf}{S6\_09, Paper~#1}, label~\nolinkurl{#2}}\else \href{https://github.com/KokunoYumeto/s6-complex-structure-workbench/blob/\SsixRepositoryCommit/pro_returns/S6_09_PAPERS_final_20261117T0601Z/S6_09_PAPERS_#1.pdf}{S6\_09, Paper~#1}, label~\nolinkurl{#2}\fi}
\newcommand{\SeriesRefText}[1]{\ifcsname SRef@#1\endcsname\csname SRef@#1\endcsname\else\textbf{UNRESOLVED}\fi}
\newcommand{\SeriesRef}[1]{\ifmmode\text{\SeriesRefText{#1}}\else\SeriesRefText{#1}\fi}


\newcommand{\CurrentPaper}{V}
\expandafter\def\csname SRef@mapsI:src:A.1.1\endcsname{\href{https://github.com/KokunoYumeto/s6-complex-structure-workbench/blob/\SsixRepositoryCommit/pro_returns/S6_SERIES_GRIND_checkpoint_10_20261120T0211Z/S6_SERIES_GRIND_I.pdf\#nameddest=tag-02PA}{I.22.1 (p.~295)} }
\expandafter\def\csname SRef@mapsI:src:A.1.3\endcsname{\href{https://github.com/KokunoYumeto/s6-complex-structure-workbench/blob/\SsixRepositoryCommit/pro_returns/S6_SERIES_GRIND_checkpoint_10_20261120T0211Z/S6_SERIES_GRIND_I.pdf\#nameddest=tag-02PC}{I.22.3 (p.~297)} }
\expandafter\def\csname SRef@def:common-conventions\endcsname{\href{https://github.com/KokunoYumeto/s6-complex-structure-workbench/blob/\SsixRepositoryCommit/pro_returns/S6_SERIES_GRIND_checkpoint_10_20261120T0211Z/S6_SERIES_GRIND_I.pdf\#nameddest=tag-01O1}{I.2 (p.~21)} }
\expandafter\def\csname SRef@mapsI:src:A.2.1\endcsname{\href{https://github.com/KokunoYumeto/s6-complex-structure-workbench/blob/\SsixRepositoryCommit/pro_returns/S6_SERIES_GRIND_checkpoint_10_20261120T0211Z/S6_SERIES_GRIND_I.pdf\#nameddest=tag-02PD}{I.22.4 (p.~298)} }
\expandafter\def\csname SRef@maps:src:B.5.2\endcsname{\href{https://github.com/KokunoYumeto/s6-complex-structure-workbench/blob/\SsixRepositoryCommit/pro_returns/S6_SERIES_GRIND_checkpoint_10_20261120T0211Z/S6_SERIES_GRIND_III.pdf\#nameddest=tag-02JC}{III.3.10 (p.~23)} }
\expandafter\def\csname SRef@maps:src:B.1.2\endcsname{\href{https://github.com/KokunoYumeto/s6-complex-structure-workbench/blob/\SsixRepositoryCommit/pro_returns/S6_SERIES_GRIND_checkpoint_10_20261120T0211Z/S6_SERIES_GRIND_III.pdf\#nameddest=tag-02J5}{III.3.3 (p.~14)} }
\expandafter\def\csname SRef@maps:src:B.2.1\endcsname{\href{https://github.com/KokunoYumeto/s6-complex-structure-workbench/blob/\SsixRepositoryCommit/pro_returns/S6_SERIES_GRIND_checkpoint_10_20261120T0211Z/S6_SERIES_GRIND_III.pdf\#nameddest=tag-02J6}{III.3.4 (p.~15)} }
\expandafter\def\csname SRef@maps:src:B.5.1\endcsname{\href{https://github.com/KokunoYumeto/s6-complex-structure-workbench/blob/\SsixRepositoryCommit/pro_returns/S6_SERIES_GRIND_checkpoint_10_20261120T0211Z/S6_SERIES_GRIND_III.pdf\#nameddest=tag-02JB}{III.3.9 (p.~20)} }
\expandafter\def\csname SRef@maps:src:D.1.2\endcsname{\href{https://github.com/KokunoYumeto/s6-complex-structure-workbench/blob/\SsixRepositoryCommit/pro_returns/S6_SERIES_GRIND_checkpoint_10_20261120T0211Z/S6_SERIES_GRIND_III.pdf\#nameddest=tag-02KB}{III.5.2 (p.~32)} }
\expandafter\def\csname SRef@maps:src:D.1.4\endcsname{\href{https://github.com/KokunoYumeto/s6-complex-structure-workbench/blob/\SsixRepositoryCommit/pro_returns/S6_SERIES_GRIND_checkpoint_10_20261120T0211Z/S6_SERIES_GRIND_III.pdf\#nameddest=tag-02KD}{III.5.4 (p.~34)} }
\expandafter\def\csname SRef@maps:src:D.1.9\endcsname{\href{https://github.com/KokunoYumeto/s6-complex-structure-workbench/blob/\SsixRepositoryCommit/pro_returns/S6_SERIES_GRIND_checkpoint_10_20261120T0211Z/S6_SERIES_GRIND_III.pdf\#nameddest=tag-02KI}{III.5.9 (p.~39)} }
\expandafter\def\csname SRef@maps:src:E.2.2\endcsname{\href{https://github.com/KokunoYumeto/s6-complex-structure-workbench/blob/\SsixRepositoryCommit/pro_returns/S6_SERIES_GRIND_checkpoint_10_20261120T0211Z/S6_SERIES_GRIND_III.pdf\#nameddest=tag-02MK}{III.9.3 (p.~317)} }
\expandafter\def\csname SRef@maps:src:B.6.24\endcsname{\href{https://github.com/KokunoYumeto/s6-complex-structure-workbench/blob/\SsixRepositoryCommit/pro_returns/S6_SERIES_GRIND_checkpoint_10_20261120T0211Z/S6_SERIES_GRIND_IV.pdf\#nameddest=tag-02JT}{IV.2.54 (p.~102)} }

\newcommand{\SsixRepositoryCommit}{83d76022bcd442278b760db9ff438afceed6c6ea}
\begin{document}
\title{The $(3,4,\infty)$ family of complex six-spheres\\[.5em]V. Complex fourfolds from the family}
\author{AI-produced mathematical text}
\date{21 November 2026}
\maketitle
\begin{abstract}
We construct compact complex fourfolds diffeomorphic, with their complex
orientation, to $3\mathbb{HP}^2\#(-\mathbb{HP}^2)$.
They form a family of complex three-tori over $\mathbb P^1$ with
monodromy in $GL_6(\mathbb Z)$: two multiple fibres, of multiplicities
five and four, and one Kuhn cusp. Its periods come from the exterior
square of the genus-two pencil $y^2=x^5+\lambda^3x+\lambda^4$.
The family is parametrized exactly
by the half-plane $\operatorname{Im}\xi>c_*$, where
$c_*\approx1.09454$ is minus the minimum of an explicit real-analytic
period function.  We compute the Hodge numbers, the tangent
cohomology (the Kuranishi germ is smooth of dimension one) and the
automorphism group $\mathbb C\times\mathbb Z/2$, and we show that the
family extends analytically across no boundary point of its parameter
domain.  We also construct a compact complex fourfold outside Fujiki's
class $\mathcal C$ as a small crepant resolution of a Weierstrass model,
bimeromorphic to the fibre product of a complex six-sphere with an
elliptic surface; give an effective
classification of the integral monodromy data of this kind that yield
three-connected eight-manifolds, all of which are diffeomorphic to
$3\mathbb{HP}^2\#(-\mathbb{HP}^2)$; and describe shifted fibre products
of complex dimension five.
\end{abstract}
\clearpage
\tableofcontents
\clearpage
\section{Results and their proofs in outline}\label{sec:results}
\subsection{The series and this paper}
This paper is the fifth of a series on the holomorphic families built
from the $(3,4,\infty)$ family of torus-fibred complex structures on
$S^6$ \citep{AC26,E26}.  The table gives the main result of each paper
of the series and the papers it uses.
\begin{center}\small
\begin{tabularx}{\textwidth}{@{}lXl@{}}\toprule
Paper & Main result & Uses\\\midrule
I & Exact periods, moduli, fibre strata and the two ends & ---\\
II & Torus-fibred threefolds with $b_1=b_2=0$: the Campana--Demailly--Peternell argument and the sign theorem for all degenerations & I, IV\\
III & Maps from the six-spheres and the manifolds they give & I, II, IV, V\\
IV & Smooth monodromy in $\Theta_7$ & I, II, III, V\\
V & Complex fourfolds on $3\mathbb{HP}^2\#(-\mathbb{HP}^2)$ and a resolved fibre product & I, III, IV\\
\bottomrule\end{tabularx}\end{center}
\subsection{The main theorems}
\hypertarget{tag-03KQ}{}\label{series:V:main-theorems}

\begin{definition}[The period parameter]
\hypertarget{tag-PV14}{}\label{pub:V:parameter}
\StatementRecord{pub:V:parameter}{proved}{PV14}{}
We use the algebraic splitting $\sigma_b$, invariant vector $e$ and
functional $\alpha$ with the exact normalization of
\cref{def:hb,night:P1.15}. For the rank-three bundle $P$ of
\cref{night:P1.5}, put
\[
 F_\xi=(\sigma_b+\xi\alpha e)(P),\qquad
 c_*=-\min h_b,\qquad
 \mathcal U=\{\xi\in\mathbb C:\operatorname{Im}\xi>c_*\},
\]
where $h_b$ is the real-intersection function of \cref{def:hb}.
We write $M_\xi$ for the compact filling of these periods with the
base and fibre conjugation of \cref{rem:V:conjugation-parameter}.
For the imaginary ray we write $M_c=M_{ic}$, with $c>c_*$.
The positive Kuhn fan is fixed throughout this family.
\end{definition}

\hypertarget{main-V-D}{}
\begin{maintheorem}[The oriented diffeomorphism type]
\hypertarget{tag-03KU}{}\label{series:V:D}
\StatementRecord{series:V:D}{proved}{03KU}{\cref{maps:src:F.3.7}, \cref{maps:src:F.4.1}}
For every $\xi\in\mathcal U$, the genus-two periods of
\cref{series:V:C}, with the integral datum of
\cref{pub:V:rank-six-data}, extend over the finite root charts
and the positive marked Kuhn cusp. They give a compact complex
fourfold $M_\xi$ and an orientation-preserving diffeomorphism
\[
 M_\xi\ \cong\ 3\mathbb{HP}^2\#(-\mathbb{HP}^2).
\]
The fourfolds form a holomorphic family over $\mathcal U$.
Their smooth middle lattice and spin characteristic vector are
isometric to $\bigl(I_{3,1},(1,1,1,1)\bigr)$.
\end{maintheorem}

\textit{Mechanism.} The middle-lattice isometry transports both Euler and Pontryagin clutching data, after which the cited absorption theorem removes the homotopy sphere.
\textit{Relation to literature.} The compact complex realization is new; the smooth classification and absorption results used to identify its underlying manifold are listed in \S\ref{pub:V:known-new}.
\hyperlink{outline-V-D}{Outline of the proof}.

\hypertarget{main-V-E}{}
\begin{maintheorem}[The exact family and its two ends]
\hypertarget{tag-03KV}{}\label{series:V:E}
\StatementRecord{series:V:E}{proved}{03KV}{\cref{night:P1.18}, \cref{night:P1A.1}, \cref{night:P1A.2}, \cref{night:P1A.3}, \cref{night:P1A.4}, \cref{night:P1A.5}, \cref{night:P1.17}, \cref{night:P1A.6}}
Put
\[
 c_*=-h_{b,\min}=\frac{96}{625\pi}\operatorname{Re}(XH_\infty)_{03},
 \qquad \mathcal U=\{\xi\in\C:\Im\xi>c_*\},
\]
where $X$ is the normalized integral intertwiner of
\cref{night:P1.15} and $H_\infty$ is its regular infinity-root
fundamental-solution limit, normalized at $\lambda=1$ as in
\cref{night:P1.18}. Thus
\[c_*\approx1.09454.\]

The marked family is parametrized exactly by $\mathcal U$:
its members are pairwise nonisomorphic, have algebraic dimension one
and automorphism group $\mathbb C\times\mathbb Z/2$, and admit no
antiholomorphic isomorphism for this fan.  The real-intersection
function determines the parameter domain; the invariant line and
one half-period determine the automorphisms, while the reflected cusp map fails to preserve a ray of the chosen fan.
On the effective root curve put
$L=\OO((4/5)[\infty]-(3/4)[0])$. Then
\[
 R^qh_*\OO=\bigwedge^q(\OO\oplus L^{-1}\oplus L^{-2}),
 \qquad K_M=h^*L^{-6}.
\]
With $p$ indexing rows and $q$ indexing columns, the Hodge table is
\[
\begin{array}{c|ccccc}
h^{p,q}&0&1&2&3&4\\\hline
0&1&1&0&0&0\\
1&0&2&1&0&0\\
2&0&1&6&1&0\\
3&0&0&1&2&0\\
4&0&0&0&1&1
\end{array}
\]
The tangent tuple is $h^q(T_M)=(1,1,1,2,0)$ and its Kuranishi germ
is smooth of dimension one. A full parameter disc across the
transversality boundary has the proper-filling obstruction of
\cref{night:P1A.5}. A proper smooth family on a punctured disc
whose fibres belong to $\mathcal U$ has the smooth filling after
the explicit degree-two torsor cover of \cref{night:P1A.6}.


\end{maintheorem}

\textit{Mechanism.} The normalized Frobenius period minimum fixes the exact domain, and the marked period and cap maps determine the family invariants and isomorphisms.
\textit{Relation to literature.} The exact parameter domain and the family invariants are new; the cited direct-image and torus-deformation inputs are listed in \S\ref{pub:V:known-new}.
\hyperlink{outline-V-E}{Outline of the proof}.

\hypertarget{main-V-C}{}
\begin{maintheorem}[The genus-two period construction]
\hypertarget{tag-03KT}{}\label{series:V:C}
\StatementRecord{series:V:C}{proved}{03KT}{\cref{night:P1.4}, \cref{night:P1.5}, \cref{maps:src:F.3.5}, \cref{maps:src:F.3.6}, \cref{night:P1.18}}
The pencil $y^2=x^5+\lambda^3x+\lambda^4$ has one nodal
fibre away from $0,\infty$, at $\lambda_*=-256/3125$.
Its good root charts have orders eight and five, and on that root
curve its Hodge bundle is
\[
 \mathcal H=T^3\oplus T,\qquad
 T=\OO((2/5)[\infty]-(3/8)[0]),\qquad \deg T=1/40.
\]
The meromorphic differential $dx/y$ generates $\ell_0=T^3$.
In the full rank-four de Rham bundle, projection of
$\ell_0\wedge\mathcal H_{\rm dR}$ to
$\bigwedge^2\mathcal H_{\rm dR}/\mathbb C\omega$
is a rank-three bundle $P$.  Its holomorphic lift to the rank-six extension is $(\sigma_b+\xi\alpha e)(P)$, with the
normalization of $e,\alpha,\sigma_b$ in \cref{night:P1.15}.
It is transverse to its conjugate on the completed base exactly
for the parameters $\xi\in\mathcal U$ of Theorem~\ref{series:V:E}.
The exterior-square and lift maps descend the finite root
orders to five and four; the cusp has period
$Z(q)=(\log q/(2\pi i))I_3+C(q)$ with $C$ holomorphic at zero
(\cref{night:P1.4,night:P1.5,maps:src:F.3.5,maps:src:F.3.6,night:P1.18}).

\end{maintheorem}

\textit{Mechanism.} The wedge image of the genus-two Hodge line has rank three, and its split lift becomes transverse after the normalized imaginary shift.
\textit{Relation to literature.} The pencil, its integral exterior-square comparison and its normalized lifted planes are new.
\hyperlink{outline-V-C}{Outline of the proof}.

\hypertarget{main-V-B}{}
\begin{maintheorem}[The integral rank-six construction]
\hypertarget{tag-03KS}{}\label{series:V:B}
\StatementRecord{series:V:B}{proved}{03KS}{\cref{maps:src:F.2.5}, \cref{maps:src:F.2.6}, \cref{night:P1.13}}
The rank-six integral monodromy of \cref{pub:V:rank-six-data}, with free affine fillings of orders five and four and the index-one Kuhn cusp, gives a closed smooth eight-manifold that is three-connected, has free integral homology, and has middle intersection form $3\langle1\rangle\oplus\langle-1\rangle$.

\end{maintheorem}

\textit{Mechanism.} In the integral Leray spectral sequence over $\mathbb P^1$, each order-two suspension class from the order-four filling transgresses to the corresponding nonzero global coinvariant, so the resulting homology is free; Todd integrality and the signature bound then force the middle form.
\textit{Relation to literature.} The integral rank-six datum and its global homology calculation are new; the toroidal filling is the cited construction listed in \S\ref{pub:V:known-new}.
\hyperlink{outline-V-B}{Outline of the proof}.

\hypertarget{main-V-F}{}
\begin{maintheorem}[The integral input classification]
\hypertarget{tag-03KW}{}\label{series:V:F}
\StatementRecord{series:V:F}{proved}{03KW}{\cref{maps:src:F.2.39}, \cref{maps:src:F.3.11}, \cref{maps:src:F.3.4}}
Let $k\in\{2,3\}$ and $\Lambda=\mathbb Z^6$.
Consider ordered integral tuples $(M_1,\ldots,M_k,M_0)$ with each
$M_i$ of finite order and having a nonzero fixed vector, with
$M_0=(M_1\cdots M_k)^{-1}$, and with
$N=M_0-I$ satisfying $N^2=0$, rank three, and
$\operatorname{im}N=\ker N$ as integral sublattices of $\Lambda$.
Every such tuple has a computable canonical representative and an
explicit simultaneous unimodular conjugator; equality of the
representatives is equivalent to simultaneous integral conjugacy.
For input entry height $H$, the finite stopping bound is
$(2H+1)^{36(k+1)}$ candidates before the proved rejection filters.
The ordered sector with both finite linear orders two has exactly
56 classes.  Each has free affine lifts of affine order four,
and every free affine filling of this sector has nontrivial
fundamental group.

Specify the free affine actions
$\tau_i(x)=M_ix+v_i/m_i$, with $v_i\in\Lambda$ and affine order $m_i$, an integral translation clutch $\kappa\in\Lambda$,
and a primitive Kuhn cusp marking.  The algorithms compute the
full fundamental group and every integral homology group of the
resulting marked gluing for every such input.  Their constructive
parametrization and finite tests identify exactly the inputs with
$\pi_1=1$ and $H_2=H_3=0$; there is no fixed bound on the
orders, twists or lattice indices (\cref{maps:src:F.2.39}).
Every resulting three-connected manifold is orientation-preservingly
diffeomorphic to $3\mathbb{HP}^2\#(-\mathbb{HP}^2)$, by
\cref{maps:src:F.2.30}.

On $B^*=\mathbb P^1\setminus\{p_0,\ldots,p_k\}$ with $k\ge1$,
every representation $\rho:\pi_1(B^*)\to GL^+(6,\mathbb R)$
whose $p_0$ monodromy satisfies $(M_0-I)^2=0$ and
$\operatorname{rank}(M_0-I)=3$ has a holomorphic period family
transverse to its conjugate everywhere on $B^*$. In the marking
$M_0=\left(\begin{smallmatrix}I&I\\0&I\end{smallmatrix}\right)$,
its cusp period is
\[
 Z(t)=\frac{\log t}{2\pi i}I_3+C(t),\qquad C\text{ holomorphic at }0,
\]
with any prescribed finite jet of $C$ at zero. An integral primitive
cusp admits an integral marking. The other $p_i$ remain punctures
in this assertion (\cref{maps:src:F.3.11}).

For a polarized algebraic abelian-surface family with toric-rank-two
cusp, the construction that takes a fibrewise degree-zero line bundle
and quotients its punctured bundle by a lifted translation cannot
have $b_1=b_2=0$ with the stipulated rank-three square-zero Kuhn
cusp; this holds for both proportional and nonproportional source
sections. The filtration and Euler-parity proof is
\cref{maps:src:F.3.4}.

\end{maintheorem}

\textit{Mechanism.} Exact intertwiner lattices decide simultaneous conjugacy, finite integral complexes compute each marked gluing, and the cusp-period construction realizes the punctured analytic data.
\textit{Relation to literature.} The classification and bounded-cusp period theorem are new; their cited lattice-isomorphism and Oka inputs are listed in \S\ref{pub:V:known-new}.
\hyperlink{outline-V-F}{Outline of the proof}.


\hypertarget{main-V-A}{}
\begin{maintheorem}[The resolved transverse elliptic fibre product]
\hypertarget{tag-03KR}{}\label{series:V:A}
\StatementRecord{series:V:A}{proved}{03KR}{\cref{maps:src:E.1.2}, \cref{maps:src:E.2.1}}
Let $X$ be a canonical marked member of the $(3,4,\infty)$
family, with reduced multiple fibres $S_1,S_2$ of multiplicities
three and four, and put $L=\mathcal O_X(S_1-S_2)$.
Let $S\to\mathbb P^1$ be the elliptic surface in
\cref{maps:src:E.1.2}, with its cusp over the cusp of $X$.
The Weierstrass model $\mathscr W\to X$ of that theorem has a
small crepant resolution $r:Z\to\mathscr W$.
Its composition $\Pi:Z\to X$ is proper, and $Z$ is
bimeromorphic to $X\times_{\mathbb P^1}S$.
It is a compact complex fourfold with
$K_Z=\Pi^*L^{-5}$, algebraic dimension two,
$b_1(Z)=0$, and $h^{0,1}(Z)=1$.
Its signature is two and its middle intersection form is
$\langle1\rangle\oplus\langle1\rangle$.
It lies outside Fujiki's class $\mathcal C$.
\end{maintheorem}

\textit{Mechanism.} The explicit period comparison and a small crepant Kuhn resolution preserve the root canonical calculation and the integral characteristic pairings.
\textit{Relation to literature.} The global small crepant resolution and its computed invariants are new; the cited six-sphere, elliptic and toroidal constructions are the inputs listed in \S\ref{pub:V:known-new}.
\hyperlink{outline-V-A}{Outline of the proof}.

\subsection{The proofs in outline}
\hypertarget{outline-V-D}{}
\subsubsection*{Theorem~\ref{series:V:D}: The oriented diffeomorphism type}
\begin{enumerate}[label=\arabic*.,leftmargin=*]
\item\hypertarget{step-V-D-1}{} The lifted periods commute with the finite root actions and
  have the prescribed bounded cusp correction, so the same quotient
  and toroidal charts are holomorphic and proper.  Their collar
  identifications are the marked maps used in the smooth construction
  (\cref{maps:src:F.3.7}, \cref{night:P1.9}).

\item\hypertarget{step-V-D-2}{} The preceding integral calculation gives the intersection
  lattice $I_{3,1}$.  The spin characteristic vector $p_1/2$ has
  square two, by the Chern-number calculation
  (\cref{night:P1.13}).

\item\hypertarget{step-V-D-3}{} The characteristic vector has odd coordinates and norm two in $I_{3,1}$. The integral reflection in $(1,1,1,1)$ strictly decreases the absolute negative coordinate until all four coordinates have absolute value one. The resulting explicit isometry is proved in \cref{night:P1.13}.

\item\hypertarget{step-V-D-4}{} In the quaternionic clutching coordinates, the map is
  $(a,b)\mapsto(e,p_1/2)=(a-b,a+b)$, with inverse
  $((e+p_1/2)/2,(p_1/2-e)/2)$.  Thus the pairing and characteristic
  vector determine the normal datum; the Poincar\'e-duality map
  transports it with the middle lattice isometry
  (\cref{maps:src:F.4.1}).

\item\hypertarget{step-V-D-5}{} Each quaternionic projective-plane summand has $p_1=2u$.
  The oriented connected-sum pairing and characteristic vector are
  therefore the same pair $\bigl(I_{3,1},(1,1,1,1)\bigr)$.
  Wall's stated extended-form classification gives the oriented
  diffeomorphism up to a homotopy eight-sphere
  (\cref{maps:src:F.4.1}).

\item\hypertarget{step-V-D-6}{} Kramer--Stolz's Theorem~A absorbs that sphere in one positive
  $\mathbb{HP}^2$ summand.  Its diffeomorphism has degree $+1$ because
  the square of a middle generator evaluates to $1$ on both sides.
  Oriented connected-sum maps therefore give the required oriented
  diffeomorphism of the full fourfold
  (\cref{maps:src:F.4.1}).

\end{enumerate}

\hypertarget{outline-V-E}{}
\subsubsection*{Theorem~\ref{series:V:E}: The exact family and its two ends}
\begin{enumerate}[label=\arabic*.,leftmargin=*]
\item\hypertarget{step-V-E-1}{} The infinity-root symmetry makes the differential of the normalized height vanish at zero. Its positive Hessian proves a strict local minimum. Convergent root and cusp expansions, with interval continuation over finite closed covers, prove the height is strictly larger on their complements; the positive cusp graph treats the remaining inner disc. This proves the unique global minimum and the exact domain in \cref{night:P1.18}.

\item\hypertarget{step-V-E-2}{} On the effective root curve of orders five and four put
  $L=\mathcal O((4/5)[\infty]-(3/4)[0])$.  The direct images are
  $R^qh_*\mathcal O=\bigwedge^q
  (\mathcal O\oplus L^{-1}\oplus L^{-2})$, and
  $K_M=h^*L^{-6}$.  The toric normalization complex proves the
  exterior-algebra assertion at the cusp, and the root-chart
  invariant cohomology proves it at the finite fibres
  (\cref{night:P1A.1}).

\item\hypertarget{step-V-E-3}{} The logarithmic cotangent sequence and its explicit residue
  matrices give the Hodge table.  Serre duality applied to
  the twisted logarithmic and residue sequences gives
  $h^q(T_M)=(1,1,1,2,0)$.  The normal tangent map and the period lift give a nonzero Kodaira--Spencer class; the local
  analytic-ring comparison makes the Kuranishi germ smooth of
  dimension one (\cref{night:P1A.2}, \cref{night:P1A.3}).

\item\hypertarget{step-V-E-4}{} Restriction of a horizontal divisor to a regular torus fibre
  gives an invariant alternating class.  The computed invariant
  space is zero, while positive pairing with a torus Kähler form
  rules out a nonzero effective divisor of zero class.  Thus the
  meromorphic field is the base field $\mathbb C(\lambda)$.
  Every isomorphism therefore preserves the base and its three
  distinguished values.  The simultaneous linear centralizer is
  $\mathbb Q I$; the order-five affine norm excludes its negative
  integral unit.  The translation sheaf has identity component
  $\mathbb C$ and an integral parabolic quotient $\mathbb Z/2$,
  represented by the displayed half-period, proving the
  automorphism group and equality of parameters for isomorphic
  members (\cref{night:P1A.4}).

\item\hypertarget{step-V-E-5}{} The two integral anti-intertwiners are $\pm R$.  One is excluded
  by the finite cap norm.  For the other, the based cusp
  correction acts on rays by $-P$, and takes $e_2$ to $(0,-1,1)$,
  outside the component's rays $\pm(\{0,1\}^3\setminus\{0\})$.
  The corresponding monomial map contracts a divisor, contradicting
  a biholomorphism (\cref{night:P1.17}).

\item\hypertarget{step-V-E-6}{} A proposed proper Hausdorff filling across a boundary point of
  a full parameter disc recovers the marked base from the moving
  pencil of $K_M^{-4}$, whose two sections have ratio $\lambda$.
  Closed relative forms give torus fibres outside a proper
  analytic subset.  Their marked periods continue as $F_\xi$,
  while the real hypersurface $h_b+\operatorname{Im}\xi=0$
  cannot be contained in that analytic subset.  A torus fibre
  there violates real transversality (\cref{night:P1A.5}).  For a family on a punctured disc whose
  fibres belong to $\mathcal U$, the unique classifying map is
  holomorphic; its Cayley transform extends and the maximum principle
  keeps its limit inside $\mathcal U$.  After the explicit degree-two
  removal of the finite torsor, that family has a smooth filling
  (\cref{night:P1A.6}).

\end{enumerate}

\hypertarget{outline-V-C}{}
\subsubsection*{Theorem~\ref{series:V:C}: The genus-two period construction}
\begin{enumerate}[label=\arabic*.,leftmargin=*]
\item\hypertarget{step-V-C-1}{} The pencil has one node at $\lambda_*=-256/3125$ and good
  reductions after root orders eight and five.  Its Hodge bundle on
  that root curve is $T^3\oplus T$, with
  $T=\mathcal O((2/5)[\infty]-(3/8)[0])$ and
  $\deg T=1/40$ (\cref{night:P1.4}).

\item\hypertarget{step-V-C-2}{} In the exterior-square lattice, projection induces
  $j:\omega^\perp\to(\bigwedge^2\mathbb Z^4)/\mathbb Z\omega$
  with Smith factors $(1,1,1,1,2)$.  This exact index-two map identifies
  the rational systems, and the quotient is the dual of the primitive
  integral lattice.  The lower five-dimensional blocks of the given
  rank-six tuple are the quotient blocks under this map
  (\cref{night:P1.2}).

\item\hypertarget{step-V-C-3}{} Put $P=p(\ell_0\wedge\mathcal H_{\mathrm{dR}})$, where
  $\ell_0=\langle dx/y\rangle$ and $p$ is projection modulo $\omega$.
  The wedge map has rank three and its image meets $\mathbb C\omega$
  trivially: a decomposable two-vector has square zero, whereas
  $\omega^2\ne0$.  This proves the required rank before a lift is
  chosen (\cref{night:P1.5}, \cref{maps:src:F.3.5}).

\item\hypertarget{step-V-C-4}{} The root-line decomposition of $P$ has $H^1(P^*)=0$ and a one-dimensional space of functionals by \cref{night:P1.5}. Its extension therefore splits holomorphically. The exact contraction and period comparison in \cref{def:hb} identify the normalization of $\alpha$ and $\sigma_b$.

\item\hypertarget{step-V-C-5}{} For the real-intersection line, the normalized lift has transversality exactly when $h_b+\operatorname{Im}\xi\ne0$. The cusp has $Z=(\log q)/(2\pi i)I_3+C(q)$ and $h_b\to+\infty$. Thus the complete parameter domain is the half-plane above $-\min h_b$, by \cref{night:P1.8,night:P1.18}.

\end{enumerate}

\hypertarget{outline-V-B}{}
\subsubsection*{Theorem~\ref{series:V:B}: The integral rank-six construction}
\begin{enumerate}[label=\arabic*.,leftmargin=*]
\item\hypertarget{step-V-B-1}{} The matrices of \cref{pub:V:rank-six-data} have orders five and four. The square-zero cusp has primitive rank-three image, and its explicit unimodular marking is proved in \cref{night:P1.1}. The invariant covector $x\mapsto x_0$ detects every nonidentity affine power, proving the freeness used in \cref{pub:V:smooth-gluing}.

\item\hypertarget{step-V-B-2}{} The cusp and finite quotient charts give the marked smooth gluing. Its van Kampen relations reduce to $5a=h$, $4b=-h$ and $a+b=0$; their integral relation matrix is unimodular. Hence the manifold is simply connected, by \cref{pub:V:smooth-gluing}.

\item\hypertarget{step-V-B-3}{} The signed-orbit bases at the finite cores, their norm maps, and the primitive cusp collapse give the integral restriction sheaves and their point kernels. Determinant minors compute their cohomology in \cref{pub:V:integral-restrictions}, using the collapse of \cref{night:P1.11}.

\item\hypertarget{step-V-B-4}{} For the order-four core, the loops $v=g^4h^{-1}$ and $t=g$ form an integral basis. A torsion representative pulls back to $(0,w)$, so its transgression is $\pm[w]$ by \cref{mapsaudit:order-four-torsion-transgression}. The three nonzero global images are proved in \cref{pub:V:integral-restrictions}.

\item\hypertarget{step-V-B-5}{} The three torsion differentials are isomorphisms onto $\mathbb Z/2$. The remaining cusp kernel is free of rank one. The Leray filtration, universal coefficients, duality and Hurewicz then give free integral homology and three-connectivity in \cref{maps:src:F.2.6}.

\item\hypertarget{step-V-B-6}{} The almost complex structure has $c_1=c_3=0$. The index and signature identities of \cref{sec:V-known-inputs} force $c_2^2=2$, $\sigma=2$ and the middle form $I_{3,1}$. The Pontryagin identities and characteristic vector are computed in \cref{night:P1.13}.

\end{enumerate}

\hypertarget{outline-V-F}{}
\subsubsection*{Theorem~\ref{series:V:F}: The integral input classification}
\begin{enumerate}[label=\arabic*.,leftmargin=*]
\item\hypertarget{step-V-F-1}{} Primitive square-zero cusp matrices have the explicit integral
  normal form and centralizer
  $\left(\begin{smallmatrix}P&Q\\0&P\end{smallmatrix}\right)$.
  Eigenvalue one is necessary for a free affine lift, leaving the
  finite cyclotomic list and its complex-compatibility filter
  (\cref{maps:src:F.2.1}, \cref{maps:src:F.2.2}).

\item\hypertarget{step-V-F-2}{} The rational intertwiner algebra, its radical and its
  semisimple quotient reduce the simultaneous integral-conjugacy
  problem to the cited lattice-isomorphism algorithm over the resulting
  product of matrix algebras over number fields.  Searching by height and lexicographic order terminates
  by the input height, yielding a canonical representative and an
  explicit unimodular conjugator (\cref{maps:src:F.2.36}, \cref{maps:src:F.2.37}).

\item\hypertarget{step-V-F-3}{} For that conjugator $R$, the coordinate map $x\mapsto Rx$
  transports $(M_i,v_i,\kappa,P)$ to
  $(RM_iR^{-1},Rv_i,R\kappa,RP)$.  Since
  $(RP)^{-1}(Rx)=P^{-1}x$, it intertwines the cusp charts
  as well as the affine caps and clutch.  The resulting finite
  cellular cone and group presentation compute all integral homology
  and the full fundamental group (\cref{maps:src:F.2.31}, \cref{maps:src:F.2.35}, \cref{maps:src:F.2.39}).

\item\hypertarget{step-V-F-4}{} The exact two-point tests and the three-point
  orbifold quotient give the complete constructive parametrization;
  the separate involution branch has 56 classes.  The displayed
  unbounded trace family shows why the terminating bound depends on
  the input (\cref{maps:src:F.2.30}, \cref{maps:src:F.2.38}, \cref{maps:src:F.2.39}).

\item\hypertarget{step-V-F-5}{} Nine explicitly given complex-line directions and a uniform
  singular-value bound put a spanning tube in the unpolarized period
  domain.  The local-spray approximation theorem, with the
  relative extension-point argument proved here, produces a
  holomorphic transverse period section and retains any prescribed
  bounded cusp jet (\cref{maps:src:F.3.8}, \cref{maps:src:F.3.10}, \cref{maps:src:F.3.11}).

\item\hypertarget{step-V-F-6}{} For a polarized abelian-surface source, semisimplicity and the
  maximal square-zero cusp force the middle rank-four system to have
  no invariant or coinvariant vectors.  The filtration
  $0\subset\mathbb Q\subset E\subset T$, with middle quotient that
  system and top quotient $\mathbb Q$, then forces both relevant
  invariant ranks to be one; the integral Euler-parity relation
  excludes $b_1=b_2=0$ (\cref{maps:src:F.3.4}).

\end{enumerate}

\hypertarget{outline-V-A}{}
\subsubsection*{Theorem~\ref{series:V:A}: The resolved transverse elliptic fibre product}
\begin{enumerate}[label=\arabic*.,leftmargin=*]
\item\hypertarget{step-V-A-1}{} Pull back the Weierstrass family to the root curve of the
threefold fibration.  The integral matrices $A,B$ give the period comparison $AZ_3B^{-1}=\operatorname{diag}(Z_2,s+c_E)$.
See \cref{maps:src:E.1.2}.
\item\hypertarget{step-V-A-2}{} The unimodular Kuhn fan resolves the local fibre product
without adding rays.  Its period action is proper and its marked
section extends face by face.  The resolution is therefore small
and crepant.
See \cref{maps:src:E.1.2}.
\item\hypertarget{step-V-A-3}{} The canonical formula on the root curve gives
$K=\Pi^*L^{-5}$.  The split relative elliptic complex and Leray
give $(h^{0,q})=(1,1,0,0,0)$; the explicit meromorphic product
comparison gives algebraic dimension two.
See \cref{maps:src:E.2.1}.
\item\hypertarget{step-V-A-4}{} A K\"ahler modification in class $\mathcal C$ preserves
$b_1$ and $h^{0,1}$.  Here $b_1=0$ and $h^{0,1}=1$, contradicting
the K\"ahler equality $b_1=2h^{0,1}$.  This proves the stated
exclusion through these invariants.
See \cref{maps:src:E.2.1}.
\item\hypertarget{step-V-A-5}{} The universal Todd and Pontryagin formulas give
$c_2^2=2$, $c_4=6$ and signature two.  The middle lattice is
positive definite unimodular of rank two, hence has the displayed
form.  
See \cref{maps:src:E.2.1}.
\end{enumerate}

\subsection{Known and new}\label{pub:V:known-new}
\begin{center}\small
\begin{tabularx}{\textwidth}{@{}lX@{}}\toprule
Theorem & Status and principal cited inputs\\\midrule
\ref{series:V:D} & New compact complex realization. Smooth classification: \cite[Definition~3.6 and Theorem~3.7(1), arXiv v3, pp.~12--14]{IVVWallForm}; sphere absorption: \cite[Theorem~A, p.~178]{IVVKramerStolz}; connected-sum criterion and tangent classes: \cite[Corollary~1.4, p.~4, Lemma~2.1, p.~5, and Theorem~2.4, p.~7]{ffMandal}.\\
\ref{series:V:E} & New parameter domain and invariants. Weighted-base and direct-image method: \cite[Lemmas~6.1--6.2, pp.~74--75, Proposition~6.10, pp.~79--80]{DGKX2026v1}; torus deformations: \cite[Theorem~2.1, p.~264]{Catanese2004}.\\
\ref{series:V:C} & New genus-two pencil and integral lifted-period construction. The cusp model is \cite[Theorem~8.7, arXiv v1, pp.~57--59]{Yang2026}.\\
\ref{series:V:B} & New integral rank-six gluing and homology. Toroidal filling and Euler number: \cite[Theorem~8.7 and Corollary~8.8, arXiv v1, pp.~57--59]{Yang2026}; integral specialization: \cite[\S7.7, pp.~60--63]{AC26}.\\
\ref{series:V:F} & New integral classification and period theorem. Lattice isomorphism: \cite[Theorem~3.4, p.~6, Algorithm~8.3, p.~21]{rootBHJ22}; local-spray approximation: \cite[Proposition~7.13 and preceding paragraph, arXiv v5, pp.~23--24]{IVVOkaOneV5}.\\
\ref{series:V:A} & New small crepant resolution and its invariants. Six-sphere: \cite[Theorems~7.17 and~7.22, pp.~57 and~60]{AC26}; elliptic surface: \cite[\S2, arXiv v1, p.~3]{E26}; characteristic identities: \cite[(1.2)--(1.4), pp.~3--4]{ffEsoleKang}; class-$\mathcal C$ criterion: \cite[Theorem~1.1, Remark~1.2, p.~2]{ffJiaMeng}.\\
\bottomrule\end{tabularx}\end{center}
\subsection{Methodology}
This text is AI-produced; no part of it was written by a human.
GPT-5.6 Sol (OpenAI), in Codex at reasoning effort ultra, wrote the
working edition on 24 August--9 September 2026. GPT-6 Astra (OpenAI),
in ChatGPT Pro mode, carried out the research of 12--16 November,
the independent re-derivation of 16--18 November, and the writing of
16--20 November 2026. The publication revision and further checks of
this paper were undertaken by GPT-6 Astra on 20--21 November 2026.
The reasoning-effort setting of this publication session was inherited;
the interface exposed no separate numerical setting.

Claude Fable 5.2 (Anthropic), at effort high and at effort max, and
Claude Opus 5.5 (Anthropic), at effort max, undertook readings, intakes,
checks and recomputations on 12--20 November 2026. Their objections
and calculations were compared with the statements, proofs and cited
sources, and the resulting corrections were incorporated into the text.
The records specify the passages and calculations examined by each model.

For the statements printed in this paper, the record contains
334 examinations of proof steps, 45 independent re-derivations, and
166 computations or formal certificates. It contains another eight
proof examinations and three re-derivations associated specifically
with steps of the introductory outlines. Each total sums recorded
links from a check to a statement, so one calculation applied to several
statements can contribute several links. Each outcome states the part
examined. The 301 readings associated with the current statements are
recorded separately. Numerical interval enclosures are accompanied by
an independently implemented numerical comparison and exact checks of
the finite covering and sign inequalities. Their respective scopes are
given in the records. No human has checked the mathematics.

The historical sources, programs and check records are available in the
\href{https://github.com/KokunoYumeto/s6-complex-structure-workbench/tree/\SsixRepositoryCommit/pro_returns/S6_SERIES_GRIND_checkpoint_10_20261120T0211Z/}{project repository}.
The additional records and programs of this publication revision
accompany its source distribution.
The text and accompanying work are dedicated to the public domain under
\href{https://creativecommons.org/publicdomain/zero/1.0/}{CC0 1.0 Universal}.

\subsection{How the paper is organized}
Part~\ref{pub:V:part-rank-six} constructs the marked rank-six family.
The local monodromy and cusp in \S\ref{pub:V:sec-local} precede the
period planes in \S\ref{pub:V:sec-periods}, the integral homology in
\S\ref{pub:V:sec-integral}, and the oriented smooth classification
in \S\ref{pub:V:sec-diffeomorphism}. Sections
\ref{pub:V:sec-parameter}--\ref{pub:V:sec-boundary} determine its
parameter domain, invariants and boundary.
Part~\ref{pub:V:part-classification} classifies the integral data and
constructs punctured-base periods. Part~\ref{pub:V:part-resolved}
gives the resolved elliptic fibre product. Part~\ref{pub:V:part-shifted}
uses the sphere fibrations to construct shifted fivefolds and determines
their integral homotopy and smooth topology.

\begin{figure}[htbp]\centering
\begin{tikzpicture}[>=stealth,every node/.style={font=\small,align=center}]
\node[draw,rounded corners,text width=3cm] (local) at (0,0) {Shared marked Kuhn filling\\\cref{night:P1.9,pub:V:cusp-smooth-dependence}};
\node[draw,rounded corners,text width=2.7cm] (period) at (-4,-1.6) {Periods\\Theorem~\ref{series:V:C}};
\node[draw,rounded corners,text width=2.7cm] (topology) at (0,-1.6) {Integral topology\\Theorem~\ref{series:V:B}};
\node[draw,rounded corners,text width=2.7cm] (resolved) at (4,-1.6) {Resolved product\\Theorem~\ref{series:V:A}};
\node[draw,rounded corners,text width=3cm] (compact) at (-2,-3.3) {Compact complex fourfold\\Theorem~\ref{series:V:D}};
\node[draw,rounded corners,text width=3cm] (family) at (-2,-4.8) {Exact family and boundary\\Theorem~\ref{series:V:E}};
\draw[->] (local)--(period);\draw[->] (local)--(topology);\draw[->] (local)--(resolved);
\draw[->] (period)--(compact);\draw[->] (topology)--(compact);\draw[->] (compact)--(family);
\end{tikzpicture}
\caption{The shared marked cusp and the main construction. The general classification of Theorem~\ref{series:V:F} is developed separately in Part~\ref{pub:V:part-classification}.}
\end{figure}

\subsubsection*{Further results}
\begin{itemize}[leftmargin=*]
\item \Cref{maps:src:E.2.3,maps:src:E.2.6}: the shifted fivefold has a torus-knot peripheral presentation and explicitly computed integral homology.
\item \Cref{maps:src:E.2.7,maps:src:E.2.9}: all transverse shifts have one oriented smooth type and are framed null-bordant over their fundamental group.
\item \Cref{maps:src:E.2.11,maps:src:E.2.17}: the cusp and peripheral maps determine the full second-homotopy module and its coinvariants.
\item \Cref{maps:src:E.2.13,maps:src:E.2.14}: a circle fibration has an infinite-order return map given by the separating torus twist.
\item \Cref{maps:src:E.2.19,maps:src:E.2.21}: the relative cusp model gives an explicit smooth framed-handle presentation.
\end{itemize}

\clearpage
\part{The marked rank-six construction and its family}\label{pub:V:part-rank-six}
\section{The monodromy and the marked cusp}\label{pub:V:sec-local}

\begin{definition}[The marked rank-six datum]\hypertarget{tag-PV02}{}\label{pub:V:rank-six-data}\hypertarget{tag-034O}{}\label{sec:fourfold}
\StatementRecord{pub:V:rank-six-data}{proved}{PV02}{}
Let $\Lambda=\mathbb Z^6$ be the marked homology lattice, with column vectors. We use the integral cohomology matrices
\[
A=\begin{pmatrix}
1&0&0&1&-1&1\\0&1&-1&-1&1&0\\0&1&0&-1&0&0\\
0&2&-1&-1&0&1\\0&2&0&-1&0&0\\0&1&1&-1&0&0
\end{pmatrix},\quad
B=\begin{pmatrix}
1&0&0&0&0&0\\0&0&-1&0&1&0\\0&0&-1&0&0&1\\
0&0&-1&0&0&0\\0&1&-1&0&0&0\\0&0&-1&1&0&0
\end{pmatrix}.
\]
Put $M_1=A^t$, $M_2=B^t$, $M_0=(M_1M_2)^{-1}$ and
\[
 v_1=(1,3,-2,-1,2,0)^t,\qquad v_2=(-1,0,0,0,0,0)^t.
\]
The affine orders are $(m_1,m_2)=(5,4)$, and $\tau_i(x)=M_ix+v_i/m_i$ on $\mathbb R^6/\mathbb Z^6$.
Positive based meridians satisfy $g_1g_2g_0=1$. The exterior-square basis of $\bigwedge^2\mathbb Z^4$ is $(01,02,03,12,13,23)$.
\end{definition}

\begin{proposition}[The cusp normal form]\hypertarget{tag-02MY}{}\hypertarget{tag-02R7}{}\label{maps:src:F.2.1}
\StatementRecord{maps:src:F.2.1}{proved}{02MY}{the definitions and constructions stated in the proof}


Let $N=M_0-I$ satisfy $N^2=0$, $\operatorname{rank}N=3$,
and suppose that $\operatorname{im}N$ is primitive in
$\Lambda=\mathbb Z^6$. Then $M_0$ is conjugate in
$\operatorname{GL}_6(\mathbb Z)$ to
\[
 U=\begin{pmatrix}I_3&I_3\\0&I_3\end{pmatrix}.
\]
Its integral centralizer consists exactly of
$\left(\begin{smallmatrix}P&Q\\0&P\end{smallmatrix}\right)$
with $P\in\operatorname{GL}_3(\mathbb Z)$ and
$Q\in\operatorname{Mat}_{3\times3}(\mathbb Z)$.

\end{proposition}


\begin{proof}
\begin{enumerate}[label=\arabic*.]
\item Let $N=M_0-I$.  The index-one condition says that the image
is the saturated kernel lattice and that
$\Lambda/\ker N\to\operatorname{im}N$ is an isomorphism.
Choose a basis $a_1,a_2,a_3$ of this image and lifts $b_i$ with
$Nb_i=a_i$.  Saturation shows that the six vectors form an
integral basis.  It gives the displayed matrix.  Commuting a
block matrix with $N$ forces its lower-left block to vanish
and its two diagonal blocks to agree.  Invertibility over
$\mathbb Z$ is then equivalent to unimodularity of that block.

\end{enumerate}
\end{proof}

\begin{theorem}[The rank-six monodromy]\hypertarget{tag-02QI}{}\label{night:P1.1}
\StatementRecord{night:P1.1}{proved}{02QI}{the definitions and constructions stated in the proof}
The finite homology monodromies have orders $5,4$, respectively, and
characteristic polynomials
$(T-1)^2\Phi _5(T)$ and $(T-1)^2(T+1)^2(T^2+1)$.
For $N=M_0-I$ we have $N^2=0$, $\operatorname{rank}N=3$, and
nonzero Smith factors $1,1,1$.  Moreover
\[
 \Lambda^{\langle M_1,M_2\rangle}=0,\qquad
 \Lambda_{\langle M_1,M_2\rangle}=\mathbb Z,
\]
where the quotient homomorphism is $\ell(x)=x_0$.
On the cohomology lattice the corresponding groups are $\mathbb Ze_0$
and $\mathbb Z/2$.  Every nonidentity power of each specified affine
action $\tau_i$ is free.
\end{theorem}
\begin{proof}
\begin{enumerate}[label=\arabic*.]
\item Exact multiplication gives
\[
N=\begin{pmatrix}
0&0&0&0&0&0\\0&0&0&1&0&-1\\0&0&0&0&1&1\\
-1&0&0&0&0&0\\1&0&0&0&0&0\\-1&0&0&0&0&0
\end{pmatrix},\qquad
S=\begin{pmatrix}
0&0&0&1&0&0\\0&1&0&0&0&0\\0&0&1&0&0&0\\
-1&0&0&0&1&0\\1&0&0&0&0&1\\-1&0&0&0&0&0
\end{pmatrix}.
\]
The matrix $S$ is unimodular and
$S^{-1}M_0S=\left(\begin{smallmatrix}I_3&I_3\\0&I_3\end{smallmatrix}\right)$.
The columns of $S$ are henceforth denoted
$(a_1,a_2,a_3,b_1,b_2,b_3)$, so that $Nb_j=a_j$.
This integral change of basis also proves the Smith assertion.
The Hermite basis of
$(M_1-I)\Lambda+(M_2-I)\Lambda$ is $(e_1,\ldots,e_5)$;
the stacked matrix of the two fixed-vector equations has rank six.
Transposing the matrices gives a fixed lattice $\mathbb Ze_0$ and
coinvariant Smith factors $(1,1,1,1,1,2)$.  The pairing
$\Lambda^\vee\times\Lambda\longrightarrow\mathbb Z$ identifies the
fixed cohomology vector $e_0$ with the covector $\ell$.

\item We have $M_iv_i=v_i$, $\ell(v_1)=1$, and $\ell(v_2)=-1$.
Thus $\tau_i^k(x)=M_i^kx+kv_i/m_i$, and $\tau_i^{m_i}=1$ on the
torus.  A fixed point of $\tau_i^k$ would give
$k\ell(v_i)/m_i=0$ in $\mathbb R/\mathbb Z$.  For
$0<k<m_i$ this is impossible.

\end{enumerate}
\end{proof}

\begin{theorem}[The Kuhn filling at the cusp]\hypertarget{tag-02R0}{}\label{night:P1.9}
\StatementRecord{night:P1.9}{proved}{02R0}{the definitions and constructions stated in the proof}
Let $C(q)$ be a holomorphic $3\times3$ matrix on a disc about
zero, and put $s=\log q/(2\pi i)$ on its punctured universal cover.
After shrinking the disc, the periodic Kuhn fan in the marking $S$
gives a smooth proper holomorphic filling of the periods
$Z(q)=sI_3+C(q)$.  Its central fibre is reduced with normal crossings.
Modulo the index-one period lattice it has respectively $1,7,12,6$
strata of dimensions $3,2,1,0$, each isomorphic to the corresponding
complex torus $(\mathbb C^*)^d$, and Euler number $6$.
An index-$d$ period sublattice gives $d$ times these stratum counts and
Euler number $6d$.
\end{theorem}
\begin{proof}
The fan consists of the cones over the tetrahedra
\begin{equation}
 [m,m+e_{\pi_1},m+e_{\pi_1}+e_{\pi_2},m+(1,1,1)],
 \quad m\in\mathbb Z^3,\quad\pi\in S_3.
 \hypertarget{tag-029R}{}\label{maps:eq:E.4}
\end{equation}
The period-lattice action is
\begin{equation}
 (x,q)\longmapsto
 \bigl(e^{2\pi iC(q)\lambda}q^\lambda x,q\bigr),
 \qquad\lambda\in\mathbb Z^3,
 \hypertarget{tag-029S}{}\label{maps:eq:E.5}
\end{equation}

\begin{enumerate}[label=\arabic*.]
\item Apply \cite[Theorem~8.7 and Corollary~8.8, arXiv v1, pp.~57--59]{Yang2026}
with $n=3$, $t=q$, $D(q)=C(q)$, $B=I_3$, and the displayed
periodic Kuhn triangulation.  These substitutions give exactly the
action
\[
 (z,q)\longmapsto
 \bigl(q^\lambda\exp(2\pi iC(q)\lambda)z,q\bigr)
\]
and the period order $[I_3\mid sI_3+C(q)]$, where
$q=\exp(2\pi i s)$.  Thus the identity on the dense torus, with
the same integral monomial charts, identifies our quotient with
the cited quotient.  The source theorem supplies its smoothness,
Hausdorffness, properness, and the reduced local equations
$q=x_0\cdots x_k$.  Flatness follows since this nonconstant
function is a nonzero divisor in each smooth chart.
\item For an index-$d$ period sublattice $\Gamma=B\mathbb Z^3$,
the same substitutions give the source's quotient map induced
by $\Gamma\subset\mathbb Z^3$.  It has degree $d$, and
Corollary~8.8 gives $\chi(W)=3!d=6d$ for this marked Kuhn fan.
\end{enumerate}
\end{proof}

\begin{lemma}[Smooth dependence on the bounded cusp correction]
\hypertarget{tag-PV01}{}\label{pub:V:cusp-smooth-dependence}
\StatementRecord{pub:V:cusp-smooth-dependence}{proved}{PV01}{}
Let $n\in\{2,3\}$ and let $C_t(q)$, $0\leq t\leq1$, be a
smooth family of holomorphic $n\times n$ matrix functions on a disc,
uniformly bounded on a smaller closed disc. After choosing a sufficiently
small common disc, the marked Kuhn fillings with periods
$[I_n\mid (\log q/(2\pi i))I_n+C_t(q)]$ are diffeomorphic over
that common disc. The diffeomorphisms preserve the central-fibre
strata and the marked angular and deck lattices. They can be chosen
on the outer collar to extend a smooth marked trivialization of the
regular torus bundles.
\end{lemma}
\begin{proof}
For the linear deformation $C_t=tC$,
\cite[Theorem~8.7, arXiv v1, pp.~57--58]{Yang2026} applies directly
in dimension $n$ to $D(q,s)=sC(q)$ for $|s|<2$.
For a smooth bounded family the same local charts give the following
comparison. The maximal toric charts have coordinates
$(z_0,\ldots,z_n)$ with $q=z_0\cdots z_n$. Their changes of
coordinates depend smoothly on $t$ and are toric monomials multiplied
by the holomorphic units $\exp(2\pi iC_t(q)\lambda)$.
They therefore define a smooth family of manifolds with the same
coordinate-stratum charts. The parameter vector field in each chart
lifts $\partial/\partial t$, fixes $q$, and is tangent to every
coordinate stratum. Differences of these lifts are vertical fields
with the same properties. A partition of unity gives a global lift;
on an outer collar take the prescribed lift from the regular marked
torus bundle and use a partition subordinate to that collar and the
interior charts.

For $y=\operatorname{Im}(\log q/(2\pi i))$, uniform boundedness
of the corrections gives, on a sufficiently small common disc,
\[
 \tfrac12\|\lambda\|\leq
 \|(I_n+y^{-1}\operatorname{Im}C_t(q))\lambda\|
 \leq\tfrac32\|\lambda\|,
 \qquad 0\leq t\leq1.
\]
The same fundamental-domain estimate as in the local filling theorem
therefore gives a proper family over the closed disc times $[0,1]$.
The lifted field has a flow for the entire parameter interval.
That flow fixes $q$, preserves each stratum, and has the specified
collar restriction. Since the parameter interval is contractible
and the lattice marking is constant there, its induced lattice map
is the identity.
\end{proof}

\begin{theorem}[The central cusp fibre and its collapse map]\hypertarget{tag-02QK}{}\label{night:P1.11}\hypertarget{tag-02M6}{}\label{maps:src:E.1.3}
\StatementRecord{night:P1.11}{proved}{02QK}{\cref{night:P1.9}}
The index-one central cusp fibre $C$ has torsion-free integral homology
with ranks
\begin{equation}
 (b_0,\ldots,b_6)=(1,3,9,9,10,3,1).
 \hypertarget{tag-029T}{}\label{maps:eq:E.6}
\end{equation}
The collapse $\kappa:T^6\to C$ has primitive image in every degree;
it is onto except in degree four, where its cokernel is $\mathbb Z$.
Its image ranks are $(1,3,9,9,9,3,1)$.
In every degree its kernel is the saturation of
$(\bigwedge^kM_0-I)\bigwedge^k\Lambda$.
Dually, the specialization
$\kappa^*:H^k(C;\mathbb Z)\to H^k(T^6;\mathbb Z)$
has image equal to the full monodromy-invariant lattice.
Its kernel is zero except in degree four, where it is $\mathbb Z$.
\end{theorem}
\begin{proof}
\begin{enumerate}[label=\arabic*.]
\item Apply \cref{pub:V:cusp-smooth-dependence} to $C_t(q)=tC(q)$.
The comparison preserves the central strata, the integral marking,
and the real oriented logarithmic boundary; it therefore intertwines
the collapse maps. We calculate with $C(q)=0$.

\item Here is a concrete topological model.  Put $G=4I_3-\mathbf1\mathbf1^t$.
Its positive eigenvalues are $1,4,4$.  The Voronoi polytope in dual
coordinates is cut out by
\[
|v^ty|\leq \tfrac12v^tGv
 \quad (0\ne v\in\{0,1\}^3).
\]
The value on a binary vector with $k$ nonzero entries is $k(4-k)$.
These fourteen primitive facet normals give the star of a vertex in
the Kuhn triangulation.  One checks its Delaunay tetrahedra by the
successive chains
$0,e_i,e_i+e_j,e_1+e_2+e_3$ for the six permutations $(i,j,k)$;
solving their three perpendicular-bisector equations and the fourteen
inequalities verifies precisely their common faces.  The translates
of this Voronoi polytope by $G\mathbb Z^3$ tile $\mathbb R^3$.

\item For a simplex $\sigma$ with vertices $v_0,\ldots,v_r$, put
$L_\sigma=\sum_i\mathbb Z(v_i-v_0)$ and
$D_\sigma=\bigcap_{v\in\sigma}P_v$, where $P_v=P+Gv$.
The component torus has lattice $A=\mathbb Z^3$; the common-stratum
marking is the explicit map
\[
 (A\oplus\mathbb Z)/\sum_i\mathbb Z(v_i,1)
 \longrightarrow A/L_\sigma,\qquad
 [(a,k)]\longmapsto a-kv_0.
\]
The Kuhn unimodularity makes $L_\sigma$ primitive.

We choose compatible moment maps as follows.  The polytope $2P_v$
is integral, since its vertex facet normals are unimodular and its
facet constants are integral.  Embed the component $X_v$ using the
characters $\chi^{-m}$ for
$m\in\mathcal A_v=2P_v\cap A^\vee$, and define
\[
 \mu_v([Z_m])=
 \frac{\sum_{m\in\mathcal A_v}m|Z_m|^2}
 {2\sum_{m\in\mathcal A_v}|Z_m|^2}.
\]
The sign is specified: the ray $w-v$ tends to the face on which
$\langle w-v,-m\rangle$ is minimal, hence to the upper common
Voronoi face of $P_v$ and $P_w$.  Thus the character polytope is
$-2P_v$, and $\mu_v$ is minus one-half its usual moment map.
The monomial moment theorem
\cite[equations (12.2.3)--(12.2.4) and Theorem 12.2.5, p.~571]{CoxLittleSchenck2011}
identifies the nonnegative toric part with $P_v$.
For the squared weights here, the intervening map on that part is
$[a_m]\mapsto[a_m^2]$, with inverse $[b_m]\mapsto[\sqrt{b_m}]$:
both preserve the homogeneous binomial relations and positive common
projective rescaling.  Write $s_v:P_v\to X_v$ for the resulting
continuous positive section.

On a shared face put
$\mathcal A_\sigma=2D_\sigma\cap A^\vee$ and choose
$m_0\in\mathcal A_\sigma$.
For $m\in\mathcal A_\sigma$, the difference $m-m_0$ annihilates
$L_\sigma$.  The ratio $Z_m/Z_{m_0}$ is consequently the restriction
of the same character
\[
 \bigl(-(m-m_0),\langle m-m_0,v\rangle\bigr)
       \in(A\oplus\mathbb Z)^\vee
\]
for every vertex $v$ of $\sigma$.  At $C=0$ all deck unit multipliers are one, and the toric gluing
identifies these character functions exactly.  On either component
the restricted moment map is therefore the identical expression
\[
 \frac{\sum_{m\in\mathcal A_\sigma}
       m|\chi^{-(m-m_0)}|^2}
 {2\sum_{m\in\mathcal A_\sigma}|\chi^{-(m-m_0)}|^2}.
\]
The equality persists on every smaller face by restricting the index
set.  Positive character values also agree, so the sections $s_v$
glue on the entire shared orbit closure.  The angular maps are the
displayed quotient homomorphisms $A\to A/L_\sigma$.

The period deck map is $(a,k)\mapsto(a+k\nu,k)$, sending $v$ to
$v+\nu$.  Its component marking is unchanged because
$a+k\nu-k(v+\nu)=a-kv$.
Moreover $\mathcal A_{v+\nu}=\mathcal A_v+2G\nu$; the common
character factor cancels in projective ratios.  Hence
\[
 \mu_{v+\nu}(T_\nu p)=\mu_v(p)+G\nu,
 \qquad s_{v+\nu}(y+G\nu)=T_\nu s_v(y).
\]
These give an explicit well-defined map
\[
 \begin{gathered}
 \mathcal Q=
 \bigl((\mathbb R^3/G\mathbb Z^3)\times
       (\mathbb R^3/\mathbb Z^3)\bigr)/\!\sim
 \longrightarrow C,\\
 [y,\theta]\longmapsto e^{2\pi i\theta}\,s_v(y)
 \quad(y\in P_v).
 \end{gathered}
\]
where over $\operatorname{relint}D_\sigma$ the relation collapses
the subtorus $L_\sigma\otimes(\mathbb R/\mathbb Z)$.
On each stratum the inverse takes its moment value and its angular
character class in $T_A/T_{L_\sigma}$; the latter is obtained by
$z\mapsto z/|z|$ on its nonzero characters.
The map is a continuous bijection, and its source is compact and
$C$ Hausdorff, so it is a homeomorphism.
The angular identity $[y,0]$ gives the continuous section, with
the same face and deck compatibilities just proved.

\item The nearby-fibre collapse is the quotient map in this model.
Give the semistable family its divisorial logarithmic structure.
The local equations $q=x_0x_1x_2x_3$ and their face restrictions
make it fine, saturated, log smooth and vertical over the standard
logarithmic disc.  The family is proper by \cref{night:P1.9}.
In its logarithmic central fibre, use all arguments of the $x_j$
with their sum fixed by $\arg q$.  The integral angular lattice over
a face with vertices $v_0,\ldots,v_r$ is the kernel of the height
map on $\mathbb Z^3\oplus\mathbb Z$.  Forgetting the arguments of
the zero coordinates induces exactly
\[
 \mathbb Z^3\longrightarrow
 (\mathbb Z^3\oplus\mathbb Z)/
       \sum_i\mathbb Z(v_i,1)
 \ \cong\ \mathbb Z^3/L_\sigma.
\]
The last isomorphism sends the class of $(a,k)$ to
$a-kv_0\pmod {L_\sigma}$; its inverse sends $a$ to $(a,0)$.

The polar maps agree under the integral monomial overlaps.
At $C(q)=0$ the deck multipliers have positive real values, so the
logarithmic fibre at angle one is the product of the base and angular
three-tori in item~3, with their given markings.

The mapping-cylinder comparison for a proper vertical log-smooth
family identifies this logarithmic forgetful map with the nearby-fibre collapse, compatibly with the central strata
\cite[formula (3-3-1), p.~478; Theorem 4.1.1 and Lemmas
4.1.2--4.1.3, pp.~480--482]{ffAchingerOgus20}.
Its angular cycles are $a_j$ and its base cycles are $b_j$.
In these column-basis conventions monodromy fixes each $a_j$
and sends $b_j$ to $b_j+a_j$, as specified by $Nb_j=a_j$.
The stratified deformation in item~1 transports this map
back to the original bounded period matrix.

\item Filter $\mathcal Q$ by the dimension $d$ of its dual base cells.
Over the interior of the dual cell of a $p$-simplex, $d=3-p$ and the
angular quotient has lattice
$Q_\sigma=\mathbb Z^3/L_\sigma$ of rank $d$.
Excision for the cell and its boundary gives the relative chain
groups of the resulting homology spectral sequence:
\[
E^1_{d,q}=\bigoplus_{\dim\sigma=3-d}\bigwedge^qQ_\sigma.
\]
Orient dual cells by the orientation of the base and the ordered
simplex.  The differential is the signed quotient map for each
codimension-one coface.  More explicitly, with integral projection
matrices $Q_\sigma$ and chosen right inverses $S_\sigma$,
\begin{equation}
D_{p,q}[\tau,\sigma]=
 \sum_{\sigma=\partial_j\tau\ {\rm mod}\ \mathbb Z^3}
 (-1)^j\bigwedge^q(Q_\tau S_\sigma).
\label{pub:V:eq:1-11-1}
\end{equation}
All $L_\sigma$ are primitive because the tetrahedra are unimodular.
Thus these projections and right inverses are integral.  Formula
\eqref{pub:V:eq:1-11-1} specifies every integer entry, including multiple incidences
under translation identifications. There are $26$ simplex classes and $74$ codimension-one incidences;
the signed quotient formula above determines their integer matrices.

\item The nonzero $E^2_{d,q}$ ranks are
\[
\begin{array}{c|rrrr}
q\backslash d&0&1&2&3\\\hline
0&1&3&3&1\\
1&0&6&8&3\\
2&0&0&7&3\\
3&0&0&0&1
\end{array}
\]
and every group is free. The reconstructed incoming Smith matrices
have only unit nonzero factors and reproduce every entry of this table.
Support leaves only the possible higher differential
$d_2:E^2_{3,0}\to E^2_{1,1}$.
The angular identity section splits the base projection as a map of
filtered spaces, so the base fundamental class survives.  This kills
that differential.  All subsequent differentials have zero source or
target.  The filtration extensions are extensions of free abelian
groups and therefore split, proving the asserted integral homology.

\item The product-to-quotient map induces on $E^1$ the exterior powers of
the quotient projections.  On integral cycle and quotient bases, the face projections and the
Poincar\'e-dual Alexander--Whitney formula in the following item give the induced maps. Their images are primitive, and the only nonzero graded cokernel
is $\mathbb Z$ at $(d,q)=(2,2)$.  The summed image ranks are
\[
             r_k=(1,3,9,9,9,3,1).
\]
The same exact exterior-power calculation gives these ranks for
$\ker(\bigwedge^kM_0-I)$.
The collapse satisfies
$\kappa_*(\bigwedge^kM_0-I)=0$, because a monodromy circuit
has the same retraction to the central fibre in the mapping-cylinder
comparison.  Therefore its rank is at most $r_k$.
A filtered map has rank at least the sum of the ranks of its
graded maps; hence the rank is exactly $r_k$.

Choose integral bases compatible with the source and target
filtrations; their graded groups are free.  In each graded block
choose the unit Smith pivots of the computed map.  The corresponding
$r_k$ rows and columns of the filtered map form a block
triangular submatrix with unimodular diagonal pivot blocks.
Its determinant is $\pm1$.  Since the full map has rank $r_k$,
these unit pivots prove that its image is primitive.
Thus the map is onto in every degree except four, where its
cokernel is the free group of rank $10-9=1$.
This rank argument also rules out additional lower-filtration images
of source classes in the kernels of the graded maps.

\item We represent a base cycle $b_J$ by
Poincar\'e duality from the Alexander--Whitney cup cocycle on the
complementary coordinates.  On an ordered $p$-simplex its coefficient
is
\[
\operatorname{sgn}(J,J^c)
 \prod_{r=1}^p(v_r-v_{r-1})_{(J^c)_r}.
\]
Tensoring with an angular exterior vector and applying \eqref{pub:V:eq:1-11-1}
gives the marked matrices. Their source basis is the ordered exterior
basis $(a_1,a_2,a_3,b_1,b_2,b_3)$; the output gives the integral
target homology projections and all Smith basis maps.
They annihilate $\bigwedge^kM_0-I$ exactly.  The collapse also
annihilates that operator, since turning the smoothing parameter
does not change its image in the central fibre.  The computed ranks
equal the ranks of the monodromy coinvariants; primitivity therefore
identifies both kernels with the saturation of its image.
Consequently these are the matrices of the collapse in suitable integral target bases.

\item Every homology group just computed is free, so the universal
coefficient maps identify cohomology with the integral duals of homology.
The transpose of the collapse has image
\[
 \operatorname{im}\kappa^*
 =\operatorname{Ann}(\ker\kappa_*)
 =\operatorname{Ann}\bigl((\bigwedge^kM_0-I)\bigwedge^k\Lambda\bigr)
 =\ker\bigl((\bigwedge^kM_0)^t-I\bigr).
\]
The first equality uses the primitive image of $\kappa_*$: an integral
functional on that image extends to the free central-fibre group.
Passing to saturation preserves the annihilator. The last group is
the full integral cohomology invariant lattice; inverse transpose
has the same fixed subgroup. The kernel of $\kappa^*$ is the
integral dual of $\operatorname{coker}\kappa_*$, giving precisely the
stated degree-four kernel.
\end{enumerate}
\end{proof}

\begin{remark}[The branch resolution of the Kuhn fibre]
\hypertarget{tag-PV08}{}\label{pub:V:cusp-branch-resolution}
\StatementRecord{pub:V:cusp-branch-resolution}{proved}{PV08}{}
The branch resolution and the moment-cell complex have the following
augmentation. Over a dual cell $D_\eta$ of the Kuhn decomposition,
the angular lattice is $Q_\eta=\mathbb Z^3/L_\eta$ and the branch
indices form the simplex on the vertices of $\eta$. Tensor its
augmented integral simplex complex with $\bigwedge^qQ_\eta$ and
sum over dual cells. The degree-zero augmentation adds branch
coefficients; the higher branch terms map to zero. The incidence
maps $Q_\eta\to Q_\tau$ for cofaces $\tau\supset\eta$ commute
with this augmentation. Inserting the least vertex contracts each
augmented simplex complex. Thus the augmentation is a
quasi-isomorphism to the moment-cell complex of
\cref{night:P1.11}. Taking cohomology of the compact toric branch
intersections first gives the normalization table below.

The normalization is the three-dimensional permutohedral toric variety.
Its paired boundary surfaces consist of four copies of $\mathrm{dP}_6$
and three copies of $\mathbb P^1\times\mathbb P^1$.  There are twelve
normalized triple curves and six quadruple points.  The normalization
complex for the constant sheaf is the alternating branch complex.
Local exactness follows from the augmented simplex on the branches
of $z_0z_1z_2z_3=0$.  Its nonzero rows after taking cohomology are
\begin{equation}
\begin{array}{c|rrrr|rr}
q& p=0&p=1&p=2&p=3&\operatorname{rank}d^0&
 \operatorname{rank}d^1\ \text{or }d^2\\ \hline
0&1&7&12&6&0&4,5\\
2&11&22&12&0&5&9\\
4&11&7&0&0&4&0\\
6&1&0&0&0&0&0
\end{array}
\hypertarget{tag-029U}{}\label{maps:eq:E.7}
\end{equation}
Every nonzero Smith invariant in these matrices is one.


The integral matrices are as follows.
The fourteen rays are $\pm v$, with
$v\in\{0,1\}^3\setminus\{0\}$.  The twenty-four maximal cones
are the translated Kuhn tetrahedra containing zero.  A unit cube contains zero exactly when its lower vertex lies in
$\{-1,0\}^3$, so these eight cubes give the complete enumeration.  The Picard group of the normalization
is the group of ray coefficients modulo the three character relations.
We set the coefficients at $e_1,e_2,e_3$ to zero, giving an integral
basis of rank eleven.  On a boundary surface $D_v$, choose adjacent
rays $u_1,u_2$ for which $(v,u_1,u_2)$ is unimodular and set the two
remaining character coefficients to zero.  Restriction is evaluated
by subtracting the unique integral linear function with the prescribed
values at these three rays.  The paired restriction is
$\operatorname{res}_{D_v}-\operatorname{res}_{D_{-v}}$, with the
second adjacent ray corresponding to $u-v$.


For a triangle $(0,v,u)$ with $0<v<u$ coordinatewise, the next map
is the alternating sum of the degrees on the three boundary curves
of $D_{u-v},D_u,D_v$.  In a surface fan, if the neighbours of a ray
$w$ are $h_1,h_2$ and $h_1+h_2=bw$ modulo $v$, that degree is
$a_{h_1}+a_{h_2}-ba_w$.  In degree four, the restriction map is dual
to the seven divisor differences $[D_v]-[D_{-v}]$.  In degree zero
the matrices are the cellular incidence matrices of the one-vertex
Kuhn triangulation of $T^3$.  These rules define \eqref{maps:eq:E.7} over
$\mathbb Z$, and both compositions are zero.
Its exact Smith reductions give the ranks and primitivity just stated.


\end{remark}

\begin{lemma}[Integral degree-two classes on the Kuhn filling]
\hypertarget{tag-PV13}{}\label{pub:V:cusp-chern-generators}
\StatementRecord{pub:V:cusp-chern-generators}{proved}{PV13}{\cref{pub:V:cusp-smooth-dependence}}
For the standard Kuhn filling with bounded correction zero, in dimension
$n=2$ or $3$, every integral symmetric matrix $S$ gives a line bundle
whose Chern class on a smooth fibre is
$\sum_{i,j}S_{ij}a_i\wedge b_j$.
Here $a_i$ are the angular classes and $b_i$ the invariant deck
classes in the marked product description. The corresponding classes
on the central fibre specialize to these classes on the smooth fibre.
The deck classes $b_i$ also lie in the specialization image, and
the marked bounded-correction comparison preserves these specialization images.
\end{lemma}
\begin{proof}
At a fan vertex $m\in\mathbb Z^n$ put
\[
 q_S(m)=\frac{m^tSm-\sum_iS_{ii}m_i}{2}.
\]
This is integral because $m_i^2-m_i$ is even. Interpolation on
each unimodular Kuhn simplex gives an integral piecewise-linear
support function. Translation by $\lambda$ changes it by the
integral affine function $m^tS\lambda+q_S(\lambda)$.
The associated Cartier data consequently have deck linearization
\[
 f_\lambda(x,q)=q^{-q_S(\lambda)}x^{-S\lambda}.
\]
The identity
$q_S(\lambda+\mu)=q_S(\lambda)+q_S(\mu)+\lambda^tS\mu$
gives $f_\lambda(q^\mu x,q)f_\mu(x,q)=f_{\lambda+\mu}(x,q)$,
so the line bundle descends to the period quotient.

For the Chern-class convention write $q=e^{2\pi i s}$,
$y=\operatorname{Im}s>0$, and $x_i=e^{2\pi iz_i}$ with
$z=u+sv$. Then $a_i=du_i$ and $b_i=dv_i$.
The Hermitian squared norm
\[
 h(z)=\exp\!\left(-\frac{2\pi}{y}
 (\operatorname{Im}z)^tS(\operatorname{Im}z)
 +2\pi\sum_iS_{ii}\operatorname{Im}z_i\right)
\]
satisfies $h(z+s\lambda)|f_\lambda(z)|^2=h(z)$.
Its Chern form is
\[
 \frac{i}{2y}\sum_{i,j}S_{ij}\,dz_i\wedge d\bar z_j
 =\sum_{i,j}S_{ij}\,du_i\wedge dv_j,
\]
where symmetry of $S$ cancels the angular--angular and deck--deck
terms. Restriction of the line bundle to the central fibre gives
the stated specialization classes. The projection to the deck torus
in the moment-cell model supplies the classes $b_i$.
Finally \cref{pub:V:cusp-smooth-dependence} identifies the marked
collapse maps under a bounded correction, and hence transports
these integral classes and their specialization images.
\end{proof}

\begin{lemma}[The marked smooth gluing]
\hypertarget{tag-PV11}{}\label{pub:V:smooth-gluing}
\StatementRecord{pub:V:smooth-gluing}{proved}{PV11}{\cref{night:P1.1}, \cref{night:P1.9}, \cref{night:P1.11}; \cite[Theorem 4.1.1 and Lemmas 4.1.2--4.1.3, pp.~480--482]{ffAchingerOgus20}}
The datum of \cref{pub:V:rank-six-data}, with its two affine caps,
the primitive Kuhn cusp in the marking $S$, and zero additional
translation clutch, gives a smooth map $f:M\to S^2$ from a
closed connected oriented smooth eight-manifold $M$. It is simply
connected, has Euler number six,
and admits an almost complex structure with this orientation.
\end{lemma}
\begin{proof}
\begin{enumerate}[label=\arabic*.,leftmargin=*]
\item Use the matrices and affine vectors of \cref{pub:V:rank-six-data}.
Theorem~\ref{night:P1.1} gives their orders, the primitive cusp
marking, the coinvariant homomorphism $\ell(x)=x_0$, and freeness of
every nonidentity power of the affine actions. Consequently
\[
 (T^6\times\Delta_u)/
 \langle(x,u)\mapsto(M_ix+v_i/m_i,\zeta_{m_i}u)\rangle,
 \qquad (m_1,m_2)=(5,4),
\]
are smooth oriented finite caps. On the universal angular cover
$u=r e^{2\pi ia}$ their marked collar maps are
\[
 \alpha_i(x,u)=x-av_i,
 \qquad
 \alpha_i(M_ix+v_i/m_i,\zeta_{m_i}u)=M_i\alpha_i(x,u).
\]
At the cusp we take the index-one Kuhn cap of
Theorem~\ref{night:P1.9} with $C(q)=0$, in the explicit marking
$S$ of Theorem~\ref{night:P1.1}. The three displayed boundary bundles
therefore glue to the marked torus bundle over the intervening pair
of pants, with zero additional clutch translation. Every boundary
is paired. The result is a closed connected smooth eight-manifold,
oriented by the cusp complex structure and the continued fibre/base
orientation.

\item The small cusp cap retracts onto its central fibre $C$ through
the mapping-cylinder homeomorphism
\cite[Theorem 4.1.1 and Lemmas 4.1.2--4.1.3, pp.~480--482]{ffAchingerOgus20}.
Indeed its divisorial log chart is
$\mathbb N\to\mathbb N^r$, $1\mapsto(1,\ldots,1)$; it is log smooth,
saturated and vertical, and the cap is proper by
Theorem~\ref{night:P1.9}. In the resulting cylinder the homotopy
$[x,r]\mapsto[x,(1-t)r]$ fixes $C$ and retracts onto it.

Before the period quotient, the components of $C$ and their
nonempty intersections are smooth complete toric strata. Each is
simply connected: a loop moves into its dense torus by general
position, and the one-parameter loops along the rays of a smooth
maximal cone bound the coordinate discs; those rays are an integral
basis. Van Kampen sends component-crossing paths to edges and triple
intersections to triangles of the Kuhn nerve in $\mathbb R^3$.
Thus the covering core has trivial fundamental group. Its period
quotient has deck group $\mathbb Z^3$, so $\pi_1(C)=\mathbb Z^3$.
The nearby-torus map is the quotient of $\Lambda$ by the saturated
image of $M_0-I$, by the $H_1$ collapse in
Theorem~\ref{night:P1.11}. It is onto, and the cusp meridian dies
because its zero section extends.

We also record the invariant normal closure needed below. Let
$K$ be the smallest $\langle M_1,M_2\rangle$-invariant subgroup
containing $(M_0-I)\Lambda$, and let
$J=(M_1-I)\Lambda+(M_2-I)\Lambda=\ker\ell$.
The equality $\ell M_i=\ell$ gives $K\subseteq\ker\ell$.
On the abelian quotient $\Lambda/K$, the induced maps $a,b$
satisfy $ab=1$, $a^5=1$ and $b^4=1$. Thus $a^4=a^5=1$, so
$a=b=1$ and $J\subseteq K$. Consequently $K=J=\ker\ell$ as
integral subgroups.


Let $a_i$ be the two finite peripheral lifts. The universal covering
deck transformation of the $i$th cap has $m_i$th power equal to
translation by $v_i$, since $M_iv_i=v_i$. The collar map in
Step~1 sends $a_i$ to that transformation. Van Kampen therefore gives
\[
 a_i\lambda a_i^{-1}=M_i\lambda,
 \qquad a_1^5=v_1,\quad a_2^4=v_2,\quad a_1a_2=1.
\]
Write $a_1=t$ and $a_2=t^{-1}$. Then
$t=t^5t^{-4}=v_1+v_2$ lies in the abelian fibre subgroup. Conjugation
is consequently trivial, and that subgroup is a quotient of
$\Lambda_{\langle M_1,M_2\rangle}=\mathbb Z$, with coordinate $\ell$.
Now $\ell(v_1+v_2)=0$, so $t=0$, while $t^5=v_1$ kills its generator
$\ell(v_1)=1$. Every generator vanishes, proving $\pi_1=1$.
Equivalently, after this abelian reduction, the relations on
$(\ell,a_1,a_2)$ have matrix
\[
 \begin{pmatrix}-1&5&0\\1&0&4\\0&1&1\end{pmatrix},
 \qquad\det=-1.
\]

\item Euler additivity over the base strata gives $\chi(M)=6$:
the regular tori and finite free torus quotients have Euler number
zero, and \cref{night:P1.9} gives six for the cusp.
\item Each finite monodromy admits an invariant real linear complex
structure: its $1$ and $-1$ real eigenspaces have even dimension, and
each remaining conjugate eigenpair gives a real two-plane. Choose
a complex structure on each summand. Reversing it on one fixed
two-plane selects the required orientation. These structures commute
with $M_i$, so their direct sums with the root-disc structure descend
to the affine caps. The cusp cap already has its complex
structure. Their vertical complex structures on the boundary collars
are translation-invariant in the marked torus fibres.

The space $\mathcal J=\operatorname{GL}_6^+(\mathbb R)/
\operatorname{GL}_3(\mathbb C)$ of allowable vertical structures
retracts to $\operatorname{SO}(6)/\operatorname{U}(3)$ by the explicit
conjugation
\[
 J\longmapsto P_tJP_t^{-1},
 \qquad P_t=\bigl((1-t)I+t(I+J^tJ)\bigr)^{1/2}.
\]
At $t=1$ the structure is orthogonal, since $J$ preserves the metric
$I+J^tJ$; an orthogonal $J$ is fixed throughout this homotopy.
The inclusion $\operatorname{U}(3)\to\operatorname{SO}(6)$ sends
$\operatorname{diag}(e^{i\theta},1,1)$ to one-plane rotation and
induces the surjection $\mathbb Z\to\mathbb Z/2$ on fundamental
groups. Indeed the determinant fibration and
$\operatorname{SU}(2)=S^3\to\operatorname{SU}(3)\to S^5$ give the
first group; the quaternion cover of $\operatorname{SO}(3)$ and
the sphere fibrations for $\operatorname{SO}(n)$ give the second,
with a full rotation lifting from $1$ to $-1$. The homogeneous-space
exact sequence therefore gives $\pi_0\mathcal J=\pi_1\mathcal J=0$.

Over the two-dimensional pair of pants the prescribed boundary
sections of the associated $\mathcal J$-bundle extend over relative
one-cells by connectedness and over relative two-cells by simple
connectivity. Smooth approximation relative to the collars gives
a smooth vertical structure. Horizontal splittings form an affine
space, so a partition of unity extends the prescribed complex
splittings from those collars. Combining the vertical structure with
the pulled-back base structure and gluing the caps gives the required
almost complex structure on $M$.

\end{enumerate}
\end{proof}

\section{The genus-two pencil and the lifted period planes}\label{pub:V:sec-periods}



\begin{theorem}[The genus-two pencil and its Hodge bundle]\hypertarget{tag-02QV}{}\label{night:P1.4}
\StatementRecord{night:P1.4}{proved}{02QV}{the definitions and constructions stated in the proof}
The smooth projective models of the pencil
\begin{equation}
 C_\lambda:\quad y^2=x^5+\lambda^3x+\lambda^4
 \hypertarget{tag-02FB}{}\label{maps:eq:F.7}
\end{equation}
have one nodal fibre away from $0,\infty$, at $\lambda_*=-256/3125$.
The minimal good root orders at $0,\infty$ are $8,5$.
On the root curve with reduced root divisors $D_0=[0]/8$ and $D_\infty=[\infty]/5$, the Hodge bundle is
\begin{equation}
 \mathcal E=\mathcal T^3\oplus\mathcal T,\qquad
 \mathcal T=\mathcal O_{\mathcal B'}(2D_\infty-3D_0),\qquad
 \deg\mathcal T=\frac1{40}.
 \hypertarget{tag-02FC}{}\label{maps:eq:F.8}
\end{equation}
Thus $T=\mathcal T=\mathcal O((2/5)[\infty]-(3/8)[0])$, and $\ell_0=T^3$ is generated meromorphically by $dx/y$.
The extension at $\lambda_*$ is the dualizing Hodge extension.
\end{theorem}
\begin{proof}
\begin{enumerate}[label=\arabic*.]
\item The discriminant of $x^5+\lambda^3x+\lambda^4$ is
$\lambda^{15}(3125\lambda+256)$.
At $\lambda_*$ the double root is $64/625$;
the second $x$-derivative and first $\lambda$-derivative there are
nonzero.  The remaining roots are simple.  The odd-degree
hyperelliptic compactification is smooth at its branch point at
infinity.  This gives one node, a smooth total surface at that node,
and genus-one normalization.

\item At zero set $\lambda=u^8$, $x=u^6X$, $y=u^{15}Y$; at infinity set
$\lambda=u^{-5}$, $x=u^{-4}X$, $y=u^{-10}Y$.
The resulting smooth equations and differentials are
\[
\begin{array}{c|c|c|c}
 &\text{equation}&dx/y&x\,dx/y\\\hline
0&Y^2=X^5+X+u^2&u^{-9}dX/Y&u^{-3}X\,dX/Y\\
\infty&Y^2=X^5+uX+1&u^6dX/Y&u^2X\,dX/Y.
\end{array}
\]
Their central polynomials are square-free.  The deck maps at zero and
infinity are $(X,Y)\mapsto(iX,\zeta_8Y)$ and
$(X,Y)\mapsto(\zeta_5^{-1}X,Y)$, respectively.
The eigenvalues on $H^{1,0}$ are $(\zeta_8,\zeta_8^3)$ and
$(\zeta_5^{-1},\zeta_5^{-2})$.  Including conjugates gives characteristic
polynomials $T^4+1$ and $\Phi_5$, proving minimality of the root orders.
The fourth power at zero is the hyperelliptic involution, acting as
$-I$ on $H^1$ and as $I$ on its exterior square.

\item The two meromorphic differential generators have orbifold divisors
$(6/5)[\infty]-(9/8)[0]$ and
$(2/5)[\infty]-(3/8)[0]$.
At the node they are independent sections of the dualizing Hodge
extension: one has opposite nonzero residues at the two points of the
normalization, and the linear combination vanishing at the node is a
regular nonzero elliptic differential.  Thus there is no additional
divisor at $\lambda_*$.  These computations prove the asserted direct
sum and its degrees.

\end{enumerate}
\end{proof}









\begin{theorem}[The rational rank-five monodromy]\hypertarget{tag-02O7}{}\label{maps:src:F.3.5}
\StatementRecord{maps:src:F.3.5}{proved}{02O7}{\cref{maps:src:F.2.5}}


For $U=\mathbb P^1\setminus\{0,\infty,\lambda_*\}$, let
$\mathcal H_{\rm dR}$ be the polarized rational cohomology system
of the pencil in \cref{night:P1.4}. Its primitive exterior square
is rationally isomorphic to the rank-five quotient of the system
with matrices $A,B$ in \cref{pub:V:rank-six-data} by its invariant
first-coordinate line. The identification uses the explicit
companion and quotient maps below and the stated based meridians.
\end{theorem}


\begin{proof}
\begin{enumerate}[label=\arabic*.]
\item By \cref{night:P1.4}, the cusp is an ordinary node with
genus-one normalization. Cutting and regluing an annular neighborhood
gives a Dehn twist in the primitive nonseparating vanishing circle.
Its logarithm $N_H$ has square zero and rank one.
The finite characteristic polynomials, from the same theorem, are
$f(T)=T^4+T^3+T^2+T+1$ at infinity and $g(T)=T^4+1$ at zero.

\item We identify the representation.
Choose positive peripheral generators $T_1,T_2$ at infinity and
zero, in that order, with $T_1T_2$ inverse to nodal monodromy.
Set $R=T_1^{-1}$ and $S=T_2$.  Then $R-S$ has rank one and their
characteristic polynomials are $f,g$.  The companion presentation
is Levelt's theorem in \cite[Theorem 3.5, pp.~331--332;
Corollary 3.6, p.~332]{BeukersHeckman1989}.  Its generators are
related to the present pair by
\[
 h_\infty=R=T_1^{-1},\qquad h_0=S^{-1}=T_2^{-1},\qquad
 h_1=R^{-1}S=T_1T_2,\qquad h_\infty h_1h_0=I.
\]
The following calculation supplies the rational change
of basis for this presentation.
If $u$ spans $\operatorname{im}(R-S)$, it is cyclic for $R$:
the span of its iterates is stable under both $R$ and $S$, and
the irreducibility of $f$ makes that span all of $\mathbb Q^4$.
In the rational basis $u,Ru,R^2u,R^3u$, the matrix of $R$ is
the companion matrix of $f$, and $S=R-u\varphi$ for a covector
$\varphi$.  The determinant identity
\begin{equation}
 g(T)=f(T)+\varphi\operatorname{adj}(TI-R)u
 \hypertarget{tag-02FD}{}\label{maps:eq:F.9}
\end{equation}
determines $\varphi$ uniquely: the four coordinates of
$\operatorname{adj}(TI-R)u$ are a basis of the polynomials of
degree at most three, as is seen from their successively distinct
leading degrees in the cyclic basis.  It follows that any two
such pairs are simultaneously rationally conjugate.  Apply the
lemma to the explicit companion pair
\begin{equation}
 C_f=\begin{pmatrix}
 0&0&0&-1\\1&0&0&-1\\0&1&0&-1\\0&0&1&-1
 \end{pmatrix},\qquad
 C_g=\begin{pmatrix}
 0&0&0&-1\\1&0&0&0\\0&1&0&0\\0&0&1&0
 \end{pmatrix}.
 \hypertarget{tag-02EG}{}\label{maps:eq:F.10}
\end{equation}
Their difference has rank one and their characteristic polynomials
are $f,g$.  More explicitly, put
\[
 \begin{aligned}
 \mathsf B&=(u,Ru,R^2u,R^3u),\\
 d&=(C_f-C_g)e_4=(0,-1,-1,-1)^t,\\
 \mathsf B_c&=(d,C_fd,C_f^2d,C_f^3d).
 \end{aligned}
\]
Here $\det\mathsf B_c=1$.  Uniqueness of the covector in
\eqref{maps:eq:F.9} gives the intertwiner
$\mathsf P=\mathsf B\mathsf B_c^{-1}$, with
$R\mathsf P=\mathsf P C_f$ and $S\mathsf P=\mathsf P C_g$.
Thus $(T_1,T_2)$ is rationally conjugate to
$(C_f^{-1},C_g)$, with the chosen peripheral conventions.

\item Use the lexicographic basis in $\bigwedge^2\mathbb Q^4$.
Its common invariant vector is $(1,1,1,1,1,1)^t$.
Let $Q$ be the identity matrix of size six with first column
replaced by this vector.  It is unimodular.  Direct multiplication
gives the exact identities
\begin{equation}
 \left(Q^{-1}\bigwedge^2(C_f^{-1})Q\right)_{2:6,2:6}
   =A_{2:6,2:6},\qquad
 \left(Q^{-1}\bigwedge^2(C_g)Q\right)_{2:6,2:6}
   =B_{2:6,2:6}.
 \hypertarget{tag-02EH}{}\label{maps:eq:F.11}
\end{equation}
The invariant line is the inverse of the polarization, up to a
nonzero rational scalar.  Projection along that line identifies
the quotient with the primitive exterior square $\mathcal W$.

\item The differentials $\omega_1=dx/y$ and $\omega_2=x\,dx/y$ span
the extended Hodge bundle everywhere.  At the node they are
independent dualizing differentials: $\omega_1$ has nonzero
opposite residues on the normalization and
$(x-r)\omega_1$ is a nonzero holomorphic differential there.
The extension assertion can be checked on the nodal annulus.
In coordinates $uv=t$ a relative dualizing differential has
a convergent Laurent expansion
$\sum_n a_n(t)u^n\,du/u$.  Integration on the vanishing circle
is $2\pi i a_0(t)$.  Integration across the annulus has the
single logarithmic term $a_0(t)\log t$; integration of every
nonconstant Laurent monomial and of the complementary fixed
arcs gives a holomorphic remainder.  Thus the period vector
has the form $\exp(sN_H)$ times a holomorphic vector,
$s=\log t/(2\pi i)$.  Untwisting by $\exp(-sN_H)$ is exactly
the nilpotent canonical extension used here.  The two independent
dualizing limits remain a basis of its Hodge subbundle.
Their isotropic pairing, zero away from zero, remains zero
at the limit by holomorphic continuation, so that limit is
Lagrangian.
The root-divisor calculation is \eqref{maps:eq:F.8}.
The nonzero residue of $\omega_1$ also shows
\begin{equation}
 N_H\ell_\infty\ne0,
 \qquad \ell=\mathcal T^3\subset\mathcal E,
 \hypertarget{tag-02EI}{}\label{maps:eq:F.12}
\end{equation}
where $\ell_\infty$ denotes its limit at the nodal cusp.
Indeed integration on the vanishing circle tends to
$2\pi i$ times that residue; the transvection sends the cohomology
class of a differential to that period times its vanishing
cohomology vector.

\end{enumerate}
The algebraic calculations are implemented in \hyperlink{prog:V:periods}{the programs in Table~\ref*{tab:V:programs}}.
\end{proof}




\begin{theorem}[The integral exterior-square quotient]\hypertarget{tag-02QT}{}\label{night:P1.2}
\StatementRecord{night:P1.2}{proved}{02QT}{the definitions and constructions stated in the proof}
The lower blocks of $A,B$ are precisely the quotient representations of
$\bigwedge^2(C_f^{-1}),\bigwedge^2C_g$ in the basis $Q$.
The integral quotient is the dual of the primitive exterior-square
lattice.  The natural map from the primitive lattice into that quotient
has cokernel $\mathbb Z/2$; these two integral representations are not
isomorphic.  Their rationalizations are isomorphic by that same map.
\end{theorem}
\begin{proof}
\begin{enumerate}[label=\arabic*.]
\item The two matrices obtained from their two-by-two minors are
\[
Q^{-1}\bigwedge^2(C_f^{-1})Q=
\begin{pmatrix}
1&-1&0&1&0&0\\0&1&-1&-1&1&0\\0&1&0&-1&0&0\\
0&2&-1&-1&0&1\\0&2&0&-1&0&0\\0&1&1&-1&0&0
\end{pmatrix},
\]
\[
Q^{-1}\bigwedge^2C_gQ=
\begin{pmatrix}
1&0&1&0&0&0\\0&0&-1&0&1&0\\0&0&-1&0&0&1\\
0&0&-1&0&0&0\\0&1&-1&0&0&0\\0&0&-1&1&0&0
\end{pmatrix}.
\]
Their lower five-by-five blocks are the asserted blocks, entry by entry.
Put $L=\bigwedge^2\mathbb Z^4$, let $(\ ,\ )$ be its unimodular wedge
pairing, and set $\omega=(1,1,1,1,1,1)^t$.
Then $(\omega,\omega)=2$ and the primitive sublattice
$P_{\mathbb Z}=\omega^\perp$ is the kernel of
$(1,-1,1,1,-1,1)$.  In the basis with columns
\[
\begin{pmatrix}
1&-1&-1&1&-1\\1&0&0&0&0\\0&1&0&0&0\\
0&0&1&0&0\\0&0&0&1&0\\0&0&0&0&1
\end{pmatrix},
\]
the inclusion followed by projection
$j:P_{\mathbb Z}\longrightarrow L/\mathbb Z\omega$ has matrix
\[
j=\begin{pmatrix}
0&1&1&-1&1\\-1&2&1&-1&1\\-1&1&2&-1&1\\
-1&1&1&0&1\\-1&1&1&-1&2
\end{pmatrix},\qquad \det j=2.
\]
Its Smith factors are $(1,1,1,1,2)$.  The unimodular pairing on $L$
identifies $L/\mathbb Z\omega$ with $P_{\mathbb Z}^\vee$;
under this identification $j$ is the pairing map.

\item We rule out every possible integral intertwiner.  Denote the lower blocks above by $A_5,B_5$.
The Krylov matrix $(e_3,A_5e_3,A_5^2e_3,A_5^3e_3,A_5^4e_3)$
has determinant $-1$, so $A_5$ is cyclic.  Any rational endomorphism
commuting with $A_5$ therefore has the form
$u_0I+u_1A_5+\cdots+u_4A_5^4$.
The entries $(1,2),(1,4),(2,4),(3,1)$ of its commutator with $B_5$
give, on $(u_1,u_2,u_3,u_4)^t$, the coefficient matrix
\[
\begin{pmatrix}
0&1&0&1\\1&-1&0&0\\1&-1&-1&1\\1&1&0&1
\end{pmatrix},\qquad \det=-1.
\]
Thus $u_1=u_2=u_3=u_4=0$, proving that the simultaneous rational
centralizer is exactly $\mathbb QI$.
The natural map $j$ above is equivariant and rationally invertible.
For any rational equivariant map $f:P_{\mathbb Q}\to
(L/\mathbb Z\omega)_{\mathbb Q}$, the composite $fj^{-1}$ lies
in this centralizer.  Consequently $f=cj$ for some $c\in\mathbb Q$.
An integral isomorphism would have determinant $\pm1$, whereas
$\det(cj)=2c^5$.  Its $2$-adic valuation would give
$1+5v_2(c)=0$, impossible because $v_2(c)\in\mathbb Z$.
Thus $j$ is the explicit rational isomorphism, and the determinant
valuation excludes an integral isomorphism.

\end{enumerate}
\end{proof}




\begin{theorem}[Integral marking of the pencil]\hypertarget{tag-02QY}{}\label{night:P1.7}
\StatementRecord{night:P1.7}{proved}{02QY}{\cref{maps:src:F.3.5}, \cref{night:P1.2}}
The integral $H^1$ monodromy of the pencil has a companion marking with
finite pair $(C_f^{-1},C_g)$, and cusp product $C_fC_g^{-1}$.
Its exterior-square quotient is consequently the integral quotient
representation appearing in \cref{night:P1.2}.
\end{theorem}
\begin{proof}
\begin{enumerate}[label=\arabic*.]
\item Let $L=H^1(C_\lambda,\mathbb Z)$ with its integral cup pairing,
and write $T_1,T_2$ for the finite generators of
\cref{maps:src:F.3.5}. Put $A=C_f$ and $B=C_g$.
The rational intertwiner proved there is a map
\[
 q:L_{\mathbb Q}\longrightarrow\mathbb Q^4,
 \qquad qT_1q^{-1}=A^{-1},\quad qT_2q^{-1}=B.
\]
Its relation to the Levelt presentation is
$(h_\infty,h_1,h_0)=(T_1^{-1},T_1T_2,T_2^{-1})$;
the ordered product is one. The known rigidity input and its explicit
rational change of basis are recorded in \cref{maps:src:F.3.5}.
The nodal peripheral is $c=(T_1T_2)^{-1}$. Its primitive
Picard--Lefschetz class $\delta\in L$ and coefficient one are the
geometric node input proved there: the nonseparating vanishing circle
has a transverse integral cycle meeting it once, and the annulus
$xy=t$ turns once. Thus $T_1T_2$ acts by
$v\mapsto v+\epsilon\delta\langle\delta,v\rangle$,
where $\epsilon\in\{1,-1\}$ records the fixed inverse and pairing
conventions.

\item A $B$-invariant alternating form has upper entries
$(u,v,u,u,v,u)$. Invariance under $A$, evaluated on the first and
fourth basis vectors, gives $v=-u$. The common invariant forms are
therefore exactly the multiples of
\[
 J=\begin{pmatrix}
 0&1&-1&1\\-1&0&1&-1\\1&-1&0&1\\-1&1&-1&0
 \end{pmatrix},\qquad
 d=\begin{pmatrix}1\\1\\1\\0\end{pmatrix}.
\]
The Pfaffian of $J$ is $1$, and direct multiplication gives
\[
 A^tJA=B^tJB=J,\qquad
 A^{-1}B-I=d e_4^t,\qquad d^tJ=e_4^t.
\]
Here $e_4$ is the fourth standard vector. Write the transported cup
form as $sJ$ and $q(\delta)=k d$, with $s,k\in\mathbb Q^*$.
Comparison with the Picard--Lefschetz formula gives
$\epsilon s k^2=1$.

\item The integral classes
$\delta_i=T_1^{-i}\delta$, $0\leq i\leq3$, have rational images
equal to the columns of $kD$, where
\[
 D=(d,Ad,A^2d,A^3d)=
 \begin{pmatrix}1&0&-1&0\\1&1&-1&-1\\1&1&0&-1\\0&1&0&0\end{pmatrix},
 \qquad \det D=1.
\]
Their pairing is $\epsilon G$, with
\[
 G=D^tJD=
 \begin{pmatrix}0&1&0&0\\-1&0&1&0\\0&-1&0&1\\0&0&-1&0\end{pmatrix},
 \qquad \det G=1.
\]
These four classes are a basis of $L_{\mathbb Q}$. For any $v\in L$,
write $v=\sum_i x_i\delta_i$ with $x_i\in\mathbb Q$.
The vector $(\langle\delta_i,v\rangle)_i=\epsilon Gx$ is integral.
Since $G^{-1}$ is integral, so is $x$. Thus the four classes are an
integral basis of $L$.

\item The required integral marking is the precise map
$\mu=k^{-1}q:L\longrightarrow\mathbb Z^4$.
Indeed $q(L)=kD\mathbb Z^4=k\mathbb Z^4$. Its inverse is specified
on standard vectors by
\[
 \mu^{-1}(e_j)=\sum_{i=0}^3(D^{-1})_{i+1,j}\delta_i,
 \qquad
 D^{-1}=\begin{pmatrix}1&-1&1&0\\0&0&0&1\\0&-1&1&0\\1&-1&0&1\end{pmatrix}.
\]
These two integral maps preserve the cup form
$\epsilon J$. In this marking the finite pair is $(A^{-1},B)$.

\item The exact simultaneous companion reversal is
\[
 W=\begin{pmatrix}0&0&0&1\\0&0&1&0\\0&1&0&0\\1&0&0&0\end{pmatrix}.
\]
It satisfies $W^2=I$, $\det W=1$, $WAW=A^{-1}$,
$WBW=B^{-1}$ and $W^tJW=-J$. These identities give the integral comparison of the two companion conventions, including
the pairing sign. For the based cusp used in the statement take
$c'=T_1^{-1}cT_1=(T_2T_1)^{-1}$. Then
\[
 \mu c\mu^{-1}=B^{-1}A,\qquad
 \mu c'\mu^{-1}=A(B^{-1}A)A^{-1}=AB^{-1}.
\]
The homology matrices are the inverse transposes under the evaluation pairing with $L$.

\item Apply $\bigwedge^2\mu$ to the integral wedge lattice. In the
lexicographic basis of \cref{night:P1.2}, the invariant vector is
$\omega=(1,1,1,1,1,1)^t$. Let $Q$ have first column $\omega$ and
remaining columns the last five standard vectors, as in that result.
The induced quotient is the map
\[
 [u\wedge v]\longmapsto
 \operatorname{pr}_{2,\ldots,6}
 Q^{-1}\bigl(\mu(u)\wedge\mu(v)\bigr).
\]
Its kernel before passage to the quotient is precisely the invariant
line. Its inverse is induced by $(\bigwedge^2\mu^{-1})Q$ on the last
five standard vectors. The finite quotient matrices are the lower
five-by-five blocks of
\[
 Q^{-1}(\bigwedge^2 A^{-1})Q,\qquad
 Q^{-1}(\bigwedge^2 B)Q.
\]
They are exactly the integral quotient matrices of
\cref{night:P1.2}, and the cusp follows by the same based word
$(T_2T_1)^{-1}$.
The integral saturation argument in steps 2--4 supplies the map $\mu$
for the geometric lattice.
\end{enumerate}
\end{proof}




\begin{theorem}[The rank-three bundle $P$]\hypertarget{tag-02QW}{}\label{night:P1.5}
\StatementRecord{night:P1.5}{proved}{02QW}{the definitions and constructions stated in the proof}
The wedge map from
$\ell_0\otimes(\mathcal H_{\mathrm{dR}}/\ell_0)$, followed by
projection modulo the nondegenerate polarization bivector, is
injective and has rank-three image. This image defines the bundle
$P$ used in the period construction.
Moreover
$P\cong\mathcal O\oplus T^2\oplus T^4$,
$H^1(P^\vee)=0$, and $H^0(P^\vee)$ is the one-dimensional
space of contractions. Its inclusion into the rank-five quotient
lifts holomorphically to the rank-six extension.
\end{theorem}
\begin{proof}
On the root curve of orders five and eight, let $\mathcal H_D$ be
the logarithmic de Rham bundle of the pencil. Its Hodge subbundle is
$\mathcal E=T^3\oplus T$, with $\ell=T^3$. The polarization identifies
$\mathcal H_D/\mathcal E$ with $\mathcal E^\vee$. This remains true
at the node: the dualizing differentials have independent limits,
and their isotropic pairing extends there holomorphically.
The primitive projection $p$ and the marked extension give
\begin{align}
 0&\longrightarrow\mathbb C e\longrightarrow\mathcal V
   \longrightarrow\mathcal W\longrightarrow0,
 \hypertarget{tag-02EK}{}\label{maps:eq:F.14}\\
 P&=p(\ell\wedge\mathcal H_D).
 \hypertarget{tag-02EL}{}\label{maps:eq:F.15}
\end{align}
In logarithmic frames the first sequence extends over the cusp,
since the nilpotent exponential commutes with its inclusion and
quotient. It extends equivariantly over the finite root charts.
Every nonzero vector in $\ell\wedge\mathcal H_D$ is decomposable, whereas the polarization bivector has nonzero square. Thus primitive projection is injective there. The Hodge sequence therefore gives
\begin{equation}
 0\longrightarrow T^4\longrightarrow P
   \longrightarrow\mathcal O\oplus T^2\longrightarrow0.
 \hypertarget{tag-02EM}{}\label{maps:eq:F.16}
\end{equation}
For the coarse projection $\rho$, invariant sections in the two
root charts give
\begin{equation}
 \rho_*T^j=\mathcal O_{\mathbb P^1}
 \left(\left\lfloor\frac{2j}{5}\right\rfloor+
       \left\lfloor-\frac{3j}{8}\right\rfloor\right).
 \hypertarget{tag-02EN}{}\label{maps:eq:F.17}
\end{equation}
Finite-group averaging makes descent exact. For $j=\pm2,\pm4$
the degree is $-1$, so both cohomology groups vanish. Consequently
$P=\mathcal O\oplus T^2\oplus T^4$ and
$H^1(P^\vee)=0$. The pullback of \eqref{maps:eq:F.14} to $P$
splits holomorphically. Its splittings differ by $z\alpha e$,
since $H^0(P^\vee)=\mathbb C\alpha$.
\end{proof}

\subsection{The algebraic connection and its splitting}

\begin{theorem}[The algebraic connection]\hypertarget{tag-02QO}{}\label{night:P1.15}
\StatementRecord{night:P1.15}{proved}{02QO}{\cref{night:P1.1}, \cref{night:P1.2}, \cref{night:P1.5}, \cref{night:P1.7}}
The rank-six flat extension, as a complex flat bundle with its invariant
line, has the following explicit algebraic connection model.  Set
$\omega_k=x^kdx/y$, $0\leq k\leq3$, and use the quotient frame
\[
(w_2,w_1,w_0,w_{-1},w_{-2})
 =(\omega_{01},\omega_{02},\omega_{03},\omega_{13},\omega_{23}),
\qquad \omega_{12}=-3\omega_{03}.
\]
These are meromorphic frames of
$(L^2,L,\mathcal O,L^{-1},L^{-2})$, where $L=T^2$ on the root
orbifold of orders $5,4$.  If a frame is written as a row of vectors,
its derivative is right multiplication by
\begin{equation}
\Gamma_6=\begin{pmatrix}0&b\\0&\Gamma_W\end{pmatrix},\qquad
b=\frac1{\lambda-\lambda_*}
 \left(-\frac{25}{6\lambda^2},-\frac{125}{48\lambda},0,1,0\right),
\label{pub:V:eq:1-15-1}
\end{equation}
where
\begin{equation}
\resizebox{\textwidth}{!}{$\displaystyle
\Gamma_W=\frac1{3125\lambda+256}
\begin{pmatrix}
-8(625\lambda+48)/\lambda&-8&10\lambda&75\lambda^2/2&0\\
-125/(2\lambda)&-4(625\lambda+48)/\lambda&8&0&75\lambda^2/2\\
300/\lambda^2&-375/\lambda&0&-48&60\lambda\\
120/\lambda^3&0&125/(2\lambda)&4(625\lambda+48)/\lambda&-8\\
0&120/\lambda^3&50/\lambda^2&-125/(2\lambda)&8(625\lambda+48)/\lambda
\end{pmatrix}.$}
\label{pub:V:eq:1-15-2}
\end{equation}
This model specifies a canonical algebraic splitting, denoted
$\sigma_b$.  The compatible contraction used for the real parameter is
$\alpha=(625\pi i/128)\alpha_{\rm alg}$, where
$\alpha_{\rm alg}(w_0)=4/3$ and it vanishes on $w_1,w_2$.
\end{theorem}
\begin{proof}
\begin{enumerate}[label=\arabic*.]
\item Reduction modulo exact differentials is completely explicit.  To
differentiate $\omega_k$, solve the polynomial identity
\[
-\tfrac12x^k\partial_\lambda f
 =f\sum_{j=0}^3G_{kj}x^j+R_k'f-\tfrac12R_kf',
 \qquad\deg R_k\leq4,
\]
where $f=x^5+\lambda^3x+\lambda^4$.  The last two terms are the
numerator of $d(R_k/y)$.  Matching the nine coefficients gives the
four-by-four Gauss--Manin matrix exactly.
At the branch point at infinity, a local coordinate $z=x^{-1/2}$
gives $\omega_k=-2z^{2-2k}(1+O(z^8))\,dz$.
Residues of a primitive times a differential give the alternating
pairings $J_{03}=4/3$, $J_{12}=4$, with their negatives in reverse
order and all other entries zero.  This form is nondegenerate.
The invariant bivector is therefore $3\omega_{03}+\omega_{12}$.
Taking the exterior-square connection and this quotient gives
\eqref{pub:V:eq:1-15-2}.

\item The root-chart powers computed from $x^kdx/y$ are
$u^{6k-9}$ at zero and $u^{6-4k}$ at infinity.
Their nondegenerate pairing shows that the four resulting regular
frames form the full de Rham bundle on those charts.  At the node
the computed connection has a simple pole and square-zero residue;
hence these are the nilpotent logarithmic extension frames there.
Thus the four line bundles are
$T^3,T,T^{-1},T^{-3}$, and their quotient exterior-square lines are
the stated powers of $L$.

\item For the coarse projection $\rho$ on the effective root curve,
\[
 \rho_*L^j=\mathcal O_{\mathbb P^1}
 \left(\left\lfloor\frac{4j}{5}\right\rfloor+
       \left\lfloor-\frac{3j}{4}\right\rfloor\right).
\]
The degree is zero for $j=0$ and minus one for $j=\pm1,\pm2$.
Finite stabilizer invariants are exact, so all five $H^1(L^j)$ vanish.  Consequently every holomorphic
extension of this $W$ by $\mathcal O$ splits as a bundle.
In such a splitting its flat connection is an upper block extension
by a row $b$.  Regularity at the finite root charts and a logarithmic
pole at the cusp require, entry by entry,
\[
b=\frac1{\lambda-\lambda_*}
 (a\lambda^{-2},b'\lambda^{-1},0,c,d\lambda).
\]
For example the $L^j$ entry, after $\lambda=u^4$, has exponent
$3j+3+4p$ for a term $\lambda^p/ (\lambda-\lambda_*)$;
the infinity chart requires $-4j-1-5p\geq0$.
These two inequalities leave exactly the four displayed terms.
The equation that the full cusp residue have square zero is
$b_{\rm res}(\Gamma_W)_{\rm res}=0$.
Its solution space has basis
\[
(-25/6,-125/48,1,0),\qquad(5/3,-25/3,0,1).
\]
The holomorphic splitting changes are the constant multiples of the
row $r_0=(0,0,1,0,0)$ selecting $w_0$.  Its connection coboundary
is the third row $r_0\Gamma_W$, whose parameter vector is
$(-48b_1+60b_2)/3125$.  The quotient by this exact gauge direction
is one-dimensional, with coordinate $s=c+4d/5$.  The first
vector has $s=1$ and rank-three square-zero cusp residue.

The extension is nontrivial, since its dual has no common
invariant vector by \cref{night:P1.1}; consequently its coordinate $s$ is nonzero.
Set $t=3125d/60$.  Its full row is
$\beta=s b+t r_0\Gamma_W$.  The constant bundle maps
\[
 C=\begin{pmatrix}1&-t r_0\\0&sI_5\end{pmatrix},\qquad
 C^{-1}=\begin{pmatrix}1&(t/s)r_0\\0&s^{-1}I_5\end{pmatrix}
\]
are holomorphic in the stated root frames and fix the invariant
section $e$.  Direct multiplication gives
\[
 C^{-1}\Gamma_6 C+C^{-1}dC
   =\begin{pmatrix}0&\beta\\0&\Gamma_W\end{pmatrix}.
\]
This proves the required equivalence of the complete flat bundles
with their fixed invariant line.  Evaluation at the marked basepoint
gives a common intertwiner for both specified meridians.  Thus its
existence follows from the exact connection maps.  

\item For the last normalization, the kernel of the quotient cusp residue
on $(w_2,w_1,w_0)$ is generated by
$v=(-4096/1171875,64/1875,1)^t$.
It satisfies $b_{\rm res}(v)=625/192$ and
$\alpha_{\rm alg}(v)=4/3$.
The logarithmic monodromy is $2\pi i$ times this residue.
We fix the contraction on this kernel line by the exact companion identity $N_V(\widetilde v)=\alpha(v)e$.
Since the space of global contraction forms is one-dimensional, the
ratio $(625/192)/(4/3)=625/256$ fixes its scalar everywhere and gives
the displayed $625\pi i/128$.

\end{enumerate}
\end{proof}

\begin{definition}[The normalized height function]
\hypertarget{tag-PV03}{}\label{def:hb}
\StatementRecord{def:hb}{proved}{PV03}{}
Let $e=e_0$ be the fixed real section of the rank-six system,
and let $\sigma_b:P\to V$ be its algebraic splitting.
The invariant bivector, primitive projection and contraction are
normalized in the companion marking as in \cref{night:P1.15}.
For each regular base point, including the finite root points,
let $n$ be the real vector in $P\cap\bar P$ with $\alpha(n)=1$.
We define $h_b$ by
\begin{equation}
 \sigma_b(n)-\overline{\sigma_b(n)}=2i h_b e.
 \hypertarget{tag-02ER}{}\label{maps:eq:F.20}
\end{equation}
\end{definition}
\begin{proof}[Well-definedness]
Let $u$ be a local nonzero section of $\ell_0$.
Its Hodge metric is positive, so $\operatorname{span}_{\mathbb C}\{u,\bar u\}$
is a nondegenerate symplectic plane. The inverse polarization
bivector $\omega$ has a nonzero component on the exterior square
of its symplectic orthogonal complement. Consequently an equality
$u\wedge v-\bar u\wedge w=t\omega$ forces $t=0$; with $t=0$
its common value belongs to $\mathbb C(u\wedge\bar u)$.
The primitive projection therefore gives
\begin{equation}
 P\cap\bar P=p(\ell_0\wedge\bar\ell_0).
 \hypertarget{tag-02EP}{}\label{maps:eq:F.19}
\end{equation}
We make the scalar in the primitive decomposition explicit.
In the algebraic de Rham frame put
$\omega_{\rm alg}=3\omega_{03}+\omega_{12}$ and let
$\alpha_{\rm alg}(v)$ be the contraction of its unique wedge
lift $j(v)$ with the residue pairing $J_{03}=4/3$, $J_{12}=4$. Then
$\operatorname{contr}_{J}(\omega_{\rm alg})=8$. The raw primitive
coefficient is therefore $\alpha_0=\alpha_{\rm alg}/8$, whereas
\cref{night:P1.15} fixes
\[
 \alpha=\frac{625\pi i}{128}\alpha_{\rm alg}
        =\frac{625\pi i}{16}\alpha_0,
 \qquad
 \omega_{\rm norm}=\frac{16}{625\pi i}\omega_{\rm alg}.
\]
For the primitive representative $v$ of a quotient vector, its
unique wedge lift is exactly
\begin{equation}
 j(v)=v+\alpha_0(v)\omega_{\rm alg}
     =v+\alpha(v)\omega_{\rm norm}.
 \hypertarget{tag-02EO}{}\label{maps:eq:F.18}
\end{equation}

Here $\omega_{\rm norm}$ is real in the marked flat system,
although the algebraic formula contains $i$. To see the precise
comparison, use the companion cup matrix $\varepsilon J_c$ of
\cref{night:P1.7}, with $\varepsilon\in\{1,-1\}$, and its
invariant bivector $\omega_c=(1,1,1,1,1,1)^t$.
Contraction with $J_c$ gives $\operatorname{contr}_{J_c}(\omega_c)=2$.
For the rank-four period comparison $\Pi$, the residue formula is
$\Pi^t(\varepsilon J_c)\Pi=2\pi iJ$.
Indeed, cut the curve along a symplectic polygon and small circles
around the poles; the boundary integral of a primitive times the
second differential is the alternating period pairing on the
polygon and $2\pi i$ times the sum of residues on the circles.
It follows by inversion of the two alternating forms that
$(\bigwedge^2\Pi)\omega_{\rm alg}=(8\pi i/\varepsilon)\omega_c$.

The scaled exterior-square map
$\Psi=(625\varepsilon/256)\bigwedge^2\Pi$ induces the
quotient comparison in the fixed rank-six extension. Its scalar is forced by the cusp contraction identity: on the
kernel vector of \cref{night:P1.15} the extension residue is
$625/192$, while $\alpha_{\rm alg}$ has value $4/3$; the
logarithmic monodromy contributes $2\pi i$. The ratio is
$625\pi i/128$, exactly the pullback of
$\operatorname{contr}_{J_c}$ under this scaled comparison.
The one-dimensional extension calculation in
\cref{night:P1.15} then lifts it uniquely to the normalized
rank-six comparison fixing $e$.
Consequently $\Psi(\omega_{\rm norm})=\omega_c/2$, and
the induced primitive comparison sends $\alpha$ to
$\operatorname{contr}_{J_c}$. Both are real in the companion
marking. For $r=ip(u\wedge\bar u)$, the value $\alpha(r)$
is real and nonzero by the positive polarization pairing.
Thus $n=r/\alpha(r)$ is the unique real vector specified above.
Multiplying $u$ by a nonzero complex number multiplies both
$r$ and $\alpha(r)$ by its squared modulus, and leaves $n$
unchanged.

The image of $\sigma_b(n)$ in the quotient is real. Its
difference from its conjugate is therefore an imaginary
multiple of the real vector $e$, which gives a unique real
$h_b$. Every marked monodromy is real and fixes $e$, so the
same equality defines one function on the base. The construction
is real analytic on each root chart; it is invariant under its
finite deck action and hence descends.

For its coordinate expression, let $H'=H\Gamma_6$, $H(1)=I$.
Take the positive zero meridian with stem from $1$ to $1/32$
and a counterclockwise circle of radius $1/32$, and the infinity
meridian $\lambda=\exp(-2\pi it)$, $0\le t\le1$.
Let $H_0,H_\infty^{\rm mon}$ be their based continuations.
The comparison matrix $X$ is specified by
\[
 XH_0=BX,\qquad XH_\infty^{\rm mon}=AX,\qquad Xe_0=e_0.
\]
Its existence is the marked flat-bundle identification in
\cref{night:P1.15}. The equations $RA=AR$, $RB=BR$ for the
explicit integral matrices force $R$ to be scalar. Thus any
two such comparisons differ by a scalar, and $Xe_0=e_0$ fixes
that scalar. This proves uniqueness with the based meridians
specified.

Put $Y=XH$ and $k=625\pi/96$. Since
$\alpha(w_0)=ki$ and $\alpha(w_1)=\alpha(w_2)=0$,
the vector $kn$ has coordinates
$(0,v_2,v_1,-i,0,0)^t$ in the quotient lift for a unique
pair $(v_2,v_1)$. Equivalently, the lower five coordinates
of its image under $Y$ are real. Equation
\eqref{maps:eq:F.20} gives
\[
 h_b(\lambda)=\frac{96}{625\pi}
 \operatorname{Im}\bigl(Y_{01}v_2+Y_{02}v_1-iY_{03}\bigr).
\]
This proves the existence, uniqueness and normalization of the
real system used in the height computation.
\end{proof}

\begin{remark}[Change of splitting]
\hypertarget{tag-PV04}{}\label{rem:V:splitting-change}
\StatementRecord{rem:V:splitting-change}{proved}{PV04}{}
Every holomorphic splitting of the pullback extension over $P$
has the form $\sigma_b-a\alpha e$, $a\in\mathbb C$, because
$H^0(P^\vee)=\mathbb C\alpha$. Its parameter and height are
related to the ones fixed here by the exact equalities
\[
 (\sigma_b-a\alpha e+\xi_{\rm alt}\alpha e)(P)
  =(\sigma_b+(\xi_{\rm alt}-a)\alpha e)(P),
 \qquad h_{\rm alt}=h_b-\operatorname{Im}a.
\]
Thus $\xi=\xi_{\rm alt}-a$. We use $\sigma_b$ and $\xi$
throughout.
\end{remark}

\subsection{The lifted periods and compactification}

\begin{theorem}[Lifted periods and their exact parameter domain]
\hypertarget{tag-02O8}{}\hypertarget{tag-02QZ}{}
\label{maps:src:F.3.6}\label{night:P1.8}
\StatementRecord{night:P1.8}{proved}{02QZ}{\cref{night:P1.5,night:P1.15,def:hb}}
Let $P$ be the rank-three wedge bundle of \cref{night:P1.5}, and use
its normalized contraction $\alpha$ and the algebraic splitting $\sigma_b$.
The planes
\begin{equation}
 F_\xi=(\sigma_b+\xi\alpha e)(P),\qquad \xi\in\mathbb C,
 \hypertarget{tag-02ES}{}\label{maps:eq:F.21}
\end{equation}
are holomorphic rank-three subbundles of the marked rank-six system.
They extend over its finite root charts of orders five and four.
The function $h_b$ of \cref{def:hb} has an attained finite minimum
$h_{b,\min}$, tends to $+\infty$ at the cusp, and satisfies
\begin{equation}
 F_\xi\cap\overline{F_\xi}=0
 \quad\Longleftrightarrow\quad
 h_b+\operatorname{Im}\xi\ne0
 \hypertarget{tag-02ET}{}\label{maps:eq:F.22}
\end{equation}
at each nonsingular base point, including the finite root points.
Thus the parameters transverse on the completed base are exactly
$\operatorname{Im}\xi>-h_{b,\min}$.
In an integral marking at the cusp their periods are
\begin{equation}
 Z(q)=\frac{\log q}{2\pi i}I_3+C_\xi(q),
 \qquad C_\xi\text{ holomorphic at }q=0.
 \hypertarget{tag-02EJ}{}\label{maps:eq:F.13}
\end{equation}
\end{theorem}
\begin{proof}
The holomorphic splitting and the one-dimensional space of its changes are given by \cref{night:P1.5}. We use the algebraic choice $\sigma_b$ of \cref{night:P1.15}.

The real-intersection calculation and well-definedness proof of
\cref{def:hb} give $P\cap\bar P=\mathbb Cn$ with
$\alpha(n)=1$ and
$\sigma_b(n)-\overline{\sigma_b(n)}=2ih_b e$.
The lifted vectors for parameter $\xi$ differ from their
conjugates by $2i(h_b+\operatorname{Im}\xi)e$.
An arbitrary complex multiple of $n$ gives no further
intersection direction. This proves \eqref{maps:eq:F.22}.

We calculate its cusp behavior in the fixed companion marking.
For a local nonzero $u\in\ell_0$, put $r=ip(u\wedge\bar u)$
and $a_r=\alpha(r)$, so that $n=r/a_r$.
Put $D=C_fC_g^{-1}-I$, $N_V=(BA)^{-1}-I$ and let $N_W$ be its
quotient. The only nonzero column of $D$ is
$(0,1,1,1)^t$, and the top row of $N_V$ is
$(0,0,0,-1,1,-1)$. The wedge projection and normalized contraction
satisfy, for all $a,b\in\mathbb C^4$,
\[
 (N_V)_{0,*}(0,p(a\wedge Db))
 =\operatorname{contr}(a\wedge Db)=a_0b_0,
 \qquad N_Wp(a\wedge Db)=0.
\]
These identities follow by expanding the two-by-two minors in
the companion basis. Write $u=\exp(sD)a(q)$ with
$s=\log q/(2\pi i)$ and $y=\operatorname{Im}s$.
The vector $a$ is holomorphic, and $a_0(0)\ne0$ because
$dx/y$ has nonzero nodal residue. In the logarithmic frame,
\[
 e^{-sN_W}r=ip(a\wedge\bar a)+2y\,p(a\wedge D\bar a),
 \qquad
 a_r=i\operatorname{contr}(a\wedge\bar a)+2y|a_0|^2.
\]
The splitting is bounded in that frame. On returning by
$e^{sN_V}$, its imaginary invariant-coordinate contribution is
$2|a_0|^2y^2+O(y)$. Division by
$a_r=2|a_0|^2y+O(1)$ yields
$h_b(q)=y+O(1)$, uniformly with $\operatorname{Re}s$ in a unit
interval. Thus $h_b(q)\to+\infty$ at the cusp.

The same identities give the logarithmic graph. At the limiting
plane choose a symplectic basis $a,b,d_1,d_2$ with $N_Hb=a$
and $\ell=\mathbb Cb$. Then
\begin{equation}
 P_0=\langle p(b\wedge a),b\wedge d_1,b\wedge d_2\rangle,
 \quad N_W(b\wedge d_i)=a\wedge d_i,
 \quad N_Wp(b\wedge a)=0.
 \hypertarget{tag-02EU}{}\label{maps:eq:F.23}
\end{equation}
The first generator maps under $N_V$ to the nonzero invariant
line, while the other two map to independent quotient directions.
Therefore $N_V:F_{\xi,0}\to\ker N_V$ is an isomorphism for
every $\xi$. In the primitive integral Jordan marking, this
makes the untwisted plane a graph with holomorphic coefficient
$C_\xi(q)$. Restoring $e^{sN_V}$ adds $sI_3$, giving
\eqref{maps:eq:F.13}.

The height extends continuously over both finite root charts.
Its cusp growth and compactness away from that cusp give an
attained minimum. The completed punctured base is connected, so
its image is $[h_{b,\min},\infty)$. Equation
\eqref{maps:eq:F.22} gives the exact parameter half-plane.
The fourth power of the order-eight action on $\mathcal H_D$
is $-I$; it acts trivially on $\mathcal W,P,\alpha$ and
$\mathcal V$. The equivariant lifted plane therefore descends
to the order-four chart. The other chart has order five.
\end{proof}












\begin{remark}[The parameter under base and fibre conjugation]
\hypertarget{tag-PV05}{}\label{rem:V:conjugation-parameter}
\StatementRecord{rem:V:conjugation-parameter}{proved}{PV05}{}
Put $k=625\pi/96$. In the normalized integral frame the lifted
period columns are
$W_\xi=(Y_{*,1},Y_{*,2},Y_{*,3}+ki\xi e_0)$.
For the based meridians fixed in \cref{def:hb}, complex conjugation
reverses each meridian. Since the connection is rational and
$H(1)=I$, its based continuation satisfies
$\overline{H_\gamma}=H_\gamma^{-1}$ for either meridian.
Consequently $K_0=X\bar X^{-1}$ satisfies
$K_0A=A^{-1}K_0$, $K_0B=B^{-1}K_0$ and $K_0e_0=e_0$.
Direct multiplication gives the same identities for
\[
 K=\begin{pmatrix}
 1&2&-1&-1&2&-1\\
 0&0&0&0&-1&1\\
 0&0&-1&0&0&1\\
 0&0&0&-1&0&1\\
 0&-1&0&0&0&1\\
 0&0&0&0&0&1
 \end{pmatrix},\qquad K^2=I.
\]
The common commutant of $A,B$ is scalar by the calculation in
\cref{def:hb}. Hence $K^{-1}K_0$ is scalar, and its value on
$e_0$ makes that scalar one. Thus $K_0=K$ is real.
Along conjugate continuation paths,
$Y(\bar\lambda)=K\overline{Y(\lambda)}$, so
\[
 \overline{W_{-\bar\xi}(\bar\lambda)}=K W_\xi(\lambda).
\]
This is the exact parameter map for simultaneous base and fibre
conjugation. It preserves the half-plane of \cref{night:P1.8}
and fixes every $\xi=ic$, $c\in\mathbb R$. The conjugated
Hodge plane, parametrized in this way, is holomorphic jointly
in $\lambda$ and $\xi$ because it equals the fixed linear
image $K F_\xi$.
\end{remark}

\begin{theorem}[Compact holomorphic filling]\hypertarget{tag-02O9}{}\label{maps:src:F.3.7}\hypertarget{tag-02QJ}{}\label{night:P1.10}
\StatementRecord{maps:src:F.3.7}{proved}{02O9}{\cref{night:P1.8,night:P1.9,pub:V:rank-six-data,pub:V:smooth-gluing,pub:V:cusp-smooth-dependence,rem:V:conjugation-parameter}}

In the notation of \cref{pub:V:parameter}, for every $\xi\in\mathcal U$
the marked gluing of \cref{pub:V:smooth-gluing} has a compact smooth
complex realization $M_\xi$ and a proper holomorphic map to $\mathbb P^1$.
This includes the subfamily $M_c$, $c>c_*$.
Its regular fibres are complex three-tori; its two finite reduced
fibres are the free affine quotients of orders five and four in
\cref{pub:V:rank-six-data}; its remaining fibre is the index-one
Kuhn degeneration. With its complex orientation its fundamental
group is trivial and its Euler number is six.

\end{theorem}


\begin{proof}
\begin{enumerate}[label=\arabic*.]
\item We fix the dual convention together with the parameter.
The preceding construction has cohomology monodromies $A,B$;
\cref{pub:V:rank-six-data} uses homology monodromies $A^t,B^t$.
For each admissible $\xi$ we start with the preceding member
at parameter $-\bar\xi$, pull it back by conjugation of the
base, and conjugate its complex fibres. In the fixed marking
the resulting Hodge plane is
\[
 \widehat F_\xi(\lambda)
 =\overline{F_{-\bar\xi}(\bar\lambda)}
 =K F_\xi(\lambda),
\]
by \cref{rem:V:conjugation-parameter}. Thus this family is
holomorphic in the coordinate $\xi$. The real based stems
of \cref{def:hb} are fixed by conjugation and their peripheral
meridians are reversed. Its cohomology monodromies are therefore
$A^{-1},B^{-1}$ and its homology monodromies are exactly
$A^t,B^t$.
Put $M_1=A^t$, $M_2=B^t$, $M_o=M_2M_1$, $M_0=(M_1M_2)^{-1}$ and
$J=\operatorname{diag}(I_3,-I_3)$.  The based peripheral word is
$\kappa(g_0)=g_1^{-1}g_0^{-1}g_1$, so
$M_0=M_1M_o^{-1}M_1^{-1}$.  For the target marking $S$,
direct integral multiplication gives
\[
 S^{-1}M_0S=\begin{pmatrix}I_3&I_3\\0&I_3\end{pmatrix},\qquad
 S_o=M_1^{-1}SJ,\qquad
 S_o^{-1}M_oS_o=\begin{pmatrix}I_3&I_3\\0&I_3\end{pmatrix}.
\]
Any two markings of this positive standard form differ by
$G=\left(\begin{smallmatrix}P&Q\\0&P\end{smallmatrix}\right)$,
where $P\in\operatorname{GL}_3(\mathbb Z)$ and $Q$ is integral:
this follows by multiplying $G$ with the displayed unipotent
matrix.  The period change is
\[
 (I_3,Z)G=(P,Q+ZP),\qquad
 P^{-1}(P,Q+ZP)
  =(I_3,sI_3+P^{-1}Q+P^{-1}CP).
\]
Thus the bounded holomorphic logarithmic form holds in $S_o$.
In the real-normalized cusp coordinate,
$q_o=\bar q$ and $s_o=-\bar s$.
For the member just defined, base and fibre conjugation give
$(I_3,-sI_3+C^\kappa_\xi(q))$, where
\[
 C^\kappa_\xi(q)
   =\overline{C_{-\bar\xi}(\bar q)}.
\]
The period-cycle change $J$ gives the target marking
$M_1S_oJ=S$ and the positive periods
$(I_3,sI_3-C^\kappa_\xi(q))$.


\item For a holomorphic Hodge subbundle $\mathcal F$ with
$\mathcal F\cap\bar{\mathcal F}=0$, restriction of integral
homology functionals embeds the homology lattice discretely in
$\mathcal F^*$.  Real transversality makes its real span the
whole real vector space.  The quotient
$\mathcal F^*/\mathcal V_{\mathbb Z}^*$ is consequently a
holomorphic family of compact complex three-tori over $U$,
with its zero section.  This also proves local properness,
by a smoothly varying fundamental parallelepiped.

\item On the root disc at the finite point of order $m_i$, the family
extends smoothly and its deck map has linear part $M_i$.
The invariant vector $v_i$ of \cref{pub:V:rank-six-data} gives a holomorphic
torsion section $v_i/m_i$.  Use the diagonal affine deck action
\begin{equation}
 (z,s)\longmapsto(M_i z+v_i/m_i,\zeta_{m_i}s).
 \hypertarget{tag-02EV}{}\label{maps:eq:F.24}
\end{equation}
Here $M_i z$ denotes the induced complex-linear map
$\mathcal F_s^*\to\mathcal F_{\zeta_{m_i}s}^*$ from the marked
lattice action, not multiplication of a complex three-vector
by a six-by-six matrix.  Likewise $v_i/m_i$ means its evaluation
as a point of the target torus.
It has order $m_i$, since the $m_i$th translation is integral.
Every nonidentity power is fixed-point-free on the central
torus: the invariant integral covector $\ell$ takes $v_i$ to
$1$ or $-1$, so a fixed point would require the noninteger
$\pm j/m_i$ to vanish in $\mathbb R/\mathbb Z$.
Away from the central fibre the base rotation is already free.
The quotient is therefore smooth and proper over the coarse
disc, with precisely the required multiple fibre.

\item Its marked punctured family is holomorphically identified with
the original torus family by logarithmic translation.  On the
root cover use $r(s)=(\log s/(2\pi i))v_i$.  Then
$r(\zeta_{m_i}s)-M_i r(s)=v_i/m_i$, modulo the integral lattice;
the change of branch of $\log s$ is itself an integral vector.
Thus translation by $r$ conjugates the linear deck action to
\eqref{maps:eq:F.24}.  This fixes the affine clutching, including its sign.

\item At the cusp use the integral marking of \eqref{maps:eq:F.13}, the fan \eqref{maps:eq:E.4}
and the action \eqref{maps:eq:E.5} with $C_3=C$.
Its maximal charts are $\mathbb C^4$ with map
$t=z_0z_1z_2z_3$, and its transition maps are the toric changes
of coordinates multiplied by the holomorphic units
$\exp(2\pi i C(t)\lambda)$.
The multiplicative comparison is
\[
 q=\bar q_o,\qquad x=(\bar x_o)^{-1},\qquad
 \nu=-\nu_o,\qquad
 C_{{\rm new},\xi}(q)=-C^\kappa_\xi(q).
\]
It carries the old period action to
$x\mapsto\exp(2\pi iC_{{\rm new},\xi}(q)\nu)q^\nu x$.
On cocharacters it is
$\Phi=\operatorname{diag}(-I_3,1)$ together with conjugation.
The Kuhn simplex at $v$ with ordering $\sigma$ is carried to
the simplex at $-v-(1,1,1)$ with reversed ordering.
If $N_o$ and $N$ are their height-one generator matrices and
$R$ reverses four columns, then
$N=\Phi N_oR$ and $N^{-1}\Phi=RN_o^{-1}$.
Consequently the full affine-chart map is
$z_j=\overline{z_{o,3-j}}$.  These character identities also
give equality on every monomial overlap; the product of the
four coordinates is the conjugate base coordinate.
The already displayed period-action equality makes the map
descend to the analytic period quotient.  Hence the fan and
marking in this comparison are precisely the prescribed ones.

\Cref{night:P1.9} applies to this bounded
holomorphic correction. It gives smoothness, proper discontinuity
and properness in the prescribed marking. We glue with zero cusp
translation.

\item The three punctured overlap maps are holomorphic isomorphisms
with the same marked torus family, so their cocycle condition
is literal.  Proper local models and a compact complementary
base give a Hausdorff compact total space.  Every chart is
smooth.  The affine vectors, cusp fan and clutching are exactly
those in \cref{pub:V:smooth-gluing}.  Varying their bounded corrections
and varying the fibre complex structures along a smooth path
does not change these marked smooth fillings: on the finite
root models the underlying real torus family is the fixed
local-system bundle; on the cusp use the marked, stratum-preserving comparison of
\cref{pub:V:cusp-smooth-dependence}.  Thus the resulting underlying smooth manifold
is the gluing of \cref{pub:V:smooth-gluing}. Its fundamental group is trivial and its Euler number is six by that lemma.

The finite actions and all gluing maps are holomorphic in $\xi$.
The cusp correction $C_{{\rm new},\xi}$ is jointly holomorphic;
on each relatively compact parameter neighbourhood the cusp disc
can be chosen uniformly by \cref{night:P1.9}. These constructions
therefore give a holomorphic family over the parameter half-plane.

\end{enumerate}
\end{proof}






\section{Characteristic-number conventions}\label{pub:V:sec-characteristic-conventions}
\begin{proposition}[The characteristic numbers and middle form]
\hypertarget{tag-034M}{}\label{sec:V-known-inputs}
\StatementRecord{sec:V-known-inputs}{proved}{034M}{}
Let $M$ be a closed connected almost complex manifold of real dimension
eight, with its complex orientation, with $H^2(M;\mathbb Q)=
H^6(M;\mathbb Q)=0$, $b_4(M)=4$, and $\chi(M)=6$.
Then its middle integral pairing modulo torsion is $I_{3,1}$ and
\begin{equation}
 H^4(M;\mathbb Z)/\operatorname{tors}\cong I_{3,1},\qquad
 c_4[M]=6,\qquad c_2^2[M]=2,\qquad
 \sigma(M)=2,\qquad\operatorname{Td}[M]=0.
 \hypertarget{tag-02EQ}{}\label{maps:eq:F.2}
\end{equation}
\end{proposition}
\begin{proof}
The classes $c_1,c_3$ vanish rationally. The index and signature
polynomials of \cite[equations~(1.2)--(1.4), pp.~3--4]{ffEsoleKang}
therefore reduce to
\begin{equation}
 \operatorname{Td}=\frac{3c_2^2-c_4}{720},\qquad
 \chi^1=\frac{3c_2^2-31c_4}{180},\qquad
 \sigma=\frac{3c_2^2+14c_4}{45}.
 \hypertarget{tag-02EF}{}\label{maps:eq:F.1}
\end{equation}
Here the first two numbers are elliptic-symbol indices of the
almost complex tangent bundle; their integrality follows from
\cite[Theorem~24.5.4 and its corollary, p.~184]{HirzebruchTopologicalMethods}.
Substitution of $c_4[M]=6$ gives
$\sigma(M)=16\operatorname{Td}[M]+2$.
The inequality $|\sigma(M)|\le b_4(M)=4$ forces
$\operatorname{Td}[M]=0$ and $\sigma(M)=2$, and then $c_2^2[M]=2$.
Poincar\'e duality makes the middle pairing modulo torsion
unimodular. Its signature excludes an even unimodular form,
whose signature is divisible by eight. The classification of
indefinite odd unimodular forms
\cite[Theorem~2.2.1, p.~61]{ffPetersSterk} gives $I_{3,1}$.
\end{proof}

\begin{remark}[The orientation in the connected-sum criterion]
\hypertarget{tag-PV06}{}\label{pub:V:mandal-convention}
\StatementRecord{pub:V:mandal-convention}{proved}{PV06}{}
We use the orientation in
\cite[Corollary~1.4, p.~4, and its proof in \S5]{ffMandal}:
$\ell\mathbb{HP}^2\#m(-\mathbb{HP}^2)$ admits an
orientation-compatible almost complex structure exactly for
$\ell=2m+1$ and odd $m$.
The tangent-class conventions are those of
\cite[Lemma~2.1, p.~5, and Theorem~2.4, p.~7]{ffMandal}.
\end{remark}

\section{The integral eight-manifold}\label{pub:V:sec-integral}






  






\begin{theorem}[Homology of the finite fillings]\hypertarget{tag-02QU}{}\label{night:P1.3}
\StatementRecord{night:P1.3}{proved}{02QU}{the definitions and constructions stated in the proof}
The two finite filling pieces have first integral homology
$\mathbb Z^2$ and $\mathbb Z^2\oplus\mathbb Z/2$.
\end{theorem}
\begin{proof}
\begin{enumerate}[label=\arabic*.]
\item For $K_i=T^6/\langle\tau_i\rangle$, the universal covering deck group
has commuting translations $x^u$, $u\in\Lambda$, and a generator $a_i$,
with the complete relations
\[
a_ix^ua_i^{-1}=x^{M_iu},\qquad a_i^{m_i}=x^{v_i}.
\]
Indeed the affine lifts satisfy these relations, and they reduce every
word uniquely to $x^ua_i^r$, $0\leq r<m_i$.
The abelianization relation matrix has columns
$((M_i-I)e_j,0)$ and $(-v_i,m_i)$.
Smith reduction of these explicit seven-by-seven matrices gives the
groups above.  The disk quotient
$(T^6\times\Delta)/\langle(\tau_i,\zeta_{m_i})\rangle$
retracts onto $K_i$ by equivariant radial contraction.  Freeness supplies
ordinary smooth quotient charts by disjoint translates of small balls.

\end{enumerate}
\end{proof}


\begin{lemma}[The integral restriction complex]
\hypertarget{tag-PV12}{}\label{pub:V:integral-restrictions}
\StatementRecord{pub:V:integral-restrictions}{proved}{PV12}{\cref{pub:V:smooth-gluing}, \cref{pub:V:rank-six-data}, \cref{night:P1.1}, \cref{night:P1.11}; \cite[Tag 09V5, Lemma 20.18.1]{StacksProject}}
Let $f:M\to S^2$ be the map of the marked smooth gluing in
\cref{pub:V:smooth-gluing}, and let $j:U\hookrightarrow S^2$
be the complement of its three special values $p_1,p_2,p_0$.
Put $\mathcal R_q=R^qf_*\mathbb Z$,
$\mathcal L_q=\mathcal R_q|_U$,
$\mathcal S_q=\operatorname{im}(\mathcal R_q\to j_*\mathcal L_q)$,
and $\mathcal K_q=\ker(\mathcal R_q\to j_*\mathcal L_q)$.
The fibres are connected, so $\mathcal R_0=\mathbb Z$.
For $1\leq q\leq4$ the image sheaf has
$H^2(S^2,\mathcal S_q)=\mathbb Z/2$; its only nonzero groups
in base degrees zero and one are
$H^0(S^2,\mathcal S_1)=\mathbb Z$ and
$H^1(S^2,\mathcal S_3)=\mathbb Z^3$.
Writing $i_a:\{p_a\}\hookrightarrow S^2$, its kernels are
\[
\mathcal K_1=0,\qquad
\mathcal K_2=\mathcal K_3=(i_2)_*(\mathbb Z/2),\qquad
\mathcal K_4=(i_2)_*(\mathbb Z/2)\oplus(i_0)_*\mathbb Z.
\]
For $q=2,3,4$, the torsion suspension generator at $p_2$ has
nonzero image in the global coinvariant group of
$\mathcal L_{q-1}$.
\end{lemma}
\begin{proof}
\begin{enumerate}[label=\arabic*.,leftmargin=*]
\item Use the marked smooth gluing of
\cref{pub:V:smooth-gluing}, the matrices and affine vectors of
\cref{pub:V:rank-six-data}, and the integral cusp-collapse maps
of \cref{night:P1.11}. The matrix $S$ is the explicit integral
marking in the first proof item of \cref{night:P1.1}.
Write
$m_1=5$, $m_2=4$, and put
\[
\begin{split}
C_1&=\begin{pmatrix}
-1&-1&0&-1&0\\0&1&0&1&1\\0&1&0&0&0\\
0&-1&-1&0&-1\\0&0&1&0&0
\end{pmatrix},\\[2pt]
C_2&=\begin{pmatrix}
-1&0&-1&0&-1\\0&1&0&0&0\\0&0&0&1&0\\
0&-1&0&-1&1\\0&0&1&0&0
\end{pmatrix}.
\end{split}
\]
For the translation vectors $v_i$, the matrices
$P_i=[v_i,(0;C_i)]$ are unimodular.  Direct multiplication gives
\[
 P_i^{-1}M_iP_i=1\oplus D_i,\qquad Q_i=P_i^{-t},
\]
where $D_1$ sends $e_j$ to $e_{j+1\bmod5}$, and $D_2$ sends
$e_j$ to $e_{j+1\bmod4}$ for $0\le j<4$ and $e_4$ to $-e_4$.
These equalities give the integral coordinate maps.  In these
coordinates the affine action is
\[
 (u,y,z)\longmapsto
 (u+1/m_i,D_i y,e^{2\pi i/m_i}z).
\]
Thus its core $F_i$ is the mapping torus of $D_i:T^5\to T^5$.
The function $W=e^{-2\pi i u}z$ trivializes its normal disc bundle.
With $\theta=m_i u\bmod1$, the coarse-base function on that bundle is
\[
                  q=W^{m_i}e^{2\pi i\theta}.
\]
In particular the boundary is $F_i\times S^1_W$ with this specified
map to the base circle.

\item We compute the integral restriction images on this model.
Let $V=\mathbb Z^5$ be the vertical cohomology lattice, let
$T_{i,k}=\bigwedge^kD_i^t$, and let
$N_{i,k}=\sum_{a=0}^{m_i-1}T_{i,k}^{a}$.
The signed permutations preserve the product cell structure with one
vertex and one edge on each circle.  Its cellular differential is zero.
Fibre cells and their prisms give the mapping-torus complex
\[
 C^q(F_i)=\bigwedge^qV\oplus\bigwedge^{q-1}V,
 \qquad d(a,b)=(0,(T_{i,q}-I)a).
\]
The $m_i$-fold covering by the generic six-torus pulls back
$(a,b)$ to $(a,N_{i,q-1}b)$: the horizontal cell is the union of the
$m_i$ consecutive prisms, with their transported coefficients.
Consequently its image in the marked stalk
$L_q=\bigwedge^q\mathbb Z^6$ of $\mathcal L_q$ is
\[
 J_{i,q}=\bigwedge^q Q_i
 \left(\operatorname{Inv}T_{i,q}\ \oplus\
 e_0\wedge N_{i,q-1}\bigwedge^{q-1}V\right).
\]
Here vertical coordinates occupy slots $1,\ldots,5$ in the six-lattice.
This formula specifies both the lattice and its inclusion.

We fix its column basis as follows.  Exterior coordinates are in
increasing lexicographic order, with indices starting at zero.
For each signed-permutation orbit, start at its least coordinate
and sum its successive signed images until it returns there.
A positive orbit of length $l$ gives a primitive invariant column;
the norm image has the column $m_i/l$ times that column.
A negative orbit has zero invariant lattice and one $\mathbb Z/2$
in its coinvariant group.  Indeed the successive relations identify
all its coordinates with signed copies of the first, and the last
relation is twice that first coordinate.  Order the positive orbits
by their least coordinates, with vertical columns before horizontal
columns.  For $D_1$ there is no torsion in vertical degrees $1$ through
$4$.  For $D_2$ there is one negative orbit in each such degree; the
positive-orbit counts are $1,2,2,1$.

\item At the cusp, \cref{night:P1.11} supplies a primitive collapse
image and free homology, with the only collapse cokernel equal to
$\mathbb Z$ in degree four.  Dualizing those integral maps therefore
gives
\[
 J_{0,q}=\operatorname{Inv}\bigl(\bigwedge^q M_0^t\bigr),
\]
with cohomological restriction kernel zero for $q=1,2,3$ and
$\mathbb Z$ for $q=4$.
Here are explicit primitive bases, which also fix the column order
in the calculation below.  In the cusp marking $S$, write
the cohomology basis as $x_0,x_1,x_2,y_0,y_1,y_2$; its action sends
$x_i$ to $x_i+y_i$ and fixes $y_i$.
In degree one take $(y_0,y_1,y_2)$.  In degree two take, in order,
\[
 y_i\wedge y_j\ (i<j),\qquad
 x_i\wedge y_i\ (0\le i\le2),\qquad
 x_i\wedge y_j+x_j\wedge y_i\ (i<j).
\]
In degree three take $y_0\wedge y_1\wedge y_2$, then
$x_i\wedge y_j\wedge y_k$ in lexicographic $(j,k)$ order with
$j<k$ and $i=j,k$, and then
\[
\begin{split}
 &x_0\wedge y_1\wedge y_2+x_1\wedge y_0\wedge y_2,\\
 &x_0\wedge y_1\wedge y_2-x_2\wedge y_0\wedge y_1.
\end{split}
\]
Their coordinate equations show that these bases are primitive and
have ranks $3,9,9$.  For degree four set
$J=\left(\begin{smallmatrix}0&I_3\\-I_3&0\end{smallmatrix}\right)$,
and let $H(e_I)$ be the complementary exterior basis vector, with
the sign of the permutation $(I,I^c)$.  Apply
$H\bigwedge^2J$ to the displayed degree-two basis.
The identity $JU=U^{-t}J$ gives the invariant equation, and the
signed-permutation maps $H$ and $\bigwedge^2J$ preserve primitivity.
Transport all these columns by $\bigwedge^qS^{-t}$ to obtain
$J_{0,q}$ in the original marking.  Thus its degree-four rank is nine.

\item Since $M$ is compact and the base is Hausdorff, the proper-base-change
map identifies its stalks with fibre cohomology
\cite[Tag 09V5, Lemma 20.18.1; chapter PDF p.~30]{StacksProject}.
Its three stalk images are exactly the lattices above.  The kernel
$\mathcal K_q$ is supported at those points. It is zero for $q=1$,
since the degree-zero norm on the finite suspension group is
injective and the cusp restriction is injective. For $q=2,3$ it is the
order-four $\mathbb Z/2$, and for $q=4$ it is that $\mathbb Z/2$
together with the cusp $\mathbb Z$.
Point-supported sheaves have no positive cohomology, so
$H^p(\mathcal R_q)=H^p(\mathcal S_q)$ for $p>0$.

Put $A_q=\bigwedge^qM_1^t$ and $B_q=\bigwedge^qM_2^t$.
The based marking gives
$\bigwedge^qM_0^t=(B_qA_q)^{-1}$.
Oriented cellular duality on the punctured sphere and its two-loop
spine gives the compact-support complex $L_q^2\to L_q$ in degrees
one and two.  Attaching the three stalk lattices gives the full complex
\[
\begin{gathered}
 J_{1,q}\oplus J_{2,q}\oplus J_{0,q}
 \xrightarrow{d_0}L_q^2\xrightarrow{\Delta_q}L_q,\\
 d_0(j_1,j_2,j_0)=(j_1-j_0,j_2-A_qj_0),\\
 \Delta_q(a,b)=(A_q-I)a+(B_q-I)b.
\end{gathered}
\]
For the cusp column, the two consecutive peripheral paths carry
$j_0$ and $A_qj_0$; their boundary is
$(B_qA_q-I)j_0=0$.  The reverse cusp orientation gives their negatives.
The two finite columns are the corresponding single peripheral
loops.  This is the relative cellular sequence for extension by zero
and the three stalk inclusions, so it preserves their integral indices.
It also verifies $\Delta_qd_0=0$ with the displayed transports.

We compute its cohomology by determinants in the ordered bases above.
Every matrix is specified by the marked $M_i,S,C_i$ and exterior minors.
For each $q$ select the following columns $\mathcal P_q$ of
$\Delta_q$ and call their square matrix $E_q$:
\[
\begin{aligned}
\mathcal P_1&=(1,2,5,8,9,10),\\
\mathcal P_2&=(0,1,2,8,9,11,12,13,16,17,23,24,25,26,28),\\
\mathcal P_3&=(1,2,3,5,6,10,11,12,13,15,16,18,19,\\
             &\hspace{22mm}26,29,31,32,34,35,37),\\
\mathcal P_4&=(0,1,5,7,9,10,11,13,16,17,20,23,26,28,29).
\end{aligned}
\]
Direct integer multiplication gives
\[
 \det E_q=(2,2,2,-2)_q,\qquad
 E_q^{-1}\Delta_q\in\operatorname{Mat}(\mathbb Z).
\]
Hence $E_q\mathbb Z^{\binom6q}$ is the entire image of
$\Delta_q$ and its cokernel is $\mathbb Z/2$.
Let $\mathcal F_q$ be the increasing complement of $\mathcal P_q$
in $\{0,\ldots,2\binom6q-1\}$.
These free-coordinate rows identify $\ker\Delta_q$ integrally:
the pivot coordinates of a kernel vector are
$-E_q^{-1}(\Delta_q)_{\mathcal F_q}$ times its free coordinates.
The relation matrix is therefore $R_q=(d_0)_{\mathcal F_q}$.
Its sizes, ranks and primitive minors are
\[
\begin{array}{c|c|c|l}
q&\text{size of }R_q&\operatorname{rank}R_q&
\text{maximal minors}\\ \hline
1&6\times7&6&-20,\ -5,\ 4\\
2&15\times15&15&\det R_2=1\\
3&20\times17&17&-1\\
4&15\times15&15&\det R_4=-1
\end{array}
\]
For $q=1$ the three column sets are respectively
$(0,1,2,3,5,6)$, $(0,1,2,4,5,6)$, $(0,2,3,4,5,6)$,
using all rows.  For $q=3$ use all columns and rows
$(0,1,\ldots,13,15,16,17)$.
These determinants have greatest common divisor one, so the
quotient by the relations is free.  The stalk column counts are
$(2,2,3)$, $(3,3,9)$, $(4,4,9)$ and $(3,3,9)$.
Thus the kernel dimensions and the quotient ranks give exactly
\[
\begin{array}{c|ccc}
q&H^0(\mathcal S_q)&H^1(\mathcal S_q)&H^2(\mathcal S_q)\\ \hline
1&\mathbb Z&0&\mathbb Z/2\\
2&0&0&\mathbb Z/2\\
3&0&\mathbb Z^3&\mathbb Z/2\\
4&0&0&\mathbb Z/2.
\end{array}
\]

\item For $k=1,2,3$ the negative signed orbits have representatives
$e_4$, $e_0\wedge e_2$, $e_0\wedge e_2\wedge e_4$.
Map them by the vertical inclusion followed by $\bigwedge^kQ_2$;
call the resulting columns $w_k\in L_k$.
The following rows $\rho_k$, read modulo two, annihilate
$\Delta_k$:
\[
\begin{aligned}
\rho_1&=(0,1,1,1,1,1),\\
\rho_2&=(1,1,1,1,1,1,1,0,0,1,0,1,0,1,1),\\
\rho_3&=(0,0,0,0,0,0,0,0,0,0,1,1,0,1,0,1,1,1,0,0).
\end{aligned}
\]
The explicit representatives and integer witnesses are
\[
\begin{aligned}
w_1&=(0,0,1,1,1,0)^t,\\
w_2&=(0,0,0,0,0,0,0,0,-1,0,0,-1,0,-1,-1)^t,\\
w_3&=(0,0,0,0,0,0,0,0,0,0,0,0,1,0,1,1,0,0,0,0)^t,\\
z_1&=(0,0,0,-1,-1,-1)^t,\\
z_2&=(-1,0,1,-2,-2,-2,4,2,-1,-3,2,0,-2,0,2)^t,\\
z_3&=(-2,-2,-3,0,0,-1,-1,-1,2,2,-1,3,-1,4,3,\\
   &\hspace{23mm}-3,-2,1,-4,1)^t.
\end{aligned}
\]
They satisfy $E_kz_k=2w_k$ and $\rho_kw_k=1\bmod2$.
The nonzero local torsion class therefore maps to the nonzero global
class in each degree $q=2,3,4$.

\end{enumerate}
\end{proof}

\begin{lemma}[The torsion transgression at the order-four filling]
\hypertarget{tag-02QF}{}\label{mapsaudit:order-four-torsion-transgression}
\StatementRecord{mapsaudit:order-four-torsion-transgression}{proved}{02QF}{\cref{pub:V:rank-six-data}, \cref{pub:V:integral-restrictions}}

Use the order-four filling and the marked matrices of \cref{pub:V:rank-six-data}.
Let $D$ be the bottom-right $5\times5$ block of $M_2$, and put
$T=\bigwedge^{q-1}D^t$ on the integral vertical cochains of the
five-torus. If $w$ represents a torsion element of
$\operatorname{coker}(T-I)$, its point-supported degree-$q$ class has
Leray differential equal, up to the one boundary-orientation sign, to
the image of $w$ in the global degree-$(q-1)$ coinvariants.  For each of $q=2,3,4$, the unique nonzero local torsion class maps
to the nonzero class in the global $\mathbb Z/2$ quotient.
\end{lemma}


\begin{proof}
\begin{enumerate}[label=\arabic*.]
\item Use the explicit matrix $C_2$ in the proof of
\cref{pub:V:integral-restrictions}.  In the original vertical marking $D$ is the
bottom-right block of $M_2$; multiplication gives
$DC_2=C_2D_2$, where $D_2$ is the four-cycle followed by the sign
on its fifth coordinate.  The change of cochain coordinates is
$\bigwedge^kC_2^t$, and the inverse is $\bigwedge^kC_2^{-t}$.
Thus the product cell structure of the signed-permutation torus,
transported by this integral torus isomorphism, has zero cellular
differential and the stated action $T=\bigwedge^{q-1}D^t$.
The fibre cells and their prisms give the mapping-torus
cochain complex.  The degree-$q$ suspension classes are consequently
the coinvariants of $T$ on vertical degree-$q-1$ cochains.

\item Write $u=-x_0$ in the original marking.  The affine action is
$(u,y,z)\mapsto(u+1/4,Dy,iz)$, and
$W=e^{-2\pi i u}z$ is invariant.  With core angle $a=4u\bmod1$
and normal angle $k=\arg W/(2\pi)$, the coarse angle is $a+4k$.
Let $g,h$ denote the two corresponding loops.  The integral basis
\[
 v=g^4h^{-1},\qquad t=g,
 \qquad [v,t]_{[g,h]}=\begin{pmatrix}4&1\\-1&0\end{pmatrix}
\]
has determinant one.  Its coarse degrees are zero and one; its
vertical monodromies are $T^4=I$ and $T$ respectively.
It specifies the generic six-torus fibre and the primitive
base loop in the boundary product.

\item A suspension cochain representing $[w]$ has values $(w,0)$
on $(g,h)$.  Following the four consecutive core prisms transports
their values by $I,T,T^2,T^3$, giving
\begin{equation}
 (w,0)_{(g,h)}\longmapsto(N_Tw,w)_{(v,t)},
 \qquad N_T=I+T+T^2+T^3.
 \hypertarget{tag-02J0}{}\label{maps:eq:FT.1}
\end{equation}
For a torsion coinvariant there are an integer $n\ne0$ and an
integral cochain $b$ with $nw=(T-I)b$.  Therefore
$nN_Tw=N_T(T-I)b=0$, and $N_Tw=0$ in the free cochain lattice.
The fibre restriction vanishes; its base-circle component is $[w]$.
Replacing $w$ by $w+(T-I)b$ changes that component by its cellular
coboundary.  The formula is consequently a well-defined map of the
stated integral coinvariant groups.

\item The base cover by a disc at the finite point and its
complement gives the relative \v{C}ech connecting map.  On their
annular intersection the degree-zero fibre restriction is zero and
the next component is the one-cocycle just computed.
The connecting map sends it to $\pm[w]$ in the global coinvariants;
the two ordinary peripheral loops impose the relations
$\bigwedge^{q-1}M_1^t-I$ and $\bigwedge^{q-1}M_2^t-I$.
These are the two successive boundaries in the same base--fibre
complex, so their composition is the Leray $d_2$ for $f$.  The one boundary-orientation choice accounts
for the sign.  Its coefficient is one, as witnessed by the
degree-one loop $t$ above.  Point-supported kernels or cokernels in
the comparison of restriction sheaves have no positive cohomology,
so they preserve this degree-two base quotient.

\item In signed-permutation coordinates the negative orbits in
vertical degrees $1,2,3$ are represented respectively by
$e_4$, $e_0\wedge e_2$ and $e_0\wedge e_2\wedge e_4$.
The orbit relations give exactly one local $\mathbb Z/2$ in each
case.  Transport by $\bigwedge^kC_2^{-t}$ and then include vertically
in the six-lattice.  This is the same map as the vertical inclusion
followed by $\bigwedge^kQ_2$ used in
\cref{pub:V:integral-restrictions}, because $P_2=(-1)\oplus C_2$.
The explicit $w_k,\rho_k,E_k,z_k$ there satisfy
\[
 \rho_k\Delta_k=0\pmod2,\qquad
 \rho_kw_k=1\pmod2,\qquad E_kz_k=2w_k,
 \qquad |\det E_k|=2.
\]
Thus the unique nonzero local class maps to the unique nonzero
global class in each of the three degrees.  Formula
\eqref{maps:eq:FT.1} sends the zero representative to zero as well.
\end{enumerate}
\end{proof}




\begin{theorem}[Integral homology and the middle form]
\hypertarget{tag-02NW}{}\label{maps:src:F.2.5}
\hypertarget{tag-02NX}{}\label{maps:src:F.2.6}
\hypertarget{tag-02QL}{}\label{night:P1.12}
\StatementRecord{maps:src:F.2.6}{proved}{02NX}{\cref{pub:V:smooth-gluing}, \cref{pub:V:integral-restrictions}, \cref{mapsaudit:order-four-torsion-transgression}, \cref{sec:V-known-inputs}, \cref{maps:src:F.3.7}}
The marked smooth gluing $M$ of \cref{pub:V:smooth-gluing} is
three-connected and has torsion-free integral homology with Betti numbers
\begin{equation}
 (1,0,0,0,4,0,0,0,1).
 \hypertarget{tag-02F9}{}\label{maps:eq:F.5}
\end{equation}
It admits an almost complex structure, has signature two, and its
middle integral intersection form is
$3\langle1\rangle\oplus\langle-1\rangle$.
The compact holomorphic realization of \cref{maps:src:F.3.7}
has this same marked smooth type.
\end{theorem}
\begin{proof}
The marked identification in \cref{maps:src:F.3.7} compares its
finite affine actions, cusp fan and overlap maps with the gluing
of \cref{pub:V:smooth-gluing}. The bounded-correction comparison
\cref{pub:V:cusp-smooth-dependence} preserves its collars.
We compute the integral groups of that gluing.
By \cref{pub:V:integral-restrictions} and
\cref{mapsaudit:order-four-torsion-transgression}, the point-supported
torsion classes have differential
\begin{equation}
 d_2[w]=\pm[w]\in L_{q-1}/\operatorname{im}\Delta_{q-1}.
 \hypertarget{tag-02FA}{}\label{maps:eq:F.6}
\end{equation}
The three images are the nonzero global classes, as shown by the
integer witnesses in \cref{pub:V:integral-restrictions}.
Since $H^0(\mathcal S_q)=0$ in these three degrees, the
point-supported kernels give
\[
 H^0(\mathcal R_2)=\mathbb Z/2,\quad
 H^0(\mathcal R_3)=\mathbb Z/2,\quad
 H^0(\mathcal R_4)=\mathbb Z\oplus\mathbb Z/2.
\]
The Leray sequence
\cite[Tag 01F2, Lemma 20.13.4; chapter PDF p.~22]{StacksProject}
has only base columns zero, one and two.
The maps from $(0,2)$ to $(2,1)$ and from $(0,3)$ to $(2,2)$ are
the isomorphisms just proved.  The restriction of the map from
$(0,4)$ to $(2,3)$ to its torsion summand is also an isomorphism.
Writing the image of a free cusp generator as $\epsilon\in\mathbb Z/2$,
its kernel is the free group generated by $(1,\epsilon)$.
All total-degree-three terms consequently vanish.  In total degree
four the remaining terms are $\mathbb Z$ and $\mathbb Z^3$.
Their extension is free, so
\[
                   H^3(M;\mathbb Z)=0,\qquad
                   H^4(M;\mathbb Z)=\mathbb Z^4.
\]
The degree-one row of the same Leray sequence is
\[
 H^0(\mathcal R_1)=\mathbb Z
 \xrightarrow{d_2^{0,1}}H^2(\mathcal R_0)=\mathbb Z.
\]
Its kernel is $H^1(M;\mathbb Z)=0$ by
\cref{pub:V:smooth-gluing}. Its degree is nonzero, so it becomes
an isomorphism over $\mathbb Q$. The image-sheaf groups in
\cref{pub:V:integral-restrictions} give
$H^2(M;\mathbb Q)=H^3(M;\mathbb Q)=0$.
Universal coefficients
\cite[Theorem~3.2 and Corollary~3.3, pp.~195--196]{Hatcher2002}
identify the torsion of $H_2$ and $H_3$ with the corresponding
$\operatorname{Ext}$ subgroups of $H^3$ and $H^4$; both are zero.
Their rational ranks also vanish, so $H_2=H_3=0$.
Integral Poincar\'e duality
\cite[Theorem~3.30, p.~241]{Hatcher2002} gives the remaining groups.
The successive Hurewicz homomorphisms
\cite[Theorem~4.32, p.~366]{Hatcher2002} give $\pi_2=\pi_3=0$.
Finally the almost complex structure and Euler number in
\cref{pub:V:smooth-gluing} satisfy
\cref{sec:V-known-inputs}, which gives the stated signature and
middle form.
\end{proof}






















\section{The oriented diffeomorphism type}\label{pub:V:sec-diffeomorphism}

\begin{theorem}[Characteristic numbers and the middle lattice]\hypertarget{tag-02QM}{}\label{night:P1.13}
\StatementRecord{night:P1.13}{proved}{02QM}{\cref{night:P1.12}, \cref{sec:V-known-inputs}}
With the complex orientation the fourfold has
\[
\sigma(M_c)=2,\qquad H^4(M_c;\mathbb Z)\cong I_{3,1},\qquad
w_2(M_c)=0,
\]
and there is a diagonal integral basis in which
$p_1/2=(1,1,1,1)$.
Its Chern numbers are
$c_4=6$, $c_2^2=2$, and
$c_1^4=c_1^2c_2=c_1c_3=0$; its Pontryagin numbers are
$p_1^2=8$, $p_2=14$.
\end{theorem}
\begin{proof}
\begin{enumerate}[label=\arabic*.]
\item By \cref{night:P1.12}, the integral groups in degrees two
and six vanish, so $c_1=c_3=0$.
\Cref{sec:V-known-inputs} gives
$c_4[M_c]=6$, $c_2^2[M_c]=2$, $\sigma(M_c)=2$, and the
integral middle form $I_{3,1}$.
The Chern-root identities
$p_1=-2c_2$ and $p_2=c_2^2+2c_4$ then give
$p_1^2[M_c]=8$ and $p_2[M_c]=14$.

\item Since $H^2(M_c;\mathbb Z/2)=0$, its $w_2$ is zero.
For the complex tangent bundle,
$w_4=c_2\bmod2=(p_1/2)\bmod2$.
The Wu identity in degree four, with $w_1=w_2=0$, says that $w_4$ is
characteristic for the middle pairing.  Thus the vector $p_1/2$ is
characteristic and has square $2$.

\item We give the required integral isometry of this pair explicitly as
a terminating procedure.  In a diagonal basis a characteristic vector
has four odd coordinates.  Change their signs to write it as
$(a,b,c,d)$ with all positive and
$a^2+b^2+c^2=d^2+2$.
For $d>1$,
\[
d<a+b+c<3d.
\]
The first inequality follows by squaring and using the positive cross
terms; the second follows from Cauchy--Schwarz and $d>1$.
Reflection in the integral norm-two vector $(1,1,1,1)$ sends this
vector to itself minus $(a+b+c-d)(1,1,1,1)$.
Its fourth coordinate is $2d-a-b-c$, whose absolute value is strictly
smaller than $d$.  The reflection preserves the integral lattice and
characteristic parity.  Repeating, with sign changes at each step,
terminates at $d=1$, when the norm equation forces $a=b=c=1$.
This gives an integral isometry taking $p_1/2$ to $(1,1,1,1)$.

\end{enumerate}
\end{proof}




\begin{theorem}[A complex structure on the stated connected sum]\hypertarget{tag-02OC}{}\label{maps:src:F.4.1}
\StatementRecord{maps:src:F.4.1}{proved}{02OC}{\cref{maps:src:F.3.7}, \cref{maps:src:F.2.6}, \cref{sec:V-known-inputs}, \cref{night:P1.13}; \cite[Corollary 1.4, p.~4]{ffMandal}; \cite[Definition 3.6 and Theorem 3.7(1), pp.~12--14]{IVVWallForm}; \cite[Theorem~A, p.~178]{IVVKramerStolz}}


The compact complex fourfold of Theorem \ref{maps:src:F.3.7} is three-connected,
has Euler number six, and is orientation-preservingly diffeomorphic to
\begin{equation}
 3\mathbb H\mathbb P^2\ \#\ (-\mathbb H\mathbb P^2).
 \hypertarget{tag-02EW}{}\label{maps:eq:F.25}
\end{equation}
Its integral middle lattice is $I_{3,1}$ and, in a diagonal integral
basis, $p_1/2=(1,1,1,1)$.  

\end{theorem}


\begin{proof}
\begin{enumerate}[label=\arabic*.]
\item \Cref{maps:src:F.2.6} gives three-connectivity,
Euler number six, and free homology of ranks $1,4,1$ in degrees
$0,4,8$. The complex orientation is the one in
\cref{maps:src:F.3.7}. By \cref{night:P1.13}, its middle pairing
is $Q=I_{3,1}$ and $p=p_1/2=-c_2$ is characteristic with square two.

\item Use the integral isometry supplied by the terminating descent
in \cref{night:P1.13}. In diagonal coordinates it is a matrix $U$
with $Up=(1,1,1,1)^t$. We use this specific isometry in the
comparison of normal data below.

\item We compute the complete normal clutching datum in fixed
coordinates. Orient $\mathbb H$ by $(1,i,j,k)$ and use
\[
\operatorname{Spin}(4)=S^3_L\times S^3_R,
\qquad (u,w)h=uhw^{-1}.
\]
The covering induces an isomorphism on $\pi_3$. Let $(a,b)$ be
the degrees of the two clutching maps. Evaluation on $1$ gives
$uw^{-1}$, of degree $a-b$, so $e=a-b$, with the sphere orientation
fixed by the left degree-one Euler generator. The left action is
complex linear for right multiplication by $i$, whose real complex
basis $(1,i,j,-k)$ has the opposite quaternion orientation. Its
degree-one generator therefore has $c_2=-1$, and hence $p_1=2$.
The right inverse action is complex linear for left multiplication
by $i$, with real basis $(1,i,j,k)$. Its Euler number and $c_2$ are
both $-1$, so its $p_1$ is again $2$. The determinant-one complex
actions have $c_1=0$. Consequently the exact homomorphism and its
inverse on its image are
\[
(a,b)\longmapsto(e,p_1/2)=(a-b,a+b),
\qquad
(a,b)=\left(\frac{e+p_1/2}{2},\frac{p_1/2-e}{2}\right).
\]
The image consists of pairs of the same parity, and the map is
injective. The tangent clutching class of $S^4$ has $(e,p_1)=(2,0)$,
hence coordinates $(1,-1)$. The homotopy sequence of
$SO(4)\to SO(5)\to S^4$ identifies its span with the kernel of
stabilization and gives surjectivity on $\pi_3$. Subsequent
stabilizations are isomorphisms there, since
$\pi_3(S^n)=\pi_4(S^n)=0$ for $n\geq5$.
The stable coordinate is therefore $a+b=p_1/2$, with the left
generator sent to $1$.

\item For $x\in H_4(M;\mathbb Z)$, use the embedded sphere
representative in the setup of
\cite[\S3.3, pp.~11--12]{IVVWallForm}. Orient its normal bundle
$\nu_x$ by the ambient orientation. Self-intersection and the
tangent splitting give
\[
e(\nu_x)=x\cdot x,\qquad
\langle p_1(\nu_x),[S^4]\rangle=\langle p_1(TM),x\rangle,
\]
because $TS^4$ is stably trivial. Thus its two clutching coordinates
are exactly
\[
\left(\frac{x\cdot x+\langle p,x\rangle}{2},
      \frac{\langle p,x\rangle-x\cdot x}{2}\right),
\qquad S\alpha_M(x)=\langle p,x\rangle.
\]
Their integrality is the characteristic congruence in step~1.
This is the stable normal homomorphism of
\cite[Definition~3.6, p.~12]{IVVWallForm} with its generator.
For an oriented middle-cohomology isometry $\psi$ from $M$ to a
comparison manifold $N$, the corresponding homology map is
\[
\Phi=\operatorname{PD}_N\circ\psi\circ\operatorname{PD}_M^{-1}.
\]
When $\psi(p_M)=p_N$, the cup-product evaluation identity gives
$\langle p_N,\Phi(x)\rangle=\langle p_M,x\rangle$.
The map $\Phi$ preserves the intersection pairing and both displayed
normal coordinates, so it intertwines the complete extended forms.

\item Take $N=3\mathbb H\mathbb P^2\#(-\mathbb H\mathbb P^2)$.
On each summand choose the generator $u$ with $p_1=2u$, as in
\cite[Lemma~2.1, p.~5]{ffMandal}. Reversing orientation changes
the evaluation of $u^2$ and leaves the tangent Pontryagin class
unchanged. The collapse maps of the oriented connected sum give
\[
H^4(N;\mathbb Z)=\bigoplus_{i=1}^4\mathbb Z u_i,
\qquad Q_N=I_{3,1},\qquad p_N=u_1+u_2+u_3+u_4;
\]
the tangent formula is
\cite[Theorem~2.4, p.~7]{ffMandal}.
Van Kampen, the connected-sum homology sequence and Hurewicz show
that $N$ is three-connected. The isometry $U$ of step~2 is an
isometry $\psi:H^4(M;\mathbb Z)\to H^4(N;\mathbb Z)$ carrying
$p_M$ to $p_N$. Step~4 supplies the isomorphism of their smooth
extended forms. Wall's classification in
\cite[Theorem~3.7(1), pp.~13--14]{IVVWallForm} now gives an
orientation-preserving diffeomorphism
$M\cong N\#\Sigma$ for an oriented homotopy eight-sphere $\Sigma$.

\item The theorem of Kramer and Stolz
\cite[Theorem~A, p.~178]{IVVKramerStolz} applies to
$\mathbb H\mathbb P^2$: it is smooth, simply connected, has dimension
eight, and has total integral homology $\mathbb Z\oplus\mathbb Z
\oplus\mathbb Z$. It gives
$\mathbb H\mathbb P^2\#\Sigma\cong\mathbb H\mathbb P^2$.
This diffeomorphism preserves orientation: its map on $H^4$ sends
a generator to that generator or its negative, whose square still
evaluates to $1$, forcing its degree to be $+1$.
The oriented associativity and permutation maps for connected sum
place $\Sigma$ beside one positive summand. Choose the punctured
balls compatibly with that diffeomorphism and use the identity on
the remaining punctured summands. These maps give
\[
\begin{aligned}
N\#\Sigma
&\cong(\mathbb H\mathbb P^2\#\Sigma)\#\mathbb H\mathbb P^2
       \#\mathbb H\mathbb P^2\#(-\mathbb H\mathbb P^2)\\
&\cong N.
\end{aligned}
\]
Composing with step~5 supplies the required oriented smooth
diffeomorphism. Transporting the integrable structure of
\cref{maps:src:F.3.7} along it gives the asserted complex structure.
\end{enumerate}
\end{proof}


\section{The normalized connection and the exact parameter domain}\label{pub:V:sec-parameter}















\begin{theorem}[The minimum of the height function]
\hypertarget{tag-02QR}{}\label{night:P1.18}
\StatementRecord{night:P1.18}{proved}{02QR}{\cref{night:P1.15,def:hb,night:P1.8,rem:V:conjugation-parameter}}
The function $h_b$ has its unique minimum at $\lambda=\infty$.
For the normalized fundamental solution $H(1)=I$ of the algebraic
connection and its regular infinity limit $H_\infty$, this minimum is
\[
 h_{b,\min}=-\frac{96}{625\pi}\operatorname{Re}(XH_\infty)_{03}.
\]
Its negative $c_*=-h_{b,\min}$ satisfies the exact bounds
\[
 \begin{gathered}
1.094540250311402780116268263054945703725<c_*,\\
c_*<1.094540250311402780116268263054945703726.
\end{gathered}
 \hypertarget{tag-03KZ}{}\label{layout:V:fresh-bounds}
\]
The full parameter domain is $\mathcal U=\{\xi:\operatorname{Im}\xi>c_*\}$.
\end{theorem}
\begin{proof}
Put $a=256/3125$ and $c_0=96/(625\pi)$. We use the rational
connection of \cref{night:P1.15}, its integral matrices $A,B$,
and the normalization $Y=XH$, $H(1)=I$, $Xe_0=e_0$.
The finite-root gauges are
\[
 \lambda=u^4,\quad D_0=\operatorname{diag}(1,u^6,u^3,1,u^{-3},u^{-6}),
 \qquad
 \lambda=u^{-5},\quad D_\infty=
 \operatorname{diag}(1,u^{-8},u^{-4},1,u^4,u^8).
 \hypertarget{tag-03KX}{}\label{layout:V:fresh-gauges}
\]
For either gauge the regular connection is
$D^{-1}\lambda'\Gamma_6D+D^{-1}D'$.
Let $U_0,U_\infty$ be its solutions with value $I$ at zero.
The deck characters $T_0,T_\infty$ give the exact based
monodromies $U_0(1)^{-1}T_0^{-1}U_0(1)$ and
$U_\infty(1)^{-1}T_\infty^{-1}U_\infty(1)$.
The unique common intertwiner fixing $e_0$ is $X$, by
\cref{night:P1.15,def:hb}. In these gauges the exact identities are
\[
 H(u^4)D_0(u)=U_0(1)^{-1}U_0(u),\qquad
 H(u^{-5})D_\infty(u)=U_\infty(1)^{-1}U_\infty(u).
\]
Thus the regular infinity limit in the statement means precisely
\[
 H_\infty=\lim_{u\to0}H(u^{-5})D_\infty(u)
             =U_\infty(1)^{-1}.
\]
The regular limiting frames are consequently
$F_0=XU_0(1)^{-1}$ and $F_\infty=XU_\infty(1)^{-1}$.

The height is obtained from a real linear system. Write
$q=(0,v_2,v_1,-i,0,0)^t$ and impose that the lower five
coordinates of $Yq$ be real. The four real unknowns are the
real and imaginary parts of $v_2,v_1$; existence and uniqueness
are those of \cref{def:hb}. In these coordinates
$h_b=c_0\operatorname{Im}(Yq)_0$.
Every monodromy is real and fixes $e_0$, so the imaginary top
coordinate is unchanged when the lower coordinates are real.
This proves that the calculation gives one function on the base.
The same height is preserved by every constant real invertible
change of frame $T$ fixing $e_0$. Indeed
$T=\left(\begin{smallmatrix}1&r\\0&S\end{smallmatrix}\right)$
with $S$ real and invertible: the lower coordinates of $Yq$
are real exactly when those of $TYq$ are real, and the top
coordinates differ by a real number. The differential equation
of $TY$ is still $(TY)'=(TY)G$. Such changes of frame may
therefore be used in the rectangle enclosures below.
The real normalized conjugation matrix $K=X\bar X^{-1}$ satisfies
$Y(\bar\lambda)=K\overline{Y(\lambda)}$ and $Ke_0=e_0$;
using $-\bar q$ gives $h_b(\bar\lambda)=h_b(\lambda)$.

We give the enclosure method and the resulting whole-domain
inequalities. For an analytic right connection $G$ on a disc
$|z-z_0|\le R$, with a submultiplicative norm bound $\|G\|\le M$,
its normalized solution satisfies $\|U\|\le e^{MR}$.
Gronwall's estimate \cite[Lemma~2.7, pp.~42--43]{TeschlODE2012}
and Cauchy's coefficient formula
\cite[\S1.9(iii), (1.9.30)--(1.9.31)]{DLMFCauchy2026}
therefore bound the tail after degree $N$, on $|z-z_0|\le r<R$, by
\begin{equation}
 e^{MR}\frac{(r/R)^{N+1}}{1-r/R}.
\label{pub:V:eq:1-16-1}
\end{equation}
For rational $G=P/Q$, the positive lower bound
$|Q_0|-\sum_{j\ge1}|Q_j|R^j$ gives $M$ directly.
These estimates, with outward interval operations, enclose the
root transports, the intertwiner and the four-equation real
system. Existence of the exact intertwiner follows from
\cref{night:P1.15}; an invertible square subsystem encloses it.
A matrix inverse on a rectangle and differentiation of
$Mv=b$ give
\[
 v_s=M^{-1}(b_s-M_sv),\qquad
 v_{st}=M^{-1}(b_{st}-M_{st}v-M_sv_t-M_tv_s).
\]
Together with $Y_x=YG$, $Y_y=iYG$ and
$Y_{xx}=Y(G^2+G')$, these enclose the real Hessian of $h_b$
on whole rectangles. For the value on a rectangle of
half-diagonal $\rho$, an entrywise bound $|G|\le\mathcal A$
gives
$|Y(z)-Y(z_0)|\le|Y(z_0)|_+(e^{\rho\mathcal A}-I)$.
All inputs to these enclosures are the displayed rational connection, its gauges,
the integral monodromies and the specified normalization.
The endpoint, Hessian, cusp and outer-cover enclosures are given by the
\hyperlink{prog:V:height}{executable certificate in Table~\ref*{tab:V:programs}}.

At infinity use the regular frame
$Y(u^{-5})D_\infty(u)=F_\infty U_\infty(u)$.
The normalized vector in that frame is $D_\infty(u)^{-1}q$;
its two variable coefficients are
$\widetilde v_2=u^8v_2$ and $\widetilde v_1=u^4v_1$.
The deck action sends this pair to
$(\zeta^8\widetilde v_2,\zeta^4\widetilde v_1)$,
where $\zeta=e^{2\pi i/5}$. Its value at the root point
is therefore zero. The height there is
$h_b(\infty)=-c_0\operatorname{Re}(F_\infty)_{03}$,
and the same deck action makes its gradient zero.
In the scaled coordinate $w=a^{1/5}u$, so $\lambda=aw^{-5}$,
the enclosures on the closed square $[-1/10,1/10]^2$ give
\[
 \min(h_{xx},h_{yy})-|h_{xy}|>\frac14.
\]
The real Hessian is consequently at least $I_2/4$ there.
Integrating along the segment from zero gives
\[
 h_b(aw^{-5})-h_b(\infty)\ge\frac18|w|^2.
\]
The normalized endpoint enclosure has radius below $4\cdot10^{-76}$
and lies between the negatives of the two endpoints in the statement.
\hypertarget{tag-03KY}{}\label{layout:V:fresh-minimum}

For the rest of the base use $v=a^{-1/4}u_0$, so
$\lambda=av^4$, and put $R=49/50$, $S=103/100$.
The first-quadrant $w$-disc $|w|\le R$ covers
$|\lambda|\ge aR^{-5}$ under its finite deck maps; the
first-quadrant $v$-disc covers $|\lambda|\le aR^4$.
The collar $R\le|w|\le S$ covers the intervening annulus,
since $R^4\ge S^{-5}$.
The cusp coordinate is $q=\lambda/\lambda_*-1=-w^{-5}-1$.
The same enclosure formulas on rational rectangular partitions
give the following strict lower bounds away from the central
square and the inner cusp disc:
\[
\begin{array}{c|c}
\text{region}&\text{lower bound for }h_b\\\hline
\text{infinity-root disc}&-1.094216\\
\text{zero-root disc}&-1.083056\\
\text{cusp annulus }1/100\le|q|\le1/4&-1.083777\\
\text{collar}&-1.030556
\end{array}
\]
Every endpoint in this table is an exact terminating decimal
strictly greater than $h_b(\infty)$. Rectangles meeting a
boundary are kept closed. Those lying in another of the stated
regions are handled there; those outside a specified radius are
excluded by their exact squared-distance inequality. The deck
maps and conjugation give the remaining sectors. Thus the
bounds cover all points outside $|q|<1/100$ and the central
square.

For the inner cusp disc, the exact gauge is
$D_\lambda=\operatorname{diag}(1,\lambda^2,\lambda,1,\lambda^{-1},\lambda^{-2})$.
Its connection is $R_c/q+A_0/(1+q)$ with $R_c^2=0$.
Writing $YD_\lambda=H_c(I+R_c\log q)P(q)$, $P(0)=I$,
the coefficient recurrence is
\[
 (nI-\operatorname{ad})P_n
 =\sum_{j<n}(-1)^{n-1-j}P_jA_0,
 \qquad\operatorname{ad}(Q)=QR_c-R_cQ.
\]
Here $\operatorname{ad}^3=0$, so the inverse is
$n^{-1}(I+\operatorname{ad}/n+\operatorname{ad}^2/n^2)$.
A scalar majorant $(1-|q|)^{-\gamma}$ gives the same
Cauchy tail control. Continuation to the positive cusp anchor
fixes $H_c$ and the logarithm branch. In the unimodular marking
$J$ with $J^{-1}NJ=\left(\begin{smallmatrix}0&I\\0&0\end{smallmatrix}\right)$,
the lower graph block is invertible throughout the closed
inner disc. The marking used for this anchor and the initial
cusp marking are related by $N=AN'A^{-1}$; their change has
block form $\left(\begin{smallmatrix}D&E\\0&D\end{smallmatrix}\right)$
and sends $Z$ to $DZD^{-1}+ED^{-1}$. It preserves the
real-intersection criterion.
\hypertarget{tag-03L0}{}\label{layout:V:fresh-marking}

At the parameter $\xi=ic_*$ write the resulting graph as
$Z(q)=\log q\,I_3/(2\pi i)+C(q)$, and put
$y_0=\log100/(2\pi)$ and $t=-\log|q|/(2\pi)-y_0\ge0$.
The coefficients, in increasing degree, of
$\det(tI_3+y_0I_3+\operatorname{Im}C(q))$ are bounded below
on $|q|\le1/100$ by
$3/2,9/2,15/4,1$, respectively.
The determinant is therefore positive for every nonzero $q$
in that disc. A graph vector $(Zv,v)$ is real precisely when
$v$ is real and $\operatorname{Im}Z\,v=0$. Thus the
$c_*$ plane is transverse there. The height at $q=1/100$
lies between $-0.451$ and $-0.450$; connectedness and
\eqref{maps:eq:F.22} give
$h_b(q)>h_b(\infty)$ on the entire punctured inner disc.

The strict central Hessian inequality, the four outer bounds and
the inner cusp inequality prove that infinity is the unique
minimum. The exact domain now follows from
\cref{night:P1.8}.
\hypertarget{tag-03L1}{}\label{layout:V:fresh-splitting}
\end{proof}

\begin{theorem}[The height at the marked basepoint]
\hypertarget{tag-02QP}{}\label{night:P1.16}
\StatementRecord{night:P1.16}{proved}{02QP}{\cref{night:P1.15}, \cref{night:P1.18}, \cref{night:P1.8}}
In the canonical splitting $\sigma_b$ of \cref{night:P1.15},
\[
\begin{aligned}
-1.0342766224163217994942488453783461&<h_b(1),\\
h_b(1)&<-1.0342766224163217994942488453783460.
\end{aligned}
\]
In particular $h_b(1)<-1$. The value $c=5/(3\pi)$, for which
the determinant $\det C_1$ in \cref{night:P1A.2} vanishes at
$\xi=ic$, lies outside the global transversality domain.
\end{theorem}
\begin{proof}
Use the root transports, normalized intertwiner and interval
linear system in \cref{night:P1.18}. At the marked basepoint
$H(1)=I$, so the real system in \cref{def:hb} uses $Y=X$.
Its enclosure has radius less than $7\cdot10^{-76}$ and is
contained in the following wider interval:
\[
\begin{aligned}
h_b(1)\in\bigl[&
-1.03427662241632179949424884537834605259496477790218821132\\
&\qquad\mathbin{\pm}8.11\cdot10^{-57}\bigr].
\end{aligned}
\]
This proves the two strict bounds. For the coefficient column
$q=(0,v_2,v_1,-i,0,0)^t$ of \cref{def:hb}, the normalized
vector at zero in the regular frame is $D_0(u)^{-1}q$; its variable coefficients
are $u^{-6}v_2,u^{-3}v_1$. The nontrivial deck characters
force their limiting values to vanish. With the regular limiting frame $F_0$
of \cref{night:P1.18}, the same enclosure method gives
\[
\begin{aligned}
h_b(0)&=-\frac{96}{625\pi}\operatorname{Re}(F_0)_{03}\\
&=-0.931793376356523505430795105796011114990828\ldots .
\end{aligned}
\]

The connection of this calculation with based-loop continuation
is an exact gauge identity. Put $a=256/3125$ and
\[
D(\lambda)=\operatorname{diag}
 (1,\lambda^2,\lambda,1,\lambda^{-1},\lambda^{-2}).
\]
Then $\widetilde Y=YD$ satisfies the Fuchsian equation
\[
\begin{aligned}
\widetilde\Gamma
 &=D^{-1}\Gamma_6D+D^{-1}D'
 =\frac{R_0}{\lambda}+\frac{R_*}{\lambda+a},\\
R_0&=\lim_{\lambda\to0}\lambda\widetilde\Gamma,
\qquad
R_* =\lim_{\lambda\to-a}(\lambda+a)\widetilde\Gamma.
\end{aligned}
\]
Since $D(1)=I$, the based meridians of \cref{def:hb}
retain their matrices $B,A$. The regular-root deck calculation
in \cref{night:P1.18} represents those same meridians.
Both descriptions therefore give the same normalized
intertwiner $X$, by its uniqueness in \cref{def:hb}.
If $T_{12}$ is the Fuchsian transition along $[1,2]$ and
$u=2^{-1/5}$, uniqueness of the normalized differential equation
gives the exact equality
\[
XT_{12}D(2)^{-1}
 =F_\infty U_\infty(u)D_\infty(u)^{-1}.
\]
Thus the loop and regular-root calculations refer to the same
period plane with the same scalar and based marking.

Finally $5/(3\pi)<1$ and $h_b(1)<-1$, whereas
$h_b(\lambda)\to+\infty$ at the cusp. Along a path from $1$
to that region the function $h_b+5/(3\pi)$ has a zero.
The exact transversality criterion of \cref{night:P1.8}
therefore excludes this parameter.
\end{proof}


\section{Cohomology, deformations and isomorphisms}\label{pub:V:sec-cohomology}

\begin{definition}[Root-orbifold notation]\hypertarget{tag-PV15}{}\label{def:V:root-orbifold-notation}
\StatementRecord{def:V:root-orbifold-notation}{proved}{PV15}{}
Use the complex orientation and the integral marking of \cref{night:P1.1}.
Let $\widehat B$ be the root orbifold of orders $5,4$, let
$h:M\to\widehat B$ be the map given by the two multiple divisors,
and let $\pi:\widehat B\to\mathbb P^1$ be its coarse map.
Write $L=T^2$; in the pencil coordinate it is
\[
 L=\mathcal O_{\widehat B}\bigl(\tfrac45[\infty]
                  -\tfrac34[0]\bigr),\qquad \deg L=\tfrac1{20}.
\]
The cusp point has line bundle $L^{20}$, and
$T_{\widehat B}=L^9$. These identities include the stabilizer
characters: their quotients have integral divisors of degree zero
on $\mathbb P^1$. Denote the cusp by $W$, its normalization by $Y$,
and its reduced singular locus by $D$.
\end{definition}

\begin{theorem}[Direct images and the canonical bundle]\hypertarget{tag-02R1}{}\label{night:P1A.1}
\StatementRecord{night:P1A.1}{proved}{02R1}{\cref{night:P1.10}, \cref{night:P1.8}, \cref{night:P1.9}}
For every admissible member, its vertical translation bundle on
$\widehat B$ and its first coherent direct image are
\[
 E_0=F^\vee=\mathcal O\oplus L^{-1}\oplus L^{-2},\qquad
 G=R^1h_*\mathcal O_M=\mathcal O\oplus L^{-1}\oplus L^{-2}.
\]
Moreover $R^qh_*\mathcal O_M=\bigwedge^qG$ for $0\le q\le3$,
$K_M=h^*L^{-6}$, and
\[
 (h^0(\mathcal O_M),\ldots,h^4(\mathcal O_M))=(1,1,0,0,0).
\]
\end{theorem}
\begin{proof}
\begin{enumerate}[label=\arabic*.]
\item The evaluation of the marked lattice on the lifted plane $F$
identifies the Lie bundle of the smooth tori with $F^\vee$.
\Cref{night:P1.8} gives $F=\mathcal O\oplus L\oplus L^2$; its lift as a bundle
has the same splitting. The logarithmic toric vector fields extend
this Lie bundle over the cusp, so the identification holds on the
whole orbifold, including both finite root charts.

\item On the smooth locus the constant-cocycle description of torus
cohomology gives $G=V/F$. In the rank-five orthogonal quotient,
$K=\ker(\alpha:P\to\mathcal O)=L\oplus L^2$ is the radical of
the restricted pairing on $P$. Indeed every wedge containing the
fixed line pairs to zero with another such wedge except through
their contraction components; removing that scalar component
leaves exactly $K$. The nondegenerate pairing consequently
identifies $W/P$ with $K^\vee$. Thus
\[
 0\longrightarrow\mathcal O e\longrightarrow V/F
   \longrightarrow L^{-1}\oplus L^{-2}\longrightarrow0.
\]
For every integer $j$, descent by invariant local Laurent series gives
\begin{equation}\hypertarget{tag-01O3}{}\label{eq:P1Adegree}
 \pi_*L^j=\mathcal O_{\mathbb P^1}(d(j)),\qquad
 d(j)=\lfloor4j/5\rfloor+\lfloor-3j/4\rfloor.
\end{equation}
The extension group $H^1(L\oplus L^2)$ therefore vanishes, so this
sequence splits as a holomorphic vector-bundle sequence.

\item We verify that this quotient is the coherent extension at the cusp.
The normalization resolution of $\mathcal O_W$ has one toric
threefold, seven toric surfaces, twelve toric curves and six points.
All positive coherent cohomology of those smooth complete toric
strata vanishes \cite[Theorem 9.2.3, p.~410]{CoxLittleSchenck2011}. Their constants form the cellular cochain complex of the periodic
Kuhn triangulation modulo $\mathbb Z^3$: its simplices are the faces of
$(0,e_{\pi(1)},e_{\pi(1)}+e_{\pi(2)},(1,1,1))$ for
$\pi\in S_3$, identified by integral translations. The quotient map
$x\mapsto x\bmod\mathbb Z^3$ identifies its realization with the dual
three-torus and respects every face map. Its cohomology is
$\bigwedge^q\mathbb C^3$, with the exterior cup product.
The invariant flat cohomology $\ker N$ maps to these classes by
constant Cech cocycles on the toroidal covering. A component
translation in that cover evaluates the corresponding lattice
coordinate. In the nilpotent logarithmic frame, $\ker N$ maps
isomorphically to $V/F$, since $N:F\to\ker N$ is an isomorphism.
Thus these are the extending quotient classes.

\item Flatness over each smooth root disc follows because its base is a
DVR and the smooth total space is torsion free over it. Properness was proved in \cref{night:P1.9,night:P1.10}. The coherent fibre dimensions are the
constant binomial dimensions on smooth tori, finite root fibres,
and the cusp. \cite[\S7, Satz~5]{Grauert60}, with these hypotheses,
gives local freeness and base change. The cup-product maps are
isomorphisms on every fibre, including the preceding central
cellular model, so they identify the remaining direct images with
$\bigwedge^qG$.

\item There is an exact sequence of locally free logarithmic bundles
\begin{equation}\hypertarget{tag-01O4}{}\label{eq:P1Alogcotangent}
0\longrightarrow h^*L^{11}\longrightarrow\Omega_M^1(\log W)
 \longrightarrow h^*F\longrightarrow0.
\end{equation}
In a chart $q=x_0\cdots x_r$, its first arrow sends
$dq/q$ to $\sum dx_i/x_i$ and its quotient is the height-zero
logarithmic cotangent bundle of the torus. These formulas also
prove exactness on every stratum. Taking determinants and
cancelling $\mathcal O_M(W)=h^*L^{20}$ yields
$K_M=h^*L^{11+3-20}=h^*L^{-6}$.

\item For any bundle which is a sum of powers of $L$, projection and
the exact finite averaging projector compute its direct images
by tensoring with $\bigwedge^qG$ and applying
\eqref{eq:P1Adegree}. All exponents $-1,\ldots,-5$ have coarse
degree $-1$, while exponent zero has degree zero. The Leray
sequence on the curve then gives the asserted structure-sheaf
cohomology. This includes the nonreduced coarse fibres, since
the computation was made before exact finite descent.

\end{enumerate}
\end{proof}




\begin{theorem}[The Hodge numbers]\hypertarget{tag-02R2}{}\label{night:P1A.2}
\StatementRecord{night:P1A.2}{proved}{02R2}{\cref{night:P1.15,night:P1A.1,night:P1.18}}
The complete Hodge table of every admissible fourfold is the
following, with rows indexed by $p$ and columns by $q$:
\[
\begin{array}{c|ccccc}
h^{p,q}&0&1&2&3&4\\\hline
0&1&1&0&0&0\\
1&0&2&1&0&0\\
2&0&1&6&1&0\\
3&0&0&1&2&0\\
4&0&0&0&1&1
\end{array}
\]
\end{theorem}
\begin{proof}
\begin{enumerate}[label=\arabic*.]
\item All maps needed for the two nontrivial rows can be evaluated at
the cusp. We give their matrices and the reason evaluation
determines their global ranks. In the algebraic splitting of
\cref{night:P1.15} write the three generators of $F$ as $(f_0,f_1,f_2)$ of
weights $(0,1,2)$, and those of $G$ as $(g_0,g_1,g_2)$ of
weights $(0,-1,-2)$. Put $z=(625\pi i/128)\xi$.
The residue of the quotient connection is
\[
R:F_0\longrightarrow G_0,\qquad
R=\begin{pmatrix}0&A&B\\C&0&D_0\\E&D_0&0\end{pmatrix},
\]
where
\begin{align*}
A&=-125(192z-3125)/12288,&
B&=-78125(96z+3125)/393216,\\
C&=-125/512,&E&=78125/32768,\\
D_0&=-146484375/2097152.
\end{align*}
Here the lifted first vector is $w_0+(4/3)ze$.
These entries result from the exact Gauss--Manin reduction of
\cref{night:P1.15}, with the quotient projection taken after lifting $F$.

\item The connecting map of \eqref{eq:P1Alogcotangent}, on relative
cohomology, sends $f\otimes\eta$ to $R(f)\wedge\eta$ at the
cusp. To see the extension class, trivialize logarithmic
forms on the toric covering. A deck translation changes a
relative logarithmic form by its period derivative times
$dq/q$. The constant Cech class of this derivative is exactly
the residue of the quotient connection. This proves the
formula with the indicated frames. It also proves its exterior
power version, with the usual alternating signs.

\item The floor formula gives the following absolute cohomology:
\[
\begin{array}{c|ccccc}
&h^0&h^1&h^2&h^3&h^4\\\hline
h^*F&1&3&2&0&0\\
h^*L^{11}&0&2&3&1&0\\
h^*(\bigwedge^2F)&0&2&3&1&0\\
h^*(L^{11}\otimes F)&2&6&7&3&0.
\end{array}
\]
Every group in this display comes from global sections of the
indicated direct image; none comes from $H^1$ of a line on the
coarse base. Every contributing summand has coarse degree zero,
so evaluation at the cusp is injective on its space of sections.
For the equal-dimensional source and target spaces used below
it is an isomorphism onto the displayed weight subspaces.
Thus fibrewise evaluation computes the ranks of the
absolute connecting maps.

\item For one-forms, the three middle source classes are $f_i\otimes g_i$.
In the ordered target $(g_0g_1,g_0g_2,g_1g_2)$, the matrix is
\[
C_1=\begin{pmatrix}-C&A&0\\-E&0&B\\0&-D_0&D_0\end{pmatrix},
\qquad
\det C_1=
\frac{476837158203125}{281474976710656}(384z+3125).
\]
The degree-zero connecting map is nonzero because $R(f_0)$
has nonzero $g_1,g_2$ coefficients. The degree-two connecting
map, from the classes $f_1g_0g_1,f_2g_0g_2$ to $g_0g_1g_2$,
has rank one because $D_0\ne0$.

\item We must exclude the sole possible zero of $\det C_1$ on the
whole admissible parameter domain. For $\xi=ic$ the preceding
coordinate relation gives $z=-(625\pi/128)c$.
The unique exceptional parameter is therefore
$\xi=5i/(3\pi)$. By \cref{night:P1.18}, the admissible half-plane is
$\operatorname{Im}\xi>-h_{b,\min}>1$, whereas $5/(3\pi)<1$.
Thus this zero lies outside the whole admissible parameter domain.
Consequently the three connecting ranks are $1,3,1$.
The long exact sequence of \eqref{eq:P1Alogcotangent} gives
\[
h^q(\Omega_M^1(\log W))=(0,1,1,0,0).
\]

\item The residue sequence
\[
0\to\Omega_M^1\to\Omega_M^1(\log W)\to\nu_*\mathcal O_Y\to0
\]
is exact even at self-intersections: in a local branch chart it
records independently the coefficient of each $dx_i/x_i$ on
$x_i=0$. The normalization $Y$ is smooth complete toric, so its
only coherent structure cohomology is $H^0=\mathbb C$ \cite[Theorem 9.2.3, p.~410]{CoxLittleSchenck2011}.
The one-form row follows immediately.

\item For two-forms take the second exterior power of
\eqref{eq:P1Alogcotangent}. Its short exact sequence has left
term $h^*(L^{11}\otimes F)$ and right term
$h^*(\bigwedge^2F)$. In degree two the three source classes
are $f_if_j\otimes g_ig_j$, $i<j$. The target is
$F\otimes g_0g_1g_2$. Its connecting matrix is
\[
C_2=\begin{pmatrix}-D_0&D_0&0\\E&0&-B\\0&-C&A\end{pmatrix},
\qquad \det C_2=-\det C_1.
\]
Thus it is an isomorphism, and
$h^3(\Omega_M^2(\log W))=1$.

\item The ordinary/logarithmic comparison in this degree is the
exact residue resolution
\[
0\to\Omega_M^2\to\Omega_M^2(\log W)
 \to\nu_*\Omega_Y^1(\log\partial Y)
 \to\nu^{[2]}_*\mathcal O_{W^{[2]}}\to0.
\]
The last space is the disjoint normalization of the seven
double surfaces. Local exactness is verified by writing a
logarithmic two-form as a sum of
$a_{ij}\,dx_i/x_i\wedge dx_j/x_j$ and terms with ordinary
differentials. The two residues of an $ij$ term have opposite
sign; their cancellation is the sole compatibility. Every local
double-surface coefficient lifts to one branch, which proves
surjectivity at the last term.
Both final sheaves have zero positive cohomology: the first
is the pushforward of the trivial rank-three toric logarithmic
bundle, and the second has complete toric surface components.
Writing $Q=\Omega_M^2(\log W)/\Omega_M^2$, we obtain
$H^q(Q)=0$ for $q\ge2$. Hence
$h^3(\Omega_M^2)=1$.
In degree zero the last residue map is zero: the fourteen toric
boundary divisors of $Y$ occur in seven opposite pairs; their
primitive normal vectors have opposite signs and their gluing
is a torus translation. This gives the zero $7$-by-$3$ residue
matrix on global logarithmic one-forms. Thus $H^0(Q)=\mathbb C^3$.
The two global logarithmic two-forms come from the
$L^{12},L^{13}$ summands of $L^{11}\otimes F$; their residues
are the independent $f_1,f_2$ directions at the cusp.
Therefore $h^0(\Omega_M^2)=0$.

\item Serre duality \cite[Theorem 4, p.~22]{Serre1955Duality} on this compact smooth complex fourfold gives
$h^{2,1}=h^{2,3}=1$. Finally the analytic Riemann--Roch theorem
\cite[Theorem 25.3.1, p.~187; \S25.4(a), pp.~188--190]{HirzebruchTopologicalMethods} applies to the Dolbeault elliptic
operator of $\Omega_M^p$.
Using $c_1=c_3=0$, the Chern-root product gives
\[
\chi_y(M)=
\frac{(1+y)^4}{240}\,c_2^2+
\left(\frac{(-1+10y-y^2)^2}{144}
       -\frac{(1+y)^4}{120}\right)c_4
=-y+4y^2-y^3.
\]
The displayed polynomial gives
$\chi(\Omega_M^2)=4$, hence $h^{2,2}=6$.
\Cref{night:P1A.1} and Serre duality supply all other entries of the table.

\end{enumerate}
\end{proof}




\begin{theorem}[Tangent cohomology and the Kuranishi germ]\hypertarget{tag-02R3}{}\label{night:P1A.3}
\StatementRecord{night:P1A.3}{proved}{02R3}{\cref{night:P1.9}, \cref{night:P1.10}, \cref{night:P1.15}, \cref{night:P1A.1}, \cref{night:P1A.2}}
For every admissible fourfold,
\[
 (h^0(T_M),h^1(T_M),h^2(T_M),h^3(T_M),h^4(T_M))
 =(1,1,1,2,0).
\]
The complex parameter family has a nonzero Kodaira--Spencer map
at every member and supplies a smooth one-dimensional miniversal
Kuranishi germ.
\end{theorem}
\begin{proof}
\begin{enumerate}[label=\arabic*.]
\item Put $a=-256/3125$, and denote the whole reduced cusp fibre by
$W$.  The logarithmic sequence established in the proof of
\cref{night:P1A.2}, tensored with $K_M=h^*L^{-6}$ from
\cref{night:P1A.1}, is
\[
0\longrightarrow h^*L^5
\longrightarrow\Omega^1_M(\log W)\otimes K_M
\longrightarrow h^*(F\otimes L^{-6})\longrightarrow0.
\]
Here $K_{\widehat B}(a)=L^{11}$.  More precisely, in the meromorphic
$L$-frames of \cref{night:P1.15},
$\kappa=d\lambda/(\lambda-a)$ is $\lambda^9$ times the
$L^{11}$-frame: both sides have divisor $3D_0-D_\infty$.
The logarithmic sequence is exact in each chart
$q=z_0\cdots z_r$, where its first arrow sends $dq/q$ to
$\sum_{j=0}^r dz_j/z_j$; the quotient is the relative translation
cotangent bundle $h^*F$.  These arrows agree on the monomial and deck transitions and descend through the free finite
root-chart actions.

\item For a summand $L^n$, the coarse degree is
$d(n)=\lfloor4n/5\rfloor+\lfloor-3n/4\rfloor$.
Projection formula and
$R^jh_*\mathcal O_M=\bigwedge^jG$ compute the two end terms.
For $h^*L^5$ the weights in relative degrees $j=0,1,2,3$ are
\[
\{5\},\quad\{5,4,3\},\quad\{4,3,2\},\quad\{2\}.
\]
Weights $4,5$ have degree zero and weights $2,3$ have degree $-1$.
Thus $h^q(h^*L^5)=(1,2,1,0,0)$.
The three weights of $F\otimes L^{-6}$ are $-4,-5,-6$.
After tensoring with $\bigwedge^jG$, every weight is between
$-9$ and $-4$.  We have $d(-9)=-2$ and $d(n)=-1$ for
$-8\le n\le-4$.  Weight $-9$ occurs once in each of relative
degrees $2$ and $3$, so
\[
h^q\bigl(h^*(F\otimes L^{-6})\bigr)=(0,0,0,1,1).
\]
The base has cohomological dimension one, so the two-column Leray
spectral sequences have no further differential.  Every connecting
map for the displayed logarithmic sequence has zero target.
Consequently
\[
h^q\bigl(\Omega^1_M(\log W)\otimes K_M\bigr)=(1,2,1,1,1).
\]

\item Let $\nu:Y\to W$ be the branch normalization.  The residue
sequence, with its branchwise maps, is
\[
0\longrightarrow\Omega^1_M\otimes K_M
\longrightarrow\Omega^1_M(\log W)\otimes K_M
\longrightarrow\nu_*\mathcal O_Y\otimes K_M|_W
\longrightarrow0.
\]
Since $K_M$ is pulled back from the base, its restriction to $W$ is
a constant line.  The complete smooth toric variety $Y$ has
$H^0(\mathcal O_Y)=\mathbb C$ and $H^{>0}(\mathcal O_Y)=0$
\cite[Theorem~9.2.3, p.~410, with divisor zero]{CoxLittleSchenck2011}.
The sole global section of the logarithmic middle term comes from
$h^*L^5$.  Its meromorphic expression is
$\lambda^{-5}\kappa$ tensored with the $L^{-6}$-frame: the
unique $L^5$ section is $\lambda^4$ times its frame, and
$\kappa$ contributes $\lambda^9$.  Its residue at $a$ is
$a^{-5}\ne0$.  Hence the map on global sections is an isomorphism
onto the residue line, and
\[
h^q(\Omega^1_M\otimes K_M)=(0,2,1,1,1).
\]
The pointwise perfect evaluation pairing
$T_M\otimes(\Omega^1_M\otimes K_M)\to K_M$ induces analytic
Serre duality
\cite[\S13, Theorem~4, p.~22; \S15, p.~23]{Serre1955Duality}.
It identifies $H^q(T_M)^\vee$ with
$H^{4-q}(\Omega^1_M\otimes K_M)$ and gives the stated tangent
cohomology dimensions.

\item We check the class of the parameter family on the
smooth fibre $A=h^{-1}(1)$.  In the fixed frame its period plane
has basis
\[
w_2,\quad w_1,\quad w_0+\frac{625\pi i}{96}\xi e.
\]
Differentiation in $\xi$, followed by projection to $G=H/F$,
gives the rank-one map $P:F\to G$ with sole nonzero matrix entry
$P_{13}=625\pi i/96$ in the ordered bases above and
$(e,w_{-1},w_{-2})$ of the quotient.
Differentiation in the base coordinate gives, by projecting the
connection in \cref{night:P1.15}, the matrix
\[
Q=\frac1{3125(1-a)}
\begin{pmatrix}
-78125/6+300r&-390625/48-375r&0\\
120&0&125/2\\
0&120&50
\end{pmatrix},
\qquad r=\frac{625\pi}{96}(-i\xi).
\]
Its determinant is $3/[625(1-a)^3]\ne0$.
The translation-valued \v Cech cocycles obtained by differentiating
the lattice gluing identify
$H^1(T_A)=\operatorname{Hom}(F,G)$; contraction with a relative
one-form sends the two cocycles to the displayed projected period
derivatives, in the same convention.

\item The normal tangent sequence on this fibre is
\[
0\longrightarrow T_A\longrightarrow T_M|_A
\longrightarrow\mathcal O_A\otimes T_{\widehat B,1}
\longrightarrow0.
\]
Its connecting map sends a base tangent vector to its
Kodaira--Spencer class $Q$.  As $H^0(\mathcal O_A)=\mathbb C$,
the kernel of $H^1(T_A)\to H^1(T_M|_A)$ is exactly
$\mathbb C Q$.  The nonzero matrix $P$ has determinant zero;
every nonzero scalar multiple of $Q$ has nonzero determinant.
Therefore the image of $P$ in $H^1(T_M|_A)$ is nonzero.
The global parameter Kodaira--Spencer cocycle restricts to this
image, so its class in $H^1(T_M)$ is nonzero.
The finite affine actions are fixed, and the cusp multipliers
$\exp(2\pi iC(q,\xi)\nu)$ vary holomorphically in $\xi$.
The same finite fan charts, with $C$ bounded on a small parameter
neighbourhood, give the proper local family of
\cref{night:P1.9}, \cref{night:P1.10}.  Thus the cocycle just computed is
the derivative of this compact family.

\item Douady's versality and semiuniversality theorems
\cite[\S VII.8, pp.~598--599; \S VIII.1, pp.~599--600]{Douady1974}
give a miniversal germ with tangent space $H^1(T_M)=\mathbb C$.
Its local ring has a minimal presentation $\mathbb C\{u\}/I$,
where $I\subset(u)^2$.  The parameter disc classifies by
$u\mapsto g(s)$ with $g'(0)\ne0$, by the preceding step.
Substitution $\mathbb C\{u\}\to\mathbb C\{s\}$ is therefore an
isomorphism of convergent power-series rings.  Since this map factors
through $\mathbb C\{u\}/I$, it annihilates $I$; its injectivity
gives $I=0$.  The miniversal germ is reduced and smooth of dimension
one, and the classifying map is a local isomorphism.  This proves
miniversality of the given family at each member.
\end{enumerate}
\end{proof}




\begin{theorem}[Algebraic dimension and automorphisms]\hypertarget{tag-02R4}{}\label{night:P1A.4}
\StatementRecord{night:P1A.4}{proved}{02R4}{\cref{night:P1.1}, \cref{night:P1.9}, \cref{night:P1.15}, \cref{night:P1A.1}}
The fourfold has algebraic dimension one and
\[
\Aut^0(M)=\mathbb C,\qquad
\Aut(M)=\mathbb C\times\mathbb Z/2.
\]
Two members $M_\xi,M_\eta$ of the family in
\cref{pub:V:parameter} are biholomorphic exactly when $\xi=\eta$.
\end{theorem}
\begin{proof}
\begin{enumerate}[label=\arabic*.]
\item An effective divisor dominating the base would restrict to
a nonempty effective divisor on a general complex torus.
Pairing its Chern class with the square of a Kähler form
on that torus is strictly positive. The restriction class is an alternating integral form $J$ on $\Lambda$
fixed by both monodromies:
\[
M_1^tJM_1=J,\qquad M_2^tJM_2=J.
\]
The exact linear system in its fifteen alternating entries has rank
fifteen, so $J=0$.  This contradicts the positive pairing, and every
divisor is vertical.
The pole divisor of a meromorphic function is vertical;
its restriction to a general compact fibre is holomorphic
and therefore constant. The compact reduced graph of the function maps properly to
$\mathbb P^1\times\mathbb P^1$ by $h$ on its first factor.
Remmert's proper-image theorem, in
\cite[\S7, Satz~1, p.~60]{Grauert60}, makes its image analytic.
It is generically a single graph and therefore defines the descended
meromorphic function on the normal base curve, including the three
special values.
Conversely pullbacks from the base occur. Hence the
meromorphic function field is $\mathbb C(t)$ and $a(M)=1$.
Every automorphism consequently induces a Möbius map on
the base. The two distinct multiplicities and the cusp
identify all three values separately, so this map is the
identity.

\item On smooth fibres an automorphism is affine. Its integral
linear part is locally constant and commutes with $M_1,M_2$.
The two simultaneous commutator equations have rational solution
space $\mathbb QI$ by rational row reduction. Their integral invertible
elements are therefore $\pm I$. At the order-five filling an affine lift with
linear part $\epsilon I$ would have to conjugate the same
finite affine generator. For its translation $b$ this
requires
\[
(I-M_1)b=(1-\epsilon)v_1/5\pmod\Lambda.
\]
Applying the invariant integral covector $\ell$ makes the
left side zero. For $\epsilon=-1$ the right side is
$2/5\pmod\mathbb Z$, which is nonzero. Thus every
automorphism is a translation.

\item Let $\mathcal A$ be the sheaf of translations extending to
the cusp model and finite root charts. There is
an exact analytic sheaf sequence
\[
0\to j_*\Lambda\to E_0\to\mathcal A\to0
\]
on the root orbifold. Here $j$ omits the cusp and the stalk
of $j_*\Lambda$ there is $\ker N$.
On a finite root cover the torus evaluation sequence is exact and
equivariant for the prescribed affine action; its translation part
cancels under conjugation.  We descend this sequence on the orbifold, with its finite stabilizer cohomology.

To check the last surjectivity at the cusp, lift a translation
to the toroidal cover. It preserves the fan precisely when
its valuation is an integral translation of the periodic
fan. Subtract the corresponding period; the remaining
characters are nonvanishing holomorphic functions on a
disc and possess logarithms. These give a local section
of $E_0$. The kernel consists exactly of the invariant
integral periods. This also proves that the sequence
accounts for all extendable translations, including
possible component permutations.

\item Write $A=M_1$, $B=M_2$ and
$N_1=I+A+\cdots+A^4$, $N_2=I+B+B^2+B^3$.
The orbifold fundamental group away from the cusp has generators
$g_1,g_2$ with relations $g_1^5=g_2^4=1$.  A lattice cocycle is
therefore a pair $(x_1,x_2)$ with $N_1x_1=N_2x_2=0$.
The degree-one exact sequence for $j_*$ says that it extends across
the cusp precisely when its value on $g_1g_2$ is a coboundary:
\[
x_1+Ax_2=(AB-I)u,\qquad u\in\Lambda.
\]
Subtract the coboundary of $u$.  The resulting pair is $(-Ay,y)$,
where $N_1y=N_2y=0$.  The ambiguity in $u$ is
$t\in\ker(AB-I)$, which changes $y$ by $(B-I)t$.
Thus the cohomology group, with the finite-root torsion, is
\[
H^1(\widehat B,j_*\Lambda)=
\frac{\ker N_1\cap\ker N_2}{(B-I)\ker(AB-I)}.
\]
Primitive integral column bases for its numerator and for
$\ker(AB-I)$ are respectively
\[
K=\begin{pmatrix}
0&0&0\\0&1&0\\1&0&0\\0&0&0\\0&-1&0\\0&0&1
\end{pmatrix},\qquad
T=\begin{pmatrix}
0&0&0\\0&1&0\\1&0&0\\0&0&-1\\0&0&1\\0&0&-1
\end{pmatrix}.
\]
Direct multiplication gives
\[
(B-I)T=KZ,\qquad
Z=\begin{pmatrix}-2&-1&1\\0&-1&1\\1&0&1\end{pmatrix}.
\]
The determinant of $Z$ has absolute value two; an entry and a
$2$-by-$2$ minor have absolute value one.  Its Smith factors are
$1,1,2$, proving $H^1(\widehat B,j_*\Lambda)=\mathbb Z/2$.
\item Now $H^0(j_*\Lambda)=0$ by \cref{night:P1.1},
$H^0(E_0)=\mathbb C$, and $H^1(E_0)=0$ by the floor
calculation. The exact sequence gives
\[
0\to\mathbb C\to H^0(\mathcal A)\to\mathbb Z/2\to0.
\]
The half-period
$\tau=(0,1,1,1,1,1)^t/2$ is fixed modulo $\Lambda$
by every monodromy matrix. It extends over the finite
quotients because the translations commute there.
At the cusp $N\tau$ is integral; primitiveness of $N$
provides $\lambda\in\Lambda$ with $N\lambda=N\tau$.
The representative $\tau-\lambda\in\ker N$ is a bounded
angular half-period and extends on every toric chart.
It is a nonzero order-two global section. Since the
additive group $\mathbb C$ has no nonzero torsion, it
splits the displayed exact sequence. Translations
commute, giving the asserted full group.

\item For an isomorphism between two members, the same intrinsic
base argument applies. Its integral linear part still
lies in the common centralizer. Either scalar $\pm I$
preserves a complex period plane exactly by equality
of the two planes. The graphs
$(\sigma_b+\xi\alpha e)(P)$ are equal precisely when their
parameters $\xi$ agree, since $e\notin F$ and $\alpha$
is nonzero. Translation does not change this condition.
This proves the parameter assertion.

\end{enumerate}
\end{proof}




\begin{theorem}[No antiholomorphic isomorphisms]\hypertarget{tag-02QQ}{}\label{night:P1.17}
\StatementRecord{night:P1.17}{proved}{02QQ}{\cref{night:P1.1,night:P1A.4,night:P1.9}}
For the cusp marking $S$ fixed in \cref{night:P1.1} and the corresponding Kuhn fan,
none of the completed fourfolds in this family admits an
anti-holomorphic automorphism.  In particular none admits a real
structure.  More strongly, two members with this same cusp marking
are never anti-biholomorphic.
\end{theorem}
\begin{proof}
\begin{enumerate}[label=\arabic*.]
\item By \cref{night:P1A.4}, the meromorphic function field of each
member is $\mathbb C(\lambda)$, with the given base map as its
algebraic reduction.  An anti-biholomorphism therefore induces an
anti-M\"obius map of the base.  The two distinct multiplicities $5,4$
and the singular cusp
support distinguish the three special values, so this base map fixes
$\infty,0,\lambda_*$.  As those points are real, the map is precisely
$\lambda\mapsto\bar\lambda$.

\item A holomorphic map between compact complex tori is affine: after
lifting to their vector spaces, its derivative is periodic and
therefore constant, since it descends to holomorphic functions on a
compact torus.  Applying this to an anti-holomorphic map after
conjugating its source gives the same assertion for its anti-linear
lift.  Its induced integral map on $H_1$ is constant along the
regular base.  Complex conjugation reverses the two finite meridians,
so this matrix must simultaneously conjugate $M_i$ to $M_i^{-1}$.
Solving these exact integral linear equations gives only $R$ and
$-R$ among unimodular matrices, where
\[
R=\begin{pmatrix}
1&0&0&0&0&0\\
2&0&0&0&-1&0\\
-1&0&-1&0&0&0\\
-1&0&0&-1&0&0\\
2&-1&0&0&0&0\\
-1&1&1&1&1&1
\end{pmatrix},\qquad R^2=I.
\]
Indeed the simultaneous intertwiner space over $\mathbb Q$ has
dimension one and this primitive matrix is unimodular; determinant
one in absolute value permits only the scalar factors $\pm1$.


\item The finite affine cap excludes $G=R$.  Normalize the pullback by
$t=u^{m_i}$; any base-preserving isomorphism lifts to this canonical
root cover and conjugates its deck generator to the inverse.
At its central torus the translation part $t_i$ would satisfy
\[
(I-M_i^{-1})t_i=-(v_i+Gv_i)/m_i\pmod\Lambda.
\]
Applying $\ell$, and using $\ell R=\ell$, gives
$0=-2\ell(v_i)/m_i$ in $\mathbb R/\mathbb Z$, impossible already
for $m_i=5$.  Hence $G=-R$ is necessary.

\item We now use the cusp marking, keeping track of the conjugation of based meridians.  Conjugation sends
$g_0=(g_1g_2)^{-1}$ to $g_2g_1=g_2g_0^{-1}g_2^{-1}$.
Thus the matrix in a fixed local cusp frame is
$H=M_2^{-1}G$, and it anti-conjugates $M_0$.
Exact multiplication in the marking $S$ gives
\begin{equation}
S^{-1}HS=
\begin{pmatrix}
1&0&0&-1&-1&0\\
0&1&0&-2&0&0\\
1&-1&-1&1&-1&-1\\
0&0&0&-1&0&0\\
0&0&0&0&-1&0\\
0&0&0&-1&1&1
\end{pmatrix}
=\begin{pmatrix}P&Q\\0&-P\end{pmatrix},\quad
P=\begin{pmatrix}1&0&0\\0&1&0\\1&-1&-1\end{pmatrix}.
\label{pub:V:eq:1-17-1}
\end{equation}
Changing the based cusp path by a cusp loop multiplies by a power of
$M_0$ and does not change $P$.

\item We explain why an arbitrary translation cannot remove the obstruction
in \eqref{pub:V:eq:1-17-1}.  In the additive torus coordinates whose first periods
are the $a_j$, the anti-linear map must have the form
$\zeta' =P\bar\zeta+t(q)$, because its values on these three real
periods already determine its linear part.  In multiplicative
coordinates it is
$z'=u(q)\bar z^{-P}$.
The source zero section limits to a smooth point of the central fibre.
An isomorphism takes it to another smooth point.  The central fibre
is irreducible, and its smooth locus is exactly the dense torus of
its toric normalization: every toric boundary divisor is a
normal-crossings self-identification.  In a chart about this image,
the translation therefore has the form $q^v$ times nonvanishing
holomorphic units, with $v\in\mathbb Z^3$.  Translating by the period
$v$ changes the selected component chart; the units act by torus
translation.  Neither changes the linear action $-P$ on the fan.

\item The primitive rays of the normal fan of that component are
\[
\{\pm v:0\ne v\in\{0,1\}^3\}.
\]
An extending torus monomial isomorphism must permute these rays.
This can be seen directly on a one-parameter arc $z=t^v$ meeting the
generic point of its boundary divisor: if the image vector lies in
the interior of a cone of dimension at least two, that divisor is
contracted, which is impossible for an isomorphism.  Complex
conjugation itself preserves the toric charts, so the same criterion
applies to the anti-monomial map above.
But $-Pe_2=(0,-1,1)^t$ is not one of the displayed rays.
This is a contradiction.  We have exhausted both possible integral
linear parts, so no anti-biholomorphism exists.


\end{enumerate}
\end{proof}


\section{The boundary of the family}\label{pub:V:sec-boundary}

\begin{definition}[Analytic extension at the parameter boundary]\hypertarget{tag-PV16}{}\label{def:V:analytic-boundary-extension}
\StatementRecord{def:V:analytic-boundary-extension}{proved}{PV16}{}

We use the canonical splitting $\sigma_b$, height $h_b$, parameter
domain $\mathcal U$ and family $M_\xi$ of \cref{pub:V:parameter}.
The exact coordinate change for another splitting is
\cref{rem:V:splitting-change}.
By an analytic extension across a boundary point we mean a proper
Hausdorff complex-analytic family on a full complex disc whose
restriction to the nonempty open part lying in $\mathcal U$ is
the prescribed family. At infinity we use the coordinate $u=1/\xi$.
This definition permits singular central spaces, multiplicities,
and additional central components.
\end{definition}

\begin{remark}[Deformations of complex tori]\hypertarget{tag-PV17}{}\label{rem:V:tori-deformations-input}
\StatementRecord{rem:V:tori-deformations-input}{proved}{PV17}{}
The stability of complex tori under deformation in the large is
established in \cite[definition, p.~261; Theorem 2.1, p.~264]{Catanese2004}.
The relative-form calculation below writes the torus quotient
on the regular analytic open set used in the boundary argument.
\end{remark}

\begin{theorem}[No analytic extension across the boundary]\hypertarget{tag-02R5}{}\label{night:P1A.5}
\StatementRecord{night:P1A.5}{proved}{02R5}{\cref{night:P1A.1}, \cref{night:P1.8}, \cref{night:P1.18}, \cref{rem:V:conjugation-parameter}; \cite[Theorem 1.6, preprint p.~5]{BM97}; \cite[\S6, Hauptsatz I; \S7, Satz 1 and Satz 5]{Grauert60}}
There is no such analytic extension across any finite boundary
point of $\mathcal U$, or across infinity. The same conclusion
holds after any finite ramified change of the local boundary
coordinate. In particular neither end of the real interval
$c_*<c<+\infty$ has a proper Hausdorff analytic extension in
this sense.
\end{theorem}
\begin{proof}
\begin{enumerate}[label=\arabic*.]
\item We first identify the moving pencil on each known member.
By \cref{night:P1A.1}, $K_M^{-4}=h^*L^{24}$ and the fibres of $h$
are connected.  For the coarse map
$\rho:\widehat B\to\mathbb P^1$, the root-divisor formula gives
\[
 \rho_*L^{24}
 =\mathcal O_{\mathbb P^1}
   \bigl(\lfloor96/5\rfloor\,\infty+
                \lfloor-72/4\rfloor\,0\bigr)
 =\mathcal O_{\mathbb P^1}(19\infty-18\,0).
\]
Its sections are spanned by $\lambda^{18},\lambda^{19}$.
On a root chart the invariant sections have a common local
monomial factor, and their ratio cancels that factor.  Thus
$h^0(M,K_M^{-4})=2$, and the moving part of this complete
system is exactly the marked base map, up to a projective
coordinate change.  In particular its base ideal is invertible
on every member of the prescribed family.

\item Suppose that a proper Hausdorff analytic extension over a full
parameter disc exists.  Take the reduced component containing
the given nonempty open family and resolve it.  Analytic
resolution applies on a relatively compact neighbourhood of
the compact central fibre and can be kept unchanged over the
given smooth locus
\cite[Theorem 1.6 and its following analytic-neighbourhood
paragraph, preprint p.~5]{BM97}.
Write $p:Z\to\Delta$ for the resulting proper map from a smooth
irreducible space.  The proper images of the exceptional and
critical loci are analytic
\cite[\S7, Satz 1, p.~60]{Grauert60}; they omit the given open
parameter set.  After shrinking the disc they are supported
at its centre.  Hence $p$ is a proper holomorphic submersion
over $\Delta^*$, with compact fourfold fibres.

Put $\mathcal J=\omega_{Z/\Delta}^{-4}$.
The sheaves $R^jp_*\mathcal J$ are coherent
\cite[\S6, Hauptsatz I, p.~59]{Grauert60}.
Multiplication by a nonzero parameter germ is injective on
$\mathcal J$ and on $p_*\mathcal J$, so $p_*\mathcal J$ is locally free
over the disc and free after shrinking.
The exact sequence for a fibre over $t$ gives
\[
 0\longrightarrow(p_*\mathcal J)_t/\mathfrak m_t(p_*\mathcal J)_t
 \longrightarrow H^0(Z_t,\mathcal J_t)
 \longrightarrow
 \ker\bigl((t'-t): (R^1p_*\mathcal J)_t
                       \longrightarrow(R^1p_*\mathcal J)_t\bigr)
 \longrightarrow0.
\]
The last term is supported in the torsion support of the
coherent sheaf $R^1p_*\mathcal J$, a closed discrete analytic
set.  At a point of the given open family outside this support,
the displayed comparison shows that $p_*\mathcal J$ has rank
two.  Shrinking makes the support empty away from the centre.
Thus a frame of $p_*\mathcal J$ gives the complete two-dimensional
linear system on each punctured fibre.  

\item The two evaluation sections define a coherent base ideal and
a meromorphic map $Z\dashrightarrow\mathbb P^1$ over the full
disc.  The ideal is locally invertible on the given open family
by the first step.  Its non-invertibility locus is analytic;
its proper image in $\Delta$ is therefore discrete and can be
removed away from the centre.  Blowing up that ideal and
resolving the graph consequently changes only the central
fibre.  It gives a proper holomorphic map to
$\mathbb P^1\times\Delta$.  The construction uses the graph
of the coherent evaluation map, so its coefficients have
finite meromorphic order in the parameter.

Its critical image is proper analytic.  Every component
dominating the parameter disc meets the given open family,
where there are precisely three critical values.  Their
degrees over the parameter are one.  The multiplicities five
and four distinguish two of them, and the third is the cusp.
After shrinking, their normalizations give three distinct
labelled meromorphic sections, and all vertical exceptional
components lie over the centre.  The projective transformation
taking these sections to $\infty,0,\lambda_*$ has meromorphic
coefficients and is invertible on $\Delta^*$.  Resolving its
graph gives the marked map $H$ over the disc, still
unchanged on punctured fibres.  Its regular restriction is a
proper holomorphic submersion
\[
 q:W\longrightarrow
 B=\bigl(\mathbb P^1\setminus\{0,\infty,\lambda_*\}\bigr)
          \times\Delta^*.
\]
Ehresmann transport makes the number of components of a fibre
locally constant.  The base is connected and the known fibres
are connected, so every fibre of $q$ is connected.

\item We need the torus structure on an analytically open dense
part of this regular family.  The coherent sheaf
$\mathcal E=q_*\Omega^1_{W/B}$ has generic rank three.
Generic cohomology and base change for this flat bundle gives
a proper analytic exceptional set outside which $\mathcal E$
is locally free of rank three and its fibre is the space of
holomorphic one-forms.  The constant-dimension comparison is
\cite[\S7, Satz 5, pp.~63--64]{Grauert60}.
Use the global coherent evaluation map
\[
 q^*\mathcal E\longrightarrow\Omega^1_{W/B}.
\]
Its cokernel has analytic support.  The proper image of that
support is analytic in $B$ and omits the nonempty part of the
known torus open set outside the generic exceptional set.
Enlarge the exceptional set by this image.  On its complement
evaluation is an isomorphism of rank-three bundles.

Relative exterior differentiation defines an
$\mathcal O_B$-linear map
$\mathcal E\to q_*\Omega^2_{W/B}$.  It vanishes on the known
torus open set.  The target is torsion-free, since it is the
proper direct image of a vector bundle under a submersion.
The identity theorem therefore makes this map zero.  Each
remaining compact connected fibre has three everywhere
independent closed holomorphic one-forms.  Their dual vector
fields commute and are complete by compactness.  Their flows
give a transitive holomorphic $\mathbb C^3$-action: every orbit
is open and the fibre is connected.  Its stabilizer $\Gamma$
is discrete, and compactness of $\mathbb C^3/\Gamma$ makes
$\Gamma$ a full real lattice.  This is the complex-torus
structure on each such fibre.

\item Let $E\subset B$ be the proper analytic exceptional set just
removed.  On the universal cover of $B\setminus E$, use
Gauss--Manin transport from the given open family to mark
integral $H^1$.  Integration of the closed relative forms
gives a holomorphic map to $\operatorname{Gr}(3,6)$.  On the
known open set it is
$\widehat F_\xi=K F_\xi$, by
\cref{rem:V:conjugation-parameter}. Applying the fixed real
integral change of marking $K^{-1}$ identifies it with the
canonical plane $F_\xi=(\sigma_b+\xi\alpha e)(P)$.
The identity theorem on the connected cover then identifies
these two maps everywhere.
The explicit graph is a rank-three holomorphic plane for
every finite $\xi$, including parameters for which it is not
real-transverse.  At infinity the same assertion is used with
$\xi=1/u$ on the punctured parameter disc.  All changes of
continuation marking are real integral and preserve the
condition $F\cap\bar F=0$.

\item Let $\xi_0$ be a finite boundary point, so
$\operatorname{Im}\xi_0=c_*=-\min h_b$.
The minimum occurs at $\lambda=\infty$ by \cref{night:P1.18}.
Continuity on its root chart gives regular $\lambda$ with
$h_b(\lambda)$ arbitrarily close to $\min h_b$.  Every sufficiently small full
parameter disc about $\xi_0$ therefore contains noncentral
points of
\[
 \operatorname{Im}\xi+h_b(\lambda)=0.
\]
This equation defines nonempty smooth real hypersurface
pieces in the complex surface $B$, since its derivative in
$\operatorname{Im}\xi$ is one.  A proper complex analytic
subset of $B$ has real dimension at most two; it cannot
contain those three-dimensional pieces.  Choose such a
point outside $E$.  Its fibre is a complex torus by the fourth
step, so its period plane is real-transverse.  The fifth step
identifies that plane with $F_\xi$, whereas the exact criterion
of \cref{night:P1.8} gives
$F_\xi\cap\overline{F_\xi}\ne0$ at the chosen point.
This is a contradiction.

At infinity fix a regular $\lambda$ and set
\[
 \xi=R-i h_b(\lambda),\qquad
 u=\frac1{R-i h_b(\lambda)},\qquad R\in\mathbb R.
\]
These are noncentral points of the same failure locus and
tend to $u=0$ as $|R|\to\infty$.  Varying $\lambda$ and $R$
again gives real hypersurface pieces in every punctured
$u$-disc, so one can choose a point outside $E$ and obtain the
same contradiction.  Under a finite ramified change of the
boundary coordinate, the parameter map has nonzero derivative
off the centre after shrinking.  The inverse images of these
hypersurface pieces remain nonempty real hypersurfaces there,
and the argument is unchanged.

The contradiction occurs on a noncentral compact smooth torus
fibre of the recovered marked map.  The proper reductions and
modifications preserve the prescribed open family.  Thus the
argument applies to every extension in
\cref{def:V:analytic-boundary-extension}, proving both assertions.
\end{enumerate}
\end{proof}




\begin{theorem}[Families over a punctured disc]\hypertarget{tag-02R6}{}\label{night:P1A.6}
\StatementRecord{night:P1A.6}{proved}{02R6}{\cref{night:P1A.3}, \cref{night:P1A.4}}
Let $p:\mathcal M\to\Delta^*$ be a proper holomorphic submersion
whose fibres are members of this fourfold family. Its classifying
parameter has an interior limit. After a base change of degree at most two, it is
isomorphic to the pullback of the constructed family by
a holomorphic map of the full disc to $\mathcal U$, and
therefore has a proper smooth filling. Such a family
cannot approach either parameter end.
\end{theorem}
\begin{proof}
\begin{enumerate}[label=\arabic*.]
\item At every punctured parameter, \cref{night:P1A.3} supplies local
miniversality and a local holomorphic classifying map to
$\mathcal U$. \Cref{night:P1A.4} makes this map unique on each geometric
point. Thus these local maps agree on overlaps and define
one holomorphic function $\xi:\Delta^*\to\mathcal U$.
There is no quotient by parameter translations, by \cref{night:P1A.4}.

\item The Cayley coordinate
\[
 z=\frac{\xi-i(c_*+1)}{\xi-i(c_*-1)}
\]
has $|z|<1$. It extends holomorphically across the puncture
by the removable-singularity theorem. Its value there has
modulus strictly less than one, since a boundary value
would force $z$ to be constant by the maximum-modulus
principle. The inverse Cayley map therefore extends
$\xi$ with value in $\mathcal U$.

\item Local identifications with the pulled-back constructed
family differ by its automorphisms. These form the constant
holomorphic group $\mathbb C\times\mathbb Z/2$ over the
parameter half-plane: the translation field and the
half-period in \cref{night:P1A.4} extend over the entire family.
The finite component torsor on $\Delta^*$ is killed by
the double cover. For a coordinate $z$ on $\Delta^*$, the functions $z$ and $z^{-1}$
confine the holomorphic hull of each compact set to a compact annulus;
$z$ also separates points and supplies a local coordinate. Thus
$\Delta^*$ is Stein. Cartan's Theorem~B
\cite[Expos\'e~19, \S5, pp.~19-14--19-15]{Cartan1952B_AB},
applied to its coherent structure sheaf, gives
$H^1(\Delta^*,\mathcal O)=0$, so the remaining additive
holomorphic torsor is a coboundary. Thus after that cover the local
identifications glue to a global isomorphism. Pulling
back the already constructed proper smooth family along
the extended classifying map supplies the stated filling.

\end{enumerate}
\end{proof}


\part{Integral monodromy and marked gluings}\label{pub:V:part-classification}
\section{Admissible data and local restrictions}\label{pub:V:sec-admissible}

\begin{definition}[Admissible data]\hypertarget{tag-PV18}{}\label{def:admissible}
\StatementRecord{def:admissible}{proved}{PV18}{}
Let $\Lambda=\mathbb Z^6$ and $k\in\{2,3\}$.  An \emph{admissible linear
tuple} is a tuple $(M_1,\ldots,M_k,M_0)$ of integral matrices with each
$M_i$ of finite order and with a nonzero fixed vector,
$M_0=(M_1\cdots M_k)^{-1}$, and $N=M_0-I$ satisfying $N^2=0$,
$\operatorname{rank}N=3$ and $\operatorname{im}N=\ker N$ as integral
sublattices of $\Lambda$ (the \emph{index-one condition}).  An
\emph{admissible datum} is an admissible linear tuple together with free
affine lifts $\tau_i(x)=M_ix+v_i/m_i$ of affine orders $m_i$, an
integral translation clutch $\kappa\in\Lambda$ and a primitive marking of
the Kuhn cusp.  Gluing the affine caps, the index-one Kuhn cap and the
torus bundle over the punctured sphere gives a closed smooth
eight-manifold, the \emph{gluing} of the datum.
\end{definition}




\begin{theorem}[Euler number of a Kuhn cusp]\hypertarget{tag-02MU}{}\label{maps:src:F.1.1}
\StatementRecord{maps:src:F.1.1}{proved}{02MU}{the definitions and constructions stated in the proof}


Let $g:M^8\to S^2$ be a proper triangulable map, a $T^6$-bundle
off finitely many values.  Suppose its unique non-quotient special
fibre is the rank-three Kuhn toroidal fibre with period-lattice
index $d$, and all its other special reduced fibres are free
finite quotients of $T^6$.  Then $\chi(M)=6d$.

\end{theorem}


\begin{proof}
\begin{enumerate}[label=\arabic*.]
\item The rank-three Kuhn fibre is the case $n=3$
of \cite[Corollary~8.8, arXiv v1, p.~59]{Yang2026}: the identical
height-one fan and the inclusion $\Gamma\subset\mathbb Z^3$
identify its index-$d$ quotient, and give central Euler number
$3!d=6d$.
\item The regular fibre has Euler number $\chi(T^6)=0$.
Each other reduced fibre is a free finite quotient of $T^6$,
so its Euler number is also zero.  Euler additivity for the
stated proper triangulable map therefore gives $\chi(M)=6d$.
\end{enumerate}
\end{proof}




\begin{proposition}[The diagonal obstruction]\hypertarget{tag-02MX}{}\label{maps:src:F.1.4}
\StatementRecord{maps:src:F.1.4}{proved}{02MX}{the definitions and constructions stated in the proof}


If $A\in\operatorname{GL}_n(\mathbb Z)$ has no eigenvalue one,
every affine map $x\mapsto Ax+b$ of $\mathbb R^n/\mathbb Z^n$
has a fixed point.  In particular, finite monodromies without eigenvalue one, such as the diagonal elliptic ones of orders four and six, admit no free affine lifts.

\end{proposition}


\begin{proof}
\begin{enumerate}[label=\arabic*.]
\item The real map $I-A$ is invertible.  A real solution of
$(I-A)x=b$ maps to a fixed point in the torus.  The diagonal
copies of the order-four and order-six elliptic rotations have
no eigenvalue one, so the conclusion applies to every translation
part.

\end{enumerate}
\end{proof}









\begin{proposition}[A complete finite-order filter]\hypertarget{tag-02N9}{}\hypertarget{tag-02R8}{}\label{maps:src:F.2.2}
\StatementRecord{maps:src:F.2.2}{proved}{02N9}{\cref{maps:src:F.1.4}}


The linear order of a finite matrix in $\operatorname{GL}_6(\mathbb Z)$
which admits a free affine lift belongs to
\begin{equation}
 \{1,2,3,4,5,6,8,10,12\}.
 \hypertarget{tag-02EZ}{}\label{maps:eq:F.3}
\end{equation}
There are 38 possible degree-six cyclotomic characteristic
polynomials with an eigenvalue one, and 19 of them commute
with some real complex structure.

\end{proposition}


\begin{proof}
\begin{enumerate}[label=\arabic*.]
\item Proposition \ref{maps:src:F.1.4} requires a nonzero one-eigenspace.  The other
cyclotomic factors have total degree at most five.  The elementary
identity
\[
 \frac{\phi(n)^2}{n}=
 \prod_{p^a\parallel n}p^{a-2}(p-1)^2\geq\frac12
\]
implies $n\leq50$ when $\phi(n)\leq5$.
The complete enumeration through 50 gives
$n=1,2,3,4,5,6,8,10,12$.  Enumerating their multiplicities
with total degree six and at least one factor $T-1$ gives 38
polynomials.  A commuting complex structure exists on a real
finite-order module exactly when the multiplicities of the
real eigenvalues $1$ and $-1$ are even; nonreal pairs already
carry such a structure.  Imposing those two parity conditions
leaves 19.  Taking least common multiples of their root orders
gives exactly \eqref{maps:eq:F.3}, also for the 38-polynomial list.


\end{enumerate}
\end{proof}




\begin{theorem}[The self-dual parity obstruction]\hypertarget{tag-02O1}{}\label{maps:src:F.3.1}
\StatementRecord{maps:src:F.3.1}{proved}{02O1}{\cref{maps:src:F.1.1}}


Suppose a compact smooth eight-manifold has the torus fibration
and the Kuhn cusp of \ref{maps:src:F.1.1} with index one, and its
other reduced fibres are free quotients of complex three-tori
by finite holomorphic affine actions.  If its rational
$H^1$ monodromy representation is self-dual, then it cannot
satisfy $b_1=b_2=0$.  In particular no admissible datum with symplectic monodromy has a three-connected gluing, for either $k=2$ or $k=3$ (or any number $k$ of finite
points).

\end{theorem}


\begin{proof}
\begin{enumerate}[label=\arabic*.]
\item Let $\mathcal V$ be the rank-six cohomology local system on
the punctured sphere and let $j$ be its inclusion.  On a finite
quotient fibre, averaging gives $R^1g_*\mathbb Q=(j_*\mathcal V)$
at that point.  At the Kuhn cusp, its normalization components
are simply connected and the degree-one branch complex is
the cochain complex of its dual $T^3$.  Specialization takes
its three deck classes to the covectors annihilating the
angular lattice, exactly the three-dimensional monodromy
invariant subspace.  Hence the same equality holds at the cusp.
It holds everywhere.

\item Put $a=\dim H^0(B,j_*\mathcal V)$ and
$b=\dim H^2(B,j_*\mathcal V)$.  The only differential that
can kill a class in $H^0(R^1)$ or in $H^2(R^0)$ is
$d_2:H^0(R^1)\to H^2(R^0)=\mathbb Q$.  If $b_1=b_2=0$,
it must be an isomorphism, so $a=1$.  Moreover
$H^1(B,j_*\mathcal V)$ must vanish because its position
$(1,1)$ has no possible incoming or outgoing differential.
Self-duality and compactly supported duality give $b=a=1$.

\item At the $i$th finite point, let $r_i$ be the dimension of the
fixed subspace in $H^1$.  The derivative is complex linear,
so its one-eigenspace has even real dimension.  Each $r_i$
is therefore even.  The cusp fixed dimension is three.
Euler additivity gives
\[
 \chi(B,j_*\mathcal V)=6(1-k)+3+\sum_{i=1}^k r_i,
\]
an odd integer.  Consequently
\[
 \dim H^1(B,j_*\mathcal V)=a+b-\chi(B,j_*\mathcal V)
\]
is odd and cannot vanish.  This is the required contradiction.
A preserved nondegenerate alternating form identifies a
representation with its dual, so symplectic monodromy is
included.  The explicit form
$(6\gamma\wedge\delta+u\wedge w)\oplus\omega_E$
checks this hypothesis for the fibre-product example of Part~\ref{pub:V:part-resolved}.

\end{enumerate}
\end{proof}




\begin{proposition}[The remaining rank patterns]\hypertarget{tag-02O4}{}\label{maps:src:F.3.2}
\StatementRecord{maps:src:F.3.2}{proved}{02O4}{the definitions and constructions stated in the proof}


If a possibly non-self-dual holomorphic datum with the same
local models had $b_1=b_2=0$, put
$b=\dim H^0(B,\mathcal V^*)$.  For $k=2$, the only possible
unordered pairs $(r_1,r_2;b)$ are
\[
 (2,2;0),\qquad(2,4;2).
\]
For $k=3$, they are
\begin{equation}
 (2,2,6;0),\ (2,4,4;0),\ (2,4,6;2),\ (4,4,4;2).
 \hypertarget{tag-02F7}{}\label{maps:eq:F.4}
\end{equation}

\end{proposition}


\begin{proof}
\begin{enumerate}[label=\arabic*.]
\item The preceding Leray argument still gives $a=1$ and vanishing
$H^1(B,j_*\mathcal V)$.  Euler additivity therefore gives
$b=\sum r_i+3-1-6(k-1)$.  Each $r_i$ belongs to
$\{2,4,6\}$ because a free affine lift requires an eigenvalue
one and complex linearity makes its multiplicity even.
Also $0\leq b\leq\min(3,r_1,\ldots,r_k)$, by inclusion in
each local invariant subspace.  Enumerating these three-value
lists gives exactly \eqref{maps:eq:F.4}.

\end{enumerate}
\end{proof}






\section{Finite spectra and rational extensions}\label{pub:V:sec-spectra}



\begin{theorem}[The orbifold quotient]\hypertarget{tag-02NY}{}\label{maps:src:F.2.7}
\StatementRecord{maps:src:F.2.7}{proved}{02NY}{the definitions and constructions stated in the proof}


For the torus fillings of an admissible datum, with affine
multiplicities $m_1,\ldots,m_k$, the fundamental group of the
total space has a canonical surjection onto
\begin{equation}
 \langle x_1,\ldots,x_k\mid
       x_i^{m_i}=1,\ x_1\cdots x_k=1\rangle.
 \hypertarget{tag-02EX}{}\label{maps:eq:F.27}
\end{equation}
Consequently a simply connected example with $k=2$ requires
$\gcd(m_1,m_2)=1$.  There is no simply connected example with
$k=3$ and $m_1,m_2,m_3\geq2$.  An affine multiplicity-one
point is an ordinary smooth point and reduces to the problem
with one fewer multiple fibre.

\end{theorem}


\begin{proof}
\begin{enumerate}[label=\arabic*.]
\item Over the complement of small discs at the special values,
the torus bundle projects surjectively on fundamental groups
to the punctured base.  Quotient its fundamental group by the
normal closure of all regular-fibre loops.  At a finite
filling its root cover is a smooth torus bundle over a disc,
and its free deck group is cyclic of order $m_i$.
The covering exact sequence shows that the quotient of the
filling group by its fibre subgroup is precisely this cyclic
group.  The boundary base meridian maps to its generator.
Thus adjoining that filling imposes exactly
$x_i^{m_i}=1$ after the fibre loops have been killed.
At the index-one Kuhn cusp the fundamental group is generated
by the three deck loops of its dual torus; these are regular
fibre loops.  Its base meridian contracts.  Hence its filling
imposes the ordinary meridian relation and leaves no extra
generator in this quotient.  The base relation gives the
product relation in \eqref{maps:eq:F.27}.  Van Kampen proves that \eqref{maps:eq:F.27} is
the resulting quotient, in particular the asserted surjection.

\item For two points, substitute $x_2=x_1^{-1}$; the quotient is
$\mathbb Z/\gcd(m_1,m_2)$.
For three points of orders at least two it is a nontrivial
triangle group.  Nontriviality can be seen directly.  Take a geodesic triangle
with angles $\pi/m_1,\pi/m_2,\pi/m_3$ in the sphere, Euclidean
plane or hyperbolic plane according as their sum is greater
than, equal to or less than $\pi$.
In the spherical case the possible orders are, after sorting,
$(2,2,n),(2,3,3),(2,3,4),(2,3,5)$, and the corresponding
spherical triangles exist.  The Euclidean and hyperbolic
angle conditions give the other triangles directly.
If $r_1,r_2,r_3$ are reflections in the opposite sides,
the rotations $r_2r_3,r_3r_1,r_1r_2$ have the stated
orders and product one.  They therefore give a nontrivial
representation of \eqref{maps:eq:F.27} into the orientation-preserving
isometry group.  Its quotient cannot be trivial.
Finally a free cyclic action of order one is the identity
construction on a root disc of order one, hence an ordinary
torus bundle filling and not a multiple fibre.

\end{enumerate}
\end{proof}




\begin{proposition}[Automatic local complex compatibility]\hypertarget{tag-02N4}{}\label{maps:src:F.2.15}
\StatementRecord{maps:src:F.2.15}{proved}{02N4}{\cref{maps:src:F.2.7}, \cref{maps:src:F.1.4}}


Every admissible datum with simply connected gluing and two finite affine fillings automatically has finite linear actions commuting
with real complex structures.

\end{proposition}


\begin{proof}
\begin{enumerate}[label=\arabic*.]
\item Let $m_i$ be the order of the free cyclic affine action
defining the $i$th filling.  Its linear order divides $m_i$.
By \ref{maps:src:F.2.7} the integers $m_1,m_2$ are coprime, so at least one
of them is odd.  The corresponding linear determinant is
one, because a finite real matrix of odd order cannot have
determinant minus one.  The relation
$M_1M_2M_0=I$ and the unipotence of $M_0$ imply
$\det M_1\det M_2=1$, so both determinants are one.

\item A finite-order real matrix is semisimple.  Its nonreal
eigenvalues occur in conjugate pairs, and its only possible
real eigenvalues are $1$ and $-1$.  Determinant one says
that the $-1$ eigenspace has even real dimension.  Since the
total dimension is six, the $1$ eigenspace also has even
real dimension.  Choose a complex structure on each of
these real eigenspaces by pairing a real basis.  On each
nonreal isotypical real summand, multiplication by $i$ in
its complex eigenspace supplies a commuting complex
structure.  Their direct sum commutes with $M_i$.
The fixed eigenspace is nonzero because the affine action
is free, as in \ref{maps:src:F.1.4}, so its even dimension is at least two.

\item Each chosen structure makes $\mathbb R^6/\mathbb Z^6$ a
complex torus and the given affine map holomorphic; thus it
supplies precisely the local complex-compatibility
hypothesis used in the subsequent reductions.

\end{enumerate}
\end{proof}

\begin{proposition}[The only first-cohomology pattern for two points]\hypertarget{tag-02O0}{}\label{maps:src:F.2.9}
\StatementRecord{maps:src:F.2.9}{proved}{02O0}{\cref{maps:src:F.3.1}, \cref{maps:src:F.2.15}}


For a simply connected two-multiple-fibre datum with $b_2=0$,
the rational
cohomology system $V$ has one invariant, no coinvariants,
and fixed rank two at each finite generator.  Its quotient
$W=V/\mathbb Qe$ by the invariant line has neither invariants
nor coinvariants.  If one finite characteristic polynomial
on $W$ is $(T-1)\Phi_d(T)$ with $\deg\Phi_d=4$, then $W$
is irreducible over $\mathbb Q$.

\end{proposition}


\begin{proof}
\begin{enumerate}[label=\arabic*.]
\item Write $a=\dim V^\Gamma$, $b=\dim(V^*)^\Gamma$ and $f_i$
for the finite fixed ranks.  The Leray argument in the proof of Theorem \ref{maps:src:F.3.1} gives $a=1$,
vanishing first parabolic cohomology, and
$f_1+f_2=4+b$.  By \ref{maps:src:F.2.15} each $f_i$ is positive and even, so $b$ is
even.  The maximal square-zero cusp has
$\ker N=\operatorname{im}N$.  Therefore every globally
fixed vector belongs to $\operatorname{im}N$, whereas every
globally fixed covector kills $\operatorname{im}N$.
The global invariant and coinvariant spaces pair to zero.

\item Restriction of a globally invariant covector to the fixed
space of a finite generator is injective: averaging that
generator is the projection onto its fixed space, and the
covector is unchanged by averaging.  Its restriction kills
the $a$ global invariants in that space.  Hence
$a+b\leq f_i$ for each $i$.
If $b\geq2$ is even, the even ranks satisfy
$f_i\geq b+2$, so $4+b=f_1+f_2\geq2b+4$, a contradiction.
Thus $b=0$ and $f_1=f_2=2$.

\item The quotient $W$ has no invariant covectors, since its dual
embeds into $V^*$.  If it had a fixed line, its preimage in
$V$ would be an extension of the trivial line by the trivial
line.  Every finite generator on this two-dimensional
space would be unipotent of finite order, hence identity.
This would give two invariants in $V$, a contradiction.
Thus $W$ also has no invariant vectors.

\item If one finite characteristic polynomial is
$(T-1)\Phi_d(T)$ with the second factor irreducible of degree
four, every proper nonzero rational submodule has dimension
one or four.  In dimension one that generator acts as the
identity.  Since the product of the two finite generators
is unipotent, their induced one-dimensional characters are
inverse; hence the other generator also acts trivially.
This would be an invariant line.  In dimension four the
one-dimensional quotient is trivial by the same argument,
giving a coinvariant.  Both have just been excluded.

\end{enumerate}
\end{proof}

\begin{proposition}[Exclusion of an involutory finite part]\hypertarget{tag-02NZ}{}\label{maps:src:F.2.8}
\StatementRecord{maps:src:F.2.8}{proved}{02NZ}{\cref{maps:src:F.2.7}, \cref{maps:src:F.2.15}, \cref{maps:src:F.3.1}, \eqref{maps:eq:F.3}, \cref{maps:src:F.2.9}}


A simply connected datum with $b_2=0$, $k=2$ and the stated
finite affine fillings has neither finite linear
monodromy of order one or two.  Its unordered pair of finite
linear orders must therefore belong to
\begin{equation}
 \{(3,4),(3,5),(3,8),(3,10),(4,5),
       (5,6),(5,8),(5,12)\}.
 \hypertarget{tag-02EY}{}\label{maps:eq:F.28}
\end{equation}
The affine multiplicities may still be larger than these
linear orders.

\end{proposition}


\begin{proof}
\begin{enumerate}[label=\arabic*.]
\item The linear orders divide the affine multiplicities and are
therefore coprime by \ref{maps:src:F.2.7}.  The fixed-rank calculation
of \ref{maps:src:F.2.9}, which uses only \ref{maps:src:F.3.1} and the
evenness in \ref{maps:src:F.2.15}, gives fixed rank two at each finite
generator.  An identity generator would make the product finite and
unipotent, hence the identity, contrary to the rank-three cusp.

\item Suppose the cohomology generators are $A,B$ with $A^2=I$.
The order of $B$ is odd.  The $(-1)$-eigenspace of $A$ has dimension
four, whereas $\ker(AB-I)$ has dimension three.  Their intersection
in the six-dimensional space contains a nonzero vector $v$.
Then $ABv=v$ and $Av=-v$ imply
\[
 Bv=Av=-v,
\]
which contradicts the odd order of $B$.  Interchanging the two
generators proves the same assertion for either finite site.

\item Delete one and two from the complete finite-order list
\eqref{maps:eq:F.3} and enumerate its coprime unordered pairs.
They are exactly the eight pairs in \eqref{maps:eq:F.28}.
Only the divisibility of the linear orders into the affine
multiplicities has been used.
\end{enumerate}
\end{proof}







\begin{lemma}[Primitive exterior-square comparison]
\hypertarget{tag-PV07}{}\label{pub:V:primitive-exterior-square}
\StatementRecord{pub:V:primitive-exterior-square}{proved}{PV07}{}
Let $E=\mathbb C^4$ have symplectic basis $e_1,e_2,f_1,f_2$,
put $\omega=e_1\wedge f_1+e_2\wedge f_2$, and give
$E_0=\omega^\perp\subset\bigwedge^2E$ the wedge pairing.
The morphism
$\rho:\operatorname{Sp}(E)\to\operatorname{SO}(E_0)$
is surjective with kernel $\{I,-I\}$
\cite[\S7.3, p.~129]{pubLieGroups2016}.
A unipotent element of Jordan type $(2,2,1)$ has as its two
lifts a symplectic transvection and its negative.
For a semisimple lift with eigenvalues $x,x^{-1},y,y^{-1}$,
the eigenvalues on $E_0$ are
\[
 1,\quad xy,\quad x/y,\quad y/x,\quad(xy)^{-1}.
\]
\end{lemma}
\begin{proof}
We use the cited primitive exterior-square morphism and compute
the two comparisons needed here.
The transvection fixing $e_1,e_2,f_2$ and sending $f_1$ to
$f_1+e_1$ acts on the primitive basis
\[
 e_1\wedge e_2,\quad e_1\wedge f_2,\quad e_2\wedge f_1,
 \quad f_1\wedge f_2,\quad e_1\wedge f_1-e_2\wedge f_2.
\]
Its difference from the identity sends the third vector to the
negative of the first and the fourth to the second, and kills
the others. Its Jordan type is $(2,2,1)$.
Conversely, an orthogonal unipotent $I+N$ of this type has
$N^2=0$ and $N^*=-N$, so its image is an isotropic plane.
Choose that plane, its dual isotropic plane, and the remaining
nondegenerate line. In this decomposition $N$ is determined by
a nonzero alternating form on the dual plane. All such forms
are equivalent under its general linear group; extending the
change by the inverse transpose on the image plane gives an
orthogonal change of basis. Its determinant is one, after the
choice of sign on the residual line. Thus these unipotents form
one special-orthogonal conjugacy class. Surjectivity of $\rho$
lifts this conjugacy, and its kernel gives the two stated lifts.
Finally, taking pairwise products of the four eigenvalues and
removing the invariant bivector gives the displayed five
eigenvalues.
\end{proof}

\begin{theorem}[Exclusion of three further order pairs]\hypertarget{tag-02MZ}{}\label{maps:src:F.2.10}
\StatementRecord{maps:src:F.2.10}{proved}{02MZ}{\cref{maps:src:F.2.9}, \eqref{maps:eq:F.31}, \cref{maps:src:F.2.8}, \cref{maps:src:F.3.1}}


No simply connected datum with $k=2$, $b_2=0$, and finite linear orders
$(3,5)$, $(3,8)$ or $(3,10)$ has the required maximal
square-zero cusp.  Among such data, the
remaining necessary unordered linear order pairs are
\begin{equation}
 \{(3,4),(4,5),(5,6),(5,8),(5,12)\}.
 \hypertarget{tag-02F0}{}\label{maps:eq:F.30}
\end{equation}
When one order is five, the other finite characteristic
polynomial cannot be $(T-1)^2\Phi_4(T)^2$ or
$(T-1)^2\Phi_6(T)^2$.

\end{theorem}


\begin{proof}
\begin{enumerate}[label=\arabic*.]
\item For each proposed pair, \ref{maps:src:F.2.9} gives an
irreducible rational rank-five quotient $W$.  Its cusp has Jordan
type $(2,2,1)$.  One finite polynomial is
$(T-1)\Phi_3(T)^2$ and the other is $(T-1)\Phi_d(T)$ for
$d=5,8,10$.  For the two additional assertions, replace
$\Phi_3^2$ by $\Phi_4^2$ or $\Phi_6^2$, with $d=5$.
In all these cases a commuting endomorphism of $W$ restricts to a
rational scalar on the one-dimensional fixed space of the second
generator.  This restriction is injective: a nonzero endomorphism
of an irreducible representation is invertible and cannot vanish
on that line.  Consequently
\[
 \operatorname{End}_{\Gamma}(W)=\mathbb Q.
\]

\item Consider $\mathcal H=\operatorname{Hom}(W,W^*)$.
The local fixed dimensions are $9,5,13$.  At the repeated finite
site, pairing reciprocal eigenspaces of dimensions $1,2,2$ gives
$1+4+4=9$.  At the distinct finite site it gives five.
For the cusp, the ordered pairs of blocks of sizes $2,2,1$
contribute the sum of their minimum sizes, namely thirteen.
The invariant-sheaf Euler formula therefore gives
\begin{equation}
 \chi(\mathbb P^1,j_*\mathcal H)=-25+9+5+13=2.
 \hypertarget{tag-02F1}{}\label{maps:eq:F.31}
\end{equation}
Here degree zero is $\operatorname{Hom}_{\Gamma}(W,W^*)$,
and degree two is dual to
$\operatorname{Hom}_{\Gamma}(W^*,W)$.  At least one is nonzero,
so irreducibility gives a nondegenerate invariant bilinear form
$J:W\longrightarrow W^*$, inverting the map if necessary.
Its alternating part is equivariant; if nonzero it would be
invertible, which is impossible in odd dimension.  Thus $J$ is
symmetric.

\item An invariant bivector is an alternating equivariant map
$K:W^*\longrightarrow W$.  The comparison map
\[
 K\longmapsto KJ\in\operatorname{End}_{\Gamma}(W)
\]
makes it a scalar, so $K$ is a scalar multiple of the symmetric
map $J^{-1}$.  Hence $K=0$.  Applying the same argument to an
alternating map $W\longrightarrow W^*$ gives
\[
 (\Lambda^2W)^\Gamma=(\Lambda^2W^*)^\Gamma=0.
\]
Both degree zero and degree two of $j_*\Lambda^2W$ vanish.
Its Euler characteristic must therefore be nonpositive.

\item The repeated finite spectrum has four fixed bivectors and
the distinct spectrum has two.  At the cusp, write the unipotent
module as two blocks of size two and one block of size one.
Their exterior square has two trivial blocks from the individual
wedges, blocks of sizes three and one from the product of the
two size-two blocks, and two further size-two blocks.  Its fixed
rank is six.  Consequently
\[
 \chi(\mathbb P^1,j_*\Lambda^2W)=-10+4+2+6=2,
\]
contradicting the preceding step.  This proves all five stated
exclusions.  Removing the three excluded order pairs from
\eqref{maps:eq:F.28} leaves exactly \eqref{maps:eq:F.30}.
\end{enumerate}
\end{proof}




\begin{proposition}[Rational uniqueness of the extension]\hypertarget{tag-02N0}{}\label{maps:src:F.2.11}
\StatementRecord{maps:src:F.2.11}{proved}{02N0}{\cref{maps:src:F.2.9}, \eqref{maps:eq:F.32}}


Let $W$ be a rational rank-five representation of a
twice-punctured affine line, with two finite generators
having fixed rank one and product with unipotent Jordan type
$(2,2,1)$.  Suppose $W$ has neither invariants nor
coinvariants.  Up to rescaling the invariant line, there is
exactly one rational extension
$0\to\mathbb Qe\to V\to W\to0$ in which the two finite
generators retain their finite orders and the cusp has
square-zero logarithm of rank three.  Thus, for the two-point data of \ref{maps:src:F.2.9}, the remaining rational
representation problem is exactly the problem of the
rank-five quotients.  

\end{proposition}


\begin{proof}
\begin{enumerate}[label=\arabic*.]
\item Write the finite generators of $W$ as $S_1,S_2$ and the
extension generators as
\[
 A_i=\begin{pmatrix}1&a_i\\0&S_i\end{pmatrix}.
\]
If $S_i$ has order $m_i$, retaining that finite order is equivalent
to $a_i(I+S_i+\cdots+S_i^{m_i-1})=0$.  Averaging identifies this
condition with $a_i\in\operatorname{row}(S_i-I)$, a space of
dimension four.  A change of splitting by $b\in W^*$ sends the
two rows to $a_i+b(S_i-I)$.  This coboundary map is injective
because $W$ has no invariant covectors.  Thus the extension
space subject to the two finite conditions has dimension $8-5=3$.

\item Put $U=S_1S_2=I+N$, so $N^2=0$ and $\operatorname{rank}N=2$.
The product $A_1A_2$ has upper row $c=a_1S_2+a_2$.
Changing the splitting sends $c$ to $c+bN$.  Hence cusp restriction
is the explicit map
\[
 [(a_1,a_2)]\longmapsto[c]\in W^*/\operatorname{row}N.
\]
It is an isomorphism.  Indeed, a class in its kernel can be
represented with $c=0$, so $a_2=-a_1S_2$.  The finite conditions
then say
\[
 a_1\in\operatorname{row}(S_1-I)
              \cap\operatorname{row}(S_2-I).
\]
This intersection has dimension three: the two four-dimensional
row spaces span $W^*$ because $W$ has no common fixed vector.
The rows $b$ satisfying $bN=0$ also form a three-dimensional
space.  Their map $b\mapsto b(S_1-I)$ lands in the displayed
intersection, since $bU=b$ gives
$b(S_1-I)S_2=b-bS_2$.  It is injective, because
$bS_1=b$ and $bU=b$ together give $bS_2=b$, forcing $b=0$.
It is therefore onto.  Every pair with $c=0$ is a coboundary,
proving injectivity of cusp restriction.  Both its source and
target have dimension three, so it is an isomorphism.

\item The square of the lifted product minus identity vanishes
exactly when $cN=0$.  Under this restriction isomorphism the
square-zero extensions form
\begin{equation}
 \frac{\operatorname{ann}(\operatorname{im}N)}
      {\operatorname{row}N},
 \hypertarget{tag-02F2}{}\label{maps:eq:F.32}
\end{equation}
of dimension $3-2=1$.  The lifted nilpotent has the rows of $N$
together with the row $c$, so its rank is three exactly when
$c\notin\operatorname{row}N$.  These are precisely the nonzero
classes in \eqref{maps:eq:F.32}.  Rescaling the invariant line
rescales this one-dimensional rational space and is transitive
on its nonzero classes.  This proves the asserted existence and
uniqueness with the stated finite orders and maximal cusp.
\end{enumerate}
\end{proof}




\begin{theorem}[An integral obstruction for a rational $(5,12)$ family]\hypertarget{tag-02N1}{}\label{maps:src:F.2.12}
\StatementRecord{maps:src:F.2.12}{proved}{02N1}{\eqref{maps:eq:F.33}, \cref{maps:src:F.2.1}, \cref{maps:src:F.2.7}, \cref{night:P1.11}}


Let a rational rank-six cohomology representation have generators
\[
 A=\begin{pmatrix}1&0\\0&A_5\end{pmatrix},\qquad
 B=\begin{pmatrix}1&c\\0&B_5\end{pmatrix},\qquad
 c=(0,0,0,1,0),
\]
where
\begin{equation}
 A_5=\begin{pmatrix}
 1&0&1&0&0\\1&0&0&0&0\\-1&0&0&1&1\\
 1&-1&0&-1&0\\0&1&0&0&0
 \end{pmatrix},\qquad
 B_5=\begin{pmatrix}
 0&1&0&0&0\\0&0&0&0&1\\1&-2&0&0&0\\
 0&1&0&-1&-1\\0&0&1&1&2
 \end{pmatrix}.
 \hypertarget{tag-02F3}{}\label{maps:eq:F.33}
\end{equation}
For every invariant integral lattice with an index-one maximal
square-zero cusp, and every free affine filling of multiplicities
$m_1,m_2$ divisible respectively by five and twelve, a simply
connected total space has
\begin{equation}
 \operatorname{Tor}_3 H_2
       \simeq\mathbb Z/3^{v_3(m_2)}.
 \hypertarget{tag-02F4}{}\label{maps:eq:F.34}
\end{equation}
In particular none of these data is three-connected.
The standard lattice nevertheless has a primitive cusp,
free affine fillings of orders five and twelve, and a simply
connected total space with rational Betti vector
$(1,0,0,0,4,0,0,0,1)$.

\end{theorem}


\begin{proof}
\begin{enumerate}[label=\arabic*.]
\item Put $R=\mathbb Z_{(3)}$ and use the rational coordinates of
\eqref{maps:eq:F.33}.  Let $\mathcal P_R$ be the algebra of blocks
$\left(\begin{smallmatrix}a&b\\0&D\end{smallmatrix}\right)$,
where $a\in R$, $b\in R^{1\times5}$ and $D\in\operatorname{Mat}_5(R)$.
We first determine the integral algebra generated by $A,B$.
The following ordered words have $31$ coordinate columns in the order
$(0,j)$, $0\leq j\leq5$, followed by $(i,j)$,
$1\leq i,j\leq5$:
\[
\begin{gathered}
I,A,B,A^2,AB,BA,B^2,A^3,A^2B,ABA,AB^2,BA^2,\\
BAB,B^2A,B^3,A^3B,A^2BA,A^2B^2,ABA^2,AB^2A,\\
AB^3,BA^3,BA^2B,BAB^2,B^2A^2,B^3A,A^3BA,\\
A^3B^2,A^2BA^2,BA^3B,B^2A^3.
\end{gathered}
\]
Their coordinate determinant is $-32$.  Its adjugate divided by
$-32$ therefore expresses every elementary matrix of
$\mathcal P_R$ as an $R$-linear combination of these words.
Since $A,B\in\mathcal P_R$, this proves
$R[A,B]=\mathcal P_R$ with the coordinate map.

For a stable lattice $L_R$, the elementary idempotents split its
six coordinate fractional ideals; the matrices $E_{ij}$ for
$i,j\geq1$ make the last five equal.  Thus
\[
 L_R=3^aRe_0\oplus3^bR^5,\qquad b\geq a,
\]
where the inequality follows from the matrices $E_{0j}$.
The matrix $N=AB-I$ has square zero and nonzero Smith factors
$1,1,1$, while $A_5B_5-I$ has factors $1,1$.
In this basis of $L_R$, its upper-right row is multiplied by
$3^{b-a}$.  For $b>a$ its reduction modulo three has rank two,
contradicting the primitive rank-three cusp.  Hence $b=a$.
The scalar map $3^{-a}I$ identifies $L_R$ with $R^6$
and intertwines every exterior power of the two generators.
The given cusp operator has the same elementary divisors:
$BA-I=A^{-1}(AB-I)A$, and
$(BA)^{-1}-I=-(BA)^{-1}(BA-I)$.
This proves the asserted local marking for every stable lattice.


\item Write $v_i=g_i^{m_i}$ for the finite power vectors.
They are fixed by the homology matrices $M_i$.
The base-orbifold quotient in \cref{maps:src:F.2.7} shows that
simple connectivity forces $\gcd(m_1,m_2)=1$.
After the cusp meridian relation the two base generators are
inverse up to a fibre translation.  Their actions on the abelian
fibre lattice are therefore inverse.  Their linear orders five
and twelve force both actions to be the identity in this quotient.
The exact Smith factors of $[A^t-I\ B^t-I]$ are five units,
with first row zero, so its fibre coinvariant module over $R$
is $R$ through the primitive character $\ell=x_0$.

Set $a_i=\ell(v_i)$.  An arbitrary cusp clutch contributes a scalar
$\delta$ to $g_1+g_2=\delta z$.  The remaining relation matrix is
\[
 \begin{pmatrix}m_1&0&-a_1\\0&m_2&-a_2\\1&1&-\delta\end{pmatrix}.
\]
Its determinant is a unit for a simply connected total space.
Since $3\mid m_2$ and $3\nmid m_1$, reduction modulo three makes
$a_2$ a unit in $R$.  This conclusion includes arbitrary clutch
translations and twists.

\item We determine the second finite-cap restriction over $R$.
Choose a primitive integral invariant character $\alpha$ whose
$R$-span is $\ell$.  The stable sublattice
\[
 \Lambda'=\mathbb Z v_2\oplus\ker\alpha
\]
has index $|\alpha(v_2)|$, prime to three.  It gives a finite torus cover; the affine action lifts because it preserves
$v_2$, preserves $\ker\alpha$, and translates by $v_2/m_2$.
The corresponding reduced cap $F'_2$ covers $F_2$ with the same
degree.  In its basis beginning with $v_2$, the action is
\[
 (u,y)\longmapsto(u+1/m_2,Dy).
\]
Thus $F'_2$ is the mapping torus of $D:T^5\to T^5$, and its
regular six-torus is the $m_2$-fold cover in the circle direction.
The coefficient change has determinant a unit in $R$; under the
scalar marking of item~1, the vertical cohomology action $D^t$
is the displayed $B_5$.

The matrices $D^t-I$ and $\bigwedge^2D^t-I$ have no nonunit
three-primary Smith factors.  Their Wang sequences therefore make
$H^q(F'_2;R)$ free for $q\leq3$.  Pullback to the generic torus
is injective in these degrees, because transfer kills its kernel
by $m_2$ and the source is free.  Deck transformations of
$F'_2\to F_2$ act by translations on the generic torus cover.
Their action there is trivial on cohomology, so the injectivity
just proved makes their action on $H^q(F'_2;R)$ trivial as well.
Transfer for this prime-to-three cover now identifies
$H^q(F_2;R)$ with $H^q(F'_2;R)$, compatibly with their restriction maps.  In particular the degree-two and degree-three
restrictions have zero kernel.

Put $T=D^t$.  Its polynomial is $(x-1)\Phi_{12}(x)$, and
\[
 E=\Phi_{12}(T)=T^4-T^2+I,\qquad E^2=E,\qquad
 \sum_{j=0}^{11}T^j=12E.
\]
The image of this integral rank-one idempotent is the primitive
invariant line.  Hence the norm for every stated $m_2$ is $m_2E$.
The Wang sequence in degree two has quotient
$\operatorname{Inv}(\bigwedge^2T)$ and kernel
$\operatorname{coker}(T-I)$.  Restriction to the $m_2$-fold cover
projects identically onto the former group.  On the latter it is
the horizontal wedge with the norm: a horizontal cell lifts to
$m_2$ successive transported prisms.
Consequently the restriction image $J_2$ projects onto the vertical
invariants, and its intersection with the horizontal invariant
summand is $e_0\wedge m_2\operatorname{Inv}(T)$.
The horizontal inclusion induces the isomorphism
\[
 \frac{\operatorname{Inv}(T)}{m_2\operatorname{Inv}(T)}
 \ \xrightarrow{\ \sim\ }\\
 \frac{\operatorname{Inv}(\bigwedge^2V_R)}{J_2}
 \ \cong\ R/3^rR.
\]
At the first finite cap, $m_1$ is a unit
in $R$, so transfer and averaging give the full invariant image
and zero kernel.  At the cusp, \cref{night:P1.11} gives injective
restrictions onto the invariant lattices in degrees one, two
and three, by dualizing its saturated-kernel collapse.

\item The integral specialization-sheaf approach is used in
\cite[Lemma 7.25 and Propositions 7.26--7.27, pp.~61--63]{AC26}
and \cite[section ``The integral homology'']{E26}.
Here are the coefficient maps and calculations for the present lattice.  Put $L_q=\bigwedge^qV_R$, $A_q=\bigwedge^qA$,
$B_q=\bigwedge^qB$, and let $j$ include the thrice-punctured base.
For the full invariant-extension sheaf $j_*L_q$ the oriented
cellular complex is
\[
\begin{gathered}
 J_{1,q}\oplus J_{2,q}\oplus J_{0,q}
 \xrightarrow{d_0}L_q^2\xrightarrow{\Delta_q}L_q,\\
 J_{1,q}=\ker(A_q-I),\quad J_{2,q}=\ker(B_q-I),\quad
 J_{0,q}=\ker(B_qA_q-I),\\
 d_0(j_1,j_2,j_0)=(j_1-j_0,j_2-A_qj_0),\\
 \Delta_q(a,b)=(A_q-I)a+(B_q-I)b.
\end{gathered}
\]
These are the base attachment maps: the two cusp paths
carry $j_0$ and $A_qj_0$, whose combined boundary is
$(B_qA_q-I)j_0=0$.  Primitive invariant bases and a unimodular
kernel basis of $\Delta_q$ give the following Smith forms:
\[
\begin{array}{c|c|c|c|c}
q&\operatorname{rank}H^0&\operatorname{rank}H^1&
 \operatorname{SNF}(d_0\text{ in }\ker\Delta_q)&
 \operatorname{SNF}(\Delta_q)\\\hline
1&1&0&1^6&1^5,2\\
2&0&0&1^{15}&1^{15}\\
3&0&3&1^{16},2&1^{19},2.
\end{array}
\]
Here $1^s$ denotes $s$ entries equal to one.  The row-coordinate
matrices and their unimodular HNF transformations are specified by the preceding formulas.
In particular $H^0(j_*L_2)=H^1(j_*L_2)=0$,
$H^2(j_*L_1)=0$ over $R$, and $H^0(j_*L_3)=0$.

Write $\mathcal R_q=R^qf_*R$.  The stalk/fibre
identification is proper base change
\cite[Tag 09V5, Lemma 20.18.1]{StacksProject}.
The local maps of item~3 give
\[
 0\longrightarrow\mathcal R_2\longrightarrow j_*L_2
 \longrightarrow Q\longrightarrow0,
 \qquad Q_{p_2}=R/3^rR,
\]
with no other support.  Hence its connecting homomorphism gives
$H^1(\mathcal R_2)=R/3^rR$.
The degree-one quotient is also point supported, so
$H^2(\mathcal R_1)=H^2(j_*L_1)=0$.
The degree-three sheaf map is injective, so
$H^0(\mathcal R_3)=0$.

\item The only total-degree-three term in the Leray sequence
\cite[Tag 01F2, Lemma 20.13.4]{StacksProject} is therefore
$E_2^{1,2}=R/3^rR$.  It has no possible differential because
the base has cohomological dimension two.  Thus
$H^3(M;R)=R/3^rR$.
Universal coefficients
\cite[Theorem 3.2, p.~195]{Hatcher2002} identify this torsion
module with $\operatorname{Ext}_R(H_2(M;R),R)$: the quotient
$\operatorname{Hom}_R(H_3(M;R),R)$ is free and must vanish.
The elementary divisor decomposition over the PID $R$ proves
\eqref{maps:eq:F.34}.

\item For the standard lattice take
\[
 v_1=(3,0,0,0,0,0)^t,\qquad
 v_2=(-7,7,-7,7,0,7)^t,\qquad (m_1,m_2)=(5,12).
\]
They are fixed by $A^t,B^t$ respectively.
Their $\ell$-values $3,-7$ are units modulo the respective
multiplicities; applying $\ell$ to a fixed-point equation proves
that every nonidentity affine power is free.
The primitive cusp in item~1 has the integral marking of
\cref{maps:src:F.2.1}.  Glue the two finite quotient caps and
the index-one Kuhn cap to the marked torus bundle over a pair
of pants, with zero additional clutch translation.  The finite
boundary identification is the prescribed logarithmic fibre
translation; the cusp is attached in that integral marking.
These local models give a closed smooth oriented
manifold.  Its remaining van Kampen matrix is
\[
 \begin{pmatrix}5&0&-3\\0&12&7\\1&1&0\end{pmatrix},
\]
with determinant of absolute value one.  It is simply connected.

Over $\mathbb Q$ the finite images are full invariants.
Item~4 gives vanishing $H^1(j_*V_{\mathbb Q})$,
$H^0(j_*\bigwedge^2V_{\mathbb Q})$,
$H^1(j_*\bigwedge^2V_{\mathbb Q})$ and
$H^2(j_*V_{\mathbb Q})$, and the degree-three invariant group
is zero.  The five-term Leray sequence and $H^1(M;\mathbb Q)=0$
make the map from the one-dimensional $H^0(\mathcal R_1)$ to
$H^2(\mathbb P^1;\mathbb Q)$ an isomorphism.
Thus $b_2=b_3=0$.
Euler additivity gives $\chi(M)=6$ from the Kuhn cusp;
all other fibres have Euler number zero.
Poincar\'e duality therefore gives
$(1,0,0,0,4,0,0,0,1)$, as claimed.
\end{enumerate}
\end{proof}




\begin{theorem}[Two further rational families excluded integrally]\hypertarget{tag-02N2}{}\label{maps:src:F.2.13}
\StatementRecord{maps:src:F.2.13}{proved}{02N2}{\cref{maps:src:F.2.11}, \cref{maps:src:F.2.7}, \cref{night:P1.9}, \cref{night:P1.11}; \cite[Tag~09V5, Lemma~20.18.1; Tag~01F2, Lemma~20.13.4]{StacksProject}; \cite[Theorem~3.2, p.~195]{Hatcher2002}}


Let $C_f,C_g$ denote the degree-four companion matrices.
Take $f=\Phi_5$ and either
\begin{equation}
 \text{case A: }g=\Phi_{12},\qquad
 \text{case B: }g=\Phi_3\Phi_6.
 \hypertarget{tag-02F5}{}\label{maps:eq:F.36}
\end{equation}
Let $W$ be the quotient of
$\bigwedge^2(C_f^{-1},C_g)$ by its invariant line, and
let $V$ be its unique rational extension of \ref{maps:src:F.2.11}.
For every invariant integral lattice in $V$ with primitive
maximal square-zero cusp, every simply connected free affine
filling, with multiplicities $m_1,m_2$ divisible by five and
six respectively, fails three-connectivity.
More precisely, put $r=v_3(m_2)\geq1$.
In case A, $\operatorname{Tor}_3 H_3$ has order
$3^{2r-1}$.  In case B,
$\operatorname{Tor}_3 H_2\simeq\mathbb Z/3^r$.
For $m_2=6$ the respective obstructing group is $\mathbb Z/3$.

\end{theorem}


\begin{proof}
\begin{enumerate}[label=\arabic*.]
\item We first give rational models, including the quotient
maps.  Use the exterior basis $(01,02,03,12,13,23)$ for the companion
pair.  Its common primitive invariant vectors are
\[
 \omega_{\mathrm A}=(1,2,1,2,2,1)^t,\qquad
 \omega_{\mathrm B}=(1,0,1,0,0,1)^t.
\]
Let $T$ have columns $\omega,e_1,\ldots,e_5$; its determinant is one.
The last five rows of $T^{-1}$ give the quotient map.  In that
quotient the two matrices $a,b$, and an extension row $c$, are
\[
\begin{array}{c|c|c}
 &a&b\\\hline
\mathrm A&
\begin{pmatrix}2&-1&-2&1&0\\1&0&-1&0&0\\3&-1&-2&0&1\\
3&0&-2&0&0\\1&1&-1&0&0\end{pmatrix}&
\begin{pmatrix}0&-2&0&1&0\\0&-1&0&0&1\\0&-1&0&0&0\\
1&-2&0&0&0\\0&-1&1&0&-1\end{pmatrix}\\[3pt]
\mathrm B&
\begin{pmatrix}0&-1&0&1&0\\1&0&-1&0&0\\1&-1&0&0&1\\
1&0&0&0&0\\1&1&-1&0&0\end{pmatrix}&
\begin{pmatrix}0&0&0&1&0\\0&-1&0&0&1\\0&-1&0&0&0\\
1&0&0&0&0\\0&-1&1&0&1\end{pmatrix}
\end{array}
\]
\[
 c_{\mathrm A}=(-2,-3,1,2,0),\qquad
 c_{\mathrm B}=(0,-1,1,0,0),\qquad
 A=\operatorname{diag}(1,a),\quad
 B=\begin{pmatrix}1&c\\0&b\end{pmatrix}.
\]
Exterior eigenvalue products give characteristic polynomial
$(x-1)\Phi _5$ for $a$, and $(x-1)(x+1)^2\Phi _3$ or
$(x-1)(x+1)^2\Phi _6$ for $b$.  Each has fixed rank one.
Direct multiplication gives $A^5=B^6=I$, $(AB-I)^2=0$, and
nonzero Smith factors $1,1,1$ for $AB-I$ and $1,1$ for $ab-I$.
The row conditions are $c\sum_{j=0}^5b^j=0$ and $c(ab-I)=0$.
The full quotient matrix algebra verified in the next item implies
that this rank-five quotient is irreducible, with no invariant
vectors or covectors.  Thus \cref{maps:src:F.2.11} identifies these
models with the rational extensions, up to rescaling
their invariant line.

\item Put $R=\mathbb Z_{(3)}$ and let $\mathcal P_R$ be the algebra
of blocks $\left(\begin{smallmatrix}u&v\\0&D\end{smallmatrix}\right)$,
with $u\in R$, $v\in R^{1\times5}$ and $D\in\operatorname{Mat}_5(R)$.
In the coordinate order $(0,j)$, $0\leq j\leq5$, followed by
$(i,j)$, $1\leq i,j\leq5$, the following $31$ words have coordinate
determinants $512$ and $-512$ in cases A and B:
\[
\begin{aligned}
&I,A,B,A^2,AB,BA,B^2,A^3,A^2B,ABA,AB^2,BA^2,BAB,B^2A,\\
&A^3B,A^2BA,A^2B^2,ABA^2,AB^2A,BA^3,BA^2B,BAB^2,B^2A^2,B^2AB,\\
&A^3BA,A^3B^2,A^2BA^2,ABA^3,ABA^2B,BA^3B,A^3BA^2.
\end{aligned}
\]
Their adjugates divided by the displayed determinants express
every elementary matrix of $\mathcal P_R$ as an $R$-linear
combination of these words.  Consequently $R[A,B]=\mathcal P_R$.
Its lower block also proves the irreducibility used in item~1.

For an invariant lattice, the elementary projectors and lower-block
matrix units give
\[
 L_R=3^\alpha Re_0\oplus3^\beta R^5,\qquad \beta\geq\alpha;
\]
the inequality is forced by $E_{0j}$.  In this basis the upper
nilpotent row is multiplied by $3^{\beta-\alpha}$.  If
$\beta>\alpha$, reduction modulo three has rank two, contradicting
the primitive rank-three cusp.  Thus $\beta=\alpha$ and the map $3^{-\alpha}I$ identifies $L_R$ with $R^6$, intertwining every
exterior power.  The other product conventions are related by
\[
 BA-I=A^{-1}(AB-I)A,\qquad
 (BA)^{-1}-I=-(BA)^{-1}(BA-I).
\]

\item We record the consequence of simple connectivity for the
affine translations.  Write a finite affine generator,
after a translation of the torus, as
$x\mapsto M_ix+v_i/m_i$, where $M_iv_i=v_i$ and $v_i$ is its
integral power vector.  This change has an explicit map.  If $t$
is a lift of the original translation and $d$ is the linear order,
put $t_0=d^{-1}\sum_{j=0}^{d-1}M_i^jt$ and $b_0=t-t_0$.
Then
\[
 h=\frac1d\sum_{j=0}^{d-1}(d-1-j)M_i^jb_0,
 \qquad (I-M_i)h=b_0.
\]
For $T_h(x)=x+h$, the conjugate $T_h^{-1}gT_h$ has translation $t_0$, and
$v_i=m_it_0$ is integral because the given affine action has
period $m_i$.

The base-orbifold quotient in \cref{maps:src:F.2.7} shows that
$\gcd(m_1,m_2)=1$.  In particular $m_1$ is a unit in $R$.
The homology augmentation matrix $[A^t-I\mid B^t-I]$ has first
row zero and rank-five minors $-4$ and $-1$ in its last five
rows, respectively.  Its coinvariant quotient over $R$ is
therefore $R$ through $\ell(x)=x_0$.  With $a_i=\ell(v_i)$,
the finite power relations and cusp clutch give the quotient
relation matrix
\[
 \begin{pmatrix}m_1&0&-a_1\\0&m_2&-a_2\\1&1&-\delta\end{pmatrix}.
\]
Here $\delta$ is the character value of the cusp clutch.
The character annihilates the collapsed angular cycles, so these
maps on the finite caps and cusp agree on the boundary and define
a quotient of the abelianized fundamental group.  Simple
connectivity makes its determinant a unit.  Reduction modulo
three forces $a_2$ to be a unit.  Arbitrary clutch translations
are included in $\delta$.

\item Let $\alpha$ be the primitive integral invariant character
whose $R$-span is $\ell$, and let $K=\ker\alpha$.  The invariant
sublattice $\Lambda'=\mathbb Zv_2\oplus K$ has index
$|\alpha(v_2)|$, prime to three.  It gives a torus cover
and a cover $F'_2\to F_2$ of the second reduced finite
fibre.  In coordinates beginning with $v_2$, its affine map is
\[
 (\theta,y)\longmapsto(\theta+1/m_2,Dy).
\]
Thus $F'_2$ is a five-torus mapping torus, and its generic
six-torus is the $m_2$-fold cover in the circle direction.
Restriction to $K$ is the dual quotient
$(x_0,x')\mapsto x'$ from $V_R$ to $W_R$, intertwining $B$ with $b$.
Reversing the mapping-torus circle replaces $b$ by $b^{-1}$;
the identities $b^{-1}-I=-b^{-1}(b-I)$ and
$\sum_{j=0}^5b^{-j}=\sum_{j=0}^5b^j$ give the exact comparison.

Put $b_k=\bigwedge^kb$.  The nonzero Smith factors of $b_k-I$
for $0\leq k\leq4$ are units in $R$.  The Wang sequences make
$H^q(F'_2;R)$ free for $q\leq4$.  Pullback to its generic torus
is injective, since transfer kills its kernel by $m_2$ and the
source is free.  Deck transformations of $F'_2\to F_2$ act by
translations on that generic cover.  Their action on its
cohomology is trivial, and injectivity makes it trivial on
$H^q(F'_2;R)$ as well.  Transfer for this prime-to-three cover
therefore identifies the restriction maps with those for $F_2$.

The covering map on the horizontal term of the Wang
sequence is the norm
\[
 \mathcal N_{m_2,k}=\sum_{j=0}^{m_2-1}b_k^j:
 \operatorname{coker}(b_k-I)\longrightarrow\operatorname{Inv}(b_k).
\]
Indeed, the lift is the union of the $m_2$ successive prisms
with their transported coefficients.  On the vertical invariant
quotient pullback is the identity.  Consequently horizontal
inclusion induces the isomorphism
\[
 \frac{\operatorname{Inv}(b_{q-1})}
      {\mathcal N_{m_2,q-1}\bigwedge^{q-1}R^5}
 \ \xrightarrow{\ \sim\ }\
 \frac{\operatorname{Inv}(\bigwedge^q V_R)}{J_{2,q}}.
\]
Surjectivity follows by subtracting a restriction with the same
vertical component; its kernel is exactly the horizontal norm
image.

\item Write $\lambda=m_2/6$.  Since $b^6=I$,
$\mathcal N_{m_2,k}=\lambda\mathcal N_{6,k}$.
In case A a primitive basis of $\operatorname{Inv}(b_2)$,
in lexicographic exterior order, is the column matrix
\[
 Q=\begin{pmatrix}
 1&0\\1&4\\2&4\\0&2\\0&3\\1&4\\0&3\\1&0\\0&0\\0&2
 \end{pmatrix}.
\]
The norm coordinates in this basis are
\[
 \begin{pmatrix}
 0&3&0&-3&0&0&-4&3&4&3\\
 0&0&0&0&0&0&2&0&-2&0
 \end{pmatrix}.
\]
Their image has basis $\left(\begin{smallmatrix}1&0\\4&6\end{smallmatrix}\right)$,
and multiplication by
$\left(\begin{smallmatrix}1&0\\-4&1\end{smallmatrix}\right)$
gives $\operatorname{diag}(1,6)$.  This gives the explicit quotient
map $(x,y)\mapsto(x,y-4x)$ modulo $(\lambda,6\lambda)$.
In case B the primitive invariant vector of $b$ is
$(1,0,0,1,0)^t$, and the norm coordinates are $(3,0,0,3,0)$.
After multiplication by the remaining units of $R$, these give
\[
 Q_{2,3}^{\mathrm A}\cong R/3^{r-1}R\oplus R/3^rR,
 \qquad Q_{2,2}^{\mathrm B}\cong R/3^rR.
\]
\item We use the integral specialization and base-sheaf method
of \cite[Lemma~7.25 and Propositions~7.26--7.27, pp.~61--63]{AC26}
and \cite[section ``The integral homology'']{E26}, with the
following explicitly computed coefficients.  Put
$L_q=\bigwedge^qR^6$, $A_q=\bigwedge^qA$, $B_q=\bigwedge^qB$.
The based-sphere complex for $j_*L_q$ is
\[
\begin{gathered}
 J_{1,q}\oplus J_{2,q}\oplus J_{0,q}
 \xrightarrow{d_0} L_q^2\xrightarrow{\Delta_q}L_q,\\
 J_{1,q}=\ker(A_q-I),\quad J_{2,q}=\ker(B_q-I),\quad
 J_{0,q}=\ker(B_qA_q-I),\\
 d_0(j_1,j_2,j_0)=(j_1-j_0,j_2-A_qj_0),\\
 \Delta_q(x,y)=(A_q-I)x+(B_q-I)y.
\end{gathered}
\]
The cusp paths carry $j_0$ and $A_qj_0$, so their combined
boundary is zero.  Inverse generator conventions give a chain
isomorphism: use cusp basis $A_qJ_{0,q}$ for the inverse pair,
and maps $-I$, $\operatorname{diag}(-A_q,-B_q)$, $I$ in degrees
zero, one and two.  Both chain identities follow by multiplication.

Unimodular HNF transformations give primitive invariant bases
and kernel coordinates for $d_0$ in $\ker\Delta_q$.  In both cases
the exact Smith records are
\[
\begin{array}{c|ccc|c|c}
q&\operatorname{rank}H^0&\operatorname{rank}H^1&\operatorname{rank}H^2
 &\operatorname{SNF}(d_0\text{ in }\ker\Delta_q)
 &\operatorname{SNF}(\Delta_q)\\\hline
1&1&0&0&1^6&1^5,2\\
2&0&0&0&1^{15}&1^{14},2\\
3&0&3&0&1^{17}&1^{19},2\\
4&0&0&0&1^{15}&1^{14},2
\end{array}
\]
All displayed groups are therefore free over $R$.

\item Write $\mathcal R_q=R^qf_*R$.  Proper base change
\cite[Tag~09V5, Lemma~20.18.1]{StacksProject} identifies the stalks
with the fibres.  At the first finite cap, $m_1$ is a
unit in $R$, so transfer gives the full invariant restriction
image and zero kernel.  At the second finite cap, item~4 gives
injectivity and the norm quotient.  At the prescribed cusp,
dualizing the primitive-image collapse of
\cref{night:P1.11} gives the invariant restriction image, with
zero kernel in degrees one through three and a free kernel $R$
in degree four.  The latter is dual to the free degree-four
cokernel of the collapse.

For $q=2,3$ there is consequently an exact sequence
\[
 0\longrightarrow\mathcal R_q\longrightarrow j_*L_q
 \longrightarrow Q_q\longrightarrow0,
\]
where $Q_q$ is supported at the second finite point.  Since
$H^0(j_*L_q)=0$, its connecting map embeds $Q_q$ in
$H^1(\mathcal R_q)$, with the free quotient $H^1(j_*L_q)$
computed above.  Its torsion is exactly $Q_q$.
Also $H^2(\mathcal R_1)=H^2(\mathcal R_2)=H^2(\mathcal R_3)=0$.
Injectivity gives $H^0(\mathcal R_3)=0$.  The degree-four
free cusp kernel, together with $H^0(j_*L_4)=0$, gives
$H^0(\mathcal R_4)=R$.

\item In the Leray sequence
\cite[Tag~01F2, Lemma~20.13.4]{StacksProject}, total degree three
has only $H^1(\mathcal R_2)$.  Its middle base column has no
incoming or outgoing differential.  Thus case B gives
$H^3(M;R)=R/3^rR$.
In total degree four, the surviving terms are
$H^1(\mathcal R_3)$ and $H^0(\mathcal R_4)=R$;
the degree-two base term is zero and the possible differential
from the latter term lands in $H^2(\mathcal R_3)=0$.
The extension by the free module $R$ splits.  Hence in case A
\[
 \operatorname{Tor}H^4(M;R)
   \cong R/3^{r-1}R\oplus R/3^rR.
\]
Universal coefficients
\cite[Theorem~3.2, p.~195]{Hatcher2002} identify this torsion
with the torsion of $H_3(M;R)$, and similarly identify the
case-B torsion with that of $H_2(M;R)$.  Localization retains
the finite three-primary subgroups by the identity on each
$\mathbb Z/3^s$ summand.  This proves the stated orders and
the case-B isomorphism.  At $m_2=6$ the first case-A summand
vanishes and the second is $\mathbb Z/3$; case B is also
$\mathbb Z/3$.  These groups obstruct three-connectivity.
\item The two displayed rational models have explicit smooth marked
realisations with the required free caps and simply connected filling.
Use their displayed standard cohomology lattice $\mathbb Z^6$ and its
dual homology basis, so $M_1=A^t$, $M_2=B^t$.
Direct multiplication gives $A^5=B^6=I$, exact orders five
and six, and determinant one. Put $K=BA-I$. Then $K^2=0$, and the
third minor on rows $(0,3,4)$ and columns $(1,2,3)$ is one in both
models. Hence $K$ has rank three and all its nonzero Smith factors
are one. The marked cusp is
\[
 M_0-I=((M_1M_2)^{-1}-I)=-K^t.
\]
Set $b_1=e_0,b_2=e_3,b_3=e_4$, $a_j=(M_0-I)b_j$ and
$S=(a_1,a_2,a_3,b_1,b_2,b_3)$. Direct multiplication gives
$S\in\operatorname{GL}_6(\mathbb Z)$ and
\[
 (M_0-I)S=S\begin{pmatrix}0&I_3\\0&0\end{pmatrix}.
\]
This is the exact integral marking used to transport the index-one
Kuhn cap and its collapse maps.

The affine power vectors in these same homology coordinates are
\[
\begin{array}{c|cc}
 &v_1&v_2\\\hline
 \mathrm A&(1,0,0,0,0,0)&(-1,2,0,-1,0,0)\\
 \mathrm B&(1,0,0,0,0,0)&(-1,0,0,0,0,1).
\end{array}
\]
They satisfy $M_iv_i=v_i$. The first-coordinate covector
$\ell$ is fixed by both $M_i$ and has values $1,-1$ on $v_1,v_2$.
Consequently every nonidentity power of
$x\mapsto M_1x+v_1/5$ or $x\mapsto M_2x+v_2/6$ changes the
$\ell$-circle by $j/5$ or $-j/6$, respectively. These are free
cyclic actions of the exact stated orders. Their diagonal actions
on the product with the rotating disc therefore give the smooth finite quotient caps.

Take the affine torus bundle over the three-holed sphere with these
two boundary holonomies and their inverse product on the third
boundary. Glue the two finite caps and the proper index-one Kuhn
cap of \cref{night:P1.9,night:P1.11} by the
specified marked mapping-torus collars and zero cusp clutch.
For an affine cusp monodromy $x\mapsto M_0x+c$, a path of translations
$d(s)$ with $d(0)=0$, $d(1)=c$ gives the collar map
$(x,s)\mapsto(x+d(s),s)$ from the linear mapping torus. The prescribed
zero clutch fixes its integral lift. The lattice map on this collar
is the displayed $S$.

The cap fundamental-group maps retain these choices. On the finite
cap the universal covering $\mathbb R^6\times\Delta$ has the deck
transformation
$(x,t)\mapsto(M_ix+v_i/m_i,e^{2\pi i/m_i}t)$; its $m_i$-th power
is the lattice translation $v_i$. Thus the boundary generators
$a,b$ satisfy $a^5=v_1$, $b^6=v_2$. On the Kuhn cap the polar
projection onto its base three-torus and angular-identity section
induce inverse maps on fundamental groups. In its dual-cell model
the one-skeleton consists of the base edges, and each angular loop
contracts after transport to a dual vertex where that angular
fibre is collapsed. The fibre map is therefore exactly
$\Lambda\to\Lambda/\operatorname{im}(M_0-I)$, expressed by the last
three rows of $S^{-1}$. The meridian bounds the toric coordinate
disc at a smooth point of a central component; the zero clutch
identifies it with $(ab)^{-1}$. Hence the cusp imposes $ab=1$.

For $J=[M_1-I\mid M_2-I]$, the first row is zero. Its lower five
rows have a unit minor on columns $(1,2,3,6,10)$ in case A and
$(1,2,3,4,6)$ in case B, with determinants $-1$ and $1$.
The inverse of that minor supplies integral preimages of
$e_1,\ldots,e_5$. Thus
$\operatorname{im}J=\ker\ell$ integrally. Van Kampen now gives
simple connectivity at the group level: from $b=a^{-1}$ and the
two power relations,
\[
 a^{-1}=a^5a^{-6}=v_1v_2\in\Lambda.
\]
So $a,b$ lie in the abelian image of the fibre lattice, their
conjugation actions become trivial, and $J$ is killed. The exact
remaining fibre quotient is $\mathbb Z$ via $\ell$, with generator
$h$ and images $v_1=h$, $v_2=-h$. The displayed identity forces
$a=1$, and $a^5=v_1$ then kills $h$. Every generator is trivial.
Equivalently, after this nonabelian reduction the residual relation
matrix on $(h,a,b)$ is
\[
 \begin{pmatrix}-1&5&0\\1&0&6\\0&1&1\end{pmatrix},
 \qquad\det=1.
\]

Use the rational invariant-extension table above for these models. The finite covering restrictions identify
rational cap cohomology with invariants. The primitive cusp collapse
of~\ref{night:P1.11} gives the same identification through degree
three and a one-dimensional kernel in degree four. Hence
$R^q p_*\mathbb Q=j_*\Lambda^qV$ for $1\le q\le3$, and
\[
 0\longrightarrow\mathbb Q_{\mathrm{cusp}}
 \longrightarrow R^4p_*\mathbb Q
 \longrightarrow j_*\Lambda^4V\longrightarrow0.
\]
The table therefore leaves, in total degree at most four, only
$E_2^{0,0}=E_2^{2,0}=E_2^{0,1}=\mathbb Q$,
$E_2^{1,3}=\mathbb Q^3$ and $E_2^{0,4}=\mathbb Q$.
Simple connectivity forces the $d_2:E_2^{0,1}\to E_2^{2,0}$ to be an isomorphism. There are no
other possible nonzero differentials on these terms. Thus
$(b_0,\ldots,b_4)=(1,0,0,0,4)$; oriented duality on the closed
smooth eight-manifold gives the full tuple
$(1,0,0,0,4,0,0,0,1)$.
\end{enumerate}
\end{proof}




\begin{theorem}[Four order pairs and eight local spectral pairs]\hypertarget{tag-02N3}{}\label{maps:src:F.2.14}
\StatementRecord{maps:src:F.2.14}{proved}{02N3}{\cref{maps:src:F.2.15}, \cref{maps:src:F.2.9}, \cref{maps:src:F.2.7}, \eqref{maps:eq:F.40}, \cref{maps:src:F.2.2}, \cref{night:P1.11}}


Every simply connected two-finite-point datum with $b_2=0$ has
one finite linear order equal to five and the other equal
to four, six, eight or twelve.
After suppressing the common factor $(T-1)^2$, its unordered
pair of finite characteristic polynomials is exactly one
of the following necessary possibilities:
\begin{equation}
\begin{array}{c|c|l}
 \text{orders}&\text{order-five factor}&\text{other factor}\\\hline
 (5,4)&\Phi_5&\Phi_2^2\Phi_4\\
 (5,6)&\Phi_5&\Phi_2^2\Phi_3,\ \Phi_2^2\Phi_6,\ \Phi_3\Phi_6\\
 (5,8)&\Phi_5&\Phi_8\\
 (5,12)&\Phi_5&\Phi_{12},\ \Phi_3\Phi_4,\ \Phi_4\Phi_6.
\end{array}
 \hypertarget{tag-02F6}{}\label{maps:eq:F.39}
\end{equation}
Its rank-five quotient is rationally irreducible.
In particular all order-three pairs, including $(3,4)$,
are impossible.

\end{theorem}


\begin{proof}
\begin{enumerate}[label=\arabic*.]
\item By \cref{maps:src:F.2.9}, \cref{maps:src:F.2.15}, the cohomology
system $V$ has a one-dimensional invariant space $I$, no invariant
covectors, and two-dimensional finite fixed spaces $F_1,F_2$.
The finite orders are among the coprime pairs of
\cref{maps:src:F.2.8}.  At a finite quotient, averaging identifies
the degree-two specialization with the fixed part of $\Lambda^2V$.
At the cusp, \cref{night:P1.11} gives the surjective homology
collapse with kernel the saturation of
$(\Lambda^2M_0-I)\Lambda^2\Lambda$.  Its dual has image
$\ker(\Lambda^2M_0^t-I)$ under the evaluation pairing.  Consequently
the natural specialization morphism identifies
\[
 R^2g_*\mathbb Q=j_*\Lambda^2V.
\]
The term $H^0(j_*\Lambda^2V)$ has no incoming Leray differential,
and its only possible outgoing differential targets
$H^2(j_*V)=H^0(V^*)^*=0$.  It survives to $H^2(M;\mathbb Q)=0$.
Thus $a_2':=\dim(\Lambda^2V)^\Gamma=0$.

\item Let $\beta:V\to V^*$ be an invariant alternating form.
Since $\beta(I)\subset(V^*)^\Gamma=0$, its radical contains $I$.
For the finite averaging projection $P_i:V\to F_i$, invariance gives
$\beta(u,x)=\beta(u,P_ix)$ when $u\in F_i$.
The alternating restriction to the two-dimensional space $F_i$
has $I$ in its radical and is therefore zero.  Hence
$F_1+F_2\subset\operatorname{rad}\beta$.  This sum has dimension
three because $F_1\cap F_2=I$.  The alternating rank is even,
so every such form has rank at most two.

Suppose $\beta,\gamma$ were independent invariant forms.
They and their sum have rank two, so
$\beta\wedge\gamma=0$.  Their image planes in $V^*$
are distinct invariant two-planes meeting in a rational line.
Each finite generator acts on that line by a sign.  At least one
finite order is odd, and the unipotent product acts trivially on the
line; both signs are therefore $+1$.  This contradicts
$(V^*)^\Gamma=0$.  It follows that
$b_2':=\dim(\Lambda^2V^*)^\Gamma\leq1$.

\item At the cusp, the exterior square of three size-two unipotent
blocks has fixed rank nine: three fixed vectors come from their
individual exterior squares and two from each of the three pairwise
tensor products.  If $d_i$ denotes a finite fixed rank on
$\Lambda^2V$, the invariant-sheaf Euler identity gives
\begin{equation}
 -15+d_1+d_2+9
 =a_2'-h^1(j_*\Lambda^2V)+b_2'\leq1.
 \hypertarget{tag-02F8}{}\label{maps:eq:F.40}
\end{equation}
Here the degree-two group is dual to
$(\Lambda^2V^*)^\Gamma$, by compact-support duality and the
extension-by-zero sequence.

\item The two-dimensional unit eigenspace leaves four dimensions
for the nonunit cyclotomic factors.  Local complex compatibility
requires even multiplicity of $-1$.  Partitioning four into the
cyclotomic degrees gives the following complete list:
\[
\begin{array}{c|l}
3&\Phi_3^2\\
4&\Phi_4^2,\ \Phi_2^2\Phi_4\\
5&\Phi_5\\
6&\Phi_6^2,\ \Phi_2^2\Phi_3,\ \Phi_2^2\Phi_6,\ \Phi_3\Phi_6\\
8&\Phi_8\\
10&\Phi_{10}\\
12&\Phi_{12},\ \Phi_3\Phi_4,\ \Phi_4\Phi_6.
\end{array}
\]
The unit eigenspace contributes one fixed bivector.  A repeated
nonreal quadratic pair contributes four further fixed bivectors,
so $d_i=5$ for $\Phi_3^2,\Phi_4^2,\Phi_6^2$.
Every other listed factor contributes two, so $d_i=3$.
Inequality \eqref{maps:eq:F.40} excludes every occurrence of
$d_i=5$.  In particular all order-three pairs disappear.
The coprime order list of \cref{maps:src:F.2.8} therefore reduces
to $(5,4),(5,6),(5,8),(5,12)$; their remaining factors are exactly
the eight entries of \eqref{maps:eq:F.39}.
Each pair contains the order-five factor, so
\cref{maps:src:F.2.9} proves irreducibility of $V/I$.
\end{enumerate}
\end{proof}







\section{Integral lattices and extension classes}\label{pub:V:sec-lattices}

\begin{theorem}[An unbounded integral family]\hypertarget{tag-02N5}{}\label{maps:src:F.2.16}
\StatementRecord{maps:src:F.2.16}{proved}{02N5}{\eqref{maps:eq:F.U1}, \eqref{maps:eq:F.U3}, \eqref{maps:eq:F.U4}, \eqref{maps:eq:F.U5}, \eqref{maps:eq:F.U6}, \eqref{maps:eq:F.U7}, \cref{maps:src:F.2.5}}


There are infinitely many simultaneous rational conjugacy classes
of integral rank-six tuples of the linear type $(5,4)$ satisfying \cref{def:admissible},
with primitive maximal square-zero cusp.  An explicit infinite
subfamily has free affine lifts of orders five and four.
Apart from the parameter zero, the rank-five irreducible quotient
in the family is not self-dual.  Thus no finite bound on rational
conjugacy classes is valid for the unrestricted input tuples.


\end{theorem}


\begin{proof}
\begin{enumerate}[label=\arabic*.]
\item Let $v\in\mathbb Z$, put $q=v^2-v+1$, and let
$Ae_j=e_{j+1}$ on $\mathbb Q^5$, with indices modulo five. Define
\begin{equation}
 X=(e_0,\ e_2+ve_4),\qquad
 Y=\begin{pmatrix}
 0&-q&-v^2&-1&v\\
 0&1-vq&-vq&v^2-v&q
 \end{pmatrix},\qquad N=XY.
 \hypertarget{tag-02IK}{}\label{maps:eq:F.U1}
\end{equation}
Multiplying $Y$ by the two columns of $X$ gives $YX=0$, so
$N^2=X(YX)Y=0$. Consequently
\[
 C=(I+N)A,\qquad B=A^4(I-N)=C^{-1}.
\]
The cycle gives $A^5=I$. Exact multiplication and determinant
expansion give
\begin{equation}
 N^2=0,\qquad A^5=C^4=I_5,\qquad
 \chi_C(x)=(x-1)(x+1)^2(x^2+1).
 \hypertarget{tag-02IO}{}\label{maps:eq:F.U2}
\end{equation}
For the determinant calculation one may substitute
\[
 (xI-A)^{-1}=
 \frac{x^4I+x^3A+x^2A^2+xA^3+A^4}{x^5-1}
\]
in the rank-two determinant identity
\[
 \det(xI-C)=\det(xI-A)
       \det\bigl(I_2-YA(xI-A)^{-1}X\bigr).
\]
Direct multiplication gives $C^4=I$.
The factor $x^2+1$ shows that $B$ and $C$ have exact order four.

\item Let $P$ have columns $\mathbf1/q,e_1,e_2,e_3,e_4$, and put
$A'=P^{-1}AP$, $B'=P^{-1}BP$. Since $q>0$ and $\det P=1/q$,
these are rational changes of coordinates. Define
\begin{equation}
 s=v^2+v-1,\qquad r=-\frac{(v+1)(v^2-2)}2,
 \qquad
 \beta=(s,\ s-vr,\ v^2s+vr,\ r,\ -vs-r)
 \hypertarget{tag-02IP}{}\label{maps:eq:F.U3}
\end{equation}
and
\begin{equation}
 \widehat A=\begin{pmatrix}1&0\\0&A'\end{pmatrix},\qquad
 \widehat B=\begin{pmatrix}1&\beta\\0&B'\end{pmatrix}.
 \hypertarget{tag-02IQ}{}\label{maps:eq:F.U4}
\end{equation}

The lattice $P\mathbb Z^5$ has basis
$f_0=\mathbf1/q,f_1=e_1,\ldots,f_4=e_4$.
Its images under $A$ are $f_0,f_2,f_3,f_4$ and
$qf_0-f_1-f_2-f_3-f_4$, respectively; $A^{-1}=A^4$
therefore also preserves this lattice. The columns of $X$ are
$qf_0-f_1-f_2-f_3-f_4$ and $f_2+vf_4$, while
\[
 Y':=YP=\begin{pmatrix}
 -2&-q&-v^2&-1&v\\
 2(1-v)&1-vq&-vq&v^2-v&q
 \end{pmatrix}.
\]
Thus $X$ and $Y$ are integral maps between the lattices.
The already proved $YX=0$ shows that $I+N$ and its inverse $I-N$
preserve $P\mathbb Z^5$. Since $B=A^{-1}(I-N)$, both
$A',B'$ belong to $\operatorname{GL}_5(\mathbb Z)$.
For $r$, odd $v$ makes $v+1$ even, and even $v$ makes
$v^2-2$ even; hence $r$ and all entries of $\beta$ are integers.
The hatted matrices are therefore integral automorphisms.

The exact unit needed in the hatted lattice is
\[
 z=(-v^3+v-1,0,0,2,2(1-v^2))^t.
\]
Writing $a=-v^3+v-1$, the first coordinate of $Y'z$ is
$-2a-2+2v(1-v^2)=0$. Half its second coordinate is
$(1-v)a+(v^2-v)+q(1-v^2)=0$, since the first two terms
sum to $v^4-v^3+v-1$ and the last is its negative. Furthermore
\begin{equation}
 \beta z=s(v^3-v-1)+2v^2r=1.
 \hypertarget{tag-02IR}{}\label{maps:eq:F.U5}
\end{equation}
Indeed the two summands expand to
$v^5+v^4-2v^3-2v^2+1$ and
$-v^5-v^4+2v^3+2v^2$.
Since $AB=I-N$ and $NPz=X(Y'z)=0$, the six-dimensional
map satisfies
\[
 (\widehat A\widehat B-I)(0,z)^t=e_0.
\]
This establishes the unit in the displayed hatted coordinates.

The row identity
$\beta(I+B'+B'^2+B'^3)=0$ proves
$\widehat B^4=I$; also $\widehat A^5=I$. Their quotient actions
show that these orders are exact.

Use the lattice change
\begin{equation}
 L=\frac14\begin{pmatrix}
 -2&3&-5&-1&5&-3\\
 0&2&-2&2&-2&2\\
 0&-4&0&0&4&-4\\
 0&0&0&-4&4&-4\\
 0&0&4&-4&4&-4\\
 0&0&0&-4&4&0
 \end{pmatrix},\qquad
 A_v=L^{-1}\widehat A L,\quad B_v=L^{-1}\widehat B L.
 \hypertarget{tag-02IS}{}\label{maps:eq:F.U6}
\end{equation}
Here $\det L=-1/4$. Expanding these products gives entries in
$\mathbb Z[v]$ for both matrices, with
$LA_v=\widehat A L$ and $LB_v=\widehat B L$. Thus the
defined matrices act on the stated $\mathbb Z^6$ for every integer
parameter. Their determinants are one and their inverses are
the integer powers $A_v^4,B_v^3$. Their characteristic polynomials are
\[
 (x-1)^2\Phi_5(x),\qquad (x-1)^2(x+1)^2(x^2+1),
\]
respectively.

\item Put $K=B_vA_v-I$. Multiplication gives $K^2=0$.
The following are third minors of $K$, in zero-based indices:
\[
 \begin{array}{c|c|c}
 \text{rows}&\text{columns}&\text{determinant}\\\hline
 (0,1,4)&(1,2,3)&v(v+1)\\
 (0,1,4)&(1,3,4)&v^2\\
 (0,3,4)&(1,2,3)&(v+1)^2.
 \end{array}
\]
Their integer combination $(v+1)^2-2v(v+1)+v^2$ equals one.
Hence $K$ has rank at least three, while $K^2=0$ in dimension six
gives rank at most three. The gcd of its third minors is one,
so its three nonzero Smith factors are all one. This proves
primitivity directly in the lattice, for every $v$.
For $M_1=A_v^t$ and $M_2=B_v^t$ the cusp map is
\[
 M_0=(M_1M_2)^{-1}=(I+K^t)^{-1}=I-K^t.
\]
Thus $M_0-I=-K^t$, which has the same integral Smith factors
and the required maximal square-zero type.

\item Both $A_v,B_v$ fix $e_0$. The quotient morphism to
the original rank-five pair is
\[
 R=P\,\operatorname{pr}_5 L:\mathbb Q^6\longrightarrow\mathbb Q^5,
 \qquad RA_v=AR,\quad RB_v=BR.
\]
It is surjective and has kernel $\mathbb Qe_0$, since
$Le_0=-e_0/2$. We prove that this quotient is irreducible over
$\mathbb Q$. The coprime factors $x-1$ and $\Phi_5(x)$ of
$\chi_A$ give the projector
$(I+A+A^2+A^3+A^4)/5$. Its image is $\mathbb Q\mathbf1$ and
its kernel is $S=\{\sum x_i=0\}$. Irreducibility of $\Phi_5$
shows that these are the only nonzero proper rational
$A$-invariant subspaces.

Neither is $C$-invariant. Indeed,
\[
 N\mathbf1=-2qX\binom{1}{v-1}
\]
has nonzero first entry and zero entries in positions one and three,
so $C\mathbf1=\mathbf1+N\mathbf1$ is not a multiple of
$\mathbf1$. Write $\epsilon=(1,1,1,1,1)$. If $S$ were
$C$-invariant, then $\epsilon N$ would be a multiple of
$\epsilon$; its first entry is zero, so that multiple would
vanish. But its entries in columns three and four are
$v^3-v-1$ and $v^3+v+1$. Their simultaneous vanishing would
give $v=-1$, at which the first is $-1$. This is impossible.
Since $B=C^{-1}$, the pair $(A,B)$ is irreducible, as claimed.

\item The rank-two update gives exact traces
\[
 \operatorname{tr}(AC)=\operatorname{tr}(YA^2X)=1-v^2q,
 \qquad
 \operatorname{tr}((AC)^{-1})
       =-\operatorname{tr}(YA^3X)=1+v^2q.
\]
Here $\operatorname{tr}(A^2)=\operatorname{tr}(A^3)=0$ and
$B=A^4(I-N)$ give the displayed equalities. An equivariant
isomorphism $F$ to the contragredient representation would
satisfy $FA=A^{-t}F$ and $FB=B^{-t}F$, and therefore identify
the character of the word $AC=AB^{-1}$ with the character
of its contragredient matrix. The two displayed traces would
be equal. Their difference $-2v^2q$ is nonzero when $v\ne0$,
proving the required nonselfduality.

At $v=0$ the symmetric matrix
$Q_{ij}=1$ for $j-i\equiv1,4\pmod5$ and $Q_{ij}=0$ otherwise
has determinant two and satisfies $A^tQA=Q$, $B_0^tQB_0=Q$.
It gives the equivariant isomorphism in that case.

For the rank-six matrices, the same word has trace
\[
 h(v):=\operatorname{tr}(A_vB_v^{-1})=2-v^2q,
 \qquad h(v+1)-h(v)=-4v^3-3v^2-3v-1.
\]
Thus $h$ is strictly decreasing on the nonnegative integers,
in particular on the nonnegative multiples of twenty.
A simultaneous rational conjugacy of the ordered matrices
would preserve this trace. There are consequently infinitely
many such rational conjugacy classes already in that subfamily.

\item For $v\in20\mathbb Z$, put
$a=(1,3,-2,-1,2,0)^t$, $e=(1,0,0,0,0,0)^t$, and define
\begin{equation}
 w_1=\frac15\sum_{j=0}^4(A_v^t)^j a,
 \qquad
 w_2=-\frac14\sum_{j=0}^3(B_v^t)^j e.
 \hypertarget{tag-02IT}{}\label{maps:eq:F.U7}
\end{equation}
The polynomial matrices reduce to their parameter-zero
values modulo five and four, respectively. At zero,
$M_1a=a$ and $M_2e=e$. The two sums are therefore divisible
by their indicated denominators, so $w_1,w_2$ are integral.
Their defining averages give $M_iw_i=w_i$. The common
first-coordinate covector $\ell$ is fixed by both $M_i$ and
satisfies $\ell(w_1)=1$, $\ell(w_2)=-1$.

For $(m_1,m_2)=(5,4)$ let
$g_i(x)=M_ix+w_i/m_i$ on $\mathbb R^6/\mathbb Z^6$.
Then $g_i^j(x)=M_i^jx+jw_i/m_i$, so $g_i^{m_i}$ is the
identity. Every power with $0<j<m_i$ changes the $\ell$-circle coordinate by $j/5$ or $-j/4$, respectively.
It has no fixed point. These are free actions of the exact
orders five and four, with the affine lifts.

\item Substitution of $v=0$ in the displayed construction gives
exactly the marked pair of \cref{pub:V:rank-six-data}, and
the averages give $w_1(0)=a$, $w_2(0)=-e$. This identifies
both the matrices and the affine numerators in the fixed marking.

\end{enumerate}
\end{proof}




\begin{proposition}[An odd torsion obstruction in this lattice]\hypertarget{tag-02N6}{}\label{maps:src:F.2.17}
\StatementRecord{maps:src:F.2.17}{proved}{02N6}{\eqref{maps:eq:F.U6}, \cref{maps:src:F.2.19}, \eqref{maps:eq:F.U5}, \cref{maps:src:F.2.12}}


For the integral tuple $(A_v,B_v)$ in \eqref{maps:eq:F.U6}, any simply
connected filling with free finite affine caps and the fixed
primitive Kuhn cusp has a subgroup
\begin{equation}
 \mathbb Z/(v^2-v+1)\ \subset\operatorname{Tor}H_2(M;\mathbb Z).
 \hypertarget{tag-02IU}{}\label{maps:eq:F.U9}
\end{equation}
In particular, in this specified family of lattices only
$v=0$ and $v=1$ can be three-connected, regardless of the
affine multiplicities and twists.  The stronger exclusion for every stable lattice in each
rational class is proved in Theorem \ref{maps:src:F.2.19} below.

\end{proposition}


\begin{proof}
\begin{enumerate}[label=\arabic*.]
\item First compute the global coinvariants in the hatted lattice.
In the rank-five quotient, the $A'$ relations identify
$f_1=f_2=f_3=f_4=t$ and $qf_0=5t$.
The number $q=v^2-v+1$ is relatively prime to five, as its
five residue values show.  Thus these coinvariants are
infinite cyclic, with $f_0=5h$ and $t=qh$.
The $B'$ relations become $q\kappa_j h$, where
\begin{equation}
 (\kappa_0,\ldots,\kappa_4)=
 (2v^2,\ v(v^3+v-1),\ v(v^3+v+1),
       -v^3+v+1,\ -v^3-v-1).
 \hypertarget{tag-02IL}{}\label{maps:eq:F.U10}
\end{equation}
Their gcd is one: it divides $2v$ by subtracting the second
and third entries, and it divides $v^3-v-1$ by the fourth;
the latter is odd and is $-1$ modulo $v$.
Hence the rank-five global coinvariants are $\mathbb Z/q$.
The extension line is killed in rank-six coinvariants by
\eqref{maps:eq:F.U5}, since
$(\widehat A\widehat B-I)(0,z)^t=e$.
Consequently the hatted rank-six coinvariants are also
$\mathbb Z/q$.  The lattice change \eqref{maps:eq:F.U6} is an isomorphism
over every $\mathbb Z_{(p)}$ with $p$ odd, so the odd part of
the rank-six global coinvariants is precisely
$\mathbb Z/q$.

\item Let $V$ be the degree-one fibre cohomology local system.
The only difference between $R^1g_*\mathbb Z$ and $j_*V$
is a finite cokernel at finite caps; the cusp degree-one
specialization is onto.  Thus
\[
 H^2(\mathbb P^1,R^1g_*\mathbb Z)
 =H^2(\mathbb P^1,j_*V)=V_\Gamma,
\]
including its $q$-primary part just computed.  The space of
global degree-two invariants is zero: the irreducible odd
rank-five quotient has no invariant alternating tensor,
and the extension filtration
$0\to W\to\bigwedge^2V\to\bigwedge^2W\to0$
gives $(\bigwedge^2V)^\Gamma=0$.

\item It remains to check every possible source of the differential
$d_2:H^0(\mathbb P^1,R^2g_*\mathbb Z)
\to H^2(\mathbb P^1,R^1g_*\mathbb Z)$.
The cusp has no kernel in degree two.  At a finite cap the
kernel is torsion in $H^2$ of that cap, namely torsion in its
first homology.  Simple connectivity of the total filling
forces the distinguished affine numerator to be prime to the
affine multiplicity, by the determinant argument of
\ref{maps:src:F.2.12}.  The cap is therefore a mapping torus of a torus
automorphism whose order divides the linear order, five or
four.  Its first-homology torsion is supported only at the
primes five or two, respectively: averaging over the finite
order kills the torsion of the vertical coinvariants.
Hence the source of this $d_2$ has no torsion at any prime
dividing $q$, since $q$ is odd and prime to five.

\item The subgroup $\mathbb Z/q$ in bidegree $(2,1)$ survives.
There are no longer differentials on a two-dimensional base,
and this terminal filtration term is a subgroup of
$H^3(M;\mathbb Z)$.  Universal coefficients identify its
torsion with a subgroup of $\operatorname{Tor}H_2(M;\mathbb Z)$.
This proves \eqref{maps:eq:F.U9}.  Finally, $v^2-v+1=1$ for an integer $v$
exactly when $v=0$ or $v=1$.

\end{enumerate}
\end{proof}




\begin{theorem}[A three-connected datum with non-self-dual quotient]\hypertarget{tag-02N7}{}\label{maps:src:F.2.18}
\StatementRecord{maps:src:F.2.18}{proved}{02N7}{\eqref{maps:eq:F.U6}, \cref{maps:src:F.3.11}, \cref{pub:V:smooth-gluing}, \cref{maps:src:F.2.16}, \cref{maps:src:F.2.6}, \cref{pub:V:integral-restrictions}, \eqref{maps:eq:F.6}, \cref{maps:src:F.4.1}}


The parameter $v=1$ in \eqref{maps:eq:F.U6}, with affine orders five and four
and twists
\[
 w_1=(1,3,-2,-1,2,0)^t,\qquad
 w_2=(-1,0,0,0,0,0)^t,
\]
gives a three-connected smooth eight-manifold with
$H_4=\mathbb Z^4$ and all other reduced homology zero except
$H_8=\mathbb Z$.  It is almost complex and is oriented
diffeomorphic to $3\mathbb{HP}^2\#(-\mathbb{HP}^2)$.
Its irreducible rank-five quotient has no invariant bilinear
form, so integral three-connectivity does not force that
quotient to be orthogonal.  By \cref{maps:src:F.3.11} it has holomorphic periods with bounded cusp correction on the punctured base.

\end{theorem}


\begin{proof}
\begin{enumerate}[label=\arabic*.]
\item The first cohomological matrix is the same $A$ as in
\cref{pub:V:rank-six-data}.  The second is
\begin{equation}
 B_1=\begin{pmatrix}
 1&0&0&0&0&0\\
 0&0&-1&0&1&-1\\
 0&0&-1&0&0&0\\
 0&0&-2&0&1&-2\\
 0&1&-2&0&1&-3\\
 0&0&-1&1&0&-1
 \end{pmatrix}.
 \hypertarget{tag-02IM}{}\label{maps:eq:F.U11}
\end{equation}
Take $M_1=A^t$, $M_2=B_1^t$ and
$M_0=(M_1M_2)^{-1}$.  The displayed twists are fixed integral
vectors with $\ell(w_1)=1$, $\ell(w_2)=-1$.
They therefore give free affine actions.  The cusp is
primitive by \ref{maps:src:F.2.16}.  Exact Hermite reduction gives
\[
 \operatorname{col}_{\mathbb Z}
 [\,M_0-I\mid M_1(M_0-I)\mid M_2(M_0-I)\,]
       =\langle e_1,\ldots,e_5\rangle.
\]
Van Kampen leaves only $5a=h$, $4b=-h$, $a+b=0$,
whose relation matrix is unimodular.  This proves
$\pi_1=1$ for the gluing, including zero cusp
clutching.

\item We give the integral data needed to reuse the differential
calculation of \cref{mapsaudit:order-four-torsion-transgression}.
The image-sheaf cohomology ranks are those of
\cref{pub:V:integral-restrictions}, as the following integral data verify.  The finite cap first homology
groups are $\mathbb Z^2$ and $\mathbb Z^2\oplus\mathbb Z/2$.
On integral invariant two-form bases their suspension periods
divided by the covering orders are, respectively,
$(-1,0,-1)$ and $(0,0,1)$.  The determinants of the invariant
degree-three skew pairings are $5^4$ and $4^4$, while the
degree-two/degree-four pairings have absolute determinants
$5^3$ and $4^3$.  Thus all the finite covering indices used
in degrees two, three and four are one.

\item For exterior coefficient degrees two, three and four, the
global coboundary Smith form has one factor two and all other
nonzero factors one. Its torsion class restricts to zero at the
order-five boundary and at the cusp, and to a nonzero class at
the order-four boundary. Thus the parabolic groups are torsion
free, of ranks $0,3,0$.
The degree-one coboundary is primitive. The common invariant
covector maps to $(1\bmod5,-1\bmod4)$ in the degree-one
finite covering cokernels. Its map from $\mathbb Z$ onto
$\mathbb Z/5\oplus\mathbb Z/4$ is surjective by the Chinese
remainder theorem.  The global coinvariants in degrees one through four
are each $\mathbb Z/2$.

\item Write $M_2=\operatorname{diag}(1,D)$.  This time use the
unimodular marking
\[
 P_1=\begin{pmatrix}
 -1&1&-1&0&-1\\
 -1&2&-1&0&-1\\
 -1&0&0&2&-1\\
 1&-1&0&-1&1\\
 0&0&2&-1&1
 \end{pmatrix},\qquad \det P_1=-1.
\]
It satisfies $DP_1=P_1\operatorname{diag}(\text{four-cycle},-1)$.
The cap therefore has the same cellular mapping-torus model
used to prove \eqref{maps:eq:F.6}.  With the unchanged twist $w_2=-e_0$,
the boundary coordinates remain $\theta=a+4k$, $h=k$;
the suspension representative restricts to $w$ in the base
direction and to $-\operatorname{Norm}_4(w)=0$ in the
auxiliary fibre direction.  Consequently the coefficient in
the torsion differential is again one.

\item The new vertical representatives in lexicographic exterior
bases are
\[
\begin{split}
 q=2:\ &(0,-1,-1,-1,0),\\
 q=3:\ &(0,0,1,0,0,0,0,0,0,1),\\
 q=4:\ &(0,0,0,0,0,0,0,-1,-1,0).
\end{split}
\]
Each has residue one in its global coinvariant
$\mathbb Z/2$.  Thus all three torsion differentials in \eqref{maps:eq:F.6}
are isomorphisms.  The integral Leray table and its free
degree-one differential are therefore exactly those of
\ref{maps:src:F.2.6}.  They give $H^2=H^3=0$ and $H^4=\mathbb Z^4$.
Universal coefficients, Hurewicz and Poincar\'e duality give
the stated integral homology and three-connectivity.

\item The relative almost-complex extension over the pair of pants
in \cref{pub:V:smooth-gluing} applies: the finite matrices have the same even
real eigenspace dimensions and the cusp is the same marked
Kuhn type.  The purely smooth characteristic-vector and
normal-invariant argument in \ref{maps:src:F.4.1} then gives the stated
oriented diffeomorphism type.

\item Finally, Theorem~\ref{maps:src:F.2.16} gives nonselfduality of the irreducible rank-five quotient at $v=1$.
These calculations use the explicit monomial marking displayed above.

\end{enumerate}
The corresponding exact calculations are implemented in \hyperlink{prog:V:second-pair}{the programs in Table~\ref*{tab:V:programs}}.
\end{proof}




\begin{theorem}[A residual obstruction for every lattice]\hypertarget{tag-02N8}{}\label{maps:src:F.2.19}
\StatementRecord{maps:src:F.2.19}{proved}{02N8}{\cref{maps:src:F.2.16}, \eqref{maps:eq:F.U6}, \cref{maps:src:F.2.6}, \cref{maps:src:F.2.18}, \eqref{maps:eq:F.U12}, \cref{maps:src:F.2.17}, \eqref{maps:eq:F.U1}}


Let a rational rank-six datum have finite linear orders five
and four, a primitive maximal cusp on a stable integral lattice,
and invariant line $L$, with irreducible rank-five quotient $W$
and no rational global coinvariants.  Suppose that at a prime
$p\ne2,5$ the semisimplified reduction of a stable lattice in
$W$ is the sum of the trivial module and an absolutely
irreducible module of dimension four.  Then no stable integral
lattice in that rational datum, with primitive maximal cusp,
and no free finite affine twists can give
$\pi_1=1$ and $H_2=0$ simultaneously.

For every $v\notin\{0,1\}$, the rational family in \ref{maps:src:F.2.16}
satisfies this hypothesis at each prime dividing
$v^2-v+1$.  Consequently the exclusion of those parameters
holds for every stable lattice and every twist.  The two remaining parameters give the three-connected manifolds of \cref{maps:src:F.2.6,maps:src:F.2.18}.

\end{theorem}


\begin{proof}
\begin{enumerate}[label=\arabic*.]
\item We first give the lattice-independent mechanism.  For any
stable lattice $V_\mathbb Z$, use its primitive invariant
line $L_\mathbb Z$ and quotient lattice $W_\mathbb Z$.
The semisimplified reduction of a rational representation is
independent of its stable lattice.  One elementary way to see
this is to compare two commensurable lattices by their finite
quotient $T$.  The exact sequence obtained by reduction modulo
$p$ has end terms $T[p]$ and $T/pT$, which have the same
composition factors by the filtration of $T$ by powers of $p$.
Hence the two lattice reductions have equal classes in the
Grothendieck group of finite-length modules.  Chaining such
comparisons proves the assertion for any two lattices.

\item Thus $W_\mathbb Z/pW_\mathbb Z$ has length two, with one
trivial constituent and one irreducible four-dimensional
constituent.  It has either a trivial submodule or a trivial
quotient (and possibly both).

\item Suppose first that it has a trivial submodule.  Its inverse
image in $V_\mathbb Z/pV_\mathbb Z$ is an extension of a
trivial module by a trivial module.  Each finite generator
acts on this two-dimensional space by a unipotent matrix.
Since its order is five or four, prime to $p$, it acts as
the identity.  There are therefore at least two global
invariant vectors in degree-one fibre cohomology modulo $p$.
Write their dimension as $a_p\ge2$; this is also the
dimension of the global coinvariants of fibre homology over
$\mathbb F_p$, by the perfect pairing between fibre homology
and cohomology.

\item The abelianized van Kampen presentation now shows
\begin{equation}
 b_1(M;\mathbb F_p)\ge a_p-1.
 \hypertarget{tag-02IN}{}\label{maps:eq:F.U12}
\end{equation}
Indeed, after the monodromy relations the fibre contributes
$a_p$ dimensions, and the two finite meridians contribute
two more.  The two finite power relations and the single
base relation impose at most three linear relations.
The cusp introduces no further fibre relation in this
abelianization.  On the dual homology lattice its nilpotent
operator is still primitive and maximal, so its angular lattice is
$\ker(M_0-I)=\operatorname{im}(M_0-I)$ by primitivity, and
the latter image already vanishes among the global
monodromy relations.  Its base meridian is the base relation
just counted.  Changing affine twists changes coefficients
of these three relations, not their number.  Thus \eqref{maps:eq:F.U12}
contradicts simple connectivity.

\item Suppose instead that the reduction of $W$ has a trivial
quotient.  Its integral global coinvariants have nonzero
$p$-primary torsion, since they have rational rank zero.
The natural map from the global coinvariants of $V$ to
those of $W$ is an isomorphism.  To check its kernel, a
primitive generator $e$ of $L_\mathbb Z$ lies in
$\ker(M_0-I)=\operatorname{im}(M_0-I)$, and hence is already
zero among global coinvariants.  Therefore $H^2(R^1)$ has
nonzero $p$-primary torsion.  If the total space were simply
connected, the finite-cap argument of \ref{maps:src:F.2.17} would make
every possible torsion source in $H^0(R^2)$ supported only
at two and five.  The invariant part of $\bigwedge^2V$ is
zero by irreducibility of the odd-dimensional quotient.
Consequently the nonzero $p$-primary subgroup survives to
$H^3(M;\mathbb Z)$ and then to
$\operatorname{Tor}H_2(M;\mathbb Z)$, again a contradiction.
This proves the general criterion.

\item For its application to \ref{maps:src:F.2.16}, let $p\mid q=v^2-v+1$.
We have $p\ne2,5$.  In the original cycle lattice of \eqref{maps:eq:F.U1},
$Y\mathbf1=2q(-1,1-v)^t$, so $\mathbf1$ gives a trivial
submodule after reduction.  On the four-dimensional quotient
the first finite generator has characteristic polynomial
$\Phi_5$, and $C=B^{-1}$ has polynomial
$(x+1)^2(x^2+1)$.  This quotient is absolutely irreducible.
Here is a characteristic-independent proof over
$\overline{\mathbb F}_p$.

\item A one-dimensional composition factor would make the
eigenvalues of $A$ and $C$ equal, because $CA^{-1}$ is
unipotent.  An eigenvalue of order dividing both five and
four must be one, which is absent from $\Phi_5$.
If the quotient were reducible, its only remaining possible
composition-factor dimensions would be two and two.
On each factor, the determinants of $A$ and $C$ must agree
and have orders dividing five and four, so both determinants
are one.  The two $C$ eigenvalue pairs must therefore be
$(-1,-1)$ and $(i,-i)$ in some order.  On the former
factor $C=-I$, since its order is prime to $p$.
Then the unipotence of $CA^{-1}=-A^{-1}$ would force the
eigenvalues of $A$ to be $-1$, impossible for $\Phi_5$.
This proves absolute irreducibility and establishes the
required residual composition factors at every $p\mid q$.
When $v\notin\{0,1\}$, the integer $q$ is greater than one
and has such a prime divisor.  The criterion applies.

\end{enumerate}
\end{proof}




\begin{theorem}[The complete rational $(5,4)$ normal form]\hypertarget{tag-02NA}{}\label{maps:src:F.2.20}
\StatementRecord{maps:src:F.2.20}{proved}{02NA}{\cref{maps:src:F.2.11}, \eqref{maps:eq:F.R4}, \eqref{maps:eq:F.R6}, \eqref{maps:eq:F.R7}, \eqref{maps:eq:F.R8}}


Consider ordered rational rank-five pairs $(A,C)$ with
\[
 \chi_A=(x-1)\Phi_5,
 \qquad \chi_C=(x-1)(x+1)^2(x^2+1).
\]
Put
\[
 N=CA^{-1}-I,
 \qquad N^2=0,\quad\operatorname{rank}N=2,
\]
where $A^5=C^4=I$ and the pair is irreducible over $\mathbb Q$.
Their simultaneous rational conjugacy classes are parametrized
bijectively by
\begin{equation}
 \bigl\{[r:s:t]\in\mathbb P^2(\mathbb Q):D(r,s,t)\ne0\bigr\}
 \ \sqcup\ \mathbb A^1(\mathbb Q),
 \hypertarget{tag-02HJ}{}\label{maps:eq:F.R1}
\end{equation}
where
\begin{equation}
\begin{split}
 D={}&r^3-4r^2s-r^2t+7rs^2+rst-3rt^2\\
     &{}-7s^3+11s^2t-11st^2+7t^3.
\end{split}
 \hypertarget{tag-02HP}{}\label{maps:eq:F.R2}
\end{equation}
The second summand lies over the single exceptional point
$[4:1:3]$ of $D=0$.  The construction below is explicit and
uses only a ten-variable linear system.  By \ref{maps:src:F.2.11} it also
parametrizes the rational rank-six maximal-cusp extensions
of all these pairs, with $B=C^{-1}$.
This is a rational classification; the stable integral lattices, the affine twists and the homology of the gluings are treated below.

\end{theorem}


\begin{proof}
\begin{enumerate}[label=\arabic*.]
\item Write $K=\mathbb Q(\zeta)$, $\zeta^4+\zeta^3+\zeta^2+
\zeta+1=0$, and identify the rational space with
$\mathbb Q\oplus K$.  The matrix of $A$ is multiplication by
$(1,\zeta)$:
\begin{equation}
 A=\operatorname{diag}(1,J),\qquad
 J=\begin{pmatrix}
 0&0&0&-1\\1&0&0&-1\\0&1&0&-1\\0&0&1&-1
 \end{pmatrix},\qquad R=(A+I)^{-1}.
 \hypertarget{tag-02HQ}{}\label{maps:eq:F.R3}
\end{equation}
Its rational centralizer is $\mathbb Q^*\times K^*$.
Let $I_N=\operatorname{im}N$.  If $X$ is any basis matrix
for this plane and $N=XY$, then $YX=0$.
The identity
\[
 C+I=(A+I)+XYA
\]
shows that its two-dimensional kernel has the same dimension
as the kernel of
$I_2+YARX=I_2-YRX$.  Hence
\begin{equation}
 YRX=I_2,
 \qquad I_N\cap RI_N=0.
 \hypertarget{tag-02HR}{}\label{maps:eq:F.R4}
\end{equation}
Conversely the two matrix equations in \eqref{maps:eq:F.R4}, together with
$YX=0$, give precisely that two-dimensional $-1$ eigenspace.

\item The plane $I_N$ projects nontrivially to $\mathbb Q$:
otherwise the first-coordinate covector is invariant under
both $A$ and $C$, contrary to irreducibility.
It contains no nonzero vector of $\mathbb Q\oplus0$,
because $R$ acts as $1/2$ on that line and \eqref{maps:eq:F.R4} would fail.
Write a basis of the plane as $(0,u),(1,v)$, where
$u\in K^*$.  Multiplication on $K$ by $u^{-1}$ puts it into
\[
 I_N=\langle(0,1),(1,\beta)\rangle_{\mathbb Q}.
\]
Changing the second basis vector by a rational multiple of
the first changes $\beta$ by a rational constant.
Two such planes are in the same centralizer orbit exactly
when $\beta'=a\beta+b$, with $a\in\mathbb Q^*$ and
$b\in\mathbb Q$.  Indeed a centralizer element carrying
one to the other must preserve the kernel of the projection
to $\mathbb Q$, so its $K^*$ component is itself rational;
the displayed affine change follows, and every such change
is realized.  We may therefore take
\begin{equation}
 \beta=r\zeta+s\zeta^2+t\zeta^3,
 \qquad [r:s:t]\in\mathbb P^2(\mathbb Q),
 \qquad
 X=\begin{pmatrix}0&1\\1&0\\0&r\\0&s\\0&t\end{pmatrix}.
 \hypertarget{tag-02HS}{}\label{maps:eq:F.R5}
\end{equation}
The triple is nonzero; a rational $\beta$ would make the
plane contain the forbidden first-coordinate line.
The stabilizer of such a plane in the centralizer consists
only of common scalar matrices.  Thus it introduces no
further identifications of the resulting nilpotent operators.

\item Write $Y$ as an unknown two-by-five matrix.  It is determined
by the following ten linear equations:
\begin{equation}
 YX=0,\qquad YRX=I_2,\qquad
 \operatorname{tr}(YAX)=-1,\qquad
 \operatorname{tr}(YA^{-1}X)=1.
 \hypertarget{tag-02HT}{}\label{maps:eq:F.R6}
\end{equation}
The first eight are $YX=0$ and $YRX=I_2$. The last two are exactly
$\operatorname{tr}C=\operatorname{tr}C^{-1}=-1$, since
$\operatorname{tr}A=\operatorname{tr}A^{-1}=0$ and
$C^{-1}=A^{-1}(I-XY)$.
They suffice to obtain the required characteristic
polynomial: the two $-1$ roots and determinant one leave a
monic cubic of constant term $-1$, and the two traces fix
that cubic to $(x-1)(x^2+1)$.
These remaining three roots are distinct, and the $-1$
eigenspace already has dimension two, so $C$ is semisimple
and $C^4=I$.  Also $YRX=I_2$ ensures that $N$ has rank two.
Thus every solution of \eqref{maps:eq:F.R6} gives all the asserted local
linear conditions.

\item In row-major order on the ten entries of $Y$, the determinant
of the coefficient matrix in \eqref{maps:eq:F.R6} is $D/4$.
For $D\ne0$, the column of the ten row-ordered entries of $Y$
is $E^{-1}b$, where $E$ is the coefficient matrix of the
displayed linear equations and $b$ is their right-hand side.
This gives the unique solution.

\item We now treat the exceptional locus.
Set
\[
 Q_1=r^2-rs-rt+s^2+t^2,
 \qquad Q_2=r^2-2rs-3rt+4s^2-4st+4t^2.
\]
The first column of the unique generic solution is
\begin{equation}
 y_{10}=\frac{2Q_1Q_2}{D},\qquad
 y_{20}=\frac{2(r-s-t)Q_1}{D}.
 \hypertarget{tag-02HU}{}\label{maps:eq:F.R7}
\end{equation}
Here $y_{10},y_{20}$ are the first column of $Y$.
The two equalities after multiplication by $D$ are polynomial
linear combinations of \eqref{maps:eq:F.R6}, by its adjugate matrix.
These polynomial identities remain necessary when $D=0$.
Moreover
\[
 2Q_1=(r-s)^2+(r-t)^2+s^2+t^2>0
\]
for every nonzero rational triple.  Consequently consistency
on $D=0$ forces $r=s+t$ and $Q_2=0$.
Substitution gives
\begin{equation}
 D(s+t,s,t)=(t-3s)(2t-s)^2,
 \qquad Q_2(s+t,s,t)=(3s-t)(s-2t).
 \hypertarget{tag-02HV}{}\label{maps:eq:F.R8}
\end{equation}
There are precisely two resulting projective points,
$[3:2:1]$ and $[4:1:3]$.
At the former, $[X\mid RX]$ has rank less than four, so
\eqref{maps:eq:F.R4} is impossible.  At the latter the coefficient and
augmented matrices of \eqref{maps:eq:F.R6} both have rank nine, and every
solution is
\begin{equation}
 Y_z=\begin{pmatrix}
 4z+16/5&0&-2z-1&4/5-2z&2z\\
 2z-2/5&0&-z&2/5-z&z
 \end{pmatrix},\qquad z\in\mathbb Q.
 \hypertarget{tag-02HW}{}\label{maps:eq:F.R9}
\end{equation}
Direct substitution proves both necessity and sufficiency.
This establishes the exceptional affine line in \eqref{maps:eq:F.R1}.

\item Every pair just constructed is irreducible over $\mathbb Q$.
An $A$-stable proper rational subspace has dimension one or
four.  A one-dimensional subrepresentation or quotient must
have trivial character, since its finite eigenvalues of
orders dividing five and four must agree under the
unipotent product.  The only $A$-invariant covector is the
first-coordinate covector.  Its product with $N$ is the
second row of $Y$, which is nonzero by $YRX=I_2$.
The only $A$-invariant vector is the first-coordinate vector;
its image under $N$ vanishes exactly if both entries in
\eqref{maps:eq:F.R7} vanish.  For $D\ne0$, positivity of $Q_1$ and
\eqref{maps:eq:F.R8} rule this out.  On the exceptional line the two
entries $4z+16/5$ and $2z-2/5$ never vanish simultaneously.
There are therefore neither invariant vectors nor covectors,
which excludes both possible dimensions.

\item Finally, rescaling $(r,s,t)$ is exactly the centralizer
change already described; it gives a conjugate pair.
Conversely any conjugacy of two normalized pairs centralizes
$A$ and hence identifies their projective parameters as proved
above.  At the exceptional point its stabilizer is scalar,
so distinct $z$ give distinct pairs.  This proves the stated
bijection, including completeness and uniqueness.

\end{enumerate}
The corresponding exact calculations are implemented in \hyperlink{prog:V:normal-forms}{the programs in Table~\ref*{tab:V:programs}}.
\end{proof}




\begin{corollary}[All good-prime lattices are forced]\hypertarget{tag-02NB}{}\label{maps:src:F.2.21}
\StatementRecord{maps:src:F.2.21}{proved}{02NB}{\cref{maps:src:F.2.20}, \cref{maps:src:F.2.19}}


For a $(5,4)$ datum with $\pi_1=1$ and $H_2=0$, the
rank-five quotient has absolutely irreducible reduction
at every prime $p\ne2,5$.  At each such prime its stable
$\mathbb Z_p$-lattice is unique up to homothety.
The same is true of the rank-six extension among lattices
with primitive maximal cusp.  Thus, once a rational pair
in \ref{maps:src:F.2.20} admits one three-connected datum, all remaining integral-lattice
ambiguity is supported at the primes two and five.


\end{corollary}


\begin{proof}
\begin{enumerate}[label=\arabic*.]
\item Work first over an algebraic closure of $\mathbb F_p$.
Both finite generators are semisimple because $p\ne2,5$.
A one-dimensional composition factor must be trivial:
its $A$ and $C$ eigenvalues agree by the unipotent product,
and their orders divide five and four.
There can be at most one such factor, since the $1$
eigenvalue of $A$ occurs once.
If it occurs, the remaining four-dimensional module is
absolutely irreducible by the spectral argument in \ref{maps:src:F.2.19}.
The only other possible reducible composition pattern is
$2+3$.
On each factor, determinants of $A$ and $C$ agree and
have orders dividing five and four, hence are one.
The two-dimensional $C$ eigenvalue pair is therefore
$(-1,-1)$ or $(i,-i)$.
The first would give $C=-I$ on that factor, contradicting
the unipotence of $CA^{-1}$ and the absence of $-1$
among the eigenvalues of $A$.
In the second case the three-dimensional factor has
$C$ eigenvalues $1,-1,-1$, so $C^2=I$ there; it has
an $A$-fixed vector because the other factor has determinant
one and cannot contain the simple $1$ eigenvalue of $A$.

\item Here is the required elementary involutory identity in the
present characteristic.  Put $B=A^{-1}$ on this factor.
From $(CB-I)^2=0$ and $C^2=I$ one obtains
\[
 BC+CB^{-1}=2I.
\]
For $Bv=v$ this says $(B+I)Cv=2v$.
The operator $B+I$ is invertible, since $B^5=I$ and
$p\ne2$, so $Cv=v$.  The resulting common fixed line
contradicts simplicity of the three-dimensional factor.
We have proved that any reducible residual representation
has exactly the constituents $1$ and an absolutely
irreducible four-dimensional module.  Theorem \ref{maps:src:F.2.19}
excludes that possibility for a three-connected datum, proving absolute
irreducibility at every indicated prime.

\item Let $R=\mathbb Z_p$ and $K_p=\mathbb Q_p$.
For two stable quotient lattices, multiply one by a power
of $p$ until it is contained in the other but not in its
$p$-multiple.  Its nonzero residual image must then be the
whole quotient by absolute irreducibility.  Nakayama's
lemma, or the Smith form of the inclusion, makes the
lattices equal.  This proves quotient-lattice uniqueness.

\item For completeness the extension-lattice comparison is also
integral.  Fix a primitive-cusp lattice $V_R$, with invariant
line $Re$ and quotient $W_R$.
The coboundary map
\[
 W_R^*\longrightarrow W_R^*\oplus W_R^*
\]
is primitive, because its reduction has trivial kernel.
Thus its cokernel, the degree-one cohomology of the free
group on the two generators, is torsion free, and an
integral coboundary over $K_p$ is already integral over $R$.
The extension class $\alpha$ is primitive in this cokernel.
Indeed, divisibility by $p$ would make the extra cusp map
\[
 \ker N_W/\operatorname{im}N_W\longrightarrow Re
\]
divisible by $p$, contrary to the three unit Smith factors
of the full cusp.  Here $N_W$ has primitive image:
it is the quotient of the primitive image of $N_V$ by
$Re\subset\operatorname{im}N_V$.

\item Normalize any other stable lattice by the already proved
quotient homothety so that its quotient is $W_R$.
Its invariant line is $p^rRe$ for some integer $r$.
The extension class in that basis is $p^{-r}\alpha$;
primitivity makes it integral only when $r\leq0$.
Primitivity of the coboundary map eliminates any additional
rational change of splitting.  The extra cusp map is then
multiplied by $p^{-r}$, so its Smith factor is a unit exactly
when $r=0$.  The lattices coincide after the common homothety.
This proves the assertion.

\end{enumerate}
\end{proof}




\begin{theorem}[A terminating lattice algorithm for each rational pair]\hypertarget{tag-02NC}{}\label{maps:src:F.2.22}
\StatementRecord{maps:src:F.2.22}{proved}{02NC}{\cref{maps:src:F.2.14}, \eqref{maps:eq:F.L3}, \eqref{maps:eq:F.L5}, \cref{maps:src:F.2.11}, \eqref{maps:eq:F.L6}, \eqref{maps:eq:F.L7}}


Let $(A,C)$ be any irreducible rational rank-five pair in
one of the eight spectral cases of \ref{maps:src:F.2.14}, with
$A^5=I$, $B=C^{-1}$, and $(AB-I)^2=0$ of rank two.
There is an explicit terminating algorithm which decides
whether the pair preserves an integral lattice, and, if so,
lists all such lattices up to simultaneous integral conjugacy.
It then lists all rank-six extensions with the same finite
orders and primitive maximal square-zero cusp, up to
simultaneous integral conjugacy.  The bounds below depend on the input rational pair.

\end{theorem}


\begin{proof}
\begin{enumerate}[label=\arabic*.]
\item We give the algorithm and its termination proof.
The fixed line of $A$ has dimension one.  Choose a nonzero
vector $e$ on it and a nonzero fixed covector $\varphi$.
Irreducibility implies that the vectors $ge$, with $g$ a
word in $A,C$, span the whole rational space; the analogous
statement holds for the covectors $\varphi h$.
In fact words of length at most four suffice.  The spans of
words of lengths at most $j$ increase from dimension one;
if one step does not increase the span, it is invariant
under both generators and hence already the whole space.
There are only four possible increases.

\item Choose words $g_1,\ldots,g_5$ and $h_1,\ldots,h_5$ giving
bases in this way, and put
\begin{equation}
 P=I+A+\cdots+A^4,\qquad T_{ij}=g_iPh_j.
 \hypertarget{tag-02GT}{}\label{maps:eq:F.L1}
\end{equation}
The rank-one operator $P$ is a nonzero scalar multiple of
$e\varphi$.  Hence the 25 matrices $T_{ij}$ are a rational
basis of $\operatorname{Mat}_5(\mathbb Q)$.
Each belongs to $\mathcal R=\mathbb Z[A,C]$.
Let $\mathcal R_0$ be their integral span and set
\begin{equation}
 \mathcal R_0^\#=
 \{X:\operatorname{tr}(XT)\in\mathbb Z
                     \text{ for all }T\in\mathcal R_0\}.
 \hypertarget{tag-02GV}{}\label{maps:eq:F.L2}
\end{equation}
The trace pairing on the full matrix algebra is
nondegenerate, as is seen from the matrix units, so this
is another full rational lattice, computed by inverting
the 25-by-25 Gram matrix.

\item If any stable lattice exists, every element of $\mathcal R$
acts integrally on it.  Consequently
$\operatorname{tr}(XY)\in\mathbb Z$ for all
$X,Y\in\mathcal R$, and necessarily
\begin{equation}
 \mathcal R_0\subset\mathcal R\subset\mathcal R_0^\#.
 \hypertarget{tag-02GW}{}\label{maps:eq:F.L3}
\end{equation}
A nonintegral Gram entry therefore proves nonexistence.
Otherwise let
$d=[\mathcal R_0^\#:\mathcal R_0]$, the absolute Gram
determinant.  Start with the lattice generated by
$\mathcal R_0$ and $I$, and repeatedly replace a current
lattice $L$ by $L+AL+CL$, using exact Hermite reduction.
If it leaves $\mathcal R_0^\#$, \eqref{maps:eq:F.L3} proves that there
is no stable lattice.
While it remains inside, every strict enlargement reduces
its index by a factor at least two.
There can be at most $\lfloor\log_2d\rfloor$ such
enlargements starting from $\mathcal R_0$, so a bound of
$1+\lfloor\log_2d\rfloor$ tests after adjoining $I$ is
sufficient.  Stabilization gives exactly $\mathcal R$:
the stable lattice contains $I$ and all words in $A,C$,
and every element introduced belongs to their ring.
It is then an order, and
$\Lambda_0=\mathcal R\mathbb Z^5$ is an explicitly
computed stable lattice.  This proves the decision step,
including termination in the nonexistence case.

\item Use a basis of $\Lambda_0$ from now on, so
$\mathcal R\subset\operatorname{Mat}_5(\mathbb Z)$.
Smith reduction gives an integer $c\ge1$ with
\begin{equation}
 c\operatorname{Mat}_5(\mathbb Z)\subset\mathcal R;
 \hypertarget{tag-02GX}{}\label{maps:eq:F.L4}
\end{equation}
one may take the largest Smith factor of the inclusion
of the order in the full matrix lattice.
Normalize any stable lattice by a rational homothety so
that it is contained in $\mathbb Z^5$ and the gcd of all
coordinates of a lattice basis is one.
The operators $cE_{ij}$ preserve it.  The coordinate values
of its vectors generate $\mathbb Z$, so Bezout's identity,
applied to $cE_{ij}$, shows that it contains every $ce_i$.
Therefore every normalized lattice lies in the finite interval
\begin{equation}
 c\mathbb Z^5\subset\Lambda\subset\mathbb Z^5.
 \hypertarget{tag-02GY}{}\label{maps:eq:F.L5}
\end{equation}
Conversely every lattice in this interval stable under
$A,C$ is included by the algorithm.
The only rational operators commuting with the pair are
scalars: they act by a scalar on the one-dimensional fixed
line of $A$, and then by the same scalar on the spanning
vectors $g_ie$.  Thus rational equivalence of two stable
lattices is precisely homothety.  The coordinate-gcd
normalization is unique, so \eqref{maps:eq:F.L5} gives the integral
conjugacy classification without further identifications.

\item The enumeration has the following finite bound.
For each $p^a\Vert c$, enumerate all subspaces of
$\mathbb F_p^5$ in reduced echelon form.
Beginning with $\mathbb Z^5$, pass through $a$ levels of
stable lattices between $pH$ and $H$, retaining only those
containing $p^a\mathbb Z^5$ and identifying equal Hermite
forms.  Every desired lattice occurs: for any target
$\Lambda$, the chain
$\Lambda+p^j\mathbb Z^5$, $0\le j\le a$, is such a
chain and ends at $\Lambda$.
If $s_p$ denotes the number of all subspaces of
$\mathbb F_p^5$, there are at most $s_p^a$ branches before
deduplication.  Intersect the resulting lattices for the
different primes and discard those of coordinate gcd
greater than one.  The Chinese remainder theorem proves
that this enumerates \eqref{maps:eq:F.L5} completely.

\item For each quotient lattice $W_{\mathbb Z}$ so obtained,
put $N_W=AB-I$.
Its image must be primitive: for a primitive full cusp,
$\operatorname{im}N_V=\ker N_V$ contains the primitive
invariant line, and quotienting that image by the line
gives a primitive image in $W_{\mathbb Z}$.
Reject a quotient whose two nonzero Smith factors are
not one.  For the remaining quotients use extension rows
$c_1,c_2\in\mathbb Z^{1\times5}$ and impose
\begin{equation}
\begin{gathered}
 c_1(I+A+\cdots+A^4)=0,
 \qquad c_2(I+B+\cdots+B^{m-1})=0,\\
 (c_2+c_1B)N_W=0,
 \qquad m=\operatorname{ord}(B).
\end{gathered}
 \hypertarget{tag-02GZ}{}\label{maps:eq:F.L6}
\end{equation}
These are integral linear equations.  Their integral kernel
modulo the coboundaries
\begin{equation}
 z\longmapsto\bigl(z(A-I),z(B-I)\bigr),
 \qquad z\in\mathbb Z^{1\times5},
 \hypertarget{tag-02H0}{}\label{maps:eq:F.L7}
\end{equation}
is a finitely generated group $G$ of rank one by \ref{maps:src:F.2.11}.
Compute its free and torsion factors by Smith reduction.
Choose an integral generator $w$ of the free rank-one
quotient $\ker N_W/\operatorname{im}N_W$.
The homomorphism
\begin{equation}
 \kappa:G\longrightarrow\mathbb Z,
 \qquad [(c_1,c_2)]\longmapsto(c_2+c_1B)w
 \hypertarget{tag-02H1}{}\label{maps:eq:F.L8}
\end{equation}
is well defined and nonzero: \eqref{maps:eq:F.L6} kills the image, and
\eqref{maps:eq:F.L7} kills $w$; rational nonvanishing is \ref{maps:src:F.2.11}.
The full cusp has unit Smith factors exactly when
$\kappa=\pm1$.
Indeed the quotient rank-two image is already primitive,
and the remaining invariant-line image is generated by
that one integer.
If $\operatorname{im}\kappa\ne\mathbb Z$, reject the
lattice.  Otherwise the inverse image of $1$ is a finite
coset of $\operatorname{Tor}G$, and lists all extensions.
Changing the sign of the invariant-line generator identifies
$\kappa=-1$ with $\kappa=1$.
There are no further identifications: an integral isomorphism
on the quotient is $\pm I$, on the invariant line is
$\pm1$, and, after $\kappa=1$ is fixed, their ratio is one.
The remaining row is exactly an integral change of splitting
\eqref{maps:eq:F.L7}.
This proves completeness and termination of the extension
step as well.



\end{enumerate}
The corresponding exact calculations are implemented in \hyperlink{prog:V:lattices}{the programs in Table~\ref*{tab:V:programs}}.
\end{proof}




\begin{theorem}[Lattice classes of the two three-connected members]\hypertarget{tag-02ND}{}\label{maps:src:F.2.23}
\StatementRecord{maps:src:F.2.23}{proved}{02ND}{\cref{maps:src:F.2.16}, \cref{maps:src:F.2.6}, \cref{maps:src:F.2.18}, \eqref{maps:eq:F.U1}, \eqref{maps:eq:F.L4}, \eqref{maps:eq:F.L5}, \eqref{maps:eq:F.L6}, \eqref{maps:eq:F.L7}, \cref{maps:src:F.2.12}, \eqref{maps:eq:F.L10}, \eqref{maps:eq:F.U6}, \cref{maps:src:F.2.19}}


For each of the two rational representations $v=0$ and $v=1$
of \ref{maps:src:F.2.16}, there are exactly two homothety classes of stable
rank-five quotient lattices, and exactly three integral
conjugacy classes of rank-six extensions with primitive
maximal cusp.  Exactly one of those three can have simply
connected finite affine fillings.  It is the integral class
already realized in \ref{maps:src:F.2.6}, respectively \ref{maps:src:F.2.18}.
Thus arbitrary stable-lattice changes produce no further
linear classes with three-connected gluing anywhere in the family \ref{maps:src:F.2.16}.
The affine twists and multiplicities are treated in \cref{maps:src:F.2.24}.

\end{theorem}


\begin{proof}
\begin{enumerate}[label=\arabic*.]
\item Apply the preceding algorithm in the original five-cycle
lattice \eqref{maps:eq:F.U1}.  For both parameters the order is obtained
after four strict enlargements.  Its Smith factors in the
full matrix lattice are 21 factors one and four factors two;
thus the conductor $c=2$ in \eqref{maps:eq:F.L4} is exact.
Enumerate all $374$ subspaces of $\mathbb F_2^5$, including zero
and the whole space.
After imposing stability and coordinate gcd one, the two
lattice representatives are
\begin{equation}
 \mathbb Z^5,\qquad
 \Lambda_* =2\mathbb Z^5+\mathbb Z(1,1,1,1,1)^t.
 \hypertarget{tag-02H2}{}\label{maps:eq:F.L9}
\end{equation}
The latter has index $16$.

\item For $\mathbb Z^5$, the extension group in \eqref{maps:eq:F.L6}--\eqref{maps:eq:F.L7}
is $\mathbb Z$ and $\kappa$ is a unit on its generator.
For $\Lambda_*$, that group is
$\mathbb Z\oplus\mathbb Z/2$, again with primitive
$\kappa$.  Hence there are exactly one and two primitive
extensions, respectively.
For each extension, restrict the primitive global invariant homology
covector $\ell$ to the two finite fixed lattices. The resulting
positive divisibilities are
\begin{equation}
\begin{array}{c|c|c}
 \text{quotient lattice}&\text{extension class}&
       (\operatorname{div}_1,\operatorname{div}_2)\\\hline
 \mathbb Z^5&\text{unique}&(1,2)\\
 \Lambda_*&\text{first torsion coset}&(1,1)\\
 \Lambda_*&\text{second torsion coset}&(1,2).
\end{array}
 \hypertarget{tag-02GU}{}\label{maps:eq:F.L10}
\end{equation}
The table is identical for $v=0$ and $v=1$.

\item For a simply connected filling the determinant quotient
of \ref{maps:src:F.2.12} requires $\gcd(\ell(v_i),m_i)=1$ at each
finite cap.  The second affine multiplicity is divisible
by four.  A divisibility of two in \eqref{maps:eq:F.L10} therefore
excludes that class for every choice of twists and larger
multiplicities.
For the sole remaining class, solve the integral intertwiner
equations with the corresponding pair in \eqref{maps:eq:F.U6}.
They have a solution of determinant $\pm1$. The fillings of
\cref{maps:src:F.2.6,maps:src:F.2.18} for those pairs are three-connected.  \ref{maps:src:F.2.19} excludes every other
rational parameter for all lattices.
Together these facts prove all the claims, including the
absence of any additional linear class with three-connected gluing in the
unbounded family.

\end{enumerate}
The corresponding exact calculations are implemented in \hyperlink{prog:V:lattices}{the programs in Table~\ref*{tab:V:programs}}.
\end{proof}




\begin{theorem}[All three-connected twists in the unbounded family]\hypertarget{tag-02NE}{}\label{maps:src:F.2.24}
\StatementRecord{maps:src:F.2.24}{proved}{02NE}{\cref{maps:src:F.2.12}, \cref{maps:src:F.2.16}, \cref{maps:src:F.2.23}, \cref{maps:src:F.2.20}, \cref{maps:src:F.2.19}, \eqref{maps:eq:F.U6}, \cref{pub:V:smooth-gluing}, \cref{maps:src:F.2.18}, \eqref{maps:eq:F.T4}, \cref{maps:src:F.2.6}, \cref{pub:V:integral-restrictions}, \eqref{maps:eq:F.T6}, \eqref{maps:eq:F.T7}, \eqref{maps:eq:F.T5}, \eqref{maps:eq:F.T8}, \eqref{maps:eq:F.T3}, \eqref{maps:eq:F.T2}, \eqref{maps:eq:F.T9}}


Use the fixed Kuhn cusp and the affine translation clutching
convention of \ref{maps:src:F.2.12}.  In the entire rational family \ref{maps:src:F.2.16},
allow every stable lattice, every free affine finite lift,
and every pair of finite multiplicities.
The data with $\pi_1=1$ and $H_2=H_3=0$ are precisely the
following, up to the integral linear classification already
proved and the translation identifications described below:
\begin{equation}
\begin{gathered}
 v\in\{0,1\},\qquad
 \text{the unique eligible integral class of \ref{maps:src:F.2.23}},\\
 (m_1,m_2)=(5,4),\qquad
 20\kappa+4r_1+5r_2=\pm1.
\end{gathered}
 \hypertarget{tag-02I7}{}\label{maps:eq:F.T1}
\end{equation}
Here $r_i=\ell(v_i)$ is the invariant affine numerator, and
$\kappa\in\mathbb Z$ is the invariant coordinate of the
cusp translation clutching, with signs fixed by the displayed
van Kampen relation below.  The remaining fixed-direction
components of the twists impose no condition.
Every manifold in \eqref{maps:eq:F.T1} is oriented diffeomorphic to its corresponding known example and hence to
$3\mathbb{HP}^2\#(-\mathbb{HP}^2)$.

\end{theorem}


\begin{proof}
\begin{enumerate}[label=\arabic*.]
\item \ref{maps:src:F.2.19} excludes all other rational parameters, and \ref{maps:src:F.2.23}
excludes all other integral linear classes.
We therefore work with either known integral pair, in the
marking of \eqref{maps:eq:F.U6}.  Its primitive invariant homology
covector is $\ell(x)=x_0$.
Each finite fixed lattice has an integral basis
$u_i,\lambda_i$, where $\ell(u_i)=1$,
$\ell(\lambda_i)=0$, and $\lambda_i$ is primitive in
$\ker\ell$.
We may write the affine numerator as
$v_i=r_i u_i+s_i\lambda_i$.
This covers arbitrary finite affine lifts: translating the
origin makes their translation vector invariant by averaging,
and its $m_i$th iterate is then the integral vector $v_i$.
The resulting boundary translation change is included in the
specified affine clutching.

\item For both linear pairs, the normal closure of the cusp angular
lattice in the fibre lattice is exactly $\ker\ell$, by
\cref{pub:V:smooth-gluing} and \ref{maps:src:F.2.18}.  Van Kampen therefore leaves the abelian
presentation
\begin{equation}
 m_1a=r_1h,\qquad m_2b=r_2h,\qquad
 a+b+\kappa h=0.
 \hypertarget{tag-02IC}{}\label{maps:eq:F.T2}
\end{equation}
This also defines the sign of $\kappa$.
Its determinant is, up to sign,
\begin{equation}
 p=m_1m_2\kappa+m_2r_1+m_1r_2.
 \hypertarget{tag-02ID}{}\label{maps:eq:F.T3}
\end{equation}
Thus simple connectivity is exactly $p=\pm1$.
It forces $\gcd(m_1,m_2)=1$ and
$\gcd(r_i,m_i)=1$ by reduction modulo each multiplicity.

\item We first prove that a three-connected datum has the minimal
multiplicities.  Set $n_1=5$, $n_2=4$ and $m_i=n_i k_i$.
In the splitting $\Lambda=\mathbb Zu_i\oplus\ker\ell$,
write the finite linear map as $1\oplus D_i$.
The reduced finite fibre is a mapping torus of $T^5$ over
the circle $m_i\ell$; its fibre is connected because
$r_i$ is prime to $m_i$.
Its vertical monodromy is a power of $D_i$ prime to $n_i$,
possibly followed by a translation, which changes no
integral Wang group.

\item The primitive fixed vertical loop $\lambda_i$ gives the
only additional free two-cycle beyond vertical cycles.
Consequently an invariant integral two-form $\omega$ on
the covering torus descends integrally exactly when
\begin{equation}
 \frac{\omega(v_i,\lambda_i)}{m_i}
       =\frac{r_i\omega(u_i,\lambda_i)}{m_i}
       \quad\text{is an integer}.
 \hypertarget{tag-02IE}{}\label{maps:eq:F.T4}
\end{equation}
For precision, the torus made from the affine generator
and $\lambda_i$ projects with degree $r_i$ to the mapping-
torus base, whereas the ordinary $u_i$ suspension projects
with degree $m_i$.  A Bezout combination is a primitive
suspension.  Its integrality condition is equivalent to
\eqref{maps:eq:F.T4}, because the ordinary covering periods are already
integral and $\gcd(r_i,m_i)=1$.
The Wang sequence shows that these vertical cycles and one
primitive suspension generate the free second homology.

\item The integral invariant two-form computations in
\cref{pub:V:integral-restrictions} and \ref{maps:src:F.2.18} give
\begin{equation}
 \gcd_{\omega}\omega(u_i,\lambda_i)=n_i.
 \hypertarget{tag-02IF}{}\label{maps:eq:F.T5}
\end{equation}
Hence the covering-image cokernel in degree two at the
$i$th finite point is exactly $\mathbb Z/k_i$.
The global invariant group of $\bigwedge^2\mathcal V$
and its integral first parabolic group both vanish.
The skyscraper cokernel sequence therefore gives
\begin{equation}
 H^1(B,R^2g_*\mathbb Z)
       =\mathbb Z/k_1\oplus\mathbb Z/k_2.
 \hypertarget{tag-02IG}{}\label{maps:eq:F.T6}
\end{equation}
Point-supported torsion kernels do not change this degree-one
base cohomology, and the cusp specialization is onto.
The term \eqref{maps:eq:F.T6}, in bidegree $(1,2)$, survives every Leray
differential on the two-dimensional base.
All rational terms of total degree three vanish, independently
of the multiplicities and twists, by the same exterior
local-system table as \ref{maps:src:F.2.6}.  Thus $H^3$ is finite.
If either $k_i>1$, its surviving nonzero filtration quotient
makes $H^3$ nonzero; universal coefficients then give
nonzero torsion in $H_2$.
This proves $m_1=5$, $m_2=4$ for every three-connected datum.

\item We now identify all twists with these multiplicities.
The following integral fact permits removal of each $s_i$:
\begin{equation}
 \operatorname{Norm}_{n_i}:
 \ker\ell\longrightarrow\mathbb Z\lambda_i
 \quad\text{is onto}.
 \hypertarget{tag-02IH}{}\label{maps:eq:F.T7}
\end{equation}
Here the norm is $I+D_i+\cdots+D_i^{n_i-1}$.
One can read \eqref{maps:eq:F.T7} directly from \eqref{maps:eq:F.T5}.
Let $\varphi_i$ generate the primitive invariant covector
lattice on $\ker\ell$.
The splitting by $u_i$ shows that the gcd in \eqref{maps:eq:F.T5} is
$|\varphi_i(\lambda_i)|$: its forms are
$\ell\wedge\varphi_i$ plus invariant vertical forms.
Thus that value is $n_i$.
The norm has the form
$x\mapsto d_i\varphi_i(x)\lambda_i$ for an integer
$d_i$; evaluating $\varphi_i$ gives
$|d_i|n_i=n_i$, hence $|d_i|=1$.
This proves surjectivity. Choose integral preimages
$q_i\in\ker\ell$ with
$\operatorname{Norm}_{n_i}q_i=\lambda_i$.

\item Put
\[
 w_i=-\frac{s_i}{n_i}\sum_{j=0}^{n_i-1}jM_i^j q_i.
\]
Then
\begin{equation}
 (I-M_i)w_i=-s_i\lambda_i/n_i+s_iq_i.
 \hypertarget{tag-02II}{}\label{maps:eq:F.T8}
\end{equation}
Translation by $w_i$ conjugates the finite affine action
with numerator $r_i u_i+s_i\lambda_i$ to that with
numerator $r_i u_i$: the remaining $s_iq_i$ is integral.
It gives a diffeomorphism of the entire finite cap.
In the standard linear boundary marking, its translation
section is
$w_i+\theta s_i\lambda_i/n_i$.
Its endpoint mismatch with the linear deck map is exactly
$s_iq_i$ by \eqref{maps:eq:F.T8}, so its integral boundary residue has
$\ell$-coordinate zero.

\item We explain why these local changes extend simultaneously,
including the prescribed cusp marking.
Let $P$ be the complementary pair of pants and
$\mathcal T\to P$ its flat torus group bundle.
A section on the boundary extends over the relative
one-skeleton because the torus is connected.
The only remaining obstruction is the two-cell obstruction
with values in the fibre fundamental-group local system:
\begin{equation}
 H^2(P,\partial P;\Lambda)
        \simeq H_0(P;\Lambda)=\Lambda_\Gamma.
 \hypertarget{tag-02IJ}{}\label{maps:eq:F.T9}
\end{equation}
The equality follows from the oriented relative cellular
complex, equivalently integral Poincar\'e--Lefschetz duality
with local coefficients.  Its boundary residues map to
the indicated global coinvariants.
For both linear pairs the horizontal monodromy relation
matrix has five unit Smith factors, so
$\Lambda_\Gamma=\mathbb Z$, detected exactly by $\ell$.
Thus the residues $s_iq_i$, with the identity section at the
cusp, have zero obstruction and extend over $P$.
A section may be smoothed relative to its boundary collars;
translation by it is a fibre-preserving diffeomorphism.
Together with the finite-cap maps it removes all $s_i$
without changing the cusp.
The same argument removes all components of translation
clutching whose $\ell$-coordinate is zero.

\item It remains to handle the integer numerators and the remaining
cusp coordinate.  If $p=-1$, the congruences from \eqref{maps:eq:F.T3}
give
\begin{equation}
 r_1=1+5a,\qquad r_2=-1+4b,
 \qquad a+b+\kappa=0.
 \hypertarget{tag-02I8}{}\label{maps:eq:F.T10}
\end{equation}
Adding $5a u_1$ or $4b u_2$ changes no finite affine action;
in its linear boundary marking it changes the translation
residue by $a u_1$ or $b u_2$.
The difference of the cusp marking has residue of
$\ell$-coordinate $\kappa$, by \eqref{maps:eq:F.T2}.
Their total obstruction in \eqref{maps:eq:F.T9} is therefore
$a+b+\kappa=0$.
The same extension argument identifies the glued manifold
with the original twists $(u_1,-u_2)$ and zero cusp
clutching.
If $p=1$, the identical calculation starts from
$(-u_1,u_2)$.

\item Finally simultaneous fibre inversion identifies these two
sign choices by an orientation-preserving diffeomorphism.
It conjugates every affine translation to its negative.
It extends over the Kuhn cusp: the lattice involution
$(x,h)\mapsto(-x,h)$ preserves its toric fan, since a
simplex given by a chain in a unit cube is taken to the
chain with reversed coordinate order in the opposite cube.
It normalizes the cusp lattice action and therefore descends
to its quotient, while fixing the base parameter.  The resulting toric map is holomorphic,
and in the regular real six-torus inversion has determinant
$(-1)^6=1$.  All translation maps also preserve orientation.
Thus every datum in \eqref{maps:eq:F.T1} gives the stated oriented smooth
type and has $H_2=H_3=0$, proving sufficiency as well as
necessity.

\end{enumerate}
The corresponding exact calculations are implemented in \hyperlink{prog:V:twists}{the programs in Table~\ref*{tab:V:programs}}.
\end{proof}




\section{Three-connectivity and affine twists}\label{pub:V:sec-connectivity}

\begin{theorem}[The universal degree-two suspension obstruction]\hypertarget{tag-02NF}{}\label{maps:src:F.2.25}
\StatementRecord{maps:src:F.2.25}{proved}{02NF}{\cref{maps:src:F.2.14}, \cref{maps:src:F.2.7}, \cref{maps:src:F.2.9}, \cref{maps:src:F.2.17}, \cref{maps:src:F.2.12}, \eqref{maps:eq:F.T4}, \eqref{maps:eq:F.S3}, \eqref{maps:eq:F.S6}, \cref{maps:src:E.1.3}, \eqref{maps:eq:F.S2}, \eqref{maps:eq:F.S7}, \eqref{maps:eq:F.S5}, \eqref{maps:eq:F.S8}}


Every admissible datum whose gluing has $\pi_1=1$ and $H_2=H_3=0$ has
exactly two genuine finite fillings and has minimal affine
multiplicities.  After ordering the order-five point first,
its finite characteristic polynomials are one of precisely
the following two remaining necessary possibilities:
\begin{equation}
\begin{array}{c|c|c}
 (m_1,m_2)&\chi_{M_1}&\chi_{M_2}\\\hline
 (5,4)&(x-1)^2\Phi_5&(x-1)^2\Phi_2^2\Phi_4\\
 (5,6)&(x-1)^2\Phi_5&(x-1)^2\Phi_2^2\Phi_3.
\end{array}
 \hypertarget{tag-02HX}{}\label{maps:eq:F.S1}
\end{equation}
In particular all order-eight and order-twelve spectral
loci, and the other two order-six spectral loci of \ref{maps:src:F.2.14},
are excluded for every rational relative position, every
stable lattice and every free affine twist.

\end{theorem}


\begin{proof}
\begin{enumerate}[label=\arabic*.]
\item Theorem \ref{maps:src:F.2.7} excludes three genuine finite fillings.
For two fillings, \ref{maps:src:F.2.9} and \ref{maps:src:F.2.14} give a primitive global
invariant homology covector $\ell$, finite fixed rank two,
and an irreducible rank-five cohomology quotient.
In particular
\begin{equation}
 H^0(B,j_*\bigwedge^2\mathcal V)=0.
 \hypertarget{tag-02HY}{}\label{maps:eq:F.S2}
\end{equation}
This is also immediate from the exterior filtration and
irreducibility in odd dimension, as in \ref{maps:src:F.2.17}.
Simple connectivity gives the determinant quotient of
\ref{maps:src:F.2.12}, so each invariant affine numerator $r_i=\ell(v_i)$
is prime to the multiplicity $m_i$, and the two
multiplicities are coprime.
As before one may conjugate any finite affine lift by a
translation so that its numerator $v_i$ is integral and
fixed by its linear part.

\item Consider one finite point.  Let $M$ have linear order $n$,
so $n\mid m$, and let its complementary degree-four factor
be $f(x)$, with
$\chi_M=(x-1)^2f(x)$.
The image of $\ell$ on the finite fixed lattice is $g\mathbb Z$
for some $g\mid n$: averaging $I+M+\cdots+M^{n-1}$ on
any vector of $\ell$-value one proves this divisibility.
Since $g\mid r$ and $\gcd(r,m)=1$, we have $g=1$.
Choose a fixed integral vector $u$ with $\ell(u)=1$.
It gives an integral invariant splitting
\begin{equation}
 \Lambda=\mathbb Zu\oplus H,
 \qquad H=\ker\ell,\qquad M=1\oplus D,
 \qquad\chi_D=(x-1)f(x).
 \hypertarget{tag-02HZ}{}\label{maps:eq:F.S3}
\end{equation}
The order of $D$ is $n$.
Let $\lambda$ generate its primitive fixed lattice and
let $\varphi$ generate its primitive invariant covector
lattice.  Choose their signs so that
\begin{equation}
 h=\varphi(\lambda)>0.
 \hypertarget{tag-02I0}{}\label{maps:eq:F.S4}
\end{equation}
Their rational pairing is nonzero, since the finite map
is semisimple and its averaging projection has rank one.

\item Both $\operatorname{Norm}_n(D)$ and $f(D)$ are integral
operators, have image on the line $\mathbb Q\lambda$,
and have invariant defining covector proportional to
$\varphi$.  Primitivity of $\lambda$ and $\varphi$
therefore implies
\begin{equation}
 \operatorname{Norm}_n(D)=\frac n h\lambda\varphi,
 \qquad
 f(D)=\frac{f(1)}h\lambda\varphi,
 \qquad h\mid n,\quad h\mid f(1).
 \hypertarget{tag-02I1}{}\label{maps:eq:F.S5}
\end{equation}
For example, the integer coefficient in the second formula
is seen by applying the integral operator to a vector on
which $\varphi$ is one; evaluating on $\lambda$ determines
that coefficient.  

\item Write $v=ru+s\lambda$.  The quotient finite fibre is a
connected $T^5$ mapping torus over the circle $m\ell$,
because $r$ is prime to $m$.
Its vertical cycles already have integral periods for an
invariant integral two-form on the covering torus.
The Wang sequence leaves one additional primitive free
suspension cycle, from $\lambda$.
Exactly as in \eqref{maps:eq:F.T4}, descent of an invariant form $\omega$
is equivalent to
$m\mid r\omega(u,\lambda)$, and hence to
$m\mid\omega(u,\lambda)$.
In the integral splitting \eqref{maps:eq:F.S3}, every invariant two-form
is an integral multiple of $\ell\wedge\varphi$ plus an
invariant vertical two-form.  Its possible values on
$(u,\lambda)$ are therefore exactly $h\mathbb Z$.
It follows that the degree-two specialization image at
this point has cokernel
\begin{equation}
 C^2_p\simeq\mathbb Z/(m/h).
 \hypertarget{tag-02I2}{}\label{maps:eq:F.S6}
\end{equation}
The integer $m/h$ is defined because $h\mid n\mid m$.
Torsion in the finite fibre cohomology is a specialization
kernel.
The primitive suspension and vertical cycles generate the
free second homology, so \eqref{maps:eq:F.S6} is the full cokernel.

\item Let $I^2$ be the image sheaf of
$R^2g_*\mathbb Z\to j_*\bigwedge^2\mathcal V$.
The cusp specialization is onto by \cref{maps:src:E.1.3}, and all kernels
and cokernels at finite points are skyscraper sheaves.
Thus $H^1(R^2g_*\mathbb Z)=H^1(I^2)$.
Equation \eqref{maps:eq:F.S2} and the cokernel exact sequence give an
injection
\begin{equation}
 \bigoplus_{p\text{ finite}}C^2_p
           \lhook\joinrel\longrightarrow H^1(B,R^2g_*\mathbb Z).
 \hypertarget{tag-02I3}{}\label{maps:eq:F.S7}
\end{equation}
The Leray term of bidegree $(1,2)$ has no incoming or
outgoing differential on a real two-dimensional base.
If $H_2=H_3=0$, the universal coefficient theorem gives
$H^3=0$, so that term must vanish.
Every cyclic cokernel in \eqref{maps:eq:F.S7} is therefore zero.
Combining \eqref{maps:eq:F.S5}--\eqref{maps:eq:F.S6} yields the universal necessary
condition
\begin{equation}
 m=h=n,\qquad n\mid f(1).
 \hypertarget{tag-02I4}{}\label{maps:eq:F.S8}
\end{equation}
This both fixes the multiplicity and tests the spectrum.

\item Apply \eqref{maps:eq:F.S8} to the eight-case
list of \ref{maps:src:F.2.14}.  The order-five factor gives $\Phi_5(1)=5$.
At the other point the exact values are
\begin{equation}
\begin{array}{c|c|r|r}
 n&f&f(1)&n/\gcd(n,f(1))\\\hline
 4&\Phi_2^2\Phi_4&8&1\\
 6&\Phi_2^2\Phi_3&12&1\\
 6&\Phi_2^2\Phi_6&4&3\\
 6&\Phi_3\Phi_6&3&2\\
 8&\Phi_8&2&4\\
 12&\Phi_{12}&1&12\\
 12&\Phi_3\Phi_4&6&2\\
 12&\Phi_4\Phi_6&2&6.
\end{array}
 \hypertarget{tag-02I5}{}\label{maps:eq:F.S9}
\end{equation}
The last column is a lower bound for the nonzero local
covering-cokernel order already at the smallest allowed
multiplicity.  Only the first two rows pass.
This proves \eqref{maps:eq:F.S1} and the stated universal exclusions.

\end{enumerate}
The corresponding exact calculations are implemented in \hyperlink{prog:V:suspension}{the programs in Table~\ref*{tab:V:programs}}.
\end{proof}




\begin{theorem}[The complete rational remaining $(5,6)$ locus]\hypertarget{tag-02NG}{}\label{maps:src:F.2.26}
\StatementRecord{maps:src:F.2.26}{proved}{02NG}{\eqref{maps:eq:F.S1}, \eqref{maps:eq:F.R2}, \cref{maps:src:F.2.25}, \eqref{maps:eq:F.R3}, \eqref{maps:eq:F.R5}, \cref{maps:src:F.2.20}, \eqref{maps:eq:F.R6}, \eqref{maps:eq:F.R7}, \eqref{maps:eq:F.R8}, \eqref{maps:eq:F.R12}, \cref{maps:src:F.2.22}}


The irreducible rational rank-five pairs in the second row
of \eqref{maps:eq:F.S1}, with square-zero rank-two cusp, are classified
up to simultaneous rational conjugacy by
\begin{equation}
 \{[r:s:t]\in\mathbb P^2(\mathbb Q):
             D(r,s,t)\ne0,\ Q_-(r,s,t)\ne0\}
       \ \sqcup\ \mathbb A^1(\mathbb Q),
 \qquad Q_-=r^2-rs+st,
 \hypertarget{tag-02HK}{}\label{maps:eq:F.R10}
\end{equation}
where $D$ is the same cubic as in \eqref{maps:eq:F.R2}.
The exceptional line again lies over $[4:1:3]$.
Thus the two rational spectral loci left by \ref{maps:src:F.2.25} both
have complete explicit normal forms, including their
entire rational exceptional loci.

\end{theorem}


\begin{proof}
\begin{enumerate}[label=\arabic*.]
\item Use precisely the space $\mathbb Q\oplus\mathbb Q(\zeta_5)$,
matrices $A,R,X$ and projective plane normalization in
\eqref{maps:eq:F.R3}--\eqref{maps:eq:F.R5}.  The centralizer and image-plane arguments
in \ref{maps:src:F.2.20} require only irreducibility and the two-dimensional
$-1$ eigenspace of $C$; both hold here.
Consequently every pair has a unique projective image-plane
parameter, with only scalar stabilizer.
The unknown $Y$ now satisfies
\begin{equation}
 YX=0,\qquad YRX=I_2,\qquad
 \operatorname{tr}(YAX)=-2,\qquad
 \operatorname{tr}(YA^{-1}X)=2.
 \hypertarget{tag-02HL}{}\label{maps:eq:F.R11}
\end{equation}
Its coefficient matrix is unchanged from \eqref{maps:eq:F.R6}, so has
determinant $D/4$.
The two traces and the $-1$ eigenspace imply
\[
 \chi_C=(x-1)(x+1)^2(x^2+x+1).
\]
Its other three roots are distinct, and the $-1$ eigenspace
has its full multiplicity, so $C^6=I$.
The cusp $XY$ still has square zero and rank two.
For $D\ne0$, the ten linear equations give the unique
explicit solution.

\item The two first-column adjugate identities become
\begin{equation}
 D y_{10}=2Q_-Q_2,
 \qquad D y_{20}=2(r-s-t)Q_-,
 \hypertarget{tag-02HM}{}\label{maps:eq:F.R12}
\end{equation}
with the same $Q_2$ as in \eqref{maps:eq:F.R7}.
For $D\ne0$, these entries vanish simultaneously exactly
when $Q_-=0$.
Indeed, if $Q_-\ne0$ they would instead force
$r=s+t$ and $Q_2=0$, and \eqref{maps:eq:F.R8} would give $D=0$.
The invariant-covector argument of \ref{maps:src:F.2.20} is unchanged.
Thus the pair is rationally irreducible exactly off the
conic $Q_-=0$ in the generic locus.

\item No nonzero rational triple satisfies $D=Q_-=0$.
If $s=0$, the conic forces $r=0$, and then $D=7t^3$.
If $s\ne0$, rescale to $s=1$ and put
$r=u$, $t=u-u^2$.  Direct substitution gives
\begin{equation}
 D(u,1,u-u^2)=
 -(7u^2-11u+7)(u^4-u^3+u^2-u+1).
 \hypertarget{tag-02HN}{}\label{maps:eq:F.R13}
\end{equation}
The first factor has discriminant $-75$ and positive
leading coefficient.  The second is strictly positive,
as the identity
\[
 u^4-u^3+u^2-u+1
 =\left(u^2-\frac u2\right)^2
  +\frac34\left(u-\frac23\right)^2+\frac23
\]
shows.  In particular neither has a rational zero.
On the exceptional cubic the necessary identities \eqref{maps:eq:F.R12}
therefore force the same two points $[3:2:1]$ and $[4:1:3]$
as before.  The first is again impossible because
$[X\mid RX]$ has rank less than four.
At the second, the complete solution is
\begin{equation}
 Y_z=\begin{pmatrix}
 4z+26/5&0&-2z-2&-2z-1/5&2z+1\\
 2z-2/5&0&-z&2/5-z&z
 \end{pmatrix},\qquad z\in\mathbb Q.
 \hypertarget{tag-02HO}{}\label{maps:eq:F.R14}
\end{equation}
Its first-column entries cannot vanish simultaneously,
so each pair on this line is irreducible.
There are no further identifications, by the scalar
stabilizer argument already proved in \ref{maps:src:F.2.20}.
This establishes the claimed bijection.

\item The finite lattice and extension algorithm of \ref{maps:src:F.2.22}
enumerates the integral lattices and extensions over every point of
\eqref{maps:eq:F.R10}.

\end{enumerate}
The corresponding exact calculations are implemented in \hyperlink{prog:V:normal-forms}{the programs in Table~\ref*{tab:V:programs}}.
\end{proof}




\begin{theorem}[The entire remaining order-six locus is excluded]\hypertarget{tag-02NH}{}\label{maps:src:F.2.27}
\StatementRecord{maps:src:F.2.27}{proved}{02NH}{\eqref{maps:eq:F.S1}, \cref{maps:src:F.2.25}, \eqref{maps:eq:F.S3}, \eqref{maps:eq:F.S5}, \eqref{maps:eq:F.X2}, \eqref{maps:eq:F.X3}, \eqref{maps:eq:F.S7}, \eqref{maps:eq:F.X5}}


Every simply connected datum in the second row of \eqref{maps:eq:F.S1},
at every rational relative position and every stable integral
lattice with primitive maximal cusp, has
\[
                 \operatorname{Tor}_3 H_3(M;\mathbb Z)\ne0
\]
if its degree-two suspension cokernels vanish.
Consequently no such datum is three-connected.  Together with
\ref{maps:src:F.2.25} this leaves only the minimal order pair $(5,4)$ as a possible three-connected admissible datum.
\end{theorem}


\begin{proof}
\begin{enumerate}[label=\arabic*.]
\item We use cohomology with coefficients in the discrete valuation
ring $R=\mathbb Z_3$.  Suppose the degree-two suspension
cokernels vanish; otherwise \ref{maps:src:F.2.25} already gives nonzero
$H_2$ or $H_3$.  That theorem forces $m=n=h=6$ at the
order-six point, in the notation \eqref{maps:eq:F.S3}--\eqref{maps:eq:F.S5}.
We first determine its local integral representation over $R$.

\item The splitting by the fixed vector $u$ is integral:
$\Lambda_R=Ru\oplus H_R$.  On $H_R$, the finite map $D$
has characteristic polynomial
$(x-1)(x+1)^2\Phi_3(x)$.  Since two is a unit in $R$,
the idempotents $(I\pm D^3)/2$ give an integral direct sum
\[
 H_R=L_+\oplus L_-,\qquad
 \operatorname{rank}L_+=3,\quad
 \operatorname{rank}L_-=2.
\]
On $L_-$ the map is exactly $-I$, by its rational spectrum.
On $L_+$ it has order three.  The primitive fixed vector
$\lambda$ lies in $L_+$, and the primitive invariant
covector $\varphi$ vanishes on $L_-$, so their restrictions
remain primitive.  Equation \eqref{maps:eq:F.S5} with $h=6$ gives
\begin{equation}
 2(I+D+D^2)|_{L_+}=\lambda\varphi|_{L_+}.
 \hypertarget{tag-02IV}{}\label{maps:eq:F.X1}
\end{equation}
This operator is nonzero modulo three.  But modulo three
$I+D+D^2=(D-I)^2$ and $(D-I)^3=0$.  The reduction of $D$
on the rank-three space therefore has one Jordan block
of length three.  Choose $q\in L_+$ whose image has
$(D-I)^2q\ne0$ modulo three.  Then $q,Dq,D^2q$ are a
basis modulo three and hence an $R$-basis.  In this basis
$D$ is the three-cycle permutation matrix, since $D^3=I$.
Thus the full finite representation over $R$ is
\begin{equation}
                   1\oplus\operatorname{Perm}_3\oplus(-I_2).
 \hypertarget{tag-02IW}{}\label{maps:eq:F.X2}
\end{equation}
Its dual has the same form, by reversing the three-cycle.

\item Let $u,a,b,c,d,e$ be a cohomology basis for \eqref{maps:eq:F.X2}, with
$a,b,c$ cycled and $d,e$ negated.  A full integral basis
of the invariant degree-three lattice is
\begin{equation}
 u(ab+bc+ca),\qquad abc,\qquad ude,\qquad(a+b+c)de.
 \hypertarget{tag-02IX}{}\label{maps:eq:F.X3}
\end{equation}
Here juxtaposition denotes exterior multiplication.
To check completeness, terms with an odd number of
$d,e$ are negated by the cube of the finite map and
have no invariants.  The terms with zero or two such
factors give exactly the four orbit sums in \eqref{maps:eq:F.X3},
each primitive.  With volume $uabcde$, their skew
pairing has one block with coefficient three and one
block with coefficient a unit.  Thus its determinant
has three-adic valuation two. Exterior multiplication in the displayed
integral basis gives Smith factors $(1,1,3,3)$.

\item Let $B_6$ be the reduced finite fibre.  Its torus
cover has degree six.  Its free middle cohomology has
rank four and unimodular skew pairing, by integral
Poincar\'e duality.  If $I$ denotes the index of its
pullback lattice in the invariant lattice, multiplication
of pairings by the covering degree gives
\begin{equation}
 2v_3(I)+2=4v_3(6)=4,
 \qquad v_3(I)=1.
 \hypertarget{tag-02IY}{}\label{maps:eq:F.X4}
\end{equation}
Consequently the three-primary degree-three covering
cokernel is $R/3R$.  Translations of the affine action
do not enter this calculation: the cover degree,
linear invariant lattice and oriented pairings are
those of the free quotient for every allowed
affine numerator.

\item We explain why this local group must appear in total
homology even when the global integral lattice varies.
Let $\mathcal V$ be the rational rank-six cohomology
local system.  It is an extension of the irreducible
rank-five system $W$ by a trivial line.  The determinant
of $W$ is trivial.  Irreducibility in odd dimension
implies that neither $\bigwedge^2 W$ nor its dual
has a nonzero invariant: an invariant alternating
form would have an invariant radical and would have
to be nondegenerate, impossible in odd dimension.
Also $\bigwedge^3W\simeq\bigwedge^2W^*$ by its
trivial determinant.  The exterior filtration therefore
gives
\begin{equation}
 (\bigwedge^3\mathcal V)^\Gamma=0,\qquad
 ((\bigwedge^2\mathcal V)^*)^\Gamma=0.
 \hypertarget{tag-02IZ}{}\label{maps:eq:F.X5}
\end{equation}
These rational vanishing statements remain true over
$\mathbb Q_3$, since they are kernels of rational
matrices.

\item As in \eqref{maps:eq:F.S7}, take the image of
$R^3g_*R\to j_*\bigwedge^3\mathcal V_R$.
Its kernel and cokernel are point-supported sheaves.
The first vanishing in \eqref{maps:eq:F.X5} makes the local group
$R/3R$ inject into $H^1(B,R^3g_*R)$.  Bidegree $(1,3)$
has no incoming or outgoing Leray differential on
the two-dimensional base, so this injection survives
to $E_\infty^{1,3}$.

\item The second vanishing in \eqref{maps:eq:F.X5} says
that $H^2(B,R^2g_*R)$ has rational rank zero.
Point-supported changes leave degree-two base cohomology
unchanged. Put $\mathcal L=\bigwedge^2\mathcal V_R$ on the
punctured base $B^\circ$. The quotient $j_*\mathcal L/j_!\mathcal L$
is supported at the punctures, so its exact sequence and integral
Poincar\'e duality give
\[
 H^2(B,j_*\mathcal L)=H_c^2(B^\circ,\mathcal L)
 =H_0(B^\circ,\mathcal L)=(\mathcal L_b)_\Gamma,
 \qquad b\in B^\circ.
\]
Thus the group is a finitely generated torsion $R$-module and
hence a finite group.  Therefore
$E_\infty^{2,2}$ is finite as well.  In the filtration
of $H^4(M;R)$, take the inverse image of the injected
$R/3R\subset E_\infty^{1,3}$.  It is an extension
of $R/3R$ by the finite group $E_\infty^{2,2}$.
It is thus a nonzero finite three-primary subgroup
of $H^4(M;R)$.
Universal coefficients identify this torsion with
the three-primary torsion of $H_3(M;\mathbb Z)$.
This proves the claim at every rational position and
every integral lattice, and completes the exclusion.

\end{enumerate}
The corresponding exact calculations are implemented in \hyperlink{prog:V:middle-pairing}{the programs in Table~\ref*{tab:V:programs}}.
\end{proof}




\begin{theorem}[Universal fundamental group and twist normalization]\hypertarget{tag-02NI}{}\label{maps:src:F.2.28}
\StatementRecord{maps:src:F.2.28}{proved}{02NI}{\cref{maps:src:F.2.25}, \eqref{maps:eq:F.S5}, \eqref{maps:eq:F.T8}, \cref{maps:src:F.2.24}}


In the remaining $(5,4)$ locus, fix the marked standard Kuhn
cap and allow affine translation clutching.
Let $\Lambda=\mathbb Z^6$ be fibre homology and let
$\Lambda_\Gamma$ be its global monodromy coinvariants.
Write $\Lambda_\Gamma=\mathbb Zh\oplus T$, with $T$ finite,
and let $\ell$ be the primitive invariant covector detecting
the free factor.  The following are complete reductions for
any datum that could satisfy $\pi_1=1$ and $H_2=H_3=0$.

The multiplicities must be $(5,4)$, and the local
vertical fixed pairings must be $h_1=5,h_2=4$, as in \ref{maps:src:F.2.25}.
For any affine twists and clutching the fundamental group
is abelian.  If its free presentation determinant is nonzero,
its order is
\begin{equation}
 |\pi_1|=|20\kappa+4r_1+5r_2|\,|T|,
 \qquad r_i=\ell(v_i).
 \hypertarget{tag-02I9}{}\label{maps:eq:F.T11}
\end{equation}
If that determinant is zero the group is infinite.
In particular simple connectivity is equivalent to
\begin{equation}
 T=0,\qquad 20\kappa+4r_1+5r_2=\pm1.
 \hypertarget{tag-02IA}{}\label{maps:eq:F.T12}
\end{equation}
Whenever these conditions and $h_1=5,h_2=4$ hold, every
such twist choice is oriented diffeomorphic to the normalized
filling with twists $(u_1,-u_2)$ and zero cusp clutching,
where $u_i$ is any integral fixed vector with $\ell(u_i)=1$.
Thus each eligible integral linear datum has exactly one
normalized smooth filling to test for the remaining
homology conditions.

\end{theorem}


\begin{proof}
\begin{enumerate}[label=\arabic*.]
\item The image of the fibre fundamental group is a normal
abelian subgroup in the van Kampen presentation.
After quotienting by it, the two finite meridians satisfy
$a^5=b^4=ab=1$.
That quotient is trivial, because five and four are coprime.
The entire fundamental group is therefore equal to the
image of the fibre group and is abelian.

\item In abelianization the monodromy relations give
$\Lambda_\Gamma$.
The primitive cusp angular lattice is the image of its
nilpotent operator, and this image is already zero among
those monodromy relations.
The remaining presentation is
\begin{equation}
 5a=v_1,\qquad4b=v_2,\qquad a+b+v_0=0
 \quad\text{in }\Lambda_\Gamma\oplus\mathbb Za\oplus\mathbb Zb,
 \hypertarget{tag-02IB}{}\label{maps:eq:F.T13}
\end{equation}
where $\ell(v_0)=\kappa$ fixes the sign convention.
Use any Smith presentation of $T$ and split the free part
$\mathbb Zh$.
The full relation matrix is block triangular, with the
presentation of $T$ as one diagonal block and the
three-by-three free matrix of determinant
$\pm(20\kappa+4r_1+5r_2)$ as the other.
Its determinant gives \eqref{maps:eq:F.T11}; when the latter free
matrix has rank less than three the group has positive
rank and is infinite.
For determinant $\pm1$, elimination of the three free
variables leaves precisely $T$, even if the twists have
nonzero torsion components.
This proves the exact simple-connectivity criterion.

\item We give the global twist comparison to make its scope
precise.  Under \eqref{maps:eq:F.T12}, the finite fixed covector image
is one at each point, so choose integral fixed $u_i$ with
$\ell(u_i)=1$.  The primitive vertical fixed vector
$\lambda_i$ and invariant covector $\varphi_i$ satisfy
$\varphi_i(\lambda_i)=m_i$, by the assumed maximal local
pairings.  Formula \eqref{maps:eq:F.S5} makes the finite norm onto
$\mathbb Z\lambda_i$.
Consequently every vertical numerator component
$s_i\lambda_i$ is removed by the explicit cap translation
\eqref{maps:eq:F.T8}, whose boundary residue has $\ell$-coordinate zero.

\item For the complementary pair of pants the full relative
obstruction group for extending boundary translation
sections is
$H^2(P,\partial P;\Lambda)=\Lambda_\Gamma$.
Since $T=0$, this group is $\mathbb Z$, detected by $\ell$.
Prescribing the identity translation at the cusp therefore
allows all the vertical finite changes to extend over $P$.
The same argument removes the zero-$\ell$ component of
any cusp translation clutching.
For determinant $-1$, write
$r_1=1+5a$, $r_2=-1+4b$; then \eqref{maps:eq:F.T12} says
$a+b+\kappa=0$.
The two integer changes of finite boundary marking and
the remaining cusp change thus also have zero total
relative obstruction and extend.
This gives the normalized twists $(u_1,-u_2)$.
For determinant $1$ the same argument gives their negatives.
The orientation-preserving fibre inversion extending over
the Kuhn fan, proved in \ref{maps:src:F.2.24}, identifies the two choices.
All maps used in the comparison are cap translations or
translations by smooth torus-bundle sections over $P$;
each comparison respects the prescribed cusp marking.
This proves the claimed complete twist normalization.

\end{enumerate}
\end{proof}




\begin{theorem}[A finite integral test for three-connectivity]\hypertarget{tag-02NJ}{}\label{maps:src:F.2.29}
\StatementRecord{maps:src:F.2.29}{proved}{02NJ}{\cref{maps:src:F.2.20}, \cref{maps:src:F.2.22}, \cref{maps:src:F.2.25}, \cref{maps:src:F.2.28}, \cref{maps:src:F.2.6}, \eqref{maps:eq:F.HC7}, \cref{maps:src:E.1.3}, \cref{maps:src:F.2.16}}


Take any rational pair in \ref{maps:src:F.2.20} and any of its integral,
primitive-cusp rank-six extensions given by \ref{maps:src:F.2.22}.
Write $M_1,M_2$ for geometric transport on fibre homology
$\Lambda=\mathbb Z^6$, in the peripheral convention
$M_0M_1M_2=I$.
First require the preliminary conditions
\begin{equation}
 \Lambda_\Gamma=\mathbb Z,\qquad
 \ell(\Lambda^{M_i})=\mathbb Z,\qquad h_1=5,\quad h_2=4 .
 \hypertarget{tag-02G9}{}\label{maps:eq:F.HC1}
\end{equation}
Here $\ell$ is the primitive global invariant covector, and
$h_i$ is the integral fixed-vector/fixed-covector pairing
on $\ker\ell$ from \ref{maps:src:F.2.25}.
Failure of any condition excludes the datum.
If they hold, use the unique normalized affine filling of
\ref{maps:src:F.2.28}, with twists $(u_1,-u_2)$ and zero cusp clutching.
The following six finite integer-matrix tests are necessary
and sufficient for this filling to have $H_2=H_3=0$.

Let
\begin{equation}
 T_i=M_i^{-t},\quad T_0=(T_1T_2)^{-1},\qquad
 G_q=\operatorname{coker}
 [\,\bigwedge^qT_1-I\mid\bigwedge^qT_2-I\,]
 \quad(1\le q\le4).
 \hypertarget{tag-02GE}{}\label{maps:eq:F.HC2}
\end{equation}
All $G_q$ are finite.  Let $P_q$ be the integral parabolic
cohomology group of the exterior local system
$\bigwedge^q\mathcal V$, for $q=2,4$.
Let $t_q\in G_{q-1}$, $2\le q\le5$, be the natural global
coinvariant image of the unique nonzero torsion suspension
class at the order-four cap, as specified below.
Let $c_4\in G_3$ be the natural image of the unique nonzero
torsion class in
$\operatorname{coker}(\bigwedge^3T_0-I)$.
Then the six tests are
\begin{equation}
\begin{array}{c|ccc}
 &\text{first test}&\text{second test}&\text{third test}\\ \hline
 H^3=0&G_1=\langle t_2\rangle&P_2=0&t_3\ne0\\
 H^5=0&G_3=\langle t_4,c_4\rangle&P_4=0&t_5\ne0 .
\end{array}
\hypertarget{tag-02GF}{}\label{maps:eq:F.HC3}
\end{equation}
Every group, class, and equality here is computed by
finitely many integer kernels, Smith forms, and exterior
minors.  This is a decision criterion for every rational parameter and every integral lattice in the normal form.

\end{theorem}


\begin{proof}
\begin{enumerate}[label=\arabic*.]
\item We establish the uniform local calculations for the criterion.
The fixed vectors $u_i$ split the homology lattice as
$\mathbb Zu_i\oplus H_i$, where $H_i=\ker\ell$.
Write $D_i=M_i|_{H_i}$.
The maximal pairing $h_i=m_i$ gives
\begin{equation}
 \operatorname{Norm}_{m_i}(D_i)=\lambda_i\varphi_i ,
 \hypertarget{tag-02GG}{}\label{maps:eq:F.HC4}
\end{equation}
with both factors primitive.

\item At five, work over $\mathbb Z_5$.
Modulo five the norm is $(D_1-I)^4$ and is nonzero.
Choose $q$ with $\operatorname{Norm}_5q$ primitive.
The five vectors $q,D_1q,\ldots,D_1^4q$ are independent
modulo five: otherwise the nilpotent Jordan block generated
by $q$ would have length less than five.
They are therefore a $\mathbb Z_5$ basis.
Thus $D_1$ is the regular five-cycle permutation locally.

\item At four, work over $\mathbb Z_2$.
The norm is $(D_2-I)^3$ modulo two and is nonzero.
The vectors $q,D_2q,D_2^2q,D_2^3q$ similarly generate
a primitive regular four-cycle submodule $Q$ of rank four.
The quotient is rank one with $D_2=-1$, by the
characteristic polynomial.
For a lift $z$ of its generator, put $a=(D_2+I)z\in Q$.
Since $D_2^4=I$,
\[
 (D_2^3-D_2^2+D_2-I)a=0.
\]
On the regular four-cycle module the kernel of the
displayed alternating norm is exactly $(D_2+I)Q$,
integrally.  One can see this by writing its four
coefficients; equivalently $D_2+I$ has three nonzero
Smith factors, all one, and the indicated rational image.
Subtracting a suitable element of $Q$ from $z$ makes
$D_2z=-z$.
Consequently
\begin{equation}
 M_1\sim 1\oplus\operatorname{perm}_5\quad\text{over }\mathbb Z_5,
 \qquad
 M_2\sim1\oplus\operatorname{perm}_4\oplus(-1)
 \quad\text{over }\mathbb Z_2 .
 \hypertarget{tag-02GH}{}\label{maps:eq:F.HC5}
\end{equation}
These local comparisons induce isomorphisms on cap cohomology, as follows.
Indeed the integral lattice of intertwiners is the kernel
of an integer matrix and its tensor with $\mathbb Z_p$
is the lattice of local intertwiners.
Approximate a local invertible intertwiner by an integral
one whose determinant is a $p$-adic unit.
Keeping the $u_i$ summand fixed produces an equivariant
torus isogeny of degree prime to $p$, preserving the
normalized translation $\pm u_i/m_i$.
The resulting cap and boundary coverings induce
isomorphisms on $\mathbb Z_p$ cohomology: the vertical
torus maps are cohomology isomorphisms over $\mathbb Z_p$,
so the mapping-torus Wang sequences give the assertion.
At primes not dividing $m_i$, transfer already makes
finite-cover pullback an isomorphism onto invariants.

\item Here are the full required integral results for these
two monomial models.
The ranks of their invariant exterior lattices in
degrees zero through six are
$(1,2,3,4,3,2,1)$.
The order-five cap has no cohomology torsion in degrees
two through five.
The order-four cap has exactly $\mathbb Z/2$ in each
of those four degrees.
This follows from the Wang sequence for the mapping
torus of the vertical $D_i$:
its torsion in degree $q$ is
\begin{equation}
 \operatorname{Tor}\operatorname{coker}
       (\bigwedge^{q-1}D_i^t-I).
 \hypertarget{tag-02GI}{}\label{maps:eq:F.HC6}
\end{equation}
In particular it introduces no additional torsion primes.

\item Integral pullback from either finite cap is onto the
invariant lattice in degrees two through five.
For degree two this is precisely the $m_i/h_i=1$
suspension calculation of \ref{maps:src:F.2.25}.
The determinants of the pairings between
the invariant lattices in complementary degrees
$(1,5),(2,4),(3,3)$ are respectively
$m_i,m_i^3,m_i^4$ in absolute value.
The degree-one cover image has index $m_i$.
The free downstairs complementary pairings are unimodular,
and pullback multiplies each pairing by the covering
degree $m_i$.
The middle pairing first gives degree-three index one.
The degree-two index one then gives degree-four index one;
the degree-one index $m_i$ gives degree-five index one.
For both local normal forms, exterior multiplication gives these
pairing determinants, and integer row and column operations give the
Wang Smith forms in exterior degrees zero through six.

\item We next identify the boundary maps.
Choose any unimodular matrix $P_2=[u_2,H_2]$ realizing
the splitting at four, and put
$J=P_2^{-t}|_{\mathbb Z^5}$ for the inclusion of vertical
cohomology.
Choose the unique nonzero torsion representative
\begin{equation}
 w_q\in\operatorname{coker}
            (\bigwedge^{q-1}D_2^t-I),\qquad
 t_q=[(\bigwedge^{q-1}J)w_q]\in G_{q-1}.
 \hypertarget{tag-02GJ}{}\label{maps:eq:F.HC7}
\end{equation}
A Smith transformation supplies an integral $w_q$.
Its norm is exactly zero, since it lies rationally in
the image of $\bigwedge^{q-1}D_2^t-I$.
The finite boundary coordinates of \ref{maps:src:F.2.6} satisfy
$\theta=a+4k$, $h=k$, hence $a=\theta-4h$.
The pullback of the suspension cocycle has $\theta$
component $w_q$ and $h$ component
$-\operatorname{Norm}_4(w_q)=0$.
This holds in every exterior degree, including $q=5$.
It is an integral cellular assertion for the monomial
model, where the torus cellular differentials vanish;
the prime-to-two isogeny comparison transports it to
every eligible cap.
The base connecting map therefore gives exactly the
class $t_q$ in \eqref{maps:eq:F.HC7}.
Different representative choices change it by a global
monodromy relation.  Any sign change is immaterial for
an order-two class.

\item There is one further boundary map, from the free
degree-four cusp kernel.
The primitive cusp is integrally conjugate to
$\left(\begin{smallmatrix}I&I\\0&I\end{smallmatrix}\right)$.
Its degree-three exterior difference has ten Smith
factors one and one Smith factor two; its cokernel is
$\mathbb Z^9\oplus\mathbb Z/2$.
The other exterior differences in degrees one, two,
four and five have torsion-free cokernels.
Let $C$ be the smooth cusp cap and $Y=\partial C$.
The boundary Wang sequence gives
$\operatorname{Tor}H^4(Y)=\mathbb Z/2$.
In the Poincar\'e--Lefschetz exact sequence its image in
\[
 H^5(C,Y)\cong H_3(C)=\mathbb Z^9
\]
is zero, so this torsion class lifts to $H^4(C)$.
Its restriction to an ordinary fibre is zero.
By \cref{maps:src:E.1.3} the kernel of integral specialization is
a primitive copy $K=\mathbb Z$ in degree four.
The lift is in $K$, and the image of this cyclic group
in $H^4(Y)$ contains the nonzero torsion class.
A cyclic image containing nonzero torsion cannot
also have a nonzero free part.
Its primitive generator therefore restricts to the
unique nonzero torsion class of $H^4(Y)$.
Consequently the cusp contribution to $d_2$ is exactly $c_4$
as defined above.

\item It remains to prove that these are all the tests.
For the irreducible odd-dimensional quotient $W$
of $\mathcal V$, invariant alternating forms and
bivectors vanish: a nonzero one would give an
isomorphism $W\cong W^*$ with an odd-dimensional
nondegenerate alternating pairing.
The exact exterior sequences for
$0\to\mathbb Q\to\mathcal V_{\mathbb Q}\to W\to0$,
together with $\det W=1$, give
\begin{equation}
 H^0(B,j_*\bigwedge^q\mathcal V)=0
       \quad(2\le q\le5),\qquad
 G_q\otimes\mathbb Q=0\quad(1\le q\le4).
 \hypertarget{tag-02GK}{}\label{maps:eq:F.HC8}
\end{equation}
For the endpoint assertions, the extension cannot split:
on $\mathbb Q\oplus W$ a square-zero nilpotent has rank at
most two, whereas the stipulated rank is three.
An invariant covector of $\mathcal V_{\mathbb Q}$ would
split this extension, since $W$ has no trivial quotient.
Thus $(\mathcal V_{\mathbb Q}^*)^\Gamma=0$, which gives
both $H^0(\bigwedge^5\mathcal V_{\mathbb Q})=0$ and
$G_1\otimes\mathbb Q=0$.
Here and below the first assertion is integral as well,
since the coefficient lattices are torsion free.
The local invariant ranks in exterior degrees two
and four are $3,3,9$ at the three punctures.
The Euler formula for the rank-fifteen local system
therefore gives $P_2\otimes\mathbb Q=P_4\otimes\mathbb Q=0$.
In degree one it similarly gives rational parabolic
cohomology zero.
The only global degree-one section is the invariant
line.  Simple connectivity and the universal
coefficient theorem make its free transgression
onto $H^2(B,\mathbb Z)$; equivalently the rational
second Betti number is zero, as in \ref{maps:src:F.2.6}.
Thus the normalized filling has rational $b_2=b_3=0$.

\item All specializations in exterior degrees two through
five are onto, by the local calculations and \cref{maps:src:E.1.3}.
Their only point-supported kernels are
\begin{equation}
\begin{array}{c|cccc}
q&2&3&4&5\\ \hline
\text{kernel}&\mathbb Z/2&\mathbb Z/2&
                 \mathbb Z\oplus\mathbb Z/2&\mathbb Z/2 .
\end{array}
\hypertarget{tag-02GL}{}\label{maps:eq:F.HC9}
\end{equation}
The free summand in degree four is at the cusp; the
others are at four.
Consequently $H^1(B,R^qg_*\mathbb Z)=P_q$ in the
degrees used here, while
$H^2(B,R^qg_*\mathbb Z)=G_q$.
The latter equality also holds in degree one:
changing a sheaf at isolated points changes no
degree-two cohomology.  The equality with coinvariants
follows from compactly supported top cohomology on
the punctured sphere and the sequence $j_!\to j_*$.

\item The base has real dimension two.
Thus the only possible nonzero differential in the
integral Leray spectral sequence is
$d_2:E_2^{0,q}\to E_2^{2,q-1}$.
The complete associated graded groups in total
degrees three and five are therefore
\begin{equation}
\begin{split}
 \operatorname{gr}H^3&:
   G_1/\langle t_2\rangle,\quad P_2,\quad
   \ker(\mathbb Z/2\xrightarrow{\,t_3\,}G_2),\\
 \operatorname{gr}H^5&:
   G_3/\langle t_4,c_4\rangle,\quad P_4,\quad
   \ker(\mathbb Z/2\xrightarrow{\,t_5\,}G_4).
\end{split}
\hypertarget{tag-02GA}{}\label{maps:eq:F.HC10}
\end{equation}
These are all finite.  A finite filtered group
vanishes exactly when every associated graded group
vanishes, regardless of extension classes.
Since $\pi_1=1$ and rational $b_2=b_3=0$, the universal
coefficient theorem gives
$H^3=\operatorname{Ext}(H_2,\mathbb Z)$, so
$H_2=0$ exactly when $H^3=0$.
Integral Poincar\'e duality gives $H^5\cong H_3$.
This proves \eqref{maps:eq:F.HC3}.

\item Given $A=\bigwedge^qT_1$, $B=\bigwedge^qT_2$ on a
free lattice $L$, form the saturated integer kernel
\begin{equation}
 K_q=\ker[\,A-I\mid A(B-I)\mid -(AB-I)\,]\subset L^3.
 \hypertarget{tag-02GB}{}\label{maps:eq:F.HC11}
\end{equation}
Its image under
$(x,y,z)\mapsto((A-I)x,(B-I)y)$ is precisely the
integral parabolic cocycle lattice.
Quotient by the columns of
$\left(\begin{smallmatrix}A-I\\B-I\end{smallmatrix}\right)$
to obtain $P_q$.
Use the integral images of the peripheral differences.
Smith transformations also supply every $w_q$, the
cusp order-two representative, and the tests for
their images and generated subgroups in $G_q$.

\item For the parameters $v=0,1$ of \ref{maps:src:F.2.16}, all six
conditions hold;
$G_1,\ldots,G_4$ are each $\mathbb Z/2$,
$P_2=P_4=0$, all four finite-cap images are nonzero,
and the cusp image is zero.
For $v=2$ they are each $\mathbb Z/6$,
$P_2=P_4=\mathbb Z/3$, and the first, second,
fourth and fifth conditions fail.

\item For the family matrices $A,B$ in \eqref{maps:eq:F.U6}, which
are geometric cohomology pullbacks, the geometric homology inputs
are $(A^t,B^t)$. Their left cohomology transport in
\eqref{maps:eq:F.HC2} is $(A^{-1},B^{-1})$.
For comparison with computations written using
$(A,B)$, simultaneous reversal of the two peripheral
generators sends
\begin{equation}
 g_1\mapsto g_1^{-1},\qquad
 g_2\mapsto g_2^{-1},\qquad
 g_0=(g_1g_2)^{-1}\mapsto
     g_2g_1=g_1^{-1}(g_1g_2)g_1 .
 \hypertarget{tag-02GC}{}\label{maps:eq:F.HC12}
\end{equation}
This is an orientation-reversing marking of the
pair of pants fixing its three ends: each peripheral
conjugacy class is inverted.
The image of the cusp generator is $g_1^{-1}g_0^{-1}g_1$.
Inversion preserves invariant and coinvariant
lattices, and the indicated conjugating generator
acts trivially on global coinvariants.
All boundary classes in \eqref{maps:eq:F.HC3} have order two, so
reversal signs do not change the tests.

\end{enumerate}
The corresponding exact calculations are implemented in \hyperlink{prog:V:cap-tests}{the programs in Table~\ref*{tab:V:programs}}.
\end{proof}




\begin{corollary}[Classification of the three-connected data]\hypertarget{tag-02NL}{}\label{maps:src:F.2.30}
\StatementRecord{maps:src:F.2.30}{proved}{02NL}{\cref{maps:src:F.2.20}, \cref{maps:src:F.2.22}, \eqref{maps:eq:F.HC1}, \eqref{maps:eq:F.HC3}, \eqref{maps:eq:F.T12}, \cref{maps:src:F.3.11}, \cref{maps:src:F.2.7}, \cref{maps:src:F.2.14}, \cref{maps:src:F.2.25}, \cref{maps:src:F.2.27}, \cref{maps:src:F.2.28}, \cref{maps:src:F.2.29}, \cref{maps:src:E.1.3}, \cref{maps:src:F.2.15}, \cref{pub:V:smooth-gluing}, \cref{sec:V-known-inputs}, \cref{maps:src:F.4.1}}


For the marked Kuhn cap and affine translation clutchings
used throughout this section, the following is a complete constructive
classification of the data with
$\pi_1=1$ and $H_2=H_3=0$, up to integral conjugacy of their
linear tuples and oriented diffeomorphism of their fillings.

There are no data with three genuine finite fillings.
With two fillings their affine orders are five and
four.  Order the order-five point first.
The complete, unbounded parameter set is the rational
projective normal form of \ref{maps:src:F.2.20} together with its
exceptional rational affine line.
For each parameter, run the terminating order,
conductor, stable-lattice, and primitive-extension algorithm
of \ref{maps:src:F.2.22}.  Discard precisely the integral classes failing
\eqref{maps:eq:F.HC1} or one of \eqref{maps:eq:F.HC3}.
The classes that pass these tests are exactly all
integral linear tuples with three-connected gluing.
For each such class all simply connected affine
twists are exactly the choices satisfying \eqref{maps:eq:F.T12};
their fillings are all oriented diffeomorphic to its
normalized filling.
If ordered labels with the order-four point first are
desired, interchange the finite labels and recompute
$M_0=(M_1M_2)^{-1}$.

Every such datum has
\begin{equation}
 \pi_1=1,\qquad
 H_j=\begin{cases}
  \mathbb Z,&j=0,8,\\
  \mathbb Z^4,&j=4,\\
  0,&\text{otherwise},
 \end{cases}
 \hypertarget{tag-02GD}{}\label{maps:eq:F.HC13}
\end{equation}
and its oriented smooth manifold is
$3\mathbb{HP}^2\#(-\mathbb{HP}^2)$.
It has bounded-cusp holomorphic periods in the sense of \ref{maps:src:F.3.11}.

The fundamental group and the integral homology of every other gluing are computed by \cref{maps:src:F.2.31,maps:src:F.2.35}.

\end{corollary}


\begin{proof}
\begin{enumerate}[label=\arabic*.]
\item The exclusions \ref{maps:src:F.2.7}, \ref{maps:src:F.2.14}, \ref{maps:src:F.2.25}, and \ref{maps:src:F.2.27}
leave exactly the stated minimal orders and rational
quotient problem.
The rational normal form is injective on conjugacy
classes, including the exceptional line.
The finite interval and extension-coset enumeration
in \ref{maps:src:F.2.22} gives every integral class over each rational
point, with all bounds proved for that input.
No distinct rational points can become integrally
conjugate, since an integral conjugacy induces a
rational quotient conjugacy.
Theorem \ref{maps:src:F.2.28} and the necessary-and-sufficient
criterion \ref{maps:src:F.2.29} prove both directions of the
retention rule and account for every affine twist.
This is also an exhaustive enumeration if desired:
enumerate primitive projective integer triples with
first nonzero coordinate positive by increasing height,
and enumerate reduced rational exceptional parameters
by increasing numerator and denominator height.
Every finite stage has finitely many inputs, each
terminates by \ref{maps:src:F.2.22}, and every rational input occurs.


\item Hurewicz makes these simply connected manifolds three-connected.  The only singular contribution to
Euler number is the Kuhn cusp, whose value is six
by \cref{maps:src:E.1.3}; torus bundles and the free finite caps have
Euler number zero.
Integral Poincar\'e duality and the universal coefficient
theorem then give \eqref{maps:eq:F.HC13}, including torsion freeness.

\item For the smooth classification, \ref{maps:src:F.2.15} supplies a
compatible complex structure at each finite cap, and
the cusp already has its complex model.
Choose their common fibre orientation to be the one
induced by the cusp.
Commuting complex structures at the finite caps can
be chosen in that component: their two-dimensional
fixed real summands allow the orientation to be
reversed when necessary.
The fibre of the bundle of orientation-compatible
complex structures retracts to $\mathbb{CP}^3$.
Its connectedness and simple connectivity extend
the prescribed boundary structures across the
two-dimensional pair of pants, exactly as in \cref{pub:V:smooth-gluing}.
A compatible horizontal distribution gives an
almost complex structure on the total manifold.
\Cref{sec:V-known-inputs} therefore gives signature two,
intersection lattice $I_{3,1}$, and $c_2^2=2$.
The vanishing $H^2=0$ gives $c_1=w_2=0$, so $p_1/2=-c_2$.
The spin Wu identity makes this class characteristic, and its square
is two. A nontrivial integral multiple of a lattice vector has square
divisible by the square of that integer; hence this norm-two class
is primitive.
The complete integral reflection, Wall-form and Kramer--Stolz argument in \ref{maps:src:F.4.1}
depends only on these almost complex characteristic
numbers and three-connectivity.
It applies without alteration and gives the claimed
exact oriented diffeomorphism.
Finally \ref{maps:src:F.3.11} applies to the resulting determinant-one
monodromy and primitive cusp.
It gives the punctured-base periods with bounded cusp correction.

\end{enumerate}
\end{proof}




\section{Fundamental groups and cellular models}\label{pub:V:sec-cellular}

\begin{theorem}[Complete fundamental-group data for every input]\hypertarget{tag-02NM}{}\label{maps:src:F.2.31}
\StatementRecord{maps:src:F.2.31}{proved}{02NM}{\cref{maps:src:E.1.3}; \cite[Theorems 2.3--2.4, p.~425; Euler-characteristic comparison, pp.~427--428]{rootScott83}}


Let $k=2$ or $3$, and let an integral rank-six tuple satisfy the
finite-order and primitive maximal-cusp conditions of \cref{def:admissible}.    Specify each finite affine action as
\begin{equation}
 \tau_i(x)=M_ix+v_i/m_i,\qquad v_i\in\Lambda=\mathbb Z^6,
 \quad \tau_i^{m_i}=1\text{ on }\Lambda_{\mathbb R}/\Lambda,
 \hypertarget{tag-02FO}{}\label{maps:eq:F.G1}
\end{equation}
and require its cyclic action to be free.  The numerator need not be
fixed by $M_i$.  Define its integral power vector and the cusp
lattice by
\begin{equation}
 w_i=\frac1{m_i}\sum_{j=0}^{m_i-1}M_i^jv_i\in\Lambda^{M_i},
 \qquad U=\operatorname{im}\bigl((M_1\cdots M_k)^{-1}-I\bigr).
 \hypertarget{tag-02FR}{}\label{maps:eq:F.G2}
\end{equation}
Use left conjugation transport for the peripheral generators $s_i$,
and record the integral translation clutching by $\kappa\in\Lambda$
in $s_1\cdots s_kx^\kappa=1$.  Then the full fundamental group is
\begin{equation}
 \left\langle x^\lambda,s_1,\ldots,s_k\ \left|\
 \begin{array}{l}
 x^\lambda x^\mu=x^{\lambda+\mu},\quad x^u=1\ (u\in U),\\
 s_ix^\lambda s_i^{-1}=x^{M_i\lambda},\quad
 s_i^{m_i}=x^{w_i},\quad s_1\cdots s_kx^\kappa=1
 \end{array}\right.\right\rangle .                         \hypertarget{tag-02FS}{}\label{maps:eq:F.G3}
\end{equation}
Only a basis of $\Lambda$ and of $U$ is needed in this finite
presentation. The formulas below determine the injected quotient
of the fibre subgroup.

Put
\begin{equation}
 L=\sum_{g\in\langle M_1,\ldots,M_k\rangle}gU,
 \qquad A=\Lambda/L.                                      \hypertarget{tag-02FT}{}\label{maps:eq:F.G4}
\end{equation}
Every element of $L$ is a finite integral sum of elements of the subgroups $gU$.  For two points, put
$d=\gcd(m_1,m_2)$, $M=M_1$, and
$S_j(M)=I+M+\cdots+M^{j-1}$.  Choose integers $a,b$ with
$am_1+bm_2=d$, and set, in $A$,
\begin{equation}
 \begin{split}
 w'_2&=-w_2-S_{m_2}(M)\kappa,\\
 e&=(m_2/d)w_1-(m_1/d)w'_2,\qquad
 z=aw_1+bw'_2,\qquad \overline A=A/\langle e\rangle .
 \end{split}                                             \hypertarget{tag-02FU}{}\label{maps:eq:F.G5}
\end{equation}
There is an exact sequence and a complete marked extension description
\begin{equation}
 1\longrightarrow\overline A\longrightarrow\pi_1
   \longrightarrow C_d\longrightarrow1,
 \qquad txt^{-1}=M(x),\quad t^d=z.                         \hypertarget{tag-02FV}{}\label{maps:eq:F.G6}
\end{equation}
Thus its order is $d|\overline A|$ if $\overline A$ is finite;
otherwise it is infinite.  In particular $d=1$ gives precisely
the abelian group $\overline A$.

For three genuine points, so $m_i\ge2$, let
\begin{equation}
 \Gamma=\langle q_1,q_2,q_3\mid q_i^{m_i}=1,
                         \ q_1q_2q_3=1\rangle,
 \qquad \epsilon=\sum_i1/m_i-1.                           \hypertarget{tag-02FW}{}\label{maps:eq:F.G7}
\end{equation}
The action $q_i\mapsto M_i$ is well defined on $A$.
If $\epsilon\le0$, the exact sequence is
\begin{equation}
                    1\longrightarrow A\longrightarrow\pi_1
                         \longrightarrow\Gamma\longrightarrow1. 
 \hypertarget{tag-02FX}{}\label{maps:eq:F.G8}
\end{equation}
If $\epsilon>0$, $\Gamma$ is the finite spherical triangle group.
Using its action on $A$, form the exact transfer vector
\begin{equation}
 e=\sum_{g\in\Gamma}g\kappa+
       \sum_{i=1}^3\ \sum_{g\in\Gamma/\langle q_i\rangle}g w_i.
 \hypertarget{tag-02FY}{}\label{maps:eq:F.G9}
\end{equation}
The second sums use left coset representatives and are independent
of their choice.  In this case the exact sequence has kernel
$A/\langle e\rangle$, with the same lifted relations \eqref{maps:eq:F.G3}.
Consequently \eqref{maps:eq:F.G3}, \eqref{maps:eq:F.G6}, \eqref{maps:eq:F.G8}, and \eqref{maps:eq:F.G9} give every such fundamental group by finite operations, including the precise
injected abelian kernel and all extension data.

For every input, the full first integral homology is the Smith
cokernel of the integer matrix
\begin{equation}
 \begin{pmatrix}
 M_1-I&\cdots&M_k-I&-w_1&\cdots&-w_k&\kappa\\
 0&\cdots&0&&\operatorname{diag}(m_1,\ldots,m_k)&&\boldsymbol1
 \end{pmatrix}.                                           \hypertarget{tag-02FP}{}\label{maps:eq:F.G10}
\end{equation}
Here the bottom left zero block has $k$ rows.

\end{theorem}


\begin{proof}
\begin{enumerate}[label=\arabic*.]
\item \emph{Finite and cusp caps.}
The affine lift on $\mathbb R^6$ has $m_i$th power translation by
$w_i$, as the displayed sum shows.  It commutes with that power,
so $M_iw_i=w_i$.  The group generated by lattice translations
and this lift is the fundamental group of the reduced flat
six-manifold.  Its presentation consists of the lattice, its
$s_i$ conjugation relations, and $s_i^{m_i}=w_i$.
The finite cap deformation retracts onto that reduced manifold
by the equivariant radial contraction of its root disc.
Thus these are the cap relations, even for a numerator
which is not an invariant lattice vector.

\item For the primitive Kuhn cap, the compact toric strata in its
normalization and all their intersection components are simply
connected.  Applying van Kampen successively to their tubular
neighborhoods identifies the fundamental group of the central
fibre with that of its dual three-torus.  The nearby fibre map
kills the angular lattice $U$ and sends the long-period lattice
$\Lambda/U$ isomorphically onto this dual-torus group.  This is
also the explicit toric collapse used in \cref{maps:src:E.1.3}.  The primitive
index-one condition gives $U=\ker(M_0-I)$ and makes $\Lambda/U$
free of rank three.  A section through a smooth central stratum
bounds the base meridian by its parameter disc.  The smooth cap
retracts onto the central fibre, so its filling imposes exactly these relations.

\item Over the punctured sphere the bundle has a section up to homotopy,
because the base retracts to a graph.  Its group is the appropriate
semidirect product.  Choose the cap sections and retain their
relative translation clutching as $\kappa$.  Van Kampen with the
boundary maps just described gives exactly \eqref{maps:eq:F.G3}.
The opposite transport convention is related by the following explicit isomorphism:
reverse the ordered list and replace each affine map by its
inverse, and replace $\kappa$ by $-\kappa$.  The new linear
matrices are $M_k^{-1},\ldots,M_1^{-1}$, the power vectors are
$-w_k,\ldots,-w_1$, and the cusp image lattice is unchanged.
The substitutions $s_i\mapsto s_{k+1-i}^{-1}$ give the explicit
isomorphism between these two presentations.  This supplies the
dictionary when geometric transport is recorded with opposite
path-composition order.

\item \emph{Integral closure and two-point elimination.}
In a normal abelian fibre subgroup, the normal closure of $U$
is precisely the integer sum \eqref{maps:eq:F.G4}.  Start with its column
Hermite basis and repeatedly append its images under every
$M_i$ and $M_i^{-1}$, then take Hermite normal form again.
If the rational rank increases, it can do so at most three
times.  At the first nonincrease the rational span is invariant
under every generator and its inverse, hence never increases
again.  If the current lattice has saturation index $D$ in
that span, every subsequent strict enlargement divides this
index by an integer at least two.  There can therefore be at
most $\lfloor\log_2D\rfloor$ such enlargements.  This proves
termination with a computable bound. At every step the Hermite
basis represents the integer sum of the generated columns.

\item On $A$ one has $M_2=M^{-1}$ in the two-point case.  The product
relation eliminates $s_2=t^{-1}x^{-\kappa}$, where $t=s_1$.
Multiplying this word $m_2$ times gives
\[
        (t^{-1}x^{-\kappa})^{m_2}
            =t^{-m_2}x^{-S_{m_2}(M)\kappa}.
\]
Thus the two power relations become $t^{m_1}=w_1$ and
$t^{m_2}=w'_2$.  Both right sides are fixed by $M$ on $A$.
Also $M^{m_1}=M^{m_2}=1$, hence $M^d=1$ there.
Bezout elimination gives exactly $e=0$ and $t^d=z$ in
\eqref{maps:eq:F.G5}.  Replacing the Bezout pair changes $z$ by a multiple
of $e$ and therefore changes no extension datum.

\item To prove the asserted injection, construct
$\overline A\rtimes_M\mathbb Z$ and quotient by its central
element $t^dz^{-1}$.  This cyclic subgroup has nonzero
projection $d\mathbb Z$ to the infinite cyclic factor, so it
intersects $\overline A$ trivially.  The quotient has precisely
the presentation obtained by elimination.  This proves \eqref{maps:eq:F.G6}.

\item \emph{The three-point kernel and the finite transfer.}
We compute the fibre injection by the homotopy sequence of a
fibration over the orbifold classifying space.  Let $\mathcal B$ be the orbifold classifying
space obtained from the three-holed sphere by attaching, at
its three boundary circles, the mapping cylinders of
$S^1\to BC_{m_i}$.  Over the punctured sphere take the bundle
with fibre $BA=K(A,1)$ and the action.  For each
boundary, the cyclic extension
\[
 E_i=\langle A,s_i\mid s_i a s_i^{-1}=M_i(a),\ s_i^{m_i}=w_i\rangle
\]
contains $A$ injectively: the same semidirect-product construction
as above applies because $M_i^{m_i}=1$ and $M_iw_i=w_i$.
The associated fibration $BA\to BE_i\to BC_{m_i}$ restricts
on $S^1$ to the required mapping-torus fibration.  Replace the
maps by their mapping-path fibrations before homotopy gluing,
so that the fibre is preserved on every attaching cylinder.
Glue these three fibrations, using the relative clutching
$\kappa$ on the punctured sphere.  Van Kampen gives \eqref{maps:eq:F.G3}
for the resulting total space.  Thus the homotopy sequence
of this fibration over $\mathcal B$ calculates precisely
the image of $A$ in the group under discussion.

\item When $\epsilon\le0$, the universal orbifold cover is the plane
or hyperbolic plane.  Its contractibility identifies
$\mathcal B$ with $B\Gamma$: indeed its Borel construction
has the just-described boundary-cylinder model and contractible
universal fibre.  The Euclidean and hyperbolic alternatives
are the oriented two-orbifold classification in
\cite[Theorems 2.3--2.4, p.~425; Euler-characteristic comparison, pp.~427--428]{rootScott83}.
Since $\pi_2(B\Gamma)=0$, the homotopy sequence proves \eqref{maps:eq:F.G8}.

\item When $\epsilon>0$, the universal orbifold cover is $S^2$.
Pulling the constructed $BA$ fibration to this finite cover
trivializes its local system.  Its transgression
$\pi_2(S^2)\to A$ is its integral clutching vector.
There is one lifted ordinary base cell for each $g\in\Gamma$,
contributing $g\kappa$, and one lifted cone cap for each left
coset of $\langle q_i\rangle$, contributing $gw_i$.
To check the sign, the base relation has fibre term $+\kappa$
and each cone relation has term $-w_i$; the oriented sphere
chain subtracts the cone cells from the base cell.  The result
is therefore the plus sign in both sums of \eqref{maps:eq:F.G9}.
Changing a coset representative by $q_i^j$ does not change
its contribution because $w_i$ is fixed.  The result is
$\Gamma$ invariant. The long
homotopy sequence now says that the only additional kernel
is exactly $\langle e\rangle$.  This proves the spherical
assertion, with the orientation fixed by this cell calculation.

\item After sorting its orders, the spherical possibilities are
$(2,2,n),(2,3,3),(2,3,4),(2,3,5)$, with groups of orders
$2n,12,24,60$, respectively.  These follow directly from
the positive spherical-triangle angle condition.
Construct the dihedral, tetrahedral, octahedral, or icosahedral
permutation group and find the marked peripheral generators in that
finite group. Breadth-first words supply their actions on $A$;
the generator-edge equalities are evaluated modulo the lattice $L$.
Compute the two finite sums in \eqref{maps:eq:F.G9} and their integer Smith form.

\item An order-one site is not a genuine cone.  Its matrix is the
identity and its relation is $s_i=x^{w_i}$.  Eliminate it
from \eqref{maps:eq:F.G3}; moving this vector to the right of the remaining
peripheral product replaces the clutching by
\begin{equation}
       \kappa+(M_{i+1}\cdots M_k)^{-1}w_i.                 \hypertarget{tag-02FQ}{}\label{maps:eq:F.G11}
\end{equation}
This reduces the three-point case with such a site to the
two-point computation.  A primitive rank-three cusp cannot
remain with only one finite linear monodromy: a finite-order
unipotent matrix is the identity.  

\item \emph{Freeness and abelianization.}
Freeness of \eqref{maps:eq:F.G1} is itself decidable exactly.  A nontrivial
stabilizer in a cyclic group contains a subgroup of prime
order, so it is enough to check $\tau_i^{m_i/p}$ for every
prime $p\mid m_i$.  For an affine map $Rx+b$ of a torus,
a fixed point exists precisely when every integral covector
in $\ker(R^t-I)$ takes an integral value on $b$.
Smith form proves this criterion: the image of $R-I$ on
the torus is the connected subtorus annihilated by those
primitive covectors. Compute the powers in the rational affine
matrix group and apply these finite tests.

\item Abelianizing \eqref{maps:eq:F.G3} yields \eqref{maps:eq:F.G10}.  The inclusion
\[
 \operatorname{im}(M_0-I)
       \subset\sum_i\operatorname{im}(M_i-I)
\]
follows by telescoping the product on the right; inversion
preserves the image of its difference from the identity.
The Smith form therefore computes every free and torsion
summand of $H_1$ exactly.

The bounds in this construction depend on the input.

\end{enumerate}
The corresponding exact calculations are implemented in \hyperlink{prog:V:fundamental-group}{the programs in Table~\ref*{tab:V:programs}}.
\end{proof}




\Needspace{5\baselineskip}
\begin{theorem}[A complete three-point logarithmic family]\hypertarget{tag-02NN}{}\label{maps:src:F.2.32}
\StatementRecord{maps:src:F.2.32}{proved}{02NN}{\cref{maps:src:F.3.7}, \cref{maps:src:F.2.28}, \eqref{maps:eq:F.LT4}, \eqref{maps:eq:F.LT5}, \cref{pub:V:smooth-gluing}, \eqref{maps:eq:F.HC8}, \eqref{maps:eq:F.HC9}, \eqref{maps:eq:F.LT8}, \cref{maps:src:F.2.29}, \eqref{maps:eq:F.LT6}, \eqref{maps:eq:F.LT9}}


Start with the compact complex fourfold of \ref{maps:src:F.3.7} and its
normalized order-five and order-four affine fillings.
Let $\ell$ be the primitive invariant covector on fibre
homology, with the sign chosen in \ref{maps:src:F.2.28}.
At any additional ordinary fibre choose
$v\in\mathbb Z^6$ and an integer $m\ge2$ such that
\begin{equation}
 r=\ell(v),\qquad \gcd(r,m)=1.
\hypertarget{tag-02H3}{}\label{maps:eq:F.LT1}
\end{equation}
Perform the marked logarithmic transformation with
translation $v/m$ described below.
This gives a smooth compact complex fourfold $X_{m,v}$
with Euler number six and three genuine finite affine
fillings, whose linear orders are five, four and one
and whose affine orders are five, four and $m$.
Its fundamental group is
\begin{equation}
 \pi_1(X_{m,v})=
 \langle a,b,c,h\mid h\text{ central},\
   a^5=h,\ b^4=h^{-1},\ c^m=h^r,\ abc=1\rangle .
\hypertarget{tag-02H4}{}\label{maps:eq:F.LT2}
\end{equation}
This is the group of the oriented Seifert manifold
with invariants
$(0;(5,-1),(4,1),(m,-r))$, in the convention
$x_i^{\alpha_i}h^{\beta_i}=1$ and product $x_i=1$.
Its base orbifold is hyperbolic, and its central fibre
$h$ has infinite order.

Put $C_{m,r}=\mathbb Z/(20r-m)\mathbb Z$, with
$\mathbb Z/0\mathbb Z=\mathbb Z$, and let
$\varepsilon=1$ when $20r=m$ and zero otherwise.
All the integral homology groups are
\begin{equation}
\begin{array}{c|c}
j&H_j(X_{m,v};\mathbb Z)\\ \hline
0,8&\mathbb Z\\
1,6&C_{m,r}\\
2&\mathbb Z^\varepsilon\oplus(\mathbb Z/m)^5\\
3&(\mathbb Z/m)^{10}\\
4&\mathbb Z^4\oplus(\mathbb Z/m)^{10}\\
5&(\mathbb Z/m)^5\\
7&\mathbb Z^\varepsilon .
\end{array}
\hypertarget{tag-02H5}{}\label{maps:eq:F.LT3}
\end{equation}
Under \eqref{maps:eq:F.LT1} the exceptional free-rank case is exactly
$(m,r)=(20,1)$.
In particular $(m,r)=(19,1)$ and $(21,1)$ have
$H_1=0$ but nontrivial, infinite fundamental groups.
More generally this occurs whenever $|20r-m|=1$.
The result holds for every vertical component of $v$ in $\ker\ell$.

\end{theorem}


\begin{proof}
\begin{enumerate}[label=\arabic*.]
\item We first specify the complex construction and its marking.
Over a small disc at the ordinary value use a holomorphic
period marking
$T_t=\mathbb C^3/\Lambda_t$.
Let $\omega_v(t)$ be the holomorphic period vector
corresponding to $v$.
On the root disc $t=s^m$ quotient the pulled-back family by
\begin{equation}
 (s,Z)\longmapsto
   (\zeta_m s,\ Z+\omega_v(s^m)/m).
\hypertarget{tag-02H6}{}\label{maps:eq:F.LT4}
\end{equation}
Condition \eqref{maps:eq:F.LT1} makes $v$ have exact order $m$
in $\mathbb Z^6/m\mathbb Z^6$.
Every nontrivial power therefore acts freely, including
over $s=0$.
The quotient is smooth, has a proper map to the disc,
and has central fibre of multiplicity $m$.
On the punctured disc its holomorphic identification with
the original family is
\begin{equation}
 (s,Z)\longmapsto
 \left(s^m,\ Z-\frac{\log s}{2\pi i}\omega_v(s^m)\right).
\hypertarget{tag-02H7}{}\label{maps:eq:F.LT5}
\end{equation}
Changing the branch of $\log s$ changes the second
coordinate by a lattice period.
Changing $s$ by $\zeta_m s$ changes its logarithm by
$2\pi i/m$ and cancels the translation in \eqref{maps:eq:F.LT4}.
Thus \eqref{maps:eq:F.LT5} is a well-defined holomorphic isomorphism
of the quotient with the original torus family off zero.
Gluing on a fixed annulus supplies the asserted compact
complex manifold: its charts are the original charts
and the free finite quotient charts \eqref{maps:eq:F.LT4}, and it is
the union of compact pieces over closed base regions.
The gluing is an isomorphism on collars, so the total
space is Hausdorff and smooth.
The Euler number is unchanged, since the removed
ordinary torus cap, the new cap, and their boundary
all have Euler number zero.
This marking fixes the new relation $c^m=v$ and
zero extra cusp clutching.

\item In the original datum the invariant normal closure
of the cusp angular lattice is exactly $\ker\ell$,
as proved in \cref{pub:V:smooth-gluing}.
Adding the identity linear generator does not change
that lattice.
The fibre image in van Kampen therefore reduces to
the central generator $h$, while the original two
finite relations remain $a^5=h$, $b^4=h^{-1}$.
The new relation is $c^m=h^{\ell(v)}$ and the base
relation is $abc=1$.
This proves \eqref{maps:eq:F.LT2}; the presentation depends on $v$ only through $r=\ell(v)$.

\item One may verify the Seifert identification directly:
take a product circle bundle on a three-holed sphere
and attach its three solid tori with primitive meridian
slopes $(5,-1),(4,1),(m,-r)$.
They are primitive by \eqref{maps:eq:F.LT1}, and van Kampen gives
exactly \eqref{maps:eq:F.LT2}.
Its orbifold base has
$1/5+1/4+1/m<1$ for every $m\ge2$.
Pulling back the circle fibration to its contractible
hyperbolic orbifold cover gives a circle bundle on
the plane, hence the trivial circle bundle.
The subsequent universal cover is the plane times
$\mathbb R$, and the fibre subgroup is infinite cyclic.
This also proves directly that none of these groups
is trivial.
The rational Euler number in these conventions is
$r/m-1/20$.

\item In abelianization the relations $5a=h$, $4b=-h$
give $a=4x$, $b=-5x$, $h=20x$ because five and four
are coprime.
The product relation gives $c=x$.
The remaining relation is $(m-20r)x=0$.
Thus $H_1=C_{m,r}$ exactly, including its integral
generator and the exceptional infinite-cyclic case.

\item We now compute the higher integral groups, including
all extensions.
Write $\mathcal R^q$ for the higher direct-image
sheaf of the original fourfold and
$\mathcal R_m^q$ for that of $X_{m,v}$, after identifying
the bases and their ordinary local systems.
At the new point the reduced finite quotient is a
torus with lattice
$\mathbb Z^6+\mathbb Z(v/m)$.
Its covering by the original torus is cyclic of
degree $m$.
The degree-one pullback has Smith factors
$1,1,1,1,1,m$, so its degree-$q$ cokernel is
\begin{equation}
 C^q=(\mathbb Z/m)^{\binom5{q-1}},
 \qquad 1\le q\le6 .
\hypertarget{tag-02H8}{}\label{maps:eq:F.LT6}
\end{equation}
The integral cohomology of this torus cap is free.
The covering pullback at the new point gives the sheaf comparison
\begin{equation}
 0\longrightarrow\mathcal R_m^q
  \longrightarrow\mathcal R^q
  \longrightarrow C^q_{\{p\}}\longrightarrow0 .
\hypertarget{tag-02H9}{}\label{maps:eq:F.LT7}
\end{equation}
It is the identity away from the new point and the
covering image at that point.

\item For $2\le q\le5$, all original global sections of
$\mathcal R^q$ are supported at the original finite
point or the cusp, by \eqref{maps:eq:F.HC8}--\eqref{maps:eq:F.HC9}.
Their map to $C^q_{\{p\}}$ is zero.
Consequently their $H^0$ groups are unchanged,
their $H^2$ groups are unchanged, and
\begin{equation}
 0\longrightarrow C^q
  \longrightarrow H^1(\mathcal R_m^q)
  \longrightarrow H^1(\mathcal R^q)
  \longrightarrow0 .
\hypertarget{tag-02HA}{}\label{maps:eq:F.LT8}
\end{equation}
The original degree-one base cohomology of these
sheaves is zero for $q=2,4,5$ and is $\mathbb Z^3$
for $q=3$.
The assertion for $q=5$ also follows immediately
from the original $H^6=0$: its bidegree $(1,5)$
has no incoming or outgoing differential.
Thus \eqref{maps:eq:F.LT8} gives $C^q$ in degrees two, four and
five, and $\mathbb Z^3\oplus C^3$ in degree three.

\item The original four order-two torsion transgressions
remain isomorphisms.
Indeed their boundary suspension representatives and
their images in global coinvariants are the ones in
\ref{maps:src:F.2.29}; adding the identity monodromy changes none
of those coinvariant quotients.
These point-supported boundary maps do not depend
on the free degree-one transgression, which changes
under the new logarithmic gluing.
The free degree-four cusp column likewise remains
zero, since its boundary torsion representative has
zero image in that same unchanged global quotient.
It follows that the full filtrations in degrees
three and five have just one nonzero term:
\begin{equation}
 H^3=C^2=(\mathbb Z/m)^5,\qquad
 H^5=C^4=(\mathbb Z/m)^{10}.
\hypertarget{tag-02HB}{}\label{maps:eq:F.LT9}
\end{equation}
In degree four there is an exact sequence
\[
 0\longrightarrow
       \mathbb Z^3\oplus(\mathbb Z/m)^{10}
  \longrightarrow H^4
  \longrightarrow\mathbb Z\longrightarrow0.
\]
It splits as an abelian-group extension because
its quotient is free, giving the stated middle group.

\item For the remaining free ranks,
$H^1(X_{m,v})=\operatorname{Hom}(C_{m,r},\mathbb Z)$
has rank $\varepsilon$.
The degree-one global section group has rank one,
and the only rational differential involving it is
$d_2:H^0(\mathcal R_m^1)\to H^2(\mathcal R_m^0)$.
Its kernel has rank $\varepsilon$, hence its
cokernel also has rank $\varepsilon$.
The degree-one parabolic group is rationally zero,
and there are no rational degree-two global sections.
Therefore $b_2=\varepsilon$.
The finite sheaf cokernel in this degree-one row is
$\mathbb Z/\gcd(20,m)$: the original invariant
section restricts as $20\ell$, and its value at
the new covering cokernel is $20r$ modulo $m$.
This finite group changes no rational rank.

\item Poincar\'e duality gives
$H_3=H^5$, $H_4=H^4$, and $H_5=H^3$.
The universal coefficient theorem now identifies
the torsion of $H_2$ with $H^3$ and its free rank
with $\varepsilon$.
It also gives $H^2=C_{m,r}$: when $20r\ne m$ this
is $\operatorname{Ext}(H_1,\mathbb Z)$; when
$20r=m$ it is the free rank-one group.
Duality gives $H_6=C_{m,r}$ and
$H_7=\mathbb Z^\varepsilon$.
This proves the complete table \eqref{maps:eq:F.LT3}.



\end{enumerate}
Exact calculations for the displayed finite examples appear in \hyperlink{prog:V:three-points}{the programs in Table~\ref*{tab:V:programs}}.
\end{proof}




\begin{theorem}[Arbitrary finite affine caps]\hypertarget{tag-02NO}{}\label{maps:src:F.2.33}
\StatementRecord{maps:src:F.2.33}{proved}{02NO}{\cref{maps:src:F.2.31}}


Let $g(x)=Mx+v/m$ be a free cyclic affine action of affine order
$m\geq2$ on $T^r$, with $M\in\mathrm{GL}(r,\mathbb Z)$ and
$v\in\mathbb Z^r$.  For every specified degree bound $d$, there is
an explicit terminating integer algorithm producing chains and
boundary maps through degree $d+1$ for the pair
\begin{equation}
 B=(D^2\times T^r)/\langle(s,x)\mapsto
 (e^{2\pi i/m}s,g(x))\rangle,
 \qquad Y=\partial B.                                      \hypertarget{tag-02FE}{}\label{maps:eq:F.CA1}
\end{equation}
The cap complex has at most $2^r$ generators in each degree and
the boundary complex at most $2^{r+1}$.  These are complexes of
finite rank in each degree; only their indicated finite prefixes
are required to compute $H_j(B)$, $H_j(Y)$ and $H_j(B,Y)$ for
$j\leq d$.  Their maps include the affine translations and all integral higher coherence terms.

\end{theorem}
\begin{proof}
\begin{enumerate}[label=\arabic*.]
\item We give the complete recursive algorithm.  All chains below are
finite integer dictionaries, so the displayed operations are
effective over $\mathbb Z$.

\item First consider the normalized bar resolution of $\mathbb Z^r$
over $R=\mathbb Z[t_1^{\pm1},\ldots,t_r^{\pm1}]$ and its Koszul
resolution $K=R\otimes\bigwedge^*\mathbb Z^r$.
The latter differential is contraction with $(t_i-1)_i$.
Write $t^a=t_1^{a_1}\cdots t_r^{a_r}$ and put
$S_n(t)=(t^n-1)/(t-1)$, interpreted as its finite integral
Laurent sum for positive, zero or negative $n$.
The following $\mathbb Z$-linear contraction of the augmented
Koszul complex fixes a completely explicit choice:
\begin{equation}
 s_K(t^a e_I)=
 \sum_{j<\min I}
 t_{j+1}^{a_{j+1}}\cdots t_r^{a_r}S_{a_j}(t_j)
 e_j\wedge e_I,                                           \hypertarget{tag-02FF}{}\label{maps:eq:F.CA2}
\end{equation}
where $\min\varnothing=r+1$ and indices in $I$ are increasing.
The bar contraction is
$s_B(t^a[g_1|\cdots|g_q])=[a|g_1|\cdots|g_q]$; terms with a
zero bar entry are zero.  Set $p_0=1$ and define the
$R$-linear projection and homotopy on coefficient-one bar
generators recursively by
\begin{equation}
 p_q=s_Kp_{q-1}d_B,
 \qquad
 h_q=s_B(1-ip_q-h_{q-1}d_B).                               \hypertarget{tag-02FG}{}\label{maps:eq:F.CA3}
\end{equation}
Here $i(e_I)$ is the alternating sum of the bars of the ordered
unit vectors in $I$.  The recursion decreases degree.  The
identities for the two contractions give, by induction,
$pi=1$, $dh+hd=1-ip$, $ph=hi=h^2=0$.
For example $s_K^2=s_B^2=0$, $ps_B=s_Kp$, and the alternating
unit-vector formula gives $pi=1$; substituting these into
\eqref{maps:eq:F.CA3} proves the three side identities in the same induction.
Augmenting $R\to\mathbb Z$ gives a strong chain contraction
of the ordinary torus bar complex to $\bigwedge^*\mathbb Z^r$.

\item Translations must enter before this contraction.  Let
$b_j=\sum_{a=0}^{j-1}M^av/m$ and let $p_j$ be its coordinatewise
fractional part.  The $m$ points $p_j$ are distinct because the
action is free.  Use the full fundamental groupoid of $T^r$
on these points: an arrow from object $a$ to object $b$ is an
integer vector $\lambda$, represented by the affine path from
$p_a$ to $p_b+\lambda$.  Composition adds the integer vectors.
The generator acts strictly by shifting objects and, for its
$s$th power, by
\begin{equation}
 \lambda\longmapsto M^s\lambda+z_s(b)-z_s(a),
 \qquad z_s(a)=M^sp_a+b_s-p_{a+s}\in\mathbb Z^r.             \hypertarget{tag-02FH}{}\label{maps:eq:F.CA4}
\end{equation}
The subscripts on objects are taken modulo $m$.
The nerve of this groupoid is equivariantly homotopy equivalent
to the torus: the simplex with successive arrow labels
$\lambda_i$ maps affinely to the torus with lifted vertices
$p_{a_j}+\sum_{i\leq j}\lambda_i$.
Forgetting the objects contracts its chains to the ordinary
torus bar complex.  An explicit homotopy is the negative
alternating prism which, at each possible position, inserts
the zero-labelled arrow to object zero and replaces all
subsequent objects by zero.  Combining this with \eqref{maps:eq:F.CA3}
gives a contraction $(i,p,h)$ of the groupoid chains to
$\bigwedge^*\mathbb Z^r$ with the same side identities.

\item Use the normalized cyclic bar construction on this strict
groupoid action.  Its differential is $d_v+\delta$, where
$\delta$ is the ordinary group-bar boundary, multiplied by
$(-1)^q$ on vertical degree $q$; its first face applies \eqref{maps:eq:F.CA4}.
The vertical contraction transfers this to cyclic bar chains
with exterior coefficients by the finite formulas
\begin{equation}
 I=(1+h\delta)^{-1}i,\qquad
 P=p(1+\delta h)^{-1},\qquad D=p\delta I.                  \hypertarget{tag-02FI}{}\label{maps:eq:F.CA5}
\end{equation}
Inverses mean their alternating geometric series.  On a chain
of cyclic bar degree $p$ each series has at most $p+1$ terms,
since $\delta$ lowers that degree.  Multiplication and the
contraction identities directly give $(d_v+\delta)I=ID$ and
$P(d_v+\delta)=DP$,
and a chain homotopy inverse.  Equivalently these are the
finite homological perturbation identities, here derived by
expanding the finite series.  The first term of $D$ is the
cyclic bar differential for the exterior action of $M$.
Every remaining term lowers cyclic bar degree by at least two.

\item The second contraction used in the algorithm is as follows.  The cyclic periodic resolution has one
$\mathbb Z[C_m]$ generator in each degree, with differential
$g-1$ in odd degree and $1+g+\cdots+g^{m-1}$ in even degree.
On $g^a$, $0\leq a<m$, its augmented contraction is
$\sum_{j=0}^{a-1}g^j$ in even degree and is $1$ precisely when
$a=m-1$ in odd degree, zero otherwise.
Apply \eqref{maps:eq:F.CA3} with this contraction to compare it with the
cyclic bar resolution.  Tensor its coefficients with the
exterior action of $M$.  Transferring the higher terms of $D$
once more by \eqref{maps:eq:F.CA5} terminates: each higher term followed by
the horizontal homotopy decreases cyclic degree by at least
one.  The resulting basis in total degree $n$ is
$\{(p,I):p+|I|=n,\ I\subseteq\{1,\ldots,r\}\}$.
This proves the claimed $2^r$ bound and termination in every
specified degree.

\item \par\medskip\noindent\textbf{Explicit comparison with the periodic cyclic resolution.}\quad 
\hypertarget{tag-02J1}{}\label{maps:mapsF:repair:cyclic-inclusion}
Write $C_m=\langle g\rangle$, $R_m=\mathbb Z[C_m]$, and let
$e_n$ generate the periodic resolution in degree $n$.
Let $B_*$ be the normalized $R_m$-bar resolution, with the
insertion contraction $s_B$.  Besides the periodic contraction
$s_P$ just specified, fix the inclusion recursively by
\[
 i_0(e_0)=[],\qquad
 i_n(e_n)=s_B\bigl(i_{n-1}(d_Pe_n)\bigr),
 \quad n>0,
\]
and extend it $R_m$-linearly.  This is a finite formula.  More
explicitly, with every summation index between $1$ and $m-1$,
\[
\begin{split}
 i_{2j}(e_{2j})
   &=\sum_{a_1,\ldots,a_j}
       [g^{a_1}|g|\cdots|g^{a_j}|g],\\
 i_{2j+1}(e_{2j+1})
   &=\sum_{a_1,\ldots,a_j}
       [g|g^{a_1}|g|\cdots|g^{a_j}|g].
\end{split}
\]
The empty sum of indices gives $i_0(e_0)=[]$ and
$i_1(e_1)=[g]$.  The contraction identity for $s_B$ proves
$d_Bi=id_P$ by induction: the chain to which $s_B$ is applied
is a cycle, since the preceding comparison commutes with
differentials and $d_P^2=0$.

\item Define $p$ and $h$ on coefficient-one bar generators by the
recursion already used for the torus comparison, substituting
$s_P$ for $s_K$, and extend them $R_m$-linearly.  Their values
are given by the following formulas.
For $1\leq a_i<m$, put
$\epsilon(a,b)=\lfloor(a+b)/m\rfloor$.  The projection is
\[
\begin{split}
 p_{2j}[g^{a_1}|\cdots|g^{a_{2j}}]
   &=\prod_{\nu=1}^{j}\epsilon(a_{2\nu-1},a_{2\nu})\,e_{2j},\\
 p_{2j+1}[g^{a_1}|\cdots|g^{a_{2j+1}}]
   &=\left(\sum_{b=0}^{a_1-1}g^b\right)
       \prod_{\nu=1}^{j}\epsilon(a_{2\nu},a_{2\nu+1})\,e_{2j+1}.
\end{split}
\]
These expressions follow by induction from
$p_n=s_Pp_{n-1}d_B$: the terms other than the first bar face
are annihilated after the preceding contraction identity is
used.  In $pi(e_{2j})$ only the summand with every
$a_\nu=m-1$ survives.  The same observation, with initial
factor $1$, proves $pi(e_{2j+1})=e_{2j+1}$.
Thus $pi=1$ in every degree.

\item The recursion
$h=s_B(1-ip-hd_B)$ gives $d_Bh+hd_B=1-ip$ by induction.
Moreover $p s_B=s_Pp$ follows from the displayed projection
formulas, while $s_P^2=s_B^2=0$ and $s_Bi=0$ follow directly
from their definitions.  Substitution in the same recursion
gives $ph=hi=h^2=0$.  We therefore have the strong comparison
required for the second finite perturbation series, including
its inclusion, projection, homotopy and side identities.
After tensoring with the exterior coefficient module, the
same identities hold.  Every higher perturbation followed by the horizontal homotopy lowers cyclic degree, proving termination with the $2^r$ generator bound.

\item The free action makes the Borel construction equivalent to
$T^r/C_m$, and radial contraction identifies this with $B$.
For $Y$ use exactly the same algorithm on $T^{r+1}$ with
matrix $\operatorname{diag}(1,M)$ and numerator $(1,v)$.
The projection forgetting its first circle is equivariant and
represents $Y\to B$ under the radial contraction.
More generally an equivariant integral matrix $L$ acts on
groupoid arrows as
$\lambda\mapsto L\lambda+k_b-k_a$, where
$k_a=Lp_a-p'_a\in\mathbb Z^{r'}$.
Transferring this strict bar map by $P'LI$ gives the required
integer chain map.  Its mapping cone computes the pair.
Every series is finite with the proved degree bound, and
integer kernels and Smith normal forms finish the calculation.
\end{enumerate}
The corresponding exact calculations are implemented in \hyperlink{prog:V:chains}{the programs in Table~\ref*{tab:V:programs}}.
\end{proof}











\begin{theorem}[The marked Kuhn cusp pair]\hypertarget{tag-02NP}{}\label{maps:src:F.2.34}
\StatementRecord{maps:src:F.2.34}{proved}{02NP}{\cref{maps:src:F.2.33}; \cite[Theorem 4.1.1 and Lemmas 4.1.2--4.1.3, pp.~480--482]{ffAchingerOgus20}; \cite[Theorem 8.7, pp.~57--59, and Corollary 8.8, p.~59, arXiv v1]{Yang2026}}


The fixed primitive rank-three Kuhn cusp admits an explicit
finite integer chain complex for its cap $W$, its boundary
$Y$, and the inclusion $Y\to W$.  It is obtained from a
category with 26 objects, 170 nonidentity morphisms, and
respectively $26,170,288,144$ strict flags of lengths
$0,1,2,3$.  The construction records the full integral
collapse map.  In particular it gives
\begin{equation}
\begin{array}{c|rrrrrrrrr}
j&0&1&2&3&4&5&6&7&8\\\hline
\operatorname{rank}H_j(W)&1&3&9&9&10&3&1&0&0\\
\operatorname{rank}H_j(Y)&1&4&12&18&18&12&4&1&0\\
\operatorname{rank}H_j(W,Y)&0&0&1&3&10&9&9&3&1
\end{array}                                                \hypertarget{tag-02FJ}{}\label{maps:eq:F.CK1}
\end{equation}
The only torsion in these three rows is
$\operatorname{Tor}H_3(Y)=\mathbb Z/2$.

\end{theorem}
\begin{proof}
\begin{enumerate}[label=\arabic*.]
\item Let $\mathcal K$ be the periodic Kuhn triangulation of
$\mathbb R^3$.  Normalize every simplex by translating its
lexicographically first vertex to zero.  The objects are its
simplex orbits; a morphism $\tau\to\sigma$ is every strict
translated face inclusion $\tau+w\subset\sigma$.
Keep distinct translated inclusions as distinct morphisms.
Composition adds the translation vectors.  Enumerating all
nonempty faces of the six chains in the unit cube gives
$(1,7,12,6)$ simplex orbits in dimensions zero through three,
and the displayed morphism and flag counts.  The category
nerve is the barycentric cell model of the dual decomposition
of $\mathbb R^3/\mathbb Z^3$; the direction of the face order
is reversed when interpreted as dual cells.

\item For an object $\tau$ let $L_\tau\subset\mathbb Z^3$ be the
primitive span of its vertex differences.  Its primitivity
follows from unimodularity of every Kuhn simplex.  Choose by
Smith reduction a quotient map
$Q_\tau:\mathbb Z^3\to\mathbb Z^{3-\dim\tau}$ with kernel
$L_\tau$, and an integral section $R_\tau$.
Define two strict diagrams of tori on this category:
\begin{equation}
\begin{array}{c|c|c}
 &\text{object at }\tau&\text{map for }\tau+w\subset\sigma\\\hline
\mathcal W&T^{3-\dim\tau}&Q_\sigma R_\tau\\
\mathcal Y&T^4&S_w=\begin{pmatrix}1&0\\ w&I_3\end{pmatrix}.
\end{array}                                                \hypertarget{tag-02FK}{}\label{maps:eq:F.CK2}
\end{equation}
The first maps compose because $Q_\sigma$ vanishes on
$L_\tau$, and the second compose because translations add.
The maps $[\,0\ Q_\tau\,]$ form a natural transformation
$\mathcal Y\to\mathcal W$: the only additional term is
$Q_\sigma w=0$, since $w$ is a vertex of $\sigma$.

\item Here these diagrams describe the local topology.
The toric moment quotient on the dual cell to $\tau$ has
angular torus $T^3/(L_\tau\otimes\mathbb R/L_\tau)$;
passing to a face collapses exactly the next angular
subtorus.  Barycentrically subdividing the dual cells gives
the homotopy colimit of $\mathcal W$ for the central fibre.
For its real oriented logarithmic boundary, retain the
argument $\theta_0$ of the base parameter and the three
angular coordinates $\theta$ of the toric chart.  Translating
the chart by $w$ changes them to
$(\theta_0,\theta+w\theta_0)$, which is precisely $S_w$.
Thus the logarithmic boundary is the homotopy colimit of
$\mathcal Y$, and forgetting the logarithmic phases on each
stratum is exactly the natural transformation just given.
This can also be read directly in the unimodular charts
$x_0\cdots x_k=t$: the moment coordinates form the dual
simplex, the four logarithmic phases remain, and the ordinary
stratum retains their quotient by the vertex differences.
The finite face identifications give \eqref{maps:eq:F.CK2}.

\item The Kuhn map with its divisor log structures is proper,
log smooth and vertical after restricting to a small disc.
The logarithmic comparison theorem \cite[Theorem 4.1.1 and Lemmas 4.1.2--4.1.3, pp.~480--482]{ffAchingerOgus20} identifies its topological germ with the mapping cylinder of this logarithmic collapse.  Consequently the preceding
homotopy-colimit diagrams and natural transformation model
the pair $(W,Y)$ itself.  Fix the marking by letting $b_i$
be the category loop with translation label $e_i$.  A
homotopy-colimit edge with map $S_{e_i}$ has the relation
$b_i^{-1}\theta_0b_i=\theta_0+a_i$; equivalently
$\theta_0b_i\theta_0^{-1}=b_i+a_i$.  This specifies the
required primitive cusp
$U=\begin{pmatrix}I_3&I_3\\0&I_3\end{pmatrix}$.

\item \par\medskip\noindent\textbf{Hypotheses and marking of the logarithmic comparison.}\quad 
\hypertarget{tag-02J2}{}\label{maps:mapsF:repair:log-cylinder-scope}
Give the small disc its standard divisorial log structure and
the Kuhn total space the divisorial structure of the reduced
central fibre.  In a unimodular chart the characteristic
monoid map is
\[
 \mathbb N\longrightarrow\mathbb N^{k+1},\qquad
 1\longmapsto(1,\ldots,1).
\]
Both monoids are fine and saturated.  The group cokernel is
free, and the underlying chart is the monomial semistable
map $x_0\cdots x_k=t$ with additional smooth coordinates;
this is log smooth.  The quotient monoid is a group, since
the negative of each coordinate generator equals the sum of
the other generators after the diagonal generator is killed.
The morphism is therefore vertical.  The periodic toric
construction is proper over the sufficiently small disc.

\item These are precisely the hypotheses of
\cite[Theorem~4.1.1, pp.~480--482]{ffAchingerOgus20}.
Its proof also notes that exactness is automatic over the
standard log disc.  It identifies the germ with the mapping
cylinder of the logarithmic collapse.  This homeomorphism
depends on a trivialization; Lemma~4.1.3 fixes that
trivialization on the logarithmic central fibre.  We keep
the resulting radial boundary identification and its stated
angular and tropical generators.  This suffices to model
the pair and its inclusion up to marked homotopy.  

\item We specify finite integer chains.  Replace every
torus by its ordinary normalized bar complex.  The affine
simplex map from the bar construction to the torus is natural
for all the integer homomorphisms in \eqref{maps:eq:F.CK2}, so this replacement
does not change either homotopy colimit or its collapse map.
Take the simplicial replacement over all strict flags of the
finite category.  Its vertical differential is the torus bar
boundary; its horizontal differential composes adjacent
arrows, with the first face applying the diagram map, and
has the sign $(-1)^q$ in vertical degree $q$.
Apply the first contraction and the finite series \eqref{maps:eq:F.CA5}
object by object.  There are no flags longer than three,
so every series has at most four terms.  The resulting
finite basis consists of a strict flag beginning at $\tau$
and an exterior basis vector of its torus.  Transfer the
natural transformation by $P_{\mathcal W}FI_{\mathcal Y}$.
This proves both the chain-map identity and completeness of
the collapse map, including all integral coherence corrections.

\item The raw cap chain ranks in degrees zero through six are
\[
 (26,199,688,1231,1226,648,144),
\]
and the raw boundary ranks in degrees zero through seven are
\[
 (26,274,1124,2420,3010,2186,864,144).
\]
Enumerate the objects, translated incidences and flags. Composition
and the alternating boundary formulas give $d^2=0$ and the
collapse-chain identity on their generators.  Integral cancellation of a
unit pivot is a chain homotopy equivalence: clear its row by
integer column operations, then remove the resulting
contractible two-generator summand.  The adjacent differential
updates are forced by $d^2=0$.  After cancelling unit pairs for the cap, boundary and collapse cone, Smith reduction gives \eqref{maps:eq:F.CK1}. The reduced cap and
relative differentials vanish, and the boundary has one
nonzero Smith factor, equal to two in degree four to three.
This computes the integral groups.
\end{enumerate}
The corresponding exact calculations are implemented in \hyperlink{prog:V:chains}{the programs in Table~\ref*{tab:V:programs}}.
\end{proof}


\begin{remark}[Cusp quotient and global chains]\hypertarget{tag-PV19}{}\label{rem:V:cusp-input-and-global-chains}
\StatementRecord{rem:V:cusp-input-and-global-chains}{proved}{PV19}{}
The finite-pair and cusp-pair chain models enter the global
gluing cone of \cref{maps:src:F.2.35}. Its maps to the common
torus-bundle complement include the peripheral word and clutching
translation.

The periodic toric quotient used for the cusp is the construction in
Yang's theorem \emph{Higher-dimensional periodic quotient}, and its
Kuhn-grid Euler number is the formula in his corollary
\emph{The standard grid family}
\cite[Theorem 8.7, pp.~57--59, and Corollary 8.8, p.~59, arXiv v1]{Yang2026}.
For the present primitive rank-three grid, this gives Euler number
$3!=6$.  The full cap and boundary complexes, their integral inclusion,
and the torsion groups are calculated in the proof above.
\end{remark}







\begin{theorem}[Complete integral homology of a marked gluing]\hypertarget{tag-02NQ}{}\label{maps:src:F.2.35}
\StatementRecord{maps:src:F.2.35}{proved}{02NQ}{\cref{maps:src:F.2.31}, \cref{maps:src:F.2.33}, \cref{maps:src:F.2.34}}


For every admissible datum there is a terminating exact
integer algorithm computing every integral homology group of
the resulting eight-manifold, including every torsion invariant
factor.  The input consists of the finite homology matrices
$M_1,\ldots,M_k$, free affine numerators $v_i/m_i$ of their
affine orders, the integral clutch vector $\kappa$, and the
chosen primitive cusp marking
\begin{equation}
 M_0=(M_1\cdots M_k)^{-1},\qquad
 M_0P=P\begin{pmatrix}I_3&I_3\\0&I_3\end{pmatrix},
 \qquad P\in\mathrm{GL}(6,\mathbb Z).                       \hypertarget{tag-02FZ}{}\label{maps:eq:F.GC1}
\end{equation}
Here $k=2$ or $3$ and the genuine orders are at least two;
order-one sites are first eliminated as in \ref{maps:src:F.2.31}.
If no marking is given, the algorithm chooses one by Smith reduction.
The complete chain prefix through degree nine has at most
4003 generators in any degree.  

\end{theorem}
\begin{proof}
\begin{enumerate}[label=\arabic*.]
\item We construct both boundary maps for the global gluing.  Let $F_k$ be the free group with ordered
generators $s_1,\ldots,s_k$, and let $M(w)$ be the product
specified by a reduced word $w$.  The torus-bundle complement
over the punctured sphere retracts to the corresponding bundle
over a graph.  Its fundamental group and multiplication are
\begin{equation}
 \Gamma=\mathbb Z^6\rtimes F_k,
 \qquad (x,w)(y,u)=(x+M(w)y,wu).                            \hypertarget{tag-02G0}{}\label{maps:eq:F.GC2}
\end{equation}
The fibre and base are aspherical and the base has no
two-cells.  Hence a chosen section identifies the complement
up to homotopy with $B\Gamma$.  These section and peripheral
choices are fixed by the input clutch convention below.

\item \emph{A finite complement complex with its full bar projection.}
Use reduced words for $F_k$.  The tree resolution of the free
group has $K_0=\mathbb Z[F_k]$ and
$K_1=\bigoplus_i\mathbb Z[F_k]e_i$, with
$de_i=s_i-1$, and no higher terms.  Its explicit augmented
contraction sends a word to its signed path in the Cayley
tree: for each positive letter $s_i$ add its preceding prefix
times $e_i$, and for a negative letter $s_i^{-1}$ subtract
the prefix including that letter times $e_i$.
This is the Fox derivative, and its boundary is $w-1$.
Together with the bar insertion contraction it gives, by
the recursion \eqref{maps:eq:F.CA3}, a strong contraction of the free-group
bar resolution to this tree resolution.  Every recursion
decreases bar degree and every word is finite.

\item For a bar simplex $[(x_1,w_1)|\cdots|(x_n,w_n)]$ of $\Gamma$,
put $y_j=M(w_1\cdots w_{j-1})x_j$.
The exact Alexander--Whitney chain map to the torus/free-group
double bar complex is
\begin{equation}
 \sum_{p+q=n}(-1)^{pq}
 [w_1|\cdots|w_p]\otimes[y_{p+1}|\cdots|y_n].              \hypertarget{tag-02G1}{}\label{maps:eq:F.GC3}
\end{equation}
Terms with an identity bar entry are zero.  The horizontal
first face acts on torus vectors by
$\rho(w)=M(w)^{-1}$, and the horizontal boundary has sign
$(-1)^q$.  Notice that $\rho$ is a right action:
$\rho(wu)=\rho(u)\rho(w)$.  This order is essential when
the matrices do not commute.  Formula \eqref{maps:eq:F.GC3} follows from
the diagonal bisimplicial bar construction; alternatively,
its chain-map identity follows by pairing adjacent interior
faces, with the first face using
$M(w_1)^{-1}y_j=M(w_2\cdots w_{j-1})x_j$.
The usual shuffle inverse is finite in each degree, so the
map is a chain equivalence.

\item Apply the torus contraction and \eqref{maps:eq:F.CA5}, then the free-group
tree contraction and the same finite perturbation formula.
Higher terms decrease group bar degree at least two, whereas
the tree complex only has degrees zero and one.  Thus the
resulting complement complex is
\begin{equation}
 C_n(E)=\bigwedge^n\mathbb Z^6\ \oplus\
 \bigoplus_{i=1}^k\bigwedge^{n-1}\mathbb Z^6,
 \qquad
 d(e_i\otimes\xi)=(-1)^{n-1}
   (\bigwedge^{n-1}M_i^{-1}-1)\xi.                         \hypertarget{tag-02G2}{}\label{maps:eq:F.GC4}
\end{equation}
Its differential alone has this short form; its projection
from group bars retains the finite higher perturbation terms.
Those projection terms are used for every boundary map.

\item \emph{Finite affine boundary maps.}
For the $i$th finite cap use the orbit groupoid of \ref{maps:src:F.2.33}.
Its one-edge homotopy colimit for the affine generator is
the boundary mapping torus.  Vertically contracting it gives
a bounded complex with
$C_n(Y_i)=\bigwedge^n\mathbb Z^6\oplus
\bigwedge^{n-1}\mathbb Z^6$ and differential
$(-1)^{n-1}(\bigwedge^{n-1}M_i-1)$ on the second summand.
Its inclusion into the cyclic bar construction, followed by
the cap projection in \eqref{maps:eq:F.CA5}, is the full map $Y_i\to B_i$.


\item Here is an explicit group-bar map for its other inclusion.
Send a torus-groupoid arrow labelled $\lambda$ to the
translation $x^\lambda\in\Gamma$.  With the integral offsets
$z_1(a)$ in \eqref{maps:eq:F.CA4}, put
$\eta_a=s_i^{-1}x^{z_1(a)}$.
For an arrow $f:a\to b$, the exact identity
\begin{equation}
 \eta_a\,\psi(gf)=\psi(f)\,\eta_b                        \hypertarget{tag-02G3}{}\label{maps:eq:F.GC5}
\end{equation}
is the natural-transformation identity: it is just
$s_i^{-1}x^{M_i\lambda+z_1(b)}
 =x^\lambda s_i^{-1}x^{z_1(b)}$.
For a $q$-simplex, the edge map is $(-1)^q$ times the
alternating prism which inserts $\eta$ at every position,
uses the original arrows before the insertion and the
transformed arrows afterwards.  Its boundary is precisely
the signed difference of the two end maps, by \eqref{maps:eq:F.GC5}.
Compose this prism with the finite vertical lift of the
one-edge complex, then with \eqref{maps:eq:F.GC3} and the complement
projection.  This proves the full chain-map identity.
The offsets record the integer affine lift: replacing
$v_i$ by $v_i+m_i\lambda$ changes the offsets and the map,
even if the affine map on the torus has not changed.
An optional unimodular local basis simply applies its matrix
to every translation in \eqref{maps:eq:F.GC5}.

\item We make the distinction between a bar cell and the geometric
positive root explicit.  Put
$a_i=z_1(0)=\lfloor v_i/m_i\rfloor$, coordinatewise, and
let $\tau_i$ be the positive one-edge bar cell closed using
the zero-labelled return arrow to object zero.  With the
first-face convention above its geometric loop is
\begin{equation}
 \tau_i=g_i^{-1}x^{a_i},\qquad
 F_E(\tau_i)=-e_i+M_i^{-1}a_i\in C_1(E),\qquad
 g_i=x^{a_i}\tau_i^{-1}.                                  \hypertarget{tag-02G4}{}\label{maps:eq:F.GC5a}
\end{equation}
Thus the bar cell has the opposite peripheral orientation,
with the displayed integer correction when the numerator
has a nonzero integer part.  Formula \eqref{maps:eq:F.GC5} sends its
group element to $s_i^{-1}x^{a_i}$; consequently the
geometric positive root in the last expression of \eqref{maps:eq:F.GC5a}
maps to $s_i$ exactly.  This is the stated conversion between
the boundary and complement chain conventions.  In a local basis
$Q$, replace $a_i$ in the group and complement formulas by
$Q\lfloor Q^{-1}v_i/m_i\rfloor$.
Writing $\mathcal N_i=\sum_{j=0}^{m_i-1}M_i^j$, the cyclic
degree-two cell has boundary
$m_i\tau_i+w_i-\mathcal N_i a_i$.
Indeed $\tau_i^{m_i}=x^{\mathcal N_i a_i-w_i}$ follows
directly from $g_i^{m_i}=x^{w_i}$.
This establishes the affine power relation and the degree-one
prism sign for every integer numerator.

\item There is an optional exact geometric shortcut.  If a verified
local integral basis gives $M_i=\operatorname{diag}(1,D)$ and
$v_i=\epsilon e_0$, $\epsilon\in\{1,-1\}$, eliminate the
first fibre generator from the cap group using
$g^{m_i}=u^\epsilon$.  The cap is the mapping torus with
group $\mathbb Z^5\rtimes_D\mathbb Z$.
The fibre vector $(u,a)$ maps to $x^a g^{\epsilon m_i u}$,
and the root generator maps to $g$.  Substituting this
homomorphism for $\psi$ in the same prism gives the boundary
map to this bounded geometric cap model.  The shortcut is
used only after checking the displayed integral normal form;
the general periodic construction remains available for
every other input.

\item To check the peripheral convention directly, the universal
cover of $Y_i$ has the root generator $g_i$ with linear part
$M_i$ and affine displacement $v_i/m_i$.  Its integral
translations and $g_i$ map to the same fibre translations
and $s_i$ in \eqref{maps:eq:F.GC2}.  The full base-circle translation obeys
$g_i^{m_i}=x^{w_i}t$, where $w_i$ is the power vector of
\ref{maps:src:F.2.31}.  Filling the disc kills $t$, giving exactly
$s_i^{m_i}=x^{w_i}$.  The straight paths between the chosen
orbit points give \eqref{maps:eq:F.GC5}, so the prism is this marked map,
including numerators with nonzero integer parts.

\item \emph{The marked cusp boundary map.}
Write $P=[A\ B]$ in \eqref{maps:eq:F.GC1}, and set
\begin{equation}
 w_0=(s_1\cdots s_k)^{-1},\qquad
 c=w_0x^{-\kappa}\in\Gamma.
                                                               \hypertarget{tag-02G5}{}\label{maps:eq:F.GC6}
\end{equation}
Send the three angular generators to the columns of $A$,
the three tropical generators to the columns of $B$, and
the base-phase generator to $c$.  Relation \eqref{maps:eq:F.GC1} gives
$c b_i c^{-1}=b_i+a_i$ and all the other nilpotent boundary
relations.  Filling kills the angular generators and $c$,
so \eqref{maps:eq:F.GC6} gives exactly
$s_1\cdots s_kx^\kappa=1$, in agreement with \ref{maps:src:F.2.31}.

\item This homomorphism has a direct chain implementation on the
full Kuhn diagram.  For an incidence with translation label
$w$, assign the group element $x^{Bw}$.  A vertical phase
$(n,a)\in\mathbb Z\oplus\mathbb Z^3$ maps to $c^nx^{Aa}$.
The incidence shear is $S_w(n,a)=(n,a+n w)$, and
\begin{equation}
 (c^nx^{Aa})x^{Bw}=x^{Bw}(c^nx^{A(a+n w)}).                \hypertarget{tag-02G6}{}\label{maps:eq:F.GC7}
\end{equation}
For a flag of $p$ incidence arrows and a vertical bar
simplex of length $q$, sum over all $(p,q)$ shuffles.
Each incidence step contributes its assigned group element;
each vertical step contributes its phase transformed by
all preceding incidence shears.  Give the shuffle the sign
$(-1)^{\#\{\text{incidence steps before vertical steps}\}}$.
Interior face cancellations use exactly \eqref{maps:eq:F.GC7}, and adjacent
incidence faces use addition of their translation labels.
This is therefore a chain map from the double bar complex
to $C_*(B\Gamma)$.  Precompose with the finite Kuhn lift of
\ref{maps:src:F.2.34} and postcompose with \eqref{maps:eq:F.GC3} and the complement
projection.  The other cusp map is the full logarithmic
collapse already computed in \ref{maps:src:F.2.34}.
Both the cusp boundary and the complement are aspherical,
so the displayed marked homomorphism also verifies the
homotopy class of this geometric inclusion.

\item \emph{The final finite integer calculation.}
Take the mapping cone of the just constructed maps
\begin{equation}
 \bigoplus_{i=0}^k C_*(Y_i)\longrightarrow
 C_*(E)\oplus\bigoplus_{i=0}^k C_*(B_i),
 \qquad y\longmapsto(f_E(y),-f_B(y)),                     \hypertarget{tag-02G7}{}\label{maps:eq:F.GC8}
\end{equation}
where $B_0=W$ is the Kuhn cap.  Gluing collars makes this
the homotopy pushout model of the closed manifold.
Consequently its cone computes integral homology.

\item The chain ranks give an explicit finite size bound.  Put
$t_n=\sum_{q=0}^{\min(n,6)}\binom6q$, with binomial
coefficients outside their range interpreted as zero.
Let $w_n$ and $y_n$ be the raw cap and boundary ranks of
\ref{maps:src:F.2.34}, zero outside their listed ranges.  Then the cone rank
in degree $n$ is
\begin{equation}
 \binom6n+k\binom6{n-1}+k t_n+w_n+
 k\binom7{n-1}+y_{n-1}.                                 \hypertarget{tag-02G8}{}\label{maps:eq:F.GC9}
\end{equation}
For $0\leq n\leq9$ its maximum is 3890 when $k=2$ and
4003 when $k=3$.  Every bar comparison and perturbation
series used to produce these matrices terminates by its
degree filtration.  Word reductions and affine orbit enumerations terminate.
Computing integer kernels followed by Smith normal forms
now gives $H_0,\ldots,H_8$; higher groups vanish by dimension.
Unit cancellations are optional exact integral chain
equivalences which substantially reduce the matrices before
Smith reduction.  This proves the algorithm and the bound.
\end{enumerate}
The corresponding exact calculations are implemented in \hyperlink{prog:V:chains}{the programs in Table~\ref*{tab:V:programs}}.
\end{proof}




\begin{remark}[An orientation-reversing gluing]\hypertarget{tag-PV20}{}\label{rem:V:orientation-reversing-gluing}
\StatementRecord{rem:V:orientation-reversing-gluing}{proved}{PV20}{}
Take the gluing with
$A=I_3$, $B=\operatorname{diag}(0,1,0)$ in \eqref{maps:eq:F.IV2},
affine orders $4,4$, numerators $a_1,a_2$, and zero
clutching.  These two linear actions reverse orientation.
Integral unit cancellation and Smith reduction of its gluing cone give free ranks
\[
                  (1,0,4,11,10,1,6,3,0).
\]
The nonzero torsion groups are
\[
\begin{aligned}
 \operatorname{Tor}H_1=\operatorname{Tor}H_5
       &=(\mathbb Z/2)^3\oplus\mathbb Z/4,&
 \operatorname{Tor}H_2&=(\mathbb Z/2)^2,\\
 \operatorname{Tor}H_3=\operatorname{Tor}H_7
       &=\mathbb Z/2,&
 \operatorname{Tor}H_4&=(\mathbb Z/2)^4\oplus(\mathbb Z/4)^5.
\end{aligned}
\]
Abelianizing its van Kampen presentation gives the same $H_1$.  Its Euler number
is six, its mod-two Betti vector is
$(1,4,10,14,20,14,10,4,1)$, and its integral top groups
are those of a closed nonorientable eight-manifold.
\end{remark}










\section{Canonical representatives of integral monodromy}\label{pub:V:sec-canonical}

\begin{theorem}[Simultaneous integral conjugacy for every input]\hypertarget{tag-02NR}{}\label{maps:src:F.2.36}
\StatementRecord{maps:src:F.2.36}{proved}{02NR}{\cite[Algorithm 8.3, p.~21]{rootBHJ22}; \cite[Algorithm 8.3 and its correctness proof, p.~21]{rootBHJ22}; \cite[Theorem 3.4, p.~6; Propositions 8.1--8.2, p.~21]{rootBHJ22}}


There is an unconditional exact algorithm which decides simultaneous
$\mathrm{GL}_6(\mathbb Z)$ conjugacy of any two admissible linear tuples and returns a conjugating matrix when one exists.  More generally it applies to any integral
tuple containing a matrix $1+N$ with $N^2=0$ and $\operatorname{rank}N=3$.

\end{theorem}


\begin{proof}
\begin{enumerate}[label=\arabic*.]
\item We use column vectors and seek $P$ satisfying $PX_i=Y_iP$ for every
marked index, including the cusp index. Put
\begin{equation}
\begin{aligned}
 H&=\{P\in\operatorname{Mat}_6(\mathbb Z):PX_i=Y_iP\ (\forall i)\},\\
 \mathcal O&=\{C\in\operatorname{Mat}_6(\mathbb Z):CX_i=X_iC\ (\forall i)\},
 \qquad E=\mathbb Q\mathcal O.
\end{aligned}                                                    \hypertarget{tag-02GM}{}\label{maps:eq:F.IC1}
\end{equation}
Both lattices are saturated integer kernels of explicitly displayed
linear equations in thirty-six unknowns. Thus Hermite or Smith reduction computes bases of these saturated
integer kernels.

\item We first decide rational isomorphism. For a basis $h_1,\ldots,h_s$ of
$H$, form the polynomial $f(c)=\det(\sum c_jh_j)$. If it is zero,
there is no rational isomorphism. Otherwise successively substitute
one of $0,1,\ldots,6$ for each variable while keeping the remaining
polynomial nonzero. At each stage one of these seven values works,
because the degree in that variable is at most six. The final matrix
$D\in H$ is invertible over $\mathbb Q$. This gives a terminating exact test for rational isomorphism.

\item The full lattice $L=D^{-1}H\subset E$ is a right $\mathcal O$-module.
Let $J$ be the Jacobson radical of $E$, let $B=E/J$, and use bars
for images under the quotient map. We next verify the splitting
hypothesis needed for the principal-ideal algorithm. In a rational
$N$-adapted basis,
\begin{equation}
 N=\begin{pmatrix}0&I_3\\0&0\end{pmatrix},\qquad
 E\subseteq
 \left\{\begin{pmatrix}A&T\\0&A\end{pmatrix}:A,T\in
                                      \operatorname{Mat}_3(\mathbb Q)\right\}.
                                                               \hypertarget{tag-02GN}{}\label{maps:eq:F.IC2}
\end{equation}
Projection to $A$ has square-zero kernel and image an algebra
$S\subseteq\operatorname{Mat}_3(\mathbb Q)$. A nilpotent ideal is
contained in the radical, so $E/J\cong S/J(S)$.

\item Choose a composition series of the faithful $S$-module $\mathbb Q^3$.
The kernel of its action on the associated graded module lowers this
filtration and is therefore nilpotent. Conversely the radical kills
every simple factor. Thus this kernel is exactly $J(S)$ and every
simple factor of $S/J(S)$ acts nontrivially on a composition factor.
For a simple algebra $\operatorname{Mat}_n(D)$, where $D$ has centre
$K$ and degree $d$ over that centre, every nonzero simple module has
$\mathbb Q$-dimension $n d^2[K:\mathbb Q]$. If $D$ were
noncommutative this would be at least four, which is impossible in
$\mathbb Q^3$. Consequently
\begin{equation}
 B\cong\prod_j\operatorname{Mat}_{n_j}(K_j),\qquad
 n_j[K_j:\mathbb Q]\le3.                                    \hypertarget{tag-02GO}{}\label{maps:eq:F.IC3}
\end{equation}
Every noncommutative factor is therefore
$\operatorname{Mat}_2(\mathbb Q)$ or
$\operatorname{Mat}_3(\mathbb Q)$; the remaining factors are number
fields of degree at most three. This proves hypothesis (H) of \cite[hypothesis (H), p.~2; Algorithm~8.3, p.~21]{rootBHJ22} for every tuple under consideration.

\item The decomposition can be computed as follows. Rational linear algebra computes $J$ as the radical of the
regular trace pairing and computes the quotient multiplication table.
Its simple factors are obtained from its centre. A commutative factor
is a number field. In a noncommutative factor of dimension $n^2$, we
enumerate integral combinations $e$ of a rational basis until the
left ideal $Be$ has dimension $n$. Such an $e$ exists because the
factor is $\operatorname{Mat}_n(\mathbb Q)$, so its rank-one matrices
are rational points; clearing denominators expresses one among the
enumerated combinations. Left multiplication on $Be$ then supplies
an explicit isomorphism to $\operatorname{Mat}_n(\mathbb Q)$.
The multiplication table verifies the algebra homomorphism on every
pair of basis elements.

\item The images $\overline{\mathcal O}$ and $\overline L$ are a full order
and a full right lattice in $B$. If the right order of
$\overline L$ differs from $\overline{\mathcal O}$, this lattice
cannot be principal over that order. Otherwise Algorithm 8.3 of
\cite[Algorithm 8.3 and its correctness proof, p.~21]{rootBHJ22}, in its right-ideal form, decides principality and
returns a generator $\overline\alpha$ when one exists. The cited
algorithm passes through a maximal order and its central conductor,
tests local freeness and principality over the maximal order, and
solves the remaining reduced-norm congruence in a finite quotient.
Its correctness proof uses its Theorem 4.5(b) and Theorem 7.9; its
Theorem 8.4 specifies the number-field, finite-field and integer
factorization subproblems. These have unconditional exact algorithms.
Use unconditional class-group and unit-group computations for every
number-field component before applying the principal-ideal decision.

\item If $\overline L$ is not principal, $L$ cannot be principal. If it is
principal, solve the integer linear equations expressing
$\overline\alpha$ as a combination of the images of the chosen basis
of $L$. This produces a lift $\alpha\in L$ and an integral
intertwiner $P=D\alpha\in H$. We claim that a simultaneous integral
conjugacy exists if and only if $|\det P|=1$.

\item The sufficient direction is immediate. For the necessary direction,
suppose $P_0\in H$ is unimodular. Every $h\in H$ has
$P_0^{-1}h\in\mathcal O$, so
$L=\alpha_0\mathcal O$, where $\alpha_0=D^{-1}P_0$.
Write $\alpha=\alpha_0u$ with $u\in\mathcal O$.
Because both image elements generate $\overline L$,
$\overline u$ is a unit of $\overline{\mathcal O}$.
Lift its inverse to $v\in\mathcal O$. Then $uv=1+j$ with
$j\in J\cap\mathcal O$, and the finite geometric series for
$(1+j)^{-1}$ shows that $u^{-1}=v(1+j)^{-1}\in\mathcal O$.
Thus $u$ is an integral unimodular matrix and
$P=P_0u$ is unimodular. This also proves that any chosen lift of the
image generator gives the same yes-or-no answer. It is the
right-module form of \cite[Theorem 3.4, p.~6; Propositions 8.1--8.2, p.~21]{rootBHJ22},
with the needed determinant condition written out.

\item Every branch either proves nonexistence or returns a matrix whose
integrality, determinant and all intertwining equations are verified
exactly. This completes the decision procedure.

\end{enumerate}
\end{proof}




\begin{theorem}[A complete canonical-representative algorithm]\hypertarget{tag-02NS}{}\label{maps:src:F.2.37}
\StatementRecord{maps:src:F.2.37}{proved}{02NS}{\cref{maps:src:F.2.2}, \cref{maps:src:F.2.36}, \cref{maps:src:F.2.38}, \cref{maps:src:F.2.16}}


Every admissible linear tuple, for either $k=2$ or $k=3$,
has a computable canonical integral representative and a
returned simultaneous unimodular conjugator. Two inputs have
the same representative if and only if they are simultaneously
$\mathrm{GL}_6(\mathbb Z)$ conjugate. The algorithm uses a proved
finite bound depending on the input. No fixed global entry bound
can represent every class.

This is an effective classification. The
ordered involution sector has the separate complete finite
normal form \ref{maps:src:F.2.38}, consisting of $56$ classes. 

\end{theorem}


\begin{proof}
\begin{enumerate}[label=\arabic*.]
\item We remove the existential free-lift condition by an explicit
construction. Let $d>1$ be the order of $M$ and choose a primitive
$u\in\ker(M-I)\cap\mathbb Z^6$. Such a vector is obtained from a
nonzero rational kernel vector by clearing denominators and dividing
by the greatest common divisor of its integer coordinates; both
operations retain the equation $(M-I)u=0$. Put
$g(x)=Mx+u/d^2$ and $S_d=\sum_{a=0}^{d-1}M^a$.
Then $g^j(x)=M^jx+ju/d^2$, $S_dM^j=S_d$ and $S_du=du$.
A fixed point of $g^j$ gives
$(M^j-I)x+ju/d^2=\lambda\in\mathbb Z^6$.
Applying $S_d$ gives $ju/d=S_d\lambda\in\mathbb Z^6$.
A Bezout row pairing to one with $u$ therefore gives $d\mid j$.
Writing $j=dr$, the original equation becomes $ru/d=\lambda$,
so the same row gives $d\mid r$. Thus only multiples of $d^2$
can have fixed points. Since $g^{d^2}$ is translation by the
integer vector $u$, its torus action is the identity. This proves
both its affine order $d^2$ and freeness of every nonidentity power.
For $M=I$, translation by $u/2$ for primitive $u$ has affine order
two and is free. Conversely, when $I-M$ is invertible over
$\mathbb R$, every affine map $Mx+a$ fixes the torus class of
$(I-M)^{-1}a$. The free-lift test is consequently exactly
$\ker(M-I)\ne0$.

\item The finite linear orders here are the already proved list of
\cref{maps:src:F.2.2}:
\begin{equation}
 d\in\{1,2,3,4,5,6,8,10,12\}.
 \hypertarget{tag-02FL}{}\label{maps:eq:F.CN1}
\end{equation}
Taking these finitely many powers, checking the fixed kernel,
the product relation and the cusp Smith factors therefore gives
an exact admissibility test. The affine order is the
specified gluing datum; the construction in item~1 gives
its exact relation to the chosen finite lift.

\item For an admissible full tuple $T=(M_1,\ldots,M_k,M_0)$ let
$H(T)$ be its maximum absolute entry. There is at most one identity
finite slot in a three-point tuple: two such identities would make
$M_0$ finite order, whereas $M_0=I+N$, $N^2=0$ and
$\operatorname{rank}N=3$ give $(I+N)^d=I+dN\ne I$ for every
positive integer $d$. At a fixed identity slot, delete that matrix
and reinsert $I$ at the same slot after classifying the remaining
ordered tuple. These two maps are inverse, and both commute with
$T\mapsto (QM_1Q^{-1},\ldots,QM_0Q^{-1})$ for the same
$Q\in\mathrm{GL}(6,\mathbb Z)$. They therefore retain the returned simultaneous conjugator. The remaining invertible integral
matrices have height at least one, so inserting $I$ leaves height
unchanged. Every entry in the inserted block is constant; the first
unequal row-major entry of two tuples is therefore the same before
and after insertion. This proves preservation of the lexicographic comparison as well. The affine order and
numerator at that slot are retained as gluing data and are transported
by this same $Q$. In the two-involution sector use the canonical
representative of \ref{maps:src:F.2.38}, with its own specified
convention. In all remaining sectors, order full integer tuples
first by $H$, then lexicographically by their row-major entries in
the displayed marked order, and take the first admissible member
of the simultaneous conjugacy class.

\item To compute the representative, enumerate heights
$1,\ldots,H(T)$ and all integer entries in the stated order.
There are at most
\begin{equation}
 (2H(T)+1)^{36(k+1)}
 \hypertarget{tag-02FM}{}\label{maps:eq:F.CN2}
\end{equation}
integer lists in this box. Retain the invertible tuples satisfying
the ordered product relation and the stated maximal square-zero
primitive cusp condition, which are finite integer tests, and
apply the decision-and-witness algorithm of
\cref{maps:src:F.2.36} to each candidate and the tuple.
Its general rank-three square-zero input permits precisely this
call. A positive answer transfers all finite-order, orientation
and admissibility conditions through its integral
conjugator. The tuple itself occurs, with the identity
as a possible witness, so a positive answer is reached in
finitely many tests. The first one is the required representative.
Record the returned convention as
$P_T^{-1}T_iP_T=\operatorname{Can}(T)_i$.
The candidate bound depends on $H(T)$; \cref{maps:src:F.2.36} constructs the integral witness $P_T$.

\item Every simultaneous orbit has a least element in this order.
Its height is at most the height of each given member, so
\eqref{maps:eq:F.CN2} includes that global least element.
Equivalent inputs have exactly the same ordered orbit and
therefore the same representative. Conversely, if $T,T'$ have
the same representative, the matrix
$Q=P_TP_{T'}^{-1}$ satisfies $Q^{-1}T_iQ=T'_i$ for every slot.
The separate involution convention has the same property by
\cref{maps:src:F.2.38}, with the product of its returned
witnesses giving the corresponding map. Conjugation preserves
the ordered finite orders, so the involution sector is preserved
by those same maps. This proves the classification assertion
and its finite stopping rule.

\item Finally, \ref{maps:src:F.2.16} already gives an unbounded admissible integral
family. In its homology convention direct symbolic
multiplication gives
\begin{equation}
             \operatorname{tr}(M_1M_2^3)
                       =-v^4+v^3-v^2+2.                 \hypertarget{tag-02FN}{}\label{maps:eq:F.CN3}
\end{equation}
This is unbounded on the free affine subfamily
$v\in20\mathbb Z$ of \cref{maps:src:F.2.16}. It is invariant under
simultaneous conjugacy. A fixed finite entry box contains
only finitely many tuples and hence only finitely many such
traces; it cannot contain a representative of every class.
The input-dependent nature of \eqref{maps:eq:F.CN2} is therefore necessary
for this particular exhaustive procedure.

\end{enumerate}
\end{proof}




\begin{theorem}[The complete ordered involution stratum]\hypertarget{tag-02NT}{}\label{maps:src:F.2.38}
\StatementRecord{maps:src:F.2.38}{proved}{02NT}{\eqref{maps:eq:F.IV1}, \eqref{maps:eq:F.IV3}, \eqref{maps:eq:F.IV4}, \eqref{maps:eq:F.IV2}, \cref{maps:src:F.2.31}}


There are exactly 56 simultaneous $\mathrm{GL}_6(\mathbb Z)$
conjugacy classes of ordered tuples $(M_1,M_2,M_0)$ satisfying
the primitive maximal-cusp conditions of \cref{def:admissible} with both finite linear orders equal to two.  An explicit integral normalizer
returns a canonical representative and an explicit unimodular
conjugating matrix for every such integral input.  Every representative has free affine lifts of
affine order four at both finite points.  None can give a
simply connected filling, for any choices of free affine
orders or twists.

\end{theorem}


\begin{proof}
\begin{enumerate}[label=\arabic*.]
\item First choose an integral cusp basis, so that
\begin{equation}
 M_0=I_6+N,\qquad N=\begin{pmatrix}0&I_3\\0&0\end{pmatrix}.
                                                               \hypertarget{tag-02GP}{}\label{maps:eq:F.IV1}
\end{equation}
This is effective by Smith reduction.  In fact choose vectors
$\beta_i$ whose images $\alpha_i=N\beta_i$ form a primitive
basis of $\operatorname{im}N=\ker N$.
The index-one hypothesis says that their classes form a basis
of the quotient, so $[\alpha_1,\alpha_2,\alpha_3,
\beta_1,\beta_2,\beta_3]$ is unimodular.

\item Since $M_2=M_1(I-N)$ and $M_1^2=M_2^2=I$, we have
$M_1NM_1=-N$.  Solving this block equation gives precisely
\begin{equation}
 M_1=\begin{pmatrix}A&B\\0&-A\end{pmatrix},\qquad
 M_2=\begin{pmatrix}A&B-A\\0&-A\end{pmatrix},\qquad
 A^2=I_3,\quad AB=BA.                                     \hypertarget{tag-02GQ}{}\label{maps:eq:F.IV2}
\end{equation}
Every integral conjugator preserving \eqref{maps:eq:F.IV1} is
$\begin{pmatrix}C&D\\0&C\end{pmatrix}$, with
$C\in\mathrm{GL}_3(\mathbb Z)$ and $D$ arbitrary integral.
Thus one first classifies $A$ integrally; after fixing it,
$B$ is considered modulo $AD+DA$ and modulo conjugation by
the integral centralizer of $A$.

\item We include the integral involution classification needed here.
For the lattice $L=\mathbb Z^3$, let $L_+$ and $L_-$ be the
primitive kernels of $A-I$ and $A+I$.  They satisfy
$2L\subseteq L_+\oplus L_-\subseteq L$.
The quotient $L/(L_+\oplus L_-)$ is a vector space over
$\mathbb F_2$.  Its projections into $L_+/2L_+$ and
$L_-/2L_-$ are injective: a glue vector with only one
nonzero projection would give a new integral vector in the
corresponding eigensublattice, contradicting its definition.
It is consequently the graph of an isomorphism between two
subspaces of equal dimension $r$.
Unimodular changes of bases in $L_\pm$ put this graph in
coordinate form.  For each glued pair $e_+,e_-$ the vectors
$(e_++e_-)/2$ and $(e_+-e_-)/2$ form a swap block.
The unglued vectors give trivial and sign blocks.  Hence
there is a unique type
\begin{equation}
 A\cong I_p\oplus(-I_q)\oplus
 \begin{pmatrix}0&1\\1&0\end{pmatrix}^{\oplus r},
 \qquad p+q+2r=3.                                         \hypertarget{tag-02GR}{}\label{maps:eq:F.IV3}
\end{equation}
The dimensions of the two rational eigenspaces and the
index $[L:L_+\oplus L_-]=2^r$ prove uniqueness.
Here $r$ is zero or one, so this normalization only requires
integer kernels, one Smith reduction and, in the glued case,
one explicit half-sum pair whose integrality is checked.

\item For fixed \eqref{maps:eq:F.IV3}, compute the quotient of the commuting
matrices by the subgroup $\{AD+DA\}$.
On the trivial and sign diagonal blocks the norm is
respectively $2D$ and $-2D$.  The commuting off-diagonal
blocks between the trivial and sign summands are zero.
Every commuting block involving a swap summand is a norm.
For example a trivial-to-swap block is a row $(a,a)$ and
is obtained as $(d,e)+(e,d)$; a sign-to-swap block has the
same calculation with a difference.  For two swap summands,
the commuting block is
$\begin{pmatrix}a&b\\b&a\end{pmatrix}$, and applying the
norm to $\begin{pmatrix}d&e\\f&g\end{pmatrix}$ gives
$\begin{pmatrix}e+f&d+g\\d+g&e+f\end{pmatrix}$.
The reverse mixed blocks have the transposed calculations.
Therefore the full integral quotient is exactly
\begin{equation}
 \operatorname{End}_{\mathbb Z[A]}(L)/\{AD+DA\}
 \cong \operatorname{Mat}_p(\mathbb F_2)
       \oplus\operatorname{Mat}_q(\mathbb F_2).             \hypertarget{tag-02GS}{}\label{maps:eq:F.IV4}
\end{equation}
The displayed block formulas are surjective over $\mathbb Z$, so
they compute the full norm image.

\item The centralizer action on \eqref{maps:eq:F.IV4} is precisely independent
similarity in the two matrix factors.  Indeed, in a product
of commuting block matrices, every contribution factoring
through a swap block is a norm.  The remaining diagonal
blocks are just the reductions of the trivial and sign
blocks of the centralizer matrix.  If that matrix is an
integral unit, its inverse shows that both reductions are
invertible.  Conversely block-diagonal unimodular changes
on the trivial and sign summands induce every element of
$\mathrm{GL}_p(\mathbb F_2)\times
\mathrm{GL}_q(\mathbb F_2)$; elementary matrices give
unimodular lifts.  An intertwiner of a trivial and a sign lattice is zero over $\mathbb Z$, which proves independence of the two factors.

\item The remaining classification is over the finite field.  Enumerate
all $2^{n^2}$ matrices of size $n=0,1,2,3$ over $\mathbb F_2$
and conjugate by every invertible matrix of that size.
Their respective numbers of similarity classes are
$1,2,6,14$.  Taking the smallest bit-coded representative
in each orbit gives a fixed canonical choice.
The six integral involution types and their counts are
\[
\begin{array}{c|rrrrrr}
 (p,q,r)&(0,3,0)&(1,2,0)&(2,1,0)&(3,0,0)&(0,1,1)&(1,0,1)\\\hline
 \text{classes}&14&12&12&14&2&2
\end{array}
\]
and their sum is 56.  A representative is \eqref{maps:eq:F.IV2} with
$B$ equal to the chosen $0$--$1$ representatives in its
trivial and sign blocks, and zero in every other block.

\item The construction also gives an effective normalizer for every input.  Normalize its
cusp and then its involution by the preceding Smith and
glue constructions.  Use the finite-field orbit lookup to
normalize its two reductions.  In dimensions at most three
the binary lift of an invertible binary matrix is already
unimodular: its determinant is odd and has absolute value
at most two.  The latter bound follows by adjoining a row
and column of ones and applying the $4\times4$ Hadamard
bound.  Finally solve the full integral linear equation
$AD+DA=B-B_{\rm can}$ by Smith reduction.  The preceding
norm calculation proves solvability.  Conjugating by
$\begin{pmatrix}I&D\\0&I\end{pmatrix}$ subtracts this norm
from the upper-right block.  The product of these explicitly
recorded conjugators is integral unimodular and carries all
three marked matrices to their canonical representative.

\item Both matrices in \eqref{maps:eq:F.IV2} have a three-dimensional fixed
space.  Choose a primitive integral fixed vector $v_i$.
Then $g_i(x)=M_ix+v_i/4$ has affine order four:
$g_i^2$ is the nontrivial translation $v_i/2$.
This order-two subgroup is free, and every nontrivial
subgroup of a cyclic group of order four contains it, so
the full action is free.  For any other finite affine lift
the affine order is still even, since its linear part has
order two.  The two-point orbifold quotient of \ref{maps:src:F.2.31} is
therefore the nontrivial cyclic group of order
$\gcd(m_1,m_2)$.  Simple connectivity is impossible for
every choice of twists and affine orders.

\end{enumerate}
\end{proof}




\begin{theorem}[The full effective classification and all marked gluings]\hypertarget{tag-02NU}{}\label{maps:src:F.2.39}
\StatementRecord{maps:src:F.2.39}{proved}{02NU}{\cref{maps:src:F.2.37}, \cref{maps:src:F.2.38}, \cref{maps:src:F.2.31}, \cref{maps:src:F.2.35}, \cref{maps:src:F.2.30}, \cref{maps:src:F.2.1}, \cref{maps:src:F.2.2}, \cref{maps:src:F.2.36}, \eqref{maps:eq:F.CN3}, \cref{maps:src:F.2.33}, \cref{maps:src:F.2.7}}


The admissible data and the topology of their gluings are classified as follows.
For $k=2$ and $k=3$, every admissible ordered integral tuple
has the canonical representative of \ref{maps:src:F.2.37}, with an explicit
simultaneous unimodular conjugator.  Equality of these
representatives is necessary and sufficient for simultaneous
integral conjugacy.  The finite stopping bound for an input of height $H$ is
$(2H+1)^{36(k+1)}$ candidates before the proved rejection
filters.  The separate involution normal form has exactly
56 classes and is computed completely in \ref{maps:src:F.2.38}.

For every choice of free affine orders and integer
numerators, integral clutch vector, and specified primitive
Kuhn cusp marking, \ref{maps:src:F.2.31} and \ref{maps:src:F.2.35} compute the full
fundamental group and every integral homology group of the
gluing.    The complete list of data
with $\pi_1=1$ and $H_2=H_3=0$ is the constructive
parametrization of \ref{maps:src:F.2.30}, with its exact lattice enumeration
and necessary-and-sufficient finite tests.  Each such filling has the homology and oriented smooth type stated
there.

\end{theorem}


\begin{proof}
\begin{enumerate}[label=\arabic*.]
\item The input is finite, so the prime-subgroup test in \ref{maps:src:F.2.31}
decides whether every claimed affine cyclic action is free
and has the stated affine order.  The integral cusp test and
its adapted basis are \ref{maps:src:F.2.1}.  The complete linear admissibility
filter and the canonical representative, including the proved
finite stopping bound, are \ref{maps:src:F.2.2} and \ref{maps:src:F.2.36}--38.  The
characteristic-polynomial filters used by that algorithm are
conjugacy invariants, so none can remove a possible
representative.  The unbounded trace family \eqref{maps:eq:F.CN3} proves
that no fixed finite box could replace the input-dependent
bound.

\item We spell out how the linear classification is used with a
specified gluing.  If $R$ is its returned conjugator and
$P$ is the cusp marking, change the data by
\begin{equation}
 M_i'=RM_iR^{-1},\qquad v_i'=Rv_i,\qquad
 \kappa'=R\kappa,\qquad P'=RP.                           \hypertarget{tag-02I6}{}\label{maps:eq:F.SY1}
\end{equation}
The affine orders are unchanged.  The coordinate map
$x\mapsto Rx$ on the torus intertwines the finite quotient
charts, since it sends $M_ix+v_i/m_i$ to
$M_i'x'+v_i'/m_i$.  At the cusp, its coordinates in the
marked Kuhn model are unchanged:
$(P')^{-1}x'=P^{-1}x$.  Hence it intertwines every toric
chart, translated face identification and collapse map of
that fixed filling.  It also carries the clutch path to
the path with vector $\kappa'$.  These maps agree on the
collars and give a diffeomorphism of the two fully marked
gluings.  Thus homology can be computed in the original
coordinates, while the canonical data and their group
presentation are also returned. The cusp marking in this
comparison is $P'=RP$.

\item An affine order-one site is ordinary, even if its integer
numerator has been recorded.  Delete it by the explicit comparison of \cref{maps:src:F.2.31}: if it is the $i$th site, replace
$\kappa$ by
$\kappa+(M_{i+1}\cdots M_k)^{-1}v_i$ before deletion.
The cusp matrix and its given marking are unchanged.
There cannot be an order-one site in an admissible two-point
tuple: its other finite matrix would be the inverse of a
nontrivial unipotent matrix.  Nor can a three-point tuple
have two such sites.  Thus this reduction leaves precisely
the genuine two- or three-point inputs of \ref{maps:src:F.2.35} and loses
no case from the input class.

\item For each remaining input, the exact unsaturated normal
closure, cyclic or triangle-group extension and exceptional
spherical Euler relation of \ref{maps:src:F.2.31} give the full fundamental
group.  The finite cap, complement and marked Kuhn chain
maps of \cref{maps:src:F.2.33,maps:src:F.2.34,maps:src:F.2.35} give a finite integer Mayer--Vietoris cone
through degree nine.  Integer kernels and Smith reduction
therefore give all $H_0,\ldots,H_8$ with their complete
torsion factors. Its degree-one homology agrees with the
abelianization of the van Kampen presentation.

\item The results from \cref{maps:src:F.2.7} through \cref{maps:src:F.2.30} prove both directions of the classification.  Genuine three-point orbifold quotients
exclude simple connectivity.  The two-point filters reduce
every three-connected datum to the minimal order-$(5,4)$ spectrum; the
rational normal form, finite stable-lattice and extension
algorithms include every integral representation above it.
The universal affine normalization includes every twist,
and the six exact integral tests decide three-connectivity. Conversely every datum that passes them satisfies all three
topological conditions by those tests.  \ref{maps:src:F.2.30} computes
its common integral homology and applies the almost-complex
and smooth-classification arguments to identify its oriented
type. Applying the canonical-representative algorithm identifies
precisely the simultaneously integrally conjugate inputs in this
parametrization.

\end{enumerate}
\end{proof}




\section{Source obstructions and the punctured-base period decision}\label{pub:V:sec-oka}



\begin{theorem}[Exclusion of the abelian-surface source]\hypertarget{tag-02O6}{}\label{maps:src:F.3.4}
\StatementRecord{maps:src:F.3.4}{proved}{02O6}{\cref{maps:src:F.3.1}; \cite[Definition 4.2.4 and Proposition 4.2.5, p.~44; Theorem 4.2.6, p.~45]{ffDeligne}}


The construction of \cite{E26}, starting with a polarized
algebraic family of abelian surfaces with a toric-rank-two cusp,
taking a fibrewise degree-zero line bundle, and quotienting its
punctured bundle by a lifted translation, cannot have $b_1=b_2=0$
with the stipulated rank-three square-zero Kuhn cusp.
The two sections need not be proportional.

\end{theorem}


\begin{proof}
\begin{enumerate}[label=\arabic*.]
\item Let $W$ be the rank-four rational homology local system of the
abelian surfaces.  It is semisimple by
\cite[Definition 4.2.4 and Proposition 4.2.5, p.~44; Theorem 4.2.6, p.~45]{ffDeligne}.
Its toric-rank-two cusp has
$N_W^2=0$ and $\ker N_W=\operatorname{im}N_W$ of dimension two.
Semisimplicity splits $W=W^\Gamma\oplus W'$.
The image of $N_W$ lies in $W'$ whereas $W^\Gamma$ lies in its
kernel.  The equality of image and kernel forces $W^\Gamma=0$.
The coinvariants vanish as well.

\item The rank-six rational homology system $T$ after the two Engel
operations has a natural filtration
\[
 0\subset L\subset E\subset T,\qquad
 L\simeq\mathbb Q,\quad E/L\simeq W,\quad T/E\simeq\mathbb Q.
\]
The lower line is the circle in the punctured line bundle and
the upper quotient is the lifted-translation generator.
Here polarization is used through the semisimplicity of the middle term $W$.
The triviality of $T/E$ gives $\operatorname{im}N_T\subset E$.
Maximal square-zero cusp monodromy gives
$\ker N_T=\operatorname{im}N_T$; thus every globally fixed vector
lies in $E$.  Since $W^\Gamma=0$, the fixed vectors of $E$ are
exactly $L$.  Hence $\dim T^\Gamma=1$.
The dual filtration, whose middle term is $W^*$, gives
$\dim(T^*)^\Gamma=1$ by the identical argument.

\item Consequently the two ranks $a,b$ in Theorem \ref{maps:src:F.3.1} are both one.
Its Euler-parity calculation applies without requiring self-duality
of $T$, and contradicts $b_1=b_2=0$.
This excludes the algebraic abelian-surface construction and,
more generally, every such filtration with semisimple middle term.

\end{enumerate}
\end{proof}




\begin{theorem}[The unpolarized period domain has LSAP]\hypertarget{tag-02OA}{}\label{maps:src:F.3.8}
\StatementRecord{maps:src:F.3.8}{proved}{02OA}{\eqref{maps:eq:F.O3}, \eqref{maps:eq:F.O4}; \cite[Definition 2.1, p.~4]{IVVOkaOneV5}; \cite[Proposition 7.13 and preceding paragraph, pp.~23--24]{IVVOkaOneV5}}


Let $D_3$ be the component of
\[
 \{F\in\operatorname{Gr}(3,\mathbb C^6):F\cap\overline F=0\}
\]
containing the standard complex structure, so that
$D_3\cong GL^+(6,\mathbb R)/GL(3,\mathbb C)$.
At every point, $D_3$
is dominated by a spanning tube of nine complex lines.  In particular,
it has the local spray approximation property (LSAP) and is Oka--1.

\end{theorem}


\begin{proof}
\begin{enumerate}[label=\arabic*.]
\item Use the graph chart $F_Z=\operatorname{im}\binom Z{I_3}$.
The open set
\begin{equation}
 \Omega=\{Z\in M_3(\mathbb C):\det(\operatorname{Im}Z)>0\}
 \hypertarget{tag-02HC}{}\label{maps:eq:F.O1}
\end{equation}
lies in $D_3$ and contains $iI_3$.  For a real column $w$, write
$K_w x=w\times x$, put $q_w=w^tw$, and define
\begin{equation}
 A_w=q_wI_3-ww^t+K_w.
 \hypertarget{tag-02HD}{}\label{maps:eq:F.O2}
\end{equation}
Take the six vectors
\[
 e_1,\ e_2,\ e_3,\ e_1+e_2,\ e_1+e_3,\ e_2+e_3
\]
and the corresponding six matrices $A_w$, followed by the three
matrices $K_{e_1},K_{e_2},K_{e_3}$.  Call these nine matrices
$A_1,\ldots,A_9$.  Their symmetric parts are a basis of the
six-dimensional space of real symmetric matrices: the first three
give all diagonal matrices and the last three add the three distinct
off-diagonal entries.  The final three skew matrices form a basis
of the skew matrices.  Thus the nine directions are linearly
independent over $\mathbb R$ and also over $\mathbb C$.
For reference, the determinant of their nine row-major coordinate
columns is $-16$.

\item Each $A_w$ is real normal, is zero on $\mathbb Rw$, and has
eigenvalues $q_w\pm i\sqrt{q_w}$ on $w^\perp\otimes\mathbb C$.
For every real $s$ its two transverse squared singular values are
\begin{equation}
 (1+q_ws)^2+q_ws^2
 =q_w(q_w+1)\left(s+\frac1{q_w+1}\right)^2
       +\frac1{q_w+1}\ge\frac13,
 \hypertarget{tag-02HE}{}\label{maps:eq:F.O3}
\end{equation}
since $q_w\in\{1,2\}$; the remaining squared singular value is one.
For a pure skew direction $K_{e_j}$ all singular values of
$I_3+sK_{e_j}$ are at least one.  Their determinants are positive
at $s=0$ and never vanish, so they are positive for all real $s$.

\item Let $\Lambda=\bigcup_{j=1}^9(iI_3+\mathbb C A_j)$ and let
$B_{1/4}$ be the open Frobenius ball of radius $1/4$ in
$M_3(\mathbb C)$.  Then
\begin{equation}
 T=\Lambda+B_{1/4}\subset\Omega.
 \hypertarget{tag-02HF}{}\label{maps:eq:F.O4}
\end{equation}
Indeed, the imaginary part of $iI_3+tA_j+E$ is
$I_3+(\operatorname{Im}t)A_j+\operatorname{Im}E$.
Its unperturbed least singular value is at least $1/\sqrt3$,
whereas $\|\operatorname{Im}E\|_{\rm op}<1/4$.
It therefore stays invertible throughout the straight perturbation
from $E=0$, preserving the positive determinant.
The set $T$ is open and connected and is a union of affine
translates of the spanning tree $\Lambda$.  Thus it is a spanning tube in the sense of
\cite[Definition 2.1, p.~4]{IVVOkaOneV5}.  Its inclusion in
$D_3$ is a local biholomorphism at $iI_3$.

\item The real group $GL^+(6,\mathbb R)$ acts holomorphically and
transitively on $D_3$.  Transporting this tube gives the same
domination at every point.  The observation preceding
\cite[Proposition 7.13 and preceding paragraph, pp.~23--24]{IVVOkaOneV5} gives LSAP from
this domination, and its part (b), or Theorem 2.2 of that paper,
gives Oka--1.

\end{enumerate}
The corresponding exact calculations are implemented in \hyperlink{prog:V:tubes}{the programs in Table~\ref*{tab:V:programs}}.
\end{proof}




\begin{theorem}[Holomorphic periods on every open base]\hypertarget{tag-02OB}{}\label{maps:src:F.3.9}
\StatementRecord{maps:src:F.3.9}{proved}{02OB}{\cref{maps:src:F.3.8}; \cite[Proposition 7.13(d), p.~24]{IVVOkaOneV5}}


Let $B$ be an open connected Riemann surface and
$\rho:\pi_1(B)\to GL^+(6,\mathbb R)$ any representation.
Its associated flat $D_3$-bundle admits a holomorphic section.
More generally, every continuous section which is holomorphic
near a compact Runge set can be deformed through sections to a
holomorphic section, with arbitrary uniform approximation there
and any prescribed finite-order interpolation at finitely many
points of that set.  For an integral representation, this gives
a holomorphic family of complex three-tori with the prescribed
lattice local system and a zero section over $B$.

\end{theorem}


\begin{proof}
\begin{enumerate}[label=\arabic*.]
\item The flat bundle
$\mathscr D=\widetilde B\times_\rho D_3\to B$ is a holomorphic
fibre bundle: its transition maps are constant real linear
transformations, acting holomorphically on the Grassmannian.
Its fibre has LSAP by Theorem \ref{maps:src:F.3.8}, so the bundle map is an
Oka--1 map by \cite[Proposition 7.13(d), p.~24]{IVVOkaOneV5}.
The lifting property in Definition 7.7 of that paper, applied
to the identity map of $B$, gives precisely the stated section
deformation and approximation, including the finite jets.


\item There is a continuous section to start with.  An open Riemann
surface has the homotopy type of a one-dimensional CW complex;
a bundle with connected fibre admits a section over its vertices
and edges, and the homotopy lifting property transports this
section to $B$.  Connectedness holds for the chosen component
$D_3$.  Taking the Runge set empty therefore proves existence.

\item For an integral representation let $\mathbb V_\mathbb Z$ be its
lattice local system and let $\mathcal F\subset
\mathbb V_\mathbb Z\otimes\mathcal O_B$ be the resulting
holomorphic rank-three subbundle.  Real transversality says that
projection identifies the real vector bundle with the underlying
real bundle of
$\mathcal Q=(\mathbb V_\mathbb Z\otimes\mathcal O_B)/\mathcal F$.
Hence the image of $\mathbb V_\mathbb Z$ is a full lattice in
every fibre of $\mathcal Q$.  Locally it is uniformly discrete
and acts by holomorphic translations.  The quotient
$\mathcal Q/\mathbb V_\mathbb Z\to B$ is consequently a smooth
proper holomorphic family of compact complex three-tori,
with its evident zero section and the required marking.

\end{enumerate}
\end{proof}




\begin{lemma}[Isolated extension points]\hypertarget{tag-02O2}{}\label{maps:src:F.3.10}
\StatementRecord{maps:src:F.3.10}{proved}{02O2}{\eqref{maps:eq:F.O5}; \cite[Theorem 1.2, preprint p.~3]{ffForstnericPrezelj02}; \cite[Theorem 2.13, p.~9]{IVVOkaOneV5}}


Let $R$ be an open Riemann surface, $V\to R$ a holomorphic vector
bundle, and $\mathscr U\subset\operatorname{Gr}_r(V)$ an open
subset whose projection onto $R$ is surjective.  Let $S\subset R$
be finite.  Suppose that over $R\setminus S$ this projection is
a holomorphic fibre bundle with an LSAP fibre.  Let $K\subset R$
be a compact Runge set containing $S$ in its interior.
Every continuous section of $\mathscr U$ which is holomorphic
near $K$ can be deformed through sections to a holomorphic
section, with arbitrary uniform approximation on $K$ and
interpolation of any given finite jets at finitely many points
of its interior.

\end{lemma}


\begin{proof}
\begin{enumerate}[label=\arabic*.]
\item We give the relative version of the LSAP argument
in \cite[arXiv v5, Lemma~5.1, p.~14; Theorem~6.1, p.~17; Proposition~7.13(d), p.~24]{IVVOkaOneV5}.    Enlarge $K$ slightly inside
the domain of holomorphicity, retaining the Runge property.
Choose a strongly subharmonic Morse exhaustion of $R$ starting
with this compact set.  All subsequent boundary arcs, attached
discs and critical one-handles can and will be chosen disjoint
from $S$.  On them $\mathscr U$ is an ordinary holomorphic
bundle with LSAP fibre.

\item The global spray on the old compact set can be taken to be
vertical even though the fibres at $S$ differ.  Indeed, let $f$
be the current holomorphic section on a Stein neighborhood $W$
of that set.  Put $A$ equal to the finite interpolation set,
including $S$, and choose an ideal sheaf $\mathcal I$ vanishing
to one more than the prescribed jet orders at its points.
Cartan's theorem A supplies finitely many holomorphic sections
$T_1,\ldots,T_N$ of
$\operatorname{End}(V)|_W\otimes\mathcal I$ which generate
$\operatorname{End}(V)$ on the compact attaching arcs, after
shrinking $W$ if needed.  The formula
\begin{equation}
 F(z,a)=\exp\left(\sum_{j=1}^N a_jT_j(z)\right)f(z)
 \hypertarget{tag-02HG}{}\label{maps:eq:F.O5}
\end{equation}
defines a holomorphic spray of sections of the Grassmannian
bundle.  Its differential in $a$ surjects onto the vertical
tangent space at every attaching point, since
$\operatorname{End}(V_z)\to
\operatorname{Hom}(f(z),V_z/f(z))$ is onto.
For $a$ sufficiently small it has range in $\mathscr U$ on
a neighborhood of the old compact set, by openness and
compactness.  Every parameter value preserves all the chosen
jets.  

\item For a noncritical extension, subdivide the new annular regions
into small discs lying in bundle trivializations over
$R\setminus S$.  Restrict \eqref{maps:eq:F.O5} to a thin attaching disc.
LSAP of the fibre extends this spray across the attached disc,
approximating the old spray as closely as desired on the
intersection.  The spray gluing used in the proof of
\cite[arXiv v5, Lemma~5.1, equations~(5.6)--(5.8), pp.~16--17]{IVVOkaOneV5} now applies verbatim to vertical sprays:
the inverse-function step is needed only on the attaching
intersection, where \eqref{maps:eq:F.O5} is dominating.  It produces
reparametrizations of the two sprays and a glued section.
The old part remains in $\mathscr U$ because it is evaluated
in \eqref{maps:eq:F.O5}; the new part remains in $\mathscr U$ because the
LSAP spray takes values in the bundle fibre.  The fixed jets
survive the reparametrizations, being independent of the
spray parameters on the old side.

\item For a critical extension, first extend over the attached real
arc using the given continuous section and the homotopy already
constructed on the old compact set.  Apply Mergelyan
approximation with the same finite jets to a neighborhood of
the union of the old set and this arc, and then apply the
noncritical step.  To see explicitly that this approximation
is by sections, trivialize $V$ holomorphically on the open
Riemann surface: its underlying complex bundle is trivial over
a graph spine, and the frame bundle therefore has a continuous
section.  Its fibre $GL(\operatorname{rank}V,\mathbb C)$ has
the dominating spray $(g,A)\mapsto g\exp A$, so
\cite[Theorem 1.2, preprint p.~3]{ffForstnericPrezelj02} gives a holomorphic
frame.  The section is then a map to the
fixed Grassmannian, to which
\cite[Theorem 2.13, p.~9]{IVVOkaOneV5} applies.  The base
coordinate is kept unchanged.  Close approximation on the
compact arc and old set, followed by shrinking its
neighborhood, preserves membership in the open subset
$\mathscr U$.  This is the critical case of Theorem 6.1 of that paper. The
projection is a submersion everywhere because $\mathscr U$ is
open in a Grassmannian bundle.

\item Choose successive errors summably small on the preceding
compact sets and smaller than their distances to the complement
of $\mathscr U$.  The sections converge uniformly on every
compact set to a holomorphic section with range in
$\mathscr U$.  The finite jets persist.  The close homotopies
on successive stages concatenate in the usual exhaustion
construction.  All new attachment regions lie outside $S$;
the old section remains holomorphic on a full neighborhood of
each point of $S$ at every stage.  This proves the lemma.

\end{enumerate}
\end{proof}




\begin{theorem}[The bounded-cusp period decision]\hypertarget{tag-02O3}{}\label{maps:src:F.3.11}
\StatementRecord{maps:src:F.3.11}{proved}{02O3}{\cref{maps:src:F.2.1}, \eqref{maps:eq:F.O7}, \cref{maps:src:F.3.8}, \cref{maps:src:F.3.10}, \cref{maps:src:F.2.7}, \cref{maps:src:F.2.15}}


Let
$B^*=\mathbb P^1\setminus\{p_0,p_1,\ldots,p_k\}$,
where $k\ge1$, and let
$\rho:\pi_1(B^*)\to GL^+(6,\mathbb R)$ be a representation.
Assume that the monodromy $M_0$ at $p_0$ satisfies
$(M_0-I)^2=0$ and $\operatorname{rank}(M_0-I)=3$.
Then there is a holomorphic period family with this monodromy
and nondegenerate real transversality everywhere on $B^*$,
equivalently a nondegenerate imaginary part in each local
normalized period chart.
In a marking with $M_0=\left(\begin{smallmatrix}I&I\\0&I\end{smallmatrix}\right)$
its cusp period is
\begin{equation}
 Z(t)=\frac{\log t}{2\pi i}I_3+C(t),
 \qquad C(t)\text{ holomorphic at }t=0.
 \hypertarget{tag-02HH}{}\label{maps:eq:F.O6}
\end{equation}
Any prescribed finite jet of $C$ at zero may be imposed.
For an integral primitive cusp the marking can be integral.
Consequently every admissible datum with three-connected gluing
admits such a period family on $B^*$. The points
$p_1,\ldots,p_k$ remain punctures of its base.

\end{theorem}


\begin{proof}
\begin{enumerate}[label=\arabic*.]
\item Write $N=M_0-I$ and choose a real basis with
$N=\left(\begin{smallmatrix}0&I\\0&0\end{smallmatrix}\right)$.
Its integral version is \cref{maps:src:F.2.1}.  Choose the component
of the period domain selected by the positive logarithmic
model in \eqref{maps:eq:F.O6}.  This choice matters only when an orientation
has been fixed independently of the marking; every element
of $GL^+(6,\mathbb R)$ preserves either component.

\item Adjoin the cusp to obtain the open surface
\[
 B=\mathbb P^1\setminus\{p_1,\ldots,p_k\}.
\]
The flat complex vector bundle on $B^*$ has a holomorphic
extension $V_D$ across $p_0$: in a local coordinate $t$ use
the single-valued frame obtained by the change
$\exp(-sN)$, where $s=(\log t)/(2\pi i)$.
Changing the logarithm by $2\pi i$ is cancelled by $M_0$.
This specifies the extension directly.

\item Over $B^*$ let $\mathscr U$ be the flat bundle of real-transverse
planes in the chosen component.  At $p_0$ adjoin only the
planes transverse to $\ker N_\mathbb C$, an affine space of
matrices $C\in M_3(\mathbb C)$.  This gives an \emph{open}
subset $\mathscr U\subset\operatorname{Gr}_3(V_D)$.
Indeed, in the chart $F_C=\operatorname{im}\binom C{I_3}$
the flat period is
$\exp(sN)F_C=F_{sI_3+C}$.  For $t\ne0$ membership is exactly
\begin{equation}
 \det\bigl(y(t)I_3+\operatorname{Im}C\bigr)>0,
 \qquad y(t)=-\frac{\log|t|}{2\pi}.
 \hypertarget{tag-02HI}{}\label{maps:eq:F.O7}
\end{equation}
For every fixed $C_0$, and uniformly for $C$ in a sufficiently
small bounded neighborhood of $C_0$, the matrix in \eqref{maps:eq:F.O7} is
invertible with positive determinant when $|t|$ is sufficiently
small.  This follows, for example, from
$y(t)>\|\operatorname{Im}C\|_{\rm op}$ and a straight-line
homotopy to $y(t)I_3$.  Thus a full neighborhood of
$(0,C_0)$ lies in the adjoined set, proving openness.
The projection is consequently a holomorphic submersion. Its fibre
at $p_0$ is the affine chart $\mathbb C^9$ just described; over
$B^*$ its fibre is $D_3$.

\item Start with a polynomial $C_0(t)$ realizing the desired finite
jet.  Formula \eqref{maps:eq:F.O7} makes it a holomorphic section on a small
closed cusp disc $K$, which is Runge in $B$.
This section extends continuously over $B$.  To justify the
relative extension, $D_3$ is connected and simply connected:
it retracts onto $SO(6)/U(3)$, and
$\pi_1(U(3))\to\pi_1(SO(6))$ is onto, since a unitary rotation
in one complex coordinate is a generator of the real
two-plane rotation loop.  A cellular decomposition of the
surface outside the cusp disc now extends its boundary
section over edges by connectedness and over two-cells by
simple connectivity.  The flat bundle is trivial over each cell.

\item Over $B\setminus\{p_0\}$ the fibre is the LSAP manifold
$D_3$ from Theorem \ref{maps:src:F.3.8}.  Lemma \ref{maps:src:F.3.10}, with $S=\{p_0\}$,
produces a holomorphic section of $\mathscr U$ over all of
$B$, with the chosen jet at $p_0$.  Since its value at $p_0$
lies in the affine chart just adjoined, it has a holomorphic
matrix $C(t)$ there.  Undoing the canonical frame gives
exactly \eqref{maps:eq:F.O6}, with a bounded holomorphic correction.
Everywhere else membership in $\mathscr U$ is precisely
nondegeneracy of the complex-torus periods.
More explicitly, choose three members of a real lattice basis
whose projected periods are complex-linearly independent, and
normalize them to $I_3$.  The remaining three periods form a
matrix $Z$, and real transversality is exactly
$\det\operatorname{Im}Z\ne0$.  This normalization persists
locally.  

\item For admissible data with three-connected gluing, Theorem \ref{maps:src:F.2.7} excludes $k=3$ and
Proposition \ref{maps:src:F.2.15} proves orientation preservation of the
finite matrices for $k=2$.  The cusp determinant is also one.
Thus all three-connected data satisfy the hypotheses.  

\end{enumerate}
\end{proof}





\clearpage
\part{The resolved fibre product}\label{pub:V:part-resolved}
\section{The resolved fibre-product fourfold}\label{pub:V:sec-resolved}

\begin{definition}[The resolved-product datum]
\hypertarget{tag-PV21}{}\label{def:V:resolved-product-data}
\StatementRecord{def:V:resolved-product-data}{proved}{PV21}{}
Let $f:X\to B=\mathbb P^1_t$ be a canonical marked member of the
$(3,4,\infty)$ family, with $p_1=0$, $p_2=1$, $p_0=\infty$,
sphere twists $v_1=(1,2,-4,0)^t$, $v_2=(-1,-3,3,0)^t$,
and the untwisted cusp collar. Put $B^\circ=B\setminus\{p_0,p_1,p_2\}$.
Write $S_1,S_2$ for the reduced multiple fibres, of multiplicities
three and four, and $W$ for the cusp divisor.
Let $S\to B$ be the smooth relatively minimal elliptic surface
whose Weierstrass model is
\[
 y^2=x^3-3t^3(t-1)x+2t^4(t-1)^2.
\]
It has $j=1728t$ and singular fibres of types $IV^*$, $III$, and
$I_1$ at $p_1,p_2,p_0$, respectively \cite[\S2, p.~3]{E26}.

Use the homology bases $(\widehat\gamma,\widehat u,
\widehat w,\widehat\delta)$ for $X$ and $(e_1,e_2)$ for $S$,
with dual basis $(\gamma,u^*,w^*,\delta^*,e_1^*,e_2^*)$ for
their direct sum. The marked homology monodromies are
\[
 A_1=\begin{pmatrix}
 1&0&0&0\\6&0&1&0\\-6&-1&-1&0\\-2&1&0&1
 \end{pmatrix},\qquad
 A_2=\begin{pmatrix}
 1&0&0&0\\0&0&-1&0\\-6&1&0&0\\3&0&1&1
 \end{pmatrix},\qquad A_0=(A_1A_2)^{-1},
\]
\[
 C_1=\begin{pmatrix}-1&-1\\1&0\end{pmatrix},\qquad
 C_2=\begin{pmatrix}0&1\\-1&0\end{pmatrix},\qquad
 C_0=(C_1C_2)^{-1}=\begin{pmatrix}1&1\\0&1\end{pmatrix}.
\]
Write $\Pi_X$ for the period matrix of $X$ in this canonical
marking; its last two columns are $I_2$, as in
\SeriesRef{def:common-conventions}.
Orient complex manifolds by their complex structures and use
integral cohomology unless a coefficient field is specified.
Put $L=\mathcal O_X(S_1-S_2)$.
The meromorphic functions $t$, $t-1$ and their quotient give
$3S_1\sim4S_2\sim W$. Tensoring these divisor isomorphisms gives
\[
 L^3\cong\mathcal O_X(S_2),\qquad
 L^4\cong\mathcal O_X(S_1),\qquad
 L^6\cong\mathcal O_X(2S_2),\qquad
 L^{12}\cong\mathcal O_X(W).
\]
\end{definition}

\begin{theorem}[A smooth compact model of the fibre product]\hypertarget{tag-02M5}{}\label{maps:src:E.1.2}
\StatementRecord{maps:src:E.1.2}{proved}{02M5}{\cref{def:V:resolved-product-data}, \cref{night:P1.9}, \eqref{maps:eq:E.1}, \eqref{maps:eq:E.4}, \eqref{maps:eq:E.3}}


There is a smooth compact complex fourfold $Z$ with proper holomorphic
maps $g:Z\to B$ and $\Pi:Z\to X$, with $g=f\Pi$, with the local models described below, agreeing with the marked fibre product over the punctured base.  The map $\Pi$
has a holomorphic section.  The construction is a small crepant
resolution of a Weierstrass model over $X$, and is bimeromorphic to
$X\times_B S$.

\end{theorem}


\begin{proof}
\begin{enumerate}[label=\arabic*.,leftmargin=*]
\item Let $a$ be a nonzero scalar multiple of the divisor section of $L^4$
and let $b$ be a nonzero scalar multiple of the square of the divisor
section of $\mathcal O_X(S_2)=L^3$.  Choose the scalars so that
$4a^3+27b^2$ is the pullback of the section of $\mathcal O_B(1)$
vanishing at $p_0$.  This is possible because $a^3$ and $b^2$ have
divisors $3S_1$ and $4S_2$, which correspond to a basis of
$H^0(B,\mathcal O_B(1))$.  The Weierstrass equation
\begin{equation}
 y^2z=x^3+axz^2+bz^3
 \hypertarget{tag-029F}{}\label{maps:eq:E.1}
\end{equation}
in the projective bundle associated to
$\mathcal O_X\oplus L^2\oplus L^3$ defines a proper complex space
$\mathscr W\to X$.  Its section at infinity is disjoint from its
singular locus.  Its $j$-function is $1728t$.

\item Over $B^\circ$, equation \eqref{maps:eq:E.1} is the product with Engel's elliptic
curve: both have that marked Weierstrass family.  At $p_1$, pass to the
root atlas $t=s^3$.  The coefficient $a$ has a simple zero in $s$
and $b$ is a unit.  At $p_2$, on $t-1=s^4$, $a$ is a unit and $b$ has
order two.  These are the smooth good-reduction families $E$.
The deck action on $F\times E$ is diagonal, with the $A_j\oplus C_j$ linear part and translation $(v_j,0)/m_j$.
It is free because its projection to $F$ is free.  Its quotient is
therefore smooth and is exactly the local model asserted in the theorem.
The section at infinity is invariant.  Notice that the full fibres here
have multiplicities three and four; only their reductions are smooth.

\item At the cusp, use the angular--deck marking of the toroidal model, and write the two-torus period as
$Z_2=sB_0+C(t_c)$, where
\[
 B_0=\begin{pmatrix}0&1\\-1&0\end{pmatrix},\quad
 P=\begin{pmatrix}1&-1&0\\0&1&-1\end{pmatrix},\quad
 Q=B_0^{-1}P,\quad
 S=\begin{pmatrix}1&1\\0&1\\0&0\end{pmatrix}.
\]
Here $PS=I_2$.  Put $k=(1,1,1)^t$, $r=(0,0,1)$, and write the Tate
period of the elliptic surface as $s+c_E(t_c)$, with $c_E$ holomorphic
and bounded.  Define
\begin{equation}
 Z_3=sI_3+C_3(t_c),\qquad
 C_3=SCQ+kc_Er.
 \hypertarget{tag-029O}{}\label{maps:eq:E.2}
\end{equation}
For $A=\binom{P}{r}$ and $B=\binom{Q}{r}$ we have
\begin{equation}
 \det A=\det B=1,\qquad
 AZ_3B^{-1}=\operatorname{diag}(Z_2,s+c_E).
 \hypertarget{tag-029Q}{}\label{maps:eq:E.3}
\end{equation}
This is an integral holomorphic product identification on the punctured overlap, including the bounded correction.

\item Apply \cref{night:P1.9} to the bounded holomorphic correction $C_3$ in \eqref{maps:eq:E.2}.
The fan is \eqref{maps:eq:E.4}; its action is \eqref{maps:eq:E.5} with $(q,C)=(t_c,C_3)$.
This identifies the quotient $N_3^C$ by the identity on its dense torus and gives a smooth proper filling.

\item The map $(P,r)$ takes \eqref{maps:eq:E.4} to a subdivision of the product of an
$A_2$ triangle and an interval.  There are three tetrahedra in each
triangular prism.  No ray is added.  Consequently the induced map is
a small toric resolution of the local equation
$xy=w_0w_1w_2$ in \eqref{maps:eq:E.1}; over a double curve it resolves
$xy=w_0w_1$ times a smooth factor.  All rays have height one, so the
resolution is crepant.  The formula $PZ_3=Z_2Q$ proves equivariance
of its map to the cusp of $X$.  The slice $x_3=1$ gives the zero
section: its lattice map is $S$, and $S$ takes each $A_2$ triangle
to a face of a Kuhn tetrahedron.  Thus the section extends over every
double curve and both triple points.

\item We glue $N_3^C$ and the two finite models to the punctured product by
\eqref{maps:eq:E.3} and their root-atlas identifications.  These overlap maps are
holomorphic isomorphisms and satisfy the cocycle condition because
they are identifications with the same marked product.  Properness
over the three closed discs and compactness over the remaining compact
base give compactness of the glued Hausdorff space.  It is smooth in all its charts.  The maps to $X$ and the zero section
agree on their overlaps.  Since both this model and the literal
fibre product agree on a dense open set, they are bimeromorphic.
\end{enumerate}
\end{proof}



\begin{theorem}[Fundamental group and rational homology]\hypertarget{tag-02M7}{}\label{maps:src:E.1.4}
\StatementRecord{maps:src:E.1.4}{proved}{02M7}{\cref{maps:src:E.1.2}, \cref{maps:src:E.1.3}, \eqref{maps:eq:E.10}, \eqref{maps:eq:E.11}}


The fourfold in Theorem \ref{maps:src:E.1.2} has
\begin{equation}
 \pi_1(Z)=1,\qquad \chi(Z)=6,\qquad
 (b_0,\ldots,b_8)=(1,0,2,1,2,1,2,0,1).
 \hypertarget{tag-029V}{}\label{maps:eq:E.8}
\end{equation}
The first homology groups of its two finite reduced fibres are,
respectively, $\mathbb Z^2\oplus\mathbb Z/3$ and
$\mathbb Z^2\oplus\mathbb Z/2$.

\end{theorem}


\begin{proof}
\begin{enumerate}[label=\arabic*.,leftmargin=*]
\item The section of $\Pi$ splits the fundamental-group sequence over its
smooth locus $X\setminus W$.  Every loop in $Z$ can be moved off
the discriminant inverse image and is a product of a loop in the
elliptic fibre and a section loop.  A section loop becomes trivial
in $Z$ because it bounds in the embedded simply connected $X$.
Thus $\pi_1(Z)$ is a quotient of the coinvariants of
$\mathbb Z^2$ under $C_1,C_2$.  Loops realizing these two
monodromies are the root-circle paths closed by the affine
identifications in $X\setminus W$.  The Smith form of
$(C_1-I\mid C_2-I)$ is $(1,1)$, so these coinvariants vanish.

\item Euler characteristic localizes on the exceptional fibres: all
regular torus fibres and the finite free torus quotients have Euler
characteristic zero.  The cusp has six zero-dimensional torus
strata, and every positive-dimensional torus stratum has Euler
characteristic zero.  Hence $\chi(Z)=6$.

\item Write $V=R^1g_*\mathbb Q$ on the punctured base, and $j$ for its
inclusion.  Finite quotient specialization is an isomorphism onto
invariants over $\mathbb Q$.  Theorem \ref{maps:src:E.1.3} gives
$R^qg_*\mathbb Q=j_*\bigwedge^qV$ except for one kernel
$\mathbb Q_{p_0}$ in degree four.  For any local system $L'$ on
the punctured sphere,
\begin{equation}
 H^2(B,j_*L')=(L')_{\pi_1},\qquad
 \chi(B,j_*L')=-\operatorname{rank}L'+
       \sum_{j=0}^2\dim (L')^{M_j}.
 \hypertarget{tag-029W}{}\label{maps:eq:E.9}
\end{equation}
The first equality is compactly supported duality and the second is
Euler additivity on the open stratum and the three point stalks.
The exact exterior-matrix computation gives the following rows of
the rational Leray $E_2$ page:
\begin{equation}
\begin{array}{c|rrrrrrr}
q&0&1&2&3&4&5&6\\ \hline
\dim E_2^{0,q}&1&1&2&3&3&1&1\\
\dim E_2^{1,q}&0&1&0&1&0&1&0\\
\dim E_2^{2,q}&1&1&2&3&2&1&1.
\end{array}
\hypertarget{tag-029G}{}\label{maps:eq:E.10}
\end{equation}
  The differential ranks
$r_q=\operatorname{rank}(E_2^{0,q}\to E_2^{2,q-1})$ are
\begin{equation}
 (r_1,\ldots,r_6)=(1,1,2,2,1,1).
 \hypertarget{tag-029H}{}\label{maps:eq:E.11}
\end{equation}
Indeed $\Pi^*$ transports the four nonzero transgressions of the
sphere fibration to the corresponding direct summands here.
These four transgressions are forced over $\mathbb Q$ by the
one-dimensional rows and $H^*(X;\mathbb Q)=H^*(S^6;\mathbb Q)$;
this also checks the particular reduction of Engel's Lemma 7.6
which we use.  The divisor class of the zero section gives a
permanent class $u$ restricting to the elliptic orientation class.
Multiplication by $u$ gives another independent image in degree
three, and proves $r_3=2$, $r_5=1$, $r_6=1$, and $r_4\geq2$.
The target summands are independent by the marked decomposition
$\Lambda_X\oplus\Lambda_E$.  The resulting formulas for $b_3$
and $b_5$ are $b_3=1$ and $b_5=3-r_4$.  Poincar\'e duality of
the compact smooth fourfold gives $r_4=2$.  No later differential
exists because the base has real dimension two.  Taking diagonal
sums of \eqref{maps:eq:E.10} with \eqref{maps:eq:E.11} proves \eqref{maps:eq:E.8}.

\item Finally, for a finite quotient its affine fundamental group has
generators $\Lambda\oplus\mathbb Z a$ and abelianized relations
$(A_j\oplus C_j-I)\Lambda$ and $m_ja-(v_j,0)$. The complete relation matrix is obtained by stacking these integral relations; see \hyperlink{prog:V:resolved}{Table~\ref*{tab:V:programs}}.
Its nonzero Smith factors are $(1,1,1,1,3)$ and
$(1,1,1,1,2)$, respectively.  This proves the integral assertions
about the finite fibres.
\end{enumerate}
\end{proof}

\begin{theorem}[Integral homology of the fourfold]\hypertarget{tag-02M8}{}\label{maps:src:E.1.5}
\StatementRecord{maps:src:E.1.5}{proved}{02M8}{\cref{maps:src:E.1.2}, \cref{maps:src:E.1.4}, \eqref{maps:eq:E.I1}, \cref{maps:src:E.1.3}, \SeriesRef{maps:src:D.1.4}, \eqref{maps:eq:E.I2}, \eqref{maps:eq:E.I3}, \eqref{maps:eq:E.I5}, \eqref{maps:eq:E.I4}}


The fourfold of Theorem \ref{maps:src:E.1.2} has torsion-free integral homology:
\[
 H_i(Z;\mathbb Z)=\mathbb Z^{b_i},\qquad
 (b_0,\ldots,b_8)=(1,0,2,1,2,1,2,0,1).
\]
In particular the finite-fibre torsion in \ref{maps:src:E.1.4} does not survive
as torsion in the homology of the total space.

\end{theorem}


\begin{proof}
\begin{enumerate}[label=\arabic*.,leftmargin=*]
\item We compute the integral covering cokernels. Write
\(\mathcal V=R^1g_*\mathbb Z|_{B^\circ}\), and put
\(\mathcal J_q=j_*\bigwedge^q\mathcal V\).
The cohomology monodromies are the inverse transposes of
\(A_1\oplus C_1,A_2\oplus C_2,A_0\oplus C_0\).
All lattices below are in this marking.

\item First we describe a finite filling. Since the first coordinate of each twist
\(v_j\) is \(\pm1\), the integral change of basis with first
column \(v_j/(v_j)_\gamma\) and remaining columns the last five
standard basis vectors changes its linear monodromy to
\(\operatorname{diag}(1,D_j)\) and its translation to
\(\pm e_0/m_j\). The reduced finite fibre is the mapping torus
of the linear automorphism \(D_j:T^5\to T^5\).
The covering \(T^6\to B_j\) is the \(m_j\)-fold cover of its
suspension circle. Set \(T_j=D_j^{-t}\). The two vertical matrices
on homology, in coordinates \((u,w,\delta,e_1,e_2)\), are
\[
D_1=\begin{pmatrix}
0&1&0&0&0\\-1&-1&0&0&0\\1&0&1&0&0\\
0&0&0&-1&-1\\0&0&0&1&0
\end{pmatrix},\qquad
D_2=\begin{pmatrix}
0&-1&0&0&0\\1&0&0&0&0\\0&1&1&0&0\\
0&0&0&0&1\\0&0&0&-1&0
\end{pmatrix}.
\]
The Wang sequences of the mapping torus and its cover give the
following exact description of the cokernel of covering pullback:
\begin{equation}
 C_j^q\simeq
 \frac{\ker(\bigwedge^{q-1}T_j-I)}
 {\operatorname{im}\mathcal N_j^{q-1}},\qquad
 \mathcal N_j^{q-1}=\sum_{a=0}^{m_j-1}
                  (\bigwedge^{q-1}T_j)^a .
 \hypertarget{tag-02B5}{}\label{maps:eq:E.I1}
\end{equation}
Here the target of pullback is the invariant lattice in
\(H^q(T^6;\mathbb Z)\). To verify \eqref{maps:eq:E.I1}, its restriction to
the vertical torus surjects onto
\(\ker(\bigwedge^qT_j-I)\), and restriction of a mapping-torus
class has exactly the same image, by Wang exactness. Thus every
cokernel class has a representative containing the suspension
covector. On this kernel the covering map is the displayed norm:
cut the covered circle into its \(m_j\) consecutive intervals
and transport the vertical cochain through them. This proves both
surjectivity and the precise kernel of the horizontal description.


\item For \(q=2\), the two quotients in \eqref{maps:eq:E.I1} are respectively zero
and \(\mathbb Z/2\). For \(q=3\), they are
\begin{equation}
 C_1^3=\mathbb Z/3,\qquad
 C_2^3=(\mathbb Z/2)^2\oplus\mathbb Z/4.
 \hypertarget{tag-02B6}{}\label{maps:eq:E.I2}
\end{equation}
The same Wang sequences show that the kernels of finite
specialization in degree two are \(\mathbb Z/3\) and
\(\mathbb Z/2\), and that the degree-three and degree-four
kernels vanish. The cusp specialization is integrally onto,
with the sole free degree-four kernel of \ref{maps:src:E.1.3}.

\item We give the integral global maps needed to use these local groups.
The global invariant two-form
\(q_X=u^*w^*+6\gamma\delta^*\) maps onto \(C_2^2\).
Indeed, if its pullback descended to the six-dimensional finite
fibre, restriction to the elliptic zero section would make it
descend to the original bielliptic surface \(S_2\). Its square
there would be \(\pm12/4=\pm3\), which is impossible in the
even intersection form of \(S_2\). That form is even because
\(H^2(S_2;\mathbb Z)\) is torsion free and its canonical class
is torsion, as checked in \SeriesRef{maps:src:D.1.4}. Thus the global-to-local map in
degree two is onto its only nonzero covering cokernel.

\item For degree three all three global invariant classes contain
\(\gamma\), and remain purely horizontal in both adapted
finite markings. They are
\begin{equation}
 \gamma e_1^*e_2^*,\qquad
 \gamma(-u^*e_1^*+w^*e_2^*),\qquad
 \gamma u^*w^* .
 \hypertarget{tag-02B7}{}\label{maps:eq:E.I3}
\end{equation}
Juxtaposition denotes exterior product. Consequently their images
in \eqref{maps:eq:E.I2} are exactly the vertical two-forms in \eqref{maps:eq:E.I3} modulo the
norm lattices.
In integral invariant bases,
the two norm matrices and this common image matrix are
\[
\begin{split}
N_3&=\begin{pmatrix}
0&0&0&0&0&0&0&0&0&3\\
0&0&-2&1&0&1&1&1&-1&0\\
0&0&-1&2&0&2&-1&0&-1&0\\
3&-1&0&0&2&0&0&0&0&0
\end{pmatrix},\\
N_4&=\begin{pmatrix}
0&0&0&0&0&0&0&0&0&4\\
0&0&-2&0&0&0&2&1&-1&0\\
0&0&0&2&0&2&0&-1&-1&0\\
4&-2&0&0&2&0&0&0&0&0
\end{pmatrix},\qquad
G=\begin{pmatrix}1&0&0\\0&1&0\\0&0&0\\0&0&1\end{pmatrix}.
\end{split}
\]
Every nonzero Smith factor of
\(\bigl(\operatorname{diag}(N_3,N_4)\mid(G,G)^t\bigr)\)
is one; there are eight. Thus the global classes surject onto
the direct sum of the two degree-three covering cokernels.

\item We also compute the integral parabolic groups, with their torsion.
The degree-two global coboundary matrix, obtained by stacking
the two monodromy-minus-identity matrices, has a single nontrivial
Smith factor six. Its torsion generator restricts with orders
three, two and one at the three boundary circles. Hence no
nonzero torsion class is parabolic. Its rational parabolic rank
is zero, so \(H^1(B,\mathcal J_2)=0\). In degree three the
global coboundary image is primitive and the rational parabolic
rank is one; hence \(H^1(B,\mathcal J_3)=\mathbb Z\).
These assertions use the cocycle description as the kernel of
restriction from the free-group cohomology to the three
peripheral circle groups. The matrices and restriction orders are obtained by integer kernels and Smith reduction; see \hyperlink{prog:V:resolved}{Table~\ref*{tab:V:programs}}.

\item Passing from \(\mathcal J_q\) to \(R^qg_*\mathbb Z\)
adds the point-supported kernels just described and replaces
the finite stalks by the pullback sublattices. A point-supported
sheaf has no positive cohomology. The long exact sequences,
together with the two surjections onto covering cokernels,
therefore give
\begin{equation}
 H^1(B,R^2g_*\mathbb Z)=0,\qquad
 H^1(B,R^3g_*\mathbb Z)=\mathbb Z.
 \hypertarget{tag-02B8}{}\label{maps:eq:E.I4}
\end{equation}
They also identify degree-two base cohomology with the global
coinvariants. In the relevant degrees these are torsion-free,
with
\begin{equation}
 H^2(B,R^1g_*\mathbb Z)=\mathbb Z,
 \qquad H^2(B,R^2g_*\mathbb Z)=\mathbb Z^2.
 \hypertarget{tag-02B9}{}\label{maps:eq:E.I5}
\end{equation}
The first is the corresponding group from \(X\). The second
is the direct sum of that from \(\bigwedge^2\mathcal V_X\)
and the elliptic orientation line; the mixed coinvariants are
zero. Smith forms show that these identifications are primitive.

\item The integral differential onto the first group in \eqref{maps:eq:E.I5} is
surjective by naturality from the sphere fibration: its own
corresponding Leray quotient must vanish because \(H^3(X)=0\).
The differential onto the second group is surjective as well.
The \(X\) summand is killed by the differential transported
from \(H^4(X)=0\); the elliptic orientation summand is killed
by multiplying the primitive degree-one transgression of \(X\)
by the permanent class of the elliptic zero-section divisor.
That divisor has degree one on each elliptic fibre. Thus these are two primitive independent target generators.

\item Now \(H^0(B,R^3g_*\mathbb Z)\) and
\(H^0(B,R^4g_*\mathbb Z)\) are free: in the first degree
specialization has no kernel, and in the second its only
kernel is the free cusp class. There is no later Leray
differential, since the base has cohomological dimension two.
Equation \eqref{maps:eq:E.I4} and the first surjectivity in \eqref{maps:eq:E.I5} show that
\(H^3(Z;\mathbb Z)\) is a free group, of rank one by \ref{maps:src:E.1.4}.
Equation \eqref{maps:eq:E.I4} and the second surjectivity show that
\(H^4(Z;\mathbb Z)\) is an extension of two free groups of
rank one, and hence is \(\mathbb Z^2\).
Simple connectivity makes \(H^2(Z;\mathbb Z)\) torsion free.
The universal coefficient theorem now proves successively that
\(H_2(Z;\mathbb Z)=\mathbb Z^2\) and
\(H_3(Z;\mathbb Z)=\mathbb Z\). Poincar\'e duality determines
the other degrees and gives exactly the displayed groups.
\end{enumerate}
\end{proof}


\begin{theorem}[Complex invariants and middle form]\hypertarget{tag-02M9}{}\label{maps:src:E.2.1}
\StatementRecord{maps:src:E.2.1}{proved}{02M9}{\cref{maps:src:E.1.2}, \eqref{maps:eq:E.14}, \eqref{maps:eq:E.13}, \eqref{maps:eq:E.1}, \eqref{maps:eq:E.8}; \cite[Theorem 1.1 and Remark 1.2]{ffJiaMeng}; \cite[(1.2)--(1.4)]{ffEsoleKang}; \cite[Theorem 2.2.1]{ffPetersSterk}}


For the fourfold of Theorem \ref{maps:src:E.1.2},
\begin{equation}
 K_Z=\Pi^*L^{-5},\qquad a(Z)=2,\qquad
 (h^{0,0},\ldots,h^{0,4})=(1,1,0,0,0).
 \hypertarget{tag-029I}{}\label{maps:eq:E.12}
\end{equation}
It is not in Fujiki's class $\mathcal C$.  Moreover
$c_1(Z)=0$ integrally, $\int_Zc_2^2=2$, $\sigma(Z)=2$, and
the intersection form on $H^4(Z;\mathbb Z)/\mathrm{torsion}$ is
$\langle1\rangle\oplus\langle1\rangle$.

\end{theorem}


\begin{proof}
\begin{enumerate}[label=\arabic*.,leftmargin=*]
\item We first compute the direct images of the structure sheaf.  The orbifold base has roots of orders three and four.
Write its line generator as $\mathcal O(1)$; it pulls back to $L$.
For the coarse map $\rho$,
\begin{equation}
 \rho_*\mathcal O(k)=\mathcal O_B(d(k)),\qquad
 d(k)=\frac{k}{12}-\frac{k\bmod3}{3}
                    -\frac{-k\bmod4}{4}.
 \hypertarget{tag-029J}{}\label{maps:eq:E.13}
\end{equation}
This follows by dividing a local invariant section by its necessary
root-coordinate power.  Finite-group invariants over $\mathbb C$
are exact.  In particular $\mathcal O(-1)$ and $\mathcal O(-2)$
are acyclic, since their coarse degree is $-1$.

\item For $\bar f:X\to\mathfrak B$, the direct images are
\begin{equation}
 R^0\bar f_*\mathcal O_X=\mathcal O,\quad
 R^1\bar f_*\mathcal O_X=\mathcal O\oplus\mathcal O(-1),
 \quad R^2\bar f_*\mathcal O_X=\mathcal O(-1).
 \hypertarget{tag-029K}{}\label{maps:eq:E.14}
\end{equation}
We verify \eqref{maps:eq:E.14}.
The cusp branch resolution of $\mathcal O_W$ is the cellular
complex of its dual two-torus, and has cohomology dimensions
$(1,2,1)$ with nonzero exterior cup pairing.  Proper flat base
change therefore gives vector bundles with those ranks.  The
relative logarithmic vertical bundle has a fixed line generated
by $\widehat\delta$, and its determinant has twelfth power
$\mathcal O(-p_0)$, by the modular discriminant.  Thus the
determinant is $\mathcal O(-1)$.  Relative duality gives the last
term of \eqref{maps:eq:E.14}.  The invariant covector $\gamma$ gives a nowhere
zero section of $R^1\bar f_*\mathcal O_X$: on a torus its class
in $\Lambda^*_{\mathbb C}/\operatorname{rowspan}\Pi_X$ is
nonzero by the last two columns of $\Pi_X$; at the cusp it is
the nonzero deck cocycle on the dual two-torus.  Its determinant
and the nonzero cusp cup pairing then give the extension
$0\to\mathcal O\to R^1\bar f_*\mathcal O_X
\to\mathcal O(-1)\to0$.  It splits because
$H^1(\mathfrak B,\mathcal O(1))=0$, also from \eqref{maps:eq:E.13}.
The same logarithmic determinant gives $K_X=L^{-6}$.
These computations use the root atlases and the reduced cusp.

\item The Weierstrass adjunction calculation gives
$K_{\mathscr W}=\Pi^*(K_X\otimes L)$, and the small toric
resolution is crepant.  This proves the canonical formula in
\eqref{maps:eq:E.12}.  Every fibre of $\Pi$ is a connected genus-one curve;
over the deeper strata its toric resolution is a cycle of rational
curves.  Thus $\Pi$ is flat, by equidimensionality between smooth
spaces, and relative duality gives
$R^1\Pi_*\mathcal O_Z=L^{-1}$ and
$R^0\Pi_*\mathcal O_Z=\mathcal O_X$.
Restriction to the section splits the unit map in the derived
direct image, so
\[
 R\Pi_*\mathcal O_Z\simeq
 \mathcal O_X\oplus L^{-1}[-1].
\]
Projection formula, \eqref{maps:eq:E.13}, and \eqref{maps:eq:E.14} show that
$H^*(X,L^{-1})=0$ and $h^*(X,\mathcal O_X)=(1,1,0,0)$.
This proves the asserted Hodge numbers.

\item A meromorphic trivialization of $L$ makes the coordinate $x$ in
\eqref{maps:eq:E.1} a meromorphic function over $X$, of degree two on each
elliptic fibre.  It gives a generically finite meromorphic map to
$X\times\mathbb P^1$.  Traces and elementary symmetric functions
over its finite generic fibre show that the meromorphic function
field is an algebraic extension of that of the target.  Since
$a(X)=1$, this proves $a(Z)=2$.  If $Z$ were in class
$\mathcal C$, a sequence of blowups along smooth centres would
give a compact K\"ahler model, by the form of the model theorem
in \cite[Theorem 1.1 and Remark 1.2]{ffJiaMeng}.
Both $b_1$ and $h^{0,1}$ are unchanged by such blowups: the
exceptional fibres have no first homology and no positive-degree
structure-sheaf cohomology.  The K\"ahler identity
$b_1=2h^{0,1}$ would give $0=2$, a contradiction.

\item Since $H^2(X;\mathbb Z)=0$, we have $c_1(L)=0$.
The canonical formula gives $c_1(Z)=5\Pi^*c_1(L)=0$.
Riemann--Roch and the signature
formula, in the explicit normalization of
\cite[(1.2)--(1.4)]{ffEsoleKang}, yield
\[
 0=\chi(\mathcal O_Z)=\frac{3c_2^2-c_4}{720},\qquad
 \sigma(Z)=\frac{3c_2^2+14c_4}{45}.
\]
Here $c_4[Z]=6$, hence $c_2^2[Z]=2$ and $\sigma(Z)=2$.
The integral middle pairing modulo torsion is unimodular by
Poincar\'e duality, has rank two by \eqref{maps:eq:E.8}, and is positive definite
by its signature.  The rank-at-most-four lattice classification
\cite[Theorem 2.2.1]{ffPetersSterk} therefore gives precisely
$\langle1\rangle\oplus\langle1\rangle$.
\end{enumerate}
\end{proof}

\clearpage


\part{Shifted fibre products}\label{pub:V:part-shifted}
\section{The shifted product and its bordism}\label{sec:shifted}

The same sphere fibrations give compact complex fivefolds by placing
their exceptional fibres over disjoint subsets of the base. Their topology
is governed by a torus-knot peripheral group and the integral cusp maps.

\begin{definition}[Shifted fibre product]\hypertarget{tag-PV22}{}\label{def:V:shifted-product}
\StatementRecord{def:V:shifted-product}{proved}{PV22}{}
Let $X$ and $X'$ be canonical marked members of the $(3,4,\infty)$
family, and choose $\varphi\in\operatorname{PGL}_2(\mathbb C)$ so that
the exceptional values of $X$ and $\varphi^*X'$ are disjoint.
Their \emph{shifted fibre product} is $P=X\times_B\varphi^*X'$.
\end{definition}

\begin{theorem}[A shifted second sphere]\hypertarget{tag-02MN}{}\label{maps:src:E.2.3}
\StatementRecord{maps:src:E.2.3}{proved}{02MN}{\SeriesRef{maps:src:E.2.2}, \eqref{maps:eq:E.9}}


Let $X'$ be a second canonical member and choose $\varphi\in
\operatorname{PGL}_2(\mathbb C)$ so that its three exceptional
values are disjoint from those of $X$.  The smooth fivefold
$P=X\times_B\varphi^*X'$ has Euler characteristic zero and
an infinite fundamental group with presentation
\begin{equation}
\begin{split}
G=\langle a_1,a_2,b_1,b_2,h_1,h_2\mid\ &
 h_1,h_2\ \text{central},\quad
 a_1^3=h_1,\ a_2^4=h_1^{-1},\\
 &b_1^3=h_2,\ b_2^4=h_2^{-1},\quad
 a_1a_2b_1b_2=1\rangle .
\end{split}
\hypertarget{tag-029M}{}\label{maps:eq:E.16}
\end{equation}
Its abelianization is $\mathbb Z$, and its rational Betti numbers are
\begin{equation}
 (1,1,1,20,63,88,63,20,1,1,1).
 \hypertarget{tag-029N}{}\label{maps:eq:E.17}
\end{equation}


\end{theorem}


\begin{proof}
\begin{enumerate}[label=\arabic*.,leftmargin=*]
\item Let $f:X\to B$ and $g:\varphi^*X'\to B$ be the two maps.
At a point $(x,x')$ of the fibre product, their disjoint critical-value
sets imply that at least one of $df_x$ and $dg_{x'}$ is surjective.
Consequently the difference differential
$(v,w)\mapsto df_x(v)-dg_{x'}(w)$ maps
$T_xX\oplus T_{x'}\varphi^*X'$ onto $T_{f(x)}B$.
In a local coordinate on $B$ it is the differential of the defining
function $f-g$. The holomorphic implicit function theorem therefore
makes $P$ smooth of complex dimension $3+3-1=5$.
It is compact because it is closed in the compact product
$X\times\varphi^*X'$.  At each special value the other
factor is a smooth four-torus.  The Euler characteristic of every
fibre is therefore zero, and Euler additivity gives $\chi(P)=0$.
We apply van Kampen with the cusp relations.
Over the six-punctured sphere the fibre group is
$\Lambda_X\oplus\Lambda_{X'}$ and the base is free on five
peripheral loops.  Each monodromy acts on its own direct summand.
Filling the first cusp kills its angular generators
$\widehat w,\widehat\delta$ and its cusp meridian.  Their normal
closure under $A_1$ also kills $\widehat u$, since
$A_1\widehat w=\widehat u-\widehat w$.  The same calculation
in the other summand leaves only $h_1=\widehat\gamma$ and
$h_2=\widehat\gamma'$.  The finite monodromies now act trivially
on them.  The four finite filling relations are precisely those
in \eqref{maps:eq:E.16}, since $\gamma(v_1)=1$ and $\gamma(v_2)=-1$.
Both cusp collars are untwisted; their meridians are killed.
The product of the remaining four meridians gives the last
relation.  These are all relations introduced by the six
fillings, so van Kampen gives exactly \eqref{maps:eq:E.16}.

\item Quotienting by $h_1,h_2$ gives the orbifold group of the sphere
with cone orders $3,4,3,4$.  Its orbifold Euler characteristic
is $-5/6$, so its universal orbifold cover is the disc and the
group is infinite.  Equivalently one can build the Seifert
$T^2$-bundle with these four primitive translation vectors; over
the disc it is a trivial torus bundle.  Its exact sequence shows
that the central subgroup in \eqref{maps:eq:E.16} is $\mathbb Z^2$.
The abelianized five-by-six relation matrix has Smith factors
$(1,1,1,1,1)$ by integral row and column operations, proving $G_{\rm ab}=\mathbb Z$.

\item For the rational homology there is no extra cusp kernel: each
cusp is the rank-two cusp of the sphere fibration times a smooth torus, and its
specialization is an isomorphism onto invariants over $\mathbb Q$.
The two monodromy groups act on independent summands.  In exterior
degrees $q=0,\ldots,8$, the $H^0$ and $H^2$ rows are both
\[
 (1,2,3,4,5,4,3,2,1),
\]
and the $H^1$ row is
\[
 (0,0,18,60,86,60,18,0,0).
\]
The local invariant polynomials of the first factor are
$A(z)=1+2z+2z^2+2z^3+z^4$ at a finite point and
$C(z)=A(z)+2z^2$ at its cusp.  Hence the Euler polynomial of
the direct-image rows is
\[
 -4(1+z)^8+(4A(z)+2C(z))(1+z)^4.
\]
Formula \eqref{maps:eq:E.9}, with six point strata, gives the displayed rows.

\item Each invariant algebra of a factor has one generator in degrees
zero through four.  Its invariant-to-coinvariant map is nonzero
in even degrees and zero in odd degrees.  In the marking used
above, this follows from the invariant form
$6\gamma\wedge\delta+u\wedge w$ and its exterior products:
the odd invariant classes pair to zero with every invariant
complementary class, and the even ones have nonzero square or
top pairing.  Its four transgressions are all nonzero, as the
factor is a sphere.  The derivation rule therefore sends a tensor
of bidegree $(i,j)$ into bidegrees $(i-1,j)$ and $(i,j-1)$,
with the first term nonzero exactly when $i>0$ and $j$ is even,
and the second exactly when $j>0$ and $i$ is even.  The nonzero
column supports are disjoint except for pairs in a single row,
so their scalar normalizations do not affect ranks.  This gives
\[
 (r_1,\ldots,r_8)=(1,2,2,3,3,2,2,1).
\]
Row reduction of these exterior-product matrices gives the displayed ranks.
There are no later differentials.  Their diagonal sums give
\eqref{maps:eq:E.17}, which also satisfies Poincar\'e duality and Euler zero.

\end{enumerate}
\end{proof}





\begin{proposition}[The first integral torsion of the shifted product]\hypertarget{tag-02MP}{}\label{maps:src:E.2.5}
\StatementRecord{maps:src:E.2.5}{proved}{02MP}{\cref{maps:src:E.2.3}, \cref{pub:V:cusp-chern-generators}, \SeriesRef{maps:src:D.1.2}, \eqref{maps:eq:E.I1}, \cref{maps:src:E.1.5}, \eqref{maps:eq:E.20}}


Let \(Z_\varphi=X\times_B\varphi^*X\) be the transverse
fivefold in \ref{maps:src:E.2.3}, with the two sets of three critical values
disjoint. Its second integral homology is
\[
 H_2(Z_\varphi;\mathbb Z)
   =\mathbb Z\oplus(\mathbb Z/12)^7.
\]
\end{proposition}


\begin{proof}
\begin{enumerate}[label=\arabic*.,leftmargin=*]
\item Use subscripts one and two for the two rank-four cohomology
local systems. At a critical value of one factor the other
factor is a smooth torus bundle over a disc, hence a product
there. Each local specialization is therefore the tensor
product of the original specialization with the identity on
the exterior algebra of \(\mathbb Z^4\). The original finite
reduced surfaces and the original cusp have free cohomology.
Their specialization maps are injective; at the cusp they
are integral isomorphisms onto invariants. For the latter,
apply \cref{pub:V:cusp-chern-generators} in dimension two. It supplies the invariant deck classes
\(b_1,b_2\), the three classes
\(a_1b_1,a_2b_2,a_1b_2+a_2b_1\), and their products.
Together with \(b_1b_2\) these give an integral kernel
basis in degree two; their products give integral bases
in degrees three and four. Both the central-fibre and
invariant lattices have ranks \((1,2,4,2,1)\), as follows
also from the normalization matrix in \SeriesRef{maps:src:D.1.2} at \(d=1\).
Thus the specialization is onto between free groups of
equal ranks in every degree, proving the claimed isomorphism.
Thus the only difference between \(R^qg_*\mathbb Z\) and
\(\mathcal J_q=j_*\bigwedge^q(\mathcal V_1\oplus\mathcal V_2)\)
is their finite covering cokernel, with no point-supported kernel.

\item The finite Wang calculation \eqref{maps:eq:E.I1}, now with the original
three-dimensional vertical lattice, gives
\[
 C_3^1=\mathbb Z/3,\quad C_4^1=\mathbb Z/4,
 \qquad C_3^2=0,\quad C_4^2=\mathbb Z/2.
\]
The invariant covector \(\gamma\) generates the first two
cokernels, and the invariant two-form \(q_X\) generates the
last, by the same even-intersection argument as in \ref{maps:src:E.1.5}.
Consequently the degree-two covering cokernel at all four
finite values, combined, is
\[
 C^2=\left((\mathbb Z/3)^4\oplus
                 (\mathbb Z/4)^4\oplus\mathbb Z/2\right)^2.
\]
The two monodromy groups act independently on the two factors.
The three global invariant two-forms have integral basis
\(q_1,\gamma_1\gamma_2,q_2\). The first and third map
onto the respective \(\mathbb Z/2\) summands and vanish
in the other summands. The middle form maps to the primitive
\(\gamma\) coordinate of each of the four order-three or
order-four blocks. After the two order-two summands have
been removed, it therefore generates a direct cyclic summand
of order twelve in \((\mathbb Z/12)^8\). Hence
\begin{equation}
 \operatorname{coker}\{H^0(B,\mathcal J_2)\to C^2\}
      =(\mathbb Z/12)^7.
 \hypertarget{tag-029P}{}\label{maps:eq:E.20}
\end{equation}
The eighteen-row integer presentation of these covering cokernels has
Smith form with eleven unit factors and seven factors twelve.

\item In degree two the global coboundary image is primitive.
Thus the parabolic group, a subgroup of the free-group
cohomology on the punctured base, is torsion free. Its rank
is eighteen: the six local invariant ranks are
\((16,16,18,16,16,18)\), and its global invariant and
coinvariant ranks are both three. The Euler formula on the
six-punctured sphere gives \(3+3+4\cdot28-100=18\).
These ranks and the primitivity are computed with exact exterior matrices. The sheaf sequence
\(0\to R^2g_*\mathbb Z\to\mathcal J_2\to C^2\to0\)
and \eqref{maps:eq:E.20} now give
\[
 H^1(B,R^2g_*\mathbb Z)
     =\mathbb Z^{18}\oplus(\mathbb Z/12)^7.
\]
Here the quotient by its finite subgroup is free, so this
abelian-group extension splits.

\item The group \(H^2(B,R^1g_*\mathbb Z)\) is the free
rank-two group of coinvariants, one primitive summand from
each factor. Pulling back the integral degree-two
transgression of the original sphere fibration through each
projection kills these two summands separately and
surjectively. Thus \(E_\infty^{2,1}=0\) in the Leray
sequence of \(g\). Its other middle group
\(E_\infty^{0,3}\) is free, since \(R^3g_*\mathbb Z\)
injects into \(\mathcal J_3\). There are no higher
differentials. Consequently
\(\operatorname{Tor}H^3(Z_\varphi;\mathbb Z)
=(\mathbb Z/12)^7\). The universal coefficient theorem
identifies this with \(\operatorname{Tor}H_2(Z_\varphi;\mathbb Z)\).
Finally \ref{maps:src:E.2.3} gives \(b_2=1\), proving the assertion.
\end{enumerate}
The corresponding exact calculations are implemented in \hyperlink{prog:V:shifted-homology}{the programs in Table~\ref*{tab:V:programs}}.
\end{proof}


\begin{theorem}[Full integral homology of the shifted product]\hypertarget{tag-02MQ}{}\label{maps:src:E.2.6}
\StatementRecord{maps:src:E.2.6}{proved}{02MQ}{\cref{maps:src:E.2.3}, \eqref{maps:eq:E.MV1}, \cref{maps:src:E.2.5}, \eqref{maps:eq:E.MV3}, \eqref{maps:eq:E.MV2}, \eqref{maps:eq:E.MV5}, \eqref{maps:eq:E.MV6}, \eqref{maps:eq:E.MV7}, \eqref{maps:eq:E.MV8}}


For the transverse shifted product $P$ of \ref{maps:src:E.2.3}, write
\[
 A=(\mathbb Z/12)^7,\qquad
 B=(\mathbb Z/12)^{10}\oplus(\mathbb Z/2)^7,\qquad
 C=(\mathbb Z/12)^6\oplus\mathbb Z/24\oplus(\mathbb Z/2)^{17}.
\]
Its complete integral homology is
\[
 \begin{array}{c|l}
 i&H_i(P;\mathbb Z)\\\hline
 0,1,8,9,10&\mathbb Z\\
 2&\mathbb Z\oplus A\\
 3&\mathbb Z^{20}\oplus B\\
 4&\mathbb Z^{63}\oplus C\\
 5&\mathbb Z^{88}\oplus C\\
 6&\mathbb Z^{63}\oplus B\\
 7&\mathbb Z^{20}\oplus A.
 \end{array}
\]
In particular the higher integral filtration extensions are nontrivial:
the middle torsion includes a cyclic factor of order $24$.
\end{theorem}


\begin{proof}
\begin{enumerate}[label=\arabic*.,leftmargin=*]
\item \emph{A decomposition with free groups on every piece.}
Choose a separating circle in $B=\mathbb P^1$ so that one closed
disc contains the three critical values of $X$ and the other
contains the three critical values of $X'$. The smooth torus
family of the other factor is trivial over each disc. Set
\[
 F=T^4,\qquad E=X\setminus\operatorname{int}(F\times D^2),
 \qquad E'=X'\setminus\operatorname{int}(F'\times D^2).
\]
The product trivializations give the decomposition
\begin{equation}
 P=(E\times F')\ \cup_{F\times F'\times S^1}\ (F\times E').
 \hypertarget{tag-02BA}{}\label{maps:eq:E.MV1}
\end{equation}
The boundary identification preserves the two marked torus factors
and reverses the normal circle, because the two base discs have
opposite boundary orientations. The same regular-disc trivializations define both the
boundary identifications of $E,E'$ and \eqref{maps:eq:E.MV1}.

\item Since $X,X'$ are six-spheres, the Mayer--Vietoris sequence of
$S^6=E\cup(F\times D^2)$ gives, integrally,
\begin{equation}
 H_0(E)=\mathbb Z,\qquad H_k(E)=H_{k-1}(F)\quad(1\leq k\leq4),
 \qquad H_k(E)=0\quad(k\geq5).
 \hypertarget{tag-02BB}{}\label{maps:eq:E.MV2}
\end{equation}
The isomorphism is the normal-circle sweep
$\alpha\mapsto[S^1]\times\alpha$. Indeed, in degrees one
through four the boundary map to $H_k(E)\oplus H_k(F)$ is
an isomorphism, and its second component is the identity on
the unswept fibre summand and zero on the swept summand.
In degree five the boundary fundamental class is the image
of the sphere fundamental class, so $H_5(E)=0$.
Thus all homology groups of the two pieces and their common
boundary in \eqref{maps:eq:E.MV1} are free abelian. The same statements hold
for $E'$.

\item \emph{Determining every boundary map integrally.}
Let $V=H^1(F;\mathbb Z)$, in the basis $(\gamma,u,w,\delta)$,
and use the positively ordered volume
$v=\gamma u w\delta$. The primitive global invariant classes
of the original torus fibration are
\[
 I_0=1,\quad I_1=\gamma,\quad
 I_2=q=uw+6\gamma\delta,\quad
 I_3=\gamma uw,\quad I_4=v.
\]
The images of $H^k(E)\to H^k(F)$ are exactly
\begin{equation}
 n_k\mathbb Z I_k,\qquad
 (n_1,n_2,n_3,n_4)=(12,2,2,12).
 \hypertarget{tag-02BC}{}\label{maps:eq:E.MV3}
\end{equation}
We verify this integrally. On the disc the only local restrictions on a global
invariant class are the finite-cover pullback lattices;
the cusp specialization is an integral isomorphism onto its
invariant lattice by \ref{maps:src:E.2.5}. In degree one the finite images
contain $3\gamma$ and $4\gamma$, giving the intersection
$12\mathbb Z\gamma$. In degree two the order-three image
contains $q$, while the order-four image requires an even
coefficient of $q$, by the finite integral lattices computed in the proof of \cref{maps:src:E.1.5}.
In degree three use the perfect integral pairing downstairs:
an invariant class is a pullback precisely when its pairings
with every pulled-back degree-one class are divisible by the
covering degree. The other degree-one generators are
$2u+w+3\delta$ and $u+w+2\delta$. Their pairings with
$t\gamma uw$ are, up to signs, $3t$ and $2t$.
Thus the first cover imposes no further condition, and the
second requires $t$ even. In degree four the pullback degrees
are three and four, giving $12\mathbb Z v$.

\item These local restrictions are also sufficient for the image
from $E$. For a local system on a disc with punctures, the
restriction from its free-group $H^1$ to the direct sum of
the local coinvariants is onto: each puncture has its own
independent free generator. Consequently
$H^2(D,j_*\mathcal L)=0$. Replacing $j_*\mathcal L$ by the
direct-image sheaf changes it only by point-supported
cokernels, and so does not change this vanishing.
The Leray sequence over the disc therefore makes
$H^k(E)\to H^0(D,R^kf_*\mathbb Z)$ onto. This proves
\eqref{maps:eq:E.MV3} 

\item Write $\theta$ for the normal-circle covector. The dual of
\eqref{maps:eq:E.MV2} identifies $H^k(E)$ with $\bigwedge^{k-1}V$, and
its boundary restriction has the graph form
\begin{equation}
 r_k(\alpha)=\phi_k(\alpha)+\theta\alpha,
 \qquad
 \phi_k:\bigwedge^{k-1}V\longrightarrow\bigwedge^kV.
 \hypertarget{tag-02BD}{}\label{maps:eq:E.MV4}
\end{equation}
The image of $\phi_k$ is exactly $n_k\mathbb Z I_k$ by
\eqref{maps:eq:E.MV3}. If $k+l=5$, classes extending to $E$ have zero
product evaluated on $\partial E$. Therefore
\begin{equation}
 (-1)^k\phi_k(\alpha)\beta+
                    \alpha\phi_l(\beta)=0
 \quad\text{in }\bigwedge^4V.
 \hypertarget{tag-02BE}{}\label{maps:eq:E.MV5}
\end{equation}
Each $I_k$ is primitive, and wedge pairing between the two
complementary exterior powers is unimodular. Writing
$\phi_k=f_k I_k$, equation \eqref{maps:eq:E.MV5}, tested against a class
pairing to one with $I_k$, proves that $f_k$ is an integral
multiple of $\alpha\mapsto\langle\alpha I_{5-k},[F]\rangle$.
The exact image \eqref{maps:eq:E.MV3} fixes the absolute multiple. Thus
\begin{equation}
 \begin{split}
 \phi_k(\alpha)&=c_k
       \langle\alpha I_{5-k},[F]\rangle I_k,\\
 (c_1,c_2,c_3,c_4)&=(12\varepsilon,2\eta,-2\eta,-12\varepsilon),
 \qquad \varepsilon,\eta\in\{1,-1\}.
 \end{split}
 \hypertarget{tag-02BF}{}\label{maps:eq:E.MV6}
\end{equation}
This proves every boundary map up to two signs. We allow those
signs independently for $E$ and $E'$ below; the answer is the
same for all sixteen possibilities.

\item \emph{The full integral Mayer--Vietoris matrices.}
Let $i_p:H_p(F)\to H_p(E)$ be the transpose of $\phi_p$ in
the sweep bases, with $i_0=1$; let $j_{p+1}$ denote the sweep
isomorphism for $0\leq p\leq3$, and zero for $p=4$.
The analogous maps of the second copy are $i'_q,j'_{q+1}$.
For a fibre tensor $\alpha\otimes\beta$, with degrees $p,q$,
the matrix $D_k$ of \eqref{maps:eq:E.MV1} is
\begin{equation}
 \alpha\otimes\beta\longmapsto
  (i_p\alpha\otimes\beta,-\alpha\otimes i'_q\beta),
 \qquad p+q=k.
 \hypertarget{tag-02BG}{}\label{maps:eq:E.MV7}
\end{equation}
For a swept tensor, place the normal circle first. Then
\begin{equation}
 \theta\times\alpha\times\beta\longmapsto
 \bigl(j_{p+1}\alpha\otimes\beta,
                 (-1)^p\alpha\otimes j'_{q+1}\beta\bigr),
 \qquad p+q=k-1.
 \hypertarget{tag-02BH}{}\label{maps:eq:E.MV8}
\end{equation}
The second sign contains the Koszul sign from moving the circle
past $\alpha$, the reversal of the two boundary-circle
orientations, and the second minus in Mayer--Vietoris. The last
two cancel. This specifies every integer entry of $D_k$.

\item In the exterior bases, exact Smith reduction of the four
adjoint maps \eqref{maps:eq:E.MV6} and the full matrices
\eqref{maps:eq:E.MV7}--\eqref{maps:eq:E.MV8} gives the following table. The source ranks are $\binom9k$; a zero binomial is
understood outside $0\leq k\leq9$.
\[
 \begin{array}{c|rrrr|l}
 k&\text{target rank}&\text{source rank}&\operatorname{rank}D_k&
       \text{unit factors}&\text{nonunit factors}\\\hline
 0&2&1&1&1&--\\
 1&10&9&9&9&--\\
 2&28&36&27&20&12^{[7]}\\
 3&64&84&53&36&2^{[7]},12^{[10]}\\
 4&114&126&82&58&2^{[17]},12^{[6]},24\\
 5&138&126&94&70&2^{[17]},12^{[6]},24\\
 6&104&84&73&56&2^{[7]},12^{[10]}\\
 7&44&36&35&28&12^{[7]}\\
 8&8&9&8&8&--\\
 9&0&1&0&0&--\\
 10&0&0&0&0&--
 \end{array}
\]
The notation $a^{[r]}$ denotes $r$ entries equal to $a$.
All sixteen choices of the two signs in each copy give precisely
this table.

\item Finally the exact homology sequence gives
\[
 0\longrightarrow\operatorname{coker}D_k
 \longrightarrow H_k(P)
 \longrightarrow\ker D_{k-1}\longrightarrow0.
\]
The last group is free, since the boundary homology is free.
This sequence therefore splits as abelian groups, so the Smith
factors in the table are the torsion of $H_k(P)$.
Adding the free kernel ranks gives exactly the rational Betti
vector in \ref{maps:src:E.2.3} and the integral groups asserted here.
The resulting torsion also satisfies the required ten-dimensional
duality $\operatorname{Tor}H_k(P)\cong
\operatorname{Tor}H_{9-k}(P)$.
\end{enumerate}
The corresponding exact calculations are implemented in \hyperlink{prog:V:shifted-homology}{the programs in Table~\ref*{tab:V:programs}}.
\end{proof}

\begin{proposition}[A unique transverse family type and its stable tangent bundle]\hypertarget{tag-02MR}{}\label{maps:src:E.2.7}
\StatementRecord{maps:src:E.2.7}{proved}{02MR}{\cref{maps:src:E.2.6}}


For fixed marked sphere fibrations, all M\"obius transformations
whose critical triples are disjoint give orientation-preservingly
diffeomorphic shifted products. The same conclusion holds while
either sphere varies in its proper smooth marked period family.
For every such fivefold $P$ there is an isomorphism of real vector
bundles
\begin{equation}
 TP\oplus\mathbb R^4\simeq\mathbb R^{14}.
 \hypertarget{tag-02BI}{}\label{maps:eq:E.MV9}
\end{equation}
Thus all its positive-degree Stiefel--Whitney and Pontryagin
classes vanish. For its given complex tangent bundle,
$c_1(P)=c_2(P)=0$ integrally.
\end{proposition}


\begin{proof}
\begin{enumerate}[label=\arabic*.,leftmargin=*]
\item The allowed transformations form the complement in
$\operatorname{PGL}_2(\mathbb C)$ of the nine analytic
hypersurfaces imposing an equality of one critical value from
each triple. This complement is path connected: a piecewise
smooth path between two points can be perturbed relative to
its endpoints to avoid these sets, which have real codimension
two. Equivalently, use local complex coordinate balls and
perturb each segment there. Choose a smooth such path.

\item The corresponding fibre products form a compact smooth family
over its parameter interval. Indeed the vertical derivative of
$f_1(x)-f_2(y)$ is surjective at every point, since its two
factors cannot be critical at the same base value. The same
surjectivity solves the equation for a lift of the interval
coordinate vector field. A partition of unity supplies a
global smooth lift; compactness makes its flow exist for the
whole interval. This flow gives the claimed diffeomorphisms.
The complex orientations vary continuously and hence are
preserved. The identical argument applies to smooth proper
period families of either factor, keeping the critical triples
disjoint throughout.

\item View $P$ as the transverse inverse image of the diagonal in
$B\times B$ under $f_1\times f_2$. Its real oriented normal
two-plane bundle is the underlying bundle of $g^*T_B$,
where $g:P\to B$ is the common base map. It is trivial:
its first Chern class is
\[
 c_1(g^*T_B)=2g^*[B]^\vee
    =2\operatorname{pr}_1^*f_1^*[B]^\vee=0,
\]
because $H^2(X;\mathbb Z)=0$, and complex line bundles are
classified by their first Chern class. The tangent bundles
of the two six-spheres become trivial after adding one real
line each. Restriction of their product tangent bundle gives
\[
 (TP\oplus\mathbb R^2)\oplus\mathbb R^2
       \simeq\mathbb R^{14},
\]
proving \eqref{maps:eq:E.MV9} and the characteristic-class assertions for
the real tangent bundle. Finally the holomorphic tangent
sequence is
\[
 0\longrightarrow TP\longrightarrow
 \operatorname{pr}_1^*TX\oplus\operatorname{pr}_2^*TX'
 \longrightarrow g^*T_B\longrightarrow0.
\]
Both sphere factors have integral $c_1=c_2=0$, since their
cohomology in degrees two and four vanishes. The last line
has zero first Chern class, as just proved. Multiplying total
Chern classes in this exact sequence gives
$c_1(P)=c_2(P)=0$ integrally.
\end{enumerate}
\end{proof}

\begin{proposition}[The shifted product bounds with its natural stable framing]\hypertarget{tag-02MS}{}\label{maps:src:E.2.8}
\StatementRecord{maps:src:E.2.8}{proved}{02MS}{\eqref{maps:eq:E.MV9}, \SeriesRef{maps:src:B.2.1}, \eqref{maps:eq:E.B2}, \cref{maps:src:E.2.7}, \cref{maps:src:E.2.9}; \cite[Chapter 3, Section 2, Theorems 2.1 and 2.5]{rootHirsch76}}


Equip the two oriented six-spheres with their standard bounding
stable tangent framings. Trivialize the oriented two-plane bundle
$f_1^*T\mathbb P^1$ on the first sphere, and equip the transverse
shifted product $P$ with the induced stable framing in \eqref{maps:eq:E.MV9}.
Then there is a compact smooth oriented eleven-manifold $V$ with
\begin{equation}
 \partial V=P,\qquad TV\oplus\mathbb R^3\simeq\mathbb R^{14},
 \qquad [P]=0\quad\hbox{in }\Omega^{\mathrm{fr}}_{10}.
 \hypertarget{tag-029X}{}\label{maps:eq:E.B1}
\end{equation}
The boundary framing is the stated one. In particular its induced
spin structure bounds. 
\end{proposition}


\begin{proof}
\begin{enumerate}[label=\arabic*.,leftmargin=*]
\item Write $f_2:S^6\to\mathbb P^1$ for the shifted map of the second
factor. By \SeriesRef{maps:src:B.2.1} the first map $f_1:S^6\to\mathbb P^1$ is
null-homotopic. It therefore extends to a smooth map
$F:D^7\to\mathbb P^1$, which can be chosen to equal $f_1$ on
each sphere of a collar of the boundary. To see this with the
boundary fixed, first use the null-homotopy to obtain a continuous
extension constant in the collar direction. Smooth its interior
in coordinate charts, with cutoffs supported away from the collar;
on the overlap the extension is already smooth.

\item Let $C_2$ be the three critical values of $f_2$. The restriction
$f_1$ is transverse to each point of $C_2$, since the two critical
triples are disjoint. Perturb $F$ only in the interior to be
transverse to all these points. The applicable finite simultaneous
transversality statement is
\cite[Chapter 3, Section 2, Theorems 2.1 and 2.5]{rootHirsch76}.
For the relative assertion, the proof of that theorem uses local
perturbations multiplied by cutoff functions. Their supports may
be restricted to the complement of the collar, where
transversality has not already been secured; the open condition
on the collar persists. Compactness permits finitely many such
perturbations. 

\item Define the smooth fibre product
\begin{equation}
 V=\{(x,y)\in D^7\times S^6:F(x)=f_2(y)\}.
 \hypertarget{tag-029Y}{}\label{maps:eq:E.B2}
\end{equation}
At a point whose common value is outside $C_2$, the differential
$df_2$ is onto. At a point over $C_2$, the differential $dF$ is
onto. Consequently $dF-df_2$ is onto everywhere on \eqref{maps:eq:E.B2}.
The same argument holds on its boundary, using $df_1$ there.
The implicit-function theorem, including its boundary version,
makes $V$ a smooth eleven-manifold with boundary exactly $P$.
It is compact as a closed subset of $D^7\times S^6$.
The inverse-image orientation using the complex orientation of
$T\mathbb P^1$ induces the complex orientation of $P$ at the
boundary, because the collar keeps the first map fixed.

\item Let $g:V\to\mathbb P^1$ be the common map. The real normal
two-plane bundle of \eqref{maps:eq:E.B2} is $g^*T\mathbb P^1$, identified by
$dF-df_2$. This is the pullback from $D^7$ of $F^*T\mathbb P^1$,
which is trivial because the disc is contractible. A choice of
its oriented framing restricts on the boundary to the chosen
framing of $f_1^*T\mathbb P^1$ up to homotopy: two oriented
framings of this two-plane bundle differ by a map
$S^6\to SO(2)$, and $H^1(S^6;\mathbb Z)=0$.
The tangent exact sequence and the bounding stable framing of
the second sphere now give
\[
 \begin{split}
 TV\oplus\mathbb R^2&\simeq
       T(D^7\times S^6)|_V,\\
 TV\oplus\mathbb R^3&\simeq
       \mathbb R^7\oplus(TS^6\oplus\mathbb R)|_V
       \simeq\mathbb R^{14}.
 \end{split}
\]
At the boundary, the outward normal of $D^7$ supplies the
additional real line in $TV|_P=TP\oplus\mathbb R$.
The resulting framing is precisely the stable tangent framing
obtained from the two sphere framings and the normal framing
in \ref{maps:src:E.2.7}. This proves \eqref{maps:eq:E.B1} by the definition of framed bordism.
It also proves the assertion about the induced spin structure,
by forgetting the framing.

\end{enumerate}
\end{proof}


\begin{theorem}[Framed null-bordism over the fundamental group]\hypertarget{tag-02MT}{}\label{maps:src:E.2.9}
\StatementRecord{maps:src:E.2.9}{proved}{02MT}{\eqref{maps:eq:E.16}, \cref{maps:src:E.2.8}, \SeriesRef{maps:src:B.2.1}, \SeriesRef{maps:src:E.2.2}, \cref{maps:src:E.2.3}, \SeriesRef{maps:src:B.1.2}, \eqref{maps:eq:E.B7}; \cite[Theorems 2.3--2.4 and the Euler-number discussion in
Section 2]{rootScott83}; \cite[Definition 8.16 and Theorem 8.17]{rootDavisKirk01}; \cite[Corollary~1.3 and Table~1, pp.~112--113]{IsaksenWangXu2023}}


Let $G$ be the group in \eqref{maps:eq:E.16}. There is a closed smooth
aspherical four-manifold $N$ with
\begin{equation}
 \pi_1(N)=G,\qquad \widetilde N\simeq\mathbb H^2\times\mathbb R^2,
 \qquad H_*(N;\mathbb Z)\cong H_*(S^1\times S^3;\mathbb Z),
 \hypertarget{tag-029Z}{}\label{maps:eq:E.B3}
\end{equation}
and a smooth classifying map $c:P\to N$ inducing the displayed
identification of fundamental groups. With the natural stable
framing of \ref{maps:src:E.2.8}, one has
\begin{equation}
                 [P,c]=0\quad\hbox{in }\Omega^{\mathrm{fr}}_{10}(N).
 \hypertarget{tag-02A0}{}\label{maps:eq:E.B4}
\end{equation}
Thus $P$ has a compact stably framed eleven-dimensional filling
over $BG$. This is stronger than the filling of \ref{maps:src:E.2.8}.
\end{theorem}


\begin{proof}
\begin{enumerate}[label=\arabic*.,leftmargin=*]
\item \emph{The aspherical target.}
The quotient maps of \SeriesRef{mapsI:src:A.2.1} are smooth maps
$e_i:S^6\to S^3$ over the base, and \SeriesRef{maps:src:B.2.1} gives the exact
factorizations $f_i=h_i e_i$. Here $h_i:S^3\to\mathbb P^1$
is the weighted Hopf map of orders $3,4$, with the second
map shifted by the chosen M\"obius transformation. Define
\begin{equation}
 N=\{(u,v)\in S^3\times S^3:h_1(u)=h_2(v)\},\qquad
 c(x,y)=(e_1(x),e_2(y)).
 \hypertarget{tag-02A1}{}\label{maps:eq:E.B5}
\end{equation}
The critical values of each $h_i$ are only its two exceptional
circle values. They are disjoint, so the same transversality
argument as in \SeriesRef{maps:src:E.2.2} makes $N$ a compact smooth four-manifold.
The two circle actions on the weighted Hopf spheres give a
principal $T^2$ bundle in the orbifold sense over
\[
                  \mathcal B=S^2(3,4,3,4).
\]
At a cone of order $m$, its uniformizing chart is
$(D^2\times T^2)/C_m$, with the generator rotating the disc
and translating one torus coordinate by a primitive fraction
of denominator $m$. This local description is smooth and is
the same finite-filling description as \SeriesRef{mapsI:src:A.1.3}, with the unused circle kept.

\item The orbifold has Euler characteristic
$2-2(1-1/3)-2(1-1/4)=-5/6$. It is good and has hyperbolic
orbifold universal cover, by
\cite[Theorems 2.3--2.4 and the Euler-number discussion in
Section 2]{rootScott83}. Pulling the principal torus bundle to
$\mathbb H^2$ gives an ordinary principal $T^2$ bundle over a
contractible base, hence the trivial bundle. Its universal
cover is $\mathbb H^2\times\mathbb R^2$. The local chart above
shows that these orbifold pullbacks are ordinary covering maps
of total spaces. This proves asphericity and injectivity of
the fibre subgroup $\mathbb Z^2$.

\item Using the marked fibre generators $h_1,h_2$ and the four
positive peripheral generators gives exactly
\[
 \begin{split}
 \pi_1(N)=\langle a_1,a_2,b_1,b_2,h_1,h_2\mid {}&h_1,h_2
       \text{ central},\quad a_1^3=h_1,\quad a_2^4=h_1^{-1},\\
 &b_1^3=h_2,\quad b_2^4=h_2^{-1},\quad
 a_1a_2b_1b_2=1\rangle.
 \end{split}
\]
The map $c$ sends the surviving regular-fibre generators
$\widehat\gamma,\widehat\gamma'$ to $h_1,h_2$, respectively,
and each peripheral generator to the one with the same name.
This is because $e_i$ is the literal circle quotient on these
fibres and finite collars. Its cusp replacement extends over
a solid torus and kills precisely the cusp meridian in this
presentation. Comparing with the full van Kampen calculation
in \ref{maps:src:E.2.3} therefore shows that $c_*$ is the identity on the
displayed presentation, and in particular is an isomorphism.
As $N$ is aspherical, $c$ is a classifying map.

\item For completeness the abelian relation matrix has Smith
diagonal $(1,1,1,1,1)$; a primitive integer kernel vector in
the order $(a_1,a_2,b_1,b_2,h_1,h_2)$ is
$(4,-3,-4,3,12,-12)$. Thus $H_1(N)=\mathbb Z$.
Every stratum of its base has a torus fibre, so Euler
additivity gives $\chi(N)=0$. Integral Poincar\'e duality and
the universal coefficient theorem give $H_3(N)=\mathbb Z$
and no torsion in $H_2(N)$. Euler characteristic then forces
$H_2(N)=0$, proving the integral homology assertion in \eqref{maps:eq:E.B3}.

\item \emph{The bordism computation before taking the inverse image.}
For a based sphere, the stable splitting
$(S^3)_+\simeq S^0\vee S^3$ identifies
\begin{equation}
 \Omega^{\mathrm{fr}}_6(S^3)
       =\pi^{\mathrm S}_6\oplus\pi^{\mathrm S}_3.
 \hypertarget{tag-02A2}{}\label{maps:eq:E.B6}
\end{equation}
We use Pontryagin--Thom in its form for maps to a space,
\cite[Definition 8.16 and Theorem 8.17]{rootDavisKirk01}.
The first component of $[S^6,e_i]$ is obtained by forgetting
the map, so it is zero for the bounding sphere framing.
The second component is the stable class of $e_i$, equivalently
the framed inverse image of a regular point. It is $\eta^3$
by \SeriesRef{maps:src:B.1.2}. One can also check the splitting directly: the
Pontryagin--Thom collapse of the framed sphere into its
disjoint-basepoint suspension has degree one on the sphere
summand; composing with $e_i$ gives its stable homotopy class.
Thus
\[
                         [S^6,e_i]=(0,\eta^3).
\]
Smashing the two sphere splittings gives
\[
 (S^3\times S^3)_+\simeq
       S^0\vee S^3\vee S^3\vee S^6.
\]
Under this splitting the external product
$[S^6\times S^6,e_1\times e_2]$ has four components.
The first three each contain a factor $[S^6]=0$; the last
is $\eta^3\eta^3=\eta^6=0$. Indeed $\eta^4$ belongs to
$\pi^{\mathrm S}_4=0$, as recorded in
\cite[Corollary~1.3 and Table~1, pp.~112--113]{IsaksenWangXu2023}. Therefore
\begin{equation}
 [S^6\times S^6,e_1\times e_2]=0
      \quad\hbox{in }\Omega^{\mathrm{fr}}_{12}(S^3\times S^3).
 \hypertarget{tag-02A3}{}\label{maps:eq:E.B7}
\end{equation}


\item \emph{The comparison with the classifying map.}
By \eqref{maps:eq:E.B7} there is a compact stably framed thirteen-manifold
$Z$ with boundary $S^6\times S^6$ and a map
$E:Z\to S^3\times S^3$ extending $e_1\times e_2$.
Smooth it relative to the boundary and make it transverse
to $N$ there. Transversality on the boundary already holds:
the derivative of its defining equation is $df_1-df_2$,
which is onto. Relative transversality therefore supplies
the interior perturbation without changing that boundary.

\item The normal bundle of $N\subset S^3\times S^3$ is the
restriction of $h_1^*T\mathbb P^1$ pulled back from the first
sphere. This oriented two-plane bundle is trivial because
$H^2(S^3;\mathbb Z)=0$. Choose a framing on it before
restriction to $N$. Its pullback under $e_1\times e_2$
is a framing of $f_1^*T\mathbb P^1$, agreeing with that
in \ref{maps:src:E.2.8} up to homotopy since $H^1(S^6;\mathbb Z)=0$.
It follows that
\[
                  V=E^{-1}(N),\qquad E|_V:V\longrightarrow N
\]
is a compact stably framed eleven-manifold with boundary
$(P,c)$ and the required boundary framing. This proves \eqref{maps:eq:E.B4}.

\end{enumerate}
The finite algebra and combinatorial data are implemented in \hyperlink{prog:V:equivariant}{the programs in Table~\ref*{tab:V:programs}}.
\end{proof}


\section{The fibre complement and its homotopy coinvariants}\label{pub:V:sec-complement}

\begin{definition}[Ordinary closed spin]\hypertarget{tag-PV23}{}\label{def:V:ordinary-closed-spin}
\StatementRecord{def:V:ordinary-closed-spin}{proved}{PV23}{}
For a closed
submanifold $L\subset\operatorname{int}D^n$, its ordinary closed
spin is $L\times S^1$ in
$D^n\times S^1\subset\partial(D^n\times D^2)=S^{n+1}$.
\end{definition}

\begin{proposition}[The fibre complement and a spinning obstruction]\hypertarget{tag-02MA}{}\label{maps:src:E.2.10}
\StatementRecord{maps:src:E.2.10}{proved}{02MA}{\cref{maps:src:E.2.6}, \SeriesRef{maps:src:B.5.2}, \eqref{maps:eq:E.MV4}, \eqref{maps:eq:E.MV6}, \eqref{maps:eq:E.S5}, \eqref{maps:eq:E.S4}; \cite[Section 3, immediately before Theorem 3.4]{mapsWinter08}; \cite[Theorem 4.5]{mapsWinter08}}


Let $F=T^4$ be a regular fibre of the original sphere fibration and
$E=S^6\setminus\operatorname{int}(F\times D^2)$, with the marking of
\ref{maps:src:E.2.6}. Its group and peripheral homomorphism are
\begin{equation}
 \begin{split}
 K=\pi_1(E)&=\langle a,b,h\mid h\text{ central},\ a^3=h,
                                      \ b^4=h^{-1}\rangle,\\
 g&\longmapsto h,\qquad u,w,d\longmapsto1,
       \qquad \mu\longmapsto ab.
 \end{split}
 \hypertarget{tag-02CC}{}\label{maps:eq:E.S1}
\end{equation}
Here $\mu$ is the normal meridian, with the peripheral orientation
chosen so that $[a]=4[\mu]$, $[b]=-3[\mu]$ and $[h]=12[\mu]$.
The image of the boundary group is the embedded copy
$\langle h,\mu\rangle\simeq\mathbb Z^2$. In the unique homotopy
class of normal framings for which every pushed-off loop of $F$
has linking number zero, the image of $g$ is
\begin{equation}
             \lambda=h\mu^{-12},
 \qquad \ker\{\pi_1(F)\longrightarrow K\}
                 =\mathbb Z\langle u,w,d\rangle.
 \hypertarget{tag-02CN}{}\label{maps:eq:E.S2}
\end{equation}
In particular the peripheral structure is that of the $(3,4)$
torus knot: the marked $g$ coordinate is its regular Seifert fibre,
and \eqref{maps:eq:E.S2} gives the preferred longitude. The embedded
four-torus $F\subset S^6$ is not the ordinary threefold closed spin
of that classical knot, even allowing ambient orientation reversal
and arbitrary changes of its normal framing. More strongly, it is
not a single ordinary closed spin of any embedded $T^3\subset S^5$.

\end{proposition}


\begin{proof}
\begin{enumerate}[label=\arabic*.,leftmargin=*]
\item Apply the van Kampen presentation in \SeriesRef{maps:src:B.5.2} with the regular base disc
removed. The cusp kills $w,d$ and its peripheral loop. Conjugating
$w$ by the order-three generator also kills $u$. The two finite caps
impose $a^3=h$ and $b^4=h^{-1}$; there is now no relation $ab=1$,
since the outer boundary has not been filled. The outer meridian
is $ab$ in the stated orientation. This proves \eqref{maps:eq:E.S1}, including the peripheral map. Eliminating $h$ gives $\langle a,b\mid a^3b^4=1\rangle$.
Its quotient by $h$ is $C_3*C_4$. The word $ab$ has infinite order
in that free product, so a relation $h^r\mu^s=1$ forces $s=0$.
Abelianization then forces $r=0$. This proves peripheral injectivity.
It also shows that $\lambda$ has infinite order: its image in
$C_3*C_4$ is $(ab)^{-12}$.

\item Write $\theta\in H^1(\partial E;\mathbb Z)$ for the normal-circle
covector and use $(\gamma,u,w,\delta)$ for $H^1(F;\mathbb Z)$.
The generator $\xi\in H^1(E;\mathbb Z)$ restricts to
$\theta+12\gamma$. Homotopy classes of normal framings differ by
$H^1(F;\mathbb Z)$. Consequently there is a unique framing for
which the associated section $s:F\to\partial E$ satisfies
$s^*\xi=0$: it has $s^*\theta=-12\gamma$. Its induced peripheral
map is exactly \eqref{maps:eq:E.S2}. Simultaneously reversing the meridian
conventions changes the displayed signs and preserves all the
following assertions.

\item For this section set
\begin{equation}
 \mathcal R=\operatorname{im}\{H^2(E;\mathbb Z)
                         \longrightarrow H^2(F;\mathbb Z)\},
 \qquad
 \operatorname{rad}(\mathcal R)=
 \{v\in H_1(F;\mathbb Z):\iota_v\omega=0
                                    \text{ for all }\omega\in\mathcal R\}.
 \hypertarget{tag-02CY}{}\label{maps:eq:E.S3}
\end{equation}
This uses the canonical identification of the cohomology of a torus
with the exterior algebra of its first cohomology. Equations
\eqref{maps:eq:E.MV4}--\eqref{maps:eq:E.MV6} give classes whose boundary restrictions are
$\theta u$ and $\theta w$: their unswept terms are zero because
$u\gamma uw=w\gamma uw=0$. Restricting them along $s$ proves
\begin{equation}
                -12\gamma u,\ -12\gamma w\ \in\mathcal R.
 \hypertarget{tag-02D9}{}\label{maps:eq:E.S4}
\end{equation}
The simultaneous contraction kernel of these two forms is
$\mathbb Z d$. Indeed their contractions force successively
$\gamma(v)=u(v)=w(v)=0$. Thus
\begin{equation}
       \operatorname{rad}(\mathcal R)\subseteq\ker\gamma
        =\ker\{\pi_1(F)\longrightarrow K\}.
 \hypertarget{tag-02DK}{}\label{maps:eq:E.S5}
\end{equation}
Both the zero-linking framing and this inclusion are invariants of
the embedded pair, independently of the torus coordinates.

\item Use the ordinary closed spin of Definition~\ref{def:V:ordinary-closed-spin}.
Start with the closed classical knot $K_0\subset\operatorname{int}D^3$
and repeat this operation three times, each time placing the
previous compact pair in the interior of a disc. This is the
ordinary torus spinning construction reviewed in
\cite[Section 3, immediately before Theorem 3.4]{mapsWinter08}.
If $\Sigma$ is a Seifert surface for $K_0$, the same three spins
give the embedded Seifert hypersurface
$\Sigma\times T^3$ with boundary $K_0\times T^3$.
Pushing the latter off using the normal of this hypersurface
gives its zero-linking framing.

\item Let $v$ be the original knot-circle class and let $v_1,v_2,v_3$
be the spinning circles. Each pushed-off two-cycle $v\wedge v_i$
bounds the corresponding copy of $\Sigma\times S^1$ in a parallel
copy of that hypersurface, disjoint from the knot. Hence every
restriction of a complement two-class vanishes on $v\wedge H_1(F)$:
$v$ belongs to $\operatorname{rad}(\mathcal R)$.
On the other hand its pushed-off loop is the original preferred
longitude and is nontrivial for $K_0=T(3,4)$, whereas the spinning
circles are trivial in the complement group. These peripheral
claims also follow directly at each spin: its complement is
\[
 ((D^n\setminus\operatorname{int}\nu L)\times S^1)
       \ \cup_{S^{n-1}\times S^1}\ (S^{n-1}\times D^2),
\]
so van Kampen kills precisely the new spinning circle and preserves
the old complement group. They agree with the primary peripheral
calculation \cite[Theorem 4.5]{mapsWinter08} in the first spin.
Thus the ordinary spun torus violates \eqref{maps:eq:E.S5}. An identification
of embedded pairs would preserve its two sides, giving a
contradiction.

\item For the stronger assertion, let $L=T^3\subset S^5$ be arbitrary
and suppose that its ordinary closed spin were our four-torus.
Write $E_L$ for its exterior. Deleting a ball disjoint from
$L$ changes cohomology only in degrees at least four. In the
displayed spun-complement decomposition, with $n=5$, the cap
and intersection are $S^4\times D^2$ and $S^4\times S^1$.
Mayer--Vietoris and the surjectivity onto the latter first
cohomology therefore give, in degrees two and three, precisely
the cohomology of $(E_L\setminus D^5)\times S^1$.
Let $s$ be the spinning-circle covector and $v_s$ its primitive
homology generator. Put $\mathcal R_L^j=\operatorname{im}
\{H^j(E_L)\to H^j(L)\}$, using the zero-linking framing.
One has $\mathcal R_L^1=0$. Also $\mathcal R_L^3=0$:
the pushed-off $L$ bounds a parallel Seifert hypersurface in
its exterior, so every top-degree restriction evaluates to
zero. The product restrictions consequently satisfy
\begin{equation}
       \mathcal R^2=\mathcal R_L^2,\qquad
       \mathcal R^3=\mathcal R_L^2\wedge s.
 \hypertarget{tag-02CF}{}\label{maps:eq:E.S12}
\end{equation}
In particular $v_s\in\operatorname{rad}(\mathcal R^2)$,
and contraction by $v_s$ is injective on $\mathcal R^3$.
For our four-torus, \eqref{maps:eq:E.S4} forces $v_s=\pm d$.
But \eqref{maps:eq:E.MV4}--\eqref{maps:eq:E.MV6}, now applied to $\alpha=\gamma\delta$
in degree three, show that $\mathcal R^3$ contains
$\pm2\gamma uw$. Its swept term vanishes in the zero-linking
framing, and this is a nonzero class annihilated by $d$.
This contradicts that injectivity. No ordinary closed spin
is possible. 
\end{enumerate}
\end{proof}

\begin{theorem}[Exact second-homotopy coinvariants]\hypertarget{tag-02MB}{}\label{maps:src:E.2.11}
\StatementRecord{maps:src:E.2.11}{proved}{02MB}{\cref{maps:src:E.2.9}, \cref{maps:src:E.2.6}, \eqref{maps:eq:E.S7}, \eqref{maps:eq:E.MV7}, \eqref{maps:eq:E.MV8}, \eqref{maps:eq:E.MV6}, \eqref{maps:eq:E.MV4}, \eqref{maps:eq:E.S9}, \eqref{maps:eq:E.S10}, \eqref{maps:eq:E.S11}}


For the classifying map $c:P\to N=K(G,1)$ of \ref{maps:src:E.2.9}, one has
\begin{equation}
 c_*:H_3(P;\mathbb Z)\longrightarrow H_3(N;\mathbb Z)=\mathbb Z
       \quad\text{equal to zero},
 \qquad
 \bigl(\pi_2(P)\bigr)_G
       \simeq\mathbb Z^2\oplus(\mathbb Z/12)^6.
 \hypertarget{tag-02DR}{}\label{maps:eq:E.S6}
\end{equation}
The pullback of a generator of $H^3(N;\mathbb Z)$ has exact order
twelve in $H^3(P;\mathbb Z)$. In particular $\pi_2(P)$ is not a
projective $\mathbb ZG$-module. The fivefold cannot be diffeomorphic,
or homotopy equivalent, to $N\times S^6$ connected-summed with any
finite collection of $S^r\times S^{10-r}$, $2\leq r\leq8$, and
any homotopy ten-sphere. \end{theorem}


\begin{proof}
\begin{enumerate}[label=\arabic*.,leftmargin=*]
\item Use the separating base circle in \ref{maps:src:E.2.6} also in the target of
\ref{maps:src:E.2.9}. There $M$ denotes the exterior of a regular orbit of the
weighted Hopf map, namely the torus-knot exterior of \SeriesRef{mapsI:src:A.1.3}. The
two decompositions, and the restriction of $c$, are
\begin{equation}
 \begin{array}{ccc}
 (E\times F')\cup(F\times E')&\longrightarrow&
                  (M\times S^1)\cup(S^1\times M')=N,\\
 F\times F'\times S^1&\longrightarrow&S^1\times S^1\times S^1.
 \end{array}
 \hypertarget{tag-02DS}{}\label{maps:eq:E.S7}
\end{equation}
On the common boundary the map retains exactly the covectors
$\gamma_1,\gamma_2,\theta$. It is the literal quotient there.
Each knot exterior has the integral homology of a circle, so the
third homology and cohomology of each target piece in \eqref{maps:eq:E.S7} vanish.

\item We first prove the homological statement. In the degree-two source of \eqref{maps:eq:E.MV7}--\eqref{maps:eq:E.MV8}
write $t_1,t_2,t_{12}$ for the coefficients of
$\theta\times g_1$, $\theta\times g_2$, and $g_1\times g_2$.
The image of $i_2:H_2(F)\to H_2(E)$ lies entirely in the
$d$-sweep coordinate, by \eqref{maps:eq:E.MV6}. Consequently the $g_1$-sweep
coordinate in the first $H_2(E)$ summand of the target is just
$t_1$, and the corresponding second coordinate is just $t_2$.
Every vector in the kernel of the boundary matrix therefore has
$t_1=t_2=0$. Its first $H_1(E)\otimes\mathbb Zg_2$ coordinate
then forces $12t_{12}=0$, up to an immaterial orientation sign.
These three coordinates are the only ones retained in
$H_2(T^9)\to H_2(T^3)$. The target Mayer--Vietoris sequence
injects $H_3(N)$ into $H_2(T^3)$; its source-piece third groups
are zero. Naturality now proves $c_*H_3(P)=0$.

\item The cohomological calculation determines the extension. On the target boundary write
$v=\gamma_1\gamma_2$. The generators of the two degree-two
piece cohomology groups restrict, with consistent meridian
orientations, to
\begin{equation}
       \theta\gamma_2+12v,
       \qquad \theta\gamma_1+12v.
 \hypertarget{tag-02DT}{}\label{maps:eq:E.S8}
\end{equation}
Their image is a primitive rank-two subgroup of $H^2(T^3)$:
the minor in the last two coordinates is a unit. Its quotient
is freely generated by $v$. Since the degree-three piece groups
are zero, the connecting class $\partial v$ is a generator of
$H^3(N;\mathbb Z)$.

\item In the source decomposition both $\theta\gamma_1$ and
$\theta\gamma_2$ already belong to the image of the degree-two
piece restrictions. This is \eqref{maps:eq:E.MV4} with $\alpha=\gamma$,
since $\phi_2(\gamma)=0$; the second normal circle is reversed,
which only changes its sign. Subtracting one of these classes
from the corresponding product of a degree-one class of $E$
with $\gamma_2$ proves that $12v$ is a source boundary image.
Conversely the coefficient of $v$ in every such boundary image
is divisible by twelve. The pure degree-two terms from a single
factor have no $v$ coefficient, and the mixed terms use the
restriction $\theta+12\varepsilon\gamma$ of the degree-one
generator, with $\varepsilon=\pm1$. It follows that
\begin{equation}
                  c^*(\partial v)=\partial v
        \quad\text{has exact order }12.
 \hypertarget{tag-02DU}{}\label{maps:eq:E.S9}
\end{equation}
This uses the Mayer--Vietoris image.

\item Put $A=(\mathbb Z/12)^7$ and $L=(\pi_2(P))_G$.
The low-degree homology sequence of the universal-cover
fibration, together with \ref{maps:src:E.2.6} and \ref{maps:src:E.2.9}, gives
\begin{equation}
             0\longrightarrow\mathbb Z z\longrightarrow L
                \longrightarrow\mathbb Z\oplus A\longrightarrow0.
 \hypertarget{tag-02CD}{}\label{maps:eq:E.S10}
\end{equation}
The zero on the left uses $c_*H_3(P)=0$, just proved.
To determine this extension, replace $c$ by its mapping cylinder
and use relative groups for $(N,P)$. The induced map on fundamental
groups is an isomorphism and the target has zero second homotopy
group. On universal covers the relative pair is two-connected;
relative Hurewicz gives
$H_3(\widetilde N,\widetilde P)=\pi_2(P)$ and all lower relative
groups vanish. The relative covering homology sequence therefore
identifies $H_3(N,P)=L$ and $H_2(N,P)=0$.
The universal coefficient theorem and the relative cohomology
sequence identify the map
\[
       H^3(N,P)=\operatorname{Hom}(L,\mathbb Z)
         \longrightarrow H^3(N)=\mathbb Z
\]
with evaluation on $z$. Its image is $12\mathbb Z$ by \eqref{maps:eq:E.S9}.

\item The free summand of the quotient in \eqref{maps:eq:E.S10} splits. Lift its
seven order-twelve generators to $x_1,\ldots,x_7$ in the
remaining extension. There are integers $e_i$, defined modulo
twelve, with the complete abelian presentation
\begin{equation}
        \langle z,x_1,\ldots,x_7\mid
                      12x_i=e_i z\ (1\leq i\leq7)\rangle.
 \hypertarget{tag-02CE}{}\label{maps:eq:E.S11}
\end{equation}
Evaluation of homomorphisms to $\mathbb Z$ on $z$ has image
$\bigl(12/\gcd(12,e_1,\ldots,e_7)\bigr)\mathbb Z$.
Thus that gcd is one. The relation matrix in \eqref{maps:eq:E.S11} has
Smith factors $1,12,12,12,12,12,12$, giving
$\mathbb Z\oplus(\mathbb Z/12)^6$. Adding the split free
summand proves \eqref{maps:eq:E.S6}.

\item If $\pi_2(P)$ were projective over $\mathbb ZG$, it would be
a direct summand of a free module. Taking coinvariants preserves
that splitting, so $L$ would be a projective, hence free,
abelian group. The torsion in \eqref{maps:eq:E.S6} contradicts this. Finally
$N\times S^6$ has zero second homotopy group. Connected sum in
dimension ten with a simply connected sphere product adds its
second homotopy as a free induced $\mathbb ZG$-module, as seen
by lifting the connected-sum necks and using their simply
connected boundary spheres. Only $r=2,8$ contribute, with one
free generator each; a homotopy sphere contributes none.
Their second homotopy modules are therefore free, so none is
homotopy equivalent to $P$.

\end{enumerate}
The finite algebra and combinatorial data are implemented in \hyperlink{prog:V:equivariant}{the programs in Table~\ref*{tab:V:programs}}.
\end{proof}

\begin{theorem}[Central localizations of the homotopy module]\hypertarget{tag-02MC}{}\label{maps:src:E.2.12}
\StatementRecord{maps:src:E.2.12}{proved}{02MC}{\cref{maps:src:E.2.9}, \cref{maps:src:E.2.10}, \eqref{maps:eq:E.S17}, \cref{maps:src:E.2.3}, \eqref{maps:eq:E.S19}, \eqref{maps:eq:E.S21}, \cref{maps:src:E.2.11}}


Take the central generators $h_1,h_2$ of $G$ from \ref{maps:src:E.2.9} and put
\begin{equation}
 R=\mathbb ZG,\qquad R_i=R[(h_i-1)^{-1}],\qquad
              R_{12}=R[(h_1-1)^{-1},(h_2-1)^{-1}].
 \hypertarget{tag-02CG}{}\label{maps:eq:E.S13}
\end{equation}
These are flat central localizations. As modules over the indicated
localized group rings, the second homotopy group satisfies
\begin{equation}
 \begin{split}
 R_1\otimes_R\pi_2(P)&\simeq R_1/(h_2-1),\\
 R_2\otimes_R\pi_2(P)&\simeq R_2/(h_1-1),\\
 R_{12}\otimes_R\pi_2(P)&=0.
 \end{split}
 \hypertarget{tag-02CH}{}\label{maps:eq:E.S14}
\end{equation}
More generally, $H_j(P;R_1)$ is zero outside $2\leq j\leq7$;
in those six degrees it is respectively a direct sum of
\begin{equation}
                  1,\ 3,\ 4,\ 4,\ 3,\ 1
 \hypertarget{tag-02CI}{}\label{maps:eq:E.S15}
\end{equation}
copies of $R_1/(h_2-1)$. The analogous statement holds with the
indices interchanged. All homology with coefficients in $R_{12}$
vanishes. For the fibre complement of \ref{maps:src:E.2.10}, with
$S=\mathbb ZK[(h-1)^{-1}]$, one likewise has
\begin{equation}
             S\otimes_{\mathbb ZK}\pi_2(E)\simeq S,
       \qquad H_4(E;S)\simeq S,
 \hypertarget{tag-02CJ}{}\label{maps:eq:E.S16}
\end{equation}
and its other localized homology groups vanish. Moreover, if
$q=(h_1-1)(h_2-1)$, then
\begin{equation}
                  q^2\pi_2(P)=0,
        \qquad \dim_{\mathbb Q}\bigl(\pi_2(P)\otimes\mathbb Q\bigr)
                     =\infty.
 \hypertarget{tag-02CO}{}\label{maps:eq:E.S20}
\end{equation}
Consequently the universal cover has no finite homotopy type.
There is no fibration of $P$ over $N$ with compact manifold fibre
and inducing an isomorphism on fundamental groups.

\end{theorem}


\begin{proof}
\begin{enumerate}[label=\arabic*.,leftmargin=*]
\item \emph{The cusp module with its marked deck transformation.}
In the original cusp, the angular generators $w,d$ vanish and
$\pi_1(N_0)=\mathbb Z\langle g,u\rangle$. In $E$ their images
are respectively $h,1$; in the first factor of $P$ they are
$h_1,1$. Thus the connected cover of this cusp occurring inside
the universal cover is the cover $\widehat W$ of its central
fibre in which the $g$ deck transformation is $z$, whereas $u$
remains a loop. Write $A=\mathbb Z[z,z^{-1}]$.
Its ordinary homology modules are
\begin{equation}
 H_0(\widehat W)=H_1(\widehat W)=A/(z-1),\quad
 H_2(\widehat W)=A\oplus(A/(z-1))^2,
 \quad H_3(\widehat W)=0,\quad H_4(\widehat W)=A.
 \hypertarget{tag-02CK}{}\label{maps:eq:E.S17}
\end{equation}
Here is an integral normalization calculation of this input.
The dual complex of the cover is an infinite cylinder, whose
ordinary homology is $A/(z-1)$ in degrees zero and one and zero
in degree two. The normalization degree-four row is a free
copy of $A$. In degree two, the homology map from the three
double curves to the four component classes is the transpose,
up to the unimodular intersection marking and possible inversion
of $z$, of
\begin{equation}
 B(z)=\begin{pmatrix}
 -1&-1&-1&-1\\ 1&1&1&z\\-z&-z&-1&-z
 \end{pmatrix}
 \ \sim\
 \begin{pmatrix}1&0&0&0\\0&z-1&0&0\\0&0&z-1&0\end{pmatrix}.
 \hypertarget{tag-02CL}{}\label{maps:eq:E.S18}
\end{equation}
For example multiply on the left by
$\left(\begin{smallmatrix}-1&0&0\\1&1&0\\-z&0&1\end{smallmatrix}\right)$,
then use as new columns $e_1,e_4-e_1,e_3-e_1,e_2-e_1$.
Both operations are unimodular over $A$. The transpose is
injective and its cokernel is the displayed degree-two module.
All strata have only even homology, and the normalization
direction has length two, so no higher differential is possible.
There is no degree-two extension from the dual complex, whose
second homology is zero. This proves \eqref{maps:eq:E.S17}. Ordinary chains
in this locally finite normalization have compact support and
use direct sums.

\item \emph{Localization and the vanishing of every other piece.}
The central subgroup $\langle h_1,h_2\rangle$ is an injected
$\mathbb Z^2$ by \ref{maps:src:E.2.9}. The group ring is free over its central
Laurent subring, by choosing coset representatives. In particular
$h_i-1$ and their products are non-zero-divisors. Inverting them
is central localization and is flat; it commutes with homology
of the universal-cover chain complex. The analogous assertions
hold for $h$ in $K$.

\item Use the regular torus bundle over a graph and the six closed
cap neighborhoods in \ref{maps:src:E.2.3}. For $R_1$ coefficients, the first
regular torus has its $g_1$ loop acting by $h_1$. The torus
chain complex is the Koszul complex of its commuting loop
actions. Its $h_1-1$ term is invertible, providing a chain
contraction. Thus every regular fibre has zero homology with
these coefficients. The homology sequence of a bundle over
a graph, and the same calculation over each boundary circle,
show that the whole regular piece and every cap boundary are
acyclic. Each cap of the second factor is a product with a
first-factor regular torus over a disc, so it too is acyclic.

\item At a finite cap of the first factor, the affine generator
$a_j$ has $a_j^{m_j}=h_1^{\pm1}$ in its image in $G$.
The connected cover of the reduced flat surface corresponding
to its map onto $\langle a_j\rangle$ is $T^3\times\mathbb R$.
The $m_j$-th power acts by an integral translation with linear
part $A_j^{m_j}=I$, and therefore acts identically on all of
its homology. Thus $h_1-1$ annihilates every such homology
module before localization. Keeping the other regular torus does not change that assertion, and induction to $G$ commutes
with the central action. Both first-factor finite caps are
therefore acyclic after localization as well. 

\item \emph{The one surviving cap and induction.}
The first cusp is $N_{0,1}\times F_2$, since the two critical
triples are disjoint. Its image in $G$ is precisely
$H=\langle h_1,h_2\rangle$. A connected component of its
preimage in the universal cover is
\[
                \widehat W\times\mathbb R\times T^3
\]
up to deformation retraction: $g_1$ acts on $\widehat W$ by
$z=h_1$, and $g_2=h_2$ translates the real factor. It acts
trivially on the homology of that three-torus. All
components together give induction from $\mathbb ZH$ to
$R$. By \eqref{maps:eq:E.S17}, inverting $h_1-1$ leaves one free
$\mathbb Z[h_1^{\pm1},(h_1-1)^{-1}]$ summand in cusp degrees
two and four and nothing else. The integral product homology
with $T^3$ has no torsion correction. Inducing the result gives
\begin{equation}
 H_j(N_{0,1}\times F_2;R_1)
   \simeq
 \left(R_1/(h_2-1)\right)^{
           \binom3{j-2}+\binom3{j-4}},
 \hypertarget{tag-02CM}{}\label{maps:eq:E.S19}
\end{equation}
where a binomial coefficient is zero outside its usual range.
All other pieces and their intersections are acyclic, as
proved above, so Mayer--Vietoris identifies this with
$H_j(P;R_1)$. This proves \eqref{maps:eq:E.S15} and \eqref{maps:eq:E.S19}, and interchanging
the factors proves the second version. Inverting $h_2-1$
also kills \eqref{maps:eq:E.S19}, proving the double-localization assertion.
Since $\widetilde P$ is simply connected, Hurewicz identifies
its second homology with $\pi_2(P)$; flat localization now
gives \eqref{maps:eq:E.S14}.

\item For $E$ the same decomposition has only one cusp and two
finite caps. Every other piece is acyclic after $h-1$ is
inverted by the identical argument. Its cusp image is
$\langle h\rangle$, so inducing \eqref{maps:eq:E.S17} after localization
gives one free $S$ summand in degrees two and four. This
proves \eqref{maps:eq:E.S16}.

\item \emph{A uniform annihilator without localization.}
Put $I=(h_1-1,h_2-1)\subset R$. Before localization, the homology
of a regular eight-torus with $R$ coefficients is
\begin{equation}
                 (R/I)\otimes_{\mathbb Z}\bigwedge^j\mathbb Z^6
                      \quad(0\leq j\leq6).
 \hypertarget{tag-02CP}{}\label{maps:eq:E.S21}
\end{equation}
Indeed the two nontrivial loop images are the independent central
generators $h_1,h_2$. Their differences from one form a regular
sequence in the central Laurent ring, and $R$ is free over that
ring. The two-variable Koszul complex therefore has only its
degree-zero quotient as homology; the other six loop differentials
are zero. This proves \eqref{maps:eq:E.S21}.
The homology sequence of the regular bundle over a graph has
only base degrees zero and one. Both filtration quotients are
annihilated by $I$, by \eqref{maps:eq:E.S21}, so $I^2$ annihilates the homology
of the regular piece. The same two-stage Wang argument applies
to every boundary mapping torus. 

\item At a finite cap its connected lift is a product of six torus
directions and two real directions up to deformation retraction.
Each central deck generator acts trivially on homology: the
one belonging to that cap is the finite affine power already
checked above, and the other translates the regular real
direction. Hence $I$ annihilates its induced homology. At the
first cusp, $h_2$ translates the real factor of the other
regular torus, so $h_2-1$ annihilates all of its homology;
the second cusp is annihilated by $h_1-1$. Thus $q$ annihilates
the homology of every piece and every intersection, since
$q\in I^2$. The decomposition has a regular piece and disjoint
caps with disjoint boundary collars, so its Mayer--Vietoris
sequence gives in every degree
\[
 0\longrightarrow\operatorname{coker}D_j
   \longrightarrow H_j(P;R)
   \longrightarrow\ker D_{j-1}\longrightarrow0.
\]
Both end groups are killed by $q$, and therefore the middle
group is killed by $q^2$. Since $H_2(P;R)=\pi_2(P)$, this
proves the first assertion of \eqref{maps:eq:E.S20} integrally.

\item \emph{Infinite rational rank.}
If $\pi_2(P)\otimes\mathbb Q$ were finite dimensional,
localizing its linear automorphism $h_1$ at $h_1-1$ would
remain finite dimensional. Explicitly, this removes the
generalized eigenspace with eigenvalue one and inverts an
automorphism on the remaining finite-dimensional quotient.
But \eqref{maps:eq:E.S14}, tensored with $\mathbb Q$, contains a copy of
$\mathbb Q[h_1^{\pm1},(h_1-1)^{-1}]$, as follows by taking
the identity coset in the free module over the central
Laurent subring and setting $h_2=1$. That vector space is
infinite dimensional, a contradiction. This proves the
second assertion and rules out a finite homotopy type of
the universal cover. A fibration over the aspherical $N$
inducing a fundamental-group isomorphism would have simply
connected fibre, and pulling it to $\widetilde N$ would
identify that fibre up to homotopy with $\widetilde P$.
A compact manifold fibre has finite homotopy type, so such
a fibration is impossible.

\end{enumerate}
The corresponding exact calculations are implemented in \hyperlink{prog:V:equivariant}{the programs in Table~\ref*{tab:V:programs}}.
\end{proof}

\section{The circle fibration and its torus twist}\label{pub:V:sec-torus-twist}

\begin{theorem}[A circle fibration and infinite-order return]\hypertarget{tag-02MD}{}\label{maps:src:E.2.13}
\StatementRecord{maps:src:E.2.13}{proved}{02MD}{\cref{maps:src:E.2.9}, \eqref{maps:eq:E.S26}, \cref{maps:src:E.2.12}, \SeriesRef{maps:src:B.6.24}; \cite[Theorem 3.6]{rootScott83}}


The shifted fivefold admits a smooth fibre bundle
\begin{equation}
                  Y^9\longrightarrow P\stackrel{\rho}{\longrightarrow}S^1
 \hypertarget{tag-02CQ}{}\label{maps:eq:E.S22}
\end{equation}
with connected closed fibre, inducing the primitive character
\begin{equation}
 \chi(a_1,a_2,b_1,b_2,h_1,h_2)
                    =(4,-3,-4,3,12,-12)
 \hypertarget{tag-02CR}{}\label{maps:eq:E.S23}
\end{equation}
on the presentation of \ref{maps:src:E.2.9}. Its fibre group is the fundamental
group of the oriented principal circle bundle of Euler number
$-1$ over a surface of genus six. Here the base has its complex
orientation and the positive circle is $h_1h_2$; reversing that
circle changes the Euler number to $+1$. In particular
\begin{equation}
 \pi_1(Y)=\ker\chi,\qquad H_1(Y;\mathbb Z)=\mathbb Z^{12}.
 \hypertarget{tag-02CS}{}\label{maps:eq:E.S24}
\end{equation}
The return diffeomorphism of \eqref{maps:eq:E.S22} has infinite order even
among unbased homotopy classes of self-equivalences of $Y$,
although its induced outer automorphism of $\pi_1(Y)$ has
exact order twelve. No connected finite cover of $P$ is
homotopy equivalent to a product $S^1\times Z$.


\end{theorem}


\begin{proof}
\begin{enumerate}[label=\arabic*.,leftmargin=*]
\item \emph{A closed one-form with a nonzero derivative at every point.}
Use the smooth map $c=(e_1,e_2):P\to N$ constructed in \ref{maps:src:E.2.9}. The manifold $N$ has the locally free $T^2$ action
induced by the two weighted Hopf actions. Take a closed
one-form representing the primitive integral class $\chi$
on $N$ and average it over this compact torus. Its average
$\alpha$ is closed, represents the same integral class,
and is invariant. If $V_i$ is the fundamental vector field
with time-one orbit $h_i$, Cartan's formula gives
$d(\alpha(V_i))=0$. Integrating along an ordinary orbit
therefore gives
\begin{equation}
                  \alpha(V_1)=12,\qquad \alpha(V_2)=-12.
 \hypertarget{tag-02CT}{}\label{maps:eq:E.S25}
\end{equation}
The integrality of the periods integrates $\alpha$ to a
smooth map $\rho_N:N\to\mathbb R/\mathbb Z$.

\item Choose the two cusp replacements defining $e_i$ in disjoint
base discs, themselves disjoint from the other factor's
critical values. At every point of $P$, at least one factor
is a regular torus at which its $e_i$ is the literal character
quotient of \SeriesRef{mapsI:src:A.1.1}. In a local flat marking of that factor,
the constant vertical vector $g_i$ maps to $V_i$. Holding
the other factor fixed gives a tangent vector to $P$, and
\[
                     (c^*\alpha)(g_i)=12\quad\hbox{or }-12.
\]
Thus $c^*\alpha$ is nowhere zero.  The map $\rho=\rho_Nc$ is a
proper submersion. Equivalently, choose a vector field $W$
with $(c^*\alpha)(W)=1$; its complete flow supplies local
trivializations and a return diffeomorphism. The fibre is
connected because $\chi$ is primitive: the image in
$\pi_1(S^1)$ of a connected component of a bundle's total
space has index equal to the number of fibre components.
The homotopy sequence now proves $\pi_1(Y)=\ker\chi$.

\item \emph{The exact surface cover and its circle Euler class.}
Write $\Gamma=G/\langle h_1,h_2\rangle$ for the oriented
orbifold group of $S^2(3,4,3,4)$. The character modulo
twelve descends to the surjection
\begin{equation}
 \overline\chi:\Gamma\longrightarrow\mathbb Z/12,
       \qquad(a_1,a_2,b_1,b_2)\longmapsto(4,-3,-4,3).
 \hypertarget{tag-02CU}{}\label{maps:eq:E.S26}
\end{equation}
Its restriction to every cone group is injective. Hence its
kernel $\Gamma_0$ is torsion free and defines a degree-twelve
ordinary surface cover of the orbifold. Its Euler characteristic
is $12(-5/6)=-10$, so that surface is $\Sigma_6$. The image
of $\ker\chi$ in $\Gamma$ is exactly $\Gamma_0$: a lift of
an element of $\Gamma_0$ has character a multiple of twelve,
which can be cancelled by a power of $h_1$. Its central
kernel is the primitive subgroup
\begin{equation}
 \ker\chi\cap\langle h_1,h_2\rangle
                =\mathbb Z\langle h_1h_2\rangle.
 \hypertarget{tag-02CV}{}\label{maps:eq:E.S27}
\end{equation}

\item We compute the circle Euler number in the following orientation convention. On either weighted Hopf sphere put
$z_1=\cos r\,e^{i\varphi_1}$ and
$z_2=\sin r\,e^{i\varphi_2}$. For the positive action with
weights $(4,3)$ the smooth invariant connection is
\[
 \vartheta=\tfrac14\cos^2r\,d\varphi_1
                    +\tfrac13\sin^2r\,d\varphi_2.
\]
It evaluates to one on the generator with period $2\pi$.
With the boundary orientation of the unit ball in $\mathbb C^2$,
\begin{equation}
             \frac1{(2\pi)^2}\int_{S^3}\vartheta\wedge d\vartheta
                         =\frac1{12}.
 \hypertarget{tag-02CW}{}\label{maps:eq:E.S28}
\end{equation}
For the associated complex line whose unit circle is this
positive circle bundle, the convention agreeing with the
tautological Hopf line is $c_1=-[d\vartheta/(2\pi)]$.
Thus each orbifold circle Euler number is $-1/12$.
Both shifted base identifications preserve the complex
orientation. Pullback to $\Sigma_6$ multiplies the two Euler
numbers by twelve, giving $(-1,-1)$; this is also the
covering formula of \cite[Theorem 3.6]{rootScott83}.
In the integral fibre basis
\[
                  v=(1,1),\qquad w=(0,1)
\]
the determinant is one, the Euler vector is $-v$, and
$\chi(v)=0$, $\chi(w)=-12$. Consequently the circle kernel
has Euler number $-1$ with generator $h_1h_2$.

\item One can see this bundle directly in $N$. Invariance of
$\alpha$ implies
$\rho_N((s_1,s_2)x)=\rho_N(x)+12(s_1-s_2)$, where
$s_i\in\mathbb R/\mathbb Z$. A level contains twelve
diagonal circles over an ordinary base point. The peripheral
permutations \eqref{maps:eq:E.S26} give exactly the connected cover
$\Sigma_6$; they have no stabilizers at the cone points.
The diagonal circle acts freely, since at every cone only
one of the two Hopf actions has a nontrivial stabilizer.
The level is therefore the circle bundle just computed.
This identifies its fundamental group with $\ker\chi$
and gives a three-dimensional aspherical model.
The homology sequence of a circle bundle of Euler number
$\pm1$ gives its first homology $\mathbb Z^{12}$, proving
the second assertion in \eqref{maps:eq:E.S24}.

\item \emph{Faithful central action and the return obstruction.}
The centre of $G$ is exactly $\langle h_1,h_2\rangle$:
the quotient is a nonelementary oriented hyperbolic orbifold
group with trivial centre. This centre acts faithfully on
$\pi_2(P)$. Indeed if $h_1^a h_2^b$ acts as the identity,
localize first at $h_1-1$ and apply \ref{maps:src:E.2.12}. On
$R_1/(h_2-1)$ its action is multiplication by $h_1^a$.
This module is free over the central ring
$\mathbb Z[h_1^{\pm1},(h_1-1)^{-1}]$ after choosing coset
representatives, so this is the identity only if $a=0$.
The symmetric localization then gives $b=0$.

\item If $H<G$ has finite index, its image in $\Gamma$ is again
nonelementary. Its centralizer in $\Gamma$ is trivial:
commuting with two hyperbolic elements with distinct axes
fixes at least three boundary points and hence is the
identity. It follows that
$Z(H)=H\cap Z(G)$. A homotopy product $S^1\times Z$ for
the corresponding finite cover would supply an infinite
central element acting trivially on $\pi_2$, namely the
circle factor; this contradicts the faithful action just
proved. If a positive power of the return of \eqref{maps:eq:E.S22} were
homotopic to the identity, its mapping torus would be
homotopy equivalent to $S^1\times Y$. That mapping torus
is the corresponding finite cyclic cover of $P$, so this
too is impossible.

\item Finally take $t=a_1a_2$, for which $\chi(t)=1$, as a
stable letter. The image of $\ker\chi$ in $\Gamma$ is
$\Gamma_0$, so the same centralizer argument gives
$C_G(\ker\chi)=Z(G)$. The automorphism
$\operatorname{Ad}(t^n)|_{\ker\chi}$ is inner precisely
when $t^n$ differs from an element of $\ker\chi$ by a
central element. Since $\chi(Z(G))=12\mathbb Z$, this
holds precisely for $12\mid n$. Thus the outer action has
order twelve, while the return has infinite homotopy
order. 

\end{enumerate}
The corresponding exact calculations are implemented in \hyperlink{prog:V:cyclic-cover}{the programs in Table~\ref*{tab:V:programs}}.
\end{proof}

\begin{theorem}[The cyclic-cover splice and its separating torus twist]\hypertarget{tag-02ME}{}\label{maps:src:E.2.14}
\StatementRecord{maps:src:E.2.14}{proved}{02ME}{\cref{maps:src:E.2.13}, \cref{maps:src:E.2.6}, \cref{maps:src:E.2.10}, \eqref{maps:eq:E.S34}, \eqref{maps:eq:E.S35}, \eqref{maps:eq:E.S37}, \eqref{maps:eq:E.S39}, \eqref{maps:eq:E.S38}, \eqref{maps:eq:E.S41}, \eqref{maps:eq:E.S43}, \eqref{maps:eq:E.S36}, \eqref{maps:eq:E.S40}, \eqref{maps:eq:E.S44}, \eqref{maps:eq:E.S45}, \cref{maps:src:E.2.7}, \eqref{maps:eq:E.S47}, \eqref{maps:eq:E.S46}}


Let $E_{12}\to E$ be the connected cyclic cover defined by
reduction modulo twelve of the positive meridian character.
The circle fibration of \ref{maps:src:E.2.13} may be chosen with fibre
\begin{equation}
 Y=(E_{12}\times T^3)\ \cup_{T^8}\ (T^3\times E_{12}),
 \hypertarget{tag-02CX}{}\label{maps:eq:E.S29}
\end{equation}
with the following exact boundary marking. Write
$\theta_i'$ for the lifted normal circle and use the
two regular-torus markings. Then
\begin{equation}
       \theta_1'=\gamma_2-\gamma_1,\qquad
       \theta_2'=\gamma_1-\gamma_2.
 \hypertarget{tag-02CZ}{}\label{maps:eq:E.S30}
\end{equation}
There is a return representative $\psi$ whose twelfth power
is the separating torus twist
\begin{equation}
 \begin{gathered}
 T(x,s)=(x+\eta(s)(g_1+g_2),s)\quad\hbox{on }T^8\times[-1,1],
 \\
 T=\operatorname{id}\text{ elsewhere}.
 \end{gathered}
 \hypertarget{tag-02D0}{}\label{maps:eq:E.S31}
\end{equation}
Here $\eta$ is zero near $-1$ and one near $1$, and circle
coordinates have period one. In particular \eqref{maps:eq:E.S31} is the
identity near both ends of its collar. Its rational
cohomological variation satisfies
\begin{equation}
 (T^*-I)^2=0,\qquad
 \operatorname{rank}(T^*-I)|_{H^j(Y;\mathbb Q)}=
 \begin{cases}1&j=4,5,\\0&j\ne4,5.\end{cases}
 \hypertarget{tag-02D1}{}\label{maps:eq:E.S32}
\end{equation}
The rational Betti numbers of the fibre are
\begin{equation}
       (b_0,\ldots,b_9)(Y)
          =(1,12,121,407,693,693,407,121,12,1).
 \hypertarget{tag-02D2}{}\label{maps:eq:E.S33}
\end{equation}
Every eigenvalue of $\psi^*$ is a twelfth root of unity.
At eigenvalue one there is exactly one Jordan block of
size two in each of degrees four and five; all its other
blocks, and all blocks at other eigenvalues, have size one.
Thus the infinite order is already visible on ordinary
rational cohomology. \end{theorem}


\begin{proof}
\begin{enumerate}[label=\arabic*.,leftmargin=*]
\item \emph{Adapting the fibration to the splice.}
Cut the base along the circle used in \ref{maps:src:E.2.6}, with the two
critical triples on opposite sides. Over the first disc the
second factor is a regular marked $T^4$, and over the second
disc the first factor is regular. Let
$\xi_E:E\to\mathbb R/\mathbb Z$ represent the positive
meridian class. Its restriction to the boundary can be
chosen exactly
\begin{equation}
                         \xi_E=\theta+12\gamma.
 \hypertarget{tag-02D3}{}\label{maps:eq:E.S34}
\end{equation}
Indeed this is its integral restriction by \ref{maps:src:E.2.10}; homotopy
extension makes the representative equal to that linear
map on a collar. The two maps
\begin{equation}
        \rho_A=\xi_{E_1}-12\gamma_2,
             \qquad \rho_B=12\gamma_1-\xi_{E_2}
 \hypertarget{tag-02D4}{}\label{maps:eq:E.S35}
\end{equation}
agree exactly on the common boundary, where
$\theta_2=-\theta_1$. Smooth them with this common product
collar. The regular-factor circle derivative is $-12$ on
the first piece and $12$ on the second, so they give a
proper submersion throughout, including both cusp pieces.
These are also permissible choices in \cref{maps:src:E.2.13}: on the
quotient $N$, an invariant closed one-form differs from
these representatives only by the differential of a
single-valued function on each of the two pieces. Absorb
that function by translation in the unused regular circle.
Thus this construction uses the same quotient and
primitive character.

\item A level of $\xi_E-12\gamma_2$ is the pullback of the
degree-twelve circle cover along $\xi_E$, times the three
remaining regular circles. The pullback is $E_{12}$.
Its boundary is connected because the normal meridian has
character one. The covering of that boundary has
$\theta=12\theta'$ and unchanged $\gamma,u,w,\delta$.
Equation \eqref{maps:eq:E.S34} in the level condition gives precisely
\eqref{maps:eq:E.S30}. This proves \eqref{maps:eq:E.S29} with its smooth
boundary maps.

\item Choose the vector field integrating \eqref{maps:eq:E.S35} to be
$-g_2/12$ on the first piece off its collar and $g_1/12$
on the second. On the common regular collar interpolate by
\begin{equation}
                     W_s=\frac{\eta(s)g_1-(1-\eta(s))g_2}{12}.
 \hypertarget{tag-02D5}{}\label{maps:eq:E.S36}
\end{equation}
It has $d\rho(W_s)=1$, and it preserves the collar
coordinate $s$ and the original base point. Its flow
therefore defines a return $\psi$. At time twelve it is
the identity off the collar, since the two indicated
regular circles have period one. On the collar it
translates by $\eta(s)g_1-(1-\eta(s))g_2$; removing the
integral translation $-g_2$ gives \eqref{maps:eq:E.S31}. This comparison uses a globally regular factor on each piece.

\item \emph{The variation map and the invariant boundary calculation.}
Put $D=T^8$ and let $r:H^j(Y;\mathbb Q)\to H^j(D;\mathbb Q)$
be restriction. Denote the Mayer--Vietoris connecting map
by $\partial:H^{j-1}(D;\mathbb Q)\to H^j(Y;\mathbb Q)$.
Writing $v=g_1+g_2$, the collar pullback of a translation
invariant form is
\[
       T^*\omega=\omega+d\eta\wedge\iota_v\omega.
\]
Average its restriction to the torus and change the global
representative by an exact form to make this equality
applicable on the collar. The difference represents the
connecting class, and hence, with the connecting sign
chosen by the increasing collar coordinate,
\begin{equation}
                          T^*-I=\partial\,\iota_v\,r.
 \hypertarget{tag-02D6}{}\label{maps:eq:E.S37}
\end{equation}
As $r\partial=0$, its square is zero.

\item The cyclic deck action on $E_{12}$ is a translation on the
boundary torus, hence acts trivially on boundary cohomology.
Every noninvariant deck eigenspace therefore restricts to
zero. Rational transfer identifies the invariant part
with $H^*(E)$, whose marked restriction is \ref{maps:src:E.2.6}. Its
pullback to \eqref{maps:eq:E.S30} is
\begin{equation}
 R_{k,1}(a)=\phi_k(a)+12(\gamma_2-\gamma_1)a,
 \qquad
 R_{k,2}(b)=\phi_k(b)+12(\gamma_1-\gamma_2)b.
 \hypertarget{tag-02D7}{}\label{maps:eq:E.S38}
\end{equation}
Use the other factor's three covectors $u,w,\delta$,
and take exterior products to obtain the two subspaces
$U_j^k\subset\bigwedge^k\mathbb Q^8$.
The image of $r$ is exactly $U_1^k\cap U_2^k$, and the
kernel of $\partial$ is $U_1^{k-1}+U_2^{k-1}$.
Consequently the rank in \eqref{maps:eq:E.S32} is exactly the rank of
\begin{equation}
 U_1^k\cap U_2^k\ \stackrel{\iota_v}{\longrightarrow}
 \bigwedge^{k-1}\mathbb Q^8/(U_1^{k-1}+U_2^{k-1}).
 \hypertarget{tag-02D8}{}\label{maps:eq:E.S39}
\end{equation}

\item Here are the full dimensions for this small rational
calculation. Let $d_k$ be twice the degree-$k$ dimension
of $H^*(E)\otimes H^*(T^3)$ and let
$r_k=\dim(U_1^k+U_2^k)$:
\begin{equation}
 \begin{array}{c|rrrrrrrrr}
 k&0&1&2&3&4&5&6&7&8\\ \hline
 d_k&2&8&20&44&70&68&36&8&0\\
 r_k&1&8&19&34&47&46&27&8&0
 \end{array}
 \hypertarget{tag-02DA}{}\label{maps:eq:E.S40}
\end{equation}
Contraction in \eqref{maps:eq:E.S39} has rank one only for
$k=4,5$ and is zero in every other degree. These entries result from exterior products of \eqref{maps:eq:E.S38}
and row reduction of the induced contraction maps. One can normalize the remaining sign in
$\phi_2$ by the ring structure: with
$r_1(1)=\theta+12\gamma$, the equality
$r_1(1)r_2(\delta)=r_3(a)$ forces the coefficient
$c_2=-2$, and then $c_3=2$, in the notation of \ref{maps:src:E.2.6}.



\item \emph{The finite-cover cohomology needed for the full Betti numbers.}
The full rational cohomology of $E_{12}$ is computed as follows.
Let $\mathcal L_\lambda$ be the rank-one system on $E$
whose meridian acts by $\lambda$, with $\lambda^{12}=1$.
On a regular torus it is trivial. Its three base transports
in vertical degree $q$ are
\begin{equation}
 \lambda^4\bigwedge^q A_1^{-t},\qquad
 \lambda^{-3}\bigwedge^q A_2^{-t},\qquad
                         \bigwedge^q M_0^{-t}.
 \hypertarget{tag-02DB}{}\label{maps:eq:E.S41}
\end{equation}
The exponents carry the affine meridian character.

\item At a finite cap, averaging on its finite root cover
identifies the cap cohomology with the invariant part of
\eqref{maps:eq:E.S41}. At the unchanged cusp the system is trivial,
because $g$ has character twelve and $u$ has character zero.
Its rational specialization is an isomorphism onto the
monodromy invariants in every degree.  The cusp Betti numbers are
$(1,2,4,2,1)$ by its normalization calculation. In degree
two, the integral surjectivity of the specialization on homology (proof of \cref{maps:src:E.2.5}) gives injectivity on cohomology;
the monodromy-invariant subspace has the same dimension
four. Its image therefore includes $uw$ and $\gamma\delta$.
In degree one the image is $\langle\gamma,u\rangle$;
their products with those two degree-two classes give the
two independent invariant classes in degree three and
the top class in degree four. Thus all the specialization
maps have the asserted ranks, with no additional
skyscraper summand. The proof of \cref{maps:src:E.2.5} gives these specialization assertions integrally in all degrees.

\item It follows that $R^q f_*\mathcal L_\lambda=j_*L_q$ over
the disc containing the three critical values. On this
disc $H^p(j_*L_q)=0$ for $p\geq2$. To check this directly,
retract the punctured disc to a free graph on three loops.
Restriction of a graph cocycle to the three local
coinvariant spaces is onto: its values on the three free
generators can be chosen independently. The exact
restriction sequence therefore kills $H^2(j_*L_q)$.
If $d_q=\binom4q$, $g_q$ is the common fixed dimension of
\eqref{maps:eq:E.S41}, and $k_{i,q}$ are their three individual fixed
dimensions, this same graph calculation gives
\begin{equation}
 \dim H^0(j_*L_q)=g_q,\qquad
 \dim H^1(j_*L_q)=2d_q+g_q-\sum_i k_{i,q}.
 \hypertarget{tag-02DC}{}\label{maps:eq:E.S42}
\end{equation}
The Leray sequence has only these two base degrees and
degenerates. Exact exterior-matrix ranks give the following
table; a row applies to each character of the indicated
order:
\begin{equation}
 \begin{array}{c|rrrrrr}
 \operatorname{ord}(\lambda)&b_0&b_1&b_2&b_3&b_4&b_5\\ \hline
 1 &1&1&4&6&4&0\\
 2 &0&0&4&6&4&0\\
 3 &0&0&3&4&3&0\\
 4 &0&0&3&4&3&0\\
 6 &0&1&5&6&5&1\\
 12&0&1&4&4&4&1
 \end{array}
 \hypertarget{tag-02DD}{}\label{maps:eq:E.S43}
\end{equation}
Here $b_j=\dim H^j(E;\mathcal L_\lambda)$ and $b_6=0$.
For the nontrivial characters all $g_q$ are zero; for the
trivial one they are all one. For \eqref{maps:eq:E.S41}--\eqref{maps:eq:E.S43}, represent the
cyclotomic fields by rational companion matrices and compute each local
fixed space by exact row reduction.
Summing the rows with multiplicities $1,1,2,2,2,4$ gives
\begin{equation}
             (b_0,\ldots,b_6)(E_{12})=(1,7,46,56,46,6,0).
 \hypertarget{tag-02DE}{}\label{maps:eq:E.S44}
\end{equation}
Both sides have Euler characteristic $24$, as required by
the degree-twelve covering of $E$.

\item \emph{The Betti and Jordan calculations on the fibre.}
The return on each of the two pieces of \eqref{maps:eq:E.S29}, up to
isotopy allowing its boundary to move, is its cyclic deck
transformation times the identity of $T^3$. This follows
also by moving the interpolation in \eqref{maps:eq:E.S36} entirely into
the other piece's collar. It acts trivially on the
cohomology of $D$. The generalized eigenvalue-one part of
the Mayer--Vietoris sequence therefore uses just the
invariant part $H^*(E)$ in \eqref{maps:eq:E.S40}. Its dimensions are
\begin{equation}
       (1,0,1,19,45,45,19,1,0,1).
 \hypertarget{tag-02DF}{}\label{maps:eq:E.S45}
\end{equation}
Every other eigenspace is the direct sum of the two piece
eigenspaces, because the boundary has none. In particular
those eigenspaces are semisimple. Their combined Poincar\'e
polynomial is
\begin{equation}
     2(1+t)^3(6t+42t^2+50t^3+42t^4+6t^5),
 \hypertarget{tag-02DG}{}\label{maps:eq:E.S46}
\end{equation}
by \eqref{maps:eq:E.S43}--\eqref{maps:eq:E.S44}. Adding \eqref{maps:eq:E.S45} proves \eqref{maps:eq:E.S33}.

\item Finally the Wang sequence for $P=M_\psi$ and the rational
Betti numbers in \cref{maps:src:E.2.3} give the numbers of Jordan blocks
at eigenvalue one as
\begin{equation}
          (1,0,1,19,44,44,19,1,0,1).
 \hypertarget{tag-02DH}{}\label{maps:eq:E.S47}
\end{equation}
The only possible block sizes are one and two: the
twelfth power is \eqref{maps:eq:E.S31}, and \eqref{maps:eq:E.S37} squares to zero.
Subtracting \eqref{maps:eq:E.S47} from \eqref{maps:eq:E.S45} gives exactly one block
of size two in each of degrees four and five. This agrees with the direct variation ranks \eqref{maps:eq:E.S39}.
All other eigenvalues and blocks are already described
by \eqref{maps:eq:E.S43} and \eqref{maps:eq:E.S46}. \end{enumerate}
The corresponding exact calculations are implemented in \hyperlink{prog:V:cyclic-cover}{the programs in Table~\ref*{tab:V:programs}}.
\end{proof}

\begin{theorem}[Integral fibre groups and the coefficient of the torus twist]\hypertarget{tag-02MF}{}\label{maps:src:E.2.15}
\StatementRecord{maps:src:E.2.15}{proved}{02MF}{\eqref{maps:eq:E.S44}, \eqref{maps:eq:E.S33}, \eqref{maps:eq:E.S42}, \eqref{maps:eq:E.S53}, \eqref{maps:eq:E.S29}, \eqref{maps:eq:E.S38}, \eqref{maps:eq:E.S40}, \eqref{maps:eq:E.S37}, \eqref{maps:eq:E.S55}}


The groups $H^*(E_{12};\mathbb Z)$ are torsion free, with
the ranks in \eqref{maps:eq:E.S44}. For the fibre $Y$, the only
torsion in cohomology is
\begin{equation}
       H^3(Y;\mathbb Z)_{\rm tors}
       =H^7(Y;\mathbb Z)_{\rm tors}=(\mathbb Z/6)^2.
 \hypertarget{tag-02DI}{}\label{maps:eq:E.S48}
\end{equation}
Equivalently, its integral homology has the ranks \eqref{maps:eq:E.S33}
and torsion $(\mathbb Z/6)^2$ exactly in degrees two and six.
The twelfth return $T=\psi^{12}$ acts as the identity on
all integral cohomology except in degrees four and five.
In each of those two free groups it has, in an integral
basis, the form
\begin{equation}
             \begin{pmatrix}1&72\\0&1\end{pmatrix}
                                  \oplus I_{691}.
 \hypertarget{tag-02DJ}{}\label{maps:eq:E.S49}
\end{equation}

For $n\geq1$, let $P_{12n}$ be the connected cyclic cover
of $P$ of degree $12n$ defined by $\chi$. Its Betti numbers,
independent of $n$, are
\begin{equation}
       (1,13,133,528,1099,1384,1099,528,133,13,1).
 \hypertarget{tag-02DL}{}\label{maps:eq:E.S50}
\end{equation}
Its integral homology torsion is exactly
\begin{equation}
 H_j(P_{12n};\mathbb Z)_{\rm tors}=
 \begin{cases}
       (\mathbb Z/6)^2,&j=2,3,6,7,\\
       \mathbb Z/(72n),&j=4,5,\\
       0,&\text{otherwise}.
 \end{cases}
 \hypertarget{tag-02DM}{}\label{maps:eq:E.S51}
\end{equation}
In particular these covers are pairwise nonhomotopy-equivalent.
\end{theorem}


\begin{proof}
\begin{enumerate}[label=\arabic*.,leftmargin=*]
\item \emph{The integral cohomology of the cover.}
Apply the proof of \eqref{maps:eq:E.S42} with the coefficient lattice
$\mathbb Z[C_{12}]$, using its regular permutation matrix
$P_{12}$. In vertical degree $q$ set
\begin{equation}
 \begin{aligned}
 T_{1,q}=P_{12}^4\otimes\bigwedge^q A_1^{-t},\quad
 &T_{2,q}=P_{12}^{-3}\otimes\bigwedge^q A_2^{-t},\\
 T_{0,q}=I_{12}\otimes\bigwedge^q M_0^{-t},
 &V_q=\mathbb Z^{12\binom4q}.
 \end{aligned}
 \hypertarget{tag-02DN}{}\label{maps:eq:E.S52}
\end{equation}
This coefficient system computes $H^*(E_{12};\mathbb Z)$.
 At the first finite
cap the cover has four components, each the ordinary
degree-three torus root cover; at the second there are
three ordinary degree-four torus root covers. An invariant
vector is determined by its value on one sheet of such
an orbit, with the other values its integral monodromy
translates. This identifies cap cohomology with the full
primitive invariant lattice. The cusp has twelve unchanged
components and the integral specialization of the proof of \cref{maps:src:E.2.5}.
Consequently its direct-image sheaf is exactly the integral
$j_*L_q$, including at every critical point.

\item The three graph-cocycle values again surject onto the
three local coinvariant groups, integrally. Thus $H^2(j_*L_q)=0$,
and its remaining cohomology is the kernel and cokernel
of the exact lattice map
\begin{equation}
 \begin{gathered}
 V_q\longrightarrow
        \bigoplus_{i=0}^2\operatorname{im}(T_{i,q}-I),\\
           v\longmapsto((T_{i,q}-I)v)_i.
 \end{gathered}
 \hypertarget{tag-02DO}{}\label{maps:eq:E.S53}
\end{equation}
Use Hermite column bases for the three image lattices. The source and target ranks,
and the rank of \eqref{maps:eq:E.S53}, are respectively
\begin{equation}
 \begin{array}{c|rrrrr}
 q&0&1&2&3&4\\ \hline
 \operatorname{rank}V_q&12&48&72&48&12\\
 \operatorname{rank}\bigoplus\operatorname{im}(T_{i,q}-I)
                      &17&92&126&92&17\\
 \operatorname{rank}\eqref{maps:eq:E.S53}&11&47&71&47&11
 \end{array}
 \hypertarget{tag-02DP}{}\label{maps:eq:E.S54}
\end{equation}
Every nonzero Smith entry is one. Hence $H^0(j_*L_q)$
has rank one and $H^1(j_*L_q)$ is free with ranks
$(6,45,55,45,6)$. The two-column Leray filtration has
free end groups and therefore splits as abelian groups.
This proves the torsion-free assertion and \eqref{maps:eq:E.S44}
integrally. The integer matrices are determined by the Hermite image bases in
\eqref{maps:eq:E.S53}; their Smith reductions give the displayed invariants.

\item \emph{The integral images.}
Each piece $A=E_{12}\times T^3$ in \eqref{maps:eq:E.S29} has free
integral homology. By Poincar\'e--Lefschetz duality,
$H^{k+1}(A,\partial A;\mathbb Z)$ is free. The cokernel
of $H^k(A;\mathbb Z)\to H^k(\partial A;\mathbb Z)$
embeds in that free group, so the boundary image is a
primitive lattice. Its rational span has already been
computed in \eqref{maps:eq:E.S38}. Therefore the integral
image is precisely the saturation of that displayed
image in $\bigwedge^k\mathbb Z^8$.

\item Choose saturated column bases $\overline A_k,\overline B_k$
for the two images. The combined matrix
$[\overline A_k\ \overline B_k]$ has the ranks in the
second row of \eqref{maps:eq:E.S40}. All its nonzero Smith entries are
one except that, in each of degrees two and six, exactly
two entries are six. The Mayer--Vietoris exact sequence
has a free kernel term on its right, since all piece
cohomology is free. Thus its abelian extension splits,
and these Smith entries give exactly \eqref{maps:eq:E.S48}. The free
ranks remain \eqref{maps:eq:E.S33}.

\item \emph{The integral variation, including possible torsion components.}
Formula \eqref{maps:eq:E.S37} holds integrally. One can use the singular
prism of the loop of translations instead of differential
forms: the action map $S^1\times T^8\to T^8$ has
cohomology formula
\[
 a\longmapsto 1\otimes a+[S^1]^*\otimes\iota_v a.
\]
The difference between the two collar homotopies is
exactly this loop, so its connecting class is
$\partial\iota_v r(a)$. This proves the integral formula
also on torsion and still gives $(T^*-I)^2=0$.

\item For complete lattice control, take an integral kernel
basis of $[\overline A_k\ -\overline B_k]$. Its first
block, multiplied by $\overline A_k$, is a basis of the
common boundary lattice. Contract it by $v$ and reduce
modulo $[\overline A_{k-1}\ \overline B_{k-1}]$ using
the integral Smith transforms. In degrees four and five
the result has rank one and unique nonzero Smith entry
$72$. In degree three its two possible components
modulo six are both zero; in degree seven the common
boundary image is zero. All the remaining reductions
are zero. The reductions use the integral unimodular Smith transforms.
In degrees four and five the connecting cokernel is
free and embeds as a primitive summand of $H^k(Y)$,
because the following Mayer--Vietoris kernel is free.
Thus the divisibility in the whole cohomology group is $72$.

\item A rank-one integral map of Smith entry $72$ has the
form $72u\ell$, where $u$ and $\ell$ are primitive.
Its square is zero precisely when $\ell(u)=0$.
Extend $u$ to a basis of $\ker\ell$ and choose a vector
on which $\ell$ is one. This gives \eqref{maps:eq:E.S49}. The same
calculation proves identity on every other integral
cohomology group, including its torsion.

\item \emph{The finite-cover groups and the two potentially nontrivial extensions.}
The cover $P_{12n}$ is the mapping torus of $T^n$.
Since $(T^*-I)^2=0$, its variation is $n(T^*-I)$,
with Smith entry $72n$ in degrees four and five.
For every $k$ the integral Wang sequence gives
\begin{equation}
 0\longrightarrow\operatorname{coker}(T^{n*}-I)_{k-1}
 \longrightarrow H^k(P_{12n};\mathbb Z)
 \longrightarrow\ker(T^{n*}-I)_k\longrightarrow0.
 \hypertarget{tag-02DQ}{}\label{maps:eq:E.S55}
\end{equation}
The ranks immediately give \eqref{maps:eq:E.S50}. Whenever the right
group is free, \eqref{maps:eq:E.S55} splits as an abelian sequence.
In particular degrees four and eight have torsion
$(\mathbb Z/6)^2$, contributed by $H^3(Y)$ and $H^7(Y)$
respectively, while degrees five and six have torsion
$\mathbb Z/(72n)$. Every other such degree has no torsion.

\item The only right groups in \eqref{maps:eq:E.S55} that are not free
occur in degrees three and seven.  Instead use the
nondegenerate integral torsion linking pairing of the
closed oriented ten-manifold $P_{12n}$: cohomological
torsion in degree $k$ is Pontryagin dual to that in
degree $11-k$. The already fixed torsion in degree
eight therefore forces the degree-three torsion to
be $(\mathbb Z/6)^2$, and the fixed degree-four torsion
forces the same answer in degree seven. The free
ranks are already known. This resolves precisely
the two possibly absorbing extensions. The universal
coefficient theorem now translates these cohomology
groups to \eqref{maps:eq:E.S51}. The order $72n$ in the middle
groups distinguishes the covers for distinct $n$.
\end{enumerate}
The corresponding exact calculations are implemented in \hyperlink{prog:V:cyclic-cover}{the programs in Table~\ref*{tab:V:programs}}.
\end{proof}

\Needspace{9\baselineskip}
\section{The homotopy module and the quadratic cusp attachment}\label{pub:V:sec-homotopy-module}

\begin{theorem}[The full complement module and its peripheral map]\hypertarget{tag-02MG}{}\label{maps:src:E.2.16}
\StatementRecord{maps:src:E.2.16}{proved}{02MG}{\eqref{maps:eq:E.S1}, \cref{maps:src:E.2.17}, \cref{maps:src:E.2.21}, \cref{maps:src:E.2.12}, \eqref{maps:eq:E.C5}, \eqref{maps:eq:E.C8}, \eqref{maps:eq:E.C11}, \eqref{maps:eq:E.C7}, \eqref{maps:eq:E.C14}, \eqref{maps:eq:E.C12}, \eqref{maps:eq:E.C16}, \eqref{maps:eq:E.C17}, \cref{maps:src:E.2.6}}


Let $K=\pi_1(E)$ have the marking \eqref{maps:eq:E.S1}, and put
\[
 R=\mathbb ZK,\qquad Q=K/\langle h\rangle=C_3*C_4,
 \qquad \overline R=\mathbb ZQ.
\]
Use right deck modules; conversion to left modules applies the
involution $g\mapsto g^{-1}$ to every coefficient. Then
\begin{equation}
 \pi_2(E)\simeq R\oplus\overline R^{\,2}\oplus L,
 \qquad
 L=\langle x,t\mid t(a-1)=t(b-1)=0,
                         \ x(1+b)+3t=0\rangle_{\overline R}.
 \hypertarget{tag-02A4}{}\label{maps:eq:E.C1}
\end{equation}
Equivalently, replacing $x$ by $x+t$ and eliminating $t$ gives
\begin{equation}
 L\simeq\overline R\big/
       \big((b^2-1)\overline R+(1+b)(a-1)\overline R\big).
 \hypertarget{tag-02AG}{}\label{maps:eq:E.C2}
\end{equation}
There is a nonsplit exact sequence
\begin{equation}
 0\longrightarrow\mathbb Z\langle t\rangle
  \longrightarrow L\longrightarrow
  \mathbb Z_{-}\otimes_{\mathbb ZC_4}\overline R
  \longrightarrow0,
 \hypertarget{tag-02AJ}{}\label{maps:eq:E.C3}
\end{equation}
where $b$ acts as $-1$ on $\mathbb Z_{-}$ and the first term
has trivial $Q$ action. Its extension class is the nonzero
element of an order-two extension group. In the connected
lift of the peripheral $T^5$ to the universal cover of $E$,
the degree-two classes $uw,ud,wd$ map respectively to
\begin{equation}
                           t,\quad0,\quad0.
 \hypertarget{tag-02AK}{}\label{maps:eq:E.C4}
\end{equation}
Thus the central torsion summand of the module is annihilated
by $h-1$ itself. Formula \eqref{maps:eq:E.C1} is a full integral module
calculation, including the central value $h=1$.  There is also a chain-level peripheral comparison.
For the splitting constructed below, let $e$ generate the
free $R$ summand and let $y,z$ generate the two copies of
$\overline R$. The projection of the peripheral degree-two
chain map to \eqref{maps:eq:E.C1} is
\begin{equation}
 \begin{array}{c|rrrrrr}
 \text{cell}&gu&gw&gd&uw&ud&wd\\ \hline
 \text{value}&0&0&6t&t&0&0
 \end{array}
 \hypertarget{tag-02AE}{}\label{maps:eq:E.C19a}
\end{equation}
and, writing $\Theta$ for the outer base circle with
group image $\mu=ab$, use the cone convention: the
notation $v\Theta$ denotes the suspended cell
$\Theta\times v$, with the interval factor first.
In this convention,
\begin{equation}
 \begin{array}{c|l}
  g\Theta&eh-2x-6yb+z(5-2b)\\
  u\Theta&x-yb+zb\\
  w\Theta&-y(1+b)+z\\
  d\Theta&x+2t+z(b-1).
 \end{array}
 \hypertarget{tag-02AF}{}\label{maps:eq:E.C19b}
\end{equation}
The tables give the chain map; its induced map on lifted peripheral homology is \eqref{maps:eq:E.C4}. In
particular the free coefficient $h$ on $g\Theta$
is nonzero even though the free summand receives
zero from the lifted peripheral homology in \eqref{maps:eq:E.C4}.


\end{theorem}


\begin{proof}
\begin{enumerate}[label=\arabic*.,leftmargin=*]
\item We give the local chain comparison underlying \eqref{maps:eq:E.C1}.

\item \emph{The cusp boundary and its derived attachment.}
Write $A=\mathbb Z[h,h^{-1}]$ and $q=h-1$. Let
$\widehat W$ be the gamma cover used in \ref{maps:src:E.2.12}, and let
$\widehat M$ be the corresponding boundary cover. Ordinary
chains with the deck action give finite free $A$ chain
models. The homology of $\widehat W$ is
\begin{equation}
 H_0=H_1=A/(q),\quad H_2=A\oplus(A/(q))^2,
 \quad H_3=0,\quad H_4=A.
 \hypertarget{tag-02AL}{}\label{maps:eq:E.C5}
\end{equation}
For the boundary, start with the torus chain complex
$C= A\otimes\bigwedge\langle g,u,w,d\rangle$,
whose differential contracts $g$ and multiplies by $q$.
The cellular monodromy has degree-one matrix
\begin{equation}
 F(M_0)g=g-hd,\quad F(M_0)u=u+w,
 \quad F(M_0)w=w,\quad F(M_0)d=d.
 \hypertarget{tag-02AM}{}\label{maps:eq:E.C6}
\end{equation}
Its higher matrices are exterior powers. The mapping-torus
complex is the cone of $I-F(M_0)$, with differential
\begin{equation}
 d(c,e)=\big(dc+(I-F(M_0))e,-de\big).
 \hypertarget{tag-02AN}{}\label{maps:eq:E.C7}
\end{equation}
For example $d(g\theta)=hd-q\theta$ and
$d(u\theta)=-w$. Cancelling these two unit pairs leaves
the degree-one classes $u$ and $\theta$, with resolutions
$q$ and a unit times $q^2$. In degree two the four surviving
classes are $uw,ud,w\theta,d\theta$, each with resolution
$q$. In particular
\begin{equation}
 H_0(\widehat M)=A/(q),\quad
 H_1(\widehat M)=A/(q^2)\oplus A/(q),\quad
 H_2(\widehat M)=(A/(q))^4.
 \hypertarget{tag-02AO}{}\label{maps:eq:E.C8}
\end{equation}
Geometrically this is the Heisenberg three-manifold times
a circle: the deck transformation shears the meridian
by the $d$ circle. The homology map to \eqref{maps:eq:E.C5} is the
identity in degree zero, projection to $u$ in degree one,
and the primitive projection
$(uw,ud,w\theta,d\theta)\mapsto(U,D,0,0)$ in degree two.
The last assertion uses the integral specialization of the proof of \cref{maps:src:E.2.5} and the collapsed meridian tori.

\item The extension is determined by a relative intersection calculation. Let $S_i$ denote the normalized hexagon
component with gamma index $i$. In its standard Picard
basis $H,E_1,E_2,E_3$, the six boundary curves are
\begin{equation}
 E_1,\ H-E_1-E_2,\ E_2,\ H-E_2-E_3,
                       \ E_3,\ H-E_1-E_3.
 \hypertarget{tag-02AP}{}\label{maps:eq:E.C9}
\end{equation}
The first opposite pair is self-glued in the surviving
$u$ direction; the other two pairs join $S_i$ to its
two gamma neighbours. Put $C_0=-H+E_1$, the primitive
generator of the free kernel of the normalization matrix
$B(h)$ in \ref{maps:src:E.2.12}. The divisor $S_i$ restricts on its
own normalization to $2C_0$: the total central divisor
is principal, and the two self-gluing preimages cancel
their terms in the normal restriction. It restricts
to $-C_0$ on each neighbouring component and to zero
on every other component. Thus the free part of the
equivariant intersection map is
\begin{equation}
 H_4(\widehat W)=A\longrightarrow
 H_4(\widehat W,\widehat M)\longrightarrow A,
 \qquad 2-h-h^{-1}=-h^{-1}q^2.
 \hypertarget{tag-02A5}{}\label{maps:eq:E.C10}
\end{equation}
Poincar\'e--Lefschetz duality over $A$, with its usual
involution, gives the same coefficient, up to a unit,
for the free part of $H_2(\widehat W)\to
H_2(\widehat W,\widehat M)$. The latter relative group
is $A$: the universal-coefficient sequence in degree
four has $\operatorname{Hom}_A(H_4,A)=A$, zero $H_3$,
and no higher extension terms. Indeed every module
in \eqref{maps:eq:E.C5} has a projective resolution of length at
most one. The relative sequence therefore includes
\begin{equation}
 0\longrightarrow A\xrightarrow{\,\text{unit}\ q^2\,}A
  \longrightarrow A/(q^2)\longrightarrow0.
 \hypertarget{tag-02A6}{}\label{maps:eq:E.C11}
\end{equation}
Consequently the boundary map has a unit component in
$\operatorname{Ext}^1_A(A/(q^2),A)=A/(q^2)$.
Every unit of $A/(q^2)=\mathbb Z[\varepsilon]/(\varepsilon^2)$
lifts to a unit $\pm h^n$ of $A$.

\item This determines the low-degree chain map, as follows. All the
modules in \eqref{maps:eq:E.C5} and \eqref{maps:eq:E.C8} have projective dimension
at most one. Their complexes split in the derived
category into their two-term projective resolutions:
the obstructions to splitting consecutive truncations
are extension groups in degree at least two and vanish.
In degrees relevant to $H_2$ of a mapping cone, a map
between these split complexes consists of the displayed
homology maps and extension components from $H_i$ to
$H_{i+1}$. Normalize the component in \eqref{maps:eq:E.C11} by a unit
on the free $H_2$ summand of the target. Its components
to the two torsion $H_2$ summands can then be removed
by target maps $A\to(A/(q))^2$: reduction of the unit
modulo $q$ is $\pm1$. The remaining components from
the $u$ summand are removed by target degree-one
extension maps from $H_1(\widehat W)$ to
$H_2(\widehat W)$. Finally the degree-zero component
is removed using the isomorphism on $H_0$ and a target
extension map from $H_0$ to $H_1$.

\item Each change just described is a unit on the diagonal
of the split complex and is triangular with extension
entries off the diagonal. It is therefore invertible.
Replacing the split resolutions by the original finite
free complexes represents it by a chain homotopy
equivalence.  The change
is applied to the cusp target and its entire incoming
boundary map simultaneously. The other vertices and
all maps into the regular core are unchanged. It thus
gives a chain homotopy equivalence of the global
mapping-cone complex that is the identity on the
external peripheral complex. Tensoring with $R$ preserves these comparisons,
since $R$ is free over $A$.

\item One resulting cusp model in degrees zero through three
has ranks $(1,2,4,2)$, generators
\[
 (1);\quad(g,u);\quad(gu,E,U,D);\quad(R_U,R_D),
\]
and differentials $dg=q$, $d(gu)=qu$,
$dR_U=qU$, $dR_D=qD$. The map \eqref{maps:eq:E.C7} to this model
is, through degree two,
\begin{equation}
 \begin{array}{c|ccccc}
  &g&u&w&d&\theta\\ \hline
  i_1&g&u&0&0&0
 \end{array},\qquad
 \begin{array}{c|rrrrrr}
  &gu&gw&gd&uw&ud&wd\\ \hline
  i_2&gu&-hD&E&U&D&0
 \end{array},
 \qquad i_2(v\theta)=0.
 \hypertarget{tag-02A7}{}\label{maps:eq:E.C12}
\end{equation}
The signs and the factor $h$ satisfy the chain equation
directly. At $h=1$, choose
$E=-\overline C_1-\overline C_2$,
$U=\overline C_2+\overline C_3$, and
$D=\overline C_2$.

\item \emph{Finite caps and the global complex.}
For any marked monodromy
$A=\left(\begin{smallmatrix}1&0\\t&D\end{smallmatrix}\right)$
use the chain lift
$F(A)=\left(\begin{smallmatrix}1&0\\ht&D\end{smallmatrix}\right)$
and its exterior powers. These lifts multiply exactly
and commute with contraction by $qg^*$.

\item This exterior lift represents the equivariant
torus map. The gamma cover is
$\mathbb R_g\times T^3$, and the marked linear map is
$(g,v)\mapsto(g,Dv+gt)$. At every integer value of $g$
the translation is integral, so the vertex torus map
is exactly $D$. In a single unit strip, concentrate
the loop of translations $v\mapsto v+st$ at its
terminal endpoint by an equivariant cellular
approximation relative to the two vertex tori.
The translation prism on the minimal cellular complex
of $T^3$ is exterior multiplication by $t$: its value
on each product torus is the homology class of the
literal translation sweep, and all differentials
in that minimal torus complex are zero. The terminal
vertex has deck coefficient $h$. Consequently
\begin{equation}
 v\longmapsto\bigwedge D(v),\qquad
 g\wedge v\longmapsto
 g\wedge\bigwedge D(v)+ht\wedge\bigwedge D(v),
 \hypertarget{tag-02AH}{}\label{maps:eq:E.C20}
\end{equation}
which is exactly the exterior lift above. There is
only this one strip parameter. Composing two such
maps adds the translation loops by
$t_A+D_At_B$, the same strict multiplication law
as $F(A)F(B)=F(AB)$. 
A regular
boundary with deck image $s$ is the cone of $sI-F(A)$.
This convention gives the group relation
$s\lambda s^{-1}=A\lambda$.

\item At finite cap $j$, take its adapted first fibre vector
$n_jv_j$, where
\begin{equation}
 (m_1,n_1,v_1)=(3,1,(1,2,-4,0)),\qquad
 (m_2,n_2,v_2)=(4,-1,(-1,-3,3,0)).
 \hypertarget{tag-02A8}{}\label{maps:eq:E.C13}
\end{equation}
It has gamma coordinate one and is fixed by $A_j$.
The other three coordinates form $L_0=\langle u,w,d\rangle$,
with monodromy $D_j=A_j|_{L_0}$. The cap is the
mapping torus of this three-torus, and its complex has
differential $sI-\bigwedge^rD_j$ on the swept cells.
The fibre-to-cap map is the identity on $\bigwedge L_0$;
on the adapted long-circle factor it is the integral
norm matrix $N_{m_jn_j}(s,D_j)$ determined by
\begin{equation}
 (sI-D_j)N_{m_jn_j}(s,D_j)=(h-1)I,
 \qquad s=a\text{ or }b.
 \hypertarget{tag-02A9}{}\label{maps:eq:E.C14}
\end{equation}
For positive $n_j$ the norm is
$\sum_{r=0}^{m_j-1}s^rD_j^{m_j-1-r}$, and for
$n_j=-1$ it is minus $s^{-m_j}$ times that sum.
In phase degree $r$, the map is
$N_{m_jn_j}(s,\bigwedge^rD_j)$.
The boundary base-circle generator maps to the cap circle, while
the cell containing both that meridian and the adapted
long circle maps to zero. These are winding
maps in the mapping-torus coordinates. The linear three-torus return fixes its origin.

\item Form the core over a bouquet of three circles with
deck images $a,b,1$ and monodromies $A_1,A_2,M_0$.
Attach the two finite complexes by \eqref{maps:eq:E.C14} and the
cusp by \eqref{maps:eq:E.C12}, using their double mapping cylinders.
Through degree three the resulting free $R$ complex
has ranks $(4,20,49,62)$. Write its group-ring entries in the integral normal form $h^r$ times
an alternating word in $a,a^2$ and $b,b^2,b^3$; exact unit
cancellations use multiplication in that normal form.

\item The final reduction is as follows. Cancelling degree-zero
and degree-one unit pairs leaves ranks $(1,2,34,62)$,
with $d_1=(a-1,b-1)$. The only nonzero columns of
$d_2$ are $F$ and $F(b-1)$, where
\begin{equation}
 F=\begin{pmatrix}1+a+a^2\\h(1+b+b^2+b^3)\end{pmatrix}.
 \hypertarget{tag-02AA}{}\label{maps:eq:E.C15}
\end{equation}
The first component of $F$ is injective on $R$:
multiplication by $a-1$ converts it to $h-1$, and
$R$ is free over the Laurent ring $A$. Thus removing
the $F$ column after subtracting its multiple gives
a free basis of $\ker d_2$ of rank $33$. The corresponding
row of $d_3$ is zero by the chain equation. Unit
reductions of the remaining presentation leave five
generators $(x,t,y,z,e)$. One generator $e$ occurs
in no relation. Three pure $t$ relations are
\begin{equation}
 tq=0,\qquad t(1-a^{-1})=0,
                  \qquad t(b-a^{-1})=0.
 \hypertarget{tag-02AB}{}\label{maps:eq:E.C16}
\end{equation}
They generate exactly the augmentation relations on
$t$: multiply the second by $a$ to get $t(a-1)=0$,
and subtract it from the third to get $t(b-1)=0$.
Every other pure $t$ coefficient has augmentation
zero and is consequently redundant. Reducing the
remaining $t$ coefficients by augmentation, three
relations give successively $yq=0,zq=0,xq=0$.
The final three relations are
\begin{equation}
 x(1+b)+3t=0,\quad x(b+b^2)+3t=0,
                   \quad x(1+b+b^2+b^3)+6t=0.
 \hypertarget{tag-02AC}{}\label{maps:eq:E.C17}
\end{equation}
The last two are respectively the first multiplied
by $b$, and the first multiplied by $1+b^2$.
These identities prove \eqref{maps:eq:E.C1}.

\item The invariant subgroup generated by $t$ injects in
$L$. One way to see this is to start with the exact
sequence of right $\mathbb ZC_4$ modules
\[
 0\longrightarrow\mathbb Z\longrightarrow
 \mathbb Z[C_4/C_2]\longrightarrow\mathbb Z_-\longrightarrow0,
\]
where the first map is the sum of the two cosets.
Induce to $Q$ and push out its first term by the
augmentation to the trivial module $\mathbb Z$.
The resulting presentation is \eqref{maps:eq:E.C1}, after the
change $x\mapsto x+t$ which replaces the coefficient
$3$ by $1$. Pushout preserves the displayed injection.
Restriction to $C_4$ shows nonsplitting: an invariant
correction of the lift $x$ changes this odd coefficient
by an even number. Equivalently the standard cyclic
resolution computes
$\operatorname{Ext}^1_{\mathbb ZC_4}(\mathbb Z_-,\mathbb Z)
=\mathbb Z/2$. Induction--restriction gives the same
extension group for \eqref{maps:eq:E.C3}.

\item Finally insert the fibre cycles $uw,ud,wd$ in
the core before any cancellation and carry their
coordinate columns through the same basis changes.
The five-generator presentation gives respectively
\begin{equation}
 t,\qquad t(a^{-1}-a),\qquad t(a^{-1}-1).
 \hypertarget{tag-02AD}{}\label{maps:eq:E.C18}
\end{equation}
Relations \eqref{maps:eq:E.C16} reduce these to \eqref{maps:eq:E.C4}. The external
boundary cover has homotopy type $T^3$ and these are
its three degree-two generators. Taking $K$ coinvariants gives $2x+3t=0$, so $t$ is twice a
primitive class, exactly the coefficient $2$ in \ref{maps:src:E.2.6}.

\item For the remaining peripheral data, use the
boundary word $abc$, where the cusp generator $c$
has deck image $1$. On swept fibre chains its map to
the three core edges has blocks
\begin{equation}
                 bI,\qquad F(A_1),\qquad F(A_1)F(A_2).
 \hypertarget{tag-02AI}{}\label{maps:eq:E.C21}
\end{equation}
The identity
\[
 (aI-F(A_1))b+F(A_1)(bI-F(A_2))
       +F(A_1)F(A_2)(I-F(M_0))=(ab-1)I
\]
checks its chain equation. This is the usual prism
subdivision of the peripheral word, with
its given ribbon marking. Tracking all ten
degree-two columns through the same reductions,
and using only $t(a-1)=t(b-1)=0$ and
$xq=yq=zq=0$, gives \eqref{maps:eq:E.C19a}--\eqref{maps:eq:E.C19b}.

\end{enumerate}
The corresponding exact calculations are implemented in \hyperlink{prog:V:quadratic}{the programs in Table~\ref*{tab:V:programs}}.
\end{proof}

\begin{theorem}[The full shifted homotopy module]\hypertarget{tag-02MH}{}\label{maps:src:E.2.17}
\StatementRecord{maps:src:E.2.17}{proved}{02MH}{\cref{maps:src:E.2.16}, \cref{maps:src:E.2.13}, \eqref{maps:eq:E.P7}, \eqref{maps:eq:E.P8}, \eqref{maps:eq:E.P11}, \cref{maps:src:E.2.9}, \cref{maps:src:E.2.10}, \eqref{maps:eq:E.P9}, \eqref{maps:eq:E.P14}, \eqref{maps:eq:E.P16}, \cref{maps:src:E.2.11}}


Use the right-module convention of \ref{maps:src:E.2.16}. Write
\begin{equation}
 \begin{gathered}
 R=\mathbb ZG,\qquad q_i=h_i-1,\qquad S=\mathbb Z\Gamma,\\
 \Gamma=G/\langle h_1,h_2\rangle
  =\langle a,b,c,d\mid a^3=b^4=c^3=d^4=abcd=1\rangle.
 \end{gathered}
 \hypertarget{tag-02BJ}{}\label{maps:eq:E.P1}
\end{equation}
where $a,b$ belong to the first complement, $c,d$ to
the second, and $\mu=ab=(cd)^{-1}$. The full integral
second-homotopy module is
\begin{equation}
 \pi_2(P)\simeq R/q_2R\ \oplus\ R/q_1R
                    \ \oplus\ T_L\oplus T_R.
 \hypertarget{tag-02BS}{}\label{maps:eq:E.P2}
\end{equation}
The last two modules are inflated from the following
explicit six-generator presentations over $S$.
Put $N_a=1+a+a^2$, $N_b=1+b+b^2+b^3$, and define
$N_c,N_d$ similarly.

For $T_L$ use generators $p_i,r_i$ for $0\leq i\leq2$,
put $v_i=p_i-r_i$, and define
\begin{equation}
 \begin{split}
 \Omega_L={}&v_0(1+d-d^2-d^3)
       +v_1(d^3-d)+v_2(d^2+d^3),\\
 p_iN_a={}&0,\qquad r_iN_b=0\quad(0\leq i\leq2),\\
 \Omega_L(a-1)={}&\Omega_L(b-1)
       =\Omega_L(c-1)=\Omega_L(d-1)=0.
 \end{split}
 \hypertarget{tag-02BT}{}\label{maps:eq:E.P3}
\end{equation}
For $T_R$ use generators $B_i,C_i$, put $w_i=B_id+C_i$,
and define
\begin{equation}
 \begin{split}
 \Omega_R={}&w_0a(1-b-b^2+b^3)
              +w_1a(b^2-1)-w_2a(1+b^3),\\
 B_iN_c={}&0,\qquad C_iN_d=0\quad(0\leq i\leq2),\\
 \Omega_R(a-1)={}&\Omega_R(b-1)
       =\Omega_R(c-1)=\Omega_R(d-1)=0.
 \end{split}
 \hypertarget{tag-02BU}{}\label{maps:eq:E.P4}
\end{equation}
These are complete presentations: six generators and
ten relations in each block. In particular
\begin{equation}
 q_1q_2\pi_2(P)=0,\qquad
 \ker(q_1:\pi_2(P)\to\pi_2(P))
  \cap\ker(q_2:\pi_2(P)\to\pi_2(P))=T_L\oplus T_R.
 \hypertarget{tag-02BV}{}\label{maps:eq:E.P5}
\end{equation}
This sharpens the annihilator of \cref{maps:src:E.2.12} and determines the central-intersection summand.
Its coinvariants are $(\mathbb Z/12)^6$, so \eqref{maps:eq:E.P2}
recovers $\pi_2(P)_G=\mathbb Z^2\oplus(\mathbb Z/12)^6$.

Let $\Pi=\pi_1(Y)$ and let $\Sigma\simeq\pi_1(\Sigma_6)$
be its quotient by $z=h_1h_2$, as in \ref{maps:src:E.2.13}. Then
\begin{equation}
 \pi_2(Y)\simeq(\mathbb Z\Pi)^{24}
          \oplus\operatorname{Res}^{\Gamma}_{\Sigma}(T_L\oplus T_R),
 \qquad
 \pi_2(Y)_\Pi\simeq\mathbb Z^{110}\oplus(\mathbb Z/6)^2.
 \hypertarget{tag-02BW}{}\label{maps:eq:E.P6}
\end{equation}
The surface-group summand is inflated to $\Pi$, so
$z$ acts trivially on it. \end{theorem}


\begin{proof}
\begin{enumerate}[label=\arabic*.,leftmargin=*]
\item \emph{A small chain presentation with the peripheral map.}
The calculation in \ref{maps:src:E.2.16} permits the following model
through degree three for the first complement. Its
bases have ranks $(1,2,6,6)$, with degree-one basis
$(a,b)$ and degree-two basis $(r,x,t,y,z,e)$.
The first two differentials are
\begin{equation}
 d_1=(a-1,b-1),\qquad
 d_2(r)=\begin{pmatrix}N_a\\h_1N_b\end{pmatrix},
 \quad d_2(x)=d_2(t)=d_2(y)=d_2(z)=d_2(e)=0.
 \hypertarget{tag-02BX}{}\label{maps:eq:E.P7}
\end{equation}
The six degree-three columns are the six relations
\begin{equation}
 t(a-1),\quad t(b-1),\quad xq_1,\quad yq_1,\quad zq_1,
                         \quad x(1+b)+3t.
 \hypertarget{tag-02BY}{}\label{maps:eq:E.P8}
\end{equation}
The degree-two norm $N_a$ in \eqref{maps:eq:E.P7} is the polynomial $1+a+a^2$ of \eqref{maps:eq:E.P3}, evaluated in $R$ with $a^3=h_1$. The second complement is the identical
model with $(a,b,h_1)$ replaced by $(c,d,h_2)$.

\item The ordinary one-relator complex $C_K$ is the subcomplex
on degrees zero and one and the cell $r$. It is an
aspherical model: it is the double mapping cylinder
of the degree-three and degree-four maps of circles,
and its universal cover is a tree of contractible
lines and strips. Equivalently it is the standard
Seifert spine of the torus-knot exterior. Projection
of \eqref{maps:eq:E.P7}--\eqref{maps:eq:E.P8} onto $C_K$ retains the peripheral map. In degree one this map is
\begin{equation}
 g\longmapsto(N_a,0),\qquad
 u,w,d\longmapsto0,\qquad
 \Theta\longmapsto(b,1).
 \hypertarget{tag-02BZ}{}\label{maps:eq:E.P9}
\end{equation}
Its degree-two $r$ coefficient is $b-1$ on the swept
cell $g\Theta$ and zero on the other nine cells.
Its remaining components are exactly \eqref{maps:eq:E.C19a}--\eqref{maps:eq:E.C19b}.
\item Here and below the small complex is only a
presentation through the degree needed for $H_2$.
Replacing redundant degree-three columns by the six
generators of their image in \eqref{maps:eq:E.P8} changes no
degree-two cycle or boundary. Reducing a peripheral
degree-two column modulo that image changes it by
a chain homotopy with values in degree three.
Products with a torus and the mapping cones below
therefore retain the same degree-two homology.
For the comparison with an aspherical target one
may first keep the original, full complex: its
projection to $C_K$ extends by zero in degrees at
least three, since the Fox column in \eqref{maps:eq:E.P7} is
injective and hence the $r$ row of $d_3$ is zero.
Only after taking the kernel is its image in degree
two shortened as in \eqref{maps:eq:E.P8}. 

\item \emph{The exact kernel complex for the splice.}
Let $T(s)$ be the four-torus chain complex whose
first circle has deck image $s$ and whose other
three circles have trivial image. Its differential
contracts that first circle and multiplies by $s-1$.
The two vertex complexes for the splice are
\begin{equation}
 V_1=C_{E_1}\otimes T(h_2),\qquad
 V_2=C_{E_2}\otimes T(h_1).
 \hypertarget{tag-02BK}{}\label{maps:eq:E.P10}
\end{equation}
We have reversed the two product factors in the
second vertex only for notation; the Koszul signs
below account for that reversal. The boundary
complex $B$ is the nine-torus complex with ordered
coordinates
\begin{equation}
 (g_1,u_1,w_1,d_1,g_2,u_2,w_2,d_2,\Theta),
 \quad\text{deck images }(h_1,1,1,1,h_2,1,1,1,\mu).
 \hypertarget{tag-02BL}{}\label{maps:eq:E.P11}
\end{equation}
Its differential is the usual exterior contraction
by those deck images minus one.

\item The source signs are as follows. Write
a boundary product cell as $\alpha\wedge\beta$
with an optional final $\Theta$, where $\alpha$
uses the first four coordinates and $\beta$ the
second four. Put $\epsilon=1$ when $\Theta$ occurs
and $0$ otherwise. Its first vertex map applies
the \ref{maps:src:E.2.16} peripheral map to $(\alpha,\Theta)$,
keeps $\beta$, and has sign
$(-1)^{\epsilon(|\alpha|+|\beta|)}$.
The second map applies that peripheral map to
$(\beta,\Theta_2)$, keeps $\alpha$, and has sign
$(-1)^{|\alpha|(|\beta|+\epsilon)+\epsilon|\beta|}$.
If $\epsilon=1$, multiply its coefficient on the
right by $-\mu$. This is the exact cellular map
of $\Theta_1\mapsto\Theta_2^{-1}$, since
$\mu_2=cd=\mu^{-1}$. The suspended cells in \eqref{maps:eq:E.C19a}--\eqref{maps:eq:E.C19b}
have the interval first; these signs convert from
the ordered exterior cells in \eqref{maps:eq:E.P11}.
Denote the two maps by $j_1,j_2$.

\item The chain complex of the splice is the cone of
$(j_1,-j_2):B\to V_1\oplus V_2$.
Project each $C_E$ to $C_K$ and each $T(s)$ to its
first circle, and project $B$ to its three coordinates
$g_1,g_2,\Theta$. The result is the complex of the
aspherical graph of spaces
\begin{equation}
 (K(K_1,1)\times S^1)\ \cup_{T^3}\
                         (K(K_2,1)\times S^1).
 \hypertarget{tag-02BM}{}\label{maps:eq:E.P12}
\end{equation}
It is the model $N$ of \ref{maps:src:E.2.9}. Both peripheral maps
on groups are injective by \ref{maps:src:E.2.10}; their common
group is $\langle h_1,h_2,\mu\rangle=\mathbb Z^3$.
Consequently its universal cover is a Bass--Serre
tree of contractible vertex and edge spaces and
is contractible. Its chain complex has ranks
$(2,7,9,5,1)$, including its degree-four cell.

\item The indicated projection is a degreewise split
surjection of the full chain complexes. The
peripheral Fox coefficient in \eqref{maps:eq:E.P9} makes the
boundary square commute; higher degrees commute
as well by the product construction and the
zero higher groups of $C_K$. Let $D$ be its kernel.
The target has zero positive-degree homology, so
the long exact sequence gives
$H_2(D)\simeq H_2(\widetilde P)=\pi_2(P)$.
Using the small models for the relevant images,
the ranks of $D$ through degree three are
\begin{equation}
                         (0,6,40,123).
 \hypertarget{tag-02BN}{}\label{maps:eq:E.P13}
\end{equation}
Its six degree-one generators are the unused
regular-factor circles at the two vertices.
The six shifted boundary circles $u_i,w_i,d_i$
map to them by unit entries, because each circle
maps to zero in its own complement and identically
to the opposite regular torus. Cancel these six
unit pairs. This leaves a presentation of $H_2(D)$
with $34$ generators and $123$ relations.

\item Here is also an explicit block definition of that
matrix. In each $V_i$ keep the cells killed by
projection to $C_{K_i}\otimes C(S^1)$; their ranks
in degrees two and three are $17$ and $45$. In $B_2$
keep the $33$ cells containing at least one of
$u_1,w_1,d_1,u_2,w_2,d_2$. In those ordered bases
the matrix is
\begin{equation}
 \mathcal M=
 \begin{pmatrix}
 (d_{V_1})_3&0&(j_1)_2\\
 0&(d_{V_2})_3&-(j_2)_2
 \end{pmatrix}:
 R^{45}\oplus R^{45}\oplus R^{33}
             \longrightarrow R^{17}\oplus R^{17}.
 \hypertarget{tag-02BO}{}\label{maps:eq:E.P14}
\end{equation}
Every entry is specified by \eqref{maps:eq:E.P7}--\eqref{maps:eq:E.P11} and
\eqref{maps:eq:E.C19a}--\eqref{maps:eq:E.C19b}; \eqref{maps:eq:E.P14} itself is a complete finite
presentation of the desired module. The formula
includes the derived peripheral comparison.


\item \emph{Reduction to the compact presentation.}
The peripheral chain equations and $d_2d_3=0$ hold by multiplication
of their group-ring matrices. Cancelling the indicated unit pairs reduces
\eqref{maps:eq:E.P14} to $16$ generators and $52$ nonzero columns.
The calculation takes place in the amalgam
\begin{equation}
 G=(K_1\times\mathbb Z)\mathop{*}_{\mathbb Z^3}
                                      (\mathbb Z\times K_2),
 \hypertarget{tag-02BP}{}\label{maps:eq:E.P15}
\end{equation}
whose normal form is specified as follows: the torus-knot factor uses $h^r$ times an
alternating word in $a,a^2$ and $b,b^2,b^3$.
Modulo the peripheral subgroup, strip initial
$ab$ and $b^3a^2=(ab)^{-1}$, allowing their last
letter to be a partial power; choose $b^3$ in
the sole length-one tie with $a$. Inverting these
left coset representatives gives right coset
representatives. The usual alternating amalgam
normal form, followed by a final element of
$\mathbb Z^3$, is unique. These word operations track both central exponents.

\item In the resulting $16$-generator presentation two
generators occur only in their respective single
relations $q_2e_1=0$ and $q_1e_2=0$. They give the
first two direct summands of \eqref{maps:eq:E.P2}. For the other
14 generators, two generators $t_0,t_1$ have
trivial action of $a,b,c,d$. There are three
generators initially killed by $q_2$ whose pure
relations include $h_1N_b$; the identity
$h_1N_b(b-1)=1-h_1$ kills them by $q_1$ as well.
There are three initially killed by $q_1$ with
pure $N_c$ relations, and three more with pure
$h_2N_d$ relations. They are killed by $q_2$.
The remaining generators already have both
central relations. Thus the whole $14$-generator summand is a module over $S$. The pure norm columns
are exact differences of the recorded columns
and are checked before setting either central
element equal to $1$.

\item The $14$ generators separate into two seven-generator blocks. The
first has generators $U_0,U_1,U_2,X,Y,Z,t_1$,
with $t_1$ invariant, $U_iN_b=0$, and
\begin{equation}
 \begin{split}
 (U_0+X-Yd+Zd)N_a&=0,\\
 (U_1-Y(1+d)+Z)N_a&=0,\\
 (U_2+X+2t_1+Z(d-1))N_a&=0,\\
 X(1+d)+3t_1&=0.
 \end{split}
 \hypertarget{tag-02BQ}{}\label{maps:eq:E.P16}
\end{equation}
The second has generators $B_i,C_i,t_0$, with
$B_iN_c=C_iN_d=0$, $t_0$ invariant, and the single
additional relation $t_0=\Omega_R$. These formulas
follow by collecting the displayed $16$-generator
columns.

\item In the first block put
\label{layout:V:1}
\begin{equation}
\begin{gathered}
v_0=X-Yd+Zd,\quad v_1=-Y(1+d)+Z,
 \\
\quad v_2=X+2t_1+Z(d-1),\qquad
 r_i=U_i,\quad p_i=U_i+v_i.
\end{gathered}
\hypertarget{tag-02BR}{}\label{maps:eq:E.P17}
\end{equation}
The linear part of this change on $(X,Y,Z)$ is
\[
 \begin{pmatrix}1&0&1\\-d&-1-d&0\\d&1&d-1\end{pmatrix},
 \quad\det=1,
\]
over $\mathbb Z[d,d^{-1}]$. It is therefore an
integral group-ring basis change. Its inverse gives
\[
 X=v_0(1-d^2)+v_1(d^2-d)+(v_2-2t_1)d^2.
\]
Substitute in the last relation of \eqref{maps:eq:E.P16}, using
the invariance of $t_1$. The result is
$t_1=\Omega_L$. Eliminate $t_1$ with this unit
coefficient. Its four invariance relations become
the four relations in \eqref{maps:eq:E.P3}. Eliminating $t_0$
in the second block gives \eqref{maps:eq:E.P4}. This proves the
compact presentation in \eqref{maps:eq:E.P2}--\eqref{maps:eq:E.P4}, including
the integral central-intersection information.

\item For the two cyclic axis summands, the other
central element is a non-zero-divisor: their
group rings are free over its Laurent polynomial
ring. Thus neither contributes to the common
kernel of $q_1,q_2$. This proves \eqref{maps:eq:E.P5}. On
coinvariants, the four invariance relations in
each block vanish, while the norm relations
become three copies of multiplication by $3$ and
three copies of multiplication by $4$. Each block
therefore has coinvariants $(\mathbb Z/12)^3$.
Taking coinvariants therefore gives the groups in \cref{maps:src:E.2.11}.

\item \emph{Restriction to the circle fibre.}
The universal covers of $P$ and $Y\times\mathbb R$
are $\Pi$-equivariantly homotopy equivalent by
\ref{maps:src:E.2.13}. The map $\Pi\to G/\langle h_i\rangle$
is injective, since $\chi(h_i)=\pm12$, and has
index $12$. Hence each of the first two group-ring
summands in \eqref{maps:eq:E.P2} is free of rank $12$ over
$\mathbb Z\Pi$. The last two factor through the
index $12$ surface subgroup $\Sigma<\Gamma$ and
have trivial $z$ action. This proves the module
formula in \eqref{maps:eq:E.P6}.

\item For its exact coinvariants use the regular
12-sheet permutation representation with
exponents $(4,-3,-4,3)$ on $(a,b,c,d)$.
Evaluating the compact $12$-by-$20$ presentation gives
an integral $144$-by-$240$ matrix. Its Smith form has
56 unit entries and two entries $6$; its free
cokernel rank is $86$. Adding the $24$ free
coinvariants proves the last formula in \eqref{maps:eq:E.P6}.

\end{enumerate}
The corresponding exact calculations are implemented in \hyperlink{prog:V:quadratic}{the programs in Table~\ref*{tab:V:programs}}.
\end{proof}


\begin{definition}[Whitehead quadratic group]\hypertarget{tag-PV24}{}\label{def:V:whitehead-quadratic}
\StatementRecord{def:V:whitehead-quadratic}{proved}{PV24}{}
For an abelian group $A$, let $\Gamma(A)$ be Whitehead's universal
quadratic group, with quadratic generator $q(a)$ and symmetric
bilinear cross effect
$[a,b]=q(a+b)-q(a)-q(b)$.  In particular $[a,a]=2q(a)$.
\end{definition}

\begin{proposition}[The cusp quadratic attachment]\hypertarget{tag-02MI}{}\label{maps:src:E.2.18}
\StatementRecord{maps:src:E.2.18}{proved}{02MI}{\cref{maps:src:E.2.16}, \eqref{maps:eq:E.P6}, \SeriesRef{maps:src:D.1.9}, \cref{maps:src:E.2.12}, \eqref{maps:eq:E.C12}, \eqref{maps:eq:E.Q4}, \eqref{maps:eq:E.Q5}, \eqref{maps:eq:E.Q3}, \eqref{maps:eq:E.Q7}, \eqref{maps:eq:E.Q6}, \eqref{maps:eq:E.C3}, \eqref{maps:eq:E.Q8}, \cref{maps:src:E.2.17}, \eqref{maps:eq:E.P2}, \eqref{maps:eq:E.P3}, \eqref{maps:eq:E.P4}, \eqref{maps:eq:E.Q11}, \eqref{maps:eq:E.Q10}, \eqref{maps:eq:E.Q12}; \cite[Theorems 1 and 5]{mapsWhitehead50}}


Keep the right-module presentation of \ref{maps:src:E.2.16} and write
\[
 R=\mathbb ZK,\qquad M=\pi_2(E)=eR\oplus T,
 \qquad T=\overline R^{\,2}\oplus L.
\]
There is a lift of the normalized cusp component
$S=\mathbb {CP}^2\#3\overline{\mathbb {CP}}^{\,2}$ to the
universal cover $\widetilde E$ whose fundamental class has image
\begin{equation}
 \pm\bigl([e(h-1),t]-12q(t)\bigr)
       \quad\hbox{under}\quad
 H_4(\widetilde E;\mathbb Z)\longrightarrow\Gamma(M).
 \hypertarget{tag-02C0}{}\label{maps:eq:E.Q1}
\end{equation}
The sign is solely the orientation chosen on the normalized
component.  Its mixed component in
\begin{equation}
 \Gamma(M)=\Gamma(eR)\oplus(eR\otimes_{\mathbb Z}T)
                                      \oplus\Gamma(T)
 \hypertarget{tag-02C4}{}\label{maps:eq:E.Q2}
\end{equation}
is nonzero.  Consequently $E$ is not homotopy equivalent to
$E_0\#(S^2\times S^4)$ for any connected compact six-manifold
$E_0$.  Thus the free $R$ summand of its second homotopy group
does not come from a sphere-product connected summand.
Likewise the nine-dimensional fibre $Y$ is not
homotopy equivalent to $Y_0\#(S^2\times S^7)$ for any
connected closed nine-manifold $Y_0$.

\end{proposition}


\begin{proof}
\begin{enumerate}[label=\arabic*.,leftmargin=*]
\item We use the natural, deck-equivariant part of Whitehead's exact
sequence
\begin{equation}
 H_4(\widetilde E)\xrightarrow{\omega_E}\Gamma(\pi_2(E))
       \longrightarrow\pi_3(E)\longrightarrow
 H_3(\widetilde E)\longrightarrow0.
 \hypertarget{tag-02C5}{}\label{maps:eq:E.Q3}
\end{equation}
The exactness, identification of the quadratic term, and direct-sum
formula \eqref{maps:eq:E.Q2} are given in
\cite[Theorem~1, p.~55; Section~2, p.~57; Theorem~5, p.~58]{mapsWhitehead50}.
The article explicitly allows infinite complexes on p.~57, so these
results apply to the simply connected universal cover.
Naturality specifies the element in \eqref{maps:eq:E.Q1}.

\item For the normalization calculation take the negatives of
the complex-oriented Picard homology generators and denote
them by $(H,E_1,E_2,E_3)$.  Their intersection matrix is
still $\operatorname{diag}(1,-1,-1,-1)$.  This common reversal matches the phase-cell orientations of the cusp basis: the toric
line with class $-(uw+6\gamma\delta)$ has complex curve
degrees $(-6,0,1)$, whereas its values on the corresponding
phase classes are $(6,0,-1)$; the support function $rs$
similarly fixes the middle sign.  The exact wall-degree
calculation in \SeriesRef{maps:src:D.1.9} gives this comparison.  With these
orientations the six conductor classes are
\[
 E_1,\ H-E_1-E_2,\ E_2,\ H-E_2-E_3,
                                      E_3,\ H-E_1-E_3.
\]
The first opposite pair is self-glued in the surviving $u$
direction.  In the gamma cover its relation is therefore
$E_1=H-E_2-E_3$.  The other two pairs connect adjacent gamma
translates.  Their relations, after eliminating $H$, are
$(h-1)E_2=(h-1)E_3=0$.  Thus the image of normalization
homology is, with these curve generators,
\begin{equation}
 A E_1\oplus(A/(h-1))E_2\oplus(A/(h-1))E_3,
 \qquad A=\mathbb Z[h,h^{-1}],\qquad H=E_1+E_2+E_3.
 \hypertarget{tag-02C6}{}\label{maps:eq:E.Q4}
\end{equation}
This also follows directly from the transpose of the matrix
$B(h)$ of \ref{maps:src:E.2.12}, with the intersection signs of the Picard basis.  It is the whole $H_2$ of the gamma cover:
the degree-zero normalization row is the cellular complex
of an infinite cylinder, whose second homology is zero.

\item We fix the comparison with the cusp basis of \eqref{maps:eq:E.C12}.  In the ordinary
central fibre the first conductor class $\overline C_1$ is the
self-$u$ pair; the other two are the images of $E_2,E_3$,
in either order.  The chosen cusp basis is
\[
 E=-\overline C_1-\overline C_2,\qquad
 U=\overline C_2+\overline C_3,\qquad D=\overline C_2.
\]
The two torsion coordinates in \eqref{maps:eq:E.Q4} are determined by
this ordinary comparison, because each is annihilated by
$h-1$.  The normalization map is onto the whole group in
\eqref{maps:eq:E.Q4}, and $E_2,E_3$ generate its entire $A$-torsion
submodule.  Consequently the coefficient of $E_1$ in the
free quotient generates $A$ and hence is a unit $\pm h^n$.  The ordinary comparison
fixes its sign.  Translating the entire chosen normalization
component by $h^{-n}$ makes this coefficient $-1$ and leaves
both torsion coordinates fixed.  The common Picard sign
convention above changes every generator simultaneously.
The resulting
normalization map in the basis used for the full peripheral
comparison is
\begin{equation}
 E_1\longmapsto-E-D,\qquad
 E_2\longmapsto D,\qquad E_3\longmapsto U-D,
 \qquad H\longmapsto-E+U-D.
 \hypertarget{tag-02C7}{}\label{maps:eq:E.Q5}
\end{equation}
Exchanging $E_2,E_3$ has no effect below.  \item The intersection form of $S$ is
$\operatorname{diag}(1,-1,-1,-1)$.  Its usual four-cell
description, or the Hopf invariant of its two-dimensional
classes, therefore gives
\[
 \omega_S[S]=q(H)-q(E_1)-q(E_2)-q(E_3).
\]
Substitute $H=E_1+E_2+E_3$ and then \eqref{maps:eq:E.Q5}.  Bilinearity
of the cross effect gives
\begin{equation}
 \omega_S[S]\longmapsto
 [E_1,E_2]+[E_1,E_3]+[E_2,E_3]
                      =-[E,U]-2q(D).
 \hypertarget{tag-02C8}{}\label{maps:eq:E.Q6}
\end{equation}
The normalization map $S\to E$ lifts to $\widetilde E$
because $S$ is simply connected.  Naturality of \eqref{maps:eq:E.Q3}
shows that the last expression is the image of its four-dimensional fundamental class.  \item The global cancellation also tracks the three cusp two-cycles in addition to the peripheral chains.  It gives, in the generators $(x,t,y,z,e)$ of
\ref{maps:src:E.2.16},
\begin{equation}
 E\longmapsto(1-h)e+6t,\qquad
 U\longmapsto t,\qquad D\longmapsto0.
 \hypertarget{tag-02C9}{}\label{maps:eq:E.Q7}
\end{equation}
For clarity, before the last relation reductions its
$t$ coefficient for $E$ is
$-3h+4ha+4ha^2+h^2$, whose augmentation is six;
its free coefficient is $1-h$.  Its coefficient for $D$
is $2h^{-1}a^2-1-a$, which has augmentation zero and hence
annihilates the invariant generator $t$.  The other coordinates vanish. These identities follow by applying
the integral group-ring cancellations of \ref{maps:src:E.2.16} to
the displayed cusp-cycle columns.  Applying $\Gamma$ to \eqref{maps:eq:E.Q7} in \eqref{maps:eq:E.Q6} gives
\[
 -[(1-h)e+6t,t]
                =[e(h-1),t]-12q(t),
\]
as asserted.  In ordinary coinvariants $t=2v$ and
$E=12v$ for a primitive class $v$.  The quadratic image
is then $-48q(v)$, equal to the degree-two cup square evaluated
on the normalization.

\item We finish with the connected-sum obstruction, allowing
every possible free summand.  Multiplication by $h-1$
is injective on $R$: a finitely supported group-ring
element constant along an infinite $h$ orbit is zero.
It annihilates $T$, so
\begin{equation}
 \operatorname{Ann}_M(h-1)=T,
                       \qquad\operatorname{Hom}_R(T,R)=0.
 \hypertarget{tag-02CA}{}\label{maps:eq:E.Q8}
\end{equation}
Suppose that a copy $fR$ is a direct summand of $M$.
Write $f=eu+t_0$ and let a retraction send $e$ to $v$.
It vanishes on $T$, and hence $vu=1$.

\item The ring $\mathbb ZK$ is directly finite.  One elementary
verification uses the index-twelve subgroup of $K$:
its quotient by the central $\mathbb Z$ is the free group
of the punctured genus-three surface, so it is
$F_6\times\mathbb Z$.  This subgroup, and hence $K$, is
residually finite.  In every finite-group quotient the
regular representation is finite dimensional over
$\mathbb Q$, so $vu=1$ implies $uv=1$.  Residual
finiteness separates the finite support of $uv-1$ and
proves the same equality in $R$.  Thus $u$ is a unit,
the kernel of the retraction is exactly $T$, and every
free splitting is obtained by replacing $e$ by
$eu+t_0$.

\item Under such a change the mixed component of \eqref{maps:eq:E.Q1} is
$f u^{-1}(h-1)\otimes t$.  It is nonzero: $t$ has
infinite order by \eqref{maps:eq:E.C3}, and $u^{-1}(h-1)$ is nonzero.
The central generator may also be reversed by a group
automorphism, which multiplies $h-1$ by a unit and
does not affect the argument.

\item If $E\simeq E_0\#(S^2\times S^4)$, its universal cover
has, through dimension four, the homotopy type of
$\widetilde E_0$ wedged with one copy of $S^2\vee S^4$
for each deck element.  This follows by lifting the
connected-sum balls; deleting or regluing their boundary
five-spheres does not alter the four-skeleton.  The new
two-spheres generate a free $R$ summand.  The four-spheres
have trivial attaching maps, and every remaining
four-dimensional Whitehead image lies in
$\Gamma(\pi_2(E_0))$.  There can therefore be no mixed
component between this free summand and its complement.
This contradicts \eqref{maps:eq:E.Q1}--\eqref{maps:eq:E.Q8} for every choice of the
summand, and proves the stated obstruction.

\item For the assertion about $Y$, first take the maps from the two complements to the splice.  Carry the four cycles $(t_0,e_0,t_1,e_1)$ through the unit
cancellations of \ref{maps:src:E.2.17}. In the resulting
sixteen-generator presentation they are four unit vectors.  The seven-generator block changes therefore give
\begin{equation}
 t_0=\Omega_R,\qquad t_1=\Omega_L,
 \qquad e_0\in R/q_2R,\quad e_1\in R/q_1R
 \hypertarget{tag-02CB}{}\label{maps:eq:E.Q9}
\end{equation}
with exactly the axes in \eqref{maps:eq:E.P2}.  
Both $t_i$ are invariant under $G$: the four displayed
relations for each $\Omega$ in \eqref{maps:eq:E.P3}--\eqref{maps:eq:E.P4} annihilate
it by $a-1,b-1,c-1,d-1$.  On $\Gamma$-coinvariants
their respective images are
\begin{equation}
 -2(B_2+C_2),\qquad 2(p_2-r_2),
 \hypertarget{tag-02C1}{}\label{maps:eq:E.Q10}
\end{equation}
and each has order exactly six: the relevant two coordinates
have independent orders three and four in \eqref{maps:eq:E.P3}--\eqref{maps:eq:E.P4}.
In particular an integer annihilating either $t_i$ in the
full module must be divisible by six.

\item Put $A_\Pi=\mathbb Z\Pi$ and $z=h_1h_2$.  Choose the twelve
coset representatives giving the free restriction of each
axis in \eqref{maps:eq:E.P6}, and denote the resulting twenty-four generators
by $f_j$.  The normalized cusp surfaces and all their chosen
deck translates are maps to $\widetilde P$.  Transfer
them to $\widetilde Y$ using the $\Pi$-equivariant homotopy
equivalence $\widetilde P\simeq\widetilde Y\times\mathbb R$.
On the first axis $h_2$ is the identity, so multiplication
by $h_1-1$ equals multiplication by $z-1$; the analogous
statement holds on the second axis.  Naturality of \eqref{maps:eq:E.Q3}
therefore supplies, for every $j$, an element in the image
of the Whitehead map of $\widetilde Y$ of the form
\begin{equation}
 \pm\bigl([f_j(z-1),t_{i(j)}]-12q(t_{i(j)})\bigr),
 \hypertarget{tag-02C2}{}\label{maps:eq:E.Q11}
\end{equation}
where $i(j)=0$ on the first twelve axes and $1$ on the
remaining twelve.  Invariance of $t_i$ under $G$ removes
any extra coset action in this formula.

\item Write $\pi_2(Y)=F\oplus T_Y$, with $F=A_\Pi^{24}$ and
$(z-1)T_Y=0$ as in \eqref{maps:eq:E.P6}.  Suppose a free copy $fA_\Pi$
is a direct summand, with a retraction $r:\pi_2(Y)\to A_\Pi$
satisfying $r(f)=1$.  Multiplication by $z-1$ is injective
on $A_\Pi$, so $r(T_Y)=0$.  Write $r(f_j)=v_j$ and
$f=\sum_j f_j u_j+t$; then
\begin{equation}
                         \sum_j v_j u_j=1.
 \hypertarget{tag-02C3}{}\label{maps:eq:E.Q12}
\end{equation}
In the splitting $fA_\Pi\oplus\ker r$, the mixed component
of \eqref{maps:eq:E.Q11} is $f v_j(z-1)\otimes t_{i(j)}$.  If all these
mixed components were zero, every integer coefficient of
$v_j(z-1)$ would be divisible by six, by \eqref{maps:eq:E.Q10}.  The
group-ring basis makes this implication immediate even if
$t_i$ has larger or infinite order.  Multiplication by
$z-1$ is also injective on $(\mathbb Z/6)[\Pi]$: along an
infinite $z$ orbit a finitely supported coefficient function
with zero difference is zero.  Hence every $v_j$ belongs
to $6A_\Pi$, contradicting \eqref{maps:eq:E.Q12}.

\item Finally, a connected sum with $S^2\times S^7$ has, through
dimension four in its universal cover, the original space
wedged with a free deck orbit of two-spheres.  All images
of its four-dimensional Whitehead map consequently have
zero mixed component with that free summand.  The preceding
argument rules this out for every free summand of the module.  This proves the assertion for $Y$.
\end{enumerate}
The corresponding exact calculations are implemented in \hyperlink{prog:V:quadratic}{the programs in Table~\ref*{tab:V:programs}}.
\end{proof}




\Needspace{9\baselineskip}
\section{Smooth plumbings and framed handles}\label{pub:V:sec-handles}

\begin{proposition}[The relative smooth cusp model]\hypertarget{tag-02MJ}{}\label{maps:src:E.2.19}
\StatementRecord{maps:src:E.2.19}{proved}{02MJ}{\cref{maps:src:E.2.21}, \cref{maps:src:E.2.16}, \eqref{maps:eq:E.T3}, \eqref{maps:eq:E.T4}, \eqref{maps:eq:E.T8}}


For the marked cusp, let
\[
 Z(q)=\frac{\log q}{2\pi i}B_0+C(q),\qquad
 B_0=\begin{pmatrix}0&1\\-1&0\end{pmatrix},
\]
where $C(q)$ is holomorphic and bounded on a disc about zero.
For a sufficiently small closed disc the six-dimensional
cusp tube is orientation-preservingly diffeomorphic, relative to
its full marked torus-bundle boundary, to the unitless toric
quotient defined by $C=0$.  The latter has the following finite
normal-crossings plumbing description.

Start with the canonical complex line bundle of
$S=\mathbb {CP}^2\#3\overline{\mathbb {CP}}^{\,2}$ and its six
toric curves in cyclic order
\begin{equation}
 \begin{aligned}
 C_0&=E_1,& C_1&=H-E_1-E_2,& C_2&=E_2,\\
 C_3&=H-E_2-E_3,& C_4&=E_3,& C_5&=H-E_1-E_3.
 \end{aligned}
 \hypertarget{tag-02DV}{}\label{maps:eq:E.T1}
\end{equation}
Pair $C_i$ with $C_{i+3}$ for $i=0,1,2$.  Both the surface-normal
line and the canonical line restricted to every $C_i$ have degree
$-1$.  In the resulting $\mathcal O(-1)\oplus\mathcal O(-1)$
coordinates $(z,a,b)$, with $z=0$ at the endpoint meeting
$C_{i+1}$ and $b$ the surface-normal coordinate, the pairing is
\begin{equation}
             (z,a,b)\longmapsto(z^{-1},zb,za).
 \hypertarget{tag-02DW}{}\label{maps:eq:E.T2}
\end{equation}
The function giving the cusp parameter is $q=zab$ in this
coordinate chart.  The full coordinate map \eqref{maps:eq:E.T2} specifies the normal
framing of each plumbing.
The two triple points are the endpoint orbits
$\{0,2,4\}$ and $\{1,3,5\}$, where endpoint $i$ means
$C_i\cap C_{i+1}$.  Ordered local complex triples undergo a
cyclic permutation around each orbit, with threefold composition
the identity.  The plumbing is formed with collars and rounded
corners; its distinguished outer boundary is $|q|=\eta$.

\end{proposition}


\begin{proof}
\begin{enumerate}[label=\arabic*.,leftmargin=*]
\item We first construct the comparison by a smooth family.  Set
$C_s(q)=(1-s)C(q)$ for $0\leq s\leq1$.  The affine toric cover
is fixed.  On its dense torus the deck translation for
$\lambda\in\mathbb Z^2$ becomes
\begin{equation}
 (x,q,s)\longmapsto
 \bigl(q^{B_0\lambda}
       \exp(2\pi i C_s(q)\lambda)x,\ q,\ s\bigr).
 \hypertarget{tag-02DX}{}\label{maps:eq:E.T3}
\end{equation}
The exponential is componentwise.  The powers of $q$ permute
the toric charts by translation of the $A_2$ triangulation,
while the exponential factors are nowhere-zero analytic units
on every chart.  Thus \eqref{maps:eq:E.T3} defines a smooth family of
holomorphic deck actions on a fixed smooth toric covering
space.  On a central stratum the lattice translation permutes
its bounded fan cells without fixed points.  Off the central
fibre, freeness and properness are exactly the nondegeneracy
of the real period lattice.

\item These assertions hold uniformly after shrinking $\eta$.
Indeed put $T=-\log|q|/(2\pi)$, and let $K$ bound
$\|\operatorname{Im}C_s(q)\|$ on the compact parameter interval
and closed disc.  For $T>K+1$ the matrix
$TB_0+\operatorname{Im}C_s(q)$ is invertible, since $B_0$ is
orthogonal.  A fundamental parallelogram has logarithmic
position
$B_0a+T^{-1}\operatorname{Im}C_s(q)a$ with
$a\in[0,1]^2$; it meets a uniformly bounded number of toric
charts.  More explicitly choose a triangle containing the
unperturbed point $B_0a$ and let $P$ be its unimodular vertex
matrix.  Its barycentric coordinates $t_j$ are nonnegative,
and the corresponding chart coordinates satisfy
$\log|z_j|=t_j\log|q|+
O(2\pi\|P^{-1}\|K\sqrt2)$.
There are only finitely many such triangles, so with
$c=2\pi\sqrt2\max\|P^{-1}\|$ every chosen chart coordinate
has modulus at most $\exp(cK)$, uniformly in $q,s$ and
all compact phases.  Their closed polydiscs have compact
closure in the affine charts.  These finitely many compact
chart images prove
compactness of the quotient over $[0,1]$ and $|q|\leq\eta$.
The central charts and their incidence relations are fixed.
Projection of this quotient to $s$ is consequently a proper
smooth submersion, including on every normal-crossings stratum.
Its only boundary is $|q|=\eta$ together with the parameter
end faces, and its boundary restriction is also a submersion.

\item On $|q|=\eta$ prescribe the marked real coordinates
$x=\exp(2\pi i(Z_s(q)a+b))$, $(a,b)\in T^4$.
They give a smooth product trivialization in $s$ and precisely
the prescribed mapping-torus marking.  Extend its horizontal
vector field over a collar, and then over the compact family.
Local toric product charts provide extensions tangent to
every central stratum; a partition of unity preserves that
property and the lift of $\partial_s$.  Integrating this
field proves the asserted relative diffeomorphism.  Its
orientation is the complex orientation on each fibre.
This is relative Ehresmann transport in a smooth
family.

\item It remains to identify the endpoint $C=0$ with the framed
plumbing data.  In the lattice $\mathbb Z^2$ use cyclic rays
\begin{equation}
 (1,0),\ (0,1),\ (-1,1),\ (-1,0),\ (0,-1),\ (1,-1).
 \hypertarget{tag-02DY}{}\label{maps:eq:E.T4}
\end{equation}
The star of the vertex zero in the $A_2$ triangulation consists
of the six triangles spanned by zero and consecutive rays.
The corresponding three-dimensional cones have primitive rays
$(v,1)$ and determinant one.  Their common divisor at the
central ray is the toric surface $S$.  A principal central
fibre has normal line
\begin{equation}
             \mathcal O_S(-C_0-\cdots-C_5)=K_S.
 \hypertarget{tag-02DZ}{}\label{maps:eq:E.T5}
\end{equation}
The cyclic surface fan in \eqref{maps:eq:E.T4} also proves this identification
directly: it is the fan of the three torus-fixed blow-ups of
$\mathbb {CP}^2$, and each neighboring-ray relation gives
$C_i^2=-1$.  In the Picard basis \eqref{maps:eq:E.T1},
$K_S=-3H+E_1+E_2+E_3$, so $K_S\cdot C_i=-1$.

\item Near a double curve the two adjacent unimodular charts give
the total space of $\mathcal O(-1)\oplus\mathcal O(-1)$.
Its two base charts have coordinates
\begin{equation}
 (z,a,b),\qquad
 (w,A,B)=(z^{-1},za,zb),\qquad q=zab=wAB.
 \hypertarget{tag-02E0}{}\label{maps:eq:E.T6}
\end{equation}
The period translation taking a neighboring component back
to the chosen component exchanges the two normal directions.
It reverses the chosen endpoint coordinate on the opposite
toric curve.  In the first chart this is exactly \eqref{maps:eq:E.T2};
in the target infinity chart it is simply $(z,b,a)$, so it
extends smoothly over the endpoint.  The map is involutive,
preserves $q$, and its complex Jacobian determinant is $1$.
These equalities specify the normal-disk attachment, including
its sign and angular twisting.

\item The endpoint $C_i\cap C_{i+1}$ is sent to
$C_{i+2}\cap C_{i+3}$.  At the first endpoint write the
ordered ambient coordinates as the surface-normal coordinate,
the $C_i$-normal coordinate and the $C_{i+1}$-normal coordinate.
The next endpoint has their order cyclically permuted.
The threefold composition is the identity on the ordered coordinates.  Thus the two endpoint orbits give
ordinary triple charts $q=z_0z_1z_2$ with trivial clutching around each triple point.

\item Equivalently, for a triangle with vertex matrix $P$ whose
columns are $(v_i,1)$, translation by $n\in\mathbb Z^2$
replaces it by $L_nP$, where
\begin{equation}
 L_n=\begin{pmatrix}1&0&n_1\\0&1&n_2\\0&0&1\end{pmatrix}.
 \hypertarget{tag-02E1}{}\label{maps:eq:E.T7}
\end{equation}
It acts as the identity in the corresponding translated
unimodular coordinates, since
$(L_nP)^{-1}L_nP=I$.  This verifies the preceding plumbing
and triple identifications in the toric charts.
More explicitly, order the star matrix as
$P_i=((0,1),(r_i,1),(r_{i+1},1))$ and reverse its
columns when using the coordinate order $(z,a,b)$.
For every $i$ the matrix of translation by $-r_i$ from
this chart to the opposite following-endpoint chart is
\begin{equation}
 \begin{pmatrix}-1&0&0\\1&0&1\\1&1&0\end{pmatrix}.
 \hypertarget{tag-02E3}{}\label{maps:eq:E.T9}
\end{equation}
Its monomials are exactly \eqref{maps:eq:E.T2}.  In normal-first order,
$P_{i+2}^{-1}L_{-r_i}P_i$ is the cyclic permutation
matrix.  Thus the framing and triple-cycle assertions hold for all six fan identifications.
Taking collars of the component core, three double-curve
annuli and two triple charts, and rounding their outer
corners, gives the asserted finite plumbing.  These are
subcharts of the smooth toric quotient, so the construction
is Hausdorff and the collar identifications are the indicated
smooth ones.

\item Finally the period translations for $g,u$ are $-e_2,e_1$.
On $q=\eta e^{-2\pi i\theta}$ the marked expression is
$x=\exp(2\pi i((-\theta+iT)B_0a+b))$.
At $\theta=1$ it agrees with $\theta=0$ and
$b$ replaced by $b-B_0a$.  This is exactly the inverse
endpoint convention for
\begin{equation}
 M_0=\begin{pmatrix}
 1&0&0&0\\0&1&0&0\\0&1&1&0\\-1&0&0&1
 \end{pmatrix}.
 \hypertarget{tag-02E2}{}\label{maps:eq:E.T8}
\end{equation}
The comparison therefore preserves the full outer marking.

\end{enumerate}
The finite algebra and combinatorial data are implemented in \hyperlink{prog:V:framings}{the programs in Table~\ref*{tab:V:programs}}.
\end{proof}


\begin{proposition}[The explicit smooth plumbing presentation]\hypertarget{tag-02ML}{}\label{maps:src:E.2.20}
\StatementRecord{maps:src:E.2.20}{proved}{02ML}{\cref{maps:src:E.2.21}, \cref{maps:src:E.2.19}, \eqref{maps:eq:E.S1}, \eqref{maps:eq:E.T2}, \eqref{maps:eq:E.U7}, \eqref{maps:eq:E.U10}, \eqref{maps:eq:E.U8}, \eqref{maps:eq:E.U9}, \SeriesRef{maps:src:B.5.1}, \cref{maps:src:E.2.6}}


The following finite smooth plumbing presentation uses only standard
pieces and specified attaching maps. Its conversion into ordinary
framed handles and the relative smooth identification are given in \ref{maps:src:E.2.21}.

Let $F=\mathbb R^4/\mathbb Z^4$, with basis $(g,u,w,d)$, and use the
two matrices $A_j$ and fixed vectors $v_j$ of \SeriesRef{mapsI:src:A.1.1}.  Set
\[
 (m_1,m_2)=(3,4),\quad(\kappa_1,\kappa_2)=(1,-1),\quad
 w_j=\kappa_jv_j,\qquad L=\ker\gamma,
\]
so that $\Lambda=\mathbb Zw_j\oplus L$ and $A_jw_j=w_j$.
On $L$ the two matrices are
\begin{equation}
 B_1=\begin{pmatrix}0&1&0\\-1&-1&0\\1&0&1\end{pmatrix},
 \qquad
 B_2=\begin{pmatrix}0&-1&0\\1&0&0\\0&1&1\end{pmatrix}.
 \hypertarget{tag-02E4}{}\label{maps:eq:E.U1}
\end{equation}
Define the standard smooth four-manifolds
\begin{equation}
 \mathcal S_j=(T^3\times\mathbb R)/
       ((l,\sigma+1)\sim(B_j^{-1}l,\sigma)),
 \qquad \mathcal N_j=\mathcal S_j\times D^2_r.
 \hypertarget{tag-02E7}{}\label{maps:eq:E.U2}
\end{equation}


The cap boundary is identified with $M_{A_j}$ by the following
explicit map.  Write $x=\phi w_j+l$ and use the interval marking
$(x,1)\sim(A_j^{-1}x,0)$ on $M_{A_j}$.  Then
\begin{equation}
 [x,\theta]\longmapsto
 \left([l,\sigma=m_j\kappa_j\phi+\theta],
                         \ r=e^{2\pi i\kappa_j\phi}\right).
 \hypertarget{tag-02E8}{}\label{maps:eq:E.U3}
\end{equation}
All angles on the right are taken in their stated mapping-torus
quotients.  In the ordered angular variables $(\phi,\theta)$
and $(\sigma,\alpha)$, where $r=e^{2\pi i\alpha}$, the map is
\begin{equation}
 \begin{pmatrix}\sigma\\\alpha\end{pmatrix}
 =
 \begin{pmatrix}m_j\kappa_j&1\\\kappa_j&0\end{pmatrix}
 \begin{pmatrix}\phi\\\theta\end{pmatrix},
 \qquad
 \begin{pmatrix}\phi\\\theta\end{pmatrix}
 =
 \begin{pmatrix}0&\kappa_j\\1&-m_j\end{pmatrix}
 \begin{pmatrix}\sigma\\\alpha\end{pmatrix}.
 \hypertarget{tag-02E9}{}\label{maps:eq:E.U4}
\end{equation}
These are integral unimodular matrices.

To define a compact six-manifold $\mathcal E_{\rm pl}$, take the
flat $F$-bundle over an oriented disc with three open interior
discs removed.  Its three marked inner loop monodromies are
$A_1,A_2,M_0$, with $A_1A_2M_0=I$.  Equivalently cut the planar
surface along three disjoint paths to its outer boundary, take
the product with $F$, and reglue the two banks by these three
integral linear diffeomorphisms.  This specifies smooth band
attachments, including their product normal coordinates.
Fill the first two boundary components by \eqref{maps:eq:E.U2}--\eqref{maps:eq:E.U4}, and
the third by the finite canonical-bundle plumbing of \ref{maps:src:E.2.19},
using its complete marked $M_0$ boundary map.  The remaining
boundary is marked as $F\times S^1_\mu$, with the outer meridian
orientation of \eqref{maps:eq:E.S1}.  The core uses its zero section.

For the regular-fibre complement there is an
orientation-preserving diffeomorphism
\begin{equation}
                  E\ \cong\ \mathcal E_{\rm pl}
                 \quad\hbox{relative to }F\times S^1_\mu .
 \hypertarget{tag-02EA}{}\label{maps:eq:E.U5}
\end{equation}
Consequently the shifted fivefold has the explicit smooth
presentation
\begin{equation}
 P\ \cong\
 (\mathcal E_{\rm pl}\times T^4)
 \ \cup_{\ (x_1,x_2,\theta)\mapsto(x_1,x_2,-\theta)\ }\
 (T^4\times\mathcal E_{\rm pl}).
 \hypertarget{tag-02EB}{}\label{maps:eq:E.U6}
\end{equation}
Every piece in this presentation is a standard torus product,
a linear three-torus mapping torus times an oriented disk, or
the canonical disk bundle of the explicitly specified toric
blow-up with the normal-disk plumbings \eqref{maps:eq:E.T2}.  All attachment
maps have been specified on their normal coordinates.

\end{proposition}


\begin{proof}
\begin{enumerate}[label=\arabic*.,leftmargin=*]
\item We verify first that \eqref{maps:eq:E.U2}--\eqref{maps:eq:E.U4} are the finite caps,
including the normal framing.  The original cap is
\begin{equation}
 (F\times D^2_s)/
 \bigl((x,s)\sim(A_jx+v_j/m_j,e^{-2\pi i/m_j}s)\bigr).
 \hypertarget{tag-02EC}{}\label{maps:eq:E.U7}
\end{equation}
The good-reduction real flat marking is smooth on the entire
root disc, so these are smooth coordinates even at $s=0$.
Because $\gamma A_j=\gamma$ and $\gamma(v_j)=\kappa_j=\pm1$,
\begin{equation}
                  r=s\,e^{2\pi i\kappa_j\gamma(x)}
 \hypertarget{tag-02ED}{}\label{maps:eq:E.U8}
\end{equation}
is invariant under the displayed action.  It is single-valued
on $F$, has the same absolute value as $s$, and is complex
linear on every normal disk.  The induced map
$[x,s]\mapsto([x],r)$ identifies \eqref{maps:eq:E.U7} smoothly with
$S_j\times D^2$.  An inverse is obtained by choosing any
representative $x$ and setting
$s=r e^{-2\pi i\kappa_j\gamma(x)}$; invariance makes this
independent of the representative in the quotient.

\item In the integral basis $(w_j,u,w,d)$ the affine generator is
\begin{equation}
                  (\phi,l)\longmapsto
                     (\phi+\kappa_j/m_j,B_jl).
 \hypertarget{tag-02EE}{}\label{maps:eq:E.U9}
\end{equation}
Thus $\sigma=m_j\kappa_j\phi$ gives exactly the mapping torus
in \eqref{maps:eq:E.U2}.  This also supplies a direct inverse to the product
description: choose a representative $(l,\sigma)$, put
$x=(\kappa_j\sigma/m_j)w_j+l$ and
$s=r e^{-2\pi i\sigma/m_j}$.  Replacing that representative
by the defining mapping-torus relation changes $(x,s)$ by
the finite affine generator and hence not its class.

\item The boundary parametrization of the finite cap is
\begin{equation}
 [x,\theta]\longmapsto
       [x+\theta v_j/m_j,e^{-2\pi i\theta/m_j}].
 \hypertarget{tag-02E5}{}\label{maps:eq:E.U10}
\end{equation}
Substitute \eqref{maps:eq:E.U10} into \eqref{maps:eq:E.U8}--\eqref{maps:eq:E.U9}.  The coefficient of
$\theta$ in the normal angle is
$-1/m_j+\kappa_j^2/m_j=0$, while the base angle is
$m_j\kappa_j\phi+\theta$.  This proves \eqref{maps:eq:E.U3} and \eqref{maps:eq:E.U4}.
At $\theta=1$ that base angle advances by one, precisely
replacing $l$ by $B_j^{-1}l$; this checks the mapping-torus
seam.  A lattice change of $\phi$ changes the base angle by $m_j\kappa_j$; the equality $B_j^{m_j}=I$ makes the map descend under this change.

\item The function of the cap that gives its original base parameter
is, in these smooth coordinates,
\begin{equation}
                         q=r^{m_j}e^{-2\pi i\sigma}.
 \hypertarget{tag-02E6}{}\label{maps:eq:E.U11}
\end{equation}
On the marked boundary it has angle
$m_j\alpha-\sigma=-\theta$, the original clockwise base
convention.  A positive normal meridian at fixed
$(l,\sigma)$ has $(\Delta\phi,\Delta\theta)=(\kappa_j,-m_j)$.
Since $\kappa_jw_j=v_j$, it is exactly $t_{v_j}\rho_j^{-m_j}$,
recovering the filling relation
$\rho_j^{m_j}=t_{v_j}$.

\item Over the punctured planar core, the family has the
smooth real flat marking with linear monodromies $A_j$ and
its zero section.  The logarithmic finite gluing has exactly
the drift in \eqref{maps:eq:E.U10}; after subtracting it, the remaining
translation has a single-valued real lift and is removed by
the collar isotopy already proved in \SeriesRef{maps:src:B.5.1}.    At the cusp use the relative smooth diffeomorphism of \ref{maps:src:E.2.19}.  Each comparison
has precisely the prescribed boundary value and can be
made product on a smaller collar.  They therefore glue to
the relative diffeomorphism \eqref{maps:eq:E.U5}.  In particular the
normal framing of the unfilled ordinary fibre is unchanged.

\item Finally use the two hemispheres in the shifted
fibre-product decomposition of \ref{maps:src:E.2.6}.  On each hemisphere
the other factor is a regular marked torus bundle over a
disc, hence a product preserving its zero section and flat
marking.  Their common boundary has unchanged torus
coordinates and opposite normal angles.  Taking the product
of \eqref{maps:eq:E.U5} with those regular factors and gluing gives
\eqref{maps:eq:E.U6}, with its induced complex orientation.  
\end{enumerate}
The finite algebra and combinatorial data are implemented in \hyperlink{prog:V:framings}{the programs in Table~\ref*{tab:V:programs}}.
\end{proof}



\begin{theorem}[Ordinary framed handles and the smooth type]\hypertarget{tag-02MM}{}\label{maps:src:E.2.21}\hypertarget{tag-02MO}{}\label{maps:src:E.2.4}
\StatementRecord{maps:src:E.2.21}{proved}{02MM}{\eqref{maps:eq:E.T4}, \cref{maps:src:E.2.20}, \eqref{maps:eq:E.U3}, \eqref{maps:eq:E.U4}, \cref{maps:src:E.2.19}, \eqref{maps:eq:E.T6}, \eqref{maps:eq:E.H13}; \cite[Section~4.1, equation~(16), p.~103]{mapsVogel09}; \cite{mapsMilnor65}}


The following finite ordinary framed-handle prescription defines a
compact oriented six-manifold $\mathcal H_6$ with marked boundary
$T^4\times S^1$, and a closed oriented ten-manifold $\mathcal H_{10}$.
There are orientation-preserving diffeomorphisms
\begin{equation}
 E\cong\mathcal H_6\quad\hbox{relative to the full marked boundary},
 \qquad P\cong\mathcal H_{10}.
 \hypertarget{tag-02AQ}{}\label{maps:eq:E.H0}
\end{equation}
The respective numbers of ordinary handles, in increasing index, are
\[
 (2,5,14,22,21,10,2),\qquad
 (4,31,119,290,485,574,485,290,119,31,4).
\]
The attaching embeddings and their normal coordinates are given below.  
The balls and complex normal charts carry their complex orientations;
orient the product blocks by their displayed coordinate embeddings.
For example the dense block below has orientation
$-d\arg X\wedge d\arg Y\wedge d\mu_1\wedge d\mu_2
\wedge d\operatorname{Re}q\wedge d\operatorname{Im}q$.
This fixes the orientation convention even when the chosen order of
angle and moment coordinates differs from the complex coordinate order.

\emph{1. The two triple-point balls.}
Use two complex coordinate balls, denoted $V_+$ and $V_-$, with
ordered coordinates $Z=(Z_0,Z_1,Z_2)$ and $W=(W_0,W_1,W_2)$.
The boundary faces are rounded coordinate faces of
\begin{equation}
 |Z_i|\leq\epsilon_v,\quad |Z_0Z_1Z_2|\leq\eta,
 \qquad
 |W_i|\leq\epsilon_v,\quad |W_0W_1W_2|\leq\eta.
 \hypertarget{tag-02AR}{}\label{maps:eq:E.H1}
\end{equation}
They are smooth six-balls.  Throughout the prescription, rounding is
in the nonnegative radii, preserves the angular coordinates, and is
fixed on the attaching collars.  Positive changes of the radii only
give the corresponding collar isotopy.  The complex products in
\eqref{maps:eq:E.H1} are the same function $q$.
The coordinate formulas may first be read on manifolds with corners;
apply the same radial collar reparametrization on the two sides of
each attaching face before smoothing their union.  This convention
also specifies the maps at a corner meeting $|q|=\eta$.

\emph{2. The three annular blocks, with their endpoint framings.}
For $i=0,1,2$ use a block $S^1\times D^1\times D^4$ with coordinates
\begin{equation}
 z=\exp(s+2\pi i\theta),\quad -R\leq s\leq R,
 \qquad (a,b)\in D^4,\qquad q=zab.
 \hypertarget{tag-02AW}{}\label{maps:eq:E.H2}
\end{equation}
The normal four-ball here is a rounded radial domain
$|a|,|b|\leq\epsilon_e$, $|ab|\leq\eta/|z|$;
its positive radial trivialization identifies it with $D^4$.
The ends attach to framed tubes about the following coordinate
circles.  A table entry $(Z_j;Z_k,Z_l)$ means that the curve
coordinate is $Z_j$ and the two ordered complex normal coordinates
are $Z_k,Z_l$; the same convention applies to $W$.
\begin{equation}
\begin{array}{c|c|c}
 i&z=0\hbox{ end: }(z;a,b)&z=\infty\hbox{ end: }(w;A,B)\\ \hline
0&(Z_2;Z_1,Z_0)&(W_0;W_1,W_2)\\
1&(W_2;W_1,W_0)&(Z_1;Z_2,Z_0)\\
2&(Z_0;Z_2,Z_1)&(W_1;W_2,W_0)
\end{array}
\hypertarget{tag-02AX}{}\label{maps:eq:E.H3}
\end{equation}
The full-dimensional endpoint map between the two end charts is
\begin{equation}
                 (w,A,B)=(z^{-1},za,zb).
 \hypertarget{tag-02AY}{}\label{maps:eq:E.H4}
\end{equation}
In particular the second attaching circle has angle $-\theta$,
and its two complex normal vectors each rotate by $+2\pi\theta$.
The factors $|z|$ are the positive radial rescalings.
  Every coordinate circle
of the two balls occurs exactly once in \eqref{maps:eq:E.H3}.  The maps are
embeddings of $S^1\times\partial D^1\times D^4$.

\emph{3. The dense torus block and its full sector maps.}
Attach $T^2\times D^2\times D^2_q$ along
$T^2\times\partial D^2\times D^2_q$.  Here is a completely explicit
choice of its base disk and the attaching map.  Use the rays
$r_i$ of \eqref{maps:eq:E.T4} and the hexagon
\begin{equation}
 \mathcal P=\{v\in\mathbb R^2:\langle v,r_i\rangle\geq-1\},
 \quad
 \operatorname{Vert}(\mathcal P)=
 \{(-1,-1),(0,-1),(1,0),(1,1),(0,1),(-1,0)\}.
 \hypertarget{tag-02AZ}{}\label{maps:eq:E.H5}
\end{equation}
On the positive real part of $(\mathbb C^*)^2$, use the coordinates
given by the toric moment map
\begin{equation}
 \mu(X,Y)=
 \frac{\displaystyle\sum_{v\in\operatorname{Vert}(\mathcal P)\cup\{0\}}
                 |X^{v_1}Y^{v_2}|^2v}
      {\displaystyle\sum_{v\in\operatorname{Vert}(\mathcal P)\cup\{0\}}
                 |X^{v_1}Y^{v_2}|^2}.
 \hypertarget{tag-02B0}{}\label{maps:eq:E.H6}
\end{equation}
Choose a small inward parallel hexagon, with its corners rounded
symmetrically, $\mathcal P_\epsilon\subset\operatorname{int}\mathcal P$.
The inverse image is $T^2\times\mathcal P_\epsilon$; its two circle
coordinates are $\arg X,\arg Y$.  On the boundary disk circle put
$(X,Y)$ equal to the positive inverse of $\mu$, multiplied by these
two phases.  On the six sectors of that boundary, the attaching
map into the preceding blocks is the following monomial map, with
the whole complex normal coordinate $q$ included:
\begin{equation}
\begin{array}{c|ccc}
 i& n_i&a_i&a_{i+1}\\ \hline
0&q/(XY)&X&Y\\
1&q/Y&XY&X^{-1}\\
2&qX&Y&(XY)^{-1}\\
3&qXY&X^{-1}&Y^{-1}\\
4&qY&(XY)^{-1}&X\\
5&q/X&Y^{-1}&XY
\end{array}
\qquad q=n_i a_i a_{i+1}.
\hypertarget{tag-02B1}{}\label{maps:eq:E.H7}
\end{equation}
Use the cyclic identifications
$(n_2,a_2,a_3)=(Z_1,Z_2,Z_0)$,
$(n_4,a_4,a_5)=(Z_2,Z_0,Z_1)$ and the analogous two cyclic
permutations of $W$ for vertices $3,5$.  Vertices $0,1$ are $Z,W$.
The chart at an edge is $(z,a,b)=(a_{i+1},a_i,n_i)$, with \eqref{maps:eq:E.H4}
at its other end.  These rules specify \eqref{maps:eq:E.H7} also where a sector
meets an annular block.  The transition on a common sector is the
monomial map $P_{i+1}^{-1}P_i$, where
$P_i=((0,1),(r_i,1),(r_{i+1},1))$; the sum of its exponent rows
is $(1,1,1)$, so it preserves $n_i a_i a_{i+1}$.  Thus the sector maps specify the
normal framing as well as the core map.  

\emph{4. Expansion into ordinary handles.}
We use the following fixed elementary rule on every parameterized
block in this prescription:
\begin{equation}
 T^r\times D^k\times D^{n-r-k}
 \quad\hbox{attached on}\quad
 T^r\times\partial D^k\times D^{n-r-k}.
 \hypertarget{tag-02B2}{}\label{maps:eq:E.H8}
\end{equation}
Cut each ordered parameter circle into its two closed semicircles,
with small endpoint collars.  The first semicircle is its zero
handle and the second is its one handle.  For a bit vector
$\varepsilon\in\{0,1\}^r$, the resulting ordinary handle is of
index $k+|\varepsilon|$.  Attach by increasing $|\varepsilon|$,
then lexicographically.  On the original attaching face use
the displayed full-dimensional attaching embedding of the block;
on a newly exposed parameter face use the identity endpoint map
in that circle.  The normal coordinates are the remaining ordered
semicircle coordinates followed by the original normal coordinates.
Corner rounding and the permutation putting negative coordinates
first complete an ordinary framed attaching embedding
$S^{k+|\varepsilon|-1}\times D^{n-k-|\varepsilon|}$.
  For one parameter
it is the literal two-piece decomposition in
\cite[Section~4.1, equation~(16), p.~103]{mapsVogel09}.


Apply \eqref{maps:eq:E.H8} first to \eqref{maps:eq:E.H2}, with $(r,k,n)=(1,1,6)$, and then
to the core, with $(r,k,n)=(2,2,6)$.  Together with the two balls
this gives cusp handle counts $(2,3,4,2,1)$.

\emph{5. The torus bands and the finite caps.}
Take the marked flat torus bundle of \ref{maps:src:E.2.20} over the disc with
three holes, regarding the cusp boundary as the incoming boundary.
Cut along the two marked disjoint arcs from the two finite boundary
circles to the outer circle.  The cut surface is an annulus;
its torus bundle is the collar of the already constructed cusp
boundary.  Restoring the two strips attaches two blocks
\begin{equation}
                 T^4\times D^1\times D^1.
 \hypertarget{tag-02B3}{}\label{maps:eq:E.H9}
\end{equation}
To fix their maps, use the oriented ribbon
polygon whose successive based inner boundary words are
$A_1,A_2,M_0$ and whose outer word is their product.  The negative
bank of the $j$th cut has coordinate $(x,t)$, the positive bank
has $(A_jx,t)$, and the strip uses the product normal coordinate.
Reversing a bank orientation replaces that matrix by its inverse.
This is the cut-polygon definition of the flat bundle in \ref{maps:src:E.2.20};
the identity $A_1A_2M_0=I$ specifies the unfilled outer boundary.
These gluing maps are linear.  Thus \eqref{maps:eq:E.H9} has its full attaching
embedding, including the transverse interval, fixed by these
linear maps.  Expand each by \eqref{maps:eq:E.H8}, with $(r,k,n)=(4,1,6)$.
The unused boundary circles are precisely the two finite boundaries
and the marked ordinary boundary.
For clarity, the initial bank coordinates on the already attached
cusp boundary are the explicit coordinates of \cref{maps:src:E.2.19}:
\begin{equation}
 \begin{gathered}
 q=\eta e^{-2\pi i\theta},\qquad T=-\log\eta/(2\pi),\\
 (X,Y)=\exp\bigl(2\pi i((-\theta+iT)B_0a+b)\bigr),
 \qquad a=(g,u),\ b=(w,d).
 \end{gathered}
 \hypertarget{tag-02B4}{}\label{maps:eq:E.H9a}
\end{equation}
The exponential is componentwise; at an interval seam use
$b\mapsto b-B_0a$.  Equations \eqref{maps:eq:E.H7} and \eqref{maps:eq:E.H9a} therefore
give the bank embeddings into the existing handle
boundary, including their derivatives.  

For the $j$th finite cap use its already explicit
$\mathcal S_j\times D^2_r$ and the entire boundary embedding
\eqref{maps:eq:E.U3}.  The mapping torus $\mathcal S_j$ has the following product
block construction: a block $T^3\times D^1$ over the first
semicircle of $\sigma$, followed by a block $T^3\times D^1$ over
the second semicircle, with endpoint maps $I$ and $B_j^{-1}$.
The first is a three-parameter zero handle, the second a
three-parameter one handle.  As a filling relative to
$\mathcal S_j\times S^1_r$, the normal disk supplies two negative
directions.  Consequently the relative cap consists of
\begin{equation}
 \begin{array}{c|c|c}
 \hbox{base-circle block}&(r,k,n)&\hbox{full attaching data}\\ \hline
 0&(3,2,6)&\hbox{\eqref{maps:eq:E.U3}, the first semicircle, and the $r$-disk}\\
 1&(3,3,6)&\hbox{\eqref{maps:eq:E.U3}, endpoint maps $I,B_j^{-1}$, and the $r$-disk}.
 \end{array}
 \hypertarget{tag-02AS}{}\label{maps:eq:E.H10}
\end{equation}
More explicitly an attaching face in the normal direction is sent
to \eqref{maps:eq:E.U3} at $r=e^{2\pi i\alpha}$, and an attaching face in a
base direction is the identical face of the preceding base block,
crossed with $D^2_r$.  At the base seam use $B_j^{-1}$ on $l$ and
the identity on $r$.  The inverse angular formula \eqref{maps:eq:E.U4} supplies
the full map back to the boundary being filled.  In particular
the normal meridian is $v_j-m_j\rho_j$.
Apply \eqref{maps:eq:E.H8} to the two blocks in \eqref{maps:eq:E.H10}.  Each cap adds the
ordinary relative counts $(1,4,6,4,1)$ in indices $2,\ldots,6$.
This completes the definition of $\mathcal H_6$.

\emph{6. The closed ten-dimensional handlebody.}
For each handle of $\mathcal H_6$, take its product with the four
ordered torus circles and apply \eqref{maps:eq:E.H8}.  This gives
$\mathcal H_6\times T^4$ with its full boundary framing.  On the
other side take the dual relative handle decomposition of
$T^4\times\mathcal H_6$: reverse the handle order and exchange
the negative and positive disks of every handle.  A handle of
index $j$ becomes one of index $10-j$, and its attaching map is
the belt embedding of the original handle.  Glue its
initial boundary by the explicit map
\begin{equation}
                 (x_1,x_2,\theta)\longmapsto
                                    (x_1,x_2,-\theta).
 \hypertarget{tag-02AT}{}\label{maps:eq:E.H11}
\end{equation}
All subsequent dual attaching maps are obtained by this fixed
exchange of the two coordinate disks, in reverse order, on the
displayed attaching embeddings. Use the factor order $E,T^4$
on the first half and $T^4,E$ on the second; their interchange
has sign $(-1)^{24}=1$.  The boundary-angle reversal in \eqref{maps:eq:E.H11}
is the orientation-reversing boundary gluing.  This defines
$\mathcal H_{10}$ entirely by ordinary handles.  If
\begin{equation}
 c(t)=2+5t+14t^2+22t^3+21t^4+10t^5+2t^6,
 \qquad
       c_{10}(t)=(c(t)+t^6c(t^{-1}))(1+t)^4,
 \hypertarget{tag-02AU}{}\label{maps:eq:E.H12}
\end{equation}
the coefficients are the counts in the statement.  Expanding these polynomials gives $76$ handles in dimension six
and $2432$ in dimension ten.

\end{theorem}


\begin{proof}
\begin{enumerate}[label=\arabic*.,leftmargin=*]
\item The main geometric point is that the first three steps give
\emph{smooth blocks of the cusp tube}.  We give the collar
argument explicitly.  Work in the smooth toric quotient of \ref{maps:src:E.2.19}.
Choose disjoint small coordinate neighborhoods of the two triple
points.  The three incident coordinate axes in each neighborhood
are exactly the six ends in \eqref{maps:eq:E.H3}.  Remove smaller endpoint disks
from each normalized double curve.  Its remaining annulus has
the two complex normal lines $\mathcal O(-1)|_{\rm annulus}$,
which are trivialized by $(a,b)$ in \eqref{maps:eq:E.T6}.  The resulting normal
domain in the tube is exactly the domain in \eqref{maps:eq:E.H2}.
Every ray in $\mathbb C^2$ meets its rounded boundary once and
transversely.  Indeed the three defining radius functions have
positive radial derivative wherever the corresponding boundary
face is active.  Positive radial rescaling therefore gives a
smooth family of four-balls as $z$ ranges over this compact
annulus.  Choose the radial change to be constant linear scaling
on a common sufficiently small ball about the origin and interpolate
monotonically only away from that ball; it is then smooth at the
origin as well as at the boundary and preserves the stated normal
frames.  The same argument for the coordinate and product
radius functions in \eqref{maps:eq:E.H1} gives six-balls at the triple points.
Rounding can be chosen constant on the end attaching collars.

\item   At
vertex $i$ the normal-first coordinates are $(n_i,a_i,a_{i+1})$.
At vertex $i+2$ the deck identification is their cyclic
permutation, and the threefold composition is the identity by
\ref{maps:src:E.2.19}.  At the two ends of $C_i$ the coordinates are
$(z,a,b)=(a_{i+1},a_i,n_i)$ and
$(w,A,B)=(a_{i-1},a_i,n_{i-1})$.  The ordinary
$\mathcal O(-1)^2$ transition is \eqref{maps:eq:E.H4}.  Substituting the two
cyclic permutations gives exactly \eqref{maps:eq:E.H3}.  This checks all six ends.

\item We next choose the stratum collars so that their remaining
interface is the entire attaching face of the dense block.
On the normalized toric surface take a small collar of its
hexagonal toric boundary, choosing the corner neighborhoods
first and the edge strips second.  Choose the strips narrower
than the corner collars and invariant under the opposite-edge
inversion and the two cyclic triple identifications.  The
uncovered moment region is the rounded interior disk
$\mathcal P_\epsilon$.  This can be done directly with six
inward parallel affine inequalities in \eqref{maps:eq:E.H5}, smoothing each
corner inside its already chosen corner neighborhood.  Thus the
common frontier in the smooth part of the normalization is
exactly $T^2\times\partial\mathcal P_\epsilon$.

\item The map \eqref{maps:eq:E.H6} gives the claimed coordinates: in logarithmic
positive coordinates it is the gradient of a log-sum-exponential
function.  Its Hessian is twice the covariance matrix of the
seven affinely spanning exponent vectors, hence positive
definite.  Strict convexity and the growth in each supporting
direction give a diffeomorphism onto
$\operatorname{int}\mathcal P$.  The inverse on the compact
disk and its frontier is smooth.  The canonical line on this
free torus locus is trivial; more specifically the toric
function $q$ is a nonzero toric factor times its normal
coordinate.  Hence \emph{$q$ itself}, including its phase,
trivializes that normal line.  The portion of the truncated
tube over this disk is therefore exactly
\begin{equation}
 (T^2\times\mathcal P_\epsilon)\times D^2_q,
 \quad\hbox{with interface}\quad
 (T^2\times\partial\mathcal P_\epsilon)\times D^2_q.
 \hypertarget{tag-02AV}{}\label{maps:eq:E.H13}
\end{equation}
In extending the preceding moment cuts off the zero fibre,
hold their equations constant in this $q$ coordinate on an
open neighborhood of the compact frontier.  On their ends
use the product cuts of the annular and triple charts.  A
collar interpolation between these choices is an interpolation
of positive radial cut functions, fixed on each already chosen
end collar.  Their radial derivatives remain positive, so its
flow is a collar isotopy.  The cuts can consequently be chosen
to agree with \eqref{maps:eq:E.H13} exactly.  The $q$ coordinate therefore gives a common normal trivialization across the six sector seams.

\item These are blocks in a sufficiently small truncated
neighborhood of the compact central fibre.  Their central
parts cover that fibre: two triple points, the three open
double annuli and the remaining normalized surface core.
Compactness from \ref{maps:src:E.2.19} permits one common sufficiently small
$\eta$: the union of their open chart neighborhoods contains
the compact zero fibre, so properness implies that it contains
the entire inverse image of a sufficiently small closed
$q$-disc.  The nested cuts partition this entire inverse image.
If a larger originally chosen tube is used, the
difference is an exterior collar: on the compact annular
region without critical points choose a vector field lifting
the radial direction of $q$, and integrate it.  On the
prescribed outer collar choose the lift in the marked real
torus coordinates of \ref{maps:src:E.2.19}.  It is then the identity on all
those marked coordinates during radial collar transport.


\item The embedding in \eqref{maps:eq:E.H13} is \eqref{maps:eq:E.H7}.  To check this,
the exponent matrix for a unimodular chart is $P_i^{-1}$,
applied to $(X,Y,q)$.  This gives each of the six displayed
rows.  Consecutive changes are $P_{i+1}^{-1}P_i$ and telescope
to the identity.  The translations identifying the two
triple orbits are exactly those already used in \eqref{maps:eq:E.H3}.
All six normal coordinate multipliers in \eqref{maps:eq:E.H7} are thereby
fixed, and the sector seams are equal maps on overlaps.  Positive rescaling of the chosen
radial domains leaves these monomials unchanged before
the collar reparametrization.  This completes the smooth,
fully framed identification of the cusp with steps 1--3.

\item The handle expansion rule follows from the product decomposition
of the parameter circles.  For one circle, write $S^1=I_0\cup I_1$.
First attach $I_0\times D^k\times D^{n-k-1}$; after rounding it
is a $k$-handle.  The remaining block is
$(I_1\times D^k)\times D^{n-k-1}$, a $(k+1)$-handle.
Its attaching face consists of the original negative face
and its two endpoint faces in the first handle.  The stated
maps on these three faces agree on their corners because
they are restrictions of the same product coordinates.
Iterate this argument on the ordered circle factors.
It proves \eqref{maps:eq:E.H8}, including all newly created attaching
faces and their normal framings.

\item The planar core relative to its incoming circle has two
ordinary surface one-handles.  Its flat bundle is a product
over each strip, so restoring the two cut strips is exactly
\eqref{maps:eq:E.H9}.  The zero section and the full lattice marking fix
the normal coordinates on these attachments.  Each cap in
\eqref{maps:eq:E.H10} is the smooth product of \ref{maps:src:E.2.20}.  Its relative
handle claim can also be checked by the function
$-|r|^2+\varepsilon f(l,\sigma)$, with a cutoff near the
incoming boundary: choose $f$ by a two-critical-point
function of $\sigma$, perturbed on its two critical tori
by the product Morse function on $T^3$.  On the remaining
arcs its base derivative dominates the perturbation, so
there are no further critical points.  The indices are
two plus the base indices, as in \eqref{maps:eq:E.H10}.    Summing the blocks
gives
\[
 (2,3,4,2,1,0,0)+2(0,1,4,6,4,1,0)
                     +2(0,0,1,4,6,4,1)
       =(2,5,14,22,21,10,2).
\]
Every comparison has its prescribed boundary value, so the
relative diffeomorphism of \ref{maps:src:E.2.20} gives the first assertion
of \eqref{maps:eq:E.H0}.  Taking the products, the dual relative handles
and the boundary flip of \ref{maps:src:E.2.20} gives the second.
The handle polynomials evaluated at $-1$ give Euler numbers $2$ and $0$.

\item The following relative uniqueness theorems determine the smooth
type from the full attaching embeddings.  The relevant relative uniqueness
results are Milnor's Theorem~3.13, which identifies an elementary
cobordism from its full characteristic embedding, Theorem~1.4,
which makes smooth composition unique with the marked boundary
fixed, and the isotopy-extension Theorem~5.8
\cite[Theorem~3.13, pp.~31--32; Theorem~1.4, p.~2 and pp.~25--26; Theorem~5.8, pp.~63--64]{mapsMilnor65}.  In the present notation the characteristic
embedding is the whole map
$S^{j-1}\times D^{n-j}\longrightarrow\partial W$.  Its normal coordinates
are part of \eqref{maps:eq:E.H3}--\eqref{maps:eq:E.H11}.  The elementary consequence needed
here can also be seen directly.  If $\phi_0,\phi_1$ are joined
by the isotopy of full attaching embeddings, extend
the velocity field on their images to a slightly larger boundary
collar and multiply it by a cutoff which is zero on the protected
marked boundary.  Its flow extends the attaching isotopy to an
ambient boundary isotopy.  Extending this flow across a product
collar of $W$, and using the identity on the new coordinate handle,
identifies $W\cup_{\phi_0}(D^j\times D^{n-j})$ with
$W\cup_{\phi_1}(D^j\times D^{n-j})$.  The same cutoff argument
applies to the small radial corner changes above.  This derives
the relative, framed version from the stated isotopy-extension
theorem.

\item Apply this comparison successively in the displayed block and
sector order.  The initial ball comparisons are their coordinate
identifications.  At each next stage the attaching map
and the prescribed attaching map agree in the computed normal
coordinates, or differ by the explicitly described positive
collar isotopy.  The comparison of that stage is consequently
the identity on the new coordinate handle and the product-collar
extension on the previous stages.  Theorems~3.13 and~1.4 give an
inductive relative diffeomorphism of the whole six-manifold.
The final ordinary boundary coordinates are unchanged.  The
same argument, with the belt embeddings and the boundary flip, applies to every handle on the dual ten-dimensional
side.  Each top handle is specified by its entire attaching map.  The computed
ordered full framed attaching data, under this relative uniqueness
theorem, are the classifying data used for \eqref{maps:eq:E.H0}.
\end{enumerate}
The quadratic attachment of the regular-fibre complement is described in
\cref{maps:src:E.2.18}.
The finite algebra and combinatorial data are implemented in \hyperlink{prog:V:framings}{the programs in Table~\ref*{tab:V:programs}}.
\end{proof}

\clearpage
\section*{Programs}
\addcontentsline{toc}{section}{Programs}
The programs below are available in the \href{https://github.com/KokunoYumeto/s6-complex-structure-workbench/tree/\SsixRepositoryCommit}{repository} at the immutable commit
\begin{center}\small\texttt{\SsixRepositoryCommit}\end{center}
Each filename links to its source at this commit.
\begingroup\small
\setlength{\tabcolsep}{4pt}
\begin{longtable}{@{}p{.39\textwidth}p{.57\textwidth}@{}}
\caption{Programs for the exact computations and interval enclosures.}\label{tab:V:programs}\\
\toprule Computation and statements & Programs \\\midrule
\endfirsthead
\toprule Computation and statements & Programs \\\midrule
\endhead
\bottomrule\endfoot
\hypertarget{prog:V:height}{}Endpoint and global height enclosures\par \cref{night:P1.18} &
\href{https://github.com/KokunoYumeto/s6-complex-structure-workbench/blob/\SsixRepositoryCommit/pro_returns/S6_SERIES_GRIND_checkpoint_10_20261120T0211Z/code/paper_v/fresh_frobenius.py}{\nolinkurl{fresh_frobenius.py}}\par \href{https://github.com/KokunoYumeto/s6-complex-structure-workbench/blob/\SsixRepositoryCommit/pro_returns/S6_SERIES_GRIND_checkpoint_10_20261120T0211Z/code/paper_v/fresh_local_domains.py}{\nolinkurl{fresh_local_domains.py}}\par \href{https://github.com/KokunoYumeto/s6-complex-structure-workbench/blob/\SsixRepositoryCommit/pro_returns/S6_SERIES_GRIND_checkpoint_10_20261120T0211Z/code/paper_v/fresh_global_cover.py}{\nolinkurl{fresh_global_cover.py}}\par \href{https://github.com/KokunoYumeto/s6-complex-structure-workbench/blob/\SsixRepositoryCommit/pro_returns/S6_SERIES_GRIND_checkpoint_10_20261120T0211Z/code/paper_v/verify_saved_certificate.py}{\nolinkurl{verify_saved_certificate.py}} \\[3pt]
\hypertarget{prog:V:resolved}{}Cohomology of the resolved fourfold\par \cref{maps:src:E.1.4,maps:src:E.1.5} &
\href{https://github.com/KokunoYumeto/s6-complex-structure-workbench/blob/\SsixRepositoryCommit/pro_returns/S6_10_PAPERS_checkpoint_15_20261119T1045Z/S6_10_PAPERS/code/V/replay_chains/code/fourfolds_checks.py}{\nolinkurl{fourfolds_checks.py}}\par \href{https://github.com/KokunoYumeto/s6-complex-structure-workbench/blob/\SsixRepositoryCommit/pro_returns/S6_10_PAPERS_checkpoint_15_20261119T1045Z/S6_10_PAPERS/code/V/replay_chains/code/fourfolds_integral_audit.py}{\nolinkurl{fourfolds_integral_audit.py}} \\[3pt]
\hypertarget{prog:V:second-pair}{}Integral calculation for the second rational pair\par \cref{maps:src:F.2.18} &
\href{https://github.com/KokunoYumeto/s6-complex-structure-workbench/blob/\SsixRepositoryCommit/pro_returns/S6_10_PAPERS_checkpoint_15_20261119T1045Z/S6_10_PAPERS/code/V/replay_chains/code/fourfolds_v1_homology.py}{\nolinkurl{fourfolds_v1_homology.py}} \\[3pt]
\hypertarget{prog:V:normal-forms}{}Relative rational normal forms\par \cref{maps:src:F.2.20,maps:src:F.2.26} &
\href{https://github.com/KokunoYumeto/s6-complex-structure-workbench/blob/\SsixRepositoryCommit/pro_returns/S6_10_PAPERS_checkpoint_15_20261119T1045Z/S6_10_PAPERS/code/V/replay_chains/code/fourfolds_relative_normalform.py}{\nolinkurl{fourfolds_relative_normalform.py}}\par \href{https://github.com/KokunoYumeto/s6-complex-structure-workbench/blob/\SsixRepositoryCommit/pro_returns/S6_10_PAPERS_checkpoint_15_20261119T1045Z/S6_10_PAPERS/code/V/replay_chains/code/fourfolds_relative_six_probe.py}{\nolinkurl{fourfolds_relative_six_probe.py}} \\[3pt]
\hypertarget{prog:V:lattices}{}Stable lattices and primitive extensions\par \cref{maps:src:F.2.22,maps:src:F.2.23} &
\href{https://github.com/KokunoYumeto/s6-complex-structure-workbench/blob/\SsixRepositoryCommit/pro_returns/S6_10_PAPERS_checkpoint_15_20261119T1045Z/S6_10_PAPERS/code/V/replay_chains/code/fourfolds_lattice_algorithm.py}{\nolinkurl{fourfolds_lattice_algorithm.py}}\par \href{https://github.com/KokunoYumeto/s6-complex-structure-workbench/blob/\SsixRepositoryCommit/pro_returns/S6_10_PAPERS_checkpoint_15_20261119T1045Z/S6_10_PAPERS/code/V/replay_chains/code/fourfolds_two_classes_lattices.py}{\nolinkurl{fourfolds_two_classes_lattices.py}} \\[3pt]
\hypertarget{prog:V:twists}{}Affine twists of the two rational classes\par \cref{maps:src:F.2.24} &
\href{https://github.com/KokunoYumeto/s6-complex-structure-workbench/blob/\SsixRepositoryCommit/pro_returns/S6_10_PAPERS_checkpoint_15_20261119T1045Z/S6_10_PAPERS/code/V/replay_chains/code/fourfolds_family_twists.py}{\nolinkurl{fourfolds_family_twists.py}} \\[3pt]
\hypertarget{prog:V:suspension}{}Finite spectral suspension calculation\par \cref{maps:src:F.2.25} &
\href{https://github.com/KokunoYumeto/s6-complex-structure-workbench/blob/\SsixRepositoryCommit/pro_returns/S6_10_PAPERS_checkpoint_15_20261119T1045Z/S6_10_PAPERS/code/V/replay_chains/code/fourfolds_suspension_filter.py}{\nolinkurl{fourfolds_suspension_filter.py}} \\[3pt]
\hypertarget{prog:V:middle-pairing}{}Order-six middle-pairing obstruction\par \cref{maps:src:F.2.27} &
\href{https://github.com/KokunoYumeto/s6-complex-structure-workbench/blob/\SsixRepositoryCommit/pro_returns/S6_10_PAPERS_checkpoint_15_20261119T1045Z/S6_10_PAPERS/code/V/replay_chains/code/fourfolds_universal_56_obstruction.py}{\nolinkurl{fourfolds_universal_56_obstruction.py}} \\[3pt]
\hypertarget{prog:V:cap-tests}{}Finite-cap and cusp tests for third and fifth cohomology\par \cref{maps:src:F.2.29} &
\href{https://github.com/KokunoYumeto/s6-complex-structure-workbench/blob/\SsixRepositoryCommit/pro_returns/S6_10_PAPERS_checkpoint_15_20261119T1045Z/S6_10_PAPERS/code/V/replay_chains/code/fourfolds_local_caps.py}{\nolinkurl{fourfolds_local_caps.py}}\par \href{https://github.com/KokunoYumeto/s6-complex-structure-workbench/blob/\SsixRepositoryCommit/pro_returns/S6_10_PAPERS_checkpoint_15_20261119T1045Z/S6_10_PAPERS/code/V/replay_chains/code/fourfolds_cusp_boundary.py}{\nolinkurl{fourfolds_cusp_boundary.py}}\par \href{https://github.com/KokunoYumeto/s6-complex-structure-workbench/blob/\SsixRepositoryCommit/pro_returns/S6_10_PAPERS_checkpoint_15_20261119T1045Z/S6_10_PAPERS/code/V/replay_chains/code/fourfolds_survivor_test.py}{\nolinkurl{fourfolds_survivor_test.py}} \\[3pt]
\hypertarget{prog:V:fundamental-group}{}Marked fundamental-group presentation\par \cref{maps:src:F.2.31} &
\href{https://github.com/KokunoYumeto/s6-complex-structure-workbench/blob/\SsixRepositoryCommit/pro_returns/S6_10_PAPERS_checkpoint_15_20261119T1045Z/S6_10_PAPERS/code/V/replay_chains/code/fourfolds_vankampen.py}{\nolinkurl{fourfolds_vankampen.py}} \\[3pt]
\hypertarget{prog:V:three-points}{}Three-finite-point family examples\par \cref{maps:src:F.2.32} &
\href{https://github.com/KokunoYumeto/s6-complex-structure-workbench/blob/\SsixRepositoryCommit/pro_returns/S6_10_PAPERS_checkpoint_15_20261119T1045Z/S6_10_PAPERS/code/V/replay_chains/code/fourfolds_three_point_family.py}{\nolinkurl{fourfolds_three_point_family.py}} \\[3pt]
\hypertarget{prog:V:chains}{}Finite cap, complement and gluing chains\par \cref{maps:src:F.2.33,maps:src:F.2.34,maps:src:F.2.35} &
\href{https://github.com/KokunoYumeto/s6-complex-structure-workbench/blob/\SsixRepositoryCommit/pro_returns/S6_10_PAPERS_checkpoint_15_20261119T1045Z/S6_10_PAPERS/code/V/replay_chains/code/fourfolds_bar_koszul.py}{\nolinkurl{fourfolds_bar_koszul.py}}\par \href{https://github.com/KokunoYumeto/s6-complex-structure-workbench/blob/\SsixRepositoryCommit/pro_returns/S6_10_PAPERS_checkpoint_15_20261119T1045Z/S6_10_PAPERS/code/V/replay_chains/code/fourfolds_affine_cap_chains.py}{\nolinkurl{fourfolds_affine_cap_chains.py}}\par \href{https://github.com/KokunoYumeto/s6-complex-structure-workbench/blob/\SsixRepositoryCommit/pro_returns/S6_10_PAPERS_checkpoint_15_20261119T1045Z/S6_10_PAPERS/code/V/replay_chains/code/fourfolds_kuhn_cap_chains.py}{\nolinkurl{fourfolds_kuhn_cap_chains.py}}\par \href{https://github.com/KokunoYumeto/s6-complex-structure-workbench/blob/\SsixRepositoryCommit/pro_returns/S6_10_PAPERS_checkpoint_15_20261119T1045Z/S6_10_PAPERS/code/V/replay_chains/code/fourfolds_global_chains.py}{\nolinkurl{fourfolds_global_chains.py}}\par \href{https://github.com/KokunoYumeto/s6-complex-structure-workbench/blob/\SsixRepositoryCommit/pro_returns/S6_10_PAPERS_checkpoint_15_20261119T1045Z/S6_10_PAPERS/code/V/replay_chains/code/fourfolds_total_homology.py}{\nolinkurl{fourfolds_total_homology.py}}\par \href{https://github.com/KokunoYumeto/s6-complex-structure-workbench/blob/\SsixRepositoryCommit/pro_returns/S6_10_PAPERS_checkpoint_15_20261119T1045Z/S6_10_PAPERS/code/V/replay_chains/code/fourfolds_marked_gluing.py}{\nolinkurl{fourfolds_marked_gluing.py}} \\[3pt]
\hypertarget{prog:V:periods}{}Algebra of the period construction\par \cref{maps:src:F.3.5} &
\href{https://github.com/KokunoYumeto/s6-complex-structure-workbench/blob/\SsixRepositoryCommit/pro_returns/S6_10_PAPERS_checkpoint_15_20261119T1045Z/S6_10_PAPERS/code/V/replay_chains/code/fourfolds_periods.py}{\nolinkurl{fourfolds_periods.py}} \\[3pt]
\hypertarget{prog:V:tubes}{}Spanning tube directions\par \cref{maps:src:F.3.8} &
\href{https://github.com/KokunoYumeto/s6-complex-structure-workbench/blob/\SsixRepositoryCommit/pro_returns/S6_10_PAPERS_checkpoint_15_20261119T1045Z/S6_10_PAPERS/code/V/replay_chains/code/fourfolds_oka_tubes.py}{\nolinkurl{fourfolds_oka_tubes.py}} \\[3pt]
\hypertarget{prog:V:shifted-homology}{}Integral homology of shifted products\par \cref{maps:src:E.2.5,maps:src:E.2.6} &
\href{https://github.com/KokunoYumeto/s6-complex-structure-workbench/blob/\SsixRepositoryCommit/pro_returns/S6_10_PAPERS_checkpoint_15_20261119T1045Z/S6_10_PAPERS/code/V/replay_chains/code/fourfolds_shifted_integral.py}{\nolinkurl{fourfolds_shifted_integral.py}}\par \href{https://github.com/KokunoYumeto/s6-complex-structure-workbench/blob/\SsixRepositoryCommit/pro_returns/S6_10_PAPERS_checkpoint_15_20261119T1045Z/S6_10_PAPERS/code/V/replay_chains/code/manifolds_shifted_mv.py}{\nolinkurl{manifolds_shifted_mv.py}} \\[3pt]
\hypertarget{prog:V:equivariant}{}Bordism and low-degree boundary calculations\par \cref{maps:src:E.2.9,maps:src:E.2.11,maps:src:E.2.12} &
\href{https://github.com/KokunoYumeto/s6-complex-structure-workbench/blob/\SsixRepositoryCommit/pro_returns/S6_10_PAPERS_checkpoint_15_20261119T1045Z/S6_10_PAPERS/code/V/replay_chains/code/shifted_group_bordism.py}{\nolinkurl{shifted_group_bordism.py}}\par \href{https://github.com/KokunoYumeto/s6-complex-structure-workbench/blob/\SsixRepositoryCommit/pro_returns/S6_10_PAPERS_checkpoint_15_20261119T1045Z/S6_10_PAPERS/code/V/replay_chains/code/maps_shifted_pi2.py}{\nolinkurl{maps_shifted_pi2.py}}\par \href{https://github.com/KokunoYumeto/s6-complex-structure-workbench/blob/\SsixRepositoryCommit/pro_returns/S6_10_PAPERS_checkpoint_15_20261119T1045Z/S6_10_PAPERS/code/V/replay_chains/code/manifolds_cusp_alexander.py}{\nolinkurl{manifolds_cusp_alexander.py}} \\[3pt]
\hypertarget{prog:V:cyclic-cover}{}Circle fibration, cyclic cover and integral fibre\par \cref{maps:src:E.2.13,maps:src:E.2.14,maps:src:E.2.15} &
\href{https://github.com/KokunoYumeto/s6-complex-structure-workbench/blob/\SsixRepositoryCommit/pro_returns/S6_10_PAPERS_checkpoint_15_20261119T1045Z/S6_10_PAPERS/code/V/replay_chains/code/maps_shifted_circle.py}{\nolinkurl{maps_shifted_circle.py}}\par \href{https://github.com/KokunoYumeto/s6-complex-structure-workbench/blob/\SsixRepositoryCommit/pro_returns/S6_10_PAPERS_checkpoint_15_20261119T1045Z/S6_10_PAPERS/code/V/replay_chains/code/maps_shifted_return.py}{\nolinkurl{maps_shifted_return.py}}\par \href{https://github.com/KokunoYumeto/s6-complex-structure-workbench/blob/\SsixRepositoryCommit/pro_returns/S6_10_PAPERS_checkpoint_15_20261119T1045Z/S6_10_PAPERS/code/V/replay_chains/code/maps_shifted_cover.py}{\nolinkurl{maps_shifted_cover.py}}\par \href{https://github.com/KokunoYumeto/s6-complex-structure-workbench/blob/\SsixRepositoryCommit/pro_returns/S6_10_PAPERS_checkpoint_15_20261119T1045Z/S6_10_PAPERS/code/V/replay_chains/code/maps_shifted_cover_integral.py}{\nolinkurl{maps_shifted_cover_integral.py}}\par \href{https://github.com/KokunoYumeto/s6-complex-structure-workbench/blob/\SsixRepositoryCommit/pro_returns/S6_10_PAPERS_checkpoint_15_20261119T1045Z/S6_10_PAPERS/code/V/replay_chains/code/maps_shifted_fibre_integral.py}{\nolinkurl{maps_shifted_fibre_integral.py}} \\[3pt]
\hypertarget{prog:V:quadratic}{}Complement and splice modules and their quadratic maps\par \cref{maps:src:E.2.16,maps:src:E.2.17,maps:src:E.2.18} &
\href{https://github.com/KokunoYumeto/s6-complex-structure-workbench/blob/\SsixRepositoryCommit/pro_returns/S6_10_PAPERS_checkpoint_15_20261119T1045Z/S6_10_PAPERS/code/V/replay_chains/code/maps_complement_module.py}{\nolinkurl{maps_complement_module.py}}\par \href{https://github.com/KokunoYumeto/s6-complex-structure-workbench/blob/\SsixRepositoryCommit/pro_returns/S6_10_PAPERS_checkpoint_15_20261119T1045Z/S6_10_PAPERS/code/III/maps_splice_module.py}{\nolinkurl{maps_splice_module.py}}\par \href{https://github.com/KokunoYumeto/s6-complex-structure-workbench/blob/\SsixRepositoryCommit/pro_returns/S6_10_PAPERS_checkpoint_15_20261119T1045Z/S6_10_PAPERS/code/III/maps_cusp_quadratic.py}{\nolinkurl{maps_cusp_quadratic.py}}\par \href{https://github.com/KokunoYumeto/s6-complex-structure-workbench/blob/\SsixRepositoryCommit/pro_returns/S6_10_PAPERS_checkpoint_15_20261119T1045Z/S6_10_PAPERS/code/III/maps_splice_quadratic.py}{\nolinkurl{maps_splice_quadratic.py}} \\[3pt]
\hypertarget{prog:V:framings}{}Plumbing, finite-cap framings and handle inventory\par \cref{maps:src:E.2.19,maps:src:E.2.20,maps:src:E.2.21} &
\href{https://github.com/KokunoYumeto/s6-complex-structure-workbench/blob/\SsixRepositoryCommit/pro_returns/S6_10_PAPERS_checkpoint_15_20261119T1045Z/S6_10_PAPERS/code/III/maps_cusp_smooth_plumbing.py}{\nolinkurl{maps_cusp_smooth_plumbing.py}}\par \href{https://github.com/KokunoYumeto/s6-complex-structure-workbench/blob/\SsixRepositoryCommit/pro_returns/S6_10_PAPERS_checkpoint_15_20261119T1045Z/S6_10_PAPERS/code/III/maps_finite_cap_framings.py}{\nolinkurl{maps_finite_cap_framings.py}}\par \href{https://github.com/KokunoYumeto/s6-complex-structure-workbench/blob/\SsixRepositoryCommit/pro_returns/S6_10_PAPERS_checkpoint_15_20261119T1045Z/S6_10_PAPERS/code/V/replay_chains/code/maps_ordinary_handles.py}{\nolinkurl{maps_ordinary_handles.py}} \\[3pt]
\end{longtable}
\endgroup

\clearpage
\begin{thebibliography}{99}
\providecommand{\natexlab}[1]{#1}
\providecommand{\doi}[1]{\url{https://doi.org/#1}}
\bibitem[Achinger and Ogus(2020)]{ffAchingerOgus20}
Piotr Achinger and Arthur Ogus.

\newblock Monodromy and log geometry.

\newblock \emph{Tunisian Journal of Mathematics}, 2\penalty0 (3):\penalty0
  455--534, 2020.

\newblock \doi{10.2140/tunis.2020.2.455}.

\newblock URL \url{https://msp.org/tunis/2020/2-3/tunis-v2-n3-p01-p.pdf}.

\bibitem[Alarc{\'o}n and Forstneri{\v c}(2025)]{IVVOkaOneV5}
Antonio Alarc{\'o}n and Franc Forstneri{\v c}.

\newblock {Oka}-1 manifolds.

\newblock \emph{Mathematische Zeitschrift}, 311, article 33, 2025.

\newblock \doi{10.1007/s00209-025-03833-4}.

\newblock Statement and page numbers refer to
\href{https://arxiv.org/pdf/2303.15855v5}{arXiv:2303.15855v5}.

\bibitem[Beukers and Heckman(1989)]{BeukersHeckman1989}
Frits Beukers and Gert Heckman.

\newblock Monodromy for the hypergeometric function {$ {}_nF_{n-1}$}.

\newblock \emph{Inventiones mathematicae}, 95:\penalty0 325--354, 1989.

\newblock \doi{10.1007/BF01393900}.

\bibitem[Bierstone and Milman(1997)]{BM97}
Edward Bierstone and Pierre~D. Milman.

\newblock Canonical desingularization in characteristic zero by blowing up the
  maximum strata of a local invariant.

\newblock \emph{Inventiones Mathematicae}, 128:\penalty0 207--302, 1997.

\newblock \doi{10.1007/s002220050141}.

\newblock Page numbers cited are those of arXiv:alg-geom/9508005.

\bibitem[Bley et~al.(2022)Bley, Hofmann, and Johnston]{rootBHJ22}
Werner Bley, Tommy Hofmann, and Henri Johnston.

\newblock Computation of lattice isomorphisms and the integral matrix
  similarity problem.

\newblock \emph{Forum of Mathematics, Sigma}, 10:\penalty0 e87, 2022.

\newblock \doi{10.1017/fms.2022.74}.

\bibitem[Borcherds et al.(2016)]{pubLieGroups2016}
R.~Borcherds, M.~Haiman, N.~Reshetikhin, V.~Serganova, and T.~Johnson-Freyd.

\newblock \emph{Part I: Lie Groups}.

\newblock Lecture notes, version of 5 October 2016.

\newblock URL \url{https://math.berkeley.edu/~serganov/261/part1.pdf}.

\bibitem[Cartan(1952)]{Cartan1952B_AB}
Henri Cartan.

\newblock Faisceaux analytiques sur les vari{\'e}t{\'e}s de {Stein}:
  d{\'e}monstration des th{\'e}or{\`e}mes fondamentaux.

\newblock \emph{S{\'e}minaire Henri Cartan}, 4:\penalty0 1--15, 1952.

\newblock URL \url{https://www.numdam.org/article/SHC_1951-1952__4__A19_0.pdf}.

\bibitem[Catanese(2004)]{Catanese2004}
Fabrizio Catanese.

\newblock Deformation in the large of some complex manifolds, {I}.

\newblock \emph{Annali di Matematica Pura ed Applicata}, 183:\penalty0
  261--289, 2004.

\newblock \doi{10.1007/s10231-003-0096-y}.

\newblock URL
  \url{https://www.mathe8.uni-bayreuth.de/de/team/prof-fabrizio-catanese-old/pdf/085.pdf}.

\bibitem[{Claude} and {L. Alp{\"o}ge}(2026)]{AC26}
{Claude} and {L. Alp{\"o}ge}.

\newblock The $(3,4,\infty)$ modular family of 2-tori, completed at its three
  special points, is a complex structure on $S^6$.

\newblock Manuscript, 108 pages, 2026.

\newblock URL \url{https://alpo.ge/s6.pdf}.

\bibitem[Conway et~al.(2023)Conway, Crowley, Powell, and Sixt]{IVVWallForm}
Anthony Conway, Diarmuid Crowley, Mark Powell, and J{\"o}rg Sixt.

\newblock Simply connected manifolds with large homotopy stable classes.

\newblock \emph{Journal of the Australian Mathematical Society}, 115\penalty0 (2):\penalty0 172--203, 2023.

\newblock \doi{10.1017/S1446788722000167}.

\newblock The page numbers cited are those of the authors' version v3 (2021), \url{https://www.maths.gla.ac.uk/~mpowell/large-homotopy-stable-class-v3.pdf}.

\bibitem[Cox et~al.(2011)Cox, Little, and Schenck]{CoxLittleSchenck2011}
David~A. Cox, John~B. Little, and Henry~K. Schenck.

\newblock \emph{Toric Varieties}, volume 124 of \emph{Graduate Studies in
  Mathematics}.

\newblock American Mathematical Society, Providence, RI, 2011.

\newblock \doi{10.1090/gsm/124}.

\bibitem[Davis and Kirk(2001)]{rootDavisKirk01}
James~F. Davis and Paul Kirk.

\newblock \emph{Lecture Notes in Algebraic Topology}, volume~35 of
  \emph{Graduate Studies in Mathematics}.

\newblock American Mathematical Society, 2001.

\newblock \doi{10.1090/gsm/035}.

\bibitem[Deligne(1971)]{ffDeligne}
Pierre Deligne.

\newblock Th{\'e}orie de Hodge, {II}.

\newblock \emph{Publications Math{\'e}matiques de l'IH{\'E}S}, 40:\penalty0
  5--57, 1971.

\newblock \doi{10.1007/BF02684692}.

\newblock URL
  \url{https://publications.ias.edu/sites/default/files/Number15.pdf}.

\bibitem[Douady(1974)]{Douady1974}
Adrien Douady.

\newblock Le probl\`eme des modules locaux pour les espaces $\mathbf{C}$-analytiques
  compacts.

\newblock \emph{Annales scientifiques de l'\'Ecole Normale Sup\'erieure},
  7\penalty0 (4):\penalty0 569--602, 1974.

\newblock \doi{10.24033/asens.1280}.

\newblock URL \url{https://numdam.org/item/10.24033/asens.1280.pdf}.

\bibitem[Du et~al.(2026)Du, Guo, Kusakabe, and Xie]{DGKX2026v1}
Yun-Heng Du, Bin Guo, Yuta Kusakabe, and Song-Yan Xie.

\newblock Oka geometry and deformations of torus-fibred complex six-spheres,
  2026.

\newblock URL \url{https://arxiv.org/abs/2610.00250v1}.

\bibitem[Engel(2026)]{E26}
Philip Engel.

\newblock Complex structures on $S^6$.

\newblock arXiv:2609.38442v1, 2026.

\newblock URL \url{https://arxiv.org/abs/2609.38442v1}.

\bibitem[Esole and Kang(2018)]{ffEsoleKang}
Mboyo Esole and Monica~Jinwoo Kang.

\newblock Characteristic numbers of crepant resolutions of Weierstrass models,
  2018.

\newblock URL \url{https://arxiv.org/pdf/1807.08755}.

\bibitem[Forstneri{\v c} and Prezelj(2002)]{ffForstnericPrezelj02}
Franc Forstneri{\v c} and Jasna Prezelj.

\newblock Oka's principle for holomorphic submersions with sprays.

\newblock \emph{Mathematische Annalen}, 322\penalty0 (4):\penalty0 633--666,
  2002.

\newblock \doi{10.1007/s002080100249}.

\newblock URL \url{https://arxiv.org/pdf/math/0101040}.

\bibitem[Grauert(1960)]{Grauert60}
Hans Grauert.

\newblock Ein Theorem der analytischen Garbentheorie und die Modulräume komplexer Strukturen.

\newblock \emph{Publications Mathématiques de l'IHÉS}, 5:\penalty0 5--64,
  1960.

\newblock \doi{10.1007/BF02684746}.

\newblock URL \url{https://www.numdam.org/article/PMIHES_1960__5__5_0.pdf}.

\bibitem[Hatcher(2002)]{Hatcher2002}
Allen Hatcher.

\newblock \emph{Algebraic Topology}.

\newblock Cambridge University Press, 2002.

\newblock URL \url{https://pi.math.cornell.edu/~hatcher/AT/AT.pdf}.

\bibitem[Hirsch(1976)]{rootHirsch76}
Morris~W. Hirsch.

\newblock \emph{Differential Topology}, volume~33 of \emph{Graduate Texts in
  Mathematics}.

\newblock Springer, 1976.

\newblock \doi{10.1007/978-1-4684-9449-5}.

\bibitem[Hirzebruch(1978)]{HirzebruchTopologicalMethods}
Friedrich Hirzebruch.

\newblock \emph{Topological Methods in Algebraic Geometry}.

\newblock Springer-Verlag, third edition, second corrected printing, 1978; reprinted in Classics in Mathematics, 1995.

\newblock URL
  \url{https://hirzebruch.mpim-bonn.mpg.de/id/eprint/114/6/Hirzebruch-TopologicalMethodsInAlgebraicGeometry.pdf}.

\bibitem[Isaksen et~al.(2023)Isaksen, Wang, and Xu]{IsaksenWangXu2023}
Daniel~C. Isaksen, Guozhen Wang, and Zhouli Xu.

\newblock Stable homotopy groups of spheres: from dimension 0 to 90.

\newblock \emph{Publications mathématiques de l'IHÉS}, 137:\penalty0
  107--243, 2023.

\newblock \doi{10.1007/s10240-023-00139-1}.

\bibitem[Jia and Meng(2024)]{ffJiaMeng}
Jia Jia and Sheng Meng.

\newblock Equivariant K{\"a}hler model for Fujiki's class.

\newblock \emph{The Journal of Geometric Analysis}, 34, article 349, 2024.

\newblock \doi{10.1007/s12220-024-01796-2}.

\newblock Statement and page numbers refer to
\href{https://arxiv.org/pdf/2201.06748v1}{arXiv:2201.06748v1}.

\bibitem[Kramer and Stolz(2007)]{IVVKramerStolz}
Linus Kramer and Stephan Stolz.

\newblock A diffeomorphism classification of manifolds which are like
  projective planes.

\newblock \emph{Journal of Differential Geometry}, 77:\penalty0 177--188, 2007.

\newblock \doi{10.4310/jdg/1191860392}.

\bibitem[Mandal(2026)]{ffMandal}
Souvik Mandal.

\newblock Complex and stable complex structures on real vector bundles over
  connected sums of quaternionic projective planes, 2026.

\newblock URL \url{https://arxiv.org/abs/2609.27929}.

\bibitem[Milnor(1965)]{mapsMilnor65}
John Milnor.

\newblock \emph{Lectures on the h-Cobordism Theorem}.

\newblock Princeton University Press, Princeton, New Jersey,
  1965.

\bibitem[{National Institute of Standards and
  Technology}(2026)]{DLMFCauchy2026}
{National Institute of Standards and Technology}.

\newblock {Digital Library of Mathematical Functions}: Cauchy's integral and
  derivative formulas, 2026.

\newblock URL \url{https://dlmf.nist.gov/1.9}.

\bibitem[Peters and Sterk(2024)]{ffPetersSterk}
Chris Peters and Hans Sterk.

\newblock Symmetric and quadratic forms, with applications to coding theory,
  algebraic geometry and topology.

\newblock Authors\textquotesingle{} June 2024 manuscript, 2024.

\newblock URL
  \url{https://www-fourier.univ-grenoble-alpes.fr/~peters/Books/QuadraticForms/QuadForms.pdf}.

\bibitem[Scott(1983)]{rootScott83}
Peter Scott.

\newblock The geometries of 3-manifolds.

\newblock \emph{Bulletin of the London Mathematical Society}, 15\penalty0
  (5):\penalty0 401--487, 1983.

\newblock \doi{10.1112/blms/15.5.401}.

\bibitem[Serre(1955)]{Serre1955Duality}
Jean-Pierre Serre.

\newblock Un th{\'e}or{\`e}me de dualit{\'e}.

\newblock \emph{Commentarii Mathematici Helvetici}, 29:\penalty0 9--26, 1955.

\newblock URL
  \url{https://www.college-de-france.fr/media/jean-pierre-serre/UPL35645220559164287_Un_theoreme_de_dualite.pdf}.

\bibitem[Teschl(2012)]{TeschlODE2012}
Gerald Teschl.

\newblock \emph{Ordinary Differential Equations and Dynamical Systems}, volume
  140 of \emph{Graduate Studies in Mathematics}.

\newblock American Mathematical Society, Providence, 2012.

\newblock URL \url{https://www.mat.univie.ac.at/~gerald/ftp/book-ode/ode.pdf}.

\bibitem[{The Stacks Project Authors}(2026)]{StacksProject}
{The Stacks Project Authors}.

\newblock The Stacks project, 2026.

\newblock URL \url{https://stacks.math.columbia.edu}.

\bibitem[Vogel(2009)]{mapsVogel09}
Thomas Vogel.

\newblock Existence of Engel structures.

\newblock \emph{Annals of Mathematics}, 169\penalty0 (1):\penalty0 79--137,
  2009.

\newblock \doi{10.4007/annals.2009.169.79}.

\newblock URL
  \url{https://annals.math.princeton.edu/wp-content/uploads/annals-v169-n1-p03.pdf}.

\bibitem[Whitehead(1950)]{mapsWhitehead50}
J.~H.~C. Whitehead.

\newblock The secondary boundary operator.

\newblock \emph{Proceedings of the National Academy of Sciences of the United
  States of America}, 36\penalty0 (1):\penalty0 55--60, 1950.

\newblock \doi{10.1073/pnas.36.1.55}.

\bibitem[Winter(2009)]{mapsWinter08}
Blake Winter.

\newblock The classification of spun torus knots.

\newblock \emph{Journal of Knot Theory and Its Ramifications},
18(9):1287--1298, 2009.

\newblock \doi{10.1142/S0218216509007476}.

\newblock Statement and page numbers refer to
\href{https://arxiv.org/abs/0711.1638v3}{arXiv:0711.1638v3}.

\bibitem[Yang(2026)]{Yang2026}
Xiufan Yang.

\newblock Compact complex threefolds fibred by two-dimensional complex tori.

\newblock arXiv:2609.29666v1, 2026.

\newblock URL \url{https://arxiv.org/abs/2609.29666v1}.

\end{thebibliography}

\end{document}
